% Options for packages loaded elsewhere
\PassOptionsToPackage{unicode}{hyperref}
\PassOptionsToPackage{hyphens}{url}
%
\documentclass{book}
\usepackage{geometry}
% \geometry{legalpaper, landscape, margin=2in} 
\usepackage{amsmath,amssymb}
\usepackage{amsthm} 
\usepackage{mathtools}
\usepackage{lmodern}
\usepackage{iftex}
\usepackage{blindtext}
\usepackage{comment}
\ifPDFTeX
  \usepackage[T1]{fontenc} 
  \usepackage[utf8]{inputenc}
  \usepackage{textcomp} % provide euro and other symbols
  %\usepackage{pgfplots}
  \usepackage[shortlabels]{enumitem}
\else % if luatex or xetex
  \usepackage{unicode-math}
  \defaultfontfeatures{Scale=MatchLowercase}
  \defaultfontfeatures[\rmfamily]{Ligatures=TeX,Scale=1}
\fi
% Use upquote if available, for straight quotes in verbatim environments
\IfFileExists{upquote.sty}{\usepackage{upquote}}{}
\IfFileExists{microtype.sty}{% use microtype if available
  \usepackage[]{microtype}
  \UseMicrotypeSet[protrusion]{basicmath} % disable protrusion for tt fonts
}{}
\makeatletter

\makeatother
\usepackage{xcolor}
\usepackage{longtable,booktabs,array}
\usepackage{multirow}
\usepackage{calc} % for calculating minipage widths
% Correct order of tables after \paragraph or \subparagraph
\usepackage{etoolbox}
\makeatletter
\patchcmd\longtable{\par}{\if@noskipsec\mbox{}\fi\par}{}{}
\makeatother
% Allow footnotes in longtable head/foot
\IfFileExists{footnotehyper.sty}{\usepackage{footnotehyper}}{\usepackage{footnote}}
\makesavenoteenv{longtable}
\usepackage{graphicx}
\makeatletter
\def\maxwidth{\ifdim\Gin@nat@width>\linewidth\linewidth\else\Gin@nat@width\fi}
\def\maxheight{\ifdim\Gin@nat@height>\textheight\textheight\else\Gin@nat@height\fi}
\makeatother
% Scale images if necessary, so that they will not overflow the page
% margins by default, and it is still possible to overwrite the defaults
% using explicit options in \includegraphics[width, height, ...]{}
\setkeys{Gin}{width=\maxwidth,height=\maxheight,keepaspectratio}
% Set default figure placement to htbp
\makeatletter
\def\fps@figure{htbp}
\makeatother
\usepackage[normalem]{ulem}
\setlength{\emergencystretch}{3em} % prevent overfull lines
\providecommand{\tightlist}{%
  \setlength{\itemsep}{0pt}\setlength{\parskip}{0pt}}
\ifLuaTeX
\usepackage{selnolig}  % disable illegal ligatures
\fi
\IfFileExists{bookmark.sty}{\usepackage{bookmark}}{\usepackage{hyperref}}
\IfFileExists{xurl.sty}{\usepackage{xurl}}{} % add URL line breaks if available
\urlstyle{same} % disable monospaced font for URLs
\hypersetup{
  hidelinks,
  pdfcreator={LaTeX via pandoc}}

\usepackage{cleveref}
\usepackage{tikz-cd}
\usepackage{enumitem}

\newtheorem{definition}{Definition}[section]
\newtheorem{lemma}[definition]{Lemma}
\newtheorem{exercise}{Exercise}[section]
\newtheorem{example}[definition]{Example}
\newtheorem{proposition}[definition]{Proposition}
\newtheorem{remark}[definition]{Remark}
\newtheorem{theorem}[definition]{Theorem}
\newtheorem{corollary}[definition]{Corollary}

\def\cstar{\(C^*\)-algebra}
\def\cstars{\(C^*\)-algebras}
\def\diag{\textnormal{diag}}
\def\Index{\textnormal{Index}}

\title{\(C^*\)-algebras}
\author{}
\date{}

\begin{document}
\maketitle

\tableofcontents
\part{\texorpdfstring{Foundations of \(C^*\)-Algebras and \(K\)-Theory}
                     {Foundations of C*-Algebras and K-Theory}}
\chapter{Basics of \texorpdfstring{\(C^*\)-algebras}{C*algebras}}
\section{Definition}
\begin{definition}
  A \textnormal{\(C^*\)-algebra} \(A\) is a complex Banach *-algebra satisfying
  \(\|a^*a\| = \|a\|^2\) for every \(a\in A\).
\end{definition}
\begin{proposition}
  If \(a\) is an element of a \(C^*\)-algebra \(A\), then \(\|a\|=\|a^*\|\).
\end{proposition}
\begin{proof}
  \(\|a\|^2=\|a^*a\|\leq\|a^*\|\|a\|\) so \(\|a\|\leq\|a^*\|\).\\
  \(\|a^*\|^2=\|aa^*\|\leq\|a\|\|a^*\|\) so \(\|a^*\|\leq\|a\|\).
\end{proof}
\begin{definition}
  Let \(A\) be a \(C^*\)-algebra.  
  \begin{enumerate}
    \item An element \(a\) is \textnormal{self-adjoint} if \(a^*=a\). We denote the subset of self-adjoint elements of \(A\) by \(A_{sa}\). 
    \item An element \(a\) is \textnormal{normal} if \(a^*a = aa^*\).
    \item An element \(p\) is a \textnormal{projection} if \(p^2 = p = p^*\); equivalently, \(p\) is a self-adjoint idempotent.
    \item Assuming that \(A\) is unital, an element \(u\) is a \textnormal{unitary} if \(u^*u = 1 = uu^*\); equivalently, \(u\) is invertible and \(u^{-1} = u^*\).
    \item Assuming that \(A\) is unital, an element \(u\) is an \textnormal{isometry} if \(u^*u = 1\).
    \item An element \(u\) is a \textnormal{partial isometry} if \(u^*u\) is a projection.
    \item An element \(a\) is \textnormal{positive} if it may be written \(a = b^*b\), for some \(b \in A\). In this case, we often write \(a \geq 0\) for brevity.
  \end{enumerate}
\end{definition}

\section{Examples}
\begin{example}\label{example:125}
  Suppose that \(A\) and \(B\) are \cstars, we form their direct
  sum
  \[A\oplus B=\{(a,b) | a \in A,b \in B\}.\]
  The algebraic operations are defined coordinate-wise and the norm is
  given by
  \[\|(a, b)\| = \max{\|a\| , \|b\| },\]
  for any \(a\) in \(A\) and \(b\) in \(B\).
  There is an obvious extension of this notion to finite direct sums. Also,
  if \(A_n,n \geq 1\) is a sequence of \cstars, their direct sum is defined as
  \[
  \bigoplus^\infty_{n=1} A_n = \{(a_1,a_2,\ldots) | a_n \in A_n\text{ for all }n, \lim_n\|a_n\| =0\},
  \]
  and we define a norm as
  \[
    \|(a_n)\|=\sup_n\|a_n\|.
  \]
\end{example}
\begin{example}
The algebra \(\mathbb C^2\), equipped with coordinatewise operations and the norm
\[
\|(z,w)\|=\max\{|z|,|w|\},
\]
is a unital \cstar.
\end{example}

\section{Spectrum}
\begin{lemma}
  \begin{enumerate}
    \item If \(a\) is an element of a unital Banach algebra \(A\) and \(\|a-1\|<1\), then \(a\) is invertible.
    \item The set of invertible elements of \(A\) is open.
  \end{enumerate}
\end{lemma}
\begin{proof}
  Let \(b=\sum_{n=0}^{\infty}{(1-a)}^n\). The series is Cauchy (proof in the exercise) thus it converges to some element.
  \[
    ab=a\lim_{N\rightarrow\infty}\sum_{n=0}^{N}{(1-a)}^n=(1-(1-a))\lim_{N\rightarrow\infty}\sum_{n=0}^{N}{(1-a)}^n
  \]
  Applying the finite geometric-series identity, for any \(N\),
  \[
    (1-r)\sum_{n=0}^{N}r^n=1-r^{N+1}
  \]
  substituting \(1-a=r\) we obtain
  \[
    ab=\lim_{N\rightarrow\infty}1-{(1-a)}^{N+1}=1.
  \]
  A similar argument shows \(ba = 1\) and so \(a\) is invertible.\par
  If \(a\) is invertible, the map \(b \mapsto ab\) is a homeomorphism of \(A\), which
  preserves the set of invertibles. It also maps the unit to \(a\), so the conclusion
  follows from the first part.
\end{proof}
\begin{definition}
  Let \(A\) be a unital algebra and let \(a\) be an element of \(A\). The \textnormal{spectrum} of \(a\) is 
  \[
    \operatorname{spec}(a) = \{\lambda\in\mathbb{C}\mid\lambda 1-a\text{ is not invertible}\}
  \]
  The \textnormal{spectral radius} of \(a\) is
  \[
    r(a)=\sup\{\left\lvert \lambda\right\rvert \mid\lambda\in \operatorname{spec}(a)\}
  \]
  which is defined provided the spectrum is non-empty, allowing the possibility of \(r(a)=\infty\).
\end{definition}
\begin{remark}
  Let us consider for the moment the \(C^*\)-algebra \(M_n(\mathbb{C})\), where
  \(n\) is some fixed positive integer. Recall the following classical result from linear algebra,
  that the following four conditions on an element \(a\) of \(M_n(\mathbb{C})\) are equivalent:
  \begin{enumerate}
    \item \(a\) is invertible,
    \item \(a: \mathbb{C}^n\rightarrow\mathbb{C}^n\) is injective,
    \item \(a: \mathbb{C}^n\rightarrow\mathbb{C}^n\) is surjective,
    \item \(\det(a)\neq 0\).
  \end{enumerate}
  This fact allows us to compute the spectrum of the element \(a\) simply by finding
  the zeros of \(\det(\lambda-a)\), which (conveniently) is a polynomial of degree \(n\).\\
  The remark we make now is that in the \(C^*\)-algebra \(\mathcal{B}(\mathcal{H})\), such a result fails
  in general when \(\mathcal{H}\) is not finite-dimensional. First of all, the determinant
  function simply fails to exist and while the first condition implies the next
  two, any other implication between the first three doesn't hold.
\end{remark}
\begin{theorem}\label{theorem:134}
  Let \(A\) be a unital Banach algebra. The spectrum of any element is non-empty and compact.
\end{theorem}
\begin{theorem}\label{theorem:135}
  Let \(a\) be an element of a unital Banach algebra \(A\). The sequence \(\|a^n\|^{1/n}\)
  is bounded by \(\|a\|\) and has limit \(r(a)\). In particular,
  \(r(a)\) is finite and \(r(a) \leq \|a\|\).
\end{theorem}
\begin{proof}
  See \cite[Theorem 1.2.7]{MGJ90}.
\end{proof}
\begin{theorem}\label{theorem:136}
   Let \(a\) be a self-adjoint element of a unital \cstar\ \(A\). Then \(\|a\| =r(a)\).
\end{theorem}
\begin{proof}
  As \(a\) is self-adjoint, we have that \(\|a^2\|=\|a^*a\|=\|a\|^2\). 
  By induction, we have that for every integer \(k\), \(\|a\|^{2^k}=\|a^{2^k}\|\). 
  Since the limit exists by Theorem \ref{theorem:135}, it agrees with the limit of every subsequence.
  \[
    r(a)=\lim_{n\rightarrow\infty}\|a^n\|^{1/n}=\lim_{k\rightarrow\infty}\|a^{2^k}\|^{2^{-k}}=\|a\|
  \]
\end{proof}
\begin{corollary}\label{cor:137}
  Let \(A\) and \(B\) be \(C^*\)-algebras and suppose that \(\rho: A \rightarrow B\)
  is a \(*\)-homomorphism. Then \(\rho\) is contractive; that is, \(\|\rho(a)\| \leq \|a\|\), for all
  \(a\) in \(A\). In particular, we have \(\|\rho\|\leq 1\).
\end{corollary}
\begin{corollary}\label{corollary:138}
  If \(A\) is a \cstar, then its norm is unique. That is, if
  a *-algebra possesses a norm in which it is a \cstar, then it possesses only one such norm.
\end{corollary}
\begin{exercise}
Let \(a\in A\) such that \(\|a-1\|<1\). Consider the sequence
\[
  s_N=\sum_{n=1}^{N}{(1-a)}^n
\]
and prove that it is a Cauchy sequence.
\begin{align*}
  \|s_N-s_M\|&=\Vert\sum_{n=1}^{N}{(1-a)}^n-\sum_{n=1}^{M}{(1-a)}^n\Vert \\
  &=\|\sum_{n=M+1}^{N}{(1-a)}^n\|\\
  &\leq\sum_{n=M+1}^{N}\|(1-a)\|^n\\
  &\leq\frac{\|1-a\|^{M+1}}{1-\|1-a\|}\rightarrow 0\mbox{ as } M\rightarrow \infty 
\end{align*}
\end{exercise}
\begin{exercise}
For any real number \(t\geq0\), consider the following element of \(M_2(\mathbb{C})\):
\[
a =
\left[\begin{pmatrix}
  1&t \\ 0&1
\end{pmatrix}\right] 
\]
Find the spectrum of \(a\), the spectral radius of \(a\) and the norm of \(a\).
\[
  \operatorname{spec}(a)=\{\lambda\in\mathbb{C}| \lambda1-a \mbox{ is not invertible}\}
\] 
\[
  M=\lambda1-a = 
  \left[\begin{pmatrix}
    \lambda-1&-t \\ 0&\lambda-1
  \end{pmatrix}\right] 
\]
\begin{align*} 
  \det(M)&={(\lambda-1)}^2-0
\end{align*}
\[\operatorname{spec}(a)=\{1\}\]
\[r(a)=1\]
\begin{align*}
  \|a\|&=\sqrt{\sup(|\lambda||\lambda1-a^*a\mbox{ is not invertible})}\\
  &=\sqrt{\sup(|\lambda||\lambda1-a^*a\mbox{ is not invertible})}
\end{align*}
\begin{align*}
  \lambda1-a^*a&=
  \begin{pmatrix}
    \lambda&0\\0&\lambda
  \end{pmatrix}-
  \begin{pmatrix}
    1&0\\t&1
  \end{pmatrix}
  \begin{pmatrix}
  1&t \\ 0&1
\end{pmatrix}\\
&=
  \begin{pmatrix}
    \lambda-1&-t\\-t&\lambda-t^2-1
  \end{pmatrix}
\end{align*}
\begin{align*}
  \det(\lambda1-a^*a)&=(\lambda-1)(\lambda-t^2-1)-t^2\\
  &=\lambda^2-\lambda(t^2+2)+1
\end{align*}
\begin{align*}
  \lambda_{1,2}=\frac{t^2+2\pm \sqrt{(t^2+2)^2-4}}{2}
\end{align*}
\begin{align*}
  \|a\|&=\sqrt{\frac{t^2+2+ t\sqrt{t^2+4}}{2}}\\
  &=\frac{t+\sqrt{t^2+4}}{2},
\end{align*}
where the last equality holds observing that:
\[
  \left(\frac{t+\sqrt{t^2+4}}{2}\right)^2=\frac{t^2+2+t\sqrt{t^2+4}}{2}.
\]
\end{exercise}

\section{Commutative \texorpdfstring{\(C^*\)}{C*}-algebras}
\begin{definition}
  Let \(A\) be an algebra over \(\mathbb{C}\). We define \(\mathcal{M}(A)\) to be the set of non-zero homomorphisms to \(\mathbb{C}\).
\end{definition}
\begin{lemma}\label{lemma:142}
  Let \(A\) be a unital commutative \(C^*\)-algebra and let \(\phi \) be in \(\mathcal{M} (A)\).
\begin{enumerate}
  \item \(\forall a\in A\), \(\phi(a)\in \operatorname{spec}(a)\).
  \item \(\phi\) is bounded and \(\|\phi\|=1\).
  \item \(\forall a\in A\), \(\phi(a^*)=\overline{\phi(a)}\).
\end{enumerate}
\end{lemma}
\begin{proof}
  \(\ker(\phi)=\{\alpha\in A|\phi(\alpha)=0\}\) is an ideal:\\
  \(\beta\in A\), \(\alpha\in\ker(\phi)\), \(\beta\alpha\in\ker(\phi)\Leftrightarrow \phi(\alpha\beta)=0\Leftrightarrow\phi(\alpha\beta)=\phi(\alpha)\phi(\beta)=0\).\\
  Ideals do not contain invertible elements.\\
  \(\phi(\phi(a)1_A-a)=\phi(a)\phi(1)-\phi(a)=0\), so \(\phi(a)1-a\in\ker(\phi)\). So it is not invertible.\\
  So \(\phi(a)\in \operatorname{spec}(a)\) and \(|\phi(a)|\leq r(a)\leq\|a\|\), so \(\phi\) is bounded by \(1\). (Norm defined as \(\|\phi\|=\sup\limits_{\|a\|\leq1}(|\phi(a)|)\)).\\
  \(\phi(1)=1\) so \(\|\phi\|\geq1\) and we conclude that \(\|\phi\|=1\).\par
  Let \(a=b+ic\), where \(b=(a+a^*)/2\) and \(c=(ia^*-ia)/2\).\\
  \(b^*=(a^*+a)/2=b\) and \(c^*=(-ia+ia^*)/2=c\) so \(b\) and \(c\) are self-adjoint.\\
  Now we prove that \(\phi(b)\in\mathbb{R}\), \(\forall b\in A\) self-adjoint.\\
  Let \(u_t=e^{ibt}\).\ \(u_t^*=e^{-ibt}=u_{-t}\). \(\|u_t u_t^*\|=\|e^{ibt}e^{-ibt}\|=\|1\|=1\). \(\|u_t\|=\|u_t^*u_t\|^{1/2}=1\).\\
  \(\phi(u_t)=\phi(e^{ibt})=e^{it\phi(b)}\).\\
  \(|e^{it\phi(b)}|=|\phi(u_t)|\leq\|\phi\|\|u_t\|=1\).\\
  The inequality \(|e^{it\phi(b)}|\leq1\) \(\forall t\in\mathbb{R}\) means that \(\phi(b)\in\mathbb{R}\) because, letting \(\phi(b)=x+iy\), we have
  \(|e^{itx}e^{-ty}|\leq1\). So we have \(e^{-ty}\leq1\).
  \begin{itemize}
    \item If \(y>0\) we choose \(t<0\) and \(e^{-ty}\leq1\) is false.
    \item If \(y<0\) we choose \(t>0\) and \(e^{-ty}\leq1\) is false.
    \item We conclude that \(y=0\) and \(\phi(b)=x\in\mathbb{R}\).
  \end{itemize}
  Now we have \({\phi(a^*)}={\phi{(b+ic)}^*}=\phi(b^*)-i\phi(c^*)=\overline{\phi(a)}\).

\end{proof}
\begin{lemma}
  Let \(A\) be a unital commutative \(C^*\)-algebra. The set \(\mathcal{M}(A)\) is a weak-\(*\) compact subset of the unit ball of the dual space \(A^*\).
\end{lemma}
\begin{proof}  
  \(\forall\phi\in\mathcal{M}(A)\) \(\|\phi\|=1\) so \(\mathcal{M}(A)\) is contained in the unit ball.\\
  We need to prove that it is closed, so consider \(\phi_\alpha\rightarrow\phi\) a net in \(\mathcal{M}(A)\). We need to prove that \(\phi\) is multiplicative:\\
  \(\phi(ab)=\lim_{\alpha}\phi_\alpha(ab)=\lim_{\alpha}\phi_\alpha(a)\phi_\alpha(b)=\lim_{\alpha}\phi_\alpha(a)\lim_{\alpha}\phi_\alpha(b)=\phi(a)\phi(b)\).\par
  Since
  \[
    \phi(1)=\lim_\alpha \phi_\alpha(1)=1,
  \]
  we conclude that \(\phi\in\mathcal{M}(A)\). It follows that \(\mathcal{M}(A)\) is closed in the unit ball,
  which is weak-* compact by Banach-Alaoglu.
\end{proof}
\begin{lemma}\label{lemma:144}
  Let \(a\) be an element in a unital commutative \cstar\ \(A\).
  Evaluation at \(a\) is a continuous map from \(\mathcal{M}(A)\) onto \(\operatorname{spec}(a)\).
\end{lemma}
\begin{proof}
  Continuity is direct consequence of the definition of the weak-* 
  topology. Secondly, for any \(\varphi\) in \(\mathcal{M}(A)\), \(\varphi(a)\) is in \(\operatorname{spec}(a)\) from Lemma \ref{lemma:142}.\par
  Finally, we need to see that the map is onto. That is, let \(\lambda\) be in \(\operatorname{spec}(a)\). A
  Zorn's Lemma argument shows that \(\lambda1-a\) is contained in a maximal
  proper ideal of \(A\), say \(I\). We claim that \(I\) is closed. Its closure is clearly
  an ideal. Moreover, it cannot contain the unit since it is invertible, the
  invertibles are open and \(I\) contains no invertible. Thus, the closure of \(I\) is
  also a proper ideal and hence by maximality, it is equal to \(I\). The quotient
  \(A/I\) is then a field and also a Banach algebra, using the quotient norm. Let \(b\)
  be any element of this Banach algebra. As its spectrum is non-empty (\ref{theorem:134}),
  we have a complex number \(\lambda\) such that \(\lambda1 - b\) is not invertible. As we are
  in a field, we have \(\lambda1-b = 0\). Hence, we see that every element of \(A/I\) is a
  scalar multiple of the unit. That is, \(A/I \cong \mathbb C\). The quotient map from \(A\) to
  \(\mathbb{C}\) is then a non-zero homomorphism which sends \(a\) to \(\lambda\), since \(\lambda1 - a\) is in
  \(I\).  
\end{proof}
\begin{lemma}
  Let \(a\) be an element in a unital, commutative \cstar,
  \(A\), and assume that \(A\) is the Banach algebra generated by \(a\), \(a^*\) and \(1\). Then
  evaluation at \(a\) is homeomorphism from \(\mathcal{M}(A)\) to \(\operatorname{spec}(a)\).
\end{lemma}
\begin{theorem}\label{theorem:146}\label{Petersen:116}
  Let \(A\) be a commutative, unital \cstar. The function
  sending \(a\) in \(A\) to \(\hat{a}\) in \(C(\mathcal{M}(A))\) defined by
  \[
    \hat{a}(\phi) = \phi(a)\text{, }\phi\in\mathcal{M}(A),
  \]
  is an isometric *-isomorphism from \(A\) to \(C(\mathcal{M}(A))\).
\end{theorem}
\begin{proof}
  The fact that $\hat a$ is a continuous function is a simple 
  consequence of the definition of the weak-$*$ topology. 
  It is easy to see that the map $a \mapsto \hat a$ is a 
  homomorphism. That it is a $*$-homomorphism is a 
  consequence of Lemma~\ref{lemma:142}. Our next aim is 
  to see that it is isometric. First, suppose that $a$ 
  is self-adjoint. From Theorem~\ref{theorem:136}, we 
  have 
  \[ \|a\| = r(a) = \sup\{ |\phi(a)| \mid \phi \in \mathcal{M}(A)\} = \sup\{ |\hat a(\phi)| \mid \phi \in \mathcal{M}(A)\} = \|\hat a\|. \] 
  For an arbitrary element, we then have 
  \[ \|a\| = \|a^*a\|^{1/2} = \|\hat{a^*a}\|^{1/2} = \|\hat a^* \hat a\|^{1/2} = \|\hat a\|. \]\par 
  We must finally show that $\hat{\cdot}$ is onto. 
  The range is clearly a unital $*$-subalgebra. 
  We will prove that the range separates the points of $\mathcal{M}(A)$. 
  Assuming this for the moment, we can apply the Stone--Weierstrass Theorem, 
  which states that a unital $*$-subalgebra of $C(\mathcal{M}(A))$ which separates points is dense. 
  Moreover, since $\hat{\cdot}$ is isometric and $A$ is complete, 
  the range is also closed, and we are done.\par 
  As for the claim that the range separates the points, 
  let $\phi \neq \psi$ be two elements of $\mathcal{M}(A)$. 
  The fact that they are unequal means that there exists 
  an $a$ in $A$ such that $\phi(a) \neq \psi(a)$. 
  In other words, \[ \hat a(\phi) \neq \hat a(\psi), \] 
  and so $\hat a$ separates $\phi$ and $\psi$.
\end{proof}
\begin{proposition}
  Let \(B\) be a unital \cstar\ and let \(A\) be a commutative
  \(C^*\)-subalgebra of \(B\) which contains the unit of \(B\). An element \(a\) of \(A\) has an
  inverse in \(A\) if and only if it has an inverse in \(B\). In consequence, we have
  \[
    \operatorname{spec}_A(a) = \operatorname{spec}_B(a).
  \]
\end{proposition}
\begin{definition}\label{Petersen:118}
  Let \(B\) be a unital \cstar\ and let \(a\) be a normal element
  of \(B\). Let \(A\) be the \(C^*\)-subalgebra of \(B\) generated by \(a\) and the unit. For each
  \(f\) in \(C(\operatorname{spec}(a))\), we let \(f(a)\) be the unique element of \(A\) such that
  \[
    \phi(f(a)) = f(\phi(a)),
  \]
  for all \(\phi\) in \(\mathcal{M}(A)\).
\end{definition}
\begin{corollary}[Continuous functional calculus]\label{corollary:149}
  Let \(B\) be a unital \cstar\ and let \(a\) be a normal element of \(B\). The map sending \(f\) to \(f(a)\) is an isometric *-isomorphism from
  \(C(\operatorname{spec}(a))\) to the \(C^*\)-subalgebra of \(B\) generated by \(a\) and the unit. Moreover, if 
  \[
    f(z) = \sum_{k,l} a_{k,l}z^k\overline{z}^l
  \]
  is any polynomial in \(z\) and \(\overline{z}\), then
  \[
   f(a) =\sum_{k,l}a_{k,l}a^k{(a^*)}^l.
  \]
\end{corollary}
\begin{lemma}\label{lemma:125}
  Let \(K\) be a non-empty compact subset of \(\mathbb{R}\), and let \(f:K\rightarrow\mathbb{C}\) be a continuous function.
  Let \(A\) be a unital \cstar, and let \(\Omega_K\) be the set of self-adjoint elements in \(A\) with spectrum contained in \(K\).
  The function
  \[f:\Omega_K\rightarrow A\qquad a\mapsto f(a),\]
  is continuous.
\end{lemma}
\begin{proof}
  The map \(A \rightarrow A\) given by \(a\mapsto a^n\) is continuous for every integer \(n\geq 0\) (because multiplication is continuous).
  It follows that every complex polynomial \(f\) induces a continuous map \(A\rightarrow A\) given by \(a\mapsto f(a)\).\par
  Now, let \(f:K\rightarrow\mathbb{C}\) be any continuous function, let \(a\) be an element in \(\Omega_K\), and let \(\epsilon>0\).
  By the Stone-Weierstrass theorem there is a complex polynomial \(g\) such that \(|f(z)-g(z)|\leq\epsilon/3\) for every \(z\) in \(K\).
  Find \(\delta>0\) such that \(\|g(a)-g(b)\|\leq\epsilon/3\) whenever \(b\) is an element in \(A\) with \(\|a-b\|\leq\delta\). Since
  \[
    \|f(c)-g(c)\|=\|(f-g)(c)\|=\sup\{|(f-g)(z)| | z\in \operatorname{spec}(c)\}\leq\epsilon/3
  \]
  for all \(c\) in \(\Omega_K\), it follows that \(\|f(a)-f(b)\|\leq\epsilon\) for all \(b\) in \(\Omega_K\) with \(\|a-b\|\leq\delta\).
\end{proof}
\begin{corollary}\label{corollary:1410}
  Let \(a\) be a normal element of the unital \cstar\ \(B\).
  Then \(a\) is self-adjoint if and only if \(\operatorname{spec}(a) \subseteq\mathbb{R}\).
\end{corollary}
\begin{corollary}
  Let \(A\) and \(B\) be unital \cstars\ and let \(\rho:A\rightarrow B\) be a
  unital \(C^*\)-homomorphism. If \(a\) is a normal element of \(A\), then \(\operatorname{spec}(\rho(a)) \subseteq
  \operatorname{spec}(a)\) and \(f(\rho(a)) = \rho(f(a))\), for any \(f\) in \(C(\operatorname{spec}(a))\).
\end{corollary}

In the non-unital case, if \(x\) is a normal element of \(A\) and
\(f\in C_0(\operatorname{spec}(x)\setminus\{0\})\), then we denote by
\(f(x)\) the element of \(A\) corresponding to \(f\) via the embedding of
\(C_0(\operatorname{spec}(x)\setminus\{0\})\) into \(A\) given by \ref{Petersen:118}.
\begin{corollary}\label{Petersen:119}
  Let \(f\) be a continuous complex-valued function on \(\mathbb C\) vanishing at \(0\).
  If \(\{x_n\}\) is a sequence of normal elements converging to \(x\), then
  \[
    f(x_n)\longrightarrow f(x).
  \]
\end{corollary}
\begin{proof}
  By the Stone--Weierstrass theorem, \(f\) can be approximated uniformly on every compact subset of \(\mathbb C\) 
  by polynomials in \(z\) and \(\overline z\). 
  Since such polynomials are continuous with respect to the norm topology on normal elements, 
  the conclusion follows by approximation.
\end{proof}
\begin{exercise}\label{ex:142}
  Suppose that \(X\) and \(Y\) are compact Hausdorff spaces and
  \(\rho : C(X) \rightarrow C(Y)\) is a unital *-homomorphism.    
  \begin{enumerate}
    \item Prove that there exists a continuous function \(h : Y \rightarrow X\) such that
    \(\rho(f) = f \circ h\), for all \(f\) in \(C(X)\).
    \item Prove that this statement may be false if \(\rho\) is not unital.
    \item Give a necessary and sufficient condition on \(h\) for \(\rho\) to be injective.
    \item Give a necessary and sufficient condition on \(h\) for \(\rho\) to be surjective.
  \end{enumerate}
\end{exercise}
\begin{exercise}
  Let \(B\) be the \cstar\ of bounded functions on \([0,1]\) (with the supremum norm). 
  Let \(Q = \{k2^{-n} | k \in \mathbb{Z},n \geq 0\}\cap [0,1)\).
  \begin{enumerate}
    \item Let \(\mathcal{A}\) be the linear span of all functions of the form \(\chi_{[s,t)},\chi_{[t,1]}\), where
    \(s, t\) are in \(Q\). Prove that \(\mathcal{A}\) is a \cstar.
    \item Let \(A\) be the closure of \(\mathcal{A}\) in the supremum norm, which is a \cstar.
    Prove that \(A\) contains \(C[0,1]\).
    \item Prove that for all \(f\) in \(A\) and \(s\) in \(Q\),
    \[
      \lim_{x\rightarrow s^+}f(x)\mbox{, }\lim_{x\rightarrow s^-}f(x)
    \]
    both exist.
    \item Let \(h : \mathcal{M}(A) \rightarrow [0,1]\) be the map given in the last exercise and the
    inclusion \(C[0,1] \subseteq A\). Describe \(\mathcal{M}(A)\) and \(h\).
  \end{enumerate}
\end{exercise}
\begin{exercise}
  Let \(B\) be a unital \cstar\ and suppose \(a\) is a self-adjoint
  element and \(\frac{1}{4}>\epsilon > 0\) such that \(\|a-a^2\| < \epsilon\).
  \begin{enumerate}
    \item Let \(f(x) = x-x^2\). Prove that
    \[
      |f(x)| \geq \min\{ \frac{|x|}{2},\frac{1-x}{2}\}
    \]
    for all \(x\) in \(\mathbb{R}\).
    \item Prove that \(\operatorname{spec}(a) \subseteq [-2 \epsilon,2\epsilon ] \bigcup  [1-2\epsilon ,1+2\epsilon ]\).
    \item Find a continuous function \(g\) on \(\operatorname{spec}(a)\) with values in \(\{0,1\}\) such that
    \(|g(x) -x| < 2\epsilon\), for all x in \(\operatorname{spec}(a)\). Also explain why there is no such
    function if \(\operatorname{spec}(a)\) is replaced by \([0,1]\).
    \item Prove that there is a projection \(p\) in \(B\) such that \(\|a-p\| < 2\epsilon\).
  \end{enumerate}
\end{exercise}

\section{Further consequences of the \texorpdfstring{\(C^*\)}{C*}-condition}
In this section, we will prove three more important results on the structure
of \cstars.\\
The first is a nice and somewhat surprising extension of our earlier result
that any *-homomorphism between \cstars\ is necessarily a contraction
(Corollary \ref{cor:137}). If we additionally assume the map is injective, then it is
actually isometric.
\begin{lemma}\label{lemma:151}
  Let \(A\) and \(B\) be unital \cstars\ and \(\rho: A \rightarrow B\) be an
  injective, unital *-homomorphism. Then for any normal element \(a\) in \(A\), we
  have \(\operatorname{spec}(a) = \operatorname{spec}(\rho(a))\).
\end{lemma}
\begin{proof}
  As \(\rho\) maps invertible elements to invertible elements, it is clear that
  \(\operatorname{spec}(\rho(a)) \subseteq \operatorname{spec}(a)\).\\
  Let us suppose that the containment \(\operatorname{spec}(\rho(a)) \subseteq \operatorname{spec}(a)\) is proper.
  Then we may find a non-zero continuous function \(f\) defined on \(\operatorname{spec}(a)\) whose
  restriction to \(\operatorname{spec}(\rho(a))\) is zero. Then we have \(f(a)\neq 0\) while, using Corollary \ref{corollary:149}, we have
  \[
    \rho(f(a)) = f(\rho(a)) = 0
  \]
  since \(f|_{\operatorname{spec}(\rho(a))} = 0\). This contradicts the hypothesis that \(\rho\) is injective.
\end{proof}
\begin{theorem}\label{theorem:152}
  If \(A\) and \(B\) are unital \cstars\ and \(\rho : A \rightarrow B\) is an
  injective, unital *-homomorphism, then \(\rho\) is an isometry; that is, \(\|\rho(a)\| = \|a\|\), for
  all \(a\) in \(A\).
\end{theorem}
\begin{proof}
  The equality \(\|\rho(a)\| = \|a\|\) for a self-adjoint element \(a\) follows from
  Theorem \ref{theorem:136} and Lemma \ref{lemma:151}. For arbitrary \(a\), we have
  \[
    \|a\|^2 = \|a^*a\| = \|\rho(a^*a)\| = \|{\rho(a)}^*\rho(a)\| = \|\rho(a)\|^2
  \]
  and we are done.
\end{proof}
\begin{theorem}\label{theorem:153}
  Let \(B\) be a unital \cstar\ and let \(A\) be a \(C^*\)-subalgebra
  of \(B\) containing its unit. If \(a\) is any element of \(A\), then its spectrum in \(B\)
  coincides with its spectrum in \(A\).
\end{theorem}
\begin{proof}
  It suffices to show that if \(a\) is any element of \(A\) which has an inverse
  in \(B\), then that inverse actually lies in \(A\). In the case that \(a\) is normal, this
  follows from Lemma \ref{lemma:151} applied to the inclusion map of \(A\) in \(B\).\\
  For arbitrary \(a\), if \(a\) is invertible, then so is \(a^*\) (its inverse is \({(a^{-1})}^*)\) and
  hence \(a^*a\) is also invertible. As \(a^*a\) is self-adjoint and hence normal and
  since it clearly lies in \(A\), \({(a^*a)}^{-1}\) also lies in \(A\). Then we observe 
  \(a^{-1} = {(a^*a)}^{-1}(a^*a)a^{-1} = {(a^*a)}^{-1}a^*\) 
  which obviously lies in \(A\).
\end{proof}
\begin{corollary}\label{corollary:154}
  Let \(a\) be a normal element of a unital \cstar\ \(B\). For
  any continuous function \(f\) on \(\operatorname{spec}(a)\), we have \(\operatorname{spec}(f(a)) = f(\operatorname{spec}(a))\).
  That is, the spectrum of \(f(a)\) is simply the range of \(f\).
\end{corollary}
\begin{proof}
  We have already noted the fact that for any compact Hausdorff space
  \(X\) and for any \(f\) in \(C(X)\), the spectrum of \(f\) is simply the range of \(f\), \(f(X)\).
  We apply this to the special case of \(f\) in \(C(\operatorname{spec}(a))\) to see that the spectrum of
  \(f\) is \(f(\operatorname{spec}(a))\). Since the map from \(C(\operatorname{spec}(a))\) to the \cstar\ generated
  by \(a\) is an isomorphism and it carries \(f\) to \(f(a)\) (by definition), we see that
  \(\operatorname{spec}(f) = \operatorname{spec}(f(a))\). The conclusion follows from the fact that the spectrum
  of \(a\) is the same in the \cstar\ \(B\) as it is in the \(C^*\)-subalgebra.
\end{proof}
\begin{exercise}
  Show that Theorem \ref{theorem:153} is false if \(A\) and \(B\) are simply
  algebras over \(\mathbb{C}\). In fact, find a maximal counter-example: \(A \subseteq B\) such
  that for every element \(a\) of \(A\) (which is not a multiple of the identity), its
  spectrum in \(A\) is \textit{as different as possible} from its spectrum in \(B\).
\end{exercise}

\section{Positivity}
\begin{lemma}\label{lemma:161}
  Let \(a\) be a self-adjoint element of a unital \cstar\ \(A\).
  \begin{enumerate}
    \item  Let \(f(x) = \max{\{x,0\}}\) and \(g(x) = \max{\{-x,0\}}\), for \(x \in \mathbb{R}\). Then
    \(f(a),g(a)\) are both positive. We have
    \begin{align*}
      a &= f(a)-g(a),\\
      af(a) &= f(a)^2,\\
      ag(a) &= -g(a)^2,\\
      f(a)g(a) &= 0.
    \end{align*}
    \item If \(\operatorname{spec}(a) \subseteq [0,\infty)\), then \(a\) is positive.
  \end{enumerate}
\end{lemma}
\begin{proof}
  It is clear that \(f\) and \(g\) are continuous functions on the spectrum of
  \(a\), each is real-valued, 
  \(f(x) - g(x) = x\mbox{, }xf(x) = {f(x)}^2\mbox{, }xg(x) = -{g(x)}^2\mbox{ and }f(x)g(x) = 0\), 
  for all \(x\) in \(\mathbb{R}\). Corollary \ref{corollary:149} then implies that \(f(a),g(a)\)
  are self-adjoint elements of \(A\), satisfying the desired equations. Since \(f(x)\)
  is positive, it has a square root and it follows that \(f(a) = {(\sqrt{f(a)})}^2 =
  {(\sqrt{f(a)})}^*(\sqrt{f(a)})\) is positive. The same argument shows \(g(a)\) is positive.
  For the second statement, it suffices to notice in the proof above that if
  \(\operatorname{spec}(a) \subseteq [0,\infty)\), then \(g(a) = 0\).
\end{proof}
\begin{lemma}\label{lemma:162}
  Let \(a\) be a self-adjoint element of a unital \cstar\ \(A\). The
  following are equivalent.
  \begin{enumerate}
    \item \(\operatorname{spec}(a) \subseteq [0,\infty)\).
    \item For all \(t \geq \|a\|\), we have \(\|t-a\| \leq t\).
    \item For some \(t \geq \|a\|\), we have \(\|t-a\| \leq t\).
  \end{enumerate}
\end{lemma}
\begin{proof}
  Since the spectrum of \(a\) is a subset of the reals and the spectral radius
  of \(a\) is its norm, \(\operatorname{spec}(a) \subseteq [-\|a\| , \|a\| ]\). We consider the function \(f_t(x) = t-x\),
  for values of \(t \geq \|a\|\). The function is positive and monotone decreasing on
  \([-\|a\| , \|a\| ]\). Hence, the norm of the restriction of \(f_t\) to the spectrum of \(a\) is
  just its value at the infimum of \(\operatorname{spec}(a)\). In view of Corollary \ref{corollary:149}, \(\|t-a\|\) is
  also the value of \(f_t\) at the infimum of \(\operatorname{spec}(a)\). Finally, notice that \(f_t(x) \leq t\)
  if and only if \(x \geq 0\). Putting this together, we see that the minimum of
  \(\operatorname{spec}(a)\) is negative if and only if \(\|t - a\| > t\). Taking the negations of both
  statements, \(\operatorname{spec}(a) \subseteq [0,\infty)\) if and only if \(\|t - a\| \leq t\). This holds for all
  \(t \geq \|a\|\).
\end{proof}
\begin{proposition}\label{prop:163}
  If \(a\) and \(b\) are self-adjoint elements of the unital \cstar\ \(A\) and \(\operatorname{spec}(a),\operatorname{spec}(b) \subseteq [0,\infty)\), 
  then \(\operatorname{spec}(a + b) \subseteq [0,\infty)\).
\end{proposition}
\begin{proof}
  Let \(t = \|a\| + \|b\|\) which is evidently at least \(\|a+b\|\). Then we have
  \begin{align*}
    \|t -(a+b)\| &= \|( \|a\| -a)+( \|b\| -b)\|\\
    &\leq \|( \|a\| -a)\| + \|( \|b\| -b)\|\\
    &\leq \|a\| + \|b\|\\
    &= t,
  \end{align*}
  where we have used the last lemma in moving from the second line to the
  third. The conclusion follows from another application of the last lemma.
\end{proof}
\begin{lemma}\label{lemma:164}
  Let \(a,b\) be two elements of a unital algebra \(A\). Then we have
  \(\operatorname{spec}(ab) \setminus \{0\} = \operatorname{spec}(ba) \setminus \{0\}\).
\end{lemma}
\begin{proof}
  It clearly suffices to prove that if \(\lambda\) is a non-zero complex number
  such that \(\lambda - ab\) is invertible, then \(\lambda - ba\) is invertible also. 
  Consider \(x = \lambda^{-1} +\lambda^{-1}b(\lambda-ab)^{-1}a\). We see that
  \begin{align*}
    x(\lambda -ba) &= (\lambda^{-1}+\lambda^{-1}b(\lambda-ab)^{-1}a)(\lambda-ba)\\
    &= 1-\lambda^{-1}ba+\lambda^{-1}b(\lambda-ab)^{-1}(\lambda a-aba)\\
    &= 1-\lambda^{-1}ba+\lambda^{-1}b(\lambda-ab)^{-1}(\lambda-ab)a\\
    &= 1-\lambda^{-1}ba+\lambda^{-1}ba\\
    &= 1.
  \end{align*}
  A similar computation which we omit shows that \((\lambda - ba)x = 1\) and we are
  done.
\end{proof}
\begin{theorem}\label{theorem:165}
  Let \(a\) be a self-adjoint element of a unital \cstar\ \(A\).
  Then \(a\) is positive if and only if \(\operatorname{spec}(a) \subseteq [0,\infty)\).
\end{theorem}
\begin{proof}
  The \textit{if} direction has already been done in Lemma~\ref{lemma:161}. Let us now
  assume that \(a = b^*b\), for some \(b\) in \(A\). Using the notation of \ref{lemma:161}, consider
  \(c = bg(a)\). Then we have \(c^*c = g(a)b^*bg(a) = g(a)ag(a) = -{g(a)}^3\). Write
  \(c = d+ie\), where \(d,e\) are self-adjoint elements of \(A\). A simple computation
  shows that
  \[
    cc^* = d^2 +e^2 +i(ed - de) = 2d^2 +2e^2 +{g(a)}^3.
  \]
  because \(c^*c=d^2+e^2+i(de-ed)\).\\
  As the functions \(x^2\) and \({g(x)}^3\) are positive, it follows from Corollary~\ref{corollary:154}
  that each of \(d^2,e^2\) and \({g(a)}^3\) has spectrum contained in \([0,\infty)\). By Proposition~\ref{prop:163}, 
  so does \(cc^*\). On the other hand, again using Corollary~\ref{corollary:154}, we have
  \(\operatorname{spec}(c^*c) = \operatorname{spec}(-{g(a)}^3) \subseteq (-\infty,0]\).
  We now appeal to Lemma \ref{lemma:164} (using \(a = c,b = c^*\)) to conclude that
  \(\operatorname{spec}(c^*c) = \operatorname{spec}(cc^*) = \{0\}\). But this means that \(-{g(a)}^3 = c^*c = 0\) and it
  follows that \(g(a) = 0\). This implies that the restriction of \(g\) to the spectrum
  of \(a\) is zero, which means \(\operatorname{spec}(a) \subseteq [0,\infty)\) and we are done.
\end{proof}
\begin{corollary}
  If \(a\) is a normal element in a unital \cstar\ \(B\) and \(f\)
  is a real-valued function on \(\operatorname{spec}(a)\), then \(f(a)\) is self-adjoint. Similarly, if \(f\)
  is positive, then \(f(a)\) is positive.
\end{corollary}
\begin{proof}
  We know that \(\operatorname{spec}(f(a)) = f(\operatorname{spec}(a))\) from Corollary~\ref{corollary:154}. Moreover,
  since \(f(a)\) lies in the \cstar\ generated by \(a\) and the unit, which
  is commutative, it itself is normal. The two statements now follow from Corollary~\ref{corollary:1410}
  and Theorem~\ref{theorem:165}.
\end{proof}

\begin{lemma}\label{Petersen:131}
  In the following conditions, when \(A\) is non-unital, \(1\) denotes the identity of the unitization \(\tilde A\),
  whose construction is given in Section~\ref{sec:adjoiningUnit}.
  The following four conditions on an element \(x\) in \(A\) are equivalent:
  \begin{enumerate}
    \item \(x\) is normal and \(\operatorname{spec}(x)\subset\mathbb{R}_+\);
    \item \(x=y^2\) with \(y\) in \(A_{sa}\);
    \item \(x=x^*\) and \(\|t1-x\|\leq t\) for any \(t\geq\|x\|\);
    \item \(x=x^*\) and \(\|t1-x\|\leq t\) for some \(t\geq\|x\|\).
  \end{enumerate}
\end{lemma}
\begin{proof}
  \begin{itemize}
    \item \((1)\Rightarrow(2)\). Using \ref{Petersen:118}, we set \(y=x^{1/2}\) and have \(y^2=x\).
    \item \((2)\Rightarrow(1)\). Embedding \(x\) and \(y\) in a commutative \(C^*\)-subalgebra, we see that \(x=x^*\) and \(\operatorname{spec}(x)\subset\mathbb{R}_+\).
    \item \((1)\Rightarrow(3)\). From \ref{Petersen:116} we have \(\|z\|=\sup\{|\lambda| | \lambda\in \operatorname{spec}(z)\}\) for each normal element \(z\) of \(\tilde{A}\).
      Applying this to \(t1-x\) with \(t\geq\|x\|\), we have
      \[
        \|t1-x\|=\sup\{|t-\lambda| | \lambda\in \operatorname{spec}(x)\}\leq t.
      \] 
    \item \((3)\Rightarrow(4)\). It is immediate.
    \item \((4)\Rightarrow(1)\). If \(\lambda\in \operatorname{spec}(x)\), then \(t-\lambda\in \operatorname{spec}(t1-x)\), whence
      \[
        |t-\lambda|\leq\|t1-x\|\leq t.
      \]
      Therefore \(\lambda\geq 0\).
  \end{itemize}
\end{proof}

\begin{theorem}\label{Petersen:133}
  The set \(A_+\) is a closed real cone in \(A_{sa}\), and \(x \in A_+\) if and only if \(x =y^* y\) for
  some \(y\) in \(A\).
\end{theorem}
\begin{proof}
  From \ref{Petersen:131} (3) it is clear that \(A_+\) is a closed subset of \(A_{sa}\) stable under
  multiplication with positive scalars. To prove that \(A_+\) is a cone, take \(x\) and \(y\) in
  \(A_+\). By \ref{Petersen:131} (3) we have
  \begin{align*}
    \|(\|x\|+\|y\|)1-(x +y)\|&=\|(\|x\|1-x)+(\|y\|1-y)\|\\
    &\leq\|\|x\|1-x\|+\|\|y\|1-y\|\leq\|x\|+\|y\|,
  \end{align*}
  whence \(x +y \in A_+\) by \ref{Petersen:131} (4) since \(\|x\|+\|y\|\geq\|x +y\|\).\par
  Assume now that \(x = y^* y\). Then \(x = x^*\), so that \(x = x_+ - x_-\) by \ref{Petersen:118}.
  Moreover,
  \[
  (yx_-^{1/2})^*(y x^{1/2}_-)=x^{1/2}_- (x_+-x_-)x^{1/2}_-=-x_-^2\in -A_+.
  \]
  Put \(yx^{1/2}_-=a+ib\) with \(a\) and \(b\) in \(A_{sa}\). Then
  \[
    (yx^{1/2}_- )(yx^{1/2}_- )^*=
    2(a^2+b^2)-(yx^{1/2}_-)^*(yx_-^{1/2}) \in A_+.
  \]
  since \(A_+\) was a cone. But, zero apart, the spectrum of a product does not depend
  on the order of the factors (A.1, Appendix), whence
  \[
    \operatorname{spec}(x^2_-) \subset \mathbb{R}_+ \cap -\mathbb{R}_+ =\{0\},
  \]
  so that \(x_- =0\) and \(x \geq 0\).
\end{proof}

\begin{proposition}\label{Petersen:135}
  If \(0\leq x\leq y\), then \(a^* x a\leq a^* ya\) for each \(a\) in \(A\). Further, \(\|x\|\leq\|y\|\).
\end{proposition}
\begin{proof}
  Since \(y-x=b^* b\), from \ref{Petersen:133} we have \(a^* ya-a^* xa={(ba)}^*(ba)\in A_+\).\par
  Adjoining a unit to \(A\), we have \(y\leq\|y\|1\) from spectral theory. Then \(x\leq\|y\|1\), whence \(\|x\|\leq\|y\|\).
\end{proof}

\begin{proposition}\label{Petersen:136}
  If \(1 \in A\) and \(x\) and \(y\) are invertible elements in \(A_+\) with \(x \leq y\), then \(y^{-1} \leq x^{-1}\).
\end{proposition}
\begin{proof}
  From \ref{Petersen:135} we have \(y^{-1/2}x y^{-1/2} \leq 1\), whence \(\|x^{1/2} y^{-1/2}\| \leq 1\), and
  thus \(\|x^{1/2}y^{-1} x^{1/2}\| \leq 1\), which implies that \(x^{1/2}y^{-1} x^{1/2} \leq 1\). By \ref{Petersen:135} 
  \(y^{-1} \leq x^{-1/2} 1 x^{-1/2} = x^{-1}\).
\end{proof}

\begin{definition}
  We say that a continuous real function \(f\) on an interval in \(\mathbb{R}\) is
  \textnormal{operator monotone (increasing)} if \(x\leq y\) implies \(f(x)\leq f(y)\)
  whenever the spectra of \(x\) and \(y\) belong to the interval of definition for \(f\).
\end{definition}
 
\begin{proposition}\label{Petersen:137}
  For each \(\alpha>0\), define \(f_\alpha\) on \(]-1/\alpha,\infty[\) by
  \[
    f_\alpha(t)={(1+\alpha t)}^{-1}t={[1-{(1+\alpha t)}^{-1}]}/\alpha.
  \]
  Since the process of taking inverses is operator monotone decreasing by \ref{Petersen:136},
  it is easy to see that \(f_\alpha\) is operator monotone increasing on \(]-1/\alpha,\infty[\).
  The family of functions \(\{f_\alpha\}\) is further used repeatedly. Note that \(f_\alpha(t)\leq\min\{t,1/\alpha\}\),
  with strict inequality when \(t\neq0\),
  and that \(\lim f_\alpha(t)=t\) uniformly on compact subsets of \(\mathbb{R}\) as \(\alpha\rightarrow 0\).
  Moreover, \(f_\alpha\geq f_\beta\) when \(\alpha\leq\beta\), and \(f_\alpha\circ f_\beta=f_{\alpha+\beta}\) on \(]-{(\alpha+\beta)}^{-1},\infty[\).
  Finally, if \(t>0\), then \(\lim\alpha f_\alpha(t)=1\) uniformly on compact subsets of \(]0,\infty[\) as \(\alpha\rightarrow \infty\).
\end{proposition}

\begin{definition}
  Let \(A\) be a \cstar\. A net \(\{u_\lambda| \lambda\in \Lambda\}\) in \(A_+\) with \(\|u_\lambda\|\leq 1\) for all \(\lambda\)
  is called an \textnormal{approximate unit} for \(A\) if \(\lambda<\mu\) implies \(u_\lambda\leq u_\mu\)
  and if \(\lim\|x(1-u_\lambda)\|=0\) for each \(x\) in \(A\).
  Then, of course, \(\lim\|(1-u_\lambda)x\|=0\) as well.
\end{definition}
\begin{theorem}\label{Petersen:142}
  Each \cstar\ contains an approximate unit.
\end{theorem}
\begin{proof}
  Consider the set \(\Lambda\) of elements \(u\) in \(A_+\) such that \(\|u\|< 1\).
  we claim that \(\Lambda\) is an increasing net in the partial ordering on \(A_{sa}\).
  To see this, take \(u\) and \(v\) in \(\Lambda\).
  Then the elements \(a={(1-u)}^{-1}u\) and \(b={(1-v)}^{-1}v\)
  belong to \(A_+\) by \ref{Petersen:119}.
  Set \(w={(1+a+b)}^{-1}(a+b)\). Then \(w\in \Lambda\). By \ref{Petersen:137} we have
  \[
    w=f_1(a+b)\geq f_1(a)={[1+{(1-u)}^{-1}]}^{-1}[{(1-u)}^{-1}u]=u,
  \]
  and similarly \(w\geq v\), which proves that \(\Lambda\) is an increasing net in \(A_+\).\par
  To show that \(\Lambda\) is an approximate unit, it suffices to prove that the net 
  \(\{\|x(1-u)x\| | u\in\Lambda\}\), which is decreasing by \ref{Petersen:135}, converges to zero for each \(x\) in \(A_+\), since
  \[
    \|(1-u)x\|^2=\|x(1-u)^2 x\|\leq\|x(1-u)x\|,
  \]
  and \(A\) is linearly spanned by \(A_+\).\par
  With \(f_\alpha\) as in \ref{Petersen:137}, note that \(\alpha f_\alpha(x)\in\Lambda\) and that
  \[
    x(1-\alpha f_\alpha(x))x=(1+\alpha x)^{-1}x^2\leq \alpha^{-1}x,
  \]
  whence \(\|x(1-\alpha f_\alpha(x))x\|\leq\alpha^{-1}\|x\|\), which tends to zero as \(\alpha\rightarrow\infty\).
\end{proof}
\begin{lemma}\label{Petersen:144}
  Let \(x,y\) and \(a\) be elements of a \cstar\ \(A\) such that \(a\geq 0\) and \(x^* x\leq a^\alpha\), \(y y^*\leq a^\beta\)
  with \(\alpha+\beta>1\). Then the sequence with elements \(u_n=x{[(1/n)+a]}^{-1/2}y\) is norm
  convergent to an element \(u\) in \(A\) with \(\|u\|\leq\|a^{(\alpha+\beta-1)/2}\|\).
\end{lemma}
\begin{proof} 
  Put \(d_{nm}={[(1/n)+a]}^{-1/2}-{[(1/m)+a]}^{-1/2}\). Then
  \begin{align*}
    \|u_n-u_m\|^2&=\|xd_{nm}y\|^2= \|y^* d_{nm}x^* x d_{nm}y\|\\
    &\leq\|y^* d_{nm}a^\alpha d_{nm}y\|=\|a^{\alpha/2}d_{nm}y\|^2\\ 
    &=\|a^{\alpha/2}d_{nm}yy^* d_{nm}a^{\alpha/2}\|\\
    &\leq \|a^{\alpha/2}d_{nm}a^\beta d_{nm}a^{\alpha/2}\|=\|d_{nm}a^{(\alpha+\beta)/2}\|^2.
  \end{align*}
  From spectral theory we see that the sequence \(\{{[(1/n)+a]}^{-1/2}a^{(\alpha+\beta)/2}\}\)
  is increasing and thus by Dini's theorem uniformly convergent to \(a^{(\alpha+\beta-1)/2}\).
  Consequently, \(\|d_{nm}a^{(\alpha+\beta)/2}\|\rightarrow 0\), so that \(\{u_n\}\)
  is norm convergent to an element \(u\) in \(A\). Reasoning as above, we have
  \[
    \|u_n\|=\|x{[(1/n)+a]}^{-1/2}y\|\leq\|a^{\alpha/2}{[(1/n)+a]}^{-1/2}a^{\beta/2}\|\leq\|a^{(\alpha+\beta-1)/2}\|,
  \]
  which shows that \(\|u\|\leq\|a^{(\alpha+\beta-1)/2}\|\).
\end{proof}
\begin{proposition}\label{Petersen:145}
  Let \(x\) and \(a\) be elements in a \cstar\ \(A\) such that \(a\geq 0\) and \(x^* x\leq a\).
  If \(0<\alpha<\frac{1}{2}\), then there is an element \(u\) in \(A\) with \(\|u\|\leq\|a^{1/2-\alpha}\|\)
  such that \(x=u a^\alpha\).
\end{proposition}
\begin{proof}
  Define \(u_n=x{[(1/n)+a]}^{-1/2}a^{1/2-\alpha}\). From \ref{Petersen:144}
  we see that \(\{u_n\}\) is convergent to an element \(u\) in \(A\) with
  \[
    \|u\|\leq\|a^{\frac{1}{2}(1+1-2\alpha-1)}\|=\|a^{\frac{1}{2}-\alpha}\|.
  \]
  Furthermore,
  \begin{align*}
    \|x-u_n a^\alpha\|^2&=\|x(1-{[(1/n)+a]}^{-1/2}a^{1/2})\|^2\\
    &\leq\|a^{1/2}(1-{[(1/n)+a]}^{-1/2}a^{1/2})\|^2\rightarrow 0
  \end{align*}
  as \(n\rightarrow\infty\) by spectral theory (Dini's theorem).
  It follows that \(x=ua^\alpha\).
\end{proof}

\begin{exercise}
  Let \(A\) be a unital \cstar\ and suppose \(a\) is in \(A\). Let
  \(f(x) = \sqrt{x},x \geq 0\) and define \(|a| = f(a^*a)\).
  \begin{enumerate}
    \item Show that if \(a\) is invertible, so is \(|a|\).
    \item Show that if \(a\) is invertible, then \(u = a|a|^{-1}\) is unitary.
    \item Prove that if \(a\) is invertible, \(u\) and \(|a|\) commute if and only if \(a\) is
      normal.
  \end{enumerate}
  The expression \(a = u|a|\) is usually called the polar decomposition of \(a\).
\end{exercise}

\section{Finite-dimensional \texorpdfstring{\(C^*\)}{C*}-algebras}\label{sec:17}
\begin{theorem}\label{theorem:171}
  Let \(A\) be a unital, finite-dimensional \cstar. Then
  there exist positive integers \(K\) and \(N_1,\ldots,N_K\) such that
  \[
    A\cong\oplus_{k=1}^K M_{N_k}(\mathbb{C}).
  \]
  Moreover, \(K\) is unique and \(N_1,\ldots,N_K\) are unique, up to a permutation.
\end{theorem}
\begin{definition}
  If \(\mathcal{H}\) is a Hilbert space and \(\xi,\eta\) are vectors in \(\mathcal{H}\), then we
  define \(\xi \otimes\eta^* : \mathcal{H}\rightarrow\mathcal{H}\) by
  \[
    \xi \otimes \eta^*(\zeta) = < \zeta,\eta > \xi,\zeta \in \mathcal{H}.
  \]
\end{definition}
\begin{lemma}
  Let \(\xi,\eta,\zeta,\omega\) be vectors in the Hilbert space \(\mathcal{H}\) and let \(a\) be in
  \(\mathcal{B}(\mathcal{H})\). We have
  \begin{enumerate}
    \item \(\xi \otimes \eta^*\) is a bounded linear operator on \(\mathcal{H}\) and \(\|\xi \otimes \eta^*\|\) = \(\|\xi\|\|\eta\|\),
    \item \((\xi \otimes \eta^*)\mathcal{H} = span\{\xi\}\), provided \(\eta\neq 0\),
    \item \({(\xi \otimes \eta^*)}^* = \eta \otimes \xi^*\),
    \item \((\xi \otimes \eta^*)(\zeta \otimes\omega^*) =< \zeta,\eta > \xi \otimes \omega^*\),
    \item \(a(\xi \otimes \eta^*) = (a\xi)\otimes \eta^*\),
    \item \((\xi \otimes \eta^*)a = \xi \otimes{(a^*\eta)}^*\).
  \end{enumerate}
  Moreover, if \(\xi_1,\xi_2,\ldots,\xi_n\) is an orthonormal basis for \(\mathcal{H}\), then
  \[
    \sum_{i=1}^{n} a\xi_i \otimes {\xi_i}^*=a
  \]
  for any \(a\) in \(\mathcal{B}(\mathcal{H})\). In particular, the linear span of 
  \(\{\xi_i \otimes \xi^*_j | 1 \leq i,j \leq n\}\) is \(\mathcal{B}(\mathcal{H})\) and
  \[
    \sum_{i=1}^{n}\xi_i \otimes\xi_i^* = 1.
  \]
\end{lemma}
\begin{lemma}
  Let \(A\) be a unital finite-dimensional \cstar.
\begin{enumerate}
  \item Every normal element in \(A\) has finite spectrum.
  \item Every normal element of \(A\) is a linear combination of projections.
\end{enumerate}
\end{lemma}
\begin{proof}
  For \(a\) normal, from \ref{corollary:149}, we know that \(C(\operatorname{spec}(a))\) is isomorphic to
  a \(C^*\)-subalgebra of \(A\), which must also be finite-dimensional. It follows that
  \(\operatorname{spec}(a)\) is finite.\\
  For each \(\lambda\) in \(\operatorname{spec}(a)\), let \(p_\lambda\) be the element of \(C(\operatorname{spec}(a))\) which is 1 on
  \(\lambda\) and zero elsewhere. As \(\operatorname{spec}(a)\) is finite, this function is continuous. As
  a function, this is clearly a self-adjoint idempotent, so the same is true of
  \(p_\lambda(a)\). Moreover, it follows that
  \[
    \sum_{\lambda\in \operatorname{spec}(a)}\lambda p_\lambda(z) = z,
  \]
  for all \(z\) in \(\operatorname{spec}(a)\). It follows from \ref{corollary:149} that
  \[
    \sum_{\lambda\in \operatorname{spec}(a)}\lambda p_\lambda(a) = a.
  \]
\end{proof}
\begin{lemma}\label{lemma:175}
  Let \(A\) be a \cstar\. The relation defined on projections
  by \(p \geq q\) if \(pq = q\) (and hence after taking adjoints \(qp = q\) also) is a partial
  order.
\end{lemma}
\begin{lemma}\label{lemma:176}
  Let \(A\) be a \cstar.
\begin{enumerate}
  \item For projections \(p,q\) in \(A\), \(p \geq q\) if and only if \(pAp \supset  qAq\). In particular,
    if \(pAp = qAq\), then \(p = q\)
  \item If \(pAp\) has finite dimension greater than \(1\), then there exists \(q\neq 0\) with
    \(p > q\).
  \item If \(A\) is unital and finite-dimensional, then it has minimal non-zero
    projections in the order \(\geq\).
\end{enumerate}
\end{lemma}
\begin{proof}
  First, suppose that \(p \geq q\). Then we have \(qAq = pqAqp \subseteq pAp\).
  Conversely, suppose that \(pAp \supset qAq\). Hence, we have \(q = qqq \in qAq \subseteq pAp\),
  so \(q = pap\), for some \(a\) in \(A\). Then \(pq = ppap = pap = q\). For the last
  statement, if \(pAp = qAq\), then it follows from the first part that \(p \geq q\) and
  \(q \geq p\), hence \(p = qp = q\).\par
  It is easily verified that \(pAp\) is a *-algebra. It contains \(ppp = p\)
  which is clearly a unit for this algebra. If it is finite-dimensional, then it is
  closed and hence is a \(C^*\)-subalgebra. We claim that if a self-adjoint element
  of \(pAp\), say \(a\), has only one point in its spectrum, then that element is a scalar
  multiple of \(p\). Let \(\operatorname{spec}(a) = \{\lambda\}\). Then functions \(f(z) = z\) and \(g(z) = \lambda\) are
  equal on \(\operatorname{spec}(a)\). So by Corollary \ref{corollary:149}, \(f(a) = g(a)\). On the other hand,
  Corollary \ref{corollary:149} also asserts that \(f(a) = a\) while \(g(a) = \lambda p\). Our claim follows
  and this means that if every self-adjoint element of \(pAp\) has only one point
  in its spectrum that \(pAp\) is spanned by \(p\) and one dimensional. So if \(pAp\) has
  dimension greater than one, we must have a self-adjoint element, \(a\), in \(pAp\)
  with at least two points in its spectrum. Choose a surjective function, \(f\),
  from \(\operatorname{spec}(a)\) to \(\{0,1\}\). It is automatically continuous on \(\operatorname{spec}(a)\), \(q = f(a)\)
  is a projection which is non-zero. As \(q\) is in \(pAp\), we have \(pq = q\).\par
  For the last statement, we first notice that since \(A\) is unital, it contains
  non-zero projections. Suppose that \(p\) is any non-zero projection in \(A\), so
  \(pAp\) has dimension at least one. If this dimension is strictly greater than
  one, we may find \(q\) as in part 2. So \(q\) is non-zero and since \(q\neq p\), \(qAq\) is
  a proper linear subspace of \(pAp\), so \(\dim(qAq) < \dim(pAp)\). Continuing in
  this fashion, we may eventually find \(q\) such that \(\dim(qAq) = 1\). Hence, \(q\) is
  minimal in the order \(\leq\) among non-zero projections.
\end{proof}
We will consider non-zero projections \(p\) which are minimal with respect
to the relation \(\geq\). From the last result, we have \(pAp = \mathbb{C}p\), for any such
projection.
\begin{lemma}\label{lemma:177}
  Suppose that \(p_1,p_2,\ldots,p_K\) are minimal, non-zero projections
  in \(A\) with \(p_i Ap_j = 0\), for \(i\neq j\). Then \(p_1,\ldots,p_K\) are linearly independent.
  A finite set of minimal projections satisfying the hypothesis will be called
  independent.
\end{lemma}
\begin{proof}
  If some \(p_i\) can be written as a linear combination of the others, then
  we would have
  \[
    \mathbb{C}p_i = p_i Ap_i = p_i A(\sum_{j\neq i} \alpha_j p_j) \subseteq \sum_{j\neq i} p_i Ap_j = 0,
  \]
  a contradiction.
\end{proof}
Lemma~\ref{lemma:176} guarantees the existence of minimal, non-zero projections
in a unital, finite-dimensional \cstar\ \(A\). Lemma~\ref{lemma:177} shows that an
independent set of these cannot contain more than \(\dim(A)\) elements and it
follows that we may take a maximal independent set of minimal, non-zero
projections and it will be finite.
\begin{theorem}\label{theorem:178}
  Let \(A\) be a finite-dimensional \cstar\ and let
  \(p_1, p_2,\ldots,p_K\) be a maximal set of independent minimal non-zero projections
  in \(A\).
  \begin{enumerate}
    \item For each \(1 \leq k \leq K\), \(Ap_k\) is a finite-dimensional Hilbert space with
    inner product \(< a,b > p_k = b^*a\), for all \(a,b \in Ap_k\).
    \item For each \(1 \leq k \leq K\), \(Ap_k A\) (meaning the linear span of elements of
    the form \(ap_k a'\) with \(a,a'\) in \(A\)) is a unital \(C^*\)-subalgebra of \(A\).
    \item \(\bigoplus_{k=1}^K Ap_k A = A\),
    \item For each \(1 \leq k \leq K\), define the map \(\pi_k : A \rightarrow \mathcal{B}(Ap_k)\) defined by
    \(\pi_k(a)b = ab\), for a in \(A\) and \(b\) in \(Ap_k\). Then \(\pi_k\) is a *-homomorphism
    and its restriction to \(Ap_l A\) is zero for \(l\neq k\) and an isomorphism for
    \(l = k\).
  \end{enumerate}
\end{theorem}
\begin{proof}
  For the first statement, if \(a,b\) are in \(Ap_k\), then \(a = ap_k,b = bp_k\) and
  so \(b^*a = {(bp_k)}^*ap_k = p_k b^*ap_k \in p_k Ap_k = \mathbb{C}p_k\), so the scalar \(< a,b >\) as
  described exists. This is clearly linear in \(a\) and conjugate linear in \(b\). It is
  clearly non-degenerate, for if \(< a,a >= 0\), then \(a^*a = 0\) which implies \(a = 0\).\par
  Choose \(B_k\) to be an orthonormal basis for \(Ap_k\) which means that for any
  a in \(Ap_k\), we have
  \[
  \sum_{b\in B_k}bb^*a =\sum_{b\in B_k}b <a,b > p_k = \sum_{b\in B_k}<a,b > b p_k=\sum_{b\in B_k}<a,b > b=a
  \]
  Next, we define
  \[
  q_k=\sum_{b\in B_k}bb^*=\sum_{b\in B_k}bp_k b^*
  \]
  It is clear that \(q_k\) is self-adjoint, in \(Ap_k A\) and, from above, \(q_k a = a\), for any
  a in \(Ap_k\).\par
  Now if \(a,a'\) are in \(A\), we have
  \[
  q_k (ap_k a') = (q_k ap_k)a'= ap_k a'
  \]
  and
  \[
  (ap_k a')q_k = {(q_k {(a')}^*p_k a^*)}^* = {({(a')}^* p_k a^*)}^* = ap_k a',
  \]
  so that \(q_k\) is a unit for \(Ap_k A\). This completes the proof of part 2. In particular,
  we note that \(q_k p_k = p_k\).\par
  Since \(p_k Ap_l = 0\) for all \(k\neq l\), we see that \(Ap_k A\cdot Ap_l A = 0\) as well. From
  this, we also see that \(q_k Ap_l = 0\) and \(q_k Ap_l A = 0\) for \(k\neq l\). In particular, the
  projections \(q_k\) are pairwise orthogonal.\par
  Consider \(q = \sum^K_{k=1} q_k\)
  which is clearly a central projection. We claim
  that it is the identity for \(A\). If not, there is an element \(b\) in \(A\) with \(bq\ne b\)
  and so \(b- bq\) is a non-zero element in the set
  \[
    q^{\perp } =\{a \in A|qa=0\}.
  \]
  It is a simple matter to verify that \(q^\perp\) is a \(C^*\)-subalgebra of \(A\) and is obviously
  finite-dimensional. If it contains a non-zero element, then it contains a non
  zero minimal projection, say \(p\). Then for any \(k = 1,2,\ldots,K\), we have
  \[pq_k=pqq_k=0\]
  and so \(pq_k = 0\). It then follows that
  \[pAp_k \subseteq pAp_k A = pq_k Ap_k A = 0.\]
  The set of projections \(p,p_1,\ldots,p_K\) is an independent set of non-zero minimal
  projections, contradicting the maximality of \(p_1,\ldots,p_K\). We conclude that \(q\)
  is the identity of \(A\), as desired. In consequence, we have
  \[
  A= \left( \sum_k q_k \right)A = \bigoplus_{k=1}^K q_k A =\bigoplus_{k=1}^K Ap_k A
  \]
  Now, we must establish the last part. We know already that \(Aq_l A\) acts
  trivially on \(Ap_k\), if \(l\neq k\). It remains to verify that \(\pi_k\) is an
  isomorphism from \(Ap_k A\) to \(\mathcal{B}(Ap_k)\). Every element of \(Ap_k\) is in the span of
  \(B_k\). It follows that every element of \(p_k A = {(Ap_k)}^*\) is a linear combination of
  the adjoints of elements of \(B_k\). Therefore \(ap_k a'= (ap_k){(a'p_k)}^*\) is in the span
  of \(bc^*,b,c \in B_k\). For any scalars \(\alpha_{b,c},b,c \in B_k\) and \(b_0,c_0 \in B_k\), we compute
  \[
    <\pi_k\left(\sum_{b,c} \alpha_{b,c}bc^*\right) c_0,b_0 > p_k = 
    \sum_{b,c} \alpha_{b,c} < bc^*c_0,b_0 > p_k = 
    \sum_{b,c} \alpha_{b,c}b^*_0bc^*c_0.
  \]
  Since \(B_k\) is an orthonormal basis, \(b^*_0 b\) is zero unless \(b = b_0\), in which case it
  is \(p_k\). Similarly, \(c^*c_0\) is zero unless \(c = c_0\), in which case it is also \(p_k\). We
  conclude that
  \[
  <\pi_k \left(\sum_{b,c} \alpha_{b,c}bc^*\right) c_0,b_0 > p_k = \alpha_{b_0,c_0}  p_k.
  \]
  If \(\pi_k\left(\sum_{b,c}  \alpha _{b,c}bc^*\right) = 0\), it follows that each coefficient \(\alpha_{b,c}\) is zero and from
  this it follows that the restriction of \(\pi_k\) to \(Ap_k A\) is injective.\par
  We now show \(\pi_k\) is onto. Let \(a,b\) be in \(Ap_k\). We claim that \(\pi_k(ab^*) =
  a \otimes b^*\). For any \(c\) in \(Ap_k\), we have
  \[
  (a \otimes b^*)c =< c,b > a =< c,b > ap_k = a(< c,b > p_k) = a(b^*c) = \pi_k(ab^*)c.
  \]
  We have shown that the rank one operator \(a \otimes b^*\) is in the range of \(\pi_k\) and
  since the span of such operators is all of \(\mathcal{B}(Ap_k)\), we are done.
\end{proof}
\begin{theorem}\label{theorem:179}
  If \(N\) is a positive integer, then the centre of \(M_N(\mathbb{C})\) is the
  scalar multiples of the identity. If \(N_1,N_2,\ldots,N_K\) are positive integers, then
  the centre \(\bigoplus ^K_{k=1}M_{N_k}(\mathbb{C})\) is isomorphic to \(\mathbb{C}^K\) and is spanned by the identity
  elements of the summands.
\end{theorem}
\begin{proof}
  We may assume that \(N \geq 2\), since the other case is trivial. We identify
  \(M_N(\mathbb{C})\) and \(\mathcal{B}(\mathbb{C}^N)\). Suppose that \(a\) is in the centre of \(\mathcal{B}(\mathbb{C}^N)\) and let \(\xi\) be
  any vector in \(\mathbb{C}^N\). Then we have
  \[
  (a\xi \otimes \xi^*) = a(\xi \otimes\xi^*) = (\xi \otimes\xi^*)a = \xi \otimes{(a^*\xi)}^*.
  \]
  Applying both sides to the vector \(\xi\), we see that
  \[
    a\xi\langle\xi,\xi\rangle
    =\langle a\xi,\xi\rangle\xi.
  \]
  It follows that, for every vector \(\xi\) there is a scalar \(r\) such that \(a\xi = r\xi\). If \(\xi\)
  and \(\eta\) are two linearly independent vectors we know that there are scalars,
  \(r, s, t\), such that
  \[r(\xi +\eta) = a(\xi +\eta) = a\xi +a\eta = s\xi+t\eta\]
  By linear independence, we see that \(r = s = t\). We have shown that for every
  vector \(\xi\), \(a\xi\) is a scalar multiple of \(\xi\). Moreover, the scalar is independent of
  \(\xi\). Thus a is a multiple of the identity.\par
  The second statement follows immediately from the first.
\end{proof}
Let us complete the proof of Theorem~\ref{theorem:171}. In fact, almost everything
is done. We know that from part 3 of Theorem~\ref{theorem:178} that \(A = \bigoplus_k Ap_k A\) and
from part 4 of the same theorem that, for each \(k\), \(\pi_k : Ap_k A \rightarrow \mathcal{B}(Ap_k)\) is an
isomorphism. Therefore, we have
\[
\bigoplus^K_{k=1} \pi_k : A = \bigoplus^K_{k=1}Ap_k A \rightarrow \bigoplus^K_{k=1} \mathcal{B}(Ap_k)
\]
is an isomorphism.\par
It only remains for us to prove the uniqueness of \(K\) and \(N_1,N_2,\ldots,N_K\).
We see from the last result that \(K\) is equal to the dimension of the centre of
\(A\). Next, the units of the summands of \(A\) are exactly the minimal non-zero
projections in the centre and if \(q_k\) is the unit of summand \(M_{N_k}\), then \(N_k\) is
the square root of the dimension of \(q_k A\).
\begin{exercise}
  Let \(A\) be a \cstar. Let \(e\) be in \(A\) and satisfy \(e^*e\) is a
  projection. Prove that
  \begin{enumerate}
    \item \(ee^*e = e\). (Hint: compute \({(ee^*e - e)}^*(ee^*e - e)\).)
    \item \(ee^*\) is also a projection.
  \end{enumerate}
\end{exercise}
\begin{exercise}
  Let \(A\) be a \cstar\ and let \(n \geq 1\). Suppose that we have
  elements of \(A\), \(e_{i,j},1 \leq i,j \leq n\) which satisfy
  \begin{equation*}
  e_{i,j}e_{k,l} =
  \begin{cases}    
  e_{i,l} &\text{ if }j =k\\
  0 &\text{ if }j\neq k
  \end{cases}    
  \end{equation*}
  and \(e^*_{i,j} = e_{j,i}\), for all \(i,j,k,l\). Assuming that at least one \(e_{i,j}\) is non-zero,
  prove that the span of the \(e_{i,j}\), \(1 \leq i,j\leq n\), is a \(C^*\)-subalgebra of \(A\) and is
  isomorphic to \(M_n(\mathbb{C})\). (Hint: first show that all \(e_{i,j}\) are non-zero, then show
  they are linearly independent.) Such a collection of elements is usually called
  a set of matrix units.
\end{exercise}
\begin{exercise}
  Suppose that \(a_1,a_2,\ldots,a_N\) are elements in a \cstar\ \(A\)
  and satisfy:
  \begin{enumerate}
    \item \(a^*_1 a_1 = a^*_2 a_2 = \cdots = a^*_N a_N\) is a projection,
    \item \(a_i a^*_i a_j a^*_j = 0\) for all \(1 \leq i\ne j \leq N\).
  \end{enumerate}
  Prove that \(e_{i,j} = a_i a^*_j\), \(1 \leq i,j \leq N\) is a set of matrix units.
\end{exercise}

\section{Unitization and non-unital
\texorpdfstring{\(C^*\)}{C*}-algebras}
\label{sec:adjoiningUnit}
We recall the standard construction of the unitization; compare
\cite[Paragraph~A.1.2]{RordamLarsenLaustsen2000}.
\begin{theorem}\label{theorem:181}
  Let \(A\) be a \cstar. There exists a \cstar\ \(\tilde{A} \) which
  is unital, contains \(A\) as a closed two-sided ideal and \(\tilde{A}/A \cong \mathbb{C}\). Moreover,
  this \cstar\ is unique up to a unique *-isomorphism that restricts to the identity on \(A\).
\end{theorem}
\begin{proof}
  Let \(\tilde{A} = \mathbb{C}\oplus A\), as a vector space. We define the product on \(\tilde{A}\) by
  \[(\lambda, a)(\mu,b) = (\lambda\mu,\lambda b + \mu a +ab),\]
  for \(a,b\) in \(A\) and \(\lambda,\mu\) in \(\mathbb{C}\). We also define
  \[{(\lambda, a)}^* = (\overline{\lambda},a^*),\]
  for \(a\) in \(A\) and \(\lambda\) in \(\mathbb{C}\). It is clear that \((1,0)\) is a unit for this algebra. The
  map sending \(a\) in \(A\) to \((0,a)\) is obviously an injective *-homomorphism of
  \(A\) into \(\tilde{A}\), whose image is an ideal. We will usually suppress this map in
  our notation; this means that, for \(a\) in \(A\), we can identify \((0,a)\) and \(a\).
  Moreover, the quotient of \(\tilde{A}\) by A is obviously isomorphic to \(\mathbb{C}\).\par
  We turn now to the issue of a norm. In fact, it is quite easy to see that
  by defining
  \[\|(\lambda, a)\|_1 = |\lambda|+ \|a\| ,\]
  for \((\lambda,a)\) in \(\tilde{A}\), we obtain a Banach algebra with isometric involution. Getting
  the \(C^*\)-condition (which does not hold for this norm) is rather trickier.\par
  Toward a definition of a good \(C^*\)-norm, let us make a nice little observation. 
  If we recall only the fact that A is a Banach space, we can study \(\mathcal{B}(A)\),
  the algebra of bounded linear operators on it. Recalling now the product on
  A, the formula \(\pi(a)b = ab\), for a,b in A, defines \(\pi(a)\) as an operator on A.
  The Banach algebra inequality \(\|ab\| \leq \|a\| \|b\|\), simultaneously shows that
  \(\pi(a)\) is in \(\mathcal{B}(A)\) and \(\|\pi(a)\| \leq \|a\|\). It is easy to see that \(\pi : A \rightarrow \mathcal{B}(A)\) is a
  homomorphism. Finally, we consider
  \[\|\pi(a)a^*\| = \|aa^*\| = {\|a\|}^2 = \|a\| \|a^*\|,\]
  which implies that \(\|\pi(a)\| = \|a\|\). In short, the norm on A can be seen as
  the operator norm of \(A\) acting on itself.\par
  We are going to use a minor variation of that: since A is an ideal in \(\tilde{A}\),
  we can think of \(\tilde{A}\) as acting on A.\par
  We define a norm by
  \[\|(\lambda, a)\| = \sup\{|\lambda|, \|(\lambda,a)b\| , \|b(\lambda,a)\|\mid b \in A, \|b\| \leq 1\},\]
  for all \((\lambda,a)\) in \(\tilde{A}\). Notice that \((\lambda,a)b = \lambda b+ab\) is in A so the norm involved
  on the right is that of \(A\). It is easy to see the set on the right is bounded
  (by \(\|(\lambda,a)\|_1\)), so the supremum exists. We claim that \(\|(0,a)\| = \|a\|\), for
  all \(a\) in \(A\). The inequality \(\leq\) follows from the fact that the norm on A is a
  Banach algebra norm. Letting \(b = \|a\|^{-1} a^*\) yields the inequality \(\geq\) (at least
  for non-zero a).\par
  We next observe that \(\|(\lambda,a)\| = 0\) implies that \(|\lambda|=0\) and, given that \(\lambda=0\), we use \(\|(0,a)\|=\|a\|\) to
  conclude that \(a=0\). It is easy to see
  that * is isometric. Next, it is trivial that \(\|(\lambda,a)\| \leq |\lambda|+\|a\|\) and since both
  \(\mathbb{C}\) and \(A\) are complete, so is \(\tilde{A}\) in this norm\footnote{To this goal, we note that \(|\lambda|\leq \|(\lambda,a)\|,\qquad\|a\|=\|(0,a)\|\leq \|(\lambda,a)\|+\|(\lambda,0)\|\leq 2\|(\lambda,a)\|\),
  hence \(\|(\lambda,a)\|_1\leq 3\|(\lambda,a)\|\)}. We leave the details that the
  norm satisfies the Banach space conditions to the reader.\par
  This leaves us to verify the \(C^*\)-condition. Put \(x=(\lambda,a)\), and set
  \[
    R(x)=\sup_{\|b\|\leq1}\|xb\|,
    \qquad
    L(x)=\sup_{\|b\|\leq1}\|bx\|.
  \]
  These two quantities are equal. Indeed, let \(\{u_\mu\}\) be a contractive approximate unit for \(A\), whose existence is guaranteed by
  Theorem~\ref{Petersen:142}. If \(b\) belongs to the unit ball of \(A\), then
  \[
    \|bx\|=\lim_\mu\|bxu_\mu\|\leq R(x),
    \qquad
    \|xb\|=\lim_\mu\|u_\mu xb\|\leq L(x).
  \]
  Hence \(L(x)=R(x)\), and therefore \(\|x\|=\max\{|\lambda|,R(x)\}\). The inequality
  \(\|x^*x\|\leq\|x\|^2\) follows from submultiplicativity and the fact that the involution is isometric. Conversely,
  \(\|x^*x\|\geq|\lambda|^2\), and, for every \(b\) in the unit ball of \(A\),
  \begin{align*}
    \|x^*x\|&\geq\|x^*xb\|\\
    &\geq\|b^*x^*xb\|\\
    &=\|xb\|^2.
  \end{align*}
  Taking the supremum over \(b\) gives \(\|x^*x\|\geq R(x)^2\), so \(\|x^*x\|\geq\|x\|^2\). Thus \(\|x^*x\|=\|x\|^2\);
  compare \cite[Theorem 9.31 and Observation 9.32]{Marcoux2021}.\par
  The last item is to prove the uniqueness of \(\tilde{A}\). Suppose that B satisfies
  the desired conditions. We define a map from \(\tilde{A}\) as above to B by \(\rho(\lambda,a) =
  \lambda 1_B + a\), for all \(\lambda \in \mathbb{C}\) and \(a \in A\). It is a simple computation to see
  that \(\rho\) is a *-homomorphism. Let \(q : B \rightarrow B/A\) be the quotient map.
  Since A is an ideal and is not all of B, it cannot contain the unit of B.
  Thus \(q(1_B)\neq 0\). Let us prove that \(\rho\) is injective. If \(\rho(\lambda,a) = 0\), then
  \(0 = q\circ \rho(\lambda,a) = q(\lambda 1_B + a) = \lambda q(1_B)\) and we conclude that \(\lambda = 0\). It
  follows immediately that \(a = 0\) as well and so \(\rho\) is injective. It is clear that
  A is contained in the image of \(\rho\) and that \(q \circ \rho\) is surjective as well. Since
  \(B/A\) is one-dimensional, we conclude that \(\rho\) is surjective. The fact that \(\rho\) is
  isometric follows from Theorem~\ref{theorem:152}.
\end{proof}
The algebra \(\tilde{A}\) is called the \textnormal{unitization} of \(A\). We use the canonical maps
\[
  \iota_A:A\to\tilde{A},\qquad \iota_A(a)=(0,a),
\]
\[
  \pi_A:\tilde{A}\to\mathbb{C},\qquad \pi_A(\lambda,a)=\lambda,
  \qquad
  \sigma_A:\mathbb{C}\to\tilde{A},\qquad \sigma_A(\lambda)=(\lambda,0).
\]
They form the split exact sequence
\[
  \begin{tikzcd}
    0\arrow[r]&A\arrow[r,"\iota_A"]&\tilde{A}\arrow[r,shift left=2pt,"\pi_A"]
    &\mathbb{C}\arrow[l,shift left=2pt,"\sigma_A"]\arrow[r]&0,
  \end{tikzcd}
\]
where \(\pi_A\circ\sigma_A=\operatorname{id}_{\mathbb{C}}\). The associated
\textnormal{scalar map} is
\[
  s_A=\sigma_A\circ\pi_A:\tilde{A}\to\tilde{A},
  \qquad s_A(\lambda,a)=(\lambda,0).
\]
For each \(n\), we denote by \(s_{A,n}:M_n(\tilde{A})\to M_n(\tilde{A})\) the
entrywise amplification of \(s_A\). Thus its image is the canonical copy of
\(M_n(\mathbb{C})\) in \(M_n(\tilde{A})\), and
\(x-s_{A,n}(x)\in M_n(A)\) for every \(x\in M_n(\tilde{A})\). We shall often
omit the subscripts when the algebra and the matrix size are clear. Notice that
\(M_n(\tilde{A})\) should not in general be confused with
\(\widetilde{M_n(A)}\): the scalar parts are \(M_n(\mathbb{C})\) and
\(\mathbb{C}1_n\), respectively.\par

If \(\varphi:A\to B\) is a *-homomorphism, not necessarily unital, its
\textnormal{unitization} is the unital *-homomorphism
\[
  \tilde{\varphi}:\tilde{A}\to\tilde{B},
  \qquad \tilde{\varphi}(\lambda,a)=(\lambda,\varphi(a)).
\]
It satisfies
\[
  \pi_B\circ\tilde{\varphi}=\pi_A,
  \qquad
  \tilde{\varphi}\circ s_A=s_B\circ\tilde{\varphi}.
\]
Moreover, unitization is functorial:
\[
  \widetilde{\psi\circ\varphi}=\tilde{\psi}\circ\tilde{\varphi},
  \qquad
  \widetilde{\operatorname{id}_A}=\operatorname{id}_{\tilde{A}}.
\]
We use this external unitization also when \(A\) is already unital; in that
case \(\tilde{A}\cong\mathbb{C}\oplus A\), rather than \(\tilde{A}=A\).

\subsection{The commutative non-unital case}
We conclude with a few remarks concerning the spectral theorem in the
case of non-unital \cstars. Particularly, we would like generalizations of
Theorem~\ref{theorem:146} and Corollary~\ref{corollary:1410}. To generalize Theorem~\ref{theorem:146}, we would
simply like to drop the hypothesis that the algebra be unital.\par
Before beginning, let us make a small remark. If \(x_0\) is a point in some
compact Hausdorff space X, there is a natural isomorphism \(\{f \in C(X) | f(
x_0) = 0\} \cong  C_0(X \setminus \{x_0\})\) which simply restricts the function to \(X \setminus \{x_0\}\).
\begin{theorem}
  Let \(A\) be a commutative \cstar, let \(\tilde{A}\) be as in Theorem~\ref{theorem:181} and let \(\pi : \tilde{A} \rightarrow \mathbb C\)
  be the quotient map with kernel A. Then \(\pi\) is in
  \(\mathcal{M}(\tilde{A})\). Moreover, the restriction of the isomorphism of~\ref{theorem:146} to \(A\) is an
  isometric *-isomorphism between \(A\) and \(C_0(\mathcal{M}(\tilde{A}) \setminus \{\pi\})\).
\end{theorem}
\begin{proof}
  We give a sketch only. The first statement is clear. Secondly, if
  \(a\) is in \(A\), then \(\hat{a} (\pi) = \pi(a) = 0\), by definition of \(\pi\). So \(\hat{\cdot}\) maps A into
  \(C_0(\mathcal{M}(\tilde{A}) \setminus \{\pi\})\) (using the identification above) and is obviously an 
  isometric *-homomorphism. The fact that it is onto follows from the facts that
  \(A\) has codimension one in \(\tilde{A}\) while the same is true of \(C_0(\mathcal{M}(\tilde{A}) \setminus \{\pi\})\) in
  \(C(\mathcal{M}(\tilde{A}))\).
\end{proof}
\begin{corollary}[Gelfand]
  If \(A\) is a commutative \cstar\ then there exists a locally
  compact Hausdorff space \(X\) such that \(A \cong  C_0(X)\). Moreover, \(X\) is compact
  if and only if \(A\) is unital.
\end{corollary}
\begin{theorem}\label{theorem:184}
  Let \(a\) be a normal element of the unital \cstar\ \(B\).
  The isomorphism of~\ref{corollary:149} restricts to an isometric *-isomorphism between
  \(C_0(\operatorname{spec}(a) \setminus \{0\})\) and the \(C^*\)-subalgebra of \(B\) generated by \(a\).
\end{theorem}
\begin{proof}
  Again, we provide a sketch. First, assume that \(0\) is not in the spectrum
  of \(a\), so \(a\) is invertible. Notice that \(C_0(\operatorname{spec}(a)\setminus\{0\}) = C(\operatorname{spec}(a))\). We would
  like to simply apply \ref{corollary:149}, but we need to see that the \cstar\ generated
  by \(a\) and the unit is the same as that generated by \(a\) alone. Now, let \(f\) be
  the continuous function on \(\operatorname{spec}(a) \cup \{0\}\) which is \(1\) on \(\operatorname{spec}(a)\) and \(0\) at \(0\).
  Let \(\epsilon>0\) be arbitrary and use Weierstrass' Theorem to approximate \(f\) by a
  polynomial \(p(z, \overline{z})\) to within \(\epsilon\). Put
  \(q(z,\overline{z})=p(z,\overline{z})-p(0,0)\). Then \(q\) is a polynomial with no constant
  term which approximates \(f\) to within \(2\epsilon\). Hence \(q(a,a^*)\) is within
  \(2\epsilon\) of \(f(a)=1\). On the other hand, the formula of~\ref{corollary:149} shows
  that \(q(a,a^*)\) is in the \(C^*\)-subalgebra generated by \(a\) alone. In this case, the
  conclusion is exactly the same as for~\ref{corollary:149}.\par
  Now we assume that \(0\) is in the spectrum of \(a\). Let \(f\) be in \(C_0(\operatorname{spec}(a) \setminus \{0\})\). 
  We may view f as a continuous function of \(\operatorname{spec}(a)\) by defining \(f(0) = 0\).
  Again let \(\epsilon> 0\) be arbitrary and use Weierstrass' Theorem to approximate
  \(f\) by a polynomial \(p(z, \overline{z})\) to within \(\epsilon\). It follows that 
  \(p(z, \overline{z}) - p(0,0)\) is a
  polynomial with no constant term which approximates f to within \(2\epsilon\). Again,
  we see that \(f(a)\) can be approximated to within \(2\epsilon\) by an element of the \cstar\ generated
  by a alone. From this we see that \(C_0(\operatorname{spec}(a) \setminus \{0\})\) is
  mapped to the \(C^*\)-subalgebra generated by a. The rest of the conclusion is
  straightforward.
\end{proof}
\begin{exercise}
  If \(A\) is a unital \cstar, then \(\tilde{A}\cong \mathbb{C}\oplus A\) as \cstars\ (see Example~\ref{example:125}). 
  Write the isomorphism explicitly, since it isn't the
  obvious one.
\end{exercise}
\begin{exercise}
  Let \(X\) be a locally compact Hausdorff space and let \(A = C_0(X)\). 
  Prove that \(M(\tilde{A})\) is homeomorphic to \(X \cup \{\infty\}\), the one-point compactification of \(X\). 
  Discuss the overlap of this exercise and the last one.
\end{exercise}

\section{Ideals and quotients}
In this section, we consider ideals in a \cstar. Now ideals come in many
forms; they can be either one-sided or two-sided and they may or may not be
closed. We will concentrate here on closed, two-sided ideals. Generally, one
sided ideals are rather difficult to describe. (For a simple tractable case, see
Exercise~\ref{ex:192} below.) Ideals that are not closed can be even worse, although
in specific situations, there are ones of interest that arise. For an example,
the set of all finite rank operators in \(\mathcal{B}(\mathcal{H})\) is an ideal which is not closed. For
another example, consider a locally compact, non-compact Hausdorff space
\(X\). The set of all compactly supported continuous functions on \(X\) is an ideal
in \(C_0(X)\).\par
If \(A\) is a \cstar\ and \(I\) is a closed subspace, then we may form the
quotient space \(A/I\). We remark that the norm on the quotient \(A/I\) is defined
by
\[\|a +I\| = \inf\{\|a+b\| |b\in I\}.\]
For the proofs that this is a well-defined norm and makes \(A/I\) into a Banach space, 
we refer the reader to any standard text on functional analysis.\par
If, in addition, \(I\) is also a two-sided ideal, then \(A/I\) is also an algebra,
just as usual in a first course in ring theory and it is an easy matter to check
that our norm above is actually a Banach algebra norm. At this point, \(A/I\)
might fall short of being a \cstar\ on two points. The first is having a
*-operation. There is an obvious candidate, \({(a+I)}^* = a^*+I\), but to see this
is well-defined, we need to know that \(I\) is closed under *. The second point
is seeing that the norm satisfies \(C^*\)-condition. This all turns out to be true
and we state our main result.
\begin{lemma}
  Let \(a\) be an element in a \cstar\ \(A\). For any \(\epsilon > 0\), there
  is a continuous function \(f\) on \([0,\infty)\) such that \(e = f(a^*a)\) is in \(A\), is positive,
  \(\|e\| \leq 1\) and \(\|a-ae\| < \epsilon\).
\end{lemma}
\begin{proof}
  We note from Theorem \ref{theorem:165} that \(\operatorname{spec}(a^*a) \subseteq [0,\infty)\). Define the
  function f on \([0,\infty)\) by \(f(t) = t{(\epsilon^2 + t)}^{-1}\). Since \(f(0) = 0\), we may apply
  Theorem \ref{theorem:184} to see that \(f(a^*a)\) is contained in the \cstar\ generated
  by \(a^*a\) and hence in \(A\). It is clear that \(0 < f(t) < 1\), for all \(0 < t < \infty\)
  and so \(e = f(a^*a)\) is well-defined, positive and norm less than or equal to
  one. Notice in what follows, when we write expressions like \(1 - e\), this can
  be regarded as an element of \(\tilde{A}\), even if \(A\) is non-unital. We have
  \begin{align*}
  \|a -ae\|^2 &= \|a(1 -e)\|^2\\
  &= \|a(1 -f(a^*a))\|^2\\
  &= \|(1 -f(a^*a))a^*a(1-f(a^*a))\|\\
  &= \|g(a^*a)\|\\
  &\leq \|g\|_\infty,
  \end{align*}
  where \(g(t) = t{(1 - f(t))}^2\). It is a simple calculus exercise to check that \(g\)
  attains its maximum (on the positive axis) at \(t =\epsilon^2\) and its maximum is
  \(\epsilon^2/4\). We can conclude, observing that
  \[
    \|a-ae\|^2\leq\epsilon^2/4,\qquad\|a-ae\|\leq\epsilon/2<\epsilon.
  \]
\end{proof}
\begin{theorem}\label{theorem:191}
  Let \(A\) be a \cstar\ and suppose that \(I\) is a closed 
  (meaning as a topological subset), two-sided ideal. Then \(I\) is also closed under the
  *-operation and \(A/I\), with the quotient norm, is a \cstar.
\end{theorem} 
\begin{proof}
  Let \(a\) be in \(I\). We will show that \(a^*\) is also. Let \(\epsilon> 0\) and apply
  the last lemma to \(a\). Since \(a\) is in \(I\), so is \(a^*a\). We claim that \(e = f(a^*a)\) is
  in \(I\), also. It requires a separate argument to apply~\ref{theorem:184} since we do not yet
  know that \(I\) is a \cstar. However, we are applying the function to the
  self-adjoint element \(a^*a\). By approximating the function by polynomials in
  \(a^*a\) we see that the \cstar\ generated by this element is the same as the
  closed algebra generated by the element and hence is contained in \(I\).
  Since \(e\) is self-adjoint and the * operation is isometric, we have
  \[\|a^* - ea^*\| = \|a-ae\| \leq \epsilon.\]
  Since \(e\) is in \(I\), so is \(ea^*\). We see then that \(a^*\) is within \(\epsilon\) of an element of \(I\).
  As \(I\) is closed and \(\epsilon\) is arbitrary, we conclude that \(a^*\) is in \(I\).\par
  We now prove that the norm satisfies the \(C^*\)-condition.\par
  Let \(a\) be in \(A\). First, we show that the * operation is isometric since
  \begin{align*}
    \|a^* +I\| &= \inf\{\|a^*+b\| |b\in I\}\\
    = \inf\{\|a^* +b^*\| | b \in I\}\\
    = \inf\{\|a+b\| |b\in I\}\\
    = \|a+I\|.
  \end{align*}
  Now it suffices to prove that \(\|a+I\|^2 \leq \|a^*a+I\|\). We claim that
  \[\|a +I\| = \inf\{\|a+b\| |b\in I\}=\inf\{\|a(1-e)\| | e \in I,e \geq 0, \|e\| \leq 1\}.\]
  For the reverse, let \(b\) be in \(I\), let \(\epsilon> 0\) and use the Lemma to obtain \(e\) as
  desired with \(\|b - be\| <\epsilon\). Then we have
  \[\|a +b\| \geq \|(a+b)(1-e)\| \geq \|a(1-e)\| - \|b-be\| \geq \|a(1-e)\| - \epsilon.\]
  As \(\epsilon>0\) was arbitrary, we have established the reverse inequality.\\
  Having the claim, we now check that for a given \(e\) with conditions as
  above, we have
  \[\|a(1 -e) \|^2 = \|(1-e)a^*a(1-e)\| \leq \|a^*a(1-e)\| .\]
  Taking infimum over all such \(e\) yields the result.
\end{proof}

\begin{definition}
  A cone \(M\) in the positive part of a \cstar\ \(A\) is called \textnormal{hereditary} if 
  \(0 \leq x \leq y\), \(y \in M\), implies \(x \in M\) for each \(x\) in \(A\). A *-subalgebra \(B\) of \(A\) is
  \textnormal{hereditary} if \(B_+\) is hereditary in \(A_+\). Given a hereditary cone \(M\) in \(A_+\), define
  \[
    L(M)=\{x\in A| x^* x\in M\}.
  \]
\end{definition}

\begin{exercise}
  \begin{enumerate}
    \item Let \(X\) be a compact, Hausdorff space and let \(A = C(X)\). Let \(Z \subseteq X\) be closed and
    \[I =\{f \in C(X)|f(z) = 0, for all z \in Z\}.\]
    First prove that \(I\) is a closed, two-sided ideal in \(A\). Next, since \(I\) and
    \(A/I\) are both clearly commutative \cstars\ (although the former may
    not be unital) find \(M(I)\) and \(M(A/I)\).
    \item Let \(A\) be as above. Prove that any closed two-sided ideal \(I\) in \(A\) arises
    as above. (Hint: \(A/I\) is commutative and unital; use Exercise~\ref{ex:142} on
    the quotient map.)
  \end{enumerate}
\end{exercise}
\begin{exercise}\label{ex:192}
  Let \(n \geq 2\) and \(A = M_n(\mathbb{C})\).
  \begin{enumerate}
    \item Suppose that \(I\) is a right ideal in \(A\). Show that \(I\mathbb{C}^n = \{a\xi | a \in I,\xi \in \mathbb{C}^n\}\) 
      is a subspace of \(\mathbb{C}^n\). (Hint: if \(\xi,\eta\) are unit vectors and \(a\) is a
      matrix, then \(a\xi = a(\xi \otimes \eta^*)\eta\).)
    \item Prove the correspondence between the set of all right ideals in \(A\) and
      the set of all linear subspaces of \(\mathbb{C}^n\) in the previous part is a bijection.
    \item Prove that \(A\) is simple; that is, it has no closed two-sided ideals except
      \(0\) and \(A\).
  \end{enumerate}
\end{exercise}
\begin{exercise}
  Let \(\rho : A \rightarrow B\) be a *-homomorphism between two \cstars. 
  Prove that \(\rho(A)\) is closed and is a \(C^*\)-subalgebra of \(B\). (Hint:
  Theorem~\ref{theorem:191} and Theorem~\ref{theorem:152}.)
\end{exercise}
\section{Traces}
\begin{definition}
  Let \(A\) be a unital \cstar. A linear functional \(\phi\) on
  \(A\) is said to be positive if \(\phi(a^*a) \geq 0\), for all \(a\) in \(A\). A trace on \(A\) is a
  positive linear functional \(\tau : A \rightarrow \mathbb{C}\) with \(\tau(1) = 1\) satisfying
  \[\tau(ab) = \tau(ba),\]
  for all \(a,b\) in \(A\). This last condition is usually called the trace property. The
  trace is said to be faithful if \(\tau(a^*a) = 0\) occurs only for \(a = 0\).
\end{definition}
Notice that the trace property is satisfied by any homomorphism, but we
will see that it is strictly weaker. For the moment, it is convenient to regard
it as something stronger than mere linearity (or even positivity) but weaker
than being a homomorphism.
\begin{example}
  If \(A\) is commutative, every positive linear functional, \(\phi\),
  with \(\phi(1) = 1\), is a trace.
\end{example}
\begin{theorem}\label{theorem:1103}
  Let \(\mathcal{H}\) be a Hilbert space of (finite) dimension \(n\). If
  \(\{\xi_1,\ldots,\xi_n\}\) is an orthonormal basis for \(\mathcal{H}\), then
  \[
  \tau(a) =n^{-1}\sum^n_{i=1}<a\xi_i,\xi_i >,
  \]
  for any \(a\) in \(\mathcal{B}(\mathcal{H})\), defines a faithful trace on \(\mathcal{B}(\mathcal{H})\). The trace is unique. If
  we identify \(\mathcal{B}(\mathcal{H})\) with \(M_n(\mathbb{C})\), then this trace is expressed as
  \[
  \tau(a) =n^{-1}\sum_{i=1}^n a_{i,i}
  \]
  for any \(a\) in \(M_n(\mathbb{C})\).
\end{theorem}
\begin{proof}
  First, it is clear that \(\tau\) is a linear functional. Next, we see that, for
  any \(a\) in \(\mathcal{B}(\mathcal{H})\), we have
  \[
  \tau(a^*a)=n^{-1}\sum_{i=1}^n <a^*a\xi_i,\xi_i > = n^{-1}\sum_{i=1}^{n} < a\xi_i,a\xi_i >  
  = n^{-1} \sum_{i=1}^n \|a\xi_i\|^2.
  \]
  It follows at once that \(\tau\) is positive. Moreover, if \(\tau(a^*a) = 0\), then \(a\xi_i = 0\),
  for all \(i\). As \(a\) is zero on a basis, it is the zero operator. We also see from
  this, using \(a = 1\), that \(\tau(1) = 1\).\par
  It remains for us to verify the trace property. We do this by first considering 
  the case of rank one operators. We observe that, for any \(\xi,\eta\) in \(\mathcal{H}\), we
  have
  \begin{align*}
    \tau(\xi \otimes \eta^*) &= n^{-1}\sum_{i=1}^n <(\xi \otimes \eta^*)\xi_i,\xi_i >\\
    &= n^{-1}\sum_{i=1}^n<\xi_i,\eta >< \xi,\xi_i >\\
    &= n^{-1}<\sum_{i=1}^n<\xi,\xi_i>\xi_i,\eta>\\
    &=n^{-1}<\xi,\eta>
  \end{align*}
  Now let \(\xi,\eta,\zeta,\omega\) be in \(\mathcal{H}\) and \(a = \xi \otimes\eta^*,b = \zeta \otimes\omega^*\). 
  We have \(ab =< \zeta,\eta> \xi \otimes\omega^*\) 
  and \(ba =< \xi,\omega >\zeta \otimes\eta^*\). It follows then from the computation above
  that
  \begin{align*}
    \tau(ab) &= <\zeta,\eta >\tau(\xi\otimes\omega^*)\\
    &= n^{-1} <\zeta,\eta ><\xi,\omega >\\
    &= <\xi,\omega>\tau(\zeta\otimes\eta^*)\\
    &= \tau(ba).
  \end{align*}
  Since the linear span of such operators is all of \(\mathcal{B}(\mathcal{H})\), the trace property holds.\par
  Hence we turn to the uniqueness of the trace. Suppose that \(\phi\) is any
  trace on \(\mathcal{B}(\mathcal{H})\). We will show that \(\phi\) and \(\tau\) agree on \(\xi_i \otimes\xi_j^*\), 
  first considering
  the case \(i\neq j\). Let \(a = \xi_i \otimes\xi^*_i\) and \(b = \xi_i \otimes\xi^*_j\). It follows that \(ab = b\) while
  \(ba = 0\). Therefore, we have \(\phi(b) = \phi(ab) = \phi(ba) = 0 = \tau(b)\). Next, for any
  \(i, j\), let \(\nu = \xi_i \otimes\xi^*_j\). Then \(\nu^*\nu = \xi_j \otimes\xi^*_j\) 
  while \(\nu\nu^* = \xi_i \otimes\xi^*_i\). Therefore, we
  have \(\phi(\xi_i \otimes\xi^*_i) = \phi(\xi_j \otimes\xi^*_j)\). Finally, we have
  \[
  1 =\phi(1) = \phi \left( \sum_{i=1}^{n}\xi_i \otimes \xi_i^* \right)= n\phi(\xi_1 \otimes \xi^*_1).
  \]
  Putting these together, we see that
  \[
  \phi(\xi_i \otimes \xi^*_i) = \phi(\xi_1 \otimes \xi^*_1) = n^{-1} = \tau(\xi_i \otimes\xi^*_i).
  \]
  We have shown that \(\tau\) and \(\phi\) agree on a spanning set, \(\xi_i \otimes \xi^*_j,1 \leq i,j \leq n\),
  hence they are equal.
\end{proof}
\begin{theorem}
  Let \(\tau\) be the unique trace on \(\mathcal{B}(\mathcal{H})\), where \(\mathcal{H}\) is a finite
  dimensional Hilbert space. If \(p\) is a projection, then \(\dim(p\mathcal{H}) = \tau(p)\dim(\mathcal{H})\).
\end{theorem}
\begin{proof}
  Choose an orthonormal basis \(\{\xi_1,\ldots,\xi_n\}\) for \(\mathcal{H}\) in such a way that
  \(\{\xi_1,\ldots,\xi_k\}\) is a basis for \(p\mathcal{H}\). We omit the remaining computation.
\end{proof}
\begin{exercise}
  Prove that, for \(n > 1\), there is no non-zero
  *-homomorphism from \(M_n(\mathbb{C})\) to \(\mathbb{C}\). In fact, give two different proofs, one
  using~\ref{theorem:1103} above, and one using~\ref{ex:192}.
\end{exercise}

\chapter{Functionals and representations}\label{Petersen:chapter3}
\section{Positive functionals}
\begin{definition}
  Let \(A^*\) denote the dual space of a \cstar\ \(A\) (this notation should not lead 
  to confusion with the set of adjoints as long as it is applied only to self-adjoint sets).
  For each \(\phi\) in \(A^*\), define \(\phi^*\) by \(\phi^*(x)=\overline{\phi(x^*)}\).
  We say that \(\phi\) is \textnormal{self-adjoint} if \(\phi=\phi^*\). This is clearly equivalent
  to the condition that \(\phi(A_{sa})\subset\mathbb{R}\).
\end{definition}
The expression 
\[
  \phi=\frac{1}{2}(\phi+\phi^*)+\frac{1}{2}(\phi-\phi^*)
\]
shows that each \(\phi\) in \(A^*\) has a unique decomposition into self-adjoint and skew-adjoint parts.
With \((A^*)_{sa}\) as the real Banach space of self-adjoint elements of \(A^*\), we note that
\((A^*)_{sa}\) is isometrically isomorphic with the dual of the real Banach space \(A_{sa}\), i.e., \((A^*)_{sa}=(A_{sa})^*\).
\begin{definition}
  A linear functional \(\phi\) on \(A\) is called \textnormal{positive} (in symbols, \(\phi\geq0\))
  if \(\phi(A_+)\subset\mathbb{R}_+\).
\end{definition}
A positive functional is self-adjoint. If \(\phi\) is positive, then for any sequence \(\{x_n\}\) in \(A_+^1\),\footnote{\(A^1\) denotes the unit ball in \(A\).} we have
\[
  \sum_{k=1}^n 2^{-k}\phi(x_k)\leq\phi\left(\sum_{k=1}^\infty 2^{-k}x_k\right),
\]
which shows that a positive functional is necessarily bounded on \(A_+^1\)\footnote{This can be proved by contradiction: if \(\phi\) were unbounded, we could take \(x_k\in A_+^1\) with \(\phi(x_k)>2^k\); for \(x=\sum 2^{-k}x_k\in A_+\) we would have \(\phi(x)>n\) for every \(n\), which is impossible.} and thus belongs to \(A^*\).
The following version of the Cauchy-Schwarz inequality provides the link between abstract \cstars\ 
and concrete algebras of operators on Hilbert spaces.
\begin{theorem}\label{Petersen:313}
  If \(\phi\) is a positive functional on a \cstar\ \(A\), then for all \(x,y\) in \(A\),
  \[
    |\phi(y^* x)|^2\leq \phi(y^* y)\phi(x^* x).
  \]
\end{theorem}
\begin{proof}
  For each complex \(\lambda\), we have \(\phi((\lambda x+y)^*(\lambda x+y))\geq 0\).
  With \(\lambda=t\phi(x^* y)|\phi(y^* x)|^{-1}\)\footnote{If \(\phi(y^*x)=0\), the theorem follows immediately.} (note that \(\overline{\phi(y^* x)}=\phi(x^* y)\)) and \(t\) in \(\mathbb{R}\), this gives
  \[
    t^2\phi(x^* x)+2t|\phi(y^* x)|+\phi(y^* y)\geq 0,
  \]
  from which the theorem follows.
\end{proof}
\begin{proposition}\label{Petersen:314}
  An element \(\phi\) in \(A^*\) is positive if and only if \(\lim\phi(u_\lambda)=\|\phi\|\)
  for some approximate unit \(\{u_\lambda\}\) in \(A\).
\end{proposition}
\begin{proof}
  If \(\phi\geq0\), then \(\{\phi(u_\lambda)\}\) is an increasing net in \(\mathbb{R}_+\) with limit \(\alpha\leq\|\phi\|\).
  For each \(x\) in \(A^1\), we have, by \ref{Petersen:313},
  \[
    |\phi(u_\lambda x)|^2\leq\phi(u_\lambda^2)\phi(x^* x)\leq\phi(u_\lambda)\|\phi\|\leq\alpha\|\phi\|.
  \]
  Since \(\{u_\lambda\}\) is an approximate unit for \(A\) and \(\phi\) is continuous, this implies that
  \(|\phi(x)|^2\leq\alpha\|\phi\|\), whence \(\|\phi\|^2\leq\alpha\|\phi\|\) and \(\alpha=\|\phi\|\).\par
  To prove the converse, suppose that \(\{\phi(u_\lambda)\}\) is convergent to \(\|\phi\|\).
  Take \(x\) in \(A_{sa}^1\) and write \(\phi(x)=\alpha+i\beta\) with \(\alpha,\beta\) in \(\mathbb{R}\).
  Adjusting the sign of \(x\), we may assume that \(\beta\geq 0\). Fix \(n\geq1\). Since
  \(\|xu_\lambda-u_\lambda x\|\to0\), there is \(\lambda_0\) such that
  \(\|xu_\lambda-u_\lambda x\|<1/n\) whenever \(\lambda\geq\lambda_0\). For every such \(\lambda\),
  \[
    \|nu_\lambda-ix\|^2=\|n^2u_\lambda^2+x^2+in(xu_\lambda-u_\lambda x)\|\leq n^2+2.
  \]
  However
  \[
    \lim_{\lambda}|\phi(nu_\lambda-ix)|^2=(n\|\phi\|+\beta)^2+\alpha^2.
  \]
  This implies that
  \[
    (n\|\phi\|+\beta)^2+\alpha^2\leq(n^2+2)\|\phi\|^2
  \]
  for all \(n\), whence \(\beta=0\). Therefore \(\phi=\phi^*\).\par
  If \(x\in A^1_+\), then \(u_\lambda-x\in A_{sa}^1\), so that \(\phi(u_\lambda-x)\leq\|\phi\|\).
  Passing to the limit, this implies that \(\phi(x)\geq 0\), whence \(\phi\geq 0\), and the proof is complete.
\end{proof}
\begin{lemma}\label{Petersen:315}
  Let \(\tilde{A}\) be the \cstar\ obtained by adjoining a unit to the \cstar\ \(A\).
  For each positive functional \(\phi\) on \(A\), define an extension \(\tilde{\phi}\) on \(\tilde{A}\) by setting
  \(\tilde{\phi}(1)=\|\phi\|\). Then \(\tilde{\phi}\) is positive on \(\tilde{A}\), and \(\|\tilde{\phi}\|=\|\phi\|\).
\end{lemma}
\begin{proof}
  If \(\{u_\lambda\}\) is an approximate unit for \(A\), then for each \(x\) in \(A\) and complex \(\lambda\),
  \begin{align*}
    \lim\sup\|\lambda u_\lambda+x\|^2&=\lim\sup\||\lambda|^2u_\lambda^2+\overline{\lambda}u_\lambda x+\lambda x^* u_\lambda+x^* x\|\\
    &\leq\lim\sup\||\lambda|^2 1+\overline{\lambda}u_\lambda x+\lambda x^* u_\lambda+x^* x\|=\|\lambda 1+x\|^2
  \end{align*}
  using \ref{Petersen:135}. It follows from \ref{Petersen:314} that
  \[
    |\tilde{\phi}(\lambda 1+x)|=\lim|\phi(\lambda u_\lambda+x)|\leq\|\lambda1+x\|\|\phi\|,
  \]
  so that \(\|\tilde{\phi}\|=\|\phi\|\). Since \(\tilde{\phi}(1)=\|\tilde{\phi}\|\), we have \(\tilde{\phi}\geq0\) by \ref{Petersen:314}.
\end{proof}
\begin{proposition}\label{Petersen:316}
  Let \(B\) be a \(C^*\)-subalgebra of \(A\). For each positive functional \(\phi\) on \(B\),
  there is a norm-preserving extension of \(\phi\) to a positive functional on \(A\). If \(B\) is 
  hereditary, then this extension is unique.
\end{proposition}
\begin{proof}
  Adjoin a unit \(1\) to \(A\) and let \(\tilde{A}\) (resp., \(\tilde{B}\)) be the \cstar\ generated by \(1\) and \(A\)
  (resp., \(B\)). Whether or not \(B\) is unital, the argument in \ref{Petersen:315} shows that the definition \(\tilde{\phi}(1)=\|\phi\|\)
  gives a positive norm-preserving extension of \(\phi\) to \(\tilde{B}\).
  By the Hahn-Banach theorem we may then extend \(\tilde{\phi}\) to a functional \(\psi\) on \(\tilde{A}\)
  such that \(\|\psi\|=\|\tilde{\phi}\|\). Since \(\|\psi\|=\psi(1)\), \(\psi\geq0\) by \ref{Petersen:314}. The restriction of \(\psi\) to
  \(A\) gives the desired extension of \(\phi\).\par
  Assume now that \(B\) is hereditary and let \(\{u_\lambda\}\) be an approximate unit for \(B\).
  Then \(\|\psi\|=\lim\phi(u_\lambda)\) by \ref{Petersen:314}, so that \(\psi(1-u_\lambda)\rightarrow0\).
  Since \(u_\lambda\tilde{A}u_\lambda\subset B\), it follows from \ref{Petersen:313} that
  \[
    \psi(x)=\lim\psi(u_\lambda x u_\lambda)=\lim\phi(u_\lambda x u_\lambda)
  \]
  for every \(x\) in \(\tilde{A}\).
\end{proof}

\section{Representations}
The study of \cstars\ is motivated by the prime example of closed * algebras 
of operators on Hilbert space. With this in mind, it is natural to
find ways that a given abstract \cstar\ may act as operators on Hilbert
space. Such an object is called a representation of the \cstar.
\begin{definition}
  Let \(A\) be a *-algebra. A \textnormal{representation} of \(A\) is a pair,
  \((\pi, \mathcal{H})\), where \(\mathcal{H}\) is a Hilbert space and \(\pi : A \rightarrow \mathcal{B}(\mathcal{H})\) 
  is a *-homomorphism.
  We also say that \(\pi\) is a representation of \(A\) on \(\mathcal{H}\).
\end{definition}
\begin{definition}
  Let \(A\) be a \cstar. Two representations of \(A\),
  \((\pi_1, \mathcal{H}_1),(\pi_2,\mathcal{H}_2)\), are \textnormal{unitarily equivalent} if there is a unitary operator \(u :
  \mathcal{H}_1 \rightarrow \mathcal{H}_2\) such that \(u\pi_1(a) = \pi_2(a)u\), for all a in \(A\). In this case, we write
  \((\pi_1, \mathcal{H}_1) \thicksim_u (\pi_2,\mathcal{H}_2)\) or \(\pi_1 \thicksim_u \pi_2\).
\end{definition}
\begin{definition}
  Let \(A\) be a \cstar\ and \((\pi_\iota,\mathcal{H}_\iota),\iota \in I\), be a collection
  of representations of \(A\). Their \textnormal{direct sum} is \((\bigoplus_{\iota\in I}\pi_\iota,\bigoplus_{\iota\in I}\mathcal{H}_\iota)\), 
  where \(\bigoplus_{\iota\in I}\mathcal{H}_\iota\)
  consists of tuples, \(\xi = (\xi_\iota)_{\iota\in I}\) satisfying \(\sum_{\iota\in I} \|\xi_\iota\|^2 < \infty\) and
  \[
  {\left(\bigoplus_{\iota\in I}\pi_\iota(a)\xi\right)}_\iota = \pi_\iota(a)\xi_\iota,\iota \in I.
  \]
\end{definition}
\begin{definition}
  Let \(A\) be a *-algebra and let \((\pi,\mathcal{H})\) be a representation
  of \(A\). A subspace \(\mathcal{N} \subset \mathcal{H}\) is said to be \textnormal{invariant} if \(\pi(a)\mathcal{N} \subset \mathcal{N}\), 
  for all \(a\) in \(A\).
\end{definition}
\begin{proposition}\label{prop:1115}
  Let \(A\) be a *-algebra and let \((\pi,\mathcal{H})\) be a representation
  of \(A\). A closed subspace \(\mathcal{N}\) is invariant if and only if \(\mathcal{N}^\bot\) is.
\end{proposition}
\begin{proof}
  For the ``only if'' direction, it suffices to consider \(\xi\) in \(\mathcal{N}^\bot\) and \(a\) in \(A\)
  and show that \(\pi(a)\xi\) is again in \(\mathcal{N}^\bot\). To this end, let \(\eta\) be in \(\mathcal{N}\). We have
  \[<\pi(a)\xi,\eta >=< \xi,{\pi(a)}^*\eta >=< \xi,\pi(a^*)\eta >= 0,\]
  since \(\pi(a^*)\mathcal{N} \subset \mathcal{N}\).\par
  The ``if'' direction follows since \({(\mathcal{N}^\bot)}^\bot = \mathcal{N}\).
\end{proof}
In the case of the Proposition above, it is possible to define two 
representations of \(A\) by simply restricting the operators to either \(\mathcal{N}\) or \(\mathcal{N}^\bot\). That is,
we define
\[\pi|_{\mathcal{N}}(a) = \pi(a)|_\mathcal{N},a \in A.\]
Moreover, the direct sum of these two representations is unitarily equivalent
to the original. That is, we have
\[(\pi, \mathcal{H}) \thicksim_u (\pi|_\mathcal{N},\mathcal{N})\oplus(\pi|_{\mathcal{N}^\bot},\mathcal{N}^\bot).\]
One can, in some sense, consider the notion of reducing a representation to
an invariant subspace and its complement as an inverse to taking direct sums.
\begin{definition}
  A representation, \((\pi,\mathcal{H})\), of a *-algebra, \(A\), is \textnormal{non-degenerate} 
  if the only vector \(\xi\) in \(\mathcal{H}\) such that \(\pi(a)\xi = 0\) for all \(a\) in \(A\),
  is \(\xi = 0\). Otherwise, the representation is \textnormal{degenerate}.
\end{definition}
\begin{proposition}
  A representation \((\pi,\mathcal{H})\) of a unital *-algebra is 
  non-degenerate if and only if \(\pi(1) = 1\).
\end{proposition}
\begin{theorem}
  Every representation of a *-algebra is the direct sum of a
  non-degenerate representation and the zero representation (on some Hilbert
  space).
\end{theorem}
\begin{definition}
  Let \((\pi,\mathcal{H})\) be a representation of a *-algebra \(A\). We say
  that a vector \(\xi\) in \(\mathcal{H}\) is \textnormal{cyclic} if the linear space \(\pi(A)\xi\) is dense in \(\mathcal{H}\). We
  say that the representation is cyclic if it has a cyclic vector.
\end{definition}
Notice that every cyclic representation is non-degenerate as follows. Let
\(\xi\) be the cyclic vector. If \(\pi(a)\eta = 0\), for all \(a\) in \(A\), then we have
\[<\pi(a)\xi,\eta >=< \xi,\pi(a^*)\eta >=< \xi,0 >= 0.\]
As \(\pi(a)\xi\) is dense in \(\mathcal{H}\), we conclude that \(\eta = 0\).
\begin{definition}
  A representation of a *-algebra is \textnormal{irreducible} if the only
  closed invariant subspaces are \(0\) and \(\mathcal{H}\). It is \textnormal{reducible} otherwise.
\end{definition}
\begin{proposition}\label{prop:11111}
  A non-degenerate representation of a *-algebra is irreducible if and only if every non-zero vector is cyclic.
\end{proposition}
\begin{proof}
  First assume that \((\pi,\mathcal{H})\) is irreducible. Let \(\xi\) be a non-zero vector,
  then \(\pi(A)\xi\) is evidently an invariant subspace and its closure is a closed
  invariant subspace. If it is \(0\), then the representation is degenerate, which is
  impossible. Otherwise, it must be \(\mathcal{H}\), meaning that \(\xi\) is a cyclic vector for \(\pi\).
  Conversely, suppose that \(\pi\) is non-degenerate, but reducible. Let \(\mathcal{N}\) be
  a closed invariant subspace which is neither \(0\) nor \(\mathcal{H}\). If \(\xi\) is any non-zero
  vector in \(\mathcal{N}\), then \(\pi(A)\xi\) is clearly contained in \(\mathcal{N}\) and hence cannot be dense
  in \(\mathcal{H}\). Hence, we have found a non-zero vector which is not cyclic.
\end{proof}
\begin{proposition}\label{prop:111112}
  A non-degenerate representation of a *-algebra is irreducible if and only if the only positive operators which commute with its
  image are scalars.
\end{proposition}
\begin{proof}
  Let \((\pi,\mathcal{H})\) be a representation of the *-algebra \(A\).\par
  First, we suppose that there is a non-trivial closed invariant subspace, \(\mathcal{N}\).
  Let \(p\) be the orthogonal projection onto \(\mathcal{N}\). That is, \(p\xi = \xi\), for all \(\xi\) in \(\mathcal{N}\)
  and \(p\xi = 0\), for all \(\xi\) in \(\mathcal{N}^\bot\). It is easy to check that \(p = p^* = p^2\), which
  means that \(p\) is positive. Moreover, as both \(\mathcal{N}\) and \(\mathcal{N}^\bot\) are non-empty, this
  operator is not a scalar. We check that it commutes with \(\pi(a)\), for any \(a\) in
  \(A\). If \(\xi\) is in \(\mathcal{N}\), we know that \(\pi(a)\xi\) is also and so
  \[
  (p\pi(a))\xi = p(\pi(a)\xi) = \pi(a)\xi = \pi(a)(p\xi) = (\pi(a)p)\xi.
  \]
  On the other hand, if \(\xi\) is in \(\mathcal{N}^\bot\), then so is \(\pi(a)\xi\) and
  \[
  (p\pi(a))\xi = p(\pi(a)\xi) = 0 = \pi(a)(0) = \pi(a)(p\xi) = (\pi(a)p)\xi.
  \]
  Since every vector in \(\mathcal{H}\) is the sum of two as above, we see that \(p\pi(a) = \pi(a)p\).\par
  Conversely, suppose that \(h\) is some positive, non-scalar operator on \(\mathcal{H}\),
  which commutes with every element of \(\pi(A)\). If the spectrum of \(h\) consists
  of a single point, then it follows from~\ref{corollary:149} that \(h\) is a scalar. As this is not
  the case, the spectrum consists of at least two points. We may then find
  non-zero continuous functions \(f,g\) on \(\operatorname{spec}(h)\) whose product is zero. Since
  \(f\) is non-zero on the spectrum of \(h\), the operator \(f(h)\) is non-zero. Let \(\mathcal{N}\)
  denote the closure of its range, which is a non-zero subspace of \(\mathcal{H}\). On the
  other hand, \(g(h)\) is also a non-zero operator, but it is zero on the range of
  \(f(h)\) and hence on \(\mathcal{N}\). This implies that \(\mathcal{N}\) is a proper subspace of \(\mathcal{H}\).\par
  Next, we claim that, as \(h\) commutes with \(\pi(A)\), so does \(f(h)\). Let \(a\) be
  in \(A\). For any \(\epsilon> 0\), we may find a polynomial \(p(x)\) such that \(\|p-f\|_\infty <\epsilon\) 
  in \(C(\operatorname{spec}(h))\) and this means that \(\|p(h) - f(h)\| < \epsilon\). On the other hand,
  it is clear that \(p(h)\) will commute with \(\pi(a)\), since h does. Finally, we claim
  that \(\mathcal{N}\) is invariant under \(\pi(a)\). In fact, it suffices to check that the range of
  \(f(h)\) is invariant. But if \(\xi\) is in \(\mathcal{H}\), we have
  \[\pi(a)(f(h)\xi) = \pi(a)f(h)\xi = f(h)\pi(a)\xi \in f(h)\mathcal{H},\]
  and we are done.
\end{proof}

\section{Multipliers}\label{section:centralizers}
\begin{definition}
  Let \(A\) be a nondegenerate \(C^*\)-subalgebra of \(\mathcal{B}(\mathcal{H})\).
  We say that an operator \(x\) in \(\mathcal{B}(\mathcal{H})\) is a \textnormal{left} (respectively, \textnormal{right})
   \textnormal{multiplier} for \(A\) if \(xA\subset A\) (respectively, \(Ax\subset A\)).
   We say that \(x\) is a \textnormal{two-sided multiplier} (or just a \textnormal{multiplier}) if \(x\) is both 
   a left and right multiplier. We say that \(x\) is a \textnormal{quasimultiplier} if \(AxA\subset A\).
\end{definition}

If \(1\in A\), then there are, of course, no left, right, or quasimultipliers outside of \(A\).
Since \(xA\subset A\) implies \(x A''\subset A''\) and \(1\in A''\) (\(A\) is nondegenerate), this implies that all multipliers of \(A\) belong to \(A''\).\par
To show that the multipliers of \(A\) do not depend on the particular representation in \(\mathcal{B}(\mathcal{H})\), we introduce the 
notion of centralizers as follows.

\begin{definition}
  A linear map \(\rho:A\rightarrow A\) is called a \textnormal{left} (respectively, \textnormal{right})
   \textnormal{centralizer} if \(\rho(xy)=\rho(x)y\) (respectively \(\rho(xy)=x\rho(y)\)) for all \(x\) and \(y\)
   in \(A\). A \textnormal{double centralizer} is a pair \((\rho_1,\rho_2)\) consisting of a right and a left centralizer such that
   \(\rho_1(x)y=x\rho_2(y)\) for all \(x\) and \(y\) in \(A\). A \textnormal{quasicentralizer} is a bilinear map \(\rho:A\times A\rightarrow A\)
   such that, for each fixed \(x\) in \(A\), the map \(\rho(x,\cdot)\) is a left centralizer, and the map \(\rho(\cdot,x)\) is a right centralizer. 
\end{definition}
\begin{lemma}\label{Petersen:3122}
  Each left or right centralizer, and each quasicentralizer of a \cstar\ \(A\) is bounded.
\end{lemma}
\begin{proof}
  Let \(\rho\) be a left centralizer of \(A\). If \(\rho\) was not bounded, then we could find
  \(x_n\) in \(A\) with \(\|x_n\|<1/n\) but \(\|\rho(x_n)\|>n\).
  Let \(a=\sum x_n x_n^*\). Then by \ref{Petersen:145} for each \(n\), there is an element \(u_n\)
  in \(A\) with \(\|u_n\|\leq\|a^{1/6}\|\) such that \(x_n=a^{1/3}u_n\), but then
  \[
    \|\rho(x_n)\|=\|\rho(a^{1/3})u_n\|\leq\|\rho(a^{1/3})\|\|a^{1/6}\|,
  \]
  a contradiction. Thus every left centralizer is bounded. If \(\rho\) is a right centralizer, then
  \(\rho^\sharp(x)=\rho(x^*)^*\) is a left centralizer; hence \(\rho^\sharp\), and therefore \(\rho\), is bounded.
  If \(\rho(\cdot,\cdot)\) is a quasicentralizer of \(A\), then for each 
  \(x\) in \(A^1\), the left centralizer \(\rho(x,\cdot)\) is a bounded map, say \(\|\rho(x,\cdot)\|=\lambda_x\), from the above.
  However, for a fixed \(y\) in \(A\), the right centralizer \(\rho(\cdot,y)\) is also bounded.
  It follows from the uniform boundedness theorem that the set \(\{\lambda_x| x\in A^1\}\) is bounded.
  Consequently, \(\rho(\cdot,\cdot)\) is a bounded map on \(A\times A\).
\end{proof}
\begin{proposition}\label{Petersen:3123}
  Let \(A\) be a nondegenerate \(C^*\)-subalgebra of \(\mathcal{B}(\mathcal{H})\).
  Then there is a bijective correspondence between left (respectively right, double, and quasi-)
  centralizers of \(A\) and left (respectively right, two-sided, and quasi-) multipliers of \(A\) in \(A''\).
\end{proposition}
\begin{proof}
  If \(\rho\) is a left centralizer of \(A\) and \(\{u_\lambda\}\) is an approximate unit for \(A\),
  then \(\{\rho(u_\lambda)\}\) is a bounded net by \ref{Petersen:3122} and has therefore a weak-* limit point \(x\) in \(A''\) by Banach-Alaoglu.
  However, for each \(y\) in \(A\), the net \(\{u_\lambda y\}\) is norm convergent to \(y\), and so, again by \ref{Petersen:3122},
  \(\{\rho(u_\lambda y)\}\) is norm convergent to \(\rho(y)\). Since \(\rho(u_\lambda y)=\rho(u_\lambda)y\),
  this implies that \(xy=\rho(y)\). It follows that \(x\) is a left multiplier corresponding to \(\rho\).
  If \(z\) was another element in \(A''\) such that \(zy=\rho(y)\), then we would have \((x-z)y=0\)
  for all \(y\) in \(A\) and hence also for all \(y\) in \(A''\), and thus \(x=z\).
  This shows that the correspondence between left centralizers and left multipliers is a bijection.\par
  If \((\rho_1,\rho_2)\) is a double centralizer of \(A\), then by the preceding there is a unique pair \((x_1,x_2)\) corresponding to
  \((\rho_1,\rho_2)\) consisting of a right and a left multiplier of \(A\). However, for all \(y\) and \(z\) in \(A\),
  \[
    yx_1 z=\rho_1(y)z=y\rho_2(z)=yx_2 z,
  \]
  whence \(x_1=x_2\).\par
  If \(\rho(\cdot,\cdot)\) is a quasicentralizer, then let \(x\) be a weak limit point in \(A''\) of the
  net \((\{\rho(u_\lambda,u_\lambda)\})\), where, as before, \(\{u_\lambda\}\) is an approximate unit for \(A\). Then for
  all \(y\) and \(z\) in \(A\), we have
  \[
    \rho(y,z)=\lim\rho(yu_\lambda,u_\lambda z)=\lim y\rho(u_\lambda,u_\lambda)z=yxz,
  \]
  which proves that \(x\) is a quasimultiplier for \(A\) in \(A''\) corresponding to \(\rho(\cdot,\cdot)\).
  The uniqueness of \(x\) is verified as above, and the proof is complete.
\end{proof}

\begin{exercise}
  Fix a positive integer \(n\) and consider \(A = M_n(\mathbb{C})\). For
  \(K \geq 1\), let
  \[
  \mathcal{H}_K =\bigoplus^K_{k=1} \mathbb{C}^n.
  \]
  Define \(\pi_K\) by
  \[\pi_K(a)(\xi_1,\xi_2,\ldots,\xi_K) = (a\xi_1,a\xi_2,\ldots,a\xi_K).\]
  This is a non-degenerate representation of \(A\) on \(\mathcal{H}_K\).
  \begin{enumerate}
    \item Give a necessary and sufficient condition on a vector \((\xi_1,\xi_2,\ldots,\xi_K)\)
    for it to be a cyclic vector for \(\pi_K\).
    \item For which values of \(K\) is the representation \(\pi_K\) cyclic?
    \item For which values of \(K\) is the representation \(\pi_K\) irreducible? (You should
    be able to give two proofs; one using the definition and one applying
    Proposition~\ref{prop:11111}.)
  \end{enumerate}
\end{exercise}
\begin{exercise}
  Fix \(n \geq 1\) and let \(A = M_n(\mathbb{C})\). There is an obvious
  representation which we call \(\rho\) of \(A\) on \(\mathbb{C}^n\):
  \[\rho(a)\xi = a\xi,a \in A,\xi \in \mathbb{C}^n.\]
  We are using the usual multiplication of matrices (or matrix and vector).
  Now let \((\pi,\mathcal{H})\) be any non-degenerate representation of \(A\). (If you like,
  you may start by assuming that \(\mathcal{H}\) is finite dimensional, although it makes
  little difference.)
  \begin{enumerate}
    \item Use the fact that \(A\) is simple (Exercise~\ref{ex:192}) to show that \(\pi(e_{1,1})\neq 0\)
    \item Let \(\xi\) be any unit vector in \(\pi(e_{1,1})\mathcal{H}\). Prove that
    \[\xi, \pi(e_{2,1})\xi, \ldots , \pi(e_{n,1})\xi\]
    is an orthonormal set.
    \item Prove that
    \[\pi(A)\xi = span\{\xi,\pi(e_{2,1})\xi,\ldots,\pi(e_{n,1})\xi\}.\]
    \item Letting \(\mathcal{H}_\xi\) denote the subspace of the last part, prove that \(\pi|_{\mathcal{H}_\xi}\) 
    is unitarily equivalent to \(\rho\).
    \item Let \(\mathcal{B}\) be an orthonormal basis for \(\pi(e_{1,1})\mathcal{H}\). Prove that, for any \(\xi\neq \eta\)
    in \(\mathcal{B}\), \(\mathcal{H}_\xi\) and \(\mathcal{H}_\eta\) are orthogonal.
    \item Prove that \(\bigoplus_{\xi\in\mathcal{B}}\mathcal{H}_\xi = \mathcal{H}\).
    \item Prove that \((\pi,\mathcal{H})\) is unitarily equivalent to \(\bigoplus_{\xi\in\mathcal{B}}(\rho,\mathbb{C}^n)\).
  \end{enumerate}
\end{exercise}
\begin{exercise}\label{ex:1113}
  Prove that if \((\pi,\mathcal{H})\) is an irreducible representation of a
  commutative \cstar, then \(\mathcal{H}\) is one-dimensional.
\end{exercise}
\begin{exercise}
  Let \(A = C[0,1]\). Consider the representation
  \((\pi, L^2([0,1],\lambda))\), where \(\lambda\) is Lebesgue measure and \((\pi(f)\xi)(x) = f(x)\xi(x)\),
  for \(f\) in \(A\), \(\xi\) in \(L^2([0,1],\lambda)\) and \(x\) in \([0,1]\). (If you like, you may assume
  that \(\xi\) actually lies in \(C[0,1]\), although we regard it as a vector in \(L^2([0,1],\lambda)\).
  This makes things somewhat easier.)
  \begin{enumerate}
    \item Give a necessary and sufficient condition on a vector \(\xi\) in \(L^2([0,1],\lambda)\)
    to be a cyclic vector.
    \item Prove that \(\pi\) is faithful; that is, it is injective.
    \item Prove that this representation is not the direct sum of a collection of
    irreducible representations.
    \item Find a countable collection of irreducible representation of \(A\) whose
    direct sum is injective.
  \end{enumerate}
\end{exercise}
\begin{exercise}
  Suppose \((\pi,\mathcal{H})\) is a finite-dimensional irreducible representation of a \cstar\ \(A\). 
  Prove that \(\pi(A) = \mathcal{B}(\mathcal{H})\). (Hint: Theorem~\ref{theorem:171},
  Theorem~\ref{theorem:179} and Proposition~\ref{prop:111112})
\end{exercise}

\section{The GNS construction}
\begin{definition}
  Let \(A\) be a \cstar. A linear functional \(\phi\) on \(A\) is
  \textnormal{positive} if \(\phi(a^*a) \geq 0\), for all \(a\) in \(A\). The linear
  functional \(\phi\) is a \textnormal{state} if it is positive and, in addition, has norm \(1\) 
  (for a unital \cstar, this is equivalent to say that \(\phi(1)=1\)).\par
  The set of states of \(A\) is denoted by \(S(A)\). 
\end{definition}

From \ref{Petersen:314} we see that the norm is an additive function on the cone of positive functionals,
which implies that \(S(A)\) is a convex set.\par
The self-adjoint part of the unit ball of \(A^*\) can be identified with the unit ball of \((A_{sa})^*\) and is therefore convex
and weak-* compact. The positive part of the unit ball of \((A_{sa})^*\), which we denote by \(Q(A)\), is convex and weak-* closed,
hence weak-* compact. It is easy to see that \(Q(A)\) is the convex span of \(S(A)\) and the zero functional.\par
If \(1\in A\), then by \ref{Petersen:314}
\[
  S(A)=\{\phi\in A^*| \|\phi\|\leq1, \phi(1)=1\},
\]
which shows that \(S(A)\) is weak-* closed in \(Q(A)\) and thus weak-* compact.
In general, this is not true, and that is the reason why we have to use the space \(Q(A)\)
(the quasistates) although our real interest lies with the state spaces \(S(A)\).

\begin{lemma}\label{Petersen:322}
  The unit ball of \((A_{sa})^*\) is the convex span of \(S(A)\) and \(-S(A)\).
\end{lemma}
\begin{proof}
  Let \(K\) be the convex span of \(Q(A)\) and \(-Q(A)\). Then \(K\) is also the convex span of \(S(A)\)
  and \(-S(A)\) and is a weak-* compact subset of the unit ball of \((A_{sa})^*\) (since \(Q(A)\) and \(-Q(A)\) are both weak-* compact).\par
  If \(x\in A_{sa}\) and \(\lambda\in \textnormal{spec}(x)\setminus\{0\}\), then \(\lambda=\omega(x)\) for some nonzero
  homomorphism \(\omega\) on the \(C^*\)-subalgebra \(B\) of \(A\) generated by \(x\).
  We may regard \(\omega\) as a state of \(B\), and then by \ref{Petersen:316} there is a state \(\phi\) of \(A\) that extends \(\omega\).
  This shows that 
  \[
    \textnormal{spec}(x)\setminus\{0\}\subset\{\phi(x)| \phi\in S(A)\},
  \]
  and thus, in particular,
  \begin{equation}\label{Petersen:eq322}
    \|x\|=\sup\{|\phi(x)|| \phi\in S(A)\}.
  \end{equation}\par
  Suppose now that \(\psi\) was an element in the unit ball of \((A_{sa})^*\) not belonging to \(K\).
  By Hahn-Banach there is an element \(x\) in \(A_{sa}\) and \(\alpha\) in \(\mathbb{R}\)
  such that \(\psi(x)>\alpha\) but \(\phi(x)\leq\alpha\) for all \(\phi\) in \(K\). Since \(K\) is symmetric, this implies
  that \(|\phi(x)|\leq\alpha\) for all \(\phi\) in \(K\), whence \(\|x\|\leq\alpha\) by (\ref{Petersen:eq322}).
  This contradicts \(\psi(x)>\alpha\), and thus \(K\) equals the unit ball of \((A_{sa})^*\).
\end{proof}

\begin{lemma}\label{Petersen:323}
  Let \(\phi\) and \(\psi\) be positive functionals on a unital \cstar\ \(A\). The following conditions are equivalent:
  \begin{enumerate}
    \item \(\|\phi-\psi\|=\|\phi\|+\|\psi\|\).
    \item For every \(\epsilon>0\), there is a \(z\) in \(A_+^1\) such that \(\phi(1-z)<\epsilon\) and \(\psi(z)<\epsilon\).
  \end{enumerate}
\end{lemma}
\begin{proof}
  \begin{itemize}
    \item \((1)\Rightarrow(2)\). Since \(\phi-\psi\in(A_{sa})^*\), there is an element \(x\) in \(A_{sa}^1\)
      such that
      \[
        \phi(x)-\psi(x)+\epsilon\geq\|\phi-\psi\|,
      \]
      but then 
      \[
        \phi(x)-\psi(x)+\epsilon\geq\|\phi\|+\|\psi\|=\phi(1)+\psi(1),
      \]
      whence \(\phi(1-x)+\psi(1+x)<\epsilon\). Since \(0\leq1-x\leq2\) and \(0\leq1+x\leq2\),
      we can choose \(z=\frac{1}{2}(1+x)\), so that \(1-z=\frac{1}{2}(1-x)\).
    \item \((2)\Rightarrow(1)\). Clearly, \(\|\phi-\psi\|\leq\|\phi\|+\|\psi\|\), but if \(\phi(1-z)<\epsilon\) and \(\psi(z)<\epsilon\), then
      \begin{align*}
        \|\phi\|+\|\psi\|&=\phi(1)+\psi(1)\leq\phi(2z-1)+\psi(1-2z)+4\epsilon\\
        &=(\phi-\psi)(2z-1)+4\epsilon\leq\|\phi-\psi\|+4\epsilon,
  \end{align*}
  since \(\|2z-1\|\leq1\). As \(\epsilon\) is arbitrary, \(\|\phi\|+\|\psi\|\leq\|\phi-\psi\|\).
  \end{itemize}
\end{proof}

\begin{definition}
  We say that two positive functionals \(\phi\) and \(\psi\) satisfying conditions in \ref{Petersen:323} are \textnormal{orthogonal}
  and write \(\phi\perp\psi\). 
\end{definition}
If \(A\) has no unit, then we must replace condition (2) with the existence of \(x\) and \(y\) in \(A_+\)
with \(\|x+y\|\leq1\) such that \(\phi(x)>\|\phi\|-\epsilon\) and \(\psi(y)>\|\psi\|-\epsilon\). We shall not need this version.

\begin{theorem}\label{Petersen:325}
  For each self-adjoint functional \(\phi\) on a \cstar\ \(A\), there is a unique pair 
  \(\phi_+\) and \(\phi_-\) of positive functionals such that \(\phi=\phi_+-\phi_-\) and \(\phi_+\perp\phi_-\). 
\end{theorem}
\begin{proof}
  By \ref{Petersen:322} we have \(\phi=\alpha\phi_1-(1-\alpha)\phi_2\) with \(\phi_1\) and \(\phi_2\) positive and of norm less or equal to \(\left\|\phi\right\|\)
  and \(0\le \alpha \le 1\).
  Put \(\phi_+=\alpha\phi_1\) and \(\phi_-=(1-\alpha)\phi_2\) and observe that
  \[
    \|\phi_+\|+\|\phi_-\|=\alpha\|\phi_1\|+(1-\alpha)\|\phi_2\|\leq\|\phi\|,
  \]
  so that \(\phi_+\perp\phi_-\).\par
  To prove the uniqueness of the decomposition, assume that 
  \(\phi_1,\phi_2,\psi_1,\psi_2\) are positive functionals on \(A\) with
  \begin{equation}\label{Petersen:eq325}
    \phi_1\perp\psi_1,\qquad\phi_2\perp\psi_2\textnormal{, and}\qquad \phi_1-\psi_1=\phi_2-\psi_2.
  \end{equation} 
  Since \(\phi_1+\psi_2=\psi_1+\phi_2\), we may adjoin a unit to \(A\) and extend the functionals as in \ref{Petersen:315}
  without destroying the relations (\ref{Petersen:eq325}). Thus we may assume that \(A\) has a unit.
  Using \ref{Petersen:323}, we have for each \(\epsilon>0\) that \(\phi_1(1-z)<\epsilon\) and \(\psi_1(z)<\epsilon\)
  with \(0\leq z\leq1\). Then
  \[
    \phi_2(z)\geq\phi_2(z)-\psi_2(z)=\phi_1(z)-\psi_1(z)>\phi_1(1)-2\epsilon.
  \]
  Likewise, \(\psi_2(1-z)>\psi_1(1)-2\epsilon\), so that
  \[
    \phi_2(z)+\psi_2(1-z)>\|\phi_1\|+\|\psi_1\|-4\epsilon=\|\phi_2\|+\|\psi_2\|-4\epsilon.
  \]
  It follows that \(\phi_2(1-z)+\psi_2(z)<4\epsilon\). Since \(\phi_1-\phi_2=\psi_1-\psi_2\), we have, for each \(x\) in \(A\),
  \begin{align*}
    \phi_1(x)-\phi_2(x)&=\phi_1(xz)-\phi_2(xz)+\phi_1(x(1-z))-\phi_2(x(1-z))\\
    &=\psi_1(xz)-\psi_2(xz)+\phi_1(x(1-z))-\phi_2(x(1-z)).
  \end{align*}
  We have
  \begin{align*}
    |\phi_1(x)-\phi_2(x)|\leq\|x\|\|\phi_1-\psi_1\|^{1/2}(\epsilon^{1/2}+(4\epsilon)^{1/2}+\epsilon^{1/2}+(4\epsilon)^{1/2})\\
    &\leq\|x\|\|\phi_1-\psi_1\|^{1/2}6\epsilon^{1/2}.
  \end{align*}
  Since \(\epsilon\) is arbitrary, \(\phi_1=\phi_2\), and thus \(\psi_1=\psi_2\).
\end{proof}
\begin{corollary}\label{Petersen:326}
  Each element in the unit sphere of \((A_{sa})^*\)  has a unique representation in \(\textnormal{Conv}(S(A)\cup(-S(A)))\).
\end{corollary}

\begin{proposition}\label{prop:1122}
  Let \(\phi\) be a linear functional on a unital \cstar\ \(A\)
  with \(\phi(1) = \|\phi\| =1\). Then \(\phi\) is a state.
\end{proposition}
\begin{proof}
  We begin by showing that if \(a\) is self-adjoint, then \(\phi(a)\) is real. Suppose
  that \(\Im(\phi(a))\neq 0\) and without loss of generality, assume it is positive.
  Then select \(0 < r < \Im(\phi(a)) \leq \|a\|\). The function \(f(s) = \sqrt{{\|a\|}^2+s^2}-s\) is
  clearly greater than \(r\) at \(s = 0\). By multiplying and dividing by 
  \(\sqrt{{\|a\|}^2+s^2}+ s\) and simplifying, we see that \(f(s)\) tends to \(0\) as \(s\) tends to infinity. Thus we
  can find \(s\) such that \(f(s) = r\). A simple geometric argument shows that the
  disc with centre \(-is\) and radius \(s + r\) then contains the spectrum of \(a\), but
  does not contain \(\phi(a)\). That is, \(\|a - is\| \leq r +s\) while \(|\phi(a) - is| > r +s\).
  But now we can compute
  \[|\phi(a - is)| = |\phi(a) - is\phi(1)| = |\phi(a) - is| > r +s,\]
  since \(\phi(1) = 1\), and this contradicts \(\|\phi\| = 1\).\par
  Now we let \(a\) be positive and show that \(\phi(a) \geq 0\). From Lemma~\ref{lemma:162},
  we know that \(\|\|a\|1 -a\| \leq \|a\|\). Applying \(\phi\) to \(\|a\|1-a\) and using the facts
  that \(\phi(1) = 1 = \|\phi\|\), we have
  \[|\|a\| -\phi(a)| \leq \|a\|,\]
  which implies that \(\phi(a) \geq 0\). We have shown that \(\phi\) is positive, as desired.
\end{proof}
\begin{lemma}
  Let \(\phi\) be a positive linear functional on a unital \cstar\ \(A\).
  \begin{enumerate}
    \item For all \(a,b\) in \(A\), we have
      \[|\phi(b^*a)|^2 \leq \phi(a^*a)\phi(b^*b).\]
    \item For all \(a\) in \(A\), we have \(\phi(a^*) = \overline{\phi(a)}\).
    \item \(\phi(1) = \|\phi\|\).
    \item For all \(a,b\) in \(A\), we have
      \[\phi(b^*a^*ab) \leq {\left\|a\right\|}^2\phi(b^*b).\]
  \end{enumerate}
\end{lemma}
\begin{proof}
  Let \(a,b\) be in \(A\). For any complex number \(\lambda\), we know that
  \[0 \leq\phi({(\lambda a+b)}^*(\lambda a+b)) = |\lambda|^2\phi(a^*a)+ \overline{\lambda}\phi(a^*b) + \lambda \phi(b^*a) + \phi(b^*b).\]
  The first and last terms are obviously positive, so the sum of the middle two
  is real. Using \(b = 1\) and \(\lambda = i\) shows that the imaginary parts of \(\phi(a)\) and
  \(\phi(a^*)\) are opposite, while \(b = 1\) and \(\lambda = 1\) shows the real parts of \(\phi(a)\) and
  \(\phi(a^*)\) are equal. The second part follows.\par
  Since the identity element of \(A\) has norm one, we have \(|\phi(1)| \leq \|\phi\|\) and
  as \(1^*1 = 1\), \(\phi(1)\) is positive. From the equation above, and the second part
  applied to \(b^*a\), we see that
  \[-2 \Re(\lambda\phi(b^*a)) \leq |\lambda|^2\phi(a^*a) + \phi(b^*b).\]
  If \(\phi(a^*a) = 0\), the fact that this holds for arbitrary \(\lambda\) means that \(\phi(b^*a) = 0\)
  also. For \(\phi(a^*a)\ne 0\), choose \(z \in \mathbb{C}\) such that \(|z| = 1, z\phi(b^*a) = |\phi(b^*a)|\) and
  using \(\lambda = -z{\phi(a^*a)}^{-1/2}{\phi(b^*b)}^{1/2}\) leads us to
  \[|\phi(b^*a)| \leq {\phi(a^*a)}^{1/2}{\phi(b^*b)}^{1/2},\]
  and the proof of the first statement is complete.\par
  For the third part, \(\|1\| = 1\), so that \(|\phi(1)| \leq \|\phi\|\). Using \(b = 1\) in part 1,
  we have
  \[|\phi(a)| \leq {\phi(a^*a)}^{1/2}\phi(1)^{1/2} \leq {\|\phi\|}^{1/2} {\|a^*a\|}^{1/2}{\phi(1)}^{1/2} = \|\phi\| ^{1/2} \|a\|{\phi(1)}^{1/2} .\]
  Taking supremum over all \(a\) of norm one, we get
  \[\|\phi\| \leq {\|\phi\|}^{1/2}{\phi(1)}^{1/2}\]
  and it follows that \(\|\phi\| \leq \phi(1)\). This completes the proof of part 3.\par
  For the last part, we consider the functional \(\psi(c) = \phi(b^*cb)\), for any
  \(c\) in \(A\). This is clearly another positive linear functional and so we have
  \(\|\psi\| =\psi(1)=\phi(b^*b)\). Then we have
  \[\phi(b^*a^*ab) = \psi(a^*a) \leq \|\psi\| \|a^*a\| = \phi(b^*b) \|a\|^2.\]
\end{proof}
\begin{theorem}\label{Petersen:333}
  Let \(A\) be a unital \cstar\ and let \(\phi\) be a state on \(A\).
\begin{enumerate}
  \item The set
    \[N_\phi =\{a\in A|\phi(a^*a) = 0\}\]
    is a closed left ideal in \(A\).
    \item The sesquilinear form \(< a+N_\phi,b+N_\phi >= \phi(b^*a)\) is well-defined and non
    degenerate on \(A/N_\phi\) and the completion of \(A/N_\phi\) is a Hilbert space,
    denoted \(\mathcal{H}_\phi\).
    \item The formula
    \[\pi_\phi(a)(b + N_\phi) = ab+N_\phi,\]
    for \(a,b\) in \(A\) extends to define \(\pi_\phi(a)\) as a bounded linear operator on
    \(\mathcal{H}_\phi\).
    \item The function \(\pi_\phi\) is a representation of \(A\) on \(\mathcal{H}_\phi\).
    \item The vector \(\xi_\phi = 1 + N_\phi \in A/N_\phi \subset \mathcal{H}_\phi\) is a cyclic vector for \(\pi_\phi\) with
    norm one.
\end{enumerate}
\end{theorem}
\begin{proof}
  The fact that \(N_\phi\) is closed follows easily from the fact that \(\phi\) is continuous. 
  Next, suppose that \(a\) is in \(A\) and \(b\) is in \(N_\phi\). The fact that \(ab\) is in
  \(N_\phi\) follows immediately from the last part of the last Lemma.\par
  For the second part, it follows from part one of the last Lemma that
  \(\phi(b^*a) = 0\) if either \(a\) or \(b\) is in \(N_\phi\) and so the sesquilinear form is well defined.
  It is also clearly positive definite.\par
  The fact that \(N_\phi\) is a left ideal means that, for any \(a\) in \(A\), the formula
  given for \(\pi_\phi(a)\) yields a well-defined linear transformation on \(A/N_\phi\). The fact
  that it is continuous (and hence extends to a bounded linear operator on \(\mathcal{H}_\phi\))
  follows immediately from the last part of the last Lemma.\par
  It is clear that \(\pi_\phi\) is linear and multiplicative. Let us check it preserves
  adjoints. For \(a,b,c\) in \(A\), we have
  \begin{align*}
    <\pi_\phi(a^*)b +N_\phi,c+N_\phi > &= <a^*b+N_\phi,c+N_\phi >\\
    &= \phi(c^*(a^*b)) = \phi({(ac)}^*b)\\
    &= <b+N_\phi,ac+N_\phi >\\
    &= <b+N_\phi,\pi_\phi(a)c+N_\phi >\\
    &= <{\pi_\phi(a)}^*b+N_\phi,c+N_\phi >.
  \end{align*}
  Since this holds for arbitrary \(b\) and \(c\) in \(A\), we conclude that \(\pi_\phi(a^*) = {\pi_\phi(a)}^*\).\par
  For the last part, if \(b\) is any element of \(A\), it is clear that
  \[\pi_\phi(b)\xi_\phi = b \cdot 1 +N_\phi = b+N_\phi.\]
  It follows then that \(\pi_\phi(A)\xi_\phi\) contains \(A/N_\phi\) and is therefore dense in \(\mathcal{H}_\phi\).\par
  Finally, we compute
  \[\|\xi_\phi\| = \|1+N_\phi\| ={\phi(1^*1)}^{1/2} ={\phi(1)}^{1/2} = 1.\]
\end{proof}
\begin{definition}
  Let \(A\) be a unital \cstar\ and let \(\phi\) be a state on \(A\).
  The triple \((\mathcal{H}_\phi,\pi_\phi,\xi_\phi)\) is called the \textnormal{GNS representation} of \(\phi\).
\end{definition}
\begin{theorem}
  Let \(A\) be a unital \cstar\ and suppose that \(\pi\) is a
  representation of \(A\) on the Hilbert space \(\mathcal{H}\) with cyclic vector, \(\xi\), of norm
  one. Then
  \[\phi(a) =< \pi(a)\xi,\xi >,\]
  for all \(a\) in \(A\), defines a state on \(A\). Moreover, the GNS representation of
  \(\phi\) is unitarily equivalent to \(\pi\) in the sense that there is a unitary operator
  \(u : \mathcal{H}\rightarrow \mathcal{H}_\phi\) satisfying \(u\pi(a)u^* = \pi_\phi(a)\), for all \(a\) in \(A\) and \(u\xi= \xi_\phi\).
\end{theorem}
We have shown a correspondence between states and cyclic representa
tions. Under this, the irreducible representations are characterized by a nice
geometric property, namely that they are extreme points among the set of
states. Such points are often called \textit{pure states}.
\begin{theorem}\label{theorem:1127}
  Let \(\phi\) be a state on the unital \cstar\ \(A\). The GNS
  representation \((\pi_\phi,\mathcal{H}_\phi)\) is irreducible if and only if \(\phi\) is not a non-trivial
  convex combination of two other states. That is, if there are states \(\phi_0\) and
  \(\phi_1\) and \(0 < t <1\) such that \(\phi = t\phi_0 +(1 -t)\phi_1\), then \(\phi_0 = \phi_1 = \phi\).
\end{theorem}
\begin{proof}
  First, we suppose that \(\phi = t\phi_0+(1-t)\phi_1\) with \(\phi_0\neq \phi_1\) and \(0 < t < 1\)
  and show that the GNS representation is not irreducible. Define a sesquilinear
  form on \(A/N_\phi\) by \((a+N_\phi,b+N_\phi) = t\phi_0(b^*a)\). We have
  \begin{align*}
    |(a + N_\phi,b+N_\phi)|^2 &= t^2|\phi_0(b^*a)|^2\\
    &\leq |t\phi_0(a^*a)||t\phi_0(b^*b)|\\
    &\leq |\phi(a^*a)||\phi(b^*b)|\\
    &= \|a+N_\phi\|^2 \|b+N_\phi\|^2.
  \end{align*}
  It follows that our sesquilinear form is well-defined on the quotient \(A/N_\phi\) and
  also extends to one on \(\mathcal{H}_\phi\). Moreover, there exists a unique bounded positive
  linear operator \(h\) on \(\mathcal{H}_\phi\) such that
  \[t\phi_0(b^*a) = (a + N_\phi,b+N_\phi) =< h(a+N_\phi),b+N_\phi > .\]
  We will show that \(h\) is not a scalar multiple of the identity and commutes
  with \(\pi_\phi(a)\), for all \(a\) in \(A\). It will then follow from Proposition~\ref{prop:111112} that
  \(\pi_\phi\) is not irreducible.\par
  If \(h\) is a multiple of the identity, then using \(b = 1\) in the formula above,
  we see that \(\phi\) is a multiple of \(\phi_0\). As both are states and take the same value
  at the unit, they are then equal. Since \(\phi = t\phi_0 + {(1 - t)}\phi_1\), we also have
  \(\phi_1 = \phi\), and this is a contradiction.\par
  Now, let \(a,b,c\) be in \(A\). We have
  \begin{align*}
    <\pi_\phi(a)h(b +N_\phi),c+N_\phi > &= <h(b+N_\phi),{\pi_\phi(a)}^*(c+N_\phi) >\\
    &= <h(b+N_\phi),a^*c+N_\phi >\\
    &= t\phi_0({(a^*c)}^*b)\\
    &= t\phi_0(c^*ab)\\
    &= <h(ab+N_\phi),c+N_\phi >\\
    &= <h\pi_\phi(a)(b+N_\phi),c+N_\phi > .
  \end{align*}
  As this equality holds on a dense set of pairs of vectors \(b + N_\phi,c + N_\phi\), we
  conclude that \(\pi_\phi(a)h = h\pi_\phi(a)\). This completes the proof.\par
  In the other direction, suppose that \(\pi_\phi\) is reducible. That is, we may find
  a non-trivial invariant subspace \(\mathcal{N} \subset \mathcal{H}_\phi\). Write \(\xi_\phi = \xi_0 + \xi_1\), where \(\xi_0\) is in
  \(\mathcal{N}\) and \(\xi_1\) is orthogonal to \(\mathcal{N}\). We claim that neither is zero, for if \(\xi_1 = 0\),
  then \(\xi_\phi = \xi_0\) is in \(\mathcal{N}\). As this subspace is invariant, we see that \(\pi_\phi(A)\xi_\phi\) is
  contained in \(\mathcal{N}\). On the other hand, since \(\xi_\phi\) is cyclic, this subspace must be
  dense in \(\mathcal{H}_\phi\) which is a contradiction. The other case is similar.\par
  Define
  \[\phi_i(a) = \|\xi_i\|^{-2} < \pi_\phi(a)\xi_i,\xi_i >,\]
  for \(a\) in \(A\) and \(i = 0,1\). Then \(\phi_0\) and \(\phi_1\) are states on \(A\). We also claim that
  \(\phi = \|\xi_0\|^2\phi_0+\|\xi_1\|^2\phi_1\). To see this, we have
  \begin{align*}
    \phi(a) =& <\pi_\phi(a)\xi_\phi,\xi_\phi >\\
    =& <\pi_\phi(a)(\xi_0 +\xi_1),(\xi_0 +\xi_1) >\\
    =& <\pi_\phi(a)\xi_0,\xi_0 > + < \pi_\phi(a)\xi_1,\xi_0 >+<\pi_\phi(a)\xi_0,\xi_1 > + < \pi_\phi(a)\xi_1,\xi_1 >\\
    =& \|\xi_0\|^2\phi_0(a)+ <\pi_\phi(a)\xi_1,\xi_0 >+<\pi_\phi(a)\xi_0,\xi_1 > +\|\xi_1\|^2\phi_1(a).
  \end{align*}
  The two central terms are both zero since \(\mathcal{N}\) and \(\mathcal{N}^\perp \) are both invariant
  under \(\pi_\phi(a)\) and \(\xi_0\) is in the former and \(\xi_1\) is in the latter.\par
  We claim that \(\phi_0\neq \phi_1\). First, let \(C = \min\{ \|\xi_0\| , \|\xi_1\| \} > 0\). Since the
  vector \(\xi_\phi\) is cyclic, we may find \(a\) in \(A\) with \(\|\pi_\phi(a)\xi_\phi -\xi_0\| < 2^{-1}C\). Writing
  \(\pi_\phi(a)\xi_\phi = \pi_\phi(a)\xi_0+\pi_\phi(a)\xi_1\) and noting the first vector is in \(\mathcal{N}\) and the second
  is in \(\mathcal{N}^\perp\), we see that
  \[\|\pi_\phi(a)\xi_0 - \xi_0\| , \|\pi_\phi(a)\xi_1\| < 2^{-1}C.\]
  Then we have
  \begin{align*}
    |\phi_0(a) - 1| &\leq | \|\xi_0\|^{-2} < \pi_\phi(a)\xi_0,\xi_0 > -\|\xi_0\|^{-2} < \xi_0,\xi_0 > |\\
    &\leq \|\xi_0\|^{-2} \|\pi_\phi(a)\xi_0-\xi_0\| \|\xi_0\|\\
    &< \frac{C}{2\|\xi_0\|}\\
    &\leq 2^{-1}.
  \end{align*}
  We also have
  \begin{align*}
    |\phi_1(a)| &\leq | \|\xi_1\|^{-2} < \pi_\phi(a)\xi_1,\xi_1 > |\\
    &\leq \|\xi_1\|^{-2} \|\pi_\phi(a)\xi_1\| \|\xi_1\|\\
    &<  \frac{C}{{2\|\xi_1\|}}\\
    &\leq 2^{-1}.
  \end{align*}
  We conclude that \(\phi_0(a)\neq \phi_1(a)\) and we are done.
\end{proof}
\begin{lemma}\label{lemma:1129}
  Let \(A\) be a unital \cstar\ and let \(a\) be a self-adjoint
  element of \(A\). There exists an irreducible representation \(\pi\) of \(A\) such that
  \(\|\pi(a)\| = \|a\| \).
\end{lemma}
\begin{proof}
  Let \(B\) be the \(C^*\)-subalgebra of \(A\) generated by \(a\) and the unit. It is
  unital and commutative. We know that \(\mathcal{M}(B)\) is homeomorphic to \(\operatorname{spec}(a)\)
  via \(\phi \rightarrow \phi(a)\) and that \(\operatorname{spec}(a)\) is a compact subset of the real numbers.
  Choose \(\phi_0\) in \(\mathcal{M}(B)\) such that \(|\phi_0(a)| = \sup\{|x| | x \in \operatorname{spec}(a)\} = \|a\|\), and put
  \(\lambda_0=\phi_0(a)\). Of course, \(\phi_0(1)=1\) and \(\lambda_0\in\{-\|a\|,\|a\|\}\).\par
  By the Hahn-Banach Theorem, \(\phi_0\) has an extension \(\phi\) to \(A\) such that \(\|\phi\|=\|\phi_0\|=1\). Since
  \(\phi(1)=\phi_0(1)=1\), Proposition~\ref{prop:1122} shows that \(\phi\) is a state. Consequently, the set
  \[
    F=\{\psi\in S(A)\mid\psi(a)=\lambda_0\}
  \]
  is non-empty. It is convex and weak-* closed in the weak-* compact state space \(S(A)\), and hence it is weak-* compact. We claim that
  \(F\) is a face of \(S(A)\). Suppose that \(\psi=t\psi_1+(1-t)\psi_2\) belongs to \(F\), where \(0<t<1\) and
  \(\psi_1,\psi_2\in S(A)\). If \(\lambda_0=\|a\|\), then \(\psi_j(a)\leq\|a\|\) for \(j=1,2\), so equality in
  \[
    \|a\|=t\psi_1(a)+(1-t)\psi_2(a)
  \]
  forces \(\psi_1(a)=\psi_2(a)=\|a\|\). The case \(\lambda_0=-\|a\|\) is analogous, using \(\psi_j(a)\geq-\|a\|\).
  Thus \(\psi_1,\psi_2\in F\), proving the claim.\par
  By the Krein-Milman theorem, \(F\) has an extreme point \(\psi\). Since \(F\) is a face of \(S(A)\), \(\psi\) is also an extreme point
  of \(S(A)\), hence a pure state. Theorem~\ref{theorem:1127} therefore implies that its GNS representation \(\pi_\psi\) is irreducible;
  see also \cite[Section 7.2]{NachtergaeleSims2016}. Finally,
  \[
    \|\pi_\psi(a)\|\geq
    |\langle\pi_\psi(a)\xi_\psi,\xi_\psi\rangle|
    =|\psi(a)|=|\lambda_0|=\|a\|.
  \]
  The reverse inequality follows because \(\pi_\psi\) is contractive, and hence \(\|\pi_\psi(a)\|=\|a\|\).
\end{proof}
\begin{theorem}\label{theorem:1128}
  Let \(A\) be a unital \cstar. Then there exists a Hilbert space
  \(\mathcal{H}\) and \(C^*\)-subalgebra \(B \subset \mathcal{B}(\mathcal{H})\) which is *-isomorphic to \(A\).
\end{theorem}
\begin{proof}
  For each \(a\) in
  the unit ball of \(A\), we find an irreducible representation \(\pi_a\) of \(A\) such that
  \(\|\pi_a(a^*a)\| = \|a^*a\| \). Of course, it follows that \(\|\pi_a(a)\| = \|a\| \), as well. We
  finally take the direct sum over all \(a\) in the unit ball of \(A\) of the representations
  \(\pi_a\). It follows that this representation is isometric and we are done by letting
  \(B\) denote its range.
\end{proof} 
\begin{remark}
  The previous theorem can be proved also for non-unital \cstars\ via unitization.
\end{remark}
\begin{exercise}
  Let \(X\) be a compact, Hausdorff space and let \(A = C(X)\).
  The Riesz representation theorem states that every linear functional \(\phi\) on \(A\)
  is given as
  \[\phi(f) = \int_X f(x)d\mu(x),\]
  for some finite Borel measure \(\mu\) on \(X\).
  \begin{enumerate}
    \item Give a necessary and sufficient condition on \(\mu\) as above for the associated \(\phi\) to be a state.
    \item Prove that \(N_\phi\) is a closed two-sided ideal in \(A\).
    \item Describe \(\mathcal{H}_\phi\) in terms of \(\mu\).
    \item Prove that \(N_\phi\) is also the kernel of \(\pi_\phi\).
  \end{enumerate}
\end{exercise}
\begin{exercise}
  Let \(A = M_N(\mathbb{C})\) and let \(\tau\) be the trace of Theorem~\ref{theorem:1103}.
  Let \(\pi\) be the obvious representation of \(A\) on \(\mathbb{C}^N: \pi(a)\xi = a\xi\). Describe the
  relationship between \(\pi\) and \(\pi_\tau\). (More explicitly, describe \(\pi_\tau\) in terms of \(\pi\),
  up to unitary equivalence. Hint: let \(\{\xi_1,\xi_2,\ldots,\xi_N\}\) be an orthonormal basis
  for \(\mathbb{C}^N\). First observing that \(\tau\) is faithful, what is \(N_\tau\)? Then use the fact
  that \(\{\xi_i \otimes \xi^*_j | 1 \leq i,j \leq N\}\) is a linear basis for \(A\). Can you find invariant
  subspaces for \(\pi_\tau\)?)
\end{exercise}

\chapter{Group \texorpdfstring{\(C^*\)-algebras}{C*algebras}}
\section{Group representations}
\begin{definition}
  Let \(G\) be a discrete group. A \textnormal{unitary representation} of
  \(G\) is a pair \((u,\mathcal{H})\), where \(\mathcal{H}\) is a complex Hilbert space and \(u\) is a group 
  homomorphism, which is continuous in the weak topology, from \(G\) to the group of unitary operators on \(\mathcal{H}\), usually written
  \(\mathcal{U}(\mathcal{H})\), with product as group operation. We usually write \(u_g\) for the image of
  an element \(g\) in \(G\) under the map \(u\). We also say that \(u\) is a unitary 
  representation of \(G\) on \(\mathcal{H}\). Unitary representations \((u,\mathcal{H})\) and \((v,\mathcal{K})\) are \textnormal{unitarily
  equivalent} if there is a unitary operator \(w : \mathcal{H} \rightarrow \mathcal{K}\) such that \(wu_g w^{-1} = v_g\),
  for all \(g\) in \(G\).
\end{definition}
Recall that a unitary operator \(u\) satisfies \(u^*u = uu^* = I\), the identity
operator. Hence, if \((u,\mathcal{H})\) is a unitary representation of \(G\), we have
\(u_{g^{-1}}=u_g^{-1}=u_g^*\), for any \(g\) in \(G\). In particular, the range of the function \(u\), while
not even a linear space, is, at least, *-closed.\\
Elementary calculations show that the weak, \(\sigma\)-weak, and strong topologies coincide on \(\mathcal{U}(\mathcal{H})\), 
so that we may use either in the definition of a unitary representation.\par
For any Hilbert space \(\mathcal{H}\), there is the trivial representation of \(G\) on \(\mathcal{H}\)
which sends each group element to the identity operator.
\begin{definition}
  Let \(G\) be a discrete group. Its \textnormal{left regular representation}
  is the unitary representation of \(G\) on the Hilbert space \(\ell^2(G)\) defined by
  \[\lambda_g \xi(h) = \xi(g^{-1}h),\]
  for all \(g,h\) in \(G\) and \(\xi\) in \(\ell^2(G)\).\par
  Let \(G\) be a discrete group. The Hilbert space \(\ell^{2}(G)\) 
  is the space of all complex-valued functions on \(G\) whose squared absolute values are summable:
  \[
  \ell^{2}(G)
  =
  \left\{
  \xi : G \to \mathbb{C}
  \;\middle|\;
  \sum_{g \in G} |\xi(g)|^{2} < \infty
  \right\}.
  \]
  We will let \(\delta_h\) be the function that is \(1\) at \(h\)
  and zero elsewhere, for some fixed group element \(h\). Notice that the set
  \(\{\delta_h | h \in G\}\) forms an orthonormal basis for \(\ell^2(G)\). Also observe that for any
  \(g, h\) in \(G\), we have \(\lambda_g\delta_h = \delta_{gh}\).
\end{definition}
\begin{definition}
  Let \((u,\mathcal{H})\) be a unitary representation of the discrete group
  \(G\). A closed subspace \(\mathcal{N} \subset \mathcal{H}\) is \textnormal{invariant} for \(u\) if \(u_g\mathcal{N} \subset \mathcal{N}\), 
  for all \(g\) in \(G\).
  We say that the representation is \textnormal{irreducible} if its only invariant subspaces
  are \(0\) and \(\mathcal{H}\) and is \textnormal{reducible} otherwise.
\end{definition}
\begin{definition}
  If \((u_\alpha,\mathcal{H}_\alpha),\alpha \in A\) is a collection of unitary 
  representations of a group \(G\), then their direct sum \((\bigoplus_{\alpha}u_\alpha,\bigoplus_\alpha \mathcal{H}_\alpha)\) 
  is the unitary representation defined by \({(\bigoplus_\alpha u_\alpha)}_g = \bigoplus_\alpha{(u_\alpha)}_g\), for all \(g\) in \(G\).
\end{definition}
\begin{proposition}
  If \(\mathcal{N}\) is an invariant subspace for the unitary 
  representation \((u,\mathcal{H})\), then so is \(\mathcal{N}^\perp\). In this case, \(u\) is unitarily equivalent to
  \(u|_\mathcal{N} \bigoplus u|_{\mathcal{N}^\perp}\).
\end{proposition}
\begin{proof}
  Same as~\ref{prop:1115}.
\end{proof}   

\section{Group algebras}
\begin{definition}
  Let \(G\) be a group. Its (complex) \textnormal{group algebra}, denoted
  \(\mathbb{C}G\), consists of all formal sums \(\sum_{g\in G} a_g g\), where each \(a_g\) is a complex number
  and \(a_g = 0\), for all but finitely many \(g\) in \(G\). Defining \(a_g g \cdot a_h h = a_g a_h gh\),
  for all \(g,h \in G\) and \(a_g,a_h \in \mathbb{C}\) and extending by linearity, \(\mathbb{C}G\) becomes a
  complex algebra. Moreover, defining \(g^* = g^{-1}\) and extending to be conjugate
  linear, \(\mathbb{C}G\) becomes a complex *-algebra.
\end{definition}
If we adopt the notation \(1g = g\), for any \(g\) in \(G\), then \(G \subset \mathbb{C}G\). The
identity element of \(G\) is also the identity of this algebra. Moreover, each
element \(g\) in \(G\) is a unitary in \(\mathbb{C}G\).
\begin{proposition}\label{prop:222}
  Let \(G\) be a group. Its group algebra is commutative if
  and only if \(G\) is abelian.
\end{proposition}
\begin{proof}
  Let us just note that
  \[
  \left(\sum_{g\in G}a_g g \right)\left( \sum_{g'\in G}b_{g'} g'\right) = \sum_{g,g'\in G} a_g g b_{g'} g' 
  = \sum_{h\in G}\left( \sum_{g\in G}a_g b_{g^{-1}h} \right)h,
  \]
  where \(h=gg'\), hence \(g'=g^{-1}h\).
\end{proof}
\begin{theorem}\label{theorem:223}
  Let \(G\) be a discrete group. If \(u : G \rightarrow \mathcal{U}(\mathcal{H})\) is a unitary
  representation of \(G\) on the Hilbert space \(\mathcal{H}\), then \(u\) has a unique extension to
  a unital representation of \(\mathbb{C}G\) on \(\mathcal{H}\) defined by
  \[
  \pi_u\left(\sum_{g\in G}a_g g \right)=\sum_{g\in G}a_g u_g
  \]
  for \(\sum_{g\in G}a_g g\) in \(\mathbb{C}G\). Moreover, if \(\pi : \mathbb{C}G \rightarrow \mathcal{B}(\mathcal{H})\) is a unital 
  representation, then its restriction to \(G\) is a unitary representation of \(G\). Finally, the
  representation \(u\) is irreducible if and only if \(\pi_u\) is.
\end{theorem}
\begin{proof}
  The first part of the statement is immediate from the definition and
  some simple calculations which we leave to the reader.\par
  For the last statement, we proceed as follows. For each group element \(g\),
  we may regard \(g\) as an element of \(\mathbb{C}G\). We define \(u_g = \pi(g)\). We must verify
  that \(u_g\) is unitary. To this end, we use the facts that \(\pi\) is a *-representation
  and that \(g^* = g^{-1}\) in \(\mathbb{C}G\) to obtain
  \[{\pi(g)}^*\pi(g) = \pi(g^*)\pi(g) = \pi(g^{-1})\pi(g) = \pi(g^{-1}g) = \pi(1) = 1.\]
  A similar computation proves that \(\pi(g){\pi(g)}^* = 1\) as well, so \(\pi(g)\) is 
  unitary. The fact that \(u\) is a homomorphism follows from the fact that \(\pi\) is
  multiplicative.\par
  For the last statement, if \(u\) is reducible, then there is a non-trivial invariant 
  subspace for all of the operators \(u_g,g \in G\). But then that subspace is also
  clearly invariant for any linear combination of these operators, hence for every
  element of \(\pi_u(\mathbb{C}G)\). Conversely, if the collection of operators \(\pi_u(a),a \in \mathbb{C}G\)
  has a common non-trivial invariant subspace, then \(u_g = \pi(g)\), for every group
  element \(g\), also leaves this space invariant.
\end{proof}
\begin{proposition}
  Let \(G\) be a discrete group. Let \((\lambda, \ell^2(G))\) be the left
  regular representation of \(G\). The associated representation of \(\mathbb{C}G\) is
  \[
  \left[ \pi_\lambda \left( \sum_{g\in G}a_g g \right)\xi \right](h)=\sum_{g\in G}a_g\xi (g^{-1}h)
  \]
  for all \(\sum_{g\in G}a_g g\) in \(\mathbb{C}G\), \(h\) in \(G\) and \(\xi\) in \(\ell^2(G)\).
\end{proposition}
\begin{theorem}\label{theorem:225}
  Let \(G\) be a group. The left regular representation \(\pi_\lambda\) of \(\mathbb{C}G\)
  is injective.
\end{theorem}
\begin{proof}
  Let \(a = \sum_{g\in G} a_g g\) be an element of \(\mathbb{C}G\) and assume that it is in the
  kernel of \(\pi_\lambda\). Fix \(g_0\) in \(G\) and let \(\delta_{g_0},\delta_e\) be as described earlier. An easy
  computation shows that
  \[
  0 =<\pi_\lambda(a)\delta_{g_0},\delta_e >=< \pi_\lambda(\sum_{g\in G} a_g g)\delta_{g_0}, \delta_e >
  =\sum_{h\in G}\sum_{g\in G}a_g\delta_{g_0}(g^{-1}h)\delta_e(h) = a_{{g_0}^{-1}}.
  \]
  As \(g_0\) was arbitrary, we conclude \(a = 0\).
\end{proof}
\begin{theorem}
  Let \(G\) be a group. The map from \(\mathbb{C}G\) to \(\mathbb{C}\) defined by
  \(\tau(\sum_{g\in G} a_g g) = a_e\), for any \(\sum_{g\in G}a_g g\) in \(\mathbb{C}G\), is a faithful trace. (Recall that
  a trace \(\tau\) is faithful if \(\tau(a^*a) = 0\) only if \(a = 0\).)
\end{theorem}
\begin{proof}
  First of all, it is clear that \(\tau\) is linear. Next, we check that
  it is positive and faithful. For any \(a = \sum_{g\in G} a_g g\) in \(\mathbb{C}G\), we have
  \begin{align*}
    a^*a&={\left(\sum_{g\in G} a_g g\right)}^*\left(\sum_{h\in G} a_h h\right)\\
    &={\left(\sum_{g\in G} \overline{a_g} g^{-1}\right)}\left(\sum_{h\in G} a_h h\right)\\
    &=\sum_{g,h\in G}\overline{a_g}a_h g^{-1}h
  \end{align*}
  Hence \(\tau(a^*a)\) is the coefficient of \(e\) which is the sum over \(g^{-1}h = e\), or \(g = h\): 
  \[
  \tau(a^*a) = \sum_{g\in G}\overline{a_g}a_g =\sum_{g\in G}|a_g|^2 \geq 0,
  \]
  and \(\tau(a^*a) = 0\) implies \(a = 0\).\par
  Finally, we verify the trace property. Consider simply \(\tau(gh)\), where \(g\) and
  \(h\) are in \(G\). This is clearly one when \(g^{-1} = h\) and zero otherwise. Exactly
  the same argument applies to \(\tau(hg)\) and so we see that \(\tau(gh) = \tau(hg)\), for
  any \(g,h\) in \(G\). The trace property follows from linearity.
\end{proof}
\begin{exercise}
  Let \(G\) be a discrete group and \(\tau\) be the trace on the group
  algebra as above. What is the result of applying the GNS construction to the
  trace? (Warning: nice answer, trivial proof.)
\end{exercise}

\section{Finite groups}
\begin{theorem}
  Let \(G\) be a finite group. Then there is a unique norm on
  \(\mathbb{C}G\) in which it is a \cstar.
\end{theorem}
\begin{proof}
  The formula \(\|a\| = \|\pi_\lambda(a)\|\) (operator norm) defines a semi-norm on
  \(\mathbb{C}G\), where \(\pi_\lambda\) is the left regular representation as above. The fact that it
  is a norm follows from the last theorem. It satisfies the \(C^*\)-condition since
  \(\pi_\lambda\) is a *-homomorphism and the operator norm satisfies the \(C^*\)-condition.
  Finally, it is complete in this norm because \(\mathbb{C}G\) is finite dimensional. The
  uniqueness is just a restatement of~\ref{corollary:138}.
\end{proof}
With the information that we already have available from Section~\ref{sec:17}
about the structure of finite-dimensional \cstars, we can completely 
describe the \cstar\ of a finite group, at least in principle. Recall that
two elements \(g_1,g_2\) of \(G\) are \textit{conjugate} if there is an element \(h\) such that
\(hg_1h^{-1} = g_2\). Conjugacy is an equivalence relation and the equivalence class
of an element \(g\), \(\{hgh^{-1} | h \in G\}\), is called a \textit{conjugacy class}.
\begin{theorem}\label{theorem:232}
  Let \(G\) be a finite group with conjugacy classes
  \(C_1,C_2,\ldots,C_K\). For each \(1 \leq k \leq K\), define
  \[
  c_k =\sum_{g\in C_k} g \in \mathbb{C}G.
  \]
  The set \(\{c_1,\ldots,c_K\}\) is linearly independent and its span is the centre of \(\mathbb{C}G\).
  In particular, \(\mathbb{C}G\) is isomorphic to \(\bigoplus^K_{k=1} M_{n_k}(\mathbb{C})\) and
  \[
  \sum^K_{k=1} n_k^2 = \#G.
  \]
\end{theorem}
\begin{proof}
  It is clear that \(G\) is a linear basis for \(\mathbb{C}G\), so \(\mathbb{C}G\) is finite dimensional.
  Hence by Theorem~\ref{theorem:171}, \(\mathbb{C}G\) is isomorphic to \(\bigoplus^L_{\ell=1}M_{n_\ell}(\mathbb{C})\) for some positive
  integers \(n_1,\ldots,n_L\).\par
  Now suppose that \(a = \sum_g a_g g\) is in the centre of \(\mathbb{C}G\). Let \(h\) be any other
  element of \(G\). Considering \(h\) as an element of \(\mathbb{C}G\), it is invertible and its
  inverse is \(h^{-1}\). As \(a\) is in the centre of \(\mathbb{C}G\), we have
  \[
  \sum_g a_g g = a = hah^{-1} = \sum_g a_g hg h^{-1} = \sum_g a_{h^{-1}g h} g.
  \]
  For any \(g\) in \(G\), comparing coefficients, we have \(a_g = a_{h^{-1}gh}\). This means
  that the function \(a\) is constant on conjugacy classes in \(G\). Hence it is a linear
  combination of the \(c_k\).\par
  Conversely, the same computation shows that each \(c_k\) commutes with
  every group element \(h\). Since the group elements span the group algebra,
  each \(c_k\) is in the centre.\par
  The fact that the set of all \(c_k\) is linearly independent is clear. Thus the centre of \(\mathbb{C}G\) has
  dimension \(K\). On the other hand, the centre of \(\bigoplus^L_{\ell=1}M_{n_\ell}(\mathbb{C})\) has dimension \(L\),
  so \(L=K\). Finally, comparing dimensions gives \(\#G=\dim(\mathbb{C}G)=\sum_{k=1}^K n_k^2\).
\end{proof}
\begin{example}\label{example:233}
  \(\mathbb{C}S_3 \thicksim = M_2 \oplus\mathbb{C}\oplus\mathbb{C}\).
\end{example}
\begin{proof}
  It is straightforward to check that there are three conjugacy classes: the
  identity, the three transpositions and the two cycles of length three. So we
  need to find \(n_1,n_2,n_3\) whose squares sum to \(6\).
\end{proof}
\begin{exercise}
  Give another proof of Example~\ref{example:233} which doesn't consider
  conjugacy classes but only linear dimension and Theorem~\ref{theorem:232}.
\end{exercise}
\begin{exercise}
  Describe \(\mathbb{C}S_4\) in terms of Theorem~\ref{theorem:171}.
\end{exercise}

\section{The \texorpdfstring{\(C^*\)-algebra}{C*algebra} of a discrete group}
\begin{definition}
  Let \(G\) be a discrete group. The \(\ell^1\)-norm on \(\mathbb{C}G\) is defined
  by
  \[\Vert {\sum_{g\in G} a_g g}\Vert_1 =\sum_{g\in G}|a_g|,\]
  where the complex numbers \(a_g\) are nonzero for only finitely many \(g\).
\end{definition}
\begin{theorem}
  Let \(G\) be a discrete group. The completion of \(\mathbb{C}G\) in the
  \(\ell^1\)-norm is
  \[
  \ell^1(G)=\{a:G\rightarrow\mathbb{C}| \sum_{g\in G}|a(g)|<\infty\}
  \]
  Moreover, the product
  \[
  (ab)(g)=\sum_{h\in G}a(h)b(h^{-1}g)
  \]
  for \(a,b\) in \(\ell^1(G)\) and \(g\) in \(G\), is well-defined, associative, extends the product
  on \(\mathbb{C}G\) and makes \(\ell^1(G)\) a Banach algebra. The involution
  \[a^*(g) =\overline{a(g^{-1})}\]
  for \(a\) in \(\ell^1(G)\) and \(g\) in \(G\), is isometric, conjugate linear and satisfies 
  \({(ab)}^* = b^*a^*\), for all \(a,b\) in \(\ell^1(G)\).
\end{theorem}
\begin{proof}
  The proof of the first statement is a standard fact in functional analysis. 
  To see that the product is well-defined, we
  note that the functions in \(\ell^1(G)\) are bounded since they are summable and
  then the sum \(\sum_{h\in G} a(h)b(h^{-1}g)\) is the product of a summable sequence and
  a bounded one and so is summable. To see it is in \(\ell^1(G)\), for any \(a,b\) in \(\ell^1(G)\),
  we have
  \[
  \sum_{g,h\in G}|a(h)b(h^{-1}g)|=\sum_{h\in G}|a(h)|\sum_{g\in G}|b(h^{-1}g)|=\sum_{h\in G}|a(h)| \|b\|_1=\|a\|_1\|b\|_1
  \]
  It follows that \(\|ab\|_1 \leq \|a\|_1 \|b\|_1\). It is clear that the involution is isometric.
  It is clearly conjugate linear and satisfies \({(ab)}^* = b^*a^*\) on functions that
  are supported on a single element. By taking linear combinations, it is also
  satisfied for finitely supported functions. It then holds on all elements of \(\ell^1(G)\) by continuity.
\end{proof}
\begin{theorem}
  Let \(G\) be a discrete group. The linear functional \(\tau(a) =
  a(e), a \in \ell^1(G)\) is bounded and is a trace on \(\ell^1(G)\).
\end{theorem}
\begin{proof}
  We note that for any \(a\) in \(\ell^1(G)\), we have \(\|a\|_1 \geq |a(e)| = |\tau(a)|\), so
  \(\tau\) is bounded. Since we have already proved that it is a trace on \(\mathbb{C}G\) and is
  continuous, the trace property extends to \(\ell^1(G)\).
\end{proof}
\begin{theorem}\label{theorem:244}
  Let \(G\) be a discrete group. There are bijective correspondences between
  \begin{enumerate}
    \item \((u,\mathcal{H})\), unitary representations of \(G\),
    \item \((\pi,\mathcal{H})\), non-degenerate representations of \(\mathbb{C}G\),
    \item \((\pi,\mathcal{H})\), non-degenerate representations of \(\ell^1(G)\).
  \end{enumerate}
  The correspondences from the third to the second and from the second to the
  first are both obtained by restriction. Moreover, these correspondences preserve 
  the irreducible representations. Finally, every representation of \(\ell^1(G)\)
  is contractive.
\end{theorem}
\begin{proof}
  We are using the facts that \(G \subset \mathbb{C}G \subset \ell^1(G)\) to take restrictions.
  Going back, unitary representations of \(G\) extend to unital representations of
  \(\mathbb{C}G\) by linearity exactly as we saw already in Theorem~\ref{theorem:223}. Unital representations of
  \(\mathbb{C}G\) extend to \(\ell^1(G)\) by continuity as follows.\par
  Suppose that \(\pi\) is a unital representation of \(\mathbb{C}G\). If \(a\) is in \(\ell^1(G)\), we define
  \[
  \pi(a) =\sum_{g\in G}a(g)\pi(g).
  \]
  Of course, the reason the sum converges is simply because the coefficients
  are absolutely summable and since \(\pi\) is a *-homomorphism, \(\pi(g)\) is unitary
  and hence has norm one, for every \(g\) in \(G\). In fact, we have
  \[
  \|\pi(a)\| = \|\sum_{g\in G} a(g)\pi(g)\| \leq
  \sum_{g\in G}|a(g)| \|\pi(g)\| =
  \sum_{g\in G}|a(g)| = \|a\|_1
  \]
  and this also proves that \(\pi\) is contractive.\par
  If \(\pi\) is a representation of \(\ell^1(G)\) having a non-trivial invariant subspace,
  then its restriction to \(\mathbb{C}G\) has the same non-trivial invariant subspace. 
  Similarly, the restriction of a reducible representation \(\mathbb{C}G\) to \(G\) will also be
  reducible. In the other direction, if a unitary representation of \(G\) has a non
  trivial invariant subspace, then extending it to \(\mathbb{C}G\) by linearity has the same
  non-trivial invariant subspace. Moreover, if \(a_n\) is any sequence in \(\mathbb{C}G\) 
  converging to \(a\) in \(\ell^1(G)\) and \(\pi(a_n),n \geq 1\) all leave \(\mathcal{N}\) invariant, then so does
  \(\pi(a)\).
\end{proof}
\begin{definition}\label{def:245}
  Let \(G\) be a discrete group. We define a norm on \(\mathbb{C}G\) by
  \[\|a\| = \sup\{\|\pi(a)\| | \pi\mbox{, a representation of }\mathbb{C}G\},\]
  for any a in \(\mathbb{C}G\). We denote the completion of \(\mathbb{C}G\) in this norm by \(C^*(G)\),
  which we refer to as the \textnormal{group \cstar} of \(G\).
\end{definition}
\begin{theorem}
  The norm given in Definition~\ref{def:245} is well-defined and
  \(C^*(G)\) is a \cstar\ containing \(\mathbb{C}G\) as a dense *-subalgebra.
\end{theorem}
\begin{proof}
  Begin by noting that every representation \(\pi\) of \(\mathbb{C}G\) has a unique 
  extension to \(\ell^1(G)\) and that every representation of \(\ell^1(G)\) is contractive. This
  implies that \(\|\pi(a)\| \leq \|a\|_1\) and so the set on the right is bounded and the
  supremum exists, for any element \(a\) of \(\mathbb{C}G\). Next, the left regular 
  representation of \(\mathbb{C}G\) is faithful and so this norm is also faithful.\par
  The fact that on \(\mathbb{C}G\) it satisfies \(\|ab\| \leq \|a\| \|b\| , a,b \in \mathbb{C}G\) and the \(C^*\)
  condition follow from the fact that it is the supremum of norms all satisfying
  these conditions. It is then easy to see its completion is a \cstar.
\end{proof}
\begin{proposition}
  If \(G\) is a discrete group, then \(C^*(G)\) is a unital \cstar\ and it is abelian if and only if \(G\) is abelian.
\end{proposition}
\begin{proof}
  The identity of the group \(G\) is also the identity of the group algebra
  \(\mathbb{C}G\) and is therefore the identity of any completion. We have also already
  seen in Proposition~\ref{prop:222} that \(G\) is abelian if and only if \(\mathbb{C}G\) is abelian. It is
  easy to see that a dense subalgebra of a \cstar\ is abelian if and only if
  the whole \cstar\ is also.
\end{proof}
\begin{theorem}
  Let \(G\) be a discrete group. For any element \(a\) in \(\mathbb{C}G\), we
  have
  \[
  \|a\| = \sup\{\|\pi_u(a)\| | u\mbox{ an irreducible representation of }G\}.
  \]
\end{theorem}
\begin{proof}
  It is a trivial consequence of Theorem~\ref{theorem:244} and
  Lemma~\ref{lemma:1129}.
\end{proof}
\begin{definition}
  Let \(G\) be a discrete group. We define the reduced norm
  on \(\mathbb{C}G\) by
  \[\|a\|_r = \|\pi_\lambda(a)\|,\]
  for any \(a\) in \(\mathbb{C}G\). We denote the completion of \(\mathbb{C}G\) in this norm by \(C^*_r(G)\),
  which we refer to as the \textnormal{reduced group \cstar} of \(G\).
\end{definition}
\begin{theorem}\label{theorem:2410}
  Let \(G\) be a discrete group. There exists a canonical 
  *-homomorphism\\ \(\rho : C^*(G)\rightarrow C^*_r(G)\) which extends the identity map on the
  group algebra \(\mathbb{C}G\).
\end{theorem}
\begin{proof}
  If \(\lambda\) denotes the left regular representation, then \(\pi_\lambda\) is among the 
  collection of representations of \(\mathbb{C}G\) considered when defining the norm. Hence,
  we have \(\|a\|_r = \|\pi_\lambda(a)\| \leq \|a\|\). The rest follows from the definitions.
\end{proof}
\begin{theorem}
  Let \(G\) be a discrete group. For each \(a\) in \(\mathbb{C}G\), we have
  \[\tau(a) =< \pi_\lambda(a)\delta_e,\delta_e > .\]
  The trace \(\tau\) extends continuously to both \(C^*(G)\) and \(C^*_r(G)\) and the same
  formula holds for \(a\) in either of these. Both extensions are traces on the
  respective \cstars.
\end{theorem}
\begin{proof}
  We compute
  \begin{align*}
    <\pi_\lambda(a)\delta_e,\delta_e > &= < \pi_\lambda(\sum_{g\in G}a(g)g)\delta_e, \delta_e >\\
    &=\sum_{g\in G}a(g) < \lambda(g)\delta_e,\delta_e >\\
    &=\sum_{g\in G} a(g) < \delta_g,\delta_e >\\
    &= a(e)\\
    &= \tau(a).
  \end{align*}
  It follows immediately from this that
  \[|\tau(a)| \leq \|\pi_\lambda(a)\| = \|a\|_r \leq \|a\| ,\]
  and so \(\tau\) extends as claimed. The fact that \(\tau\) satisfies the trace properties on
  dense subalgebras of \(C^*_r(G)\) and \(C^*(G)\) implies the extensions will also.
\end{proof}
\begin{exercise}
  Let \(G\) be a discrete group. For each group element \(g\), we
  let \(\delta_g\) be the function which is \(1\) at \(g\) and \(0\) elsewhere. We regard this as an
  element of \(\ell^2(G)\). Let \(a\) be in \(\mathcal{B}(\ell^2(G))\).
  \begin{enumerate}
    \item Prove that if \(a\) is in the closure of \(\pi_\lambda(\mathbb{C}G)\), then \((a\delta_e)(g) =< a\delta_h,\delta_g h >\),
    for all \(g,h\) in \(G\).
    \item Prove that if \(a\) is in the centre of the closure of \(\pi_\lambda(\mathbb{C}G)\), then the
    function \(a\delta_e\) is constant on conjugacy classes in \(G\).
    \item In the last two problems, which topologies can you use when taking the
    closure?
    \item Prove that if every conjugacy class in \(G\) is infinite, except that of the
    identity, then \(C^*_r(G)\) has a trivial centre.
    \item Prove that the groups \(\mathbb{F}_2\) and \(S_\infty\) have the infinite conjugacy class property.
  \end{enumerate}
\end{exercise}
\begin{exercise}
  As we noted above, the identity map on \(\mathbb{C}G\) extends to a
  *-homomorphism from \(\ell^1(G)\) to \(C^*(G)\). Prove that the composition of this
  map with \(\rho\) of~\ref{theorem:2410} is injective. (Hint: try to mimic the proof of Theorem~\ref{theorem:225}.)
\end{exercise}

\section{Abelian groups}
\begin{definition}
  Let \(G\) be a discrete abelian group. Its \textnormal{dual group}, denoted
  \(\hat{G}\), consists of all group homomorphisms \(\chi: G \rightarrow \mathbb{T}\).\footnote{\(\mathbb{T}\)
  denotes the circle group (also called the 1-torus):
  \[\mathbb{T}=\{z\in\mathbb{C}:|z|=1\}.\]}
\end{definition}
\begin{theorem}
  If \(G\) is a discrete, abelian group, then \(\hat{G}\) is exactly the set
  of irreducible representations of \(G\) (with Hilbert space \(\mathbb{C}\) in each case).
\end{theorem}
\begin{proof}
  Follows from Theorem~\ref{theorem:223} and Exercise~\ref{ex:1113}.
\end{proof}
We comment that often notationally, the value of \(\chi\) in \(\hat{G}\) on an element
\(g\) of \(G\) is written as \(< g,\chi >\). This tends to emphasize a symmetry which
we will explore further once we remove the hypothesis that \(G\) is discrete in
Section~\ref{section:28}.
\begin{theorem}
  Let \(G\) be a discrete abelian group. Its dual group, \(\hat{G}\), is a
  group when endowed with the pointwise product of functions. Moreover, the
  collection of sets
  \[U(\chi_0,F, \epsilon) = \{\chi \in \hat{G} | |\chi(g)-\chi_0(g)| < \epsilon,g \in F\},\]
  where \(\chi_0\) is in \(\hat{G}\), \(F \subset G\) is finite and \(\epsilon > 0\), forms a neighbourhood base for
  a topology on \(\hat{G}\). In this, \(\hat{G}\) is compact and Hausdorff.
\end{theorem}
\begin{proof}
  The group $\widehat G$ is naturally a subset of

  \[
  \mathbb T^G=\prod_{g\in G}\mathbb T,
  \]

  the set of all functions $G\to\mathbb T$, endowed with the product
  topology. A basic neighbourhood of $f\in\mathbb T^G$ is determined by
  finitely many coordinates, so the sets

  \[
  U(\chi_0,F,\varepsilon)
  =
  \{\chi\in\widehat G:
  |\chi(g)-\chi_0(g)|<\varepsilon,\ g\in F\}
  \]

  are precisely the basic open sets of the subspace topology on
  $\widehat G$.

  To show that $\widehat G$ is closed, for each $g,h\in G$ define

  \[
  \Phi_{g,h}:\mathbb T^G\to\mathbb T,
  \qquad
  \Phi_{g,h}(f)=f(gh)\,\overline{f(g)f(h)}.
  \]

  Each $\Phi_{g,h}$ is continuous, and

  \[
  \widehat G
  =
  \bigcap_{g,h\in G}
  \Phi_{g,h}^{-1}(\{1\}).
  \]

  Since $\{1\}$ is closed in $\mathbb T$, it follows that $\widehat G$ is
  closed in $\mathbb T^G$.

  Since $\mathbb T$ is compact and Hausdorff, Tychonoff's theorem implies
  that $\mathbb T^G$ is compact and Hausdorff. Therefore the closed subset
  $\widehat G$ is also compact and Hausdorff.
\end{proof}
\begin{proposition}
  The map which sends \(\chi\) to \(\chi(1)\) is an isomorphism and
  homeomorphism from \(\hat{\mathbb{Z}}\) to \(\mathbb{T}\).
\end{proposition}
\begin{proof}
  Since \(\mathbb{Z}\) is cyclic, any \(\chi\) in \(\hat{\mathbb{Z}}\) is uniquely determined by its value at
  \(1\), the generator of \(\mathbb{Z}\). This means our map is injective. Similarly, the fact
  that \(\mathbb{Z}\) is free means that, for any \(z\) in \(\mathbb{T}\), there is a homomorphism \(\chi\) with
  \(\chi(1) = z\): specifically, the map \(\xi(n) = z^n\). This means our map is onto. It
  follows easily from the definitions that it is a group isomorphism.\par
  The map is clearly continuous and since both spaces are compact and
  Hausdorff, it is a homeomorphism as well.
\end{proof}
\begin{theorem}\label{theorem:255}
  If \(G\) is a discrete abelian group, then \(C^*(G)\) is isomorphic
  to \(C(\hat{G})\). The isomorphism takes a group element \(g \in \mathbb{C}G \subset C^*(G)\) to the
  function \(\hat{g}(\chi) = \chi(g)\).
\end{theorem}
\begin{proof}
  We know that \(C^*(G)\) is isomorphic to \(C(\mathcal{M}(C^*(G)))\). Moreover, this
  map sends group element \(g\) to the function \(\hat{g}(\phi) = \phi(g)\). We will be done
  if we can identify \(\mathcal{M}(C^*(G))\) with \(\hat{G}\). First, if \(\phi\) is in \(\mathcal{M}(C^*(G))\), we may
  restrict it to \(G \subset \mathbb{C}G \subset C^*(G)\). The fact that \(\phi\) is multiplicative means that
  \(\phi|_G\) is a group homomorphism. This means that restriction defines a map
  from \(\mathcal{M}(C^*(G))\) to \(\hat{G}\). Two multiplicative linear functionals that agree on
  \(G\) are clearly equal since the elements of \(G\) span a dense subset of \(C^*(G)\), so
  this map is injective.\par
  Now we claim this map is surjective. Let \(\chi\) be in \(\hat{G}\). If we regard \(\chi\) as a
  unitary representation on a one-dimensional Hilbert space, then by Theorem~\ref{theorem:223}, 
  it is the restriction of a unique representation of \(\mathbb{C}G\) on this same
  Hilbert space. It is obviously continuous in the \cstar\ norm on \(\mathbb{C}G\) and
  hence extends to \(C^*(G)\). But such a map is simply a multiplicative linear
  functional. This establishes the claim.\par
  It is clear from the definition that the map is continuous, but as both
  spaces are compact and Hausdorff, it is a homeomorphism as well.
\end{proof}
\begin{exercise}\label{ex:251}
  Let \(n\) be a positive integer. Find \(\hat{\mathbb{Z}_n}\), the dual group of the
  cyclic group of order \(n\).
\end{exercise}
\begin{exercise}
  Using the answer to exercise~\ref{ex:251} above and Theorem~\ref{theorem:255},
  write each element of \(\mathbb{Z}_n\) as a continuous function on \(\hat{\mathbb{Z}_n}\). Also, find an
  element of \(\mathbb{C}\mathbb{Z}_n\) (written as a linear combination of group elements) which is
  a minimal central projection.
\end{exercise}
\begin{exercise}
  Let
  \[
    G=\mathbb{Z}[1/2] = \{\frac{n}{2^k} | n,k \in \mathbb{Z},k \geq 0 \}.
  \]
  Describe \(\hat{G}\). (Remark: the best one can do here is try to describe an element
  \(\chi\) in \(\hat{G}\) in terms of its restrictions to the subgroups
  \[
    \mathbb{Z} \subset \frac{1}{2}\mathbb{Z} \subset \frac{1}{4}\mathbb{Z} \subset \cdots
  \]
  in \(G\).)
\end{exercise}
\begin{exercise}
  With \(G\) as above, describe \(\hat{G/\mathbb{Z}}\).
\end{exercise}


\section{Locally compact groups}\label{section:28}
In this section, \(G\) denotes a locally compact group, i.e., \(G\) is a group and 
a locally compact Hausdorff space, and the map \((s,t)\mapsto s^{-1}t\) is continuous 
from \(G\times G\) to \(G\).\par
The central fact of the theory is the existence of a nonzero left-invariant Radon measure \(\mu_G\)
on \(G\), i.e., \(\mu_G(sE)=\mu_G(E)\) for each Borel set \(E\) in \(G\) and each \(s\) in \(G\).
We call \(\mu_G\) a (\textit{left}) \textit{Haar measure} on \(G\). Since \(\mu_G\) is unique 
up to scalar multiples, we talk about \textit{the} Haar measure. If \(\mu_G\) is finite
(which happens if and only if \(G\) is compact), we assume that \(\mu_G(G)=1\), and if \(G\) is discrete and infinite, we assume
that \(\mu_G(\{e\})=1\) (\(e\) denoting the identity of \(G\)).\par
There is of course also a right Haar measure on \(G\). This is linked with the left Haar measure by the existence of a continuous homomorphism
\(\Delta\) (the \textit{modular function}) of \(G\) into the positive real numbers such that \(\mu_G(Es)=\Delta(s)\mu_G(E)\).
Thus, with the customary notation \(\,ds\) instead of \(\,d\mu_G(s)\), we have
\[
  \,d(ts)=\,ds,\qquad \,d(st)=\Delta(t)\,ds,\qquad\,d(s^{-1})=\Delta(s)^{-1}\,ds.
\]\par

Let \(M(G)\) denote the Banach space of bounded complex Radon measures (signed measures) on \(G\), identified with \(C_0(G)^*\).
For \(\mu\) and \(\nu\) in \(M(G)\), define \textit{the convolution} and \textit{involution} by
\begin{align*}
  \int f(s)\,d(\mu\times \nu)(s)&=\int\int f(ts)\,d\nu(s)\,d\mu(t),\qquad f\in C_0(G);\\
  \int f(s)\,d\mu^*(s)&=\overline{\int\overline{f(s^{-1})}\,d\mu(s)},\qquad f\in C_0(G).
\end{align*}
(We use \(\times\) for convolution instead of the usual \(*\) to avoid confusion with the adjoint operator.
The product measure of \(\mu\) and \(\nu\) on \(G\times G\) is denoted by \(\mu\otimes\nu\).) 
With convolution as the product, \(M(G)\) becomes a Banach algebra with an isometric involution but not in general a \cstar.\par
We denote by \(L^1(G)\) the closed *-ideal in \(M(G)\) of complex measures that are absolutely continuous 
with respect to Haar measure and identify \(L^1(G)\) with the (classes of) \(\mu_G\)-integrable functions on \(G\). Thus, for \(f\) and \(g\) in \(L^1(G)\),
\begin{align*}
  f\times g(t)&=\int f(s)g(s^{-1}t)\,ds=\int f(ts)g(s^{-1})\,ds,\\
  f^*(t)&=\Delta(t)^{-1}\overline{f}(t^{-1}).
\end{align*}
Occasionally, we will also need the transformation \(f\mapsto\tilde{f}\), where 
\(\tilde{f}(t)=\overline{f}(t^{-1})\). This transformation is also an involution, but if \(G\) is not unimodular, we may have \(\tilde{f}\notin L^1(G)\) even though \(f\in L^1(G)\).\par
For each \(s\) in \(G\), let \(\delta_s\) be the point measure at \(s\). Then, for \(f\) in \(L^1(G)\).
\[
  (\delta_s\times f)(t)=f(s^{-1}t),\qquad (f\times \delta_s)(t)=\Delta(s)^{-1}f(t s^{-1}).
\]
\begin{proposition}\label{Petersen:714}
  There are bijective correspondences between the sets of unitary representations of \(G\),
  representations of \(M(G)\) whose restrictions to \(L^1(G)\) are nondegenerate, and nondegenerate representations of \(L^1(G)\).
\end{proposition}
\begin{proof}
  If \((u,\mathcal{H})\) is a unitary representation of \(G\), then define, for all \(\mu\) in \(M(G)\)
  and \(\xi,\eta\) in \(\mathcal{H}\),
  \[
    (\pi(\mu)\xi|\eta) = \int (u_s\xi|\eta)\,d\mu(s).
  \]
  Then \((\pi,\mathcal{H})\) is a representation of \(M(G)\), and using an approximate unit for
  \(L^1(G)\), we see that the restriction to \(L^1(G)\) is nondegenerate.\par
  Conversely, let \((\pi,\mathcal{H})\) be a nondegenerate representation of \(L^1(G)\) and
  choose an approximate unit \({f_\lambda}\) for \(L^1(G)\). Since vectors of the form \(\pi(f)\xi\),
  \(f \in L^1(G)\), \(\xi \in\mathcal{H}\), are dense in \(\mathcal{H}\), it follows that \(\{\pi(f_\lambda)\}\) converges strongly
  to \(1\). We extend \((\pi,\mathcal{H})\) to a representation \((\tilde{\pi},\mathcal{H})\) of \(M(G)\) by defining
  \[
    \tilde{\pi}(\mu)(\pi(f)\xi) =\pi(\mu\times f)\xi
  \]
  for \(f\) in \(L^1(G)\) and \(\xi\) in \(\mathcal{H}\). Equivalently,
  \[
    \tilde{\pi}(\mu)\xi =\lim\pi(\mu\times f_\lambda)\xi\qquad\textnormal{(strongly)},
  \]
  which shows that the extension is unique. The restriction of \((\tilde{\pi},\mathcal{H})\) to \(G\) gives a
  unitary representation of \(G\) whose extension to \(L^1(G)\) is precisely \((\pi,\mathcal{H})\).
\end{proof}
\begin{definition}\label{Petersen:715}
  The \textnormal{universal representation} \((\pi_u,\mathcal{H}_u)\) of \(L^1(G)\) is the direct sum of all nondegenerate
  representations of \(L^1(G)\), and the \textnormal{group \cstar} \(C^*(G)\) is the norm closure
  of \(\pi_u(L^1(G))\) in \(\mathcal{B}(\mathcal{H}_u)\).
  The extension of \((\pi_u,\mathcal{H}_u)\) to \(C^*(G)\) is thus by definition the universal representation of \(C^*(G)\).
  Applying \ref{Petersen:714}, we get an embedding of \(M(G)\) into the multiplier algebra \(M(C^*(G))\) of \(C^*(G)\).
\end{definition}
\begin{proposition}\label{Petersen:716}
  If \(G\) is abelian and \(\hat{G}\) is its dual, then \(C^*(G)\cong C_0(\hat{G})\).
\end{proposition}
\begin{proof}
  The irreducible representations of \(L^1(G)\) (hence also those of \(C^*(G)\))
  are all one-dimensional and correspond to the points in \(\hat{G}\). The topology on \(\hat{G}\)
  is the weak-* topology arising from \(L^1(G)\) or \(C^*(G)\), so that \(\hat{G}\) is homeomorphic
  with the spectrum of \(C^*(G)\). It follows that \(C^*(G)\cong C_0(\hat{G})\).
\end{proof}
\begin{theorem}[Stone]\label{Petersen:717}
  For each unitary representation \((u,\mathcal{H})\) of \(\mathbb{R}\), there is a 
  (not necessarily bounded) self-adjoint operator \(h\) on \(\mathcal{H}\) such that \(u_t=e^{ith}\) for each \(t\) in \(\mathbb{R}\).
\end{theorem}
\begin{proof}
  Let \((\pi,\mathcal{H})\) be the representation of \(L^1(\mathbb{R})\) associated with \((u,\mathcal{H})\) (cf.
  \ref{Petersen:714}). Under the Fourier-transform isomorphism
  \(C^*(\mathbb R)\cong C_0(\hat{\mathbb R})\cong C_0(\mathbb R)\), the spectral
  theorem gives a unique projection-valued measure \(E\) on the Borel subsets
  of \(\mathbb R\) such that
  \[
    \pi(f)=\int_{\mathbb R}f(s)\,dE(s),
    \qquad f\in C_0(\mathbb R);
  \]
  compare \cite[Proposition~8.3.2]{PGK79}. Define the operator \(h\) by
  \[
    \mathcal D(h)
    =
    \left\{\xi\in\mathcal H:
      \int_{\mathbb R}s^2\,d\langle E(s)\xi,\xi\rangle<\infty
    \right\},
    \qquad
    h\xi=\int_{\mathbb R}s\,dE(s)\xi.
  \]
  The spectral theorem for unbounded operators shows that \(h\) is densely
  defined and self-adjoint. The extension of \(\pi\) to bounded Borel
  functions satisfies
  \[
    e^{ith}=\int_{\mathbb R}e^{its}\,dE(s).
  \]
  On the other hand, the multiplier \(\delta_t\in M(C^*(\mathbb R))\)
  corresponds under the Fourier transform to the bounded function
  \(s\mapsto e^{its}\). Since the representation associated with \(u\)
  satisfies \(u_t=\widetilde\pi(\delta_t)\), it follows that
  \[
    u_t
    =\widetilde\pi(\delta_t)
    =\int_{\mathbb R}e^{its}\,dE(s)
    =e^{ith}.
  \]
\end{proof}

Returning to the general case, consider the (universal) unitary representation
\(t\mapsto\delta_t\) of \(G\) into \(M(C^*(G))\) obtained from \ref{Petersen:715}. For each functional \(\phi\) on
\(C^*(G)\), we define a function \(\Phi\) on \(G\) by 
\begin{equation}\label{Petersen:718}
  \Phi(t)=\phi(\delta_t).
\end{equation}
Since \(\phi\) is weakly continuous on \(C^*(G)''\), the function \(\Phi\) is continuous and bounded by \(\|\phi\|\).
If \(f\in L^1(G)\), then identifying \(L^1(G)\) with its image in \(C^*(G)\), we have
\[
  \phi(f)=\int\phi(\delta_t)f(t)\,dt=\int\Phi(t)f(t)\,dt.
\]
which shows that the map \(\phi\mapsto\Phi\) is injective. We denote by \(B(G)\) the image
of \(C^*(G)^*\) in \(C^b(G)\) under this map and let \(B_+(G)\) be the generating cone called \textit{positive definite functions}.
The real vector space spanned by \(B_+(G)\) is denoted by \(B_{sa}(G)\) and corresponds to the set of self-adjoint functionals.
An elementary computation shows that if \(\Phi\) is the image of \(\phi\), then \(\tilde{\Phi}\) is the image of \(\phi^*\).
Thus \(\Phi\in B_{sa}(G)\) if and only if \(\Phi\in B(G)\) and \(\tilde{\Phi}=\Phi\).
From \ref{Petersen:716} we see that if \(G\) is abelian, then the positive definite functions are precisely the Fourier transforms of 
bounded measures on \(\hat{G}\), so that \(B(G)=\hat{M(\hat{G})}\).
\begin{proposition}\label{Petersen:719}
  For each \(\Phi\) in \(C^b(G)\), the following conditions are equivalent:
  \begin{enumerate}
    \item \(\Phi\) is positive definite;
    \item \(\int\Phi(s)\,d(\mu^*\times\mu)(s)\geq 0\) for every \(\mu\) in \(M(G)\);
    \item For each finite set \(\{s_i\}\) in \(G\), the matrix \((a_{ij})\), with \(a_{ij}=\Phi(s^{-1}_i s_j)\), is positive semidefinite;
    \item \(\int\Phi(s)(f^*\times f)(s)\,ds\geq 0\) for every \(f\) in \(C_C(G)\).
  \end{enumerate}
\end{proposition}
\begin{proof}
  \begin{itemize}
    \item \((1)\Rightarrow (2)\). Obvious.
    \item \((2)\Rightarrow (3)\). Given \(\{s_i\}\), take any set \(\{\gamma_i\}\) of complex numbers and form the 
      signed measure \(\mu=\sum\gamma_i\delta_{s_i}\). Condition (2) applied to \(\mu\) gives
      \begin{equation}\label{Petersen:eq719}
        \sum_{ij}\Phi(s^{-1}_i s_j)\overline{\gamma_i}\gamma_j\geq 0,
      \end{equation}
      which implies that the matrix \((\Phi(s_i^{-1}s_j))\) is positive.
    \item \((3)\Rightarrow (4)\). The \((s,t)\mapsto\Phi(s^{-1}t)\overline{f}(s)f(t)\) has a compact support 
      on \(G\times G\) contained in a set \(C\times C\), where \(C\) is compact. Since \(\mu_G(C)<\infty\), we see from Krein-Milman's theorem
      that \(\mu_G|_C\) is the weak-* limit point of a net \(\{\mu_\lambda\}\) of measures
      on \(C\) with finite supports. It follows that \(\mu_G\otimes\mu_G|_{C\times C}\) is the 
      weak-* limit of the net \(\{\mu_\lambda\otimes\mu_\lambda\}\). Condition (3) in the form (\ref{Petersen:eq719}) implies that
      \[
        \int\int\Phi(s^{-1}t)\overline{f}(s)f(t)\,d\mu_\lambda(s)\,d\mu_\lambda(t)=\sum_{i,j}\Phi(s_i^{-1}t_j)\overline{f}(s_i)\overline{\gamma_i}f(t_j)\gamma_j\geq 0,
      \]
      where \(\mu_\lambda=\sum\gamma_i\delta_{t_i}\) and \(\gamma_i\geq 0\). In the weak-* limit, we therefore obtain
      \[
        0\leq\int\int\Phi(s^{-1}t)\overline{f}(s)f(t)\,ds\,dt=\int\Phi(t)(f^*\times f)(t)\,dt.
      \]
    \item \((4)\Rightarrow (1)\). Since \(C_C(G)\) is dense in \(L^1(G)\), the bounded functional \(\phi\) obtained 
      from \(\Phi\) by integration satisfies \(\phi(f^*\times f)\geq 0\) for every \(f\) in \(L^1(G)\).
      Then the Cauchy-Schwarz inequality holds on \(L^1(G)\) with the same proof as in \ref{Petersen:313}, and the Gelfand-Naimark-Segal construction
      can be carried out exactly as in \ref{Petersen:333} to yield a nondegenerate representation \((\pi,\mathcal{H})\)
      of \(L^1(G)\) with a cyclic vector \(\xi\) such that
      \[
        \int\Phi(s)f(s)\,ds=\phi(f)=(\pi(f)\xi|\xi)
      \]  
      for each \(f\) in \(L^1(G)\). With \((u,\mathcal{H})\) as the unitary representation of \(G\)
      corresponding to \((\pi,\mathcal{H})\) and \(\{f_\lambda\}\) an approximate unit for \(L^1(G)\), we obtain
      \[
        \Phi(t)=\lim\int\Phi(s)f_\lambda(t^{-1}s)\,ds=\lim(\pi(\delta_t\times f_\lambda)\xi|\xi)=(u_t\xi|\xi)
      \]
      for every \(t\) in \(G\).
  \end{itemize}
\end{proof}
\begin{proposition}\label{Petersen:7110}
  The set \(B_{sa}(G)\) is a subalgebra of \(C^b(G)\) and is a (real) Banach algebra in the norm inherited from \(C^*(G)^*\).
\end{proposition}
\begin{proof}
  We claim that \(\Phi\Psi \in B_+(G)\) for all \(\Phi\) and \(\Psi\) in \(B_+(G)\). By \ref{Petersen:719} it
  suffices to show that the pointwise product of two positive matrices \((a_{ij})\) and
  \((b_{ij} )\) is again positive. To do so, note that since \((b_{ij}) \geq 0\), there is a matrix
  \((c_{ij} )\) such that \(b_{ij} = \sum c_{ik}\overline{c_{jk}}\) for all \(i\) and \(j\). Consequently, for any set \(\{\gamma_{i}\}\) of
  complex numbers, we have
  \[
    \sum(a_{ij}b_{ij})\overline{\gamma_{j}}\gamma_i=\sum a_{ij}(c_{ik}\overline{c_{jk}})\overline{\gamma_{j}}\gamma_i=\sum a_{ij}(\overline{c_{jk}}\overline{\gamma_{j}})(c_{ik}\gamma_i)\geq 0
  \]
  because \((a_{ij})\geq 0\). Thus \((a_{ij}b_{ij})\geq 0\), as desired.\par
  From the first part of the proof it follows that \(B_{sa}(G)\) is a subalgebra of \(C^b(G)\).
  Since \(B_{sa}\) as a vector space is isomorphic to \(C^*(G)^*_{sa}\), it is a Banach space 
  in the norm inherited from \(C^*(G)^*\). It only remains to show that the product is submultiplicative.
  This is obvious for positive definite functions since they attain their norm at \(e\).
  If now \(\Phi\) and \(\Psi\) belong to \(B_{sa}(G)\), then let \(\Phi=\Phi_+-\Phi_-\)
  and \(\Psi = \Psi_+-\Psi_-\) be their decomposition into positive definite functions 
  corresponding to the decomposition of the functionals (see \ref{Petersen:325}). Then
  \begin{align*}
    \|\Phi\Psi\|&\leq \|\Phi_+\Psi_+\|+\|\Phi_+\Psi_-\|+\|\Phi_-\Psi_+\|+\|\Phi_-\Psi_-\|\\
    &=(\|\Phi_+\|+\|\Phi_-\|)(\|\Psi_+\|+\|\Psi_-\|)=\|\Phi\|\|\Psi\|.
  \end{align*}
\end{proof}
\begin{proposition}\label{Petersen:7111}
  The isomorphism \(\phi\rightarrow\Phi\) defined in \ref{Petersen:718} is a homeomorphism of the 
  state space of \(C^*(G)\) in the weak-* topology onto the set of positive definite functions \(\Phi\) 
  with \(\Phi(e)=1\), equipped with the topology of uniform convergence on 
  compact subsets of \(G\).
\end{proposition}
\begin{proof}
  If \(\{\Phi_\lambda\}\) is a net in \(B_+(G)\) converging to a function \(\Phi\) uniformly on 
  compact subsets of \(G\), then \(\Phi\in B_+(G)\) by \ref{Petersen:719}. Let \(\phi_\lambda\) and \(\phi\) denote the positive
  functionals on \(C^*(G)\) associated with \(\Phi_\lambda\) and \(\Phi\), respectively. For every \(f\) in
  \(L^1(G)\), we have
  \[
    \phi_\lambda(f)=\int f(s)\Phi_\lambda(s)\,ds\rightarrow\int f(s)\Phi(s)\,ds=\phi(f).
  \] 
  Thus \(\phi_\lambda\rightarrow\phi\) weakly on a dense set of \(C^*(G)\).
  Since the net \(\{\phi_\lambda\}\) is bounded (\(\|\phi_\lambda\|=\Phi_\lambda(e)\)), we conclude that \(\phi_\lambda\rightarrow\phi\) weak-* in \(C^*(G)^*\).\par
  Conversely, assume that \(\{\phi_\lambda\}\) is a weak-* convergent net of states of \(C^*(G)\)
  with limit \(\phi\) where \(\|\phi\|=1\), and denote by \(\Phi_\lambda\) and \(\Phi\) the corresponding positive
  definite functions. Given \(\epsilon >0\), there is an \(f\) in \(C_C(G)_+\) with \(\|f\|_1 = 1\) such
  that \(\phi((\delta_e - f)^* \times(\delta_e -f))<\epsilon^2\). Since \(\phi_\lambda \rightarrow \phi\) weak-* and 
  \(\|\phi_\lambda\|=\|\phi\|=1\) (\(=\phi_\lambda(\delta_e)\)), there is a \(\lambda_0\) such that
  \[
    \phi_\lambda((\delta_e-f)^*\times(\delta_e-f))<\epsilon^2
  \]
  for all \(\lambda>\lambda_0\). Thus, for each \(t\) in \(G\), we have
  \begin{align*}
    |\Phi(t)-(\Phi\times \tilde{f})(t)|^2&=\vert \Phi(t)-\int\Phi(s)\tilde{f}(s^{-1}t)\,ds\vert^2\\
    &=\vert\Phi(t)-\int\Phi(s) f(t^{-1}s)\,ds\vert^2=|\phi(\delta_t)-\phi(\delta_t\times f)|^2\\
    &\leq 1\cdot\phi((\delta_e-f)^*\times(\delta_e-f))<\epsilon^2.
  \end{align*}
  Similarly,
  \[
    |\Phi_\lambda(t)-(\Phi_\lambda\times\tilde{f})(t)|<\epsilon.
  \]
  Since
  \[
    (\Phi_\lambda\times\tilde{f})(t)-(\Phi\times\tilde{f})(t)=(\phi_\lambda-\phi)(\delta_t\times f),
  \]
  and \(\delta_t\times f\) depends continuously on \(t\) in \(C^*(G)\), by a standard argument,
  for every compact set \(C\), we can find \(\lambda_1>\lambda_0\) such that
  \[
    |(\Phi_\lambda\times\tilde{f})(t)-(\Phi\times\tilde{f})(t)|<\epsilon
  \]
  for all \(t\) in \(C\). Combining this with our previous results, we obtain \(|\Phi_\lambda(t)-\Phi(t)|<3\epsilon\) for all \(t\) in \(C\).
\end{proof}

\section{The regular representation}
Let \(L^2(G)\) be the Hilbert space of (classes of) square-integrable functions with
respect to Haar measure. It is well known that, for all \(\mu\) in \(M(G)\) and \(f\) in
\(L^2(G)\), we can define the convolution \(\mu \times f\) in \(L^2(G)\). It is straightforward to
check that the map
\[
  \lambda:M(G)\rightarrow\mathcal{B}(L^2(G))
\]
given by \(\lambda(\mu)f=\mu\times f\) is a representation of \(M(G)\) called the \textit{(left) regular representation}.
The regular representation is faithful, and its restriction to \(G\)
(identifying each \(s\) in \(G\) with \(\delta_s\) in \(M(G)\)) is the unitary representation \(s \mapsto \lambda_s\)
of \(G\) on \(L^2(G)\) given by
\[
  (\lambda_s f)(t)=f(s^{-1}t),\qquad f\in L^2(G).
\]
\begin{definition}
  The \textnormal{reduced group \cstar} is the norm closure of \(\lambda(L^1(G))\) in \(\mathcal{B}(L^2(G))\)
  and is denoted by \(C_r^*(G)\).
\end{definition}

Since each nondegenerate representation of \(L^1(G)\)
extends to a representation of \(C^*(G)\) (by \ref{Petersen:715}), we have \(C^*_r(G) = \lambda(C^*(G))\).
Even though \(\lambda\) is injective on \(L^1(G)\), it is not in general an isomorphism of
\(C^*(G)\). Since \(C^*_r(G)\) is a good deal more accessible than \(C^*(G)\), we devote the
next section to a discussion of those groups for which \(C^*(G)\) and \(C^*_r(G)\) are
equal.
\begin{definition}
  The \textnormal{group von Neumann algebra} is the weak closure of \(\lambda(L^1(G))\) in
  \(\mathcal{B}(L^2(G))\) and is denoted by \(\mathcal{M}(G)\). Thus \(\mathcal{M}(G)=C^*_r(G)''\).
\end{definition}
It follows from \ref{Petersen:714} and \ref{Petersen:715} that \(\mathcal{M}(G)\) is also the von Neumann algebra
generated by \(\lambda(M(G))\) or by \(\lambda(G)\), irrespectively.\par
The set of functionals on \(C^*(G)\) of the form
\[
  x\mapsto\phi(x)=\sum(\lambda(x)f_n|g_n),
\]
where \(\{f_n\}\subset L^2(G)\), \(\{g_n\}\subset L^2(G)\), \(\sum\|f_n\|_2^2<\infty\), and \(\sum\|g_n\|_2^2<\infty\),
is a norm-closed subspace of \(C^*(G)^*\) and can be identified with \(\mathcal{M}(G)_*\).
We denote by \(A(G)\) the subspace of \(B(G)\) consisting of functions \(\Phi\) whose associated functionals belong to \(\mathcal{M}(G)_*\).
Thus \(\Phi\in A(G)\) if \(\Phi=\sum f_n\times\tilde{g_n}\) with \(\{f_n\}\) and \(\{g_n\}\) as above.
Note that the series is uniformly convergent since \(\|f_n\times\tilde{g_n}\|_\infty\leq\|f_n\|_2\|g_n\|_2\).\par
\(A(G)\) is a closed subspace of \(B(G)\) in the norm inherited from \(C^*(G)^*\) and is spanned by its positive cone, denoted by \(A_+(G)\).
From the above we see that \(\Phi\in A_+(G)\) if and only if there is a sequence \(\{f_n\}\) in \(L^2(G)\)
with \(\sum\|f_n\|_2^2<\infty\) such that \(\Phi=\sum f_n\times\tilde{f_n}\). Since \(C_C(G)\) is dense in \(L^2(G)\), we see that 
each \(\Phi\) in \(A_+(G)\) can be approximated in the norm inherited from \(C^*(G)^*\) by finite sums of positive definite
functions of the form \(f\times\tilde{f}\), \(f\in C_C(G)\).\label{Petersen:722}\par
Define the map \(j\) on the functions on \(G\) by
\[
  (jf)(s)=f^*(s)\Delta(s)^{1/2}=\overline{f}(s^{-1})\Delta(s)^{-1/2}.
\]
Elementary computations show that \(j\) is an involution, and that
\[
  (jf|jg)=(g|f)\qquad\textnormal{and}\qquad j(f\times g)=(jg)\times(jf)
\]
whenever the operations make sense. In particular, \(j\) is an isometry on \(L^2(G)\).\par
For each function \(\Phi\) on \(G\) such that \(j\Phi\in L^1(G)\), we define the operator \(\rho(\Phi)\) on \(L^2(G)\) by
\[
  \rho(\Phi)f=j(\lambda(j\Phi))jf=j(j\Phi\times jf)=f\times\Phi.
\]
The map \(\Phi\mapsto\rho(\Phi)\) (the right regular representation) is defined on \(jL^1(G)\),
which contains \(C_C(G)\), and is a linear antihomomorphism of \(jL^1(G)\) into \(\mathcal{M}(G)'\).
If \(f,g\) belong to \(C_C(G)\), then we have
\[
  (\rho(\Phi)f|g)=(f\times\Phi|g)=(f|g\times\tilde{\Phi})=(f|\rho(\tilde{\Phi})g).
\]
It follows that \(\rho(\tilde{\Phi})=\rho(\Phi)^*\), and from \ref{Petersen:719}(4) we see that if \(\Phi\in C^b(G)\),
then \(\rho(\Phi)\geq0\) if and only if \(\Phi\) is positive definite.\label{Petersen:723}\par
It is convenient to define \(\lambda(f)\) and \(\rho(f)\) also in some cases where \(f\) or \(jf\)
does not belong to \(L^1(G)\). We say that a function \(f\) on \(G\) is \textit{left} (respectively,
\textit{right}) \textit{bounded} if there is a bounded operator denoted by \(\lambda(f)\) (respectively,
\(\rho(f)\))on \(L^2(G)\) such that \(\lambda(f)g = f \times g\) (respectively, \(\rho(f)g= g \times f\))for
every \(g\) in \(C_C(G)\). Note that if \(f\) is left bounded and is right bounded, then
\(\lambda(f)\) and \(\rho(\Phi)\) commute with every operator \(\rho(g)\) and \(\lambda(g)\), \(g \in C_C(G)\), re
spectively.
\begin{lemma}\label{Petersen:724}
  If \(\Phi\in B_+(G)\cap C_C(G)\), then there is a right bounded element \(\xi\) in \(L^2(G)\)
  such that \(\Phi=\xi\times\tilde{\xi}\).
\end{lemma}
\begin{proof}
  Choose an approximate unit \(\{f_\lambda\}\) for \(L^1(G)\) contained in \(C_C(G)_+\). From
  \ref{Petersen:723} we know that \(\rho(\Phi)\geq0\), so that we may define \(\xi_\lambda =\rho(\Phi)^{1/2}f_\lambda\) in \(L^2(G)\).
  Now
  \[
    \|\xi_\lambda-\xi_\mu\|_2^2=(\rho(\Phi)(f_\lambda-f_\mu)|f_\lambda-f_\mu)=
    ((f_\lambda-f_\mu)\times\Phi\times(\tilde{f_\lambda}-\tilde{f_\mu}))(e).
  \]
  which tends to zero as \(\lambda,\mu\rightarrow\infty\). There is therefore a vector \(\xi\)
  in \(L^2(G)\) such that \(\xi_\lambda\rightarrow\xi\). If \(f\in C_C(G)\), then since \(\rho(\Phi)\in\mathcal{M}(G)'\),
  \[  
    f\times\xi=\lim\lambda(f)\rho(\Phi)^{1/2}f_\lambda=\lim\rho(\Phi)^{1/2}\lambda(f)f_\lambda=\rho(\Phi)^{1/2}f,
  \]
  which proves that \(\xi\) is right bounded with \(\rho(\xi)=\rho(\Phi)^{1/2}\).
  Moreover, for each \(t\) in \(G\),
  \begin{align*}
    (\xi\times\xi)(t)&=(\xi|\lambda_t\xi)=\lim(\rho(\Phi)^{1/2}f_\lambda|\lambda_t\rho(\Phi)^{1/2}f_\lambda)\\
    &=\lim(f_\lambda\times\Phi|\lambda_t f_\lambda)=\Phi(t),
  \end{align*}
  as desired.
\end{proof}
\begin{proposition}\label{Petersen:725}
  The set \(A_+(G)\) is the closure (in \(C^*(G)^*\)) of \(B_+(G)\cap C_C(G)\).
\end{proposition}
\begin{proof}
  From \ref{Petersen:724} we see that \(B_+(G)\cap C_C(G)\subset A_+(G)\),
  and in \ref{Petersen:722}, we showed that each element in \(A_+(G)\) can be approximated by elements
  in \(B_+(G)\cap C_C(G)\), so that \(B_+(G)\cap C_C(G)\) is dense in \(A_+(G)\).
\end{proof}

\section{Amenable groups}
We need only one result (\ref{Petersen:739}) from the theory of amenable groups, 
unfortunately, a rather deep result. Since the theory is not (yet) a standard part of
harmonic analysis, we give a self-contained exposition in this section. First, we
give some definitions and explanatory remarks.
\begin{definition}
  If \(f\) is a function on a locally compact group \(G\) and \(s\in G\), then we denote by \(\lambda_s f\)
  the function \(t\mapsto f(s^{-1}t)\). We say that \(f\) is \textit{left uniformly continuous}
  if \(\|\lambda_s f-f\|_\infty\rightarrow 0\) as \(s\rightarrow e\) and denote by \(UC_l^b(G)\) the set of left uniformly 
  continuous elements of \(C^b(G)\). Similarly, we denote by \(UC_r^b(G)\) the set of 
  functions \(f\) in \(C^b(G)\) such that \(\tilde{f}\in UC_l^b(G)\) and write
  \[
    UC^b(G)=UC_l^b(G)\cap UC_r^b(G).
  \]
\end{definition}
We then have an ascending chain of commutative \cstars\ in the \(\|\cdot\|_\infty\) norm
\[
  UC^b(G)\subset UC_l^b(G)\subset C^b(G)\subset L^\infty(G),
\]
and it is easily verified that each of these algebras is invariant under left translation by elements in \(G\). Moreover, we have
\[
  L^1(G)\times L^\infty(G)\subset UC^b_l(G),\qquad L^\infty(G)\times L^1(G)^\thicksim \subset UC^b_r(G).
\]
\begin{definition}
  If \(A\) is a \cstar\ in \(L^\infty(G)\), invariant under left translation, then we say that the 
  state \(m\) on \(A\) is a \textnormal{left invariant mean} if \(m(\lambda_s f)=m(f)\) for every \(f\in A\).
  We say that \(G\) is \textnormal{amenable} if any of the conditions in the next proposition are fulfilled.
\end{definition}
\begin{proposition}\label{Petersen:734}
  The following conditions on a locally compact group \(G\) are equivalent:
  \begin{enumerate}
    \item There is a left invariant mean on \(UC^b(G)\);
    \item There is a left invariant mean on \(C^b(G)\);
    \item There is a left invariant mean on \(L^\infty(G)\);
    \item There is a state \(m\) on \(L^\infty(G)\) such that
      \[
        m(\mu\times f)=\mu(G)m(f)
      \]
      for all \(\mu\) in \(M(G)\) and \(f\) in \(L^\infty(G)\).
  \end{enumerate}
\end{proposition}
\begin{proof}
  It is trivial that \((4) \Rightarrow (3) \Rightarrow (2) \Rightarrow (1)\). We must prove that \((1) \Rightarrow (4)\).
  Assume therefore that \(m\) is a left invariant mean on \(UC^b(G)\). We claim that
  \(m(g \times f)= \int g(s)\,ds m(f)\) for all \(g\) in \(L^1(G)\) and \(f\) in \(UC^b(G)\). To see this,
  note that the functional \(g \mapsto m(g \times f)\) is bounded on \(L^1(G)\) (since 
  \(\|g \times f\|\leq \|g\|_1\|f\|_\infty\)) and invariant under left translation. Since the Haar measure is the
  unique functional satisfying this condition, there is a scalar \(\alpha\) depending only on
  \(f\) such that \(m(g \times f)= \int g(s)\,ds\alpha\). Take an approximate unit \(\{g_\lambda\}\) for \(L^1(G)\).  
  then
  \[
    \lim_{\lambda\rightarrow\infty}\|g_\lambda\times f - f\|_\infty=0
  \]
  because \(f\in UC_l^b(G)\), whence
  \[
    \alpha=m(g_\lambda\times f)\rightarrow m(f),
  \]
  as desired.\par
  Now fix an element \(g\) in \(L^1(G)\) with \(g\geq 0\) and \(\|g\|_1=1\) and define \(\overline{m}\) on \(L^\infty(G)\) by
  \[
    \overline{m}(f)=m(g\times f\times \tilde{g}),\qquad\forall f\in L^\infty(G).
  \]
  Note that \(g\times f\times \tilde{g}\in UC^b(G)\), so that the definition is meaningful.
  Clearly, \(\overline{m}\) is a positive functional on \(L^\infty(G)\), and since \(g\times 1=1=1\times\tilde{g}\), we see that
  \(\overline{m}\) is a state. With \(\{g_i\}\) as an approximate unit for \(L^1(G)\),
  for any \(\mu\) in \(M(G)\) and \(f\) in \(L^\infty(G)\), we have
  \begin{align*}
    \overline{m}(\mu\times f)&=m(g\times \mu\times f\times\tilde{g})\\
    &=\lim_{ij}m(g\times \mu\times g_i\times g_j\times f\times\tilde{g})\\
    &=\lim_{ij}m(\mu\times g_i\times g_j\times f\times\tilde{g})\\
    &=\lim_{ij}\int\mu\times g_i(s)\,ds m(g_j\times f\times\tilde{g})\\
    &=\mu(G)\lim_{j}m(g\times g_j\times f\times\tilde{g})\\
    &=\mu(G)m(g\times f\times\tilde{g})=\mu(G)\overline{m}(f),
  \end{align*}
  and the proof is complete.
\end{proof}

Every abelian group is amenable, and every compact group is amenable (with
Haar measure as the unique invariant mean). Every closed subgroup of an
amenable group is amenable. In the converse direction, if \(H\) is a closed 
normal subgroup of \(G\) such that \(H\) and \(G/H\) are amenable, then \(G\) is amenable.
Moreover, if \(G\) is the union of an increasing net of amenable groups, then it is
itself amenable.\par
The set of functions \(f\) in \(L^1(G)\) such that \(f \geq 0\) and \(\|f\|_1 =1\) play an important
role in the following. We denote it by \(S(G)\) and note that it can be identified with
the set of normal states on \(L^\infty(G)\) and that it is a semigroup under convolution.
\begin{proposition}\label{Petersen:737}
  The following conditions on \(G\) are equivalent:
  \begin{enumerate}
    \item \(G\) is amenable;
    \item There is a net \(\{g_i\}\) in \(S(G)\) such that \(h\times g_i-g_i\rightarrow0\) weak-* in
      \((L^\infty(G))^*\) for each \(h\) in \(S(G)\);
    \item There is a net \(\{g_i\}\) in \(S(G)\) such that \(\|h\times g_i-g_i\|_1\rightarrow 0\) for each \(h\) in \(S(G)\);
    \item For each compact set \(C\) and \(\epsilon>0\), there is a \(g\) in \(S(G)\) such that \(\|\lambda_s g-g\|_1<\epsilon\) for every \(s\) in \(C\).
  \end{enumerate}  
\end{proposition}
\begin{proof}  
  \begin{itemize}
    \item \((1)\Rightarrow(2)\). Take a state \(m\) on \(L^\infty(G)\) satisfying \ref{Petersen:734} (4). Since the set
      of normal states of \(L^\infty(G)\) is weak-* dense in the set of all states (being convex
      and separating), there is a net \(\{g_i\}\) in \(S(G)\) such that \((f|g_i) \rightarrow m(f)\) for every
      \(f\) in \(L^\infty(G)\). Fix \(h\) in \(S(G)\). Then
      \[
        (f|h \times g_i)=(h^* \times f|g_i)\rightarrow m(h^* \times f)=m(f). 
      \]
      Consequently, the net \(\{h\times g_i-g_i\}\) is weak-* convergent to zero in \((L^\infty(G))^*\).
    \item \((2)\Rightarrow(3)\). For each finite set \(\{h_1,\ldots,h_n\}\) in \(S(G)\) and each \(\epsilon>0\), let
      \(X\) be the \(n\)-fold direct sum of \(L^1(G)\). The set in \(X\) with elements
      \[
        \bigoplus_k(h_k\times g-g),\qquad g\in S(G),
      \]
      is convex, and by assumption it contains zero as a weak limit point. It follows from the Hahn-Banach theorem that the set contains zero as a limit point in norm.
      There is therefore a \(g\) in \(S(G)\) such that
      \[
        \sup_k\|h_k\times g-g\|_1<\epsilon,
      \]
      from which the existence of a net \(\{g_i\}\) satisfying (3) is immediate.
    \item \((3)\Rightarrow(4)\). Fix an element \(h\) in \(S(G)\) and take a net \(\{g_i\}\) in \(S(G)\) 
      satisfying (3). Then, for each \(s\) in \(G\),
      \[
        \|\lambda_s(h\times g_i)-h\times g_i\|_1\leq\|\lambda_s(h)\times g_i-g_i\|+\|h\times g_i-g_i\|\rightarrow 0.
      \]
      Since \(\lambda_s(h)\) depends continuously on \(s\) in \(L^1(G)\), by a standard argument we can,
      for each compact set \(C\) and \(\epsilon>0\), find \(g_i\) such that
      \[
        \|\lambda_s(h)\times g_i-g_i\|_1+\|h\times g_i-g_i\|_1<\epsilon
      \]
      for all \(s\) in \(C\). Taking \(g=h\times g_i\), we obtain \(\|\lambda_s g-g\|<\epsilon\) for all \(s\) in \(C\).
    \item \((4)\Rightarrow(1)\). Let \(I\) be the net of pairs \((C,\epsilon)\) consisting of a compact set \(C\)
      and \(\epsilon>0\). For each \(i\) in \(I\), choose \(g_i\) in \(S(G)\) satisfying (4) and consider the net \(\{m_i\}\)
      of states on \(C^b(G)\) given by \(m_i(f)=\int f(t)g_i(t)\,dt\). For each \(s\) in \(C\), we have
      \begin{align*}
        |m_i(\lambda_{s^{-1}}f-f)|&=\vert\int(f(st)-f(t))g_i(t)\,dt\vert\\
        &=\vert\int f(t)(g_i(s^{-1}t)-g_i(t))\,dt\vert\leq\|f\|_\infty\epsilon.
      \end{align*}
      Thus every weak-* limit point of \(\{m_i\}\) in \(C^b(G)^*\) is a left invariant mean on \(G\).
  \end{itemize}
\end{proof}
\begin{proposition}\label{Petersen:738}
  The group \(G\) is amenable if and only if there is a net \(\{f_i\}\) in the unit sphere of 
  \(L^2(G)\) such that \(\{f_i\times \tilde{f_i}\}\) converges to \(1\) uniformly on compact subsets of \(G\).
\end{proposition}
\begin{proof}
  If \(g\in S(G)\), then put \(f=g^{1/2}\) (pointwise square root). Then \(1=\|f\|^2_2=(f\times\tilde{f})(e)\). Moreover, for each \(s\) in \(G\),
  \begin{align*}
    |1-(f\times\tilde{f})(s)|^2&=|(f\times\tilde{f})(e)-(f\times\tilde{f})(s)|^2\\
    &=\left| \int f(t)(\tilde{f}(t^{-1})-\tilde{f}(t^{-1}s))\,dt\right|^2\\
    &=\left|\int f(t)(\overline{f}(t)-\tilde{f}(t^{-1}s))\,dt\right|^2\\
    &=|(f|f-\lambda_s f)|^2\leq\|f-\lambda_s f\|_2^2\\
    &=\int|g^{1/2}(t)-g^{1/2}(s^{-1}t)|^2\,dt\\
    &\leq\int|g(t)-g(s^{-1}t)|\,dt=\|g-\lambda_s g\|_1.
  \end{align*}
  Conversely, if \(f\) belongs to the unit sphere of \(L^2(G)\), then put \(g=|f|^2\). 
  Then \(g\in S(G)\). Moreover, for each \(s\) in \(G\),
  \begin{align*}
    \|g-\lambda_s g\|_1&=\int\left||f(t)|^2-|f(s^{-1}t)|^2\right|\,dt\\
    &\leq\int\bigl(|f(t)|+|f(s^{-1}t)|\bigr)|f(t)-f(s^{-1}t)|\,dt\\
    &\leq\bigl(\|f\|_2+\|\lambda_s f\|_2\bigr)\|f-\lambda_s f\|_2\\
    &=2\|f-\lambda_s f\|_2=2(2-(f|\lambda_s f)-(\lambda_s f|f))^{1/2}\\
    &\leq2\sqrt{2}|1-(\lambda_s f|f)|^{1/2}=2\sqrt{2}|1-(f\times \tilde{f})(s)|^{1/2}.
  \end{align*}
  Here the first inequality uses
  \(\bigl||z|^2-|w|^2\bigr|\leq(|z|+|w|)|z-w|\), and the second is Cauchy--Schwarz;
  this is the standard \(L^1\)--\(L^2\) estimate in Reiter's characterization of amenability
  (see \cite{Greenleaf1969}).\footnote{At this point \cite[Proposition~7.3.8]{PGK79} prints an equality
  with \(|f(t)+f(s^{-1}t)|\). That equality is false for general complex-valued \(f\): the valid relation is the
  inequality displayed above. The corrected estimate is sufficient for the argument.}
  Combining these two inequalities, we obtain
  \[
    |1-(f\times\tilde{f})(s)|^2\leq\|g-\lambda_s g\|_1\leq2\sqrt{2}|1-(f\times\tilde{f})(s)|^{1/2},
  \]
  and the result follows immediately from \ref{Petersen:737} (4).
\end{proof}
\begin{theorem}\label{Petersen:739}
  The following conditions are equivalent:
  \begin{enumerate}
    \item \(G\) is amenable;
    \item The regular representation is faithful on \(C^*(G)\).
  \end{enumerate}
\end{theorem}
\begin{proof}
  If \(G\) is amenable, then by \ref{Petersen:738} there is a net \(\{\Phi_i\}\)
  in \(A_+(G)\) with \(\Phi_i(e)=1\) such that \(\{\Phi_i\}\) converges to \(1\)
  uniformly on compact subsets of \(G\). Since \(C_C(G)\cap A_+(G)\) is dense in \(A_+(G)\)
  in the norm topology inherited from \(C^*(G)^*\) by \ref{Petersen:725}, it is a fortiori
  dense in \(A_+(G)\) in the uniform norm. We may therefore assume that \(\{\Phi_i\}\subset C_C(G)\).
  For any \(\Phi\) in \(B_+(G)\) with \(\Phi(e)=1\), we have \(\Phi\Phi_i\in C_C(G)\cap B_+(G)\)
  by \ref{Petersen:7110}, and \(\Phi\Phi_i\rightarrow \Phi\) uniformly on compact subsets of \(G\). It follows from \ref{Petersen:7111}
  that \(\mathcal{M}(G)_*\) is weak-* dense in \(C^*(G)^*\) and thus the regular representation has zero kernel in \(C^*(G)\)\par
  Conversely, if the regular representation is faithful on \(C^*(G)\), then \(\mathcal{M}(G)_*\)
  is weak-* dense in \(C^*(G)^*\), whence \(A_+(G)\) is dense in \(B_+(G)\) in the topology
  of uniform convergence on compact subsets of \(G\) by \ref{Petersen:7111}. The function \(\Phi_1\)
  given by \(\Phi_1(t) = 1\) for all \(t\) belongs to \(B_+(G)\) (and corresponds to the trivial
  representation of \(G\)). Consequently, there is a net \(\{\Phi_i\}\) in \(A_+(G)\) with \(\Phi_i(e) = 1\)
  such that \(\Phi_i \rightarrow 1(= \Phi_1)\) uniformly on compact subsets of \(G\). Moreover, by
  \ref{Petersen:725} we may assume that \(\{\Phi_i\}\subset C_C(G)\). For each \(i\), by \ref{Petersen:724} there is a vector
  \(\xi_i\) in \(L^2(G)\) such that \(\Phi_i =\xi_i \times \tilde{\xi_i}\), whence \(G\) is amenable by \ref{Petersen:738}.
\end{proof}

\section{The group \texorpdfstring{\(\mathbb{F}_2\)}{F2}}
In this section, we investigate the \cstars\ associated with the free group
on two generators, denoted \(\mathbb{F}_2\). Here, the key point is that \cstar\ is
plural: the full and reduced \cstars\ are different. This is because \(\mathbb{F}_2\) is
not amenable. Rather than go too far into that, we will mainly concentrate
on just how different these \cstars\ are.\par
We will begin with a brief review of the group \(\mathbb{F}_2\). We denote the 
generators by \(a\) and \(b\). A word in \(a,b,a^{-1},b^{-1}\) is simply a finite string of these
symbols written sequentially. There is also a special word e which is the
empty word consisting of no symbols. Two words may be concatenated by
simply writing them side by side. Of course, the order matters. We denote
the concatenation of \(w_1\) and \(w_2\) by \(w_1 w_2\). Notice that \(ew = w = we\) for all \(w\)
and that concatenation is an associative operation. We define an equivalence
relation on the set of all words by first setting
\[
  w_1 w_2 \thicksim w_1 a^{-1}aw_2 \thicksim w_1 aa^{-1}w_2 \thicksim w_1 b^{-1}bw_2 \thicksim w_1 bb^{-1}w_ 2,
\]
for all words \(w_1,w_2\) and then letting \(\thicksim\) denote the transitive closure of this
relation. We also say that a word is reduced if it contains no pair \(a,a^{-1}\) or
\(b, b^{-1}\) adjacent. Each \(\thicksim\)-equivalence class contains a unique reduced word.
The collection of \(\thicksim\)-equivalence classes of words, with concatenation is our
group \(\mathbb{F}_2\).\par
A fundamental property of the group (in fact, it is really the definition) is that,
given any group \(G\) with two elements \(g,h\) in \(G\), there is a unique group
homomorphism \(\alpha : \mathbb{F}_2 \rightarrow G\) such that \(\alpha(a) = g\) and \(\alpha(b) = h\).\par
As we said above, our aim is to show that the \cstars\ \(C^*(\mathbb{F}_2)\) and
\(C^*_r(\mathbb{F}_2)\) are very different. Compare them with the following two theorems.
\begin{theorem}
  If \(A\) is a unital \cstar\ containing unitaries \(u,v\), then
  there exists a unique unital *-homomorphism \(\rho : C^*(\mathbb{F}_2) \rightarrow A\) such that
  \(\rho(a) = u,\rho(b) = v\). In particular, for every \(n \geq 1\), there is a surjective
  *-homomorphism \(\rho : C^*(\mathbb{F}_2) \rightarrow M_n(\mathbb{C})\).
\end{theorem}
\begin{proof}
  Choose \(\pi:A\rightarrow \mathcal{B}(\mathcal{H})\) which is
  a faithful, unital representation. Then \(\pi(u)\) and \(\pi(v)\) are unitaries. That
  means there is a group homomorphism \(\alpha : \mathbb{F}_2\rightarrow \mathcal{U}(\mathcal{H})\) which maps \(a\) to
  \(\pi(u)\) and \(b\) to \(\pi(v)\). This means that \(\alpha\) is a unitary representation of \(\mathbb{F}_2\) and
  extends to \(\beta\), a unital representation of \(\mathbb{C}\mathbb{F}_2\) and also a unital representation
  of \(C^*(\mathbb{F}_2)\). It is clear that the image of the former is contained in the 
  *-algebra generated by \(\pi(u)\) and \(\pi(v)\) while the image of the latter is in the
  closure of this. In any case, both images are contained in \(\pi(A)\). We are done
  when we let \(\rho=\pi^{-1}\circ \beta\).
  For the second statement, it is a fairly simple matter to see that the
  smallest \cstar\ which contains
  \[
    u=
    \begin{pmatrix}
    1 &0 &\cdots &0\\
    0 &e^{2\pi i/n} & \cdots &0\\\
    \vdots &&  &\vdots\\
    0& 0& \cdots & e^{2\pi i(n-1)/n}
    \end{pmatrix}
  \]
  is the set of diagonal matrices. It follows that the smallest \cstar\ which
  contains this and
  \[
    v=
    \begin{pmatrix}
    0 &1 &0 &\cdots &0\\
    0 &0 &1 &\cdots &0\\
    \vdots & & & &\vdots\\
    0 &0 &&\cdots &1\\
    1 &0 &&\cdots &0\\
    \end{pmatrix}
  \]
  is \(M_n(\mathbb{C})\).\par
\end{proof}
  Now we turn to the reduced \cstar, which is trickier. We start with
  a fairly useful result that the trace is faithful.
\begin{theorem}\label{theorem:274}
  Let \(G\) be a discrete group. If \(c\) is in \(C^*_r(G)\) and \(\tau(c^*c)=0\),
  then \(c=0\). That is, the trace on \(C^*_r(G)\) is faithful.
\end{theorem}
\begin{proof}
  In this proof, we identify \(C^*_r(G)\) with its image under \(\pi_\lambda\). Let \(g\) be in
  \(G\) and compute
  \begin{align*}
    \|c\delta_g\|^2 &= < c\delta_g,c\delta_g >\\
    &= <c\lambda_g\delta_e,c\lambda_g\delta_e >\\
    &= <\lambda^*_gc^*c\lambda_g\delta_e,\delta_e >\\
    &= \tau(\lambda^*_g c^*c\lambda_g)\\
    &= \tau(\lambda_g\lambda^*_g c^*c)\\
    &= \tau(c^*c)\\
    &= 0.
  \end{align*}
  We see that \(c\) is zero on a basis for \(\ell^2(G)\) and hence is zero.
\end{proof}
\begin{lemma}[Powers' averaging estimate]\label{lemma:powers-averaging}
  If \(w\in\mathbb F_2\setminus\langle a\rangle\), then
  \[
    \lim_{M\to\infty}\left\|\frac1M\sum_{m=1}^M
    \lambda_{a^mwa^{-m}}\right\|=0.
  \]
  The analogous assertion holds with \(a\) and \(b\) interchanged.
\end{lemma}
\begin{proof}
  This is the standard free-word averaging estimate underlying Powers' proof of simplicity;
  see \cite{Powers1975}. It follows by partitioning the reduced words according to their initial
  letters and applying the resulting orthogonality estimate to the left regular representation.
\end{proof}
\begin{lemma}\label{lemma:277}
  For any \(c\) in \(C^*_r(\mathbb{F}_2)\), we have
  \[
    \lim_{M,N\rightarrow\infty}\frac{1}{MN} \sum^M_{m=1}\sum^N_{n=1}\lambda_{b^n a^m} c \lambda_{a^{-m}b^{-n}}=\tau(c).\]
\end{lemma}
\begin{proof}
  Each averaging map in the statement is contractive, and \(\tau\) is continuous. Since
  \(\mathbb C\mathbb F_2\) is dense in \(C_r^*(\mathbb F_2)\), it therefore suffices, by linearity,
  to consider \(c=\lambda_w\), with \(w\in\mathbb F_2\). If \(w=e\), the assertion is immediate.
  If \(w=a^k\) for some \(k\neq0\), then \(\lambda_w\) commutes with every \(\lambda_{a^m}\), and hence
  \[
    \frac{1}{MN}\sum_{m=1}^M\sum_{n=1}^N
    \lambda_{b^na^m}\lambda_w\lambda_{a^{-m}b^{-n}}
    =\frac1N\sum_{n=1}^N\lambda_{b^nwb^{-n}}\longrightarrow0
  \]
  by Lemma~\ref{lemma:powers-averaging}, with \(a\) and \(b\) interchanged. Finally, suppose that
  \(w\notin\langle a\rangle\), and put
  \[
    B_M=\frac1M\sum_{m=1}^M\lambda_{a^mwa^{-m}}.
  \]
  Lemma~\ref{lemma:powers-averaging} gives \(\|B_M\|\to0\), while
  \[
    \left\|\frac1N\sum_{n=1}^N\lambda_{b^n}B_M\lambda_{b^{-n}}\right\|
    \leq\|B_M\|.
  \]
  This bound is uniform in \(N\), and completes the proof.
\end{proof}
\begin{theorem}\label{theorem:272}
  The \cstar\ \(C^*_r(\mathbb{F}_2)\) is simple; that is, it has no non
  trivial ideals. In particular, it has no finite-dimensional representations.
\end{theorem}
\begin{proof}
  Suppose that \(C^*_r(\mathbb{F}_2)\) has a
  non-zero ideal \(I\) and let \(c\) be some non-zero element of that ideal. We know
  from Theorem \ref{theorem:274} that \(\tau(c^*c)>0\). We also know that \(c^*c\) is also in the
  ideal. Looking at the last Lemma, the expression on the left hand side (using
  \(c^*c\) instead of \(c\)) lies in the ideal, for every value of \(M,N\). As the ideal is
  closed, \(\tau(c^*c)>0\) is in the ideal. But this element is invertible and it follows
  that \(I=C^*_r(\mathbb{F}_2)\).
\end{proof}
\begin{theorem}\label{theorem:273}
  The trace on \(C^*_r(\mathbb{F}_2)\) is unique.
\end{theorem}
\begin{proof}
  If \(\phi\) is any other trace on \(C^*_r(\mathbb{F}_2)\), if
  we apply it to \(\frac{1}{MN} \sum^M_{m=1}\sum^N_{n=1}  \lambda_{b^n a^m} c \lambda_{a^{-m}b^{-n}}\) 
  we clearly get \(\phi(c)\) because of
  the trace property. On the other hand, applying it to the scalar \(\tau(c)\) yields
  \(\tau(c)\). The conclusion follows from Lemma~\ref{lemma:277}.
\end{proof}
\begin{exercise}
  Let \(u\) be the \(N\times N\) unitary matrix
  \[
    u=
    \begin{pmatrix}
      1 &0 &\cdots &0\\
      0 &e^{2\pi i/N} &\cdots &0\\
      \vdots&&& \vdots\\
      0 &0 &\cdots &e^{2\pi i(N-1)/N}
  \end{pmatrix}
  \]
  Define a map \(E:M_N(\mathbb{C})\rightarrow M_N(\mathbb{C})\) by
  \[E(a)= \frac{1}{N} \sum^N_{n=1} u^n au^{-n},\]
  for a in \(M_N(\mathbb{C})\).
  \begin{enumerate}
    \item Prove that the range of \(E\) is contained in the \(C^*\)-subalgebra \(D\) consisting
    of diagonal matrices.
    \item Prove that \(E(d)=d\) if \(d\) is diagonal. In consequence, \(E\circ E=E\).
    \item Prove that \(E(a^*a)\) is positive for every \(a\) in \(M_N(\mathbb{C})\).
    \item Prove that E is faithful, meaning that if \(E(a^*a)=0\), for some \(a\) in
    \(M_N(\mathbb{C})\), then \(a=0\).
    \item Prove that if \(a\) in \(M_N(\mathbb{C})\) and \(c,d\) are in \(D\), then \(E(cad)=cE(a)d\).
    \item Prove that \(\|E\| =1\).
  \end{enumerate}
\end{exercise}
\begin{exercise}
  Let \(u,v\) be the \(N\times N\) unitary matrices
  \[
    u=
    \begin{pmatrix}
      1 &0 &\cdots &0\\
      0 &e^{2\pi i/N} &\cdots &0\\
      \vdots&&& \vdots\\
      0 &0 &\cdots &e^{2\pi i(N-1)/N}
    \end{pmatrix},v=
    \begin{pmatrix}
      0 &1 &0 &\cdots& 0\\
      0 &0 &1 &\cdots& 0\\
      \vdots & & & &\vdots\\
      0 &0 & & &1\\
      1 &0 & &\cdots &0
    \end{pmatrix}.
  \]
  Prove that, for any \(c\) in \(M_N(\mathbb{C})\), we have
  \[\frac{1}{N^2}\sum^N_{m=1}\sum^N_{n=1}v^n u^m cu^{-m}v^{-n}=\tau(c),\]
  where \(\tau\) is the usual trace on \(M_N(\mathbb{C})\).
\end{exercise}
\begin{exercise}
  Show that Theorem~\ref{theorem:273} is false for \(C^*(G)\).
\end{exercise}

\chapter{Crossed Products}\label{chapter:crossedproduct}
\section{Dynamical systems}
\begin{comment}
azioni di gruppi su spazi compatti;
automorfismi di C(X);
esempi.
\end{comment}
\begin{definition}
  A \textnormal{\(C^*\)-dynamical system} \((A, G, \alpha)\) consists of a \cstar\ 
  \(A\) together with a homomorphism \(\alpha\)
  of a locally compact group \(G\) into \(Aut(A)\) (for every \(a\in A\), \(s\longmapsto\alpha_s(a)\) must be norm-continuous). 
  We will denote by 
  \(\alpha_s\) the automorphism \(\alpha(s)\) for \(s\) in \(G\). 
\end{definition}
\begin{example}
  Let \(X\) be a compact Hausdorff space and let
  \(h:X\rightarrow X\)
  be a homeomorphism.
  Define an action \(\alpha_h:\mathbb Z\to\operatorname{Aut}(C(X))\) by
  \((\alpha_h)_n(f)=f\circ h^{-n}\), for \(n\in\mathbb Z\) and \(f\in C(X)\).
  Hence \((C(X),\mathbb{Z},\alpha_h)\)
  is a \(C^*\)-dynamical system.
\end{example}
\begin{example}
  Fix \(\theta\) in \(\mathbb{R}\).
  Define
  \(R_\theta(z)=e^{2\pi i\theta}z\),
  \(z\) in \(S^1\).
  Then \(R_\theta\) is a homeomorphism of \(S^1\).
  The action \(\alpha_\theta:\mathbb Z\to\operatorname{Aut}(C(S^1))\), defined by
  \((\alpha_\theta)_n(f)=f\circ R_\theta^{-n}\),
  determines the \(C^*\)-dynamical system
  \((C(S^1),\mathbb{Z},\alpha_\theta)\).
\end{example}
\begin{exercise}
  Let \((C(S^1),\mathbb{Z},\alpha_\theta)\) be as in the previous example. If \(\theta\) is irrational, then every orbit
  \(\{R_\theta^n(z):n\in\mathbb Z\}\)
  is dense in \(S^1\).
\end{exercise}
\begin{proof}
  Let
  \(H=\overline{\{e^{2\pi i n\theta}:n\in\mathbb Z\}}\).
  Since multiplication is continuous, \(H\) is a closed subgroup of \(S^1\).
  Every closed subgroup of \(S^1\) is either finite or equal to \(S^1\).
  Assume that \(H\) is finite.
  Then \(e^{2\pi i\theta}\) has finite order and therefore there exists
  \(q\ge1\) such that
  \(e^{2\pi iq\theta}=1\).
  Hence \(q\theta\in\mathbb{Z}\),
  which contradicts the irrationality of \(\theta\).
  Therefore \(H=S^1\).
  It follows that
  \(\{e^{2\pi i n\theta}:n\in\mathbb Z\}\)
  is dense in \(S^1\).\\
  Now let \(z\in S^1\).
  Since
  \(R_\theta^n(z)=ze^{2\pi i n\theta}\),
  the orbit of \(z\) is obtained from
  \(\{e^{2\pi i n\theta}:n\in\mathbb Z\}\)
  by multiplication by \(z\).
  As multiplication by \(z\) is a homeomorphism of \(S^1\),
  the orbit
  \(\{R_\theta^n(z):n\in\mathbb Z\}\)
  is also dense in \(S^1\).
\end{proof}

\section{Covariant representations}





\begin{definition}
  Given a \(C^*\)-dynamical system \((A, G, \alpha)\), a \textnormal{covariant 
  representation} is a pair \((\pi, U)\) where \(\pi\) is a *-representation of \(A\) on a Hilbert 
  space \(\mathcal{H}\) and \(s \rightarrow U_s\) is a unitary representation of \(G\) on the same space such that 
  \[
    U_s\pi(a)U^*_s = \pi(\alpha_s(a)), 
  \]
  for all \(a\) in \(A\), \(s\) in \(G\).
\end{definition}
\begin{lemma}
  Let \((\pi,U)\) be a covariant representation.
  Then
  \[
  U_s\pi(a)
  =
  \pi(\alpha_s(a))U_s
  \]
  for every \(a\) in \(A\) and \(s\) in \(G\).
\end{lemma}
\begin{proof}
  Multiply
  \[
  U_s\pi(a)U_s^*
  =
  \pi(\alpha_s(a))
  \]
  on the right by \(U_s\).
\end{proof}
\begin{example}
  Let \(X\) be a compact Hausdorff space and let
  \(h:X\to X\)
  be a homeomorphism.
  Let \(\mathcal{H}=L^2(X,\mu)\), where \(\mu\) is an \(h\)-invariant measure.
  Define
  \[(\pi(f)\xi)(x)
  =
  f(x)\xi(x)\]
  and
  \[(U\xi)(x)
  =
  \xi(h^{-1}(x)).\]
  Then
  \[U\pi(f)U^*
  =
  \pi(f\circ h^{-1}).\]
  Hence \((\pi,U)\) is a covariant representation of
  \((C(X),\mathbb Z,\alpha_h)\).
\end{example}
In these notes, we will restrict our attention to crossed products by discrete 
groups. Because the Haar measure is then just counting measure, the notions simplify in this setting. 

\section{The crossed product \texorpdfstring{\(A\rtimes_\alpha G\)}{Ax alpha G}, discrete case}
The covariance relation shows that the algebra \(A\) and
the group \(G\) may be represented simultaneously on the
same Hilbert space.
The crossed product \(A\rtimes_\alpha G\) will be defined
so as to encode all covariant representations.\par
Let \((A,G,\alpha)\) be a \(C^*\)-dynamical system and assume
that \(G\) is discrete.
We define \(A[G]\) to be the collection of all functions
\(f:G\to A\) having finite support.
Addition and scalar multiplication are defined pointwise.
For \(f,g\in A[G]\), define
\[(f*g)(t)
=
\sum_{s\in G}
f(s)\alpha_s(g(s^{-1}t))\]
and
\[f^*(t)
=
\alpha_t(f{(t^{-1})}^*).
\]
One checks that \(A[G]\) becomes a *-algebra.
\begin{proposition}
  If \(A\) is unital, the algebra \(A[G]\) contains copies of both \(A\) and \(G\).
\end{proposition}
\begin{proof}
  For \(a\) in \(A\) define
  \[
  \tilde a(e)=a,
  \qquad
  \tilde a(t)=0
  \quad(t\neq e).
  \]
  Then
  \(\tilde a*\tilde b=\widetilde{ab}\).\par
  For \(s\) in \(G\) define
  \[
  u_s(t)
  =
  \begin{cases}
  1_A & t=s,\\
  0 & t\neq s.
  \end{cases}
  \]
  Then
  \[u_s*u_t=u_{st}
  \quad\text{and}\quad
  u_s^*=u_{s^{-1}}.\]
\end{proof}
\begin{proposition}
  For every \(a\) in \(A\) and \(s\) in \(G\),
  \[u_s\tilde a u_s^*
  =
  \widetilde{\alpha_s(a)}.\]
\end{proposition}
\begin{proof}
  Since
  \(u_s^*=u_{s^{-1}}\), we have
  \[
  u_s\tilde a u_s^*
  =
  u_s\tilde a u_{s^{-1}}.
  \]
  A direct computation using the convolution product shows that
  \[
  u_s\tilde a u_{s^{-1}}
  =
  \widetilde{\alpha_s(a)}.\]
\end{proof}
\begin{lemma}
  Let \((\pi,U)\) be a covariant representation of the \(C^*\)-dynamical system \((A,G,\alpha)\). Defining \[(\pi\times U)(f)
  =
  \sum_{s\in G}
  \pi(f(s))U_s.\] 
  we obtain a *-representation of \(A[G]\). 
\end{lemma}
\begin{proof}
  Let \(f,g\) in \(A[G]\).
  Using the covariance relation
  \(U_s\pi(a)=\pi(\alpha_s(a))U_s\),
  we obtain
  \begin{align*}
  (\pi\times U)(f)(\pi\times U)(g)
  &=
  \sum_{s,t\in G}
  \pi(f(s))U_s\pi(g(t))U_t\\
  &=
  \sum_{s,t\in G}
  \pi(f(s))
  \pi(\alpha_s(g(t)))
  U_{st}\\
  &=
  \sum_{s,t\in G}
  \pi(f(s)\alpha_s(g(t)))
  U_{st}.
  \end{align*}
  Setting \(r=st\), we obtain
  \[
  (\pi\times U)(f)(\pi\times U)(g)
  =
  \sum_{r\in G}
  \pi\!\left(
  \sum_{s\in G}
  f(s)\alpha_s(g(s^{-1}r))
  \right)
  U_r.
  \]
  Hence
  \[
  (\pi\times U)(f)(\pi\times U)(g)
  =
  (\pi\times U)(f*g).
  \]
  Similarly,
  \begin{align*}
  (\pi\times U)(f)^*
  &=
  \sum_{s\in G}
  U_s^*\pi(f(s)^*)\\
  &=
  \sum_{s\in G}
  U_{s^{-1}}\pi(f(s)^*)\\
  &=
  \sum_{s\in G}
  \pi(\alpha_{s^{-1}}(f(s)^*))U_{s^{-1}}.
  \end{align*}
  Replacing \(s^{-1}\) by \(t\), we obtain
  \[
  (\pi\times U)(f)^*
  =
  \sum_{t\in G}
  \pi(\alpha_t(f(t^{-1})^*))U_t
  =
  (\pi\times U)(f^*).
  \]
  Therefore \(\pi\times U\) is a *-representation.
\end{proof}
\begin{definition}
  For \(f\) in \(A[G]\), define
  \[
  \|f\|
  =
  \sup\{\|(\pi\times U)(f)\| |
  (\pi,U)\mbox{ a covariant representation}\}.\]
\end{definition}
The supremum is finite because every covariant representation satisfies
\[
  \|(\pi\times U)(f)\|
  \leq\sum_{s\in G}\|\pi(f(s))U_s\|
  \leq\sum_{s\in G}\|f(s)\|.
\]
To see that the resulting \(C^*\)-seminorm separates points, choose a faithful nondegenerate representation
\(\pi_0:A\rightarrow\mathcal{B}(\mathcal{H})\), and define a covariant representation on \(\ell^2(G,\mathcal{H})\) by
\[
  (\widetilde{\pi}_0(a)\xi)(t)=\pi_0(\alpha_{t^{-1}}(a))\xi(t),
  \qquad
  (\lambda_s\xi)(t)=\xi(s^{-1}t).
\]
If \((\widetilde{\pi}_0\times\lambda)(f)=0\), then, for every \(s\in G\) and \(h\in\mathcal{H}\),
\[
  0=\bigl((\widetilde{\pi}_0\times\lambda)(f)(\delta_e\otimes h)\bigr)(s)
   =\pi_0(\alpha_{s^{-1}}(f(s)))h.
\]
The faithfulness of \(\pi_0\) implies that \(f(s)=0\) for every \(s\), and hence \(f=0\). Finally, since the norm is the supremum
of norms coming from *-representations,
\[
  \|f^* * f\|=\sup_{(\pi,U)}\|(\pi\times U)(f)^*(\pi\times U)(f)\|=\|f\|^2.
\]
Thus the displayed formula defines a \(C^*\)-norm on \(A[G]\); compare \cite[Lemma 2.27]{Williams2007}.
\begin{definition}
  The \textnormal{crossed product} \(A\rtimes_\alpha G\) is the completion of \(A[G]\) in the above norm.
\end{definition}

\section{The crossed product \texorpdfstring{\(A\rtimes_\alpha G\)}{Ax alpha G}, general case}
Consider now a \(C^*\)-dynamical system \((A,G,\alpha)\) (hence \(G\) is from now locally compact).
We define involution and convolution on the linear space \(C_C(G,A)\) of continuous functions from \(G\) to \(A\) with compact supports by
\begin{align*}
  f^*(t)&={\Delta(t)}^{-1}\alpha_t({f(t^{-1})}^*),\\
  (f\times g)(t)&=\int_G f(s)\alpha_s(g(s^{-1}t))\,ds
\end{align*}
for all \(f,g\) in \(C_C(G,A)\). Straightforward computations show that \(C_C(G,A)\) becomes a *-algebra with convolution
as product. For each \(f\) in \(C_C(G,A)\), define
\[
  \|f\|_1=\int\|f(t)\|\,dt.
\]
Then \(C_C(G,A)\) is a normed algebra with an isometric involution, and we denote by \(L^1(G,A)\) its completion.\par
For all \(x\) in \(A\) and \(f\) in \(L^1(G)\), we denote by \(x\otimes f\) the element in \(L^1(G,A)\) such that
\((x\otimes f)(t)=xf(t)\). Note that the linear span of elements of the form \(x\otimes f\), \(x\in A\), \(f\in C_C(G)\), is dense in \(L^1(G,A)\).\par

As in section \ref{section:centralizers}, we define left, right, and double centralizers on the Banach algebra
\(L^1(G,A)\), the only difference being that we \textit{assume} the continuity of the maps
involved. For each \(x\) in \(M(A)\) and \(\mu\) in \(M(G)\) with compact support, define the linear maps \(L(x,\mu)\)
and \(R(x,\mu)\) on \(C_C(G,A)\) by
\begin{align}\label{Petersen:762}
  (L(x,\mu)y)(t)&=x\int\alpha_s(y(s^{-1}t))\,d\mu(s),\\
  (R(x,\mu)y)(t)&=\int y(t s^{-1})\alpha_{t s^{-1}}(x){\Delta(s)}^{-1}\,d\mu(s).
\end{align}
It is immediate that \(L(x,\mu)\) and \(R(x,\mu)\) are bounded (by \(\|x\|\|\mu\|\)) and thus extend by continuity to linear operators on \(L^1(G,A)\).
Having done this, we see that the restriction that \(\mu\) has compact support is no longer necessary.
\begin{lemma}\label{Petersen:763}
  For all \(x\) in \(M(A)\) and \(\mu\) in \(M(G)\), the operator \(L(x,\mu)\) (respectively, \(R(x,\mu)\)) is a left 
  (respectively, right) centralizer of \(L^1(G,A)\), and the pair \((R(x,\mu),L(x,\mu))\) is a double centralizer.
\end{lemma}
\begin{proof}
  Take \(y\) and \(z\) in \(C_C(G,A)\) and assume that \(\mu\) has compact support. Then
  \begin{align*}
    L(x,\mu)(y\times z)(t)&=x\int\alpha_s((y\times z)(s^{-1}t))\,d\mu(s)\\
    &=x\int\int\alpha_s(y(r)\alpha_r(z(r^{-1}s^{-1}t)))\,d\mu(s)\,dr\\
    &=x\int\int\alpha_s(y(s^{-1}r))\,d\mu(s)\alpha_r(z(r^{-1}t))\,dr\\
    &=(L(x,\mu)y\times z)(t).
  \end{align*}
  Furthermore,
  \begin{align*}
    R(x,\mu)(y\times z)(t)&=\int(y\times z)(t s^{-1})\alpha_{t s^{-1}}(x){\Delta(s^{-1})}\,d\mu(s)\\
    &=\int\int y(r)\alpha_r(z(r^{-1}t s^{-1}))\alpha_{t s^{-1}}(x)\Delta(s^{-1})\,d\mu(s)\,dr\\
    &=\int y(r)\alpha_r(\int z(r^{-1}t s^{-1})\alpha_{r^{-1}t s^{-1}}(x)\Delta(s^{-1})\,d\mu(s))\,dr\\
    &=(y\times R(x,\mu)z)(t).
  \end{align*}
  Thus \(L(x,\mu)\) and \(R(x,\mu)\) are centralizers. To show that the pair \((R(x,\mu),L(x,\mu))\) is a double centralizer, we compute
  \begin{align*}
    (R(x,\mu)y\times z)(t)&=\int R(x,\mu)y(s)\alpha_s(z(s^{-1}t))\,ds\\
    &=\int\int y(s r^{-1})\alpha_{s r^{-1}}(x)\Delta(r^{-1})(\alpha_s(z(s^{-1}t)))\,d\mu(r)\,ds\\
    &=\int\int y(s)\alpha_s(x)\alpha_{sr}(z(r^{-1}s^{-1}t))\,d\mu(r)\,ds\\
    &=(y\times L(x,\mu)z)(t).
  \end{align*}
\end{proof}

\begin{proposition}\label{Petersen:764}
  If \((\pi,u,\mathcal{H})\) is a covariant representation of \((A,G,\alpha)\), then there is a nondegenerate representation \((\pi\times u,\mathcal{H})\) of \(L^1(G,A)\) such that
  \begin{equation}\label{eq:Pet764}
    (\pi\times u)(y)=\int\pi(y(t))u_t\,dt
  \end{equation}
  for every \(y\) in \(C_C(G,A)\). Moreover, the correspondence \((\pi,u,\mathcal{H})\mapsto (\pi\times u,\mathcal{H})\)
  is a bijection onto the set of nondegenerate representations of \(L^1(G,A)\).
\end{proposition}
\begin{proof}
  Take \(y\) in \(C_C(G,A)\) and define \((\pi\times u)(y)\) on \(\mathcal{H}\) by (\ref{eq:Pet764}). Then
  \begin{align*}
    (\pi\times u)(y^*)&=\int\pi({\Delta(t)}^{-1}\alpha_t({y(t^{-1})}^*))u_t\,dt\\
    &=\int{\Delta(t)}^{-1}u_t\pi({y(t^{-1})}^{*})\,dt\\
    &=\int u_t^*\pi({y(t)}^{*})\,dt\\
    &={((\pi\times u)(y))}^*,
  \end{align*}
  \[
    \|(\pi\times u)(y)\|\leq\int\|\pi(y(t))u_t\|\,dt=\|y\|_1,
  \]
  \begin{align*} 
    (\pi\times u)(y\times z)&=\int\pi\int y(s)\alpha_s(z(s^{-1}t))\,ds u_t\,dt\\
    &=\int\int\pi(y(s))u_s \pi(z(s^{-1}t))u_s^* u_t\,ds\,dt\\
    &=\int\int\pi(y(s))u_s\pi(z(t))u_t\,ds\,dt=(\pi\times u)(y)(\pi\times u)(z),
  \end{align*}
  which show that \((\pi\times u,\mathcal{H})\) extends to a representation of \(L^1(G,A)\).\par
  Conversely, if \((\rho,\mathcal{H})\) is a nondegenerate representation of \(L^1(G,A)\), then take
  an approximate unit \(\{y_\lambda\}\) for \(L^1(G,A)\) contained in \(C_C(G,A)\).
  One such may be obtained putting \(y_\lambda=x_\lambda\otimes f_\lambda\), where \(\{x_\lambda\}\) and \(\{f_\lambda\}\)
  are approximate units for \(A\) and \(L^1(G)\) respectively. Since \(\rho,\mathcal{H}\) is nondegenerate,
  \(\{\rho(y_\lambda)\}\) converges strongly to \(1\). For all \(x\) in \(\tilde{A}\) and \(t\) in \(G\), we define
  \begin{align*}
    \pi(x)&=\lim\rho(L(x,\delta_e)y_\lambda)=\lim\rho(R(x,\delta_e)y_\lambda),\\
    u_t&=\lim\rho(L(1,\delta_t)y_\lambda)=\lim\rho(R(1,\delta_t)y_\lambda),
  \end{align*}
  where the limits are taken in the weak topology. The existence of these limits follows as in the proof of \ref{Petersen:3123}.
  Since \(L(xy,\delta_e)=L(x,\delta_e)L(y,\delta_e)\), we see that the map \(x\mapsto \pi(x)\) is a homomorphism of \(\tilde{A}\), which is *-preserving by
  the first statement in \ref{Petersen:763}. Thus \((\pi,\mathcal{H})\) is a representation of \(\tilde{A}\).
  Since \(L(1,\delta_{st})=L(1,\delta_s)L(1,\delta_t)\) and \(L(1,\delta_e)=1\), we see that the map \(t\mapsto u_t\)
  is a unitary representation of \(G\), the continuity following from the facts that the map \(t\mapsto L(1,\delta_t)y\) is continuous in the norm \(\|\cdot\|_1\)
  for every \(y\) in \(L^1(G,A)\) and that vectors of the form \(\rho(y)\xi\), \(\xi\in\mathcal{H}\), are dense in \(\mathcal{H}\).
  Finally we have
  \begin{align*}
    u_t\pi(x)u_t^*\rho(y)&=\rho(L(1,\delta_t)(L(x,\delta_e))(L(1,\delta_{t^{-1}})))\\
    &=\rho(L(\alpha_t(x),\delta_e)y)\\
    &=\pi(\alpha_t(x))\rho(y)
  \end{align*}
  for every \(y\) in \(C_C(G,A)\). Taking \(y=y_\lambda\) and passing to the limit, we see that \((\pi,u,\mathcal{H})\) is a covariant representation of \((A,G,\alpha)\).\par
  To show that the correspondence is bijective, take the covariant representation \((\pi,u,\mathcal{H})\) obtained from a nondegenerate
  representation \((\rho,\mathcal{H})\) of \(L^1(G,A)\), and for \(y,z\) in \(C_C(G,A)\), note that
  \[
    (L(y(t),\delta_e)(L(1,\delta_t)z))(s)=y(t)\alpha_t(z(t^{-1}s))
  \]
  for each \(t\) in \(G\). Therefore
  \begin{align*}
    \int\pi(y(t))u_t\,dt\rho(z)&=\int\rho(L(y(t),\delta_e)(L(1,\delta_t)z))\,dt\\
    &=\int\rho(y(t)\alpha_t(z(t^{-1}s)))\,dt\\
    &=\rho(y\times z)\\
    &=\rho(y)\rho(z).
  \end{align*}
  Taking \(z=y_\lambda\) and passing to the limit, we obtain 
  \[
    \int\pi(y(t))u_t\,dt=\rho(y),
  \]
  whence \(\rho=\pi\times u\). Conversely, if \(\rho=\pi\times u\), then take \(y_\lambda=x_\lambda\otimes f_\lambda\) as before and note that, for all \(x\) in \(A\)
  and \(t\) in \(G\), we have
  \[
    (L(x,\delta_t)y_\lambda)(s)=x\alpha_t(x_\lambda)f_\lambda(t^{-1}s).
  \]
  Consequently,
  \begin{align*}
    &\lim\rho(L(x,\delta_e)y_\lambda)=\lim\int\pi(x x_\lambda)u_s f_\lambda(s)\,ds=\pi(x),\\
    &\lim\rho(L(1,\delta_t)y_\lambda)=\lim\int\pi(\alpha_t(x_\lambda))u_s f_\lambda(t^{-1}s)\,ds=u_t.
  \end{align*}
  This completes the proof.
\end{proof}
\begin{definition}
  The \textnormal{universal representation} \((\pi_u,\mathcal{H}_u)\) of \(L^1(G,A)\) is the direct sum
  of all nondegenerate representations of \(L^1(G,A)\), and the \textnormal{crossed product} (or the \textnormal{covariance algebra})
  of the \(C^*\)-dynamical system \((A,G,\alpha)\) is the norm closure of \(\pi_u(L^1(G,A))\) in \(\mathcal{B}(\mathcal{H}_u)\),
  denoted by \(A\rtimes_\alpha G\).
  Thus by definition the identity map of \(A\rtimes_\alpha G\) into \(\mathcal{B}(\mathcal{H}_u)\) is the universal representation of \(A\rtimes_\alpha G\).
\end{definition}
\begin{theorem}\label{Petersen:766}
  For each \(C^*\)-dynamical system \((A,G,\alpha)\), there is a covariant representation \((\pi,u,\mathcal{H})\) such that 
  \[
    A\rtimes_\alpha G\subset C^*(\pi(A)\cup u_G)\subset M(A\rtimes_\alpha G),\footnote{We denote \(u_G=\{u_t| t\in G\}\)}
  \]
  and for any other nondegenerate covariant representation \((\pi',u',\mathcal{H}')\), there is a unique nondegenerate representation
  \(\rho:A\rtimes_\alpha G\rightarrow\mathcal{B}(\mathcal{H}')\) such that
  \[
    \pi'=\widetilde{\rho}\circ\pi,
    \qquad
    u'_t=\widetilde{\rho}(u_t)\quad(t\in G),
  \]
  where \(\widetilde{\rho}:M(A\rtimes_\alpha G)\rightarrow\mathcal{B}(\mathcal{H}')\) is the unique extension of \(\rho\) to the multiplier algebra;
  compare \cite[Propositions 2.34 and 2.40]{Williams2007}.
\end{theorem}
\begin{proof}
  The existence and universal properties of \((\pi,u,\mathcal{H})\) follow from the 
  definition of \(A\rtimes_\alpha G\) together with \ref{Petersen:764}. The inclusion \(\pi(A)\cup u_G \subset M(A\rtimes_\alpha G)\)
  follows from the construction of the correspondence between representations of
  \(A\rtimes_\alpha G\) and covariant representations of \((A,G,\alpha)\) in \ref{Petersen:764}. To show that \(A\rtimes_\alpha G\)
  is contained in the \cstar\ generated by \(\pi(A)\cup u_G\), take \(y\) in \(C_C(G,A)\). For
  each \(\epsilon>0\), there are a compact set \(C\) in \(G\) containing the support of \(y\), a finite
  set \(\{x_n\}\) in \(A\), and functions \(\{f_n\}\) in \(C_C(G)\) with supports in \(C\) such that 
  \[
    \left\| y(t)-\sum f_n(t)x_n\right\| <\epsilon {\mu_G(C)}^{-1}
  \]
  for all \(t\) in \(G\). Put
  \[
    h(t)=\sum_n f_n(t)x_n.
  \]
  The elements \(\int u_t f_n(t)\,dt\) belong to \(u(L^1(G))\), whence
  \[
    z=(\pi\times u)(h)
      =\sum_n\int\pi(x_n)u_t f_n(t)\,dt
      \in C^*(\pi(A)\cup u_G).
  \]
  Since \(\|y-h\|_1<\epsilon\) and the integrated representation is contractive with respect to the \(L^1\)-norm,
  \[
    \|(\pi\times u)(y)-z\|
    =\|(\pi\times u)(y-h)\|
    \leq\|y-h\|_1
    <\epsilon.
  \]
  Consequently, \((\pi\times u)(y)\in C^*(\pi(A)\cup u_G)\).
\end{proof}

\begin{proposition}[Exactness of full crossed products]
  \label{prop:exactness-full-crossed-products}
  Let \((A,G,\alpha)\) be a \(C^*\)-dynamical system, and let
  \(I\) be a closed two-sided \(\alpha\)-invariant ideal of \(A\).
  Denote by
  \[
    \alpha|_I:G\longrightarrow\operatorname{Aut}(I)
  \]
  the restricted action, and by
  \[
    \overline{\alpha}_s(a+I)
      =\alpha_s(a)+I,
    \qquad a\in A,\ s\in G,
  \]
  the induced action on \(A/I\). If
  \(\iota:I\to A\) is the inclusion and
  \(q:A\to A/I\) is the quotient map, then the induced sequence
  \[
    \begin{tikzcd}
      0\arrow[r]
      &I\rtimes_{\alpha|_I}G
        \arrow[r,"\iota\rtimes G"]
      &A\rtimes_\alpha G
        \arrow[r,"q\rtimes G"]
      &(A/I)\rtimes_{\overline{\alpha}}G
        \arrow[r]
      &0
    \end{tikzcd}
  \]
  is exact.
\end{proposition}

\begin{proof}
  Since \(\iota\) and \(q\) are equivariant, their integrated forms
  induce the indicated homomorphisms between the full crossed
  products.

  We first show that \(\iota\rtimes G\) is injective. Every
  nondegenerate representation \(\pi\) of \(I\) extends uniquely to a
  representation \(\widetilde{\pi}\) of \(A\), characterized by
  \[
    \widetilde{\pi}(a)\pi(x)\xi
      =\pi(ax)\xi,
    \qquad
    a\in A,\ x\in I.
  \]
  If \((\pi,u)\) is a covariant representation of
  \((I,G,\alpha|_I)\), then \((\widetilde{\pi},u)\) is a covariant
  representation of \((A,G,\alpha)\). Consequently, the universal
  norm of an element of \(C_C(G,I)\) is the same whether it is
  computed in \(I\rtimes_{\alpha|_I}G\) or in
  \(A\rtimes_\alpha G\). Thus \(\iota\rtimes G\) is injective.

  Its image is the closed ideal of \(A\rtimes_\alpha G\) generated
  by \(C_C(G,I)\). Indeed, convolution on either side by an element
  of \(C_C(G,A)\) maps \(C_C(G,I)\) into \(C_C(G,I)\).

  Clearly,
  \[
    (q\rtimes G)\circ(\iota\rtimes G)=0.
  \]
  Conversely, a nondegenerate representation of
  \(A\rtimes_\alpha G\) vanishes on the image of
  \(I\rtimes_{\alpha|_I}G\) precisely when its associated covariant
  representation \((\pi,u)\) satisfies \(\pi(I)=0\). In that case
  \(\pi\) factors uniquely through a representation
  \(\overline{\pi}\) of \(A/I\), and
  \((\overline{\pi},u)\) is a covariant representation of
  \((A/I,G,\overline{\alpha})\).

  The universal property of the full crossed product therefore gives
  \[
    \frac{A\rtimes_\alpha G}
         {(\iota\rtimes G)
          (I\rtimes_{\alpha|_I}G)}
    \cong
    (A/I)\rtimes_{\overline{\alpha}}G.
  \]
  Hence the sequence is exact; compare
  \cite[Proposition~3.19]{Williams2007}.
\end{proof}

We identify \(A\) with its image in \(M(A\rtimes_\alpha G)\) via the map \(x\mapsto L(x,\delta_e)\)
and let \(t\mapsto u_t\) be the unitary representation of \(G\) into \(M(A\rtimes_\alpha G)\) described in \ref{Petersen:766}.
For each functional \(\phi\) on \(A\rtimes_\alpha G\), define the function \(\Phi:G\rightarrow A^*\) by
\begin{equation}\label{Petersen:767}
  \Phi(t)(x)=\phi(x u_t),\qquad t\in G, x\in A.
\end{equation}
The set of such functions is denoted by \(B(A\rtimes_\alpha G)\), and the elements in the generating cone 
\(B_+(A\rtimes_\alpha G)\) arising from positive functionals are said to be \textit{positive definite} (with respect to \(\alpha\)).\par
From Cauchy-Schwarz inequality it follows that each \(\Phi\) in \(B_+(A\rtimes_\alpha G)\) is continuous in the norm topology
on \(A^*\) and bounded by \(\left\|\Phi(e)\right\|\). From (\ref{eq:Pet764}) we see that
\[
  \phi(y)=\int\Phi(t)(y(t))\,dt
\]
for each \(y\) in \(C_C(G,A)\), so that the correspondence \(\phi\mapsto \Phi\) is a bijection from
\(
  {(A\rtimes_\alpha G)}^*
\)
onto \(B(A\rtimes_\alpha G)\).
\begin{proposition}\label{Petersen:768}
  The following conditions on a bounded norm-continuous function \(\Phi:G\rightarrow A^*\) are equivalent:
  \begin{enumerate}
    \item \(\Phi\) is positive definite;
    \item \(\sum_{ij}\Phi(s_i^{-1}s_j)(\alpha_{s_i^{-1}}(x_i^* x_j))\geq 0\) for all finite sets \(\{s_i\}\) in \(G\) and \(\{x_i\}\) in \(A\);
    \item \(\sum_{nm}\int\Phi(t)(y_n^*\alpha_t(y_m))(f^*_n\times f_m)(t)\,dt\geq 0\) for all finite sets \(\{f_n\}\) in \(C_C(G)\) and \(\{y_n\}\) in \(A\);
    \item \(\int\Phi(t)(y^*\times y)(t)\,dt\geq 0\) for each \(y\) in \(C_C(G,A)\).
  \end{enumerate}
\end{proposition}
\begin{proof}
  \begin{itemize}
    \item (1)\(\Rightarrow\)(2). Given \(\{s_i\}\) and \(\{x_i\}\), define \(x=\sum_i x_i u_{s_i}\), in \(M(A\rtimes_\alpha G)\).
      If \(\phi\) is the functional on \(A\rtimes_\alpha G\) corresponding to \(\Phi\).
      Let \(\widetilde\phi\) denote the canonical positive extension of
      \(\phi\) to \(M(A\rtimes_\alpha G)\). Then
      \(0\leq\widetilde\phi(x^*x)=\cdots\).
      We have
      \begin{align*}
        0\leq\phi(x^* x)&=\phi\left(\sum_{ij}\alpha_{s_i^{-1}}(x_i^* x_j)u_{s_i^{-1}s_j}\right)\\
        &=\sum_{ij}\Phi(s_i^{-1}s_j)(\alpha_{s_i^{-1}}(x_i^* x_j)).
      \end{align*}
    \item (2)\(\Rightarrow\)(3). Given \(\{f_n\}\) in \(C_C(G)\) and \(\{y_n\}\) in \(A\), take a finite set \(\{s_i\}\) in \(G\)
      and apply condition (2) to the sets \(\{s_{in}\}\) and \(\{x_{in}\}\), where \(s_{in}=s_i\) and \(x_{in}=\alpha_{s_i}(y_n)f_n(s_i)\). Thus
      \[
        0\leq\sum_{ijnm}\Phi(s_i^{-1}s_j)(y^*_n\alpha_{s_i^{-1}s_j}(y_m))\overline{f_n(s_i)}f_m(s_j).
      \]
      It follows that 
      \begin{align*}
        0&\leq\sum_{nm}\int\int\Phi(s^{-1}t)(y^*_n\alpha_{s^{-1}t}(y_m))\overline{f}_n(s)f_m(t)\,ds\,dt\\
        &=\sum_{nm}\int\int\Phi(t)(y^*_n\alpha_t(y_m))f_n^*(s^{-1}){\Delta(s)}^{-1}f_m(st)\,ds\,dt\\
        &=\sum_{nm}\int\Phi(t)(y^*_n\alpha_t(y_m))(f_n^*\times f_m)(t)\,dt.
      \end{align*}
    \item (3)\(\Rightarrow\)(4). Each \(y\) in \(C_C(G,A)\) can be approximated (in \(L^1(G,A)\)) by elements of 
      the form \(y(t)=\sum f_n(t)\alpha_t(y_n)\). However, for such an element, \(y^*(t)=\sum f_n^*(t)y_n^*\), and
      \[
        y^*\times y(t)=\sum_{nm} y^*_n\alpha_t(y_m)f_n^*\times f_m(t).
      \]
      Condition (3) tells us that \(\int\Phi(t)((y^*\times y)(t))\,dt\geq 0\), and the result follows by continuity.
    \item (4)\(\Rightarrow\)(1). For each \(y\) in \(C_C(G,A)\), define
      \[  
        \phi(y)=\int\Phi(t)(y(t))\,dt.
      \] 
      From condition (4) we see that \(\phi\) is a positive functional on \(C_C(G,A)\), and as in the proof of \ref{Petersen:719}, this implies
      that \(\phi\) has a unique extension to a positive functional on \(A\rtimes_\alpha G\).
  \end{itemize}
\end{proof}
\begin{corollary}\label{Petersen:769}
  If \(\Phi\in B_+(A\rtimes_\alpha G)\) and \(\Psi\in B_+(G)\), then 
  \[
    \Psi\cdot\Phi\in B_+(A\rtimes_\alpha G).
  \]
\end{corollary}
\begin{proof}
  Condition (2) in \ref{Petersen:768} is equivalent to the statement that each matrix \((a_{ij})\) with
  \[
    a_{ij}=\Phi(s_i^{-1}s_j)(\alpha_{{s_i}^{-1}}(x_i^* x_j))
  \]
  is positive. Since the pointwise product of positive matrices is again positive (see the proof of \ref{Petersen:7110}), we conclude that
  \(\Psi\cdot\Phi\) satisfies \ref{Petersen:768} (2).
\end{proof}

\begin{proposition}\label{Petersen:7610}
  Let \(S(A\rtimes_\alpha G)\) be the state space of \(A\rtimes_\alpha G\), and let \(B_+^1(A\rtimes_\alpha G)\)
  be the set of positive definite functions \(\Phi:G\rightarrow A^*\) with \(\|\Phi(e)\|=1\).
  There is an affine homeomorphism of \(S(A\rtimes_\alpha G)\) in the weak-* topology onto
  \(B_+^1(A\rtimes_\alpha G)\) equipped with the topology of weak-* convergence uniformly on compact subsets of \(G\).
\end{proposition}
\begin{proof}
  It is clear that the correspondence \(\phi\rightarrow\Phi\) defined in \ref{Petersen:767} is an affine isomorphism 
  of \(S(A\rtimes_\alpha G)\) onto \(B_+^1(A\rtimes_\alpha G)\).\par
  Let \(\{\phi_\lambda\}\) be a net in \(S(A\rtimes_\alpha G)\) that is weak-* convergent to \(\phi\) in \(S(A\rtimes_\alpha G)\),
  and let \(\{\Phi_\lambda\}\) and \(\Phi\) be the corresponding elements in \(B_+^1(A\rtimes_\alpha G)\). Combining \ref{Petersen:768} (3) and \ref{Petersen:719} (4),
  we see that, for each \(y\) in \(\tilde{A}\), the functions
  \[
    \Psi_\lambda:t\mapsto\Phi_\lambda(t)(y^*\alpha_t(y)),\qquad\Psi:t\mapsto\Phi(t)(y^* \alpha_t(y))
  \]
  are positive definite (in the ordinary sense) and that the corresponding net of positive functionals
  \(\{\psi_\lambda\}\) is weak-* convergent to \(\psi\) in \(C^*(G)^*\).
  By \ref{Petersen:7111} this implies that \(\{\Psi_\lambda\}\) converges to \(\Psi\) uniformly on compact sets.
  Using the polarization identity
  \[
    4x^*\alpha_t(y)=\sum_{k=0}^3 i^k(y+i^k x)^*\alpha_t(y+i^k x),
  \]  
  we see that \(\{\Phi_\lambda(t)(x^*\alpha_t(y))\}\) converges to \(\Phi(t)(x^*\alpha_t(y))\) uniformly on compact sets for all
  \(x,y\) in \(\tilde{A}\). With \(y=1\), we obtain the desired result.\par
  Conversely, if \(\Phi_\lambda(t)\rightarrow\Phi(t)\) weak-*, uniformly on compact subsets of \(G\),
  then for each element \(y\) in \(A\rtimes_\alpha G\) of the form \(y=f\otimes x\), \(f\in C_C(G)\), \(x\in A\), we have
  \[
    \phi_\lambda(y)=\int\Phi_\lambda(t)(x)f(t)\,dt\rightarrow\int\Phi(t)(x)f(t)\,dt=\phi(y).
  \]
  Since finite sums of such elements form a dense set in \(A\rtimes_\alpha G\), we conclude that \(\phi_\lambda\rightarrow \phi\) weak-*.
\end{proof}

\section{Regular representations of crossed products}
\begin{definition}\label{Petersen:771}
  Given a \(C^*\)-dynamical system \((A,G,\alpha)\) and a representation
  \((\pi,\mathcal{H})\) of \(A\), define the covariant representation \((\tilde{\pi},\lambda,L^2(G,\mathcal{H}))\) by
  \[
    (\tilde{\pi}(x)\xi)(t)=\pi(\alpha_{t^{-1}}(x))\xi(t)\qquad (\lambda_s \xi)(t)=\xi(s^{-1}t)
  \]
  for all \(x\) in \(A\), \(s\) in \(G\), and \(\xi\) in \(L^2(G,\mathcal{H})\).
  The \textnormal{regular representation} of \(A\rtimes_\alpha G\) induced by \((\pi,\mathcal{H})\) is the representation 
  \((\tilde{\pi}\times\lambda,L^2(G,\mathcal{H}))\) (cf. \ref{Petersen:766}).
  From (\ref{eq:Pet764}) we see that, for all \(y\) in \(C_C(G,A)\) and \(\xi\) in \(L^2(G,\mathcal{H})\). We have
  \begin{align*}
    (((\tilde{\pi}\times\lambda)y)\xi)(t)&=\int(\tilde{\pi}(y(s))\lambda_s\xi)(t)\,ds\\
    &=\int\pi(\alpha_{t^{-1}}(y(s)))\xi(s^{-1}t)\,ds.
  \end{align*}
\end{definition}
\begin{lemma}\label{Petersen:772}
  If \(\{\xi_i\}\subset\mathcal{H}\) and \(\{f_j\}\subset L^2(G)\) are such that \(\{\xi_i\}\)
  is cyclic for \((\pi,\mathcal{H})\) (i.e., \(\sum\pi(A)\xi_i\) is dense in \(\mathcal{H}\)) and \(\{f_j\}\) is cyclic for \((\lambda,L^2(G))\),
  then the set \(\{f_j\otimes\xi_i\}\) is cyclic for \(\tilde{\pi}\times\lambda\) in \(L^2(G,\mathcal{H})\).
\end{lemma}
\begin{proof}
  Suppose that \(\xi\in L^2(G,\mathcal{H})\) such that
  \[
    ((\tilde{\pi}\times\lambda)(A\rtimes_\alpha G)f_j\otimes\xi_i|\xi)=0
  \]
  for all \(i\) and \(j\). This means that, for all \(x\) in \(A\) and \(f\) in \(C_C(G)\), we have
  \begin{align*}
    0&=((\tilde{\pi}\times\lambda)(x\otimes f)f_j\otimes\xi_i|\xi)\\
    &=\int\int(\pi(\alpha_{t^{-1}}(x))f(s)f_j(s^{-1}t)\xi_i|\xi(t))\,ds\,dt\\
    &=\int(\pi(\alpha_{t^{-1}}(x))\xi_i|\xi(t))(f\times f_j)(t)\,dt.
  \end{align*}
  Since \(C_C(G,\mathcal{H})\) is dense in \(L^2(G,\mathcal{H})\), we can extend Lusin's theorem to the vector-valued case.
  Thus, for each compact set \(E\) in \(G\) and \(\epsilon>0\), there is a 
  compact subset \(F\) of \(E\) such that \(\mu_G(E\setminus F)<\epsilon\) and 
  \(\xi|_F\in C(F,\mathcal{H})\). From the density of the subspace spanned by 
  \(\{f\times f_j\}\), \(f\in C_C(G)\), we see that the continuous function
  \(t\mapsto(\pi(\alpha_{t^{-1}}(x))\xi_i|\xi(t))\), \(t\in F\), is identically zero for all \(x\) in \(A\) and all \(\xi_i\).
  From the cyclicity of \(\{\xi_i\}\) it follows that \(\xi(t)=0\) for all \(t\) in \(F\).
  Since \(E\) and \(\epsilon\) were arbitrary, we conclude that \(\xi=0\), whence \(\{f_j\otimes \xi_i\}\)
  is cyclic for \(\tilde{\pi}\times\lambda\).
\end{proof}
\begin{proposition}\label{Petersen:773}
  Let \((A,G,\alpha)\) be a \(C^*\)-dynamical system. If \(\phi\) is a state of \(A\) with 
  cyclic representation \((\pi_\phi,\mathcal{H}_\phi,\xi_\phi)\), then for each \(z\) in \(C_C(G,A)\)
  such that \(\phi((z^*\times z)(e))=1\), there is a vector state \(\tilde{\phi}_z\) of \(A\rtimes_\alpha G\)
  in \((\tilde{\pi}\times\lambda, L^2(G,\mathcal{H}_\phi))\) such that
  \[
    \tilde{\phi}_z(y)=\phi((z^*\times y\times z)(e)),\qquad y\in C_C(G,A).
  \]
  Moreover, each vector state in \((\tilde{\pi}\times\lambda,L^2(G,\mathcal{H}_\phi))\)
  can be represented in norm by states of the form \(\tilde{\phi}_z\).
\end{proposition}
\begin{proof}
  If \(z\) in\ \(C_0(G,A)\) is such that the function \(t\mapsto\|z(t)\|\) belongs to \(L^2(G)\),
  then define \(\xi_z\) in \(L^2(G,\mathcal{H}_\phi)\) by \(\xi_z(t)=\pi_\phi(\alpha_{t^{-1}}(z(t)))\xi_\phi\).
  For \(y\) in \(C_C(G,A)\), we have
  \begin{align*}
    ((\tilde{\pi}_\phi\times\lambda)(y^*\times y)\xi_z|\xi_z)&=\left\|(\tilde{\pi}_\phi\times\lambda)(y)\xi_z\right\|^2\\
    &=\int\Vert\int\pi_\phi(\alpha_{t^{-1}}(y(s)))\xi_z(s^{-1}t)\,ds \Vert^2\,dt\\
    &=\int\Vert\int\pi_\phi(\alpha_{t^{-1}}(y(s))\alpha_{t^{-1}s}(z(s^{-1}t)))\xi_\phi\,ds \Vert^2\,dt\\
    &=\int\Vert\pi_\phi(\alpha_{t^{-1}}(y\times z)(t))\xi_\phi\Vert^2\,dt\\
    &=\int\phi(\alpha_{t^{-1}}((y\times z)(t)^*(y\times z)(t)))\,dt\\
    &=\int\phi((y\times z)^*(t^{-1})\Delta(t^{-1})\alpha_{t^{-1}}((y\times z)(t)))\,dt\\
    &=\int\phi((y\times z)^*(t)\alpha_t((y\times z)(t^{-1})))\,dt\\
    &=\phi((y\times z)^*\times (y\times z)(e))=\tilde{\phi}_z(y^*\times y).
  \end{align*}
  It follows that \(\tilde{\phi}_z\) is a vector state in \((\tilde{\pi}_\phi\times\lambda,L^2(G,\mathcal{H}_\phi))\).\par
  To show that every vector state in \((\tilde{\pi}_\phi\times\lambda,L^2(G,\mathcal{H}_\phi))\) can be
  approximated by states of the form \(\tilde{\phi}_z\), it suffices to prove that the subspace \(\{\xi_z| z\in C_C(G,A)\}\) is dense in 
  \(L^2(G,\mathcal{H}_\phi)\), a short computation shows that, for each \(f\) in \(C_C(G)\), we have 
  \[
    (\tilde{\pi}_\phi\times\lambda)(z)(f\otimes\xi_\phi)=\xi_{z\times f},
  \]
  where \(z\times f(t)=\int z(s)f(s^{-1}t)\,ds\). Since the subset \(\{f\otimes\xi_\phi| f\in C_C(G)\}\)
  is cyclic for \((\tilde{\pi}_\phi\times\lambda,L^2(G,\mathcal{H}_\phi))\) by \ref{Petersen:772}, the proof is complete.
\end{proof}
\begin{definition}\label{Petersen:774}
  Let \((A,G,\alpha)\) be a \(C^*\)-dynamical system, and let \((\pi_u,\mathcal{H}_u)\) denote the
  universal representation of \(A\). The \textnormal{reduced crossed product} of \(A\) by \(G\) is the \cstar\
  \((\tilde{\pi}_u\times\lambda)(A\rtimes_\alpha G)\) denoted by \(A\rtimes_{\alpha,r}G\).\par
  The set of functionals \(\psi\) on \(A\rtimes_\alpha G\) of the form
  \[
    \psi(y)=\sum((\tilde{\pi}_u\times\lambda)y\xi_n|\eta_n),
  \]
  where \(\{\xi_n\}\subset L^2(G,\mathcal{H}_u)\), \(\{\eta_n\}\subset L^2(G,\mathcal{H}_u)\),
  \(\sum\left\|\xi_n\right\|^2<\infty\), and \(\sum\left\|\eta_n\right\|^2<\infty\),
  is a norm-closed subspace of \((A\rtimes_\alpha G)^*\) and can be identified with
  \(((A\rtimes_{\alpha,r} G)'')_*\) (\(''\) denoting double commutant in \(\mathcal{B}(L^2(G,\mathcal{H}_u))\)).
  We denote by \(A(A\rtimes_\alpha G)\) the subspace of \(B(A\rtimes_\alpha G)\) 
  consisting of \(A^*\)-valued functions whose associated functionals belong to this set (cf.\ref{Petersen:772}).\par
\end{definition}
Since \((\pi_u,\mathcal{H}_u)\) is the direct sum of cyclic representations of \(A\),
we see from \ref{Petersen:773} that \(A_+(A\rtimes_\alpha G)\) is the 
closure (in the norm inherited from \((A\rtimes_\alpha G)^*\)) of 
finite sums of functions \(\Phi=\Phi(\phi,z)\) such that, for all \(x\) in \(A\) and \(t\) in \(G\),
\begin{align*}
  \Phi(t)(x)&=\tilde{\phi}_z(x u_t)=(\tilde{\pi}_\phi(x)\lambda_t\xi_z|\xi_z)\\
  &=\int(\pi_\phi(\alpha_{s^{-1}}(x)\alpha_{s^{-1}t}(z(t^{-1}s)))\xi_\phi|\pi_\phi(\alpha_{t^{-1}}(z(s)))\xi_\phi)\,ds\\
  &=\int\phi(\alpha_{s^{-1}}(z(s)^*x\alpha_t z(t^{-1}s)))\,ds\\
  &=\int\phi(\alpha_s(z(s^{-1})^*\Delta(s)^{-1}x\alpha_t z(t^{-1}s^{-1})))\,ds\\
  &=\phi\left( \int z^*(s)\alpha_s(x)\alpha_{st}(z((st)^{-1}))\,ds\right),
\end{align*}
where \(\phi\) is a state of \(A\), and \(z\in C_C(G,A)\).

\begin{theorem}\label{Petersen:775}
  If \((A, G, \alpha)\) is a \(C^*\)-dynamical system and \((\pi, \mathcal{H})\) is a faithful representation of \(A\),
  then \((\tilde{\pi}\times\lambda, L^2(G, \mathcal{H}))\) is a faithful representation of \(A\rtimes_{\alpha,r}G\).  
\end{theorem}
\begin{proof}
  Let \(E\) denote the set of states of \(A\) contained in \((\pi, \mathcal{H})\), and let \(F\) denote
  the set of functions in \(A^1_+(A\rtimes_\alpha G)\) whose associated states are contained in \((\tilde{\pi}\times\lambda,L^2(G,\mathcal{H}))\).
  Since \((\pi, \mathcal{H})\) is faithful on \(A\), we know that \(E\) is weak* dense in \(S(A)\).
  To show that \((\tilde{\pi}\times\lambda, L^2(G,\mathcal{H}))\) is faithful on \(A\rtimes_{\alpha,r}G\), we must
  show that the set of states contained in \((\tilde{\pi}\times\lambda,L^2(G,\mathcal{H}))\) is weak* dense in the set of states contained in \((\tilde{\pi}_u\times\lambda,L^2(G,\mathcal{H}_u))\).
  By \ref{Petersen:7610} this is equivalent to showing that \(F\) is dense in \(A^1_+(A\rtimes_\alpha G)\) in the topology of weak* convergence uniformly on compact subsets of \(G\).\par
  By \ref{Petersen:774} it suffices to approximate functions of the form
  \[
    \Phi(t)(x)=\phi\left(\int z^*(s)\alpha_s(x)\alpha_{st}(z(t^{-1}s^{-1}))\,ds\right),
  \]
  where \(\phi\in S(A)\) and \(z\in C_C(G,A)\) with \(\phi((z^*\times z)(e))=1\).
  Take a net \(\{\phi_i\}\) in \(E\) that is weak* convergent to \(\phi\) and let \(\{\Phi_i\}\) denote the corresponding net 
  of elements in \(A_+(A\rtimes_\alpha G)\) (\(z\) remains fixed). Then \(\Phi_i(t)\rightarrow \Phi(t)\) weak*, uniformly on compact subsets
  of \(G\). Moreover, each \(\Phi_i\) belongs to \(F\) by \ref{Petersen:773}, because the representation \((\pi_{\phi_i},\mathcal{H}_{\phi_i})\) is equivalent to a subrepresentation of \((\pi,\mathcal{H})\).
\end{proof}

\begin{lemma}\label{Petersen:776}
  Each element \(\Phi\) in \(B_+^1(A\rtimes_\alpha G)\) with compact support belongs to \(A_+^1(A\rtimes_\alpha G)\).
\end{lemma}
\begin{proof}
  Let \(\phi\) denote the state of \(A\rtimes_\alpha G\) corresponding to \(\Phi\) and put
  \(\psi=\Phi(e)\). From the definition of \(\Phi\) (see \ref{Petersen:767}) and the Cauchy-Schwarz
  inequality we have, for each \(x\) in \(A\),
  \begin{align*}
    |\Phi(t)(\alpha_t(x))|^2&=|\phi(\alpha_t(x)u_t)|^2=|\phi(u_t x)|^2\leq\phi(x^* x)\\
    &=\Phi(e)(x^* x)=\psi(x^* x).
  \end{align*} 
  Passing to the representation \((\pi_\psi,\mathcal{H}_\psi,\xi_\psi)\) of \(A\), it follows that, 
  for each \(t\), there is a unique vector \(\xi_t\) in \(\mathcal{H}_\psi\) such that
  \[
    \Phi(t)(\alpha_t(x))=(\pi_\psi(x)\xi_\psi|\xi_t).
  \]
  The function \(\xi:t\mapsto\xi_t\) has compact support and is weakly continuous. It therefore 
  belongs to \(L^2(G,\mathcal{H}_\psi)\).\par
  Take \(f_j\) in \(C_C(G)\) and, for each \(y\) in \(C_C(G,A)\), compute
  \begin{align*}
    ((\tilde{\pi}_\psi\times\lambda)(y)f_j\otimes\xi_\psi|\xi)&=\int\left(\int\pi_\psi(\alpha_{t^{-1}}(y(s)))f_j(s^{-1}t)\xi_\psi\,ds|\xi_t\right)\,dt\\
    &=\int(\pi_\psi(\alpha_{t^{-1}}((y\times f_j)(t)))\xi_\psi|\xi_t)\,dt\\
    &=\int\Phi(t)((y\times f_j)(t))\,dt=\phi(y\times f_j).
  \end{align*}
  Thus the functional \(\phi(\cdot f_j)\) is contained in \((\tilde{\pi}_u\times\lambda,L^2(G,\mathcal{H}_u))\).
  However, taking \(\{f_j\}\) as an approximate unit for \(L^1(G)\), we have \(y\times f_j\rightarrow y\) in the norm of \(L^1(G,A)\),
  whence \(\phi(\cdot f_j)\rightarrow \phi\) in norm. It follows that \(\Phi\in A(A\rtimes_\alpha G)\).
\end{proof}

\begin{theorem}\label{Petersen:777}
  If \((A,G,\alpha)\) is a \(C^*\)-dynamical system with \(G\) amenable, then \(A\rtimes_\alpha G\)
  is equal to \(A\rtimes_{\alpha,r}G\).
\end{theorem}
\begin{proof}
  We must show that \(A_+(A\rtimes_\alpha G)\) is dense in \(B_+(A\rtimes_\alpha G)\).
  Take \(\Phi\) in \(B_+(A\rtimes_\alpha G)\). Since \(G\) is amenable, there is a net \(\{\Psi_i\}\)
  in \(B_+(G)\cap C_C(G)\) such that \(\Psi_i(t)\rightarrow 1\) uniformly on compact subsets of \(G\) (see \ref{Petersen:738}).
  From \ref{Petersen:769} we have \(\Phi\cdot\Psi_i\in B_+(A\rtimes_\alpha G)\), and from \ref{Petersen:776} we have
  \(\Phi\cdot\Psi_i\in A_+(A\rtimes_\alpha G)\).
  Since \(\Phi\cdot\Psi_i\rightarrow \Phi\) weak-* uniformly on compact subsets of \(G\), we are done.
\end{proof}
\begin{corollary}
  If \((A,G,\alpha)\) is a (separable) \(C^*\)-dynamical system with \(G\) amenable and if \((\pi,\mathcal{H})\)
  is a faithful (separable) representation of \(A\), then \((\tilde{\pi}\times\lambda,L^2(G,\mathcal{H}))\)
  is a faithful (separable) representation of \(A\rtimes_\alpha G\)
\end{corollary}
  
\begin{proposition}\label{Petersen:7711}
  Let \((A, G, \alpha)\) be a \(C^*\)-dynamical system, and let \(\tilde{\lambda}\) denote the right regular representation of \(G\) on \(L^2(G)\),
  i.e., \((\tilde{\lambda}_t\xi)(s)=\xi(st){\Delta(t)}^{1/2}\), \(\xi\in L^2(G)\).
  Then there is a \(C^*\)-dynamical system \((A\otimes\mathcal{K}(L^2(G)), G, \alpha\otimes\operatorname{Ad}\tilde{\lambda})\) such that
  \[
    {(\alpha\otimes \textnormal{Ad}\tilde{\lambda})}_t(x\otimes y)=\alpha_t(x)\otimes\tilde{\lambda}_t y \tilde{\lambda}_t^*
  \]
  for all \(x\) in \(A\), \(y\) in \(\mathcal{K}(L^2(G))\) and \(t\) in \(G\).
\end{proposition}
\begin{proof}
  Choose a faithful covariant representation \((\pi, u, \mathcal{H})\) of \((A, G, \alpha)\).
  Then define a unitary representation \(w\) of \(G\) on \(\mathcal{H}\otimes L^2(G)\) by
  \[
    w_t(\xi\otimes\eta)=u_t\xi\otimes\tilde{\lambda}_t\eta,\qquad \xi\in \mathcal{H}, \eta\in L^2(G).
  \]
  Identifying \(A\) and \(\pi(A)\), we may consider \(A\otimes\mathcal{K}(L^2(G))\) as a subalgebra of
  \(\mathcal{B}(\mathcal{H}\otimes L^2(G))\). If \(x\in A\) and \(y\in\mathcal{K}(L^2(G))\), then
  \[
    w_t(x\otimes y)w_t^*=u_t x u_t^*\otimes\tilde{\lambda}_t y\tilde{\lambda}_t^*=\alpha_t(x)\otimes \tilde{\lambda}_t y\tilde{\lambda}_t^*.
  \]
  It follows that
  \[
    w_t(A\otimes\mathcal{K}(L^2(G)))w_t^*\subset A\otimes\mathcal{K}(L^2(G)),
  \]
  and furthermore we see that if \(y\) has finite rank, then the function \(t\mapsto\alpha_t(x)\otimes\tilde{\lambda}_t y\tilde{\lambda}_t^*\)
  is norm continuous since \(t\mapsto\tilde{\lambda}_t\) is strongly continuous.
  The elements \(z\) in \(A\otimes\mathcal{K}(L^2(G))\) for which the function \(t\mapsto w_t z w_t^*\)
  is norm continuous form a \cstar\ \footnote{
    \begin{lemma} 
      Let \(B\) be a \(C^*\)-algebra, let \(G\) be a topological group, and let
      \(\beta:G\to\operatorname{Aut}(B)\) be an action of \(G\) on \(B\).

      Then the set
      \[
      B^c=
      \{\,x\in B : t\mapsto \beta_t(x)
      \text{ is norm continuous}\,\}
      \]
      is a \(C^*\)-subalgebra of \(B\).
    \end{lemma}

    \begin{proof}
      If \(x,y\in B^c\) and \(\lambda\in\mathbb C\), then
      \[
      \beta_t(x+\lambda y)
      =
      \beta_t(x)+\lambda\beta_t(y),
      \]
      so \(t\mapsto\beta_t(x+\lambda y)\) is norm continuous.

      If \(x,y\in B^c\), then
      \[
      \beta_t(xy)
      =
      \beta_t(x)\beta_t(y),
      \]
      and since multiplication is norm continuous in \(B\), the map
      \(t\mapsto\beta_t(xy)\) is norm continuous.

      If \(x\in B^c\), then
      \[
      \beta_t(x^*)
      =
      \beta_t(x)^*,
      \]
      and since the involution is norm continuous, the map
      \(t\mapsto\beta_t(x^*)\) is norm continuous.

      Finally, let \((x_n)\) be a sequence in \(B^c\) converging in norm to
      \(x\in B\). For every \(t,s\in G\),
      \[
      \|\beta_t(x)-\beta_s(x)\|
      \le
      2\|x-x_n\|
      +
      \|\beta_t(x_n)-\beta_s(x_n)\|.
      \]
      Since \(x_n\in B^c\), the last term tends to \(0\) as \(t\to s\). It follows
      that \(t\mapsto\beta_t(x)\) is norm continuous. Hence \(B^c\) is norm closed.

      Therefore \(B^c\) is a \(C^*\)-subalgebra of \(B\).
    \end{proof}
  }. Since this algebra contains all elements \(x\otimes y\),
  where \(x\in A\) and \(y\) has finite rank, it equals \(A\otimes\mathcal{K}(L^2(G))\).
  It follows that if we define
  \[
    {(\alpha\otimes\textnormal{Ad}\tilde{\lambda})}_t(z)=w_t z w_t^*,\qquad z\in A\otimes\mathcal{K}(L^2(G)),
  \]
  then \((A\otimes\mathcal{K}(L^2(G)), G, \alpha\otimes\textnormal{Ad}\tilde{\lambda})\) is a \(C^*\)-dynamical system, as desired.
\end{proof}

\begin{theorem}\label{Petersen:7712}
  Given a \(C^*\)-dynamical system \((A, G, \alpha)\), there exist two \(C^*\)-dynamical systems 
  \((C_0(G, A), G, \rho)\) and \((C_0(G, A), G, \gamma)\), where
  \[
    (\rho_s(y))(t)=y(ts),
    \qquad
    (\gamma_s(y))(t)=\alpha_s(y(s^{-1}t))
  \]
  for \(s,t\in G\) and \(y\in C_0(G,A)\). The two actions \(\rho\) and \(\gamma\) commute, and each \(\rho_s\) extends to an automorphism of \(C_0(G,A)\rtimes_{\gamma,r}G\),
  so that we obtain a \(C^*\)-dynamical system \((C_0(G,A)\rtimes_{\gamma,r}G, G, \rho)\).
  This system is covariantly isomorphic to the system
  \((A\otimes \mathcal{K}(L^2(G)), G, \alpha\otimes\operatorname{Ad}\tilde{\lambda})\)
  defined in~\ref{Petersen:7711}; compare \cite[Section 7.1]{Williams2007}.
\end{theorem}
\begin{proof}
  It is straightforward to verify that the definitions
  \[
    (\rho_s(y))(t)=y(ts),\qquad (\gamma_s(y))(t)=\alpha_s(y(s^{-1}t)),\qquad y\in C_0(G, A),
  \]
  give \(C^*\)-dynamical systems \((C_0(G, A), G, \rho)\) and \((C_0(G, A), G, \gamma)\). Evidently, the two actions \(
  \rho\) and \(\gamma\) commute.\par
  Take now a faithful representation of \(A\) on some Hilbert space \(\mathcal{H}\) and define the faithful representation \((\pi, L^2(G, \mathcal{H}))\) of \(C_0(G, A)\) by
  \[
    (\pi(y)\xi)(t)=y(t)\xi(t),\qquad y\in C_0(G,A),\qquad \xi\in L^2(G, \mathcal{H}).
  \]
  It follows from \ref{Petersen:775} that \((\tilde{\pi}\times\lambda, L^2(G\times G, \mathcal{H}))\) is a faithful representation of \(C_0(G,A)\rtimes_{\gamma,r}G\).
  Moreover, regarding each \(z\) in \(C_C(G\times G, A)\) as an element of \(C_C(G, C_0(G, A))\), we have, for every \(\xi\) in \(L^2(G\times G, \mathcal{H})\),
  \begin{align*}
    (((\tilde{\pi}\times\lambda)z)\xi)(s,t)&=\left(\int\pi(\gamma_{s^{-1}}(z(r,\cdot)))\xi(r^{-1}s,\cdot)\,dr\right)(t)\\
    &=\int\alpha_{s^{-1}}(z(r,st))\xi(r^{-1}s,t)\,dr.
  \end{align*}\par
  Define the unitary operator on \(L^2(G\times G, \mathcal{H})\) by
  \[
    (w\xi)(s,t)=\xi(st,t){\Delta(t)}^{1/2}, \qquad \xi\in L^2(G\times G, \mathcal{H}).
  \]
  Using the formula above, we get
  \begin{align*}
    (w^*((\tilde{\pi}\times\lambda)z)w\xi)(s,t)&=(((\tilde{\pi}\times\lambda)z)w\xi)(s t^{-1},t){\Delta(t)}^{-1/2}\\
    &=\int\alpha_{t s^{-1}}(z(r,s))(w\xi)(r^{-1}s t^{-1},t)\,dr{\Delta(t)}^{-1/2}\\
    &=\int\alpha_{t s^{-1}}(z(r,s))\xi(r^{-1}s,t)\,dr.
  \end{align*}
  Take a function \(z\) in \(C_C(G\times G, A)\) of the form
  \begin{equation}\label{Petersen:eq*}
    z(r,s)=\alpha_s(x)f(r^{-1}s)g(s)\Delta(r^{-1}s),
  \end{equation}
  where \(x\in A\) and \(f,g\in C_C(G)\). Inserting this into the expression above, we get
  \[
    (w^*((\tilde{\pi}\times\lambda)z)w\xi)(s,t)=\alpha_t(x)g(s)\int f(r)\xi(r,t)\,dr.
  \]
  Thus, if we define a faithful representation \((\check{\pi}, L^2(G, \mathcal{H}))\) of \(A\) by
  \((\check{\pi}(x)\xi)(t)=\alpha_t(x)\xi(t)\), then we may realize \(\check{\pi}(A)\otimes\mathcal{K}(L^2(G))\) as operators on 
  \(L^2(G\times G, \mathcal{H})\), and taking \(z\) as in (\ref{Petersen:eq*}), we see that
  \[
    w^*((\tilde{\pi}\times\lambda)z)w=\check{\pi}(x)\otimes v_{fg},
  \]
  where \(v_{fg}\) is the one-dimensional operator on \(L^2(G)\) such that \(v_{fg}\eta=(\eta|\overline{f})g\), \(\eta\in L^2(G)\).\par
  Since operators of the form \(v_{fg}, f,g\in C_C(G)\) generate \(\mathcal{K}(L^2(G))\), we see that
  \[
    w^*((\tilde{\pi}\times\lambda)(C_0(G,A)\rtimes_{\gamma,r}G))w\supset A\otimes\mathcal{K}(L^2(G)).
  \]
  On the other hand, each function in \(C_C(G\times G, A)\) can be uniformly approximated by finite sums of functions of the form 
  \((s,t)\mapsto \alpha_t(x)f(s)g(t)\), where \(x\in A\) and \(f,g\in C_C(G)\), and therefore (applying the homeomorphism \((s,t)\mapsto (s^{-1}t,t)\) of \(G\times G\)) also by functions of
  the form (\ref{Petersen:eq*}). Since \(C_C(G\times G, A)\) generates \(C_0(G,A)\rtimes_{\gamma,r}G\), it follows that we have found a spatial isomorphism between
  \(C_0(G,A)\rtimes_{\gamma,r}G\) and \(A\otimes\mathcal{K}(L^2(G))\).\par
  Consider the natural extension of \(\rho\) from \(C_0(G,A)\) to \(C_C(G,C_0(G,A))\) such that
  \(\rho_t(z)(r,s)=z(r,st)\) for each \(z\) in \(C_C(G,C_0(G,A))\). Then \(t\mapsto\rho_t\) is a representation
  of \(G\) as *-automorphisms of the dense *-subalgebra \(C_C(G,C_0(G,A))\)
  of \(C_0(G,A)\rtimes_{\gamma,r}G\). Moreover, with \(z\) as in (\ref{Petersen:eq*}), we have
  \[
    \rho_p(z)(r,s)=\alpha_{sp}(x)f(r^{-1}sp)g(sp)\Delta(r^{-1}sp),
  \]
  whence
  \begin{align*}
    (w^*((\tilde{\pi}\times\lambda)\rho_p(z))w\xi)(s,t)&=\alpha_{tp}(x)g(sp)\Delta(p)\int f(rp)\xi(r,t)\,dr\\
    &=\alpha_{tp}(x)(\tilde{\lambda}_p g)(s)\int(\tilde{\lambda}_p f)(r)\xi(r,t)\,dr\\
    &=((\check{\pi}(\alpha_p(x))\otimes\tilde{\lambda}_p v_{fg}\tilde{\lambda}_p^*)\xi)(s,t).
  \end{align*}
  Thus, identifying \(\check{\pi}(A)\) and \(A\), we have
  \[
    w^*((\tilde{\pi}\times\lambda)\rho_p(z))w=(\alpha\otimes\textnormal{Ad}\tilde{\lambda})_p(x\otimes v_{fg}).
  \]
  It follows that each \(\rho_p\) extends to an automorphism of \(C_0(G,A)\rtimes_{\gamma,r}G\) such that we obtain a \(C^*\)-dynamical system 
  \((C_0(G,A)\rtimes_{\gamma,r}G, G, \rho)\) covariantly isomorphic to the system \((A\otimes\mathcal{K}(L^2(G)),G,\alpha\otimes\textnormal{Ad}\tilde{\lambda})\).
\end{proof}

\section{Crossed products with abelian groups}
In this section, \(G\) will denote an abelian locally compact group, and \(\hat{G}\) its dual
group, and we will use the additive notation for the group operations. For \(t\) in
\(G\) and \(\sigma\) in \(\hat{G}\), we denote by \((t,\sigma)\) the evaluation of the character \(\sigma\) at \(t\), and
we define the (inverse) Fourier transform of a measure \(\mu\) in \(M(G)\) by \(\hat{\mu}(\sigma) =\int(t,\sigma)\,d\mu(t)\). 
Note that \(\hat{\mu} \in C^b( \hat{G})\) and that \(\hat{f} \in C_0(\hat{G})\) if \(f \in L^1(G)\).
\footnote{\(C^b(X)\) defines the set of the bounded continuous functions, \(C_0(X)\) is the subset of functions vanishing at infinity.}\par
As a reminder that \(G\) is now abelian, we refer to a system \((A,G,\alpha)\) as
a \(C^*\)-dynamical \textit{a-system}.
\begin{definition}
  Let \(G\) (and \(\hat{G}\)) be given. We say that a \cstar\ \(B\) is a \textnormal{\(G\)-product} if:
  \begin{enumerate}
    \item There is a homomorphism \(t\mapsto \lambda_t\) of \(G\) into the unitary group of \(M(B)\)
      such that each function \(t\mapsto \lambda_t y\), \(y\in B\), is continuous from \(G\) to \(B\);
    \item There is a homomorphism \(\sigma\mapsto\hat{\alpha}_\sigma\) of \(\hat{G}\) into \(\textnormal{Aut}(B)\)
      such that \((B,\hat{G},\hat{\alpha})\) is a \(C^*\)-dynamical system and
      \[
        \hat{\alpha}_\sigma(\lambda_t)=(t,\sigma)\lambda_t\qquad\forall t\in G,\ \sigma\in\hat{G}.
      \]
  \end{enumerate}\par
  Given a \(G\)-product \(B\), we say that an element \(x\) in \(M(B)\) satisfies \textnormal{Landstad's conditions} if:
  \begin{enumerate}
    \item \(\hat{\alpha}_\sigma(x)=x\) for all \(\sigma\) in \(\hat{G}\);
    \item \(x\lambda_f\in B\) and \(\lambda_f x\in B\) for every \(f\) in \(L^1(G)\);
    \item The map \(t\mapsto \lambda_t x\lambda_{-t}\) is continuous on \(G\).
  \end{enumerate}
\end{definition}
\begin{proposition}\label{Petersen:783}
  Let \((A,G,\alpha)\) be a \(C^*\)-dynamical system, and for \(y\) in \(C_C(G,A)\)
  and \(\sigma\) in \(\hat{G}\), define \((\hat{\alpha}_\sigma(y))(t)=(t,\sigma)y(t)\).
  Then the \(\hat{\alpha}_\sigma\)'s extend to automorphisms such that
  \((A\rtimes_\alpha G,\hat{G},\hat{\alpha})\) becomes a \(C^*\)-dynamical system. Moreover,
  \(A\rtimes_\alpha G\) is a \(G\)-product, and each element in \(A\) satisfies Landstad's conditions.
  We say that \((A\rtimes_\alpha G,\hat{G},\hat{\alpha})\) is the \textnormal{dual system} of \((A,G,\alpha)\).
\end{proposition}
\begin{proof}
  As \(G\) is abelian, it is amenable. Thus by \ref{Petersen:777} we may consider
  \(A\rtimes_{\alpha} G\) as a subalgebra of \(\mathcal{B}(L^2(G,\mathcal{H}))\),
  where \(\mathcal{H}\) is the universal Hilbert space for \(A\) (so that \(A\subset \mathcal{B}(\mathcal{H})\)).
  The unitary representation \(t\mapsto \lambda_t\) of \(G\) on \(L^2(G,\mathcal{H})\)
  (see \ref{Petersen:771}) is by \ref{Petersen:763} a homomorphism of \(G\) into \(M(A\rtimes_\alpha G)\)
  such that each function \(t\mapsto\lambda_t y\), \(y\in A\rtimes_\alpha G\), is norm continuous. Moreover,
  again from \ref{Petersen:763}, we have an injection of \(A\) into \(M(A\rtimes_\alpha G)\).
  Note that if \(\xi\in L^2(G,\mathcal{H})\), then
  \begin{equation}\label{Petersen:eq783a}
    (\lambda_s\xi)(t)=\xi(t-s)\qquad\textnormal{and}(x\xi)(t)=\alpha_{-t}(x)\xi(t)
  \end{equation}
  for all \(s\) in \(G\) and \(x\) in \(A\).\par
  Define the unitary representation \(u\) of \(\hat{G}\) on \(L^2(G,\mathcal{H})\) by
  \begin{equation}\label{Petersen:eq783b}
    (u_\sigma\xi)(t)=(t,\sigma)\xi(t).
  \end{equation}
  For each \(y\) in \(C_C(G,A)\), we have
  \begin{align*}
    (u_\sigma y u_{-\sigma})(t)&=(t,\sigma)(y u_{-\sigma}\xi)(t)\\
    &=(t,\sigma)\int \alpha_{-t}(y(s))((u_{-\sigma}\xi)(t-s))\,ds\\
    &=\int(s,\sigma)\alpha_{-t}(y(s))\xi(t-s)\,ds=((\sigma y)\xi)(t),
  \end{align*}
  where \((\sigma y)(t)=(t,\sigma)y(t)\). It follows immediately from this that we can
  define \(\hat{\alpha}_\sigma\) in \(\textnormal{Aut}(A\rtimes_\alpha G)\) by \(\hat{\alpha}_\sigma(y)=u_\sigma y u_{-\sigma}\),
  \(y\in A\rtimes_\alpha G\), and that \((A\rtimes_\alpha G,\hat{G},\hat{\alpha})\) becomes a \(C^*\)-dynamical system. Moreover, from 
  (\ref{Petersen:eq783a}) and (\ref{Petersen:eq783b}) we see that
  \[
    \hat{\alpha}_\sigma(\lambda_s)=(s,\sigma)\lambda_s\qquad\textnormal{and}\qquad\hat{\alpha}_\sigma(x)=x
  \]
  for all \(s\) in \(G\) and \(x\) in \(A\). We have verified that \(A\rtimes_\alpha G\) is a \(G\)-product
  and that the elements of \(A\) satisfy the first and of course also the third of Landstad's conditions (as \(\lambda_t x\lambda_{-t}=\alpha_t(x)\)).\par
  A straightforward computation shows that if \(x\in A\) and \(f\in C_C(G)\), then
  \(x\lambda_f=L(x,f)\), with \(L\) as in \ref{Petersen:762}. Moreover, in this case, \(L(x,f)\in C_C(G,A)\);
  in fact, \(L(x,f)=x\otimes f\). It follows that \(x\lambda_f\in A\rtimes_\alpha G\)
  for all \(f\) in \(L^1(G)\). Since \(\lambda_f x={(x^*\lambda_{f^*})}^*\in A\rtimes_\alpha G\), we have verified also the second condition. 
\end{proof}

\section{Duality theory for crossed products}
Given a \cstar\ \(A\) and an abelian group \(G\) with dual group \(\hat{G}\), define a \(C^*\)-dynamical system
\((A,\hat{G},\iota)\) by taking \(\iota_\sigma(x)=x\) for all \(\sigma\) in \(\hat{G}\) and \(x\) in \(A\).
If \(A\subset \mathcal{B}(\mathcal{H})\), then consider the regular representation of \(A\rtimes_\iota\hat{G}\) on \(L^2(\hat{G},\mathcal{H})\), where
\[
  (x\eta)(\tau)=\int x(\sigma)\eta(\tau-\sigma)\,d\sigma,\qquad x\in C_C(\hat{G},A),\eta\in L^2(\hat{G},\mathcal{H}).
\]
Define a faithful representation of \(C_0(G,A)\) on \(L^2(G,\mathcal{H})\) by
\[
  (y\xi)(t)=y(t)\xi(t),\qquad y\in C_0(G,A), \xi\in L^2(G,\mathcal{H}).
\]
Moreover, for each \(x\) in \(C_C(\hat{G},A)\), define \(\hat{x}\) in \(C_0(G,A)\) by \(\hat{x}(t)=\int x(\sigma)(t,\sigma)\,d\sigma\).\label{Petersen:791}
Now let \(u\) be the isometry of \(L^2(G,\mathcal{H})\) onto \(L^2(\hat{G},\mathcal{H})\), where
\[
  (u\xi)(\tau)=\int\xi(t)\overline{(t,\tau)}\,dt,\qquad(u^* \eta)(t)=\int\eta(\tau)(t,\tau)\,d\tau,
\]
and the integrals converge in the \(L^2\)-sense. Then, for each \(x\) in \(C_C(\hat{G},A)\), we have
\begin{align*}
  (xu\xi)(\tau)&=\int x(\sigma)(u\xi)(\tau-\sigma)\,d\sigma\\
    &=\int\int x(\sigma)\xi(t)(t,\sigma)\overline{(t,\tau)}\,dt\,d\sigma\\
  &=\int\hat{x}\xi(t)\overline{(t,\tau)}\,dt=(u\hat{x}\xi)(\tau).
\end{align*}
Consequently, \(u^* x u=\hat{x}\). Since \(A\rtimes_\iota\hat{G}\) is the closure of \(C_C(\hat{G},A)\) in \(\mathcal{B}(L^2(\hat{G},\mathcal{H}))\),
it follows that \(A\rtimes_\iota\hat{G}\) is isomorphic to \(C_0(G,A)\).
\begin{lemma}\label{Petersen:792}
  Given a \(C^*\)-dynamical system \((A,G,\alpha)\), consider the dual system \((A\rtimes_\alpha G,\hat{G},\hat{\alpha})\)
  defined in \ref{Petersen:783} and the trivial system \((A,\hat{G},\iota)\) defined above.
  There exists a \(C^*\)-dynamical system \((A\rtimes_\iota\hat{G},G,\beta)\) such that \((\beta_t(x))(\tau)=\overline{(t,\tau)}\alpha_t(x(\tau))\)
  for each \(x\) in \(C_C(\hat{G},A)\). Moreover, the \(C^*\)-algebras \((A\rtimes_\iota\hat{G})\rtimes_\beta G\) and
  \((A\rtimes_\alpha G)\rtimes_{\hat{\alpha}}\hat{G}\) are isomorphic; compare
  \cite[Lemmas 7.2 and 7.3]{Williams2007}.
\end{lemma}
\begin{proof} For each \(t\in G\), define \(\beta_t\) on \(C_C(\widehat G,A)\) by
  \[ (\beta_t(x))(\tau) = \overline{(t,\tau)}\,\alpha_t(x(\tau)), \]   
  and define \(\gamma_t\) on \(C_0(G,A)\) by 
  \[ (\gamma_t(y))(s) = \alpha_t\bigl(y(s-t)\bigr). \] 
  It is easy to verify that each \(\gamma_t\) is an automorphism of \(C_0(G,A)\) and that 
  \[ (C_0(G,A),G,\gamma) \] is a \(C^*\)-dynamical system. Moreover, if 
  \[ u:L^2(G,\mathcal{H})\longrightarrow L^2(\widehat G,\mathcal{H}) \] is the Fourier transform as above,
  a short computation shows that 
  \[ \hat{\beta_t(x)} = \gamma_t(\hat x), \qquad x\in C_C(\hat G,A). \]
  It follows that each \(\beta_t\) extends to an automorphism of \(A\rtimes_\iota\widehat G\), and hence
  \[ (A\rtimes_\iota\widehat G,G,\beta) \] becomes a \(C^*\)-dynamical system covariantly isomorphic to
  \[ (C_0(G,A),G,\gamma). \] The regular representation of 
  \[ (A\rtimes_\iota\widehat G)\rtimes_\beta G \] induced from the universal representation of \(A\) on \(\mathcal H\) acts on
  \[ L^2(G\times\widehat G,\mathcal H) \cong L^2\bigl(G,L^2(\widehat G,\mathcal H)\bigr). \]
  For \(x\in C_C(G\times\widehat G,A)\) and \(\xi\in L^2(G\times\widehat G,\mathcal H)\), we have
  \begin{align*} (x\xi)(t,\tau) &= \int_G \beta_{-t}\bigl(x(s,\cdot)\bigr) \,\xi(t-s,\cdot)(\tau)\,ds \\ 
    &= \int_G \int_{\widehat G} \beta_{-t}\iota_{-\tau}\bigl(x(s,\sigma)\bigr) \,\xi(t-s,\tau-\sigma)\,d\sigma\,ds \\
    &= \int_G \int_{\widehat G} \alpha_{-t}\bigl(x(s,\sigma)\bigr) \,\xi(t-s,\tau-\sigma)\, (t,\sigma)\,d\sigma\,ds .
  \end{align*} 
  Similarly, the regular representation of 
  \[ (A\rtimes_\alpha G)\rtimes_{\hat\alpha}\widehat G \] acts on \(L^2(\widehat G\times G,\mathcal H)\), and for
  \(y\in C_C(\widehat G\times G,A)\) and \(\eta\in L^2(\widehat G\times G,\mathcal H)\),
  \begin{align*} 
    (y\eta)(\tau,t) &= \int_{\widehat G} \hat\alpha_{-\tau}\bigl(y(\sigma,\cdot)\bigr) \,\eta(\tau-\sigma,\cdot)(t)\,d\sigma \\
    &= \int_{\widehat G}\int_G \hat\alpha_{-\tau}\alpha_{-t} \bigl(y(\sigma,s)\bigr) \,\eta(\tau-\sigma,t-s)\,ds\,d\sigma \\
    &= \int_{\widehat G}\int_G \alpha_{-t}\bigl(y(\sigma,s)\bigr) \,\eta(\tau-\sigma,t-s)\, \overline{(s,\tau)}\,ds\,d\sigma .
  \end{align*} 
  Define 
  \[ \Phi:C_C(G\times\widehat G,A) \longrightarrow C_C(\widehat G\times G,A) \] by
  \[ (\Phi x)(\tau,t) = (t,\tau)\,x(t,\tau), \] with inverse
  \[ (\Phi^{-1}y)(t,\tau) = \overline{(t,\tau)}\,y(\tau,t). \]
  Furthermore, define an isometry 
  \[ w:L^2(G\times\widehat G,\mathcal H) \longrightarrow L^2(\widehat G\times G,\mathcal H) \] by
  \[ (w\xi)(\tau,t) = \overline{(t,\tau)}\,\xi(t,\tau), \] and 
  \[ (w^*\eta)(t,\tau) = (t,\tau)\,\eta(\tau,t). \] 
  Then for every \(x\in C_C(G\times\widehat G,A)\) and \(\xi\in L^2(G\times\widehat G,\mathcal H)\),
  \begin{align*} 
    ((\Phi x)w\xi)(\tau,t) &= \int_{\widehat G}\int_G \alpha_{-t}\bigl((\Phi x)(\sigma,s)\bigr) \,w\xi(\tau-\sigma,t-s) \,\overline{(s,\tau)} \,ds\,d\sigma \\
    &= \int_{\widehat G}\int_G \alpha_{-t}\bigl(x(s,\sigma)\bigr) \,\xi(t-s,\tau-\sigma)\, (s,\sigma)\, \overline{(t-s,\tau-\sigma)}\, \overline{(s,\tau)} \,ds\,d\sigma \\
    &= \overline{(t,\tau)} \int_{\widehat G}\int_G \alpha_{-t}\bigl(x(s,\sigma)\bigr) \,\xi(t-s,\tau-\sigma)\, (t,\sigma) \,ds\,d\sigma \\
    &= (w(x\xi))(\tau,t). 
  \end{align*}
  Hence \[ w^*(\Phi x)w=x. \] Since \(C_C(\widehat G\times G,A)\) and \(C_C(G\times\widehat G,A)\) are dense in
  \[ (A\rtimes_\alpha G)\rtimes_{\hat\alpha}\widehat G \quad\text{and}\quad (A\rtimes_\iota\widehat G)\rtimes_\beta G, \]
  respectively, it follows that the two \(C^*\)-algebras are isomorphic. 
\end{proof}


\begin{theorem}[Takai Duality]\label{theorem:Takai}\label{Petersen:793}
  Let \((A, G, \alpha)\) be a \(C^*\)-dynamical system with \(G\) abelian. Then \((A\rtimes_\alpha G)\rtimes_{\hat{\alpha}} \hat{G}\cong A\otimes \mathcal{K}(L^2(G))\).
  So the \textit{double} crossed product is stably isomorphic to \(A\), and is stable if \(G\) is infinite (if \(G\) has \(n\) elements, it is \(M_n(A)\)).
  Under this isomorphism, the second dual action \(\hat{\hat{\alpha}}\) of \(G\) on the second crossed product is \(\alpha\otimes\textnormal{Ad}\tilde{\lambda}\), 
  where \(\textnormal{Ad}\tilde{\lambda}\) is the action of \(G\)
  on \(\mathcal{K}(L^2(G))\) coming from the regular representation of \(G\).
\end{theorem}
\begin{proof}
  From \ref{Petersen:7712} we know that \((A\otimes\mathcal{K}(L^2(G)),G,\alpha\otimes \textnormal{Ad}\tilde{\lambda})\) is covariantly
  isomorphic to \((C_0(G,A)\rtimes_\gamma G,G,\rho)\), where \(\rho_t(z)(r,s)=z(r,s+t)\) for each 
  \(z\) in \(C_C(G,C_0(G,A))\), and \(\gamma\) is as the proof of \ref{Petersen:792}.
  As we saw there, the system \((C_0(G,A),G,\gamma)\) is covariantly isomorphic to \((A\rtimes_\iota\hat{G},G,\beta)\), so that
  \(C_0(G,A)\rtimes_\gamma G\) is isomorphic to \((A\rtimes_\iota\hat{G})\rtimes_\beta G\), which by \ref{Petersen:792} is isomorphic to
  \((A\rtimes_\alpha G)\rtimes_{\hat{\alpha}}\hat{G}\).\par
  Now take a function \(z\) in \(C_C(G,C_0(G,A))\) of the form \(z(s,t)=xf(s)\hat{g}(t)\),
  where \(x\in A\), \(f\in C_C(G)\), and \(g\in C_C(\hat{G})\). Regarding \(z\) as an element of \(G\times C_0(G,A)\), its image
  in \((A\rtimes_\iota\hat{G})\rtimes_\beta G\) is the function \(\check{z}\) in \(C_C(G\times \hat{G},A)\) defined by
  \(\check{z}(s,\sigma)=xf(s)g(\sigma)\) (cf. \ref{Petersen:791}). Furthermore,
  the image of \(\check{z}\) in \((A\rtimes_\alpha G)\rtimes_{\hat{\alpha}}\hat{G}\) is the function \(\Phi(\check{z})\) in \(C_C(\hat{G}\times G,A)\) (cf. the proof of \ref{Petersen:792}), where
  \[
    \Phi(\check{z})(\sigma,s)=(s,\sigma)\check{z}(s,\sigma)=xf(s)g(\sigma)(s,\sigma).
  \]
  We have \(\rho_r(z)(s,t)=xf(s)\hat{g}(t+r)\), whence \(\check{\rho_r(z)}(s,\sigma)=xf(s)g(\sigma)(r,\sigma)\),
  and consequently
  \[
    \Phi(\check{\rho_r(z)})(\sigma,s)=xf(s)g(\sigma)(s+r,\sigma)=(\Phi(\check{z})(\sigma,s))(r,\sigma).
  \]
  However, if \(y\in C_C(\hat{G}\times G,A)\), then regarding \(y\) as an element of \((A\rtimes_\alpha G)\rtimes_{\hat{\alpha}}\hat{G}\), we have by \ref{Petersen:783}
  \[
    \hat{\hat{\alpha}}_r(y)(\sigma,s)=(r,\sigma)y(\sigma,s).
  \]
  It follows that \(\hat{\hat{\alpha}}_r(\Phi(\check{z}))=\Phi(\check{\rho_r(z)})\), and since functions of the form \(z\) generate
  \(C_0(G,A)\rtimes_\gamma G\), we conclude that the systems \((C_0(G,A)\rtimes_\gamma G,G,\rho)\) and 
  \(((A\rtimes_\alpha G)\rtimes_{\hat{\alpha}}\hat{G},G,\hat{\hat{\alpha}})\) are covariantly isomorphic, completing the proof.
\end{proof}

\section{The crossed product \texorpdfstring{\(C(S^1)\rtimes \mathbb{Z}\)}{CS1xZ}}
Let \(\theta\) in \(\mathbb {R}\) and define
\[
R_\theta(z)=e^{2\pi i\theta}z,
\qquad z\in S^1.
\]
This determines a $C^*$-dynamical system
\[
(C(S^1),\mathbb Z,\alpha_\theta),
\]
where the action is given by
\[
(\alpha_\theta)_n(f)=f\circ R_\theta^{-n},\qquad n\in\mathbb Z.
\]
Let
\[
u\in C(S^1)
\]
denote the identity function
\[
u(z)=z.
\]
\begin{lemma}
\((\alpha_\theta)_1(u)=e^{-2\pi i\theta}u\).
\end{lemma}
\begin{proof}
For every $z\in S^1$,
\[
(\alpha_\theta)_1(u)(z)
=
u(e^{-2\pi i\theta}z)
=
e^{-2\pi i\theta}z
=
e^{-2\pi i\theta}u(z).
\]
\end{proof}
Let $V$ denote the unitary in the crossed product corresponding to the generator $1\in\mathbb Z$.
\begin{proposition}
For every $f\in C(S^1)$,
\[
VfV^*
=
(\alpha_\theta)_1(f).
\]
\end{proposition}
\begin{proof}
This is precisely the covariance relation defining the crossed product.
\end{proof}
\begin{corollary}
\[
VuV^*
=
e^{-2\pi i\theta}u.
\]
\end{corollary}
\begin{proof}
Apply the previous proposition to $u$ and use the lemma.
\end{proof}
\begin{corollary}\label{corollary:344}
\[
Vu=e^{-2\pi i\theta}uV.
\]
\end{corollary}
\begin{proof}
Multiply
\[
VuV^*
=
e^{-2\pi i\theta}u
\]
on the right by $V$.
\end{proof}
The relation
\[
Vu=e^{-2\pi i\theta}uV
\]
will motivate the definition of the noncommutative torus in the next chapter.





\chapter{The Noncommutative Torus}




\section{Definition}
\begin{definition}
Let \(\theta\) in \(\mathbb R\).
A pair of unitaries \(u,v\) in a \cstar\ \(B\) is called
a \textnormal{\(\theta\)-commuting pair} if
\[
vu=e^{2\pi i\theta}uv.
\]
\end{definition}
\begin{definition}
The \textnormal{noncommutative torus} \(A_\theta\) is the universal unital \(C^*\)-algebra generated by two unitaries
\(U\) and \(V\) satisfying
\[
VU=e^{2\pi i\theta}UV.
\]
Thus, if \(B\) is a unital \(C^*\)-algebra and \(u,v\in\mathcal U(B)\)
satisfy
\[
vu=e^{2\pi i\theta}uv,
\]
then there exists a unique unital *-homomorphism
\[
\pi:A_\theta\longrightarrow B
\]
such that \(\pi(U)=u\) and \(\pi(V)=v\).
\end{definition}
\begin{remark}
The existence of the universal \(C^*\)-algebra \(A_\theta\) follows
from the standard construction of universal \(C^*\)-algebras generated
by unitaries subject to norm-bounded relations; see \cite[Chapter~VI, Section~1]{Davidson1996}. Its universal property
determines it uniquely up to a canonical *-isomorphism.
A concrete realization will be obtained below as a crossed product.
\end{remark}
\begin{remark}
If \(\theta=0\), the relation becomes
\[
UV=VU.
\]
Hence \(A_0\) is commutative.
\end{remark}









\section{Realization as a crossed product}
\begin{theorem}\label{theorem:noncommutativeTorusCrossedProduct}
Let \(\alpha_\theta:\mathbb Z\to\operatorname{Aut}(C(S^1))\) be the action determined by
\[
(\alpha_\theta)_1(f)=f\circ R_\theta^{-1}.
\]
Then
\[
A_\theta
\cong
C(S^1)\rtimes_{\alpha_\theta}\mathbb Z .
\]
\end{theorem}
\begin{proof}
  Let \(z\) in \(C(S^1)\) be the identity function
  \[
  z(\lambda)=\lambda.
  \]
  Let \(W\) denote the canonical unitary in the crossed product corresponding to the generator of \(\mathbb Z\). By Corollary \ref{corollary:344},
  \[
  Wz=e^{-2\pi i\theta}zW.
  \]
  Equivalently,\[
  W^*z=e^{2\pi i\theta}zW^*.
  \]
  Since \(z\) and \(W^*\) are unitaries satisfying the defining
  relation of \(A_\theta\), the universal property yields a *-homomorphism
  \[
  \varphi:A_\theta\to C(S^1)\rtimes_{\alpha_\theta}\mathbb Z
  \]
  such that
  \[
  \varphi(U)=z,
  \qquad
  \varphi(V)=W^*.
  \]
  The crossed product is generated by \(C(S^1)\) together with \(W\). Since \(z\) generates \(C(S^1)\), the image 
  of \(\varphi\) contains both generators, and therefore \(\varphi\) is surjective. Define
  \[
  \pi:C(S^1)\to A_\theta
  \]
  by
  \[
  \pi(f)=f(U).
  \]
  Since
  \[
  V^*UV=e^{-2\pi i\theta}U,
  \]
  the pair \((\pi,V^*)\) is a covariant representation of \((C(S^1),\mathbb Z,\alpha_\theta)\). Hence there exists a *-homomorphism
  \[
  \psi:C(S^1)\rtimes_{\alpha_\theta}\mathbb Z\to A_\theta
  \]
  such that
  \[
  \psi(z)=U,
  \qquad
  \psi(W)=V^*.
  \]
  Now
  \[
  (\psi\circ\varphi)(U)=U,
  \qquad
  (\psi\circ\varphi)(V)=V.
  \]
  Since \(A_\theta\) is generated by \(U\) and \(V\), it follows that
  \[
  \psi\circ\varphi=id_{A_\theta}.
  \]
  Similarly,
  \[
  (\varphi\circ\psi)(z)=z,
  \qquad
  (\varphi\circ\psi)(W)=W.
  \]
  Since \(C(S^1)\rtimes_{\alpha_\theta}\mathbb Z\) is generated by \(z\) and \(W\),
  \[
  \varphi\circ\psi=id_{C(S^1)\rtimes_{\alpha_\theta}\mathbb Z}.
  \]
  Therefore \(\varphi\) is an isomorphism; see \cite[Sections~2.2, 2.3 and 2.6]{Williams2007}.
\end{proof}

\begin{proposition}[Gauge action]\label{prop:gaugeAction}
There is a strongly continuous action
\[
\gamma:\mathbb T^2\longrightarrow\operatorname{Aut}(A_\theta)
\]
uniquely determined by
\[
\gamma_{s,t}(U)=sU,
\qquad
\gamma_{s,t}(V)=tV
\]
for every \((s,t)\in\mathbb T^2\).
\end{proposition}

\begin{proof}
For fixed \((s,t)\in\mathbb T^2\), the unitaries \(sU\) and \(tV\)
satisfy the same commutation relation as \(U\) and \(V\). The universal
property therefore gives a unital *-homomorphism \(\gamma_{s,t}\) with the
stated values on the generators. Its inverse is
\(\gamma_{\overline{s},\overline{t}}\), so it is an automorphism. The group
law follows by checking it on \(U\) and \(V\). Finally, continuity is
immediate on finite linear combinations of the monomials \(U^mV^n\), and
then follows on all of \(A_\theta\) by density and the fact that every
\(\gamma_{s,t}\) is isometric. This is the canonical gauge action on the
rotation algebra; compare \cite[Chapter~VI, Section~1]{Davidson1996}.
\end{proof}











\section{Rational and irrational parameters}

The structure of \(A_\theta\) depends strongly on whether
\(\theta\) is rational or irrational.

\begin{proposition}
If \(\theta=0\), then
\[
A_0\simeq C(\mathbb T^2).
\]
\end{proposition}

\begin{proof}
Let \(z_1,z_2\in C(\mathbb T^2)\) be the coordinate functions,
\[
z_1(\lambda,\mu)=\lambda,
\qquad
z_2(\lambda,\mu)=\mu.
\]
They are commuting unitaries and generate \(C(\mathbb T^2)\).
Hence the universal property of \(A_0\) yields a unital
*-homomorphism
\[
\Phi:A_0\longrightarrow C(\mathbb T^2)
\]
such that
\[
\Phi(U)=z_1,\qquad \Phi(V)=z_2.
\]

Conversely, \(C(\mathbb T^2)\) is the universal unital
\(C^*\)-algebra generated by two commuting unitaries. Therefore
there exists a unital *-homomorphism
\[
\Psi:C(\mathbb T^2)\longrightarrow A_0
\]
such that
\[
\Psi(z_1)=U,\qquad \Psi(z_2)=V.
\]
Both compositions are the identity on the corresponding generators.
Consequently \(\Phi\) and \(\Psi\) are mutually inverse, and
\[
A_0\cong C(\mathbb T^2).
\]
\end{proof}

\begin{proposition}\label{prop:rationalTheta}
If \(\theta\in\mathbb Q\), then \(A_\theta\) is not simple.
\end{proposition}

\begin{proof}
Let
\[
\theta=\frac pq
\]
with \(p\) and \(q\) relatively prime. Then
\[
e^{2\pi i q\theta}=1.
\]
Using the relation
\[
VU=e^{2\pi i\theta}UV,
\]
it follows that
\[
VU^q=U^q V.
\]
Thus \(U^q\) belongs to the center of \(A_\theta\). It is not a scalar.
Indeed, if \(U^q=\lambda1\), then Proposition~\ref{prop:gaugeAction} gives
\(s^q\lambda1=\lambda1\) for every \(s\in\mathbb T\). Since \(U^q\) is
unitary, \(\lambda\ne0\), and this would imply \(s^q=1\) for every
\(s\in\mathbb T\), which is impossible. Since \(U^q\) is not
a scalar, the center is non-trivial. A simple unital \cstar\
has trivial center, therefore \(A_\theta\) cannot be simple.
\end{proof}
\begin{theorem}\label{thm:IrrationalTorusSimple}
If \(\theta\) is irrational, then \(A_\theta\) is simple.
\end{theorem}
\begin{proof}
  This result follows from the realization of \(A_\theta\) as a crossed product, see \cite[Theorem~VI.1.4]{Davidson1996};
  see also \cite[p.~415]{Rieffel1981}.
\end{proof}
\begin{corollary}
If \(\theta\) is irrational, the center of \(A_\theta\) is equal to
\[
\mathbb C1.
\]
\end{corollary}

\begin{proof}
The center of a simple unital \cstar\ is trivial.
\end{proof}













\section{The canonical trace}
Let
\[
\mathcal A_\theta
=
\operatorname{span}\{U^mV^n:m,n\in\mathbb Z\}.
\]
This is a dense unital *-subalgebra of \(A_\theta\). The gauge action
provides both the Fourier coefficients on \(\mathcal A_\theta\) and the
canonical trace on its \(C^*\)-completion. The construction below is
standard; see \cite[Chapter~VI, Section~1]{Davidson1996} and
\cite[p.~415]{Rieffel1981}.

Throughout this section, integration over \(\mathbb T^2\) is with respect
to normalized Haar measure \(\mu\). For \((k,l)\in\mathbb Z^2\), define
\[
P_{k,l}(a)
=
\int_{\mathbb T^2}s^{-k}t^{-l}\gamma_{s,t}(a)\,d\mu(s,t),
\qquad a\in A_\theta,
\]
where the integral is taken in norm.

\begin{proposition}\label{prop:uniqueFourierPolynomial}
Every element of \(\mathcal A_\theta\) has a unique expression
\[
a=\sum_{m,n}a_{m,n}U^mV^n,
\]
where only finitely many coefficients are nonzero. More precisely,
\[
P_{k,l}(a)=a_{k,l}U^kV^l.
\]
\end{proposition}

\begin{proof}
Since
\[
\gamma_{s,t}(U^mV^n)=s^mt^nU^mV^n,
\]
orthogonality of the characters of \(\mathbb T^2\) gives the formula for
\(P_{k,l}(a)\). If a finite linear combination of the monomials is zero,
applying \(P_{k,l}\) shows that its \((k,l)\)-coefficient is zero. Hence all
coefficients vanish, which proves uniqueness.
\end{proof}

Consider now the average of the gauge action,
\[
E(a)=\int_{\mathbb T^2}\gamma_{s,t}(a)\,d\mu(s,t),
\qquad a\in A_\theta.
\]
It is a positive unital contraction. On \(\mathcal A_\theta\) it satisfies
\[
E\left(\sum_{m,n}a_{m,n}U^mV^n\right)=a_{0,0}1.
\]
Since \(\mathcal A_\theta\) is dense and \(\mathbb C1\) is closed, the
range of \(E\) is exactly \(\mathbb C1\). Consequently there is a unique
linear functional \(\tau:A_\theta\to\mathbb C\) such that
\[
E(a)=\tau(a)1,
\qquad a\in A_\theta.
\]

\begin{theorem}\label{thm:canonicalTraceRotationAlgebra}
The functional \(\tau\) is a faithful tracial state on \(A_\theta\). It is
characterized on the monomials by
\[
\tau(U^mV^n)
=
\begin{cases}
1,&(m,n)=(0,0),\\
0,&(m,n)\ne(0,0).
\end{cases}
\]
Moreover, if
\(a=\sum_{m,n}a_{m,n}U^mV^n\in\mathcal A_\theta\), then
\[
\tau(a^*a)=\sum_{m,n}|a_{m,n}|^2.
\]
\end{theorem}

\begin{proof}
Because \(E\) is positive and unital, so is \(\tau\); in particular,
\(\lVert\tau\rVert=\tau(1)=1\). Thus boundedness and extension from
\(\mathcal A_\theta\) do not have to be established separately.

To prove the trace identity, it is enough by density and continuity to
check it on monomials. From
\[
U^mV^nU^pV^q
=
e^{2\pi i\theta np}U^{m+p}V^{n+q}
\]
we see that both
\(\tau(U^mV^nU^pV^q)\) and
\(\tau(U^pV^qU^mV^n)\) vanish unless
\(m+p=n+q=0\). In the remaining case both are equal to
\(e^{-2\pi i\theta mn}\). Hence \(\tau(ab)=\tau(ba)\) first on
\(\mathcal A_\theta\), and therefore on \(A_\theta\).

For a finite Fourier polynomial \(a\), expansion of \(a^*a\) and extraction
of its \((0,0)\)-coefficient give
\[
\tau(a^*a)=\sum_{m,n}|a_{m,n}|^2.
\]
This proves, in particular, faithfulness on \(\mathcal A_\theta\). To prove
faithfulness on the whole \(C^*\)-algebra, let \(a\in A_\theta\) be positive
and suppose that \(E(a)=0\). For every state \(\varphi\) on \(A_\theta\),
the function
\[
(s,t)\longmapsto\varphi(\gamma_{s,t}(a))
\]
is continuous and nonnegative, and its integral is
\(\varphi(E(a))=0\). It must therefore vanish identically. Evaluating at
\((1,1)\) gives \(\varphi(a)=0\) for every state \(\varphi\), whence
\(a=0\). Thus \(E\), and consequently \(\tau\), is faithful.
\end{proof}

\begin{corollary}\label{cor:uniqueTraceIrrationalRotation}
If \(\theta\) is irrational, then \(\tau\) is the unique tracial state on
\(A_\theta\).
\end{corollary}

\begin{proof}
Let \(\sigma\) be a tracial state on \(A_\theta\). Invariance under inner
automorphisms gives
\[
\sigma(U^mV^n)
=
\sigma(VU^mV^nV^*)
=
e^{2\pi i\theta m}\sigma(U^mV^n).
\]
If \(m\ne0\), irrationality of \(\theta\) implies
\(e^{2\pi i\theta m}\ne1\), and hence \(\sigma(U^mV^n)=0\). If \(m=0\)
and \(n\ne0\), conjugation by \(U\) gives the same conclusion. Thus
\(\sigma\) and \(\tau\) agree on every monomial, hence on the dense
subalgebra \(\mathcal A_\theta\), and therefore on \(A_\theta\). This is
the classical uniqueness property of the normalized trace on the
irrational rotation algebra; see \cite[p.~415]{Rieffel1981}.
\end{proof}












\chapter{Projections and Unitary Elements}
\section{Homotopy classes of unitary elements}
\begin{definition}
  Let X be a topological space. Two points \(a,b\) in \(X\) are 
  \textnormal{homotopic} in \(X\), written \(a \thicksim_h b\) in \(X\), if 
  there is a continuous function \(v:[0,1]\rightarrow X\) such that \(v(0)=a,v(1)=b\).\par
  The continuous function \(v\) is called a \textnormal{continuous path} from
  \(a\) to \(b\), and it is often denoted by \(t \mapsto v_t\).
\end{definition}
\begin{proposition}
  The relation \(\thicksim_h\) is an equivalence relation on \(X\).
\end{proposition}
\begin{proof}
  \begin{enumerate}
    \item If \(a\) is in \(X\), we define \(v(t)=a\) for all \(t\).
      We have \(v(0)=v(1)=a\) hence \(a\thicksim_h a\).
    \item If \(a\thicksim_h b\), we have a continuous function \(v_1\) 
    so that \(v_1(0)=a, v_1(1)=b\).\\ 
    Define \(v_2(t) = v_1(1-t)\).
    \(v_2\) is continuous because it is a composition of continuous functions, 
    \(v_2(0) =v_1(1)=b\) and \(v_2(1)=v_1(0)=a\), hence \(b\thicksim_h a\).
    \item If \(a\thicksim_h b\) and \(b\thicksim_h c\) via functions \(v_1, v_2\). Define
    \[u(t)=\begin{cases}v_1(2t), & 0\le t\le \frac12,\\v_2(2t-1), & \frac12\le t\le1\end{cases}\]
    The two cases are continuous and \(v_1(2\cdot\frac{1}{2})=b=v_2(2\cdot \frac{1}{2}-1)\) so \(u\) is continuous in \([0,1]\). 
    We have \(u(0)=v_1(0)=a\) and \(u(1)=v_2(1)=c\) so \(a\thicksim_h c\).
  \end{enumerate}
\end{proof}
\begin{definition}
  Let \(A\) be a unital \cstar\ and \(\mathcal{U}(A)\) be the subset of \(A\) of all unitary elements (\(u\) is unitary if \(u^*u=uu^*=1\)).\\
  We define \(\mathcal{U}_0(A)=\{u\in \mathcal{U}(A)|u\thicksim_h 1\mbox{ in }\mathcal{U}(A)\}\).\footnote{It is 
  important to note that in \(A\) every two elements are homotopic via the function \(v(t) = tb+(1-t)a\). 
  This is not true for \(\mathcal{U}(A)\).} 
\end{definition}
\begin{lemma}\label{lemma:514}\label{lemma:Matthes214}
  Let \(A\) be a unital \cstar.
  \begin{enumerate}
    \item\label{item:5141} For each self-adjoint element \(h\) in \(A\), \(e^{ih}\) belongs to \(\mathcal{U}_0(A)\).
    \item\label{item:5142} If \(u\) is a unitary element in \(A\) with \(\operatorname{spec}(u)\neq\mathbb{T}\), then \(u\) belongs to \(\mathcal{U}_0(A)\).
    \item If \(u,v\) are unitary elements in \(A\) with \(\|u-v\|<2\), then \(u\thicksim_h v\).
  \end{enumerate}
\end{lemma}
\begin{proof}
  \begin{enumerate}
    \item Let \(f:\mathbb R\to\mathbb T\) be a continuous function. Since
      \(
      f(x)\overline{f(x)}=1
      \)
      for every \(x\in\mathbb R\), Corollary~\ref{corollary:149} yields
      \(
      {f(h)}^*
      =
      \overline f(h).
      \)
      Therefore
      \(
      {f(h)}^*f(h)
      =
      (\overline f\,f)(h)
      =
      1
      \)
      and similarly
      \(
      f(h){f(h)}^*
      =
      1.
      \)
      Hence \(f(h)\) is unitary.
      Now, for each \(t\in[0,1]\), define
      \[
      f_t(x)=e^{itx},
      \qquad x\in \operatorname{spec}(h).
      \]
      Since \(t\mapsto f_t\) is continuous in \(C(\operatorname{spec}(h))\), the continuity of the functional calculus implies that
      \(
      t\mapsto f_t(h)
      \)
      is a continuous path in \(\mathcal{U}(A)\). Moreover,
      \(
      f_0(h)=1
      \)
      and
      \(
      f_1(h)=e^{ih}.
      \)
      Therefore
      \(
      e^{ih}\sim_h 1,
      \)
      and hence
      \(
      e^{ih}\in\mathcal U_0(A).
      \)
    \item Suppose that \(\operatorname{spec}(u)\neq\mathbb T\). Then there exists
      \(e^{i\theta}\in\mathbb T\setminus \operatorname{spec}(u)\). Let
      \(\phi:\operatorname{spec}(u)\rightarrow\mathbb R\) be defined by
      \(\phi(e^{it})=t\), where \(t\in(\theta,\theta+2\pi)\).
      Since \(e^{i\theta}\notin \operatorname{spec}(u)\), the argument function is
      well-defined and continuous on \(\operatorname{spec}(u)\). Moreover,
      \(z=e^{i\phi(z)}\) for every \(z\in \operatorname{spec}(u)\). By the continuous
      functional calculus, the element \(h=\phi(u)\) is self-adjoint.
      Applying the functional calculus again, we obtain
      \(e^{ih}=(e^{i\phi})(u)=u\). By part (\ref{item:5141}), \(e^{ih}\in\mathcal U_0(A)\).
      Therefore \(u\in\mathcal U_0(A)\).
    \item If \(\|u-v\|<2\), then \(\|v^*u-1\|=\|v^*(u-v)\|<2\) because 
      each unitary element has norm \(1\). It follows that \(-2\) is not in 
      \(\operatorname{spec}(v^*u-1)\) by Theorem~\ref{theorem:135}. This shows that \(-1\) does not belong to \(\operatorname{spec}(v^*u)\).
      We can then use (\ref{item:5142}) to conclude that \(v^*u\thicksim_h 1\),
      and hence that \(u\thicksim_h v\) as desired.
  \end{enumerate}
\end{proof}
\begin{corollary}
  The unitary group in \(M_n(\mathbb{C})\) is connected, in other words, \(\mathcal{U}_0(M_n(\mathbb{C})) = \mathcal{U}(M_n(\mathbb{C}))\).
\end{corollary}
\begin{proof}
  Each unitary element in \(M_n(\mathbb{C})\) has finite spectrum and therefore belongs to \(\mathcal{U}_0(M_n(\mathbb{C}))\) by Lemma~\ref{lemma:514}.
\end{proof}
\begin{lemma}[Whitehead]\label{lemma:516}
  Let \(A\) be a unital \cstar\ and let \(u, v\) be unitary elements in \(A\). Then
  \[
  \begin{pmatrix}
    u&0\\0&v
  \end{pmatrix}
  \thicksim_h
  \begin{pmatrix}
    uv&0\\0&1
  \end{pmatrix}
  \thicksim_h
  \begin{pmatrix}
    vu&0\\0&1
  \end{pmatrix}
  \thicksim_h
  \begin{pmatrix}
    v&0\\0&u
  \end{pmatrix}
  \qquad\mbox{ in }\mathcal{U}(M_2(A)).
  \]
  It follows in particular that
  \[
  \begin{pmatrix}
    u&0\\0&u^*
  \end{pmatrix}
  \thicksim_h
  \begin{pmatrix}
    1&0\\0&1
  \end{pmatrix}
  \qquad\mbox{ in }\mathcal{U}(M_2(A)).
  \]
\end{lemma}
\begin{proof}
  From Lemma~\ref{lemma:514} we get
  \[
  \begin{pmatrix}
    0&1\\1&0
  \end{pmatrix}
  \thicksim_h
  \begin{pmatrix}
    1&0\\0&1
  \end{pmatrix}.
  \]
  Hence
  \[
  \begin{pmatrix}
    u&0\\0&v
  \end{pmatrix}
  =
  \begin{pmatrix}
    u&0\\0&1
  \end{pmatrix}
  \begin{pmatrix}
    0&1\\1&0
  \end{pmatrix}
  \begin{pmatrix}
    v&0\\0&1
  \end{pmatrix}
  \begin{pmatrix}
    0&1\\1&0
  \end{pmatrix}
  \thicksim_h
  \begin{pmatrix}
    uv&0\\0&1
  \end{pmatrix}.
  \]
  The other claims follow in a similar way. See \cite[Lemma~2.1.5]{RordamLarsenLaustsen2000}.
\end{proof}
\begin{proposition}\label{prop:517}
  Let \(A\) be a unital \cstar.
  \begin{enumerate}
    \item\label{prop:5171} \(\mathcal{U}_0(A)\) is a normal subgroup of \(\mathcal{U}(A)\).
    \item\label{prop:5172} \(\mathcal{U}_0(A)\) is open and closed relative to \(\mathcal{U}(A)\).
    \item\label{prop:5173} An element \(u\) in \(A\) belongs to \(\mathcal{U}_0(A)\) if and only if
      \[u=e^{ih_1}\cdot e^{ih_2}\cdots e^{ih_n}\]
      for some natural number \(n\) and some self-adjoint elements \(h_1,\ldots,h_n\) in \(A\).
  \end{enumerate}
\end{proposition}
\begin{proof}
  To show~\ref{prop:5171} note first that \(\mathcal{U}_0(A)\) is closed under multiplication:
  if \(u_1,v_1,u_2,v_2\) are unitary elements in \(A\) with \(u_1\thicksim_h v_1\) and \(u_2\thicksim_h v_2\),
  then \(u_1 u_2\thicksim_h v_1 v_2\) using the multiplication of the continuous paths.
  We must also show that if \(u\) belongs to \(\mathcal{U}_0(A)\), then also do \(u^{-1}\) 
  and \(vuv^*\) for every \(v\) in \(\mathcal{U}(A)\). Let \(t\mapsto w_t\) be a continuous path in \(\mathcal{U}(A)\)
  from \(1\) to \(u\). Then \(t \mapsto w_t^{-1}\) and \(t\mapsto v w_t v^*\) are continuous paths in \(\mathcal{U}(A)\)
  from \(1\) to \(u^{-1}\) and from \(1\) to \(vu v^*\), respectively. This proves~\ref{prop:5171}\par
  Let \(G\) be the set of elements in \(A\) of the form
  \[
    e^{ih_1}\cdot e^{ih_2}\cdots e^{ih_n},
  \]
  where \(n\) is a natural number and \(h_1,\ldots,h_n\) are self-adjoint elements in \(A\).
  It follows from~\ref{prop:5171} that \(G\) is contained in \(\mathcal{U}_0(A)\).
  Since \({e^{ih}}^{-1}=e^{-ih}\) for every self-adjoint element \(h\), we see that \(G\)
  is a group.\par
  Take \(u\) in \(\mathcal{U}(A)\) and \(v\) in \(G\) with \(\|u-v\|<2\). Then \(\|1-uv^*\|=\|u-v\|<2\), 
  and by~\ref{lemma:514} and its proof, \(u v^*=e^{ih}\) for some self-adjoint element \(h\) in \(A\).
  Hence \(u=e^{ih}v\), and \(u\) therefore belongs to \(G\). This shows that \(G\) is open relative to \(\mathcal{U}(A)\).\par
  The complement \(\mathcal{U}(A)\setminus G\) is a union of right cosets of the form \(Gu\), with
  \(u\) in \(\mathcal{U}(A)\setminus G\). Right multiplication by \(u\) is a homeomorphism of
  \(\mathcal{U}(A)\), so each such coset is open because \(G\) is open.
  Consequently \(G\) is closed in \(\mathcal{U}(A)\).\par
  In conclusion \(G\) is a non-empty subset of \(\mathcal{U}_0(A)\), \(G\) is open and closed in \(\mathcal{U}(A)\),
  and \(\mathcal{U}_0(A)\) is connected. This entails that \(\mathcal{U}_0(A) = G\) and that~\ref{prop:5172} and~\ref{prop:5173} hold;
  compare \cite[Proposition~2.1.6]{RordamLarsenLaustsen2000}.
\end{proof}
\begin{lemma}\label{lemma:518}
  Let \(A\) and \(B\) be unital \cstars, and let \(\phi:A\rightarrow B\) be a surjective
  (and hence unit preserving) *-homomorphism.
  \begin{enumerate}
    \item\label{lemma:5181} \(\phi(\mathcal{U}_0(A))=\mathcal{U}_0(B)\).
    \item For each \(u\) in \(\mathcal{U}(B)\) there exists \(v\) in \(\mathcal{U}_0(M_2(A))\) such that
      \[\phi_2(v)=\begin{pmatrix}
        u&0\\0&u^*
      \end{pmatrix},\]
      where \(\phi_2: M_2(A) \rightarrow M_2(B)\) is the map induced by \(\phi\).
    \item If \(u\) is a unitary element in \(B\), and if there is a unitary element \(v\) in \(A\) such that 
      \(u\thicksim_h \phi(v)\), then \(u\) belongs to \(\phi(\mathcal{U}(A))\), i.e., \(u\) lifts to a unitary element in \(A\).
  \end{enumerate}
\end{lemma}
\begin{proof}
  \begin{enumerate}
    \item A unital *-homomorphism is continuous and maps unitary elements to unitary elements.
      Therefore \(\phi(\mathcal{U}_0(A))\) is contained in \(\mathcal{U}_0(B)\).\par
      Conversely, if \(u\) belongs to \(\mathcal{U}_0(B)\), then by Proposition~\ref{prop:517},
      \[
        u=e^{ih_1}\cdot e^{ih_2}\cdots e^{ih_n}
      \]
      for some self-adjoint elements \(h_j\) in \(B\). As \(\phi\) is surjective, there are elements \(x_j\) in \(A\)
      with \(\phi(x_j)=h_j\). Put \(k_j=(x_j+x^*_j)/2\). Then \(k_j=k_j^*\) and \(\phi(k_j)=h_j\). Set
      \[
        v=e^{ik_1}\cdot e^{ik_2}\cdots e^{ik_n}.
      \]
      Then \(v\) belongs to \(\mathcal{U}_0(A)\) by Proposition~\ref{prop:517}, and \(\phi(v)=u\).
    \item It follows from~\ref{lemma:5181} and Lemma~\ref{lemma:516}. 
    \item If \(u\thicksim_h \phi(v)\), then \(u\phi(v^*)\) belongs to \(\mathcal{U}_0(B)\), and hence
      \(u\phi(v^*)=\phi(w)\) for some \(w\) in \(\mathcal{U}_0(A)\) by~\ref{lemma:5181}. Now \(u=\phi(wv)\) as desired.
  \end{enumerate}
\end{proof}
\begin{definition}
  Let \(A\) be a unital \cstar. The group of invertible elements in \(A\) is denoted by \(GL(A)\),
  and the set of elements in \(GL(A)\) that are homotopic to \(1_A\) is denoted by \(GL_0(A)\).
  The group \(\mathcal{U}(A)\) of unitary elements of \(A\) is a subgroup of \(GL(A)\).\par
  For each element \(a\) in \(A\) set 
  \[|a|={(a^* a)}^{1/2}.\]
  The element \(|a|\) is called the \textnormal{absolute value} of \(a\).
\end{definition}
\begin{proposition}\label{prop:5111}
  Let \(A\) be a unital \cstar.
  \begin{enumerate}
    \item If \(z\) is an invertible element in \(A\), then so is \(|z|\), and \(\omega(z)=z|z|^{-1}\) belongs to \(\mathcal{U}(A)\).
      Clearly, \(z=\omega(z)|z|\).
    \item The map \(\omega:GL(A)\rightarrow\mathcal{U}(A)\) defined above is continuous, \(\omega(u)=u\) for every element in \(\mathcal{U}(A)\), 
      and \(\mathcal{U}(A)\) is a strong deformation retract of \(GL(A)\), with retraction \(\omega\).
    \item If \(u, v\) are unitary elements in \(\mathcal{U}(A)\), and if \(u\thicksim_h v\) in \(GL(A)\), then \(u\thicksim_h v\) in \(\mathcal{U}(A)\).
  \end{enumerate}
\end{proposition}
\begin{proof}
  \begin{enumerate}
    \item If \(z\) is invertible, then so are \(z^*\) and \(z^*z\). It follows that \(|z|={(z^*z)}^{1/2}\) is invertible
      with inverse \(|z|^{-1}=(z^*z)^{-1/2}\). Put \(u=z|z|^{-1}\) (\(=\omega(z)\)).
      Then \(z=u|z|\). Moreover, \(u\) is invertible, being a product of two invertible elements, and \(u^{-1}=u^*\) because
      \[u^*u=|z|^{-1}z^*z|z|^{-1}=|z|^{-1}|z|^2|z|^{-1}=1.\]
    \item Multiplication in a \cstar\ is continuous, and so is the map \(z\mapsto z^{-1}\) for \(z\) in \(GL(A)\).
      To show that \(\omega\) is continuous, it therefore suffices to show that the map \(a\mapsto |a|\) is continuous.
      This map is composed of the map \(a\mapsto a^*a\) and the map \(h\mapsto h^{1/2}\) for \(h\) in \(A^+\).
      The former of these maps is continuous because the involution and multiplication are continuous. It now suffices to show that \(h\mapsto h^{1/2}\)
      is continuous on any bounded subset \(\Omega\) of \(A^+\). But this follows from Lemma~\ref{lemma:125} 
      because each bounded subset \(\Omega\) of \(A\) is contained in \(\Omega_K\), where \(K=[0,R]\) and \(R=\sup\{\|h\||h\in\Omega\}\).\par 
      If \(u\) is a unitary element in \(A\), then \(|u|=1\), and so \(\omega(u)=u\).\par
      Let \(z\) in \(GL(A)\) be given, and put
      \[
        H(z,t)=\omega(z)((1-t)|z|+t\cdot 1_A),\qquad t\in[0,1].
      \]
      Since \(|z|\) is positive and invertible, there is \(\lambda>0\) such that \(|z|\geq\lambda\cdot 1_A\). Hence
      \[
        (1-t)|z|+t\cdot 1_A\geq ((1-t)\lambda+t)1_A>0,
      \]
      so \(H(z,t)\) is invertible for every \(t\). Moreover, \(H(z,0)=z\), \(H(z,1)=\omega(z)\), and
      \(H(u,t)=u\) for every \(u\) in \(\mathcal{U}(A)\). The continuity of \(H\) therefore shows that
      \(\mathcal{U}(A)\) is a strong deformation retract of \(GL(A)\); compare
      \cite[Section~2.1]{RordamLarsenLaustsen2000}.
    \item If \(t\mapsto z_t\) is a continuous path in \(GL(A)\) from \(u\) to \(v\), then \(t\mapsto \omega(z_t)\) is a continuous path in \(\mathcal{U}(A)\) from \(u\) to \(v\).
  \end{enumerate}
\end{proof}
\begin{definition}
  The factorization \(z=\omega(z)|z|\) from Proposition~\ref{prop:5111} of an invertible element \(z\) is called
  the \textnormal{(unitary) polar decomposition} for \(z\). We shall often write \(u\) instead of \(\omega(z)\), and the point
  of emphasis is then that each invertible element \(z\) in a unital \cstar\ \(A\) can be written as 
  \(z=u|z|\) for a (unique) unitary element \(u\) in \(A\).
\end{definition}
\begin{lemma}[Neumann series]
Let \(A\) be a unital Banach algebra. If
\[
\|1_A-a\|<1,
\]
then \(a\) is invertible and
\begin{equation}\label{CN1}
a^{-1}
=
\sum_{n=0}^{\infty}{(1_A-a)}^n.
\end{equation}
It follows in particular that the norm of \(a^{-1}\) can be estimated by
\begin{equation}\label{CN2}
  \|a^{-1}\|\leq 1+\|1_A -a\|+\|1_A-a\|^2+\cdots=(1-\|1_A-a\|)^{-1}.
\end{equation}
\end{lemma}

\begin{proof}
Since
\[
\sum_{n=0}^{\infty}\|{(1_A-a)}^n\|
\le
\sum_{n=0}^{\infty}\|1_A-a\|^n
<\infty,
\]
the series converges in norm. If
\[
s=\sum_{n=0}^{\infty}{(1_A-a)}^n,
\]
then
\[
as=(1_A-(1_A-a))s
=s-(1_A-a)s
=1_A.
\]
Similarly,
\[
sa=1_A.
\]
Hence \(s=a^{-1}\).
\end{proof}
\begin{proposition}\label{prop:invPerturbation}
  Let \(A\) be a unital \cstar. Let \(a\) be an invertible element in \(A\), 
  and take \(b\) in \(A\) with \(\|a-b\|<\|a^{-1}\|^{-1}\). Then \(b\) is invertible,
  \[\|b^{-1}\|^{-1}\geq \|a^{-1}\|^{-1}-\|a-b\|,\]
  and \(a\thicksim_h b\) in \(GL(A)\).
\end{proposition}
\begin{proof}
  From 
  \[
    \|1_A-a^{-1}b\|=\|a^{-1}(a-b)\|\leq\|a^{-1}\|\cdot\|a-b\|<1,
  \]
  we see that \(a^{-1}b\) is invertible, and by (\ref{CN2}),
  \[
    \|(a^{-1}b)^{-1}\|\leq{(1-\|1_A-a^{-1}b\|)}^{-1}.
  \]
  This implies that \(b\) is invertible with inverse \(b^{-1}=(a^{-1}b)^{-1}a^{-1}\), and
  \begin{align*}
    \|b^{-1}\|^{-1}&\geq \|(a^{-1}b)^{-1}\|^{-1}\cdot\|a^{-1}\|^{-1}\geq(1-\|1_A-a^{-1}b\|)\cdot\|a^{-1}\|^{-1}\\
    &\geq\|a^{-1}\|^{-1}-\|a-b\|.
  \end{align*}
  To prove the last claim, put \(c_t=(1-t)a+tb\) for \(t\) in \([0,1]\). Then \(\|a-c_t\|=t\|a-b\|<\|a^{-1}\|^{-1}\),
  and therefore \(c_t\) is invertible by the first part of the proof. This shows that \(a=c_0\thicksim_h c_1 = b\) in \(GL(A)\).
\end{proof}

\section{Equivalence of projections}
\begin{definition}
  Recall that an element \(p\) in a \cstar\ is a projection if \(p=p^2=p^*\).
  The set of all projections in a \cstar\ \(A\) is denoted by \(\mathcal{P}(A)\).
  Consider the following two equivalence relations on \(\mathcal{P}(A)\):
  \begin{enumerate}
    \item \(p\thicksim q\) if there exists \(v\) in \(A\) with \(p=v^* v\) and \(q=v v^*\) (\textnormal{Murray-von Neumann equivalence}),
    \item \(p\thicksim_u q\) if there exists a unitary element \(u\) in \(\mathcal{U}(\widetilde{A} )\) with \(q=u p u^*\) (\textnormal{unitary equivalence}).
  \end{enumerate}
  An element \(v\) in \(A\) for which \(v^* v\) is a projection is called a \textnormal{partial isometry}.
  As shown in Exercise~\ref{ex:25}, if \(v\) is a partial isometry, then \(v v^*\) is also a projection.
  The projection \(v^* v\) is called the \textnormal{support projection} of \(v\), and \(v v^*\) is called the \textnormal{range projection} of \(v\).
\end{definition}
Put \(p=v^* v\) and \(q=v v^*\). Then we have the following much used identities:
\begin{equation}\label{eq:23}
  v=qv=vp=qvp;
\end{equation}
see Exercise~\ref{ex:25}.\par
We can use this to show that the Murray-von Neumann relation is transitive.
Let \(p, q\) and \(r\) be equivalent projections in \(A\), and take partial isometries \(v\) and \(w\) such that
\(p=v^* v, q=v v^*=w^* w\), and \(r=w w^*\). Put \(z=wv\). Then
\[
  z^* z=v^* w^* wv=v^* qv=v^* v=p,\qquad z z^*=w v v^* w^*=w w^*=r,
\]
so that \(p\thicksim r\).
\begin{proposition}\label{prop:522}
  Let \(p, q\) be projections in a unital \cstar\ \(A\). The following conditions are equivalent:
  \begin{enumerate}
    \item\label{prop:2221} \(p\thicksim_u q\),
    \item\label{prop:2222} \(q=u p u^*\) for some unitary \(u\) in \(A\),
    \item\label{prop:2223} \(p\thicksim q\) and \(1_A-p\thicksim 1_A -q\). 
  \end{enumerate}
\end{proposition}
\begin{proof}
  Let \(1_{\widetilde{A}}\) denote the unit of \(\widetilde{A}\), and put \(f=1_{\widetilde{A}}-1_A\).
  Then \(\widetilde{A}=A\oplus\mathbb{C}f\), and \(fa=af=0\) for all \(a\) in \(A\).\par
  (\ref{prop:2221}) \(\Rightarrow\) (\ref{prop:2222}). Suppose that \(q=z p z^*\) for some \(z\) in \(\mathcal{U}(\widetilde{A})\). We have 
    \(z=u+\alpha f\) for some \(u\) in \(A\) and some \(\alpha\) in \(\mathbb{C}\). A straightforward calculation shows that \(u\) is unitary and that \(q=up u^*\).\par
  (\ref{prop:2222}) \(\Rightarrow\) (\ref{prop:2223}). Suppose that \(q=up u^*\) for some unitary \(u\) in \(\mathcal{U}(A)\). Put \(v=up\) and \(w=u(1_A-p)\). Then
  \begin{equation}\label{eq:24}
    v^* v=p,\qquad v v^*=q,\qquad w^* w=1_A-p,\qquad w w^*=1_A-q.
  \end{equation}\par
  (\ref{prop:2223}) \(\Rightarrow\) (\ref{prop:2221}). Suppose that there are partial isometries \(v\) and \(w\) satisfying (\ref{eq:24}) above. A direct calculation (using equation (\ref{eq:23})) 
  shows that \(z=v+w+f\) is a unitary element in \(\widetilde{A}\). By equation (\ref{eq:23}) we find that \(zp z^*=vp v^*=q\).
\end{proof}
\begin{lemma}\label{lemma:523}
  Let \(p\) be a projection in a \cstar\ \(A\), and let \(a\) be a self-adjoint element in \(A\). Put \(\delta=\|p-a\|\). Then
  \[
    \operatorname{spec}(a)\subseteq[-\delta,\delta]\cup [1-\delta,1+\delta].
  \]
\end{lemma}
\begin{proof}
  Since \(a\) is self-adjoint its spectrum consists of real numbers. Recall that \(\operatorname{spec}(p)\) is contained in \(\{0,1\}\).
  It suffices to show that if \(t\) is a real number whose distance, \(d\), to the set \(\{0,1\}\) is strictly greater than \(\delta\),
  then \(t\) does not belong to \(\operatorname{spec}(a)\). For such a real number \(t\) the element \(p-t\cdot 1\) is invertible in \(\widetilde{A}\) and
  \[
    \|{(p-t\cdot 1)}^{-1}\|=\max\{|\lambda-t|^{-1}\mid\lambda\in\operatorname{spec}(p)\}\leq d^{-1}.
  \]
  It follows that
  \[
    \|{(p-t\cdot 1)}^{-1}(a-t\cdot 1)-1\|=\|{(p-t\cdot 1)}^{-1}(a-p)\|\leq d^{-1}\delta<1.
  \]
  This shows that \({(p-t\cdot 1)}^{-1}(a-t\cdot 1)\) is invertible. Hence \(a-t\cdot 1\) is invertible, and \(t\) does not belong to \(\operatorname{spec}(a)\).
\end{proof}
\begin{proposition}\label{prop:Lars224}
  If \(p, q\) are projections in a \cstar\ \(A\) and \(\|p-q\|<1\), then \(p\thicksim_h q\).
\end{proposition}
\begin{proof}
  Put \(a_t=(1-t)p+tq\) for \(t\) in \([0,1]\). Then \(a_t\) is self-adjoint,
  \[
    \min\{\|a_t-p\|,\|a_t-q\|\}\leq\|p-q\|/2<1/2,
  \]
  and \(t\mapsto a_t\) is continuous. Moreover, in the notation of Lemma~\ref{lemma:125}, 
  each \(a_t\) belongs to \(\Omega_K\) by Lemma~\ref{lemma:523}, 
  where \(K=[-\delta,\delta]\cup[1-\delta,1+\delta]\) and \(\delta=\|p-q\|/2<1/2\).\par
  Define \(f:K\rightarrow\mathbb{C}\) to be the continuous function which is zero on the 
  interval \([-\delta,\delta]\) and \(1\) on the interval \([1-\delta,1+\delta]\), which are disjoint.
  Then \(f(a_t)\) is a projection for each \(t\) in \([0,1]\) because \(f=f^2=\overline{f}\). 
  The path \(t\mapsto f(a_t)\) is continuous by Lemma~\ref{lemma:125}, and hence
  \[
    p=f(p)=f(a_0)\thicksim_h f(a_1)=f(q)=q\mbox{ in }\mathcal{P}(A).
  \]
\end{proof}
\begin{proposition}\label{prop:525}
  Let \(a, b\) be self-adjoint elements in a unital \cstar\ \(A\), and suppose that \(b=za z^{-1}\) for some invertible
  element \(z\) in \(A\). Let \(z=u|z|\) be the polar decomposition for \(z\) with \(u\) in \(\mathcal{U}(A)\).
  Then \(b=uau^*\).\par
  In other words, if two self-adjoint elements in a unital \cstar\ are similar, then they are unitarily equivalent. 
\end{proposition}
\begin{proof}
  The equation \(b=zaz^{-1}\) implies that \(bz=za\) and, because \(a\) and \(b\) are self-adjoint,
  also \(z^*b=az^*\). Hence
  \[
    |z|^2a=(z^*z)a=z^*bz=a(z^*z)=a|z|^2
  \]
  and so \(a\) commutes with \(|z|^2\). Consequently \(a\) commutes with all elements in \(C^*(1,|z|^2)\)
  and in particular \(a\) commutes with \(|z|^{-1}\). It follows that
  \[
    uau^*=z|z|^{-1}au^*=za|z|^{-1}u^*=bz|z|^{-1}u^*=buu^*=b,
  \]
  as desired.
\end{proof}
\begin{proposition}\label{prop:526}
  Let \(A\) be a \cstar\, and let \(p, q\) be projections in \(\mathcal{P}(A)\).
  Then \(p\thicksim_h q\) in \(\mathcal{P}(A)\) if and only if there exists a unitary element \(u\) in \(\mathcal{U}_0(\widetilde{A})\) such that \(q=upu^*\).
\end{proposition}
\begin{proof}
  Throughout the proof, \(1\) denotes the unit of \(\widetilde{A}\).\par
  Suppose first that \(q=upu^*\) for some \(u\) in \(\mathcal{U}_0(\widetilde{A})\).
  Let \(t\mapsto u_t\) be a continuous path of unitaries in \(\widetilde{A}\) from \(1\) to \(u\).
  Then, because \(A\) is an ideal in \(\widetilde{A}\), \(t\mapsto u_t p u_t^*\) is a continuous path of projections in \(A\) from \(p\) to \(q\).\par
  Conversely, if \(p\thicksim_h q\), then there are projections \(p=p_0,p_1,p_2,\ldots,p_n=q\) in \(A\) such that \(\|p_{j+1}-p_j\|<1/2\) for each \(j\).
  It therefore suffices to prove the implication in the case where \(\|p-q\|<1/2\).\par
  Put \(z=pq+(1-p)(1-q)\). Then \(z\) belongs to \(\widetilde{A}\), \(pz=pq=zq\), and
  \[
    \|z-1\|=\|p(q-p)+(1-p)((1-q)-(1-p))\|\leq 2\|q-p\|<1.
  \]
  Hence \(z\) is invertible and \(z\thicksim_h 1\) in \(GL(\widetilde{A})\) by Proposition~\ref{prop:invPerturbation}.
  Let \(z=u|z|\) be the unitary polar decomposition for \(z\). Then \(p=uqu^*\) by Proposition~\ref{prop:525}.
  From the properties of the polar decomposition (Proposition~\ref{prop:5111}) we have \(u\thicksim_h z\thicksim_h 1\)
  in \(GL(\widetilde{A})\), and this entails that \(u\) belongs to \(\mathcal{U}_0(\widetilde{A})\) by Proposition~\ref{prop:5111}.
  Consequently \(u^*\) belongs to \(\mathcal{U}_0(\widetilde{A})\) and \(q=u^*pu\), as required; compare
  \cite[Proposition~4.6.6 and Theorem~4.6.7]{Blackadar1998}.
\end{proof}
We have three equivalence relations \(\thicksim_h, \thicksim_u\) and \(\thicksim\) on the set of projections in a \cstar.
It will be shown in Example~\ref{example:529} that these three equivalence relations are different from each other.
Proposition~\ref{prop:527} below says that homotopy is stronger than unitary equivalence which again is stronger than the Murray-von Neumann
equivalence; Proposition~\ref{prop:528} says that the three relations actually are equal modulo passing to matrix algebras.
\begin{proposition}\label{prop:527}
  Let \(p, q\) be projections in a \cstar\ \(A\).
  \begin{enumerate}
    \item\label{prop:5271} If \(p\thicksim_h q\), then \(p\thicksim_u q\).
    \item\label{prop:5272} If \(p\thicksim_u q\), then \(p\thicksim q\).
  \end{enumerate}
\end{proposition}
\begin{proof}
  Part (\ref{prop:5271}) is an immediate consequence of Proposition~\ref{prop:526}.\par
  If \(upu^*=q\) for some unitary \(u\) in \(\widetilde{A}\), then \(v=up\) belongs to \(A\), \(v^*v=p\), and \(vv^*=q\). This proves (\ref{prop:5272}).
\end{proof}

\begin{proposition}\label{prop:528}\label{matthes:229}
  Let \(p, q\) be projections in a \cstar\ \(A\).
  \begin{enumerate}
    \item If \(p\thicksim q\), then \(\begin{pmatrix}
      p&0\\0&0
    \end{pmatrix}\thicksim_u
    \begin{pmatrix}
      q&0\\0&0
    \end{pmatrix}\) in \(M_2(A)\).
    \item If \(p\thicksim_u q\), then \(\begin{pmatrix}
      p&0\\0&0
    \end{pmatrix}\thicksim_h
    \begin{pmatrix}
      q&0\\0&0
    \end{pmatrix}\) in \(M_2(A)\).
  \end{enumerate}
\end{proposition}
\begin{proof}
  \begin{enumerate}
    \item Find \(v\) in \(A\) such that \(p=v^*v\) and \(q=vv^*\). Use (\ref{eq:23}) to see that
    \[
      u=\begin{pmatrix}
        v&1-q\\1-p&v^*
      \end{pmatrix},\qquad
      w=\begin{pmatrix}
        q&1-q\\1-q&q
      \end{pmatrix}
    \]
    are unitary elements in \(M_2(\widetilde{A})\). Since
    \[
      wu\begin{pmatrix}
        p&0\\0&0
      \end{pmatrix}
      u^* w^*=w\begin{pmatrix}
        q&0\\0&0
      \end{pmatrix}
      w^*=\begin{pmatrix}
        q&0\\0&0
      \end{pmatrix},
    \]
    and since 
    \[
      wu=\begin{pmatrix}
        v+(1-q)(1-p)&(1-q)v^*\\q(1-p)&(1-q)+qv^*
      \end{pmatrix}
    \]
    belongs to \(\widetilde{M_2(A)}\), claim is proved. (Notice we must verify that the unitary element
    \(wu\), which a priori belongs to \(M_2(\widetilde{A})\), in fact belongs to the sub-\cstar\ \(\widetilde{M_2(A)}\).)
    \item Find \(u\) in \(\mathcal{U}(\widetilde{A})\) such that \(q=up u^*\). Lemma~\ref{lemma:516} provides us with a continuous path
    \(t\mapsto w_t\) in \(\mathcal{U}(M_2(\widetilde{A}))\) such that
    \[
      w_0=\begin{pmatrix}
        1&0\\0&1
      \end{pmatrix},\qquad
      w_1=\begin{pmatrix}
        u&0\\0&u^*
      \end{pmatrix}.
    \]
    Put \(e_t=w_t \text{diag}(p,0)w_t^*\). Then each \(e_t\) belongs to \(\mathcal{P}(M_2(A))\), the map \(t\mapsto e_t\) is continuous,
    \(e_0=\text{diag}(p,0)\), and \(e_1=\text{diag}(q,0)\).
  \end{enumerate}
\end{proof}
\begin{example}\label{example:529}
  Neither of the implications ``\(p\thicksim q\Rightarrow p\thicksim_u q\)'' and 
  ``\(p\thicksim_u q\Rightarrow p\thicksim_h q\)'' holds in general.\par
  An \textnormal{isometry} in a unital \cstar\ \(A\) is an element \(s\) such that \(s^*s=1\).
  Non-unitary isometries exist. Take for example the \textnormal{unilateral shift} \(S\) on the Hilbert space
  \(\ell^2(\mathbb{N})\) given by
  \[
    (S\xi)(n)=\begin{cases}
      0,& n=1\\
      \xi(n-1),& n\geq 2,
    \end{cases}
  \]
  for all \(\xi\) in \(\ell^2(\mathbb{N})\). The adjoint of \(S\) is given by \((S^*\xi)(n)=\xi(n+1)\).
  It follows that \(S^*S=I\) and hence that \(\|S\xi\|=\|\xi\|\) for all \(\xi\) in \(\ell^2(\mathbb{N})\), i.e., \(S\)
  is an isometry. Clearly, \(S\) is not surjective and therefore not unitary.\par
  Let \(A\) be a unital \cstar\ containing a non-unitary isometry \(s\). By definition, \(s^*s\thicksim ss^*\). The zero projection is not Murray-von Neumann equivalent
  to a non-zero projection, because if \(v^*v=0\), then \(v=0\) and so \(vv^*=0\).
  Hence \(1-s^*s(=0)\) is not equivalent to \(1-ss^*(\neq 0)\). By Proposition~\ref{prop:522} 
  we conclude that the projections \(s^*s\) and \(ss^*\) are not unitarily equivalent.\par
  An example of a (unital) \cstar\ \(A\) that contains projections \(p, q\) such that \(p\thicksim_u q\) and \(p\nsim_h q\) is more complicated to construct.
\end{example}

\section{Semigroups of projections}
\begin{definition}[The semigroup \(\mathcal{P}_\infty(A)\)]\label{def:PinftyA}
  Put 
  \[
    \mathcal{P}_n(A)=\mathcal{P}(M_n(A)),\qquad\mathcal{P}_\infty(A)=\bigcup_{n=1}^{\infty}\mathcal{P}_n(A),
  \]
  where \(A\) is a \cstar\ and \(n\) is a positive integer. We view the sets \(\mathcal{P}_n(A)\), \(n\in\mathbb{N}\) as being pairwise disjoint.\par
  Define the relation \(\thicksim_0\) on \(\mathcal{P}_\infty(A)\) as follows. Suppose that \(p\) is a projection in \(\mathcal{P}_n(A)\) and \(q\) in \(\mathcal{P}_m(A)\).
  Then \(p\thicksim_0 q\) if there is an element \(v\) in \(M_{m,n}(A)\) with \(p=v^*v\) and \(q=vv^*\). Here, \(M_{m,n}(A)\) is the set of all rectangular \(m\times n\) matrices
  \[
    \begin{pmatrix}
      a_{11}&a_{12}&\cdots&a_{1n}\\
      a_{21}&a_{22}&\cdots&a_{2n}\\
      \vdots&\vdots&\ddots&\vdots\\
      a_{m1}&a_{m2}&\cdots&a_{mn}
    \end{pmatrix}
  \]
  with entries \(a_{ij}\) in \(A\). The adjoint \(v^*\) of an element \(v\) in \(M_{m,n}(A)\) is the element in 
  \(M_{n,m}(A)\) obtained transposing and taking adjoints entry-wise.
  The products \(vv^*\) and \(v^*v\) are the usual matrix products.\par
  Define a binary operation \(\oplus\) on \(\mathcal{P}_\infty(A)\) by
  \[
    p\oplus q=\text{diag}(p,q)=\begin{pmatrix}
      p&0\\0&q
    \end{pmatrix}.
  \]
  so that \(p\oplus q\) belongs to \(\mathcal{P}_{m+n}(A)\) when \(p\) is in \(\mathcal{P}_n(A)\) and \(q\) is in \(\mathcal{P}_m(A)\).\\
  The relation \(\thicksim_0\) combines the Murray-von Neumann equivalence relation \(\thicksim\) with an identification
  of projections in different sized matrix algebras over \(A\). If \(p\) and \(q\) both belong to \(\mathcal{P}_n(A)\) for some \(n\), then
  \(p\thicksim_0 q\) if and only if \(p\) and \(q\) are Murray-von Neumann equivalent.\par
  The relation \(\thicksim_0\) is named after the \(K_0\)-group it will help define.
\end{definition}
\begin{lemma}
  The relation \(\thicksim_0\) is an equivalence relation on \(\mathcal{P}_\infty(A)\).
\end{lemma}
\begin{proof}
  If \(p\) belongs to \(\mathcal{P}_n(A)\), then \(p=p^*p=pp^*\), and hence \(p\thicksim_0p\).
  If \(p=v^*v\) and \(q=vv^*\) for some \(v\) in \(M_{m,n}(A)\), then
  \(q={(v^*)}^*v^*\) and \(p=v^*{(v^*)}^*\), so \(q\thicksim_0p\).
  Finally, suppose that
  \[
    p=v^*v,\qquad q=vv^*=w^*w,\qquad r=ww^*
  \]
  for suitable rectangular matrices \(v\) and \(w\). Put \(z=wv\). Since \(qv=v\) and \(wq=w\),
  \[
    z^*z=v^*w^*wv=v^*qv=p,
    \qquad
    zz^*=wvv^*w^*=wqw^*=r.
  \]
  Thus \(p\thicksim_0r\), proving transitivity. See also
  \cite[Section~2.3]{RordamLarsenLaustsen2000}.
\end{proof}
\begin{proposition}\label{prop:232}
  Let \(p,q,r,p',q'\) be projections in \(\mathcal{P}_\infty(A)\) for some \cstar\ \(A\).
  \begin{enumerate}
    \item \(p\thicksim_0 p\oplus 0_n\) for every natural number \(n\), where \(0_n\) is the zero element of \(M_n(A)\),
    \item if \(p\thicksim_0 p'\) and \(q\thicksim_0 q'\), then \(p\oplus q\thicksim_0 p'\oplus q'\),
    \item \(p\oplus q\thicksim_0 q\oplus p\),
    \item if \(p, q\) are projections in \(\mathcal{P}_n(A)\) such that \(pq=0\), then \(p+q\) is a projection and \(p+q\thicksim_0 p\oplus q\),
    \item \((p\oplus q)\oplus r=p\oplus(q\oplus r)\).
  \end{enumerate}
\end{proposition}
\begin{proof}
  \begin{enumerate}
    \item Let \(m, n\) be positive integers, and let \(p\) be a projection in \(\mathcal{P}_m(A)\).
      Put
      \[
        u_1=\begin{pmatrix}
          p\\0
        \end{pmatrix}\in M_{m+n,m}(A).
      \]
      Then \(p=u^*_1 u_1\thicksim_0 u_1 u_1^*=p\oplus 0_n\).
    \item If \(p\thicksim_0 p'\) and \(q\thicksim_0 q'\), then there exist \(v, w\) such that 
    \[
      p=v^*v,\qquad p'=vv^*,\qquad q=w^*w,\qquad q'=ww^*.
    \]
    Put \(u_2 =\text{diag}(v, w)\). Then \(p\oplus q=u_2^*u_2\thicksim_0 u_2u_2^*=p'\oplus q'\).
    \item Suppose that \(p\) belongs to \(\mathcal{P}_n(A)\) and \(q\) belongs to \(\mathcal{P}_m(A)\), and put
    \[
      u_3=\begin{pmatrix}
        0_{m,n}&q\\p&0_{n,m}
      \end{pmatrix}
    \]
    where \(0_{k,l}\) is the zero element in \(M_{k,l}(A)\). Then \(u_3\) belongs to \(M_{n+m}(A)\),
    and \(p\oplus q=u^*_3u_3\thicksim_0 u_3u^*_3=q\oplus p\).
    \item If \(n=m\) and \(pq=0\), then \(p+q\) is a projection (see Exercise~\ref{ex:24}).
    Put 
    \[
      u_4=\begin{pmatrix}
        p\\q
      \end{pmatrix}\in M_{2n,n}(A).
    \]
    Then \(p+q=u_4^*u_4\thicksim_0 u_4u^*_4=p\oplus q\).
    \item Proof is trivial.
  \end{enumerate}
\end{proof}
\begin{definition}[The semigroup \(\mathcal{D}(A)\)]
  With \((\mathcal{P}_\infty(A), \thicksim_0, \oplus)\) as in Definition~\ref{def:PinftyA}, set
  \[
    \mathcal{D}(A) = \mathcal{P}_\infty(A)/\thicksim_0.
  \]
  For each \(p\) in \(\mathcal{P}_\infty(A)\) let \({[p]}_\mathcal{D}\) in \(\mathcal{D}(A)\) denote the 
  equivalence class containing \(p\). Define addition on \(\mathcal{D}(A)\) by
  \[
    {[p]}_\mathcal{D}+{[q]}_\mathcal{D}={[p\oplus q]}_\mathcal{D},\qquad p,q\in\mathcal{P}_\infty (A).
  \]
  It follows from Proposition~\ref{prop:232} that this operation is well-defined and that \((\mathcal{D}(A),+)\) is an Abelian semigroup.
\end{definition}
In the next chapter we shall, for each unital \cstar\ \(A\), construct an Abelian group \(K_0(A)\) from the semigroup \((\mathcal{D}(A),+)\).

\section{Exercises}
\begin{exercise}\label{ex:24}
  Two projections \(p, q\) in a unital \cstar\ \(A\) are said to be \textnormal{orthogonal to each other},
  or mutually orthogonal (in symbols \(p\perp q\)), if \(pq=0\).\\
  Show that the following three conditions are equivalent:
  \begin{enumerate}
    \item \(p\perp q\),
    \item \(p+q\) is a projection,
    \item \(p+q\leq 1\).
  \end{enumerate}
  A family \(p_1, p_2,\ldots,p_n\) of projections in \(A\) is said to be mutually orthogonal if 
  \(p_i\perp p_j\) whenever \(i\neq j\). Use the first part to prove that the following are equivalent:
  \begin{enumerate}
    \item \(p_1, p_2,\ldots,p_n\) are mutually orthogonal,
    \item \(p_1+p_2+\cdots +p_n\) is a projection,
    \item \(p_1+p_2+\cdots+p_n\leq 1\).
  \end{enumerate}
\end{exercise}
\begin{proof}
  [Hint: To see that (3) implies (1), observe that (3) implies \(p(p+q)p\leq p\).]
\end{proof}
\begin{exercise}\label{ex:25}
  Let \(A\) be a \cstar, and let \(v\) in \(A\) be such that \(v^* v\) is a projection (i.e., \(v\) is a partial isometry).
  Show that \(v=v v^*v\), and conclude that \(v v^*\) is a projection.\par
  Put \(p=v^* v\) and \(q=v v^*\). Show that \(v=qv=vp=qvp\). 
\end{exercise}
\begin{exercise}\label{ex:Lars29}
  Let \(\operatorname{Tr}:M_n(\mathbb{C})\to\mathbb{C}\) be the standard trace given by
  \[
    \operatorname{Tr}\begin{pmatrix}
      \alpha_{11}&\alpha_{12}&\cdots&\alpha_{1n}\\
      \alpha_{21}&\alpha_{22}&\cdots&\alpha_{2n}\\
      \vdots&\vdots&\ddots&\vdots\\
      \alpha_{n1}&\alpha_{n2}&\cdots&\alpha_{nn}
    \end{pmatrix}
    =\sum_{j=1}^n\alpha_{jj}.
  \]
  Let \(p,q\) be projections in \(M_n(\mathbb{C})\). Show that the following are equivalent:
  \begin{enumerate}
    \item \(p\sim q\),
    \item \(\operatorname{Tr}(p)=\operatorname{Tr}(q)\),
    \item \(\dim(p(\mathbb{C}^n))=\dim(q(\mathbb{C}^n))\).
  \end{enumerate}
  Use this to show that 
  \[
    \mathcal{D}(\mathbb{C})\cong\{0,1,2,\ldots\}=\mathbb{Z}^+,
  \]
  when \(\mathbb{Z}^+\) is equipped with the usual addition.\par
  Show finally that the implications 
  \[
    ''p\sim q\Rightarrow p\sim_u q''\text{ and }''p\sim q\Rightarrow p\sim_h q''
  \]
  hold for all projections \(p,q\) in \(M_n(\mathbb{C})\). 
\end{exercise}










\chapter{The \texorpdfstring{\(K_0\)}{K0}-Group of a Unital \texorpdfstring{\(C^*\)}{C*}-Algebra}
An abelian group \(K_0(A)\) is associated to each unital \cstar\ \(A\).
The group \(K_0(A)\) arises from the abelian semigroup \((\mathcal{D}(A),+)\) and the Grothendieck construction.
We shall see that \(K_0\) is a functor from the category of unital \cstars\ to the category of abelian groups,
and some properties of \(K_0\) will be derived. Some examples of \(K_0\)-groups can be found at the end of the chapter.\par
We extend \(K_0\) to a functor from the category of all \cstars\ (unital or
not) in Chapter \ref{chap:K0Functor}.

\section{Definition of the \texorpdfstring{\(K_0\)}{K0}-group of a unital \texorpdfstring{\(C^*\)}{C*}-algebra}
One can associate an abelian group to every abelian semigroup in a way analogous to how one obtains
the integers from the natural numbers, and in much the same way as one obtains the rational numbers from the integers.
We describe here how this works.

\begin{definition}[The Grothendieck construction]
  Let \((S,+)\) be an abelian semigroup. Define an equivalence relation \(\sim\) on \(S\times S\) by \((x_1,y_1)\sim (x_2,y_2)\) 
  if there exists \(z\) in \(S\) such that \(x_1+y_2+z=x_2+y_1+z\). We leave the proof that \(\sim\) is an equivalence relation for later.\par
  Write \(G(S)\) for the quotient \((S\times S)/\sim\), and let \(\langle x,y\rangle\) denote the equivalence class in \(G(S)\) containing \((x,y)\) in \(S\times S\). The operation
  \[
    \langle x_1,y_1\rangle+\langle x_2,y_2\rangle=\langle x_1+x_2,y_1+y_2\rangle
  \]
  is well-defined and turns \((G(S),+)\) into an abelian group.
  Notice that \(-\langle x,y\rangle=\langle y,x\rangle\) and \(\langle x,x\rangle=0\) for all \(x\) in \(S\).
  The group \(G(S)\) is called the \textnormal{Grothendieck group} of \(S\).\par
  Take \(y\) in \(S\). The map
  \[
    \gamma_S:S\rightarrow G(S),\qquad x\mapsto\langle x+y,y\rangle,
  \]
  is independent of the choice of \(y\), and \(\gamma_S\) is additive. It is called the 
  \textnormal{Grothendieck map}.\par
  The semigroup \((S,+)\) is said to have the \textnormal{cancellation property} if, whenever \(x, y\) and \(z\) are elements in \(S\)
  with \(x+z=y+z\), it follows that \(x=y\).
\end{definition}
\begin{proposition}
  The relation \(\sim\) is an equivalence relation on \(S\times S\).
\end{proposition}
\begin{proof}
  It is clear that \(\sim\) is reflexive and symmetric. To see that it is transitive, suppose that \((x_1,y_1)\sim(x_2,y_2)\) and \((x_2,y_2)\sim(x_3,y_3)\). Then
  \[
    x_1+y_2+z=x_2+y_1+z,\qquad x_2+y_3+w=x_3+y_2+w
  \]
  for some \(z, w\) in \(S\). If we sum the two equations, using the commutativity of the abelian \(S\), we have
  \begin{align*}
    x_1+y_3+(y_2+z+x_2+w)&=x_3+y_1+(x_2+y_2+w+z)
  \end{align*}
  Hence we obtain \((x_1,y_1)\sim (x_3,y_3)\) using the element \(x_2+y_2+w+z\) in \(S\).
\end{proof}
\begin{proposition}\label{prop:613}
  The Grothendieck construction has the following properties.
  \begin{enumerate}
    \item \textnormal{The universal property}. If \(H\) is an abelian group, and if \(\varphi:S\rightarrow H\) is
    an additive map, then there is one and only one group homomorphism \(\psi:G(S)\rightarrow H\) making the diagram
    \begin{equation}\label{UPdiag}
    \begin{tikzcd}[column sep=scriptsize]
    S \arrow{dr}[swap]{\gamma_S} \arrow{rr}{\varphi}
    & & H \\
    & G(S) \arrow{ur}[swap]{\psi}
    \end{tikzcd}
  \end{equation}
    commutative.
    \item \textnormal{Functoriality}. To every additive map \(\varphi:S\rightarrow T\) between
    semigroups \(S\) and \(T\) there is precisely one group homomorphism \(G(\varphi):G(S)\rightarrow G(T)\)
    making the diagram
    \[
    \begin{tikzcd}
    S \arrow[r, "\varphi"] \arrow[d, swap, "\gamma_S"]
    & T \arrow[d, "\gamma_T"] \\
    G(S) \arrow[r, swap, "G(\varphi)"]
    & G(T)
    \end{tikzcd}
    \] 
    commutative.
    \item \(G(S)=\{\gamma_S(x)-\gamma_S(y)\mid x,y\in S\}\).
    \item Let \(x, y\) be elements in \(S\). Then \(\gamma_S(x)=\gamma_S(y)\) if and only if
    \(x+z=y+z\) for some \(z\) in \(S\).
    \item The Grothendieck map \(\gamma_S:S\rightarrow G(S)\) is injective if and only if \(S\) has the cancellation property.
    \item Let \((H,+)\) be an abelian group, and let \(S\) be a non-empty subset of \(H\).
    If \(S\) is closed under addition, then \((S,+)\) is an abelian semigroup with the cancellation property, \(G(S)\) is isomorphic to the
    subgroup \(H_0\) generated by \(S\), and \(H_0=\{x-y\mid x,y\in S\}\).
  \end{enumerate}
  \begin{proof}
    \begin{enumerate}[label=\arabic*.]
      \item[3.] Each element in \(G(S)\) has the form \(\langle x,y\rangle\) for some \(x, y\) in \(S\), and
      \[
      \langle x,y\rangle=\langle x+y,y\rangle-\langle x+y,x\rangle=\gamma_S(x)-\gamma_S(y).
      \]
      \item[4.] To see the ``if'' part, use that \(\gamma_S\) is additive and that \(G(S)\) is a group.
      Suppose conversely that \(\gamma_S(x)=\gamma_S(y)\). Then \(\langle x+y,y\rangle =\langle y+x,x\rangle\),
      and therefore \((x+y)+x+w=(x+y)+y+w\) for some \(w\) in \(S\). This shows that \(x+z=y+z\), where \(z=x+y+w\).
      \item[5.] It is an immediate consequence of (4).
      \item[1.] If \(\langle x_1,y_1\rangle=\langle x_2,y_2\rangle\), then \(x_1+y_2+z=x_2+y_1+z\) for some \(z\) in \(S\), and hence
      \[
        \varphi(x_1)+\varphi(y_2)+\varphi(z)=\varphi(x_2)+\varphi(y_1)+\varphi(z)
      \]
      in \(H\), which implies that \(\varphi(x_1)-\varphi(y_1)=\varphi(x_2)-\varphi(y_2)\). Therefore the map
      \(\psi:G(S)\rightarrow H\) given by \(\psi(\langle x,y\rangle)=\varphi(x)-\varphi(y)\) is well defined,
      and it makes the diagram (\ref{UPdiag}) commutative. Additivity of \(\psi\) follows easily from additivity of \(\varphi\).
      The uniqueness of \(\psi\) follows from (3).
      \item[2.] By (1), the additive map \(\gamma_T\circ\varphi:S\rightarrow G(T)\) factors uniquely through a group homomorphism
      \(G(\varphi):G(S)\rightarrow G(T)\).
      \item[6.] A non-empty subset \(S\) of an abelian group \(H\) which is closed under addition is an abelian semigroup with the cancellation property.
      By (1) applied to the inclusion map \(\iota:S\rightarrow H\), there is a group homomorphism \(\psi:G(S)\rightarrow H\) such that \(\psi(\gamma_S(x))=x\) for all \(x\) in \(S\).
      The image of \(\psi\) is precisely \(\{x-y| x,y\in S\}=H_0\) by (3). If \(\psi(\gamma_S(x)-\gamma_S(y))=0\),
      then \(x=y\) and so \(\gamma_S(x)-\gamma_S(y)=0\). Hence \(\psi\) is injective.
    \end{enumerate}
  \end{proof}
\end{proposition}
\begin{example}
  \begin{enumerate}
    \item The Grothendieck group of the abelian semigroup \((\mathbb{Z}^+,+)\) is (isomorphic to)
    \((\mathbb{Z},+)\) (Proposition~\ref{prop:613} (6)); and \((\mathbb{Z}^+,+)\) has the cancellation property.
    \item The Grothendieck group of the abelian semigroup \((\mathbb{Z}^+\cup \{\infty\},+)\), where
    \(n+\infty=\infty+n=\infty=\infty+\infty\), is \(\{0\}\), and \((\mathbb{Z}^+\cup \{\infty\},+)\) does not have the cancellation property.
  \end{enumerate}
\end{example}
\begin{definition}
  Let \(A\) be a unital \cstar, and let \((\mathcal{D}(A),+)\) be the abelian semigroup as defined above.
  Define \(K_0(A)\) to be the Grothendieck group of \(\mathcal{D}(A)\), i.e.,
  \[
    K_0(A)=G(\mathcal{D}(A)).
  \]
  Define \({[\cdot]}_0:\mathcal{P}_\infty(A)\rightarrow K_0(A)\) by
  \begin{equation}
    {[p]}_0=\gamma({[p]_\mathcal{D}})\in K_0(A),\qquad p\in \mathcal{P}_\infty(A),
  \end{equation}
  where \(\gamma:\mathcal{D}(A)\rightarrow K_0(A)\) is the Grothendieck map.
\end{definition}
\begin{definition}\label{def:Lars616}
  Define a relation \(\sim_s\) on \(\mathcal{P}_\infty(A)\) as follows. If \(p, q\) are projections in \(\mathcal{P}_\infty(A)\), then
  \(p\sim_s q\) if and only if \(p\oplus r\sim_0 q\oplus r\) for some projection \(r\) in \(\mathcal{P}_\infty(A)\). The relation \(\sim_s\) is called \textnormal{stable equivalence}.
\end{definition}
Suppose that \(A\) is unital and that \(p, q\) are projections in \(\mathcal{P}_\infty(A)\). Denote by \(1_n\) the unit of \(M_n(A)\).
Then \(p\sim_s q\) if and only if \(p\oplus 1_n\sim_0 q\oplus 1_n\) for some positive integer \(n\). Indeed, if \(p\oplus r\sim_0 q\oplus r\) for some \(r\) in \(\mathcal{P}_n(A)\), then
\[
  p\oplus 1_n\sim_0 p\oplus r\oplus (1_n-r)\sim_0 q\oplus r\oplus (1_n-r)\sim_0 q\oplus 1_n.
\]\par
The standard picture of \(K_0\), described by the two propositions below, is a concrete and useful description of the \(K_0\)-group of a unital \cstar.
Proposition~\ref{prop:618} shows that assertions (1), (2), and (3) of Proposition~\ref{prop:617} form a \textit{universal property} of \(K_0\).

\begin{proposition}\label{prop:617}
  Let \(A\) be a unital \cstar. Then
  \begin{align}\label{eq:32}
      K_0(A)&=\{{[p]}_0-{[q]}_0| p,q\in\mathcal{P}_\infty(A)\}\\
      &=\{{[p]}_0-{[q]}_0| p,q\in\mathcal{P}_n(A), n\in\mathbb{N}\}.
  \end{align}
  Moreover,
  \begin{enumerate}
    \item \({[p\oplus q]}_0={[p]}_0+{[q]}_0\) for all projections \(p, q\) in \(\mathcal{P}_\infty(A)\),
    \item \({[0_A]}_0=0\), where \(0_A\) is the zero projection in \(A\),
    \item if \(p, q\in\mathcal{P}_n(A)\) for some \(n\) and \(p\sim_h q\), then \({[p]}_0={[q]}_0\),
    \item if \(p, q\) are mutually orthogonal projections in \(\mathcal{P}_n(A)\), then \({[p+q]}_0={[p]}_0+{[q]}_0\),
    \item for all \(p, q\) in \(\mathcal{P}_\infty(A)\), \({[p]}_0={[q]}_0\) if and only if \(p\sim_s q\).
  \end{enumerate}
\end{proposition}
\begin{proof}
  The first identity follows from Proposition~\ref{prop:613} (3). Hence, if \(g\) is an element
  in \(K_0(A)\), then \(g={[p']}_0-{[q']}_0\) for some \(p'\) in \(\mathcal{P}_k(A)\) and \(q'\) in \(\mathcal{P}_l(A)\).
  Choose \(n\) greater than \(k\) and \(l\), and set \(p=p'\oplus 0_{n-k}\) and \(q=q'\oplus 0_{n-l}\).
  Then \(p, q\) are projections in \(\mathcal{P}_n(A)\) with \(p\sim_0 p'\) and \(q\sim_0 q'\) by Proposition~\ref{prop:232} (1).
  It follows that \(g={[p]}_0-{[q]}_0\).
  \begin{enumerate}
    \item We have
    \[
      {[p\oplus q]}_0=\gamma({[p\oplus q]}_\mathcal{D})=\gamma({[p]}_\mathcal{D}\oplus {[q]}_\mathcal{D})=\gamma({[p]}_\mathcal{D})+ \gamma({[q]}_\mathcal{D})={[p]}_0+{[q]}_0.
    \]
    \item Since \(0_A\oplus 0_A\sim_0 0_A\), (1) yields that \({[0_A]}_0+{[0_A]}_0={[0_A]}_0\), and hence \({[0_A]}_0=0\).
    \item It follows from the implications
    \[
      p\sim_h q\Rightarrow p\sim q\Rightarrow p\sim_0 q \Leftrightarrow {[p]}_\mathcal{D}={[q]}_\mathcal{D}\Rightarrow {[p]}_0={[q]}_0,
    \]
    where the first two relations are defined only when \(p\) and \(q\) are in the same matrix algebra over \(A\), and the other three relations are defined for all \(p, q\) in \(\mathcal{P}_\infty(A)\). 
    The first implication holds by Proposition~\ref{prop:527}.
    \item We have \(p+q\sim_0 p\oplus q\) (Proposition~\ref{prop:232}), and so \({[p+q]}_0={[p\oplus q]}_0={[p]}_0+{[q]}_0\) by (1).
    \item If \({[p]}_0={[q]}_0\), then there is a projection \(r\) in \(\mathcal{P}_\infty(A)\) with \({[p]}_\mathcal{D}+{[r]}_\mathcal{D}={[q]}_\mathcal{D}+{[r]}_\mathcal{D}\) by Proposition \ref{prop:613} (4).
    Hence \({[p\oplus r]}_\mathcal{D}={[q\oplus r]}_\mathcal{D}\), so that \(p\oplus r\sim_0 q\oplus r\), and therefore \(p\sim_s q\).\par
    Conversely, if \(p\sim_s q\), then there is a projection \(r\) in \(\mathcal{P}_\infty(A)\) with \(p\oplus r\sim_0 q\oplus r\).
    By (1) we get \({[p]}_0+{[r]}_0={[q]}_0+{[r]}_0\), and because \(K_0(A)\) is a group, we conclude that \({[p]}_0={[q]}_0\).
  \end{enumerate}
\end{proof}

\begin{proposition}[Universal property of \(K_0\)]\label{prop:618}
  Let \(A\) be a unital \cstar, let \(G\) be an abelian group, and suppose that \(\nu:\mathcal{P}_\infty(A)\rightarrow G\) is a map that satisfies
  \begin{enumerate}
    \item \(\nu(p\oplus q)=\nu(p)+\nu(q)\) for all projections \(p,q\in\mathcal{P}_\infty(A)\),
    \item \(\nu(0_A)=0\),
    \item if \(p,q\in\mathcal{P}_n(A)\) for some \(n\) and \(p\sim_h q\), then \(\nu(p)=\nu(q)\).
  \end{enumerate}
  Then there is a unique group homomorphism \(\alpha:K_0(A)\rightarrow G\) which makes the diagram
  \begin{equation}\label{diag:33}
  \begin{tikzcd}
  \mathcal{P}_\infty(A) \arrow[rd, "\nu"] \arrow[d, swap, "{{[\cdot]}_0}"]&\\
  K_0(A) \arrow[r, swap, "\alpha"]
  & G
  \end{tikzcd}
\end{equation}
commutative.
\end{proposition}
\begin{proof}
  We first show that if \(p, q\) are projections in \(\mathcal{P}_\infty(A)\)
  and if \(p\sim_0 q\), then \(\nu(p)=\nu(q)\). Find \(k, l\) in \(\mathbb{N}\) such that \(p\) belongs to \(\mathcal{P}_k(A)\)
  and \(q\) belongs to \(\mathcal{P}_l(A)\). Let \(n\) be an integer greater than \(k\) and \(l\), and put \(p'=p\oplus 0_{n-k}\)
  and \(q'=q\oplus 0_{n-l}\). Then \(p',q'\) both belong to \(\mathcal{P}_n(A)\),
  \(p'\sim_0 p\sim_0 q\sim_0 q'\), and so \(p'\sim q'\). By Proposition~\ref{prop:528}, \(p'\oplus 0_{3n}\sim q'\oplus 0_{3n}\) in \(\mathcal{P}_{4n}(A)\).
  We conclude that
  \[
    \nu(p)=\nu(p)+\underbrace{\nu(0)+\cdots+\nu(0)}_{4n-k}=\nu(p'\oplus 0_{3n})=\nu(q'\oplus 0_{3n})=\nu(q).
  \]\par
  It follows that the map \(\beta:\mathcal{D}(A)\rightarrow G\) given by \(\beta({[p]}_\mathcal{D})=\nu(p)\) is well-defined.
  Additivity of \(\beta\) follows from
  \begin{align*}
    \beta({[p]}_\mathcal{D}+{[q]}_\mathcal{D})&=\beta({[p\oplus q]}_\mathcal{D})=\nu(p\oplus q)=\nu(p)+\nu(q)\\
    &=\beta({[p]}_\mathcal{D})+\beta({[q]}_\mathcal{D}).
  \end{align*}
  Now use Proposition~\ref{prop:613} (1) to find a group homomorphism \(\alpha:K_0(A)\rightarrow G\) that makes the diagram commutative.
  The uniqueness of \(\alpha\) follows from (\ref{eq:32}).
\end{proof}

\section{Functoriality of \texorpdfstring{\(K_0\)}{K0}}\label{sec:622}
\begin{definition}
  A \textnormal{category} \(\mathbf{C}\) consists of a class \(\mathbf{\mathcal{O}(C)}\) of \textnormal{objects}
  and for each pair of objects \(A, B\) in \(\mathbf{\mathcal{O}(C)}\) a set \(\textnormal{Mor}(A,B)\) of \textnormal{morphisms} (from \(A\) to \(B\)) with an 
  associative rule of composition
  \[
    \textnormal{Mor}(A,B)\times\textnormal{Mor}(B,C)\rightarrow \textnormal{Mor}(A,C), \qquad (\varphi, \psi)\mapsto\psi\circ\varphi,
  \]
  such that for each object \(A\) there is an element \(\textnormal{id}_A\) in \(\textnormal{Mor}(A,A)\) which satisfies
  \(\textnormal{id}_B\circ\varphi=\varphi=\varphi\circ{\textnormal{id}}_A\) for every \(\varphi\) in \(\textnormal{Mor}(A,B)\).
\end{definition}
The two main examples of categories appearing in this text are the category \(\mathbf{C^*-alg}\) of \cstars\ and the category \(\mathbf{Ab}\) of abelian groups.
The class of objects, \(\mathbf{\mathcal{O}(C^*-alg)}\), in \(\mathbf{C^*-alg}\) is the class of all \cstars, and \(\textnormal{Mor}(A,B)\) is the set of *-homomorphisms 
from \(A\) to \(B\) with the usual composition. The objects in \(\mathbf{Ab}\) are abelian groups, and morphisms are group homomorphisms.
We shall in this chapter also consider the category of unital \cstars, where the morphisms are *-homomorphisms (not necessarily unit-preserving).
\begin{definition}
  A \textnormal{covariant functor} \(F\) between categories \(\mathbf{C}\) and \(\mathbf{D}\) is a map \(A\mapsto F(A)\)
  from \(\mathbf{\mathcal{O}(C)}\) to \(\mathbf{\mathcal{O}(D)}\) and a map \(\varphi\mapsto F(\varphi)\) from 
  \(\textnormal{Mor}(A,B)\) to \(\textnormal{Mor}(F(A),F(B))\) for each pair of objects \(A, B\) in \(\mathbf{\mathcal{O}(C)}\) such that
  \begin{enumerate}
    \item \(F(\textnormal{id}_A)=\textnormal{id}_{F(A)}\) for all objects \(A\) in \(\mathbf{\mathcal{O}(C)}\),
    \item\label{functordef:2} \(F(\psi\circ\varphi)=F(\psi)\circ F(\varphi)\) for all objects \(A, B, C\) in \(\mathbf{\mathcal{O}(C)}\) and all morphisms \(\varphi\) in \(\textnormal{Mor}(A,B)\) and \(\psi\) in \(\textnormal{Mor}(B,C)\).
  \end{enumerate}\par
  A \textnormal{contravariant functor} \(F\) is like a covariant functor with the exception that it reverses the direction of morphisms, i.e., it maps a morphism \(\varphi\) 
  in \(\textnormal{Mor}(A,B)\) to a morphism \(F(\varphi)\) in \(\textnormal{Mor}(F(B),F(A))\), and (\ref{functordef:2})
  is changed accordingly.\par
  An object \(N\) in a category \(\mathbf{C}\) is called a \textnormal{zero object} if \(\textnormal{Mor}(A,N)\) and \(\textnormal{Mor}(N,A)\) contain precisely
  one element for all objects \(A\) in \(\mathbf{C}\).
\end{definition}
Let \(A\) and \(B\) be unital \cstars, and let \(\varphi:A\rightarrow B\) be a *-homomorphism. Associate to \(\varphi\) a group
homomorphism \(K_0(\varphi):K_0(A)\rightarrow K_0(B)\) as follows. We know that \(\varphi\) extends entrywise to a *-homomorphism \(\varphi:M_n(A)\rightarrow M_n(B)\) for each \(n\).
A *-homomorphism maps projections to projections, and so \(\varphi\) maps \(\mathcal{P}_\infty(A)\) into \(\mathcal{P}_\infty(B)\).
Define \(\nu:\mathcal{P}_\infty(A)\rightarrow K_0(B)\) by \(\nu(p)={[\varphi(p)]}_0\) for \(p\) in \(\mathcal{P}_\infty(A)\).
Then \(\nu\) satisfies (1), (2) and (3) in Proposition~\ref{prop:618}, and \(\nu\) therefore factors uniquely through a group homomorphism
\(K_0(\varphi):K_0(A)\rightarrow K_0(B)\) given by
\begin{equation}\label{eq:34}
  K_0(\varphi)({[p]}_0)={[\varphi(p)]}_0,\qquad p\in\mathcal{P}_\infty(A).
\end{equation}
In other words, we have a commutative diagram:
\begin{equation}\label{eq:35}
\begin{tikzcd}
\mathcal{P}_\infty(A) \arrow[r, "\varphi"] \arrow[d, swap, "{[\cdot]}_0"]
& \mathcal{P}_\infty(B) \arrow[d, "{[\cdot]}_0"] \\
K_0(A) \arrow[r, swap, "K_0(\varphi)"]
& K_0(B)
\end{tikzcd}
\end{equation}
If \(A\) and \(B\) are \cstars, then denote the zero homomorphism \(A\rightarrow B\) that sends all elements in \(A\) to the zero element in \(B\) by \(0_{B,A}\) or just \(0\).
The identity map \(A\rightarrow A\) is denoted by \(\textnormal{id}_A\). We shall use a similar notation in the category of abelian groups. 
\begin{proposition}\label{prop:624}
  \begin{enumerate}
    \item For each unital \cstar\ \(A\), \(K_0(\textnormal{id}_A)=\textnormal{id}_{K_0(A)}\).
    \item If \(A, B\) and \(C\) are unital \cstars, and if \(\varphi:A\rightarrow B\) and \(\psi:B\rightarrow C\) 
    are *-homomorphisms, then \(K_0(\psi\circ\varphi)=K_0(\psi)\circ K_0(\varphi)\).
    \item \(K_0(\{0\})=\{0\}\) (we include the zero \cstar\ in the ranks of unital \cstars).
    \item For every pair of unital \cstars\ \(A\) and \(B\), \(K_0(0_{B,A})=0_{K_0(B),K_0(A)}\).
  \end{enumerate}
\end{proposition}
\begin{proof}
  Use (\ref{eq:34}) to check that
  \[
    K_0(\textnormal{id}_A)({[p]}_0)={[p]}_0,\qquad K_0(\psi\circ\varphi)({[p]}_0)=(K_0(\psi)\circ K_0(\varphi))({[p]}_0),
  \]  
  for every \(p\) in \(\mathcal{P}_\infty(A)\). Then use the standard picture of \(K_0\) (equation (\ref{eq:32})) to conclude that (1) and (2) hold.
  \begin{enumerate}
    \item[3.] We have \(\mathcal{P}_n(\{0\})=\{0_n\}\), where \(0_n\) is the zero element in \(M_n(\{0\})\).
    The zero projections \(0=0_1,0_2,\ldots\) are all equivalent, and so \(\mathcal{D}(\{0\})=\{{[0]_\mathcal{D}}\}\).
    Hence \(K_0(\{0\})=G(\{0\})=\{0\}\).
    \item[4.] Since \(0_{B,A}=0_{B,0}\circ 0_{0,A}:A\rightarrow \{0\}\rightarrow B\), (4) follows from (2) and (3).
  \end{enumerate}
\end{proof}
\begin{definition}\label{def:Lars625}
  Let \(A\) and \(B\) be \cstars. Two *-homomorphisms \(\varphi, \psi:A\rightarrow B\) are said to be \textnormal{homotopic}, in symbols \(\varphi\sim_h\psi\), if there is a path of *-homomorphisms 
  \(\varphi_t:A\rightarrow B\), \(t\in[0,1]\), such that \(t\mapsto \varphi_t(a)\) is a continuous map from \([0,1]\) to \(B\) for each \(a\) in \(A\),
  \(\varphi_0=\varphi\) and \(\varphi_1=\psi\). We say that the path \(t\mapsto \varphi_t\) is \textnormal{point-wise continuous}.\par
  The \cstars\ \(A\) and \(B\) are \textnormal{homotopy equivalent} if there are *-homomorphisms \(\varphi:A\rightarrow B\) and \(\psi:B\rightarrow A\) such that \(\psi\circ\varphi\sim_h\textnormal{id}_A\)
  and \(\varphi\circ\psi\sim_h\textnormal{id}_B\). In this case we say that
  \[
    \begin{tikzcd}
    A \arrow[r, "\varphi"]& B \arrow[r, "\psi"]&A
    \end{tikzcd}
  \]
  is a \textnormal{homotopy} (between \(A\) and \(B\)).
\end{definition}
\begin{proposition}[Homotopy invariance of \(K_0\)]\label{prop:626}
  Let \(A\) and \(B\) be unital \cstars.
  \begin{enumerate}
    \item If \(\varphi, \psi:A\rightarrow B\) are homotopic *-homomorphisms, then \(K_0(\varphi)=K_0(\psi)\).
    \item If \(A\) and \(B\) are homotopy equivalent, then \(K_0(A)\) is isomorphic to \(K_0(B)\). More specifically, if
    \[
    \begin{tikzcd}
    A \arrow[r, "\varphi"]& B \arrow[r, "\psi"]&A
    \end{tikzcd}
    \]
    is a homotopy, then \(K_0(\varphi):K_0(A)\rightarrow K_0(B)\) and \(K_0(\psi):K_0(B)\rightarrow K_0(A)\) are isomorphisms,
    and \({K_0(\varphi)}^{-1}=K_0(\psi)\).
  \end{enumerate}
\end{proposition}
\begin{proof}
  \begin{enumerate}
    \item Let \(\varphi_t:A\rightarrow B\) be a point-wise continuous path of *-homomorphisms connecting \(\varphi\) to \(\psi\). 
    Extend this path to a point-wise path of *-homomorphisms \(\varphi_t:M_n(A)\rightarrow M_n(B)\) for each \(n\) in \(\mathbb{N}\).
    For every projection \(p\) in \(\mathcal{P}_n(A)\), the path \(t\mapsto \varphi_t(p)\) is continuous, and so
    \(\varphi(p)=\varphi_0(p)\sim_h\varphi_1(p)=\psi(p)\). This shows that
    \[
      K_0(\varphi)({[p]}_0)={[\varphi(p)]}_0=K_0(\psi)({[p]}_0).
    \]
    Use the equation (\ref{eq:32}) to conclude that \(K_0(\varphi)=K_0(\psi)\).
    \item This follows from (1) and from Proposition~\ref{prop:624} (1) and (2).
  \end{enumerate}
\end{proof}
We proceed to show that \(K_0\) preserves exactness of the short exact sequence obtained by adjoining a unit
to a unital \cstar. This result will be useful when we define \(K_0\) for non-unital \cstars.\par
\begin{definition}
  Two *-homomorphisms \(\varphi,\psi:A\to B\) between \cstars\ \(A\) and \(B\) are said to be \textnormal{orthogonal to each other},
  or mutually orthogonal, in symbols, \(\varphi\perp\psi\), if \(\varphi(x)\psi(y)=0\) for all \(x,y\in A\). 
\end{definition}
\begin{lemma}\label{lemma:627}
  If \(A\) and \(B\) are unital \cstars, and if \(\varphi,\psi:A\to B\) are mutually orthogonal *-homomorphisms,
  then \(\varphi+\psi:A\to B\) is a *-homomorphism, and \(K_0(\varphi+\psi)=K_0(\varphi)+K_0(\psi)\).
\end{lemma}
\begin{proof}
  It is straightforward to check that \(\varphi+\psi\) is a *-homomorphism. The *-homomorphisms
  \(\varphi_n,\psi_n:M_n(A)\to M_n(B)\) induced by \(\varphi\) and \(\psi\) are orthogonal
  to each other for each \(n\) in \(\mathbb{N}\), and \((\varphi+\psi)_n=\varphi_n+\psi_n\).
  Using the standard picture of \(K_0\) (Proposition \ref{prop:617} (4)) we obtain for every \(p\) in \(\mathcal{P}_n(A)\)
  \begin{align*}
    K_0(\varphi+\psi)({[p]}_0)&={[(\varphi+\psi)_n(p)]}_0={[\varphi_n(p)+\psi_n(p)]}_0\\
    &={[\varphi_n(p)]}_0+{[\psi_n(p)]}_0\\
    &=K_0(\varphi)({[p]}_0)+K_0(\psi)({[p]}_0).
  \end{align*}
  This shows that \(K_0(\varphi+\psi)=K_0(\varphi)+K_0(\psi)\).
\end{proof}
\begin{lemma}\label{lemma:628}
  For every unital \cstar\ \(A\), the split exact sequence 
  \[
    \begin{tikzcd}
      0\arrow[r]&A \arrow[r, "\iota"]& \tilde{A} \arrow[r, "\pi"]&\mathbb{C}\arrow[l, shift left=2pt, "\lambda"]\arrow[r]&0,
    \end{tikzcd}
  \]
  obtained by adjoining a unit to \(A\), induces a split exact sequence
  \begin{equation}\label{diag:splitExact36}
    \begin{tikzcd}
      0\arrow[r]&K_0(A) \arrow[r, "K_0(\iota)"]& K_0(\tilde{A}) \arrow[r, "K_0(\pi)"]&K_0(\mathbb{C})\arrow[l, shift left=2pt, "K_0(\lambda)"]\arrow[r]&0.
    \end{tikzcd}
  \end{equation}
\end{lemma}
\begin{proof}
  Recall from the construction of adjoining a unit to a (unital) \cstar\ that 
  if \(f\) is the projection \(1_{\tilde{A}}-1_A\) in \(\tilde{A}\), then \(\tilde{A}=A+\mathbb{C}f\) and
  \(af=fa=0\) for all \(a\) in \(A\). Define *-homomorphisms
  \[
    \mu:\tilde{A}\to A,\qquad \lambda:\mathbb{C}\to \tilde{A}
  \]
  by \(\mu(a+\alpha f)=a\) and \(\lambda(\alpha)=\alpha f\). Then
  \[
    \operatorname{id}_A=\mu\circ\iota,\qquad\operatorname{id}_{\tilde{A}}=\iota\circ\mu+\lambda\circ\pi,\qquad\pi\circ\lambda=\operatorname{id}_{\mathbb{C}},
  \]
  and the *-homomorphisms \(\iota\circ\mu\) and \(\lambda\circ\pi\) are orthogonal to each other.
  Functoriality of \(K_0\) and Lemma \ref{lemma:627} now yield
  \begin{align*}
    0&=K_0(0)=K_0(\pi\circ\iota)=K_0(\pi)\circ K_0(\iota),\\
    \operatorname{id}_{K_0(\mathbb{C})}&=K_0(\operatorname{id}_\mathbb{C})=K_0(\pi\circ\lambda)=K_0(\pi)\circ K_0(\lambda),\\
    \operatorname{id}_{K_0(A)}&=K_0(\operatorname{id}_A)=K_0(\mu\circ\iota)=K_0(\mu)\circ K_0(\iota),\\
    \operatorname{id}_{K_0(\tilde{A})}&=K_0(\operatorname{id}_{\tilde{A}})=K_0(\iota\circ\mu+\lambda\circ\pi)\\
    &=K_0(\iota)\circ K_0(\mu)+K_0(\lambda)\circ K_0(\pi).
  \end{align*}
  Split exactness of (\ref{diag:splitExact36}) follows from these four identities: injectivity of \(K_0(\iota)\)
  follows from the third identity. If \(g\) belongs to the kernel of \(K_0(\pi)\), then \(g=K_0(\iota)(K_0(\mu)(g))\) by the last identity,
  thus showing that \(g\) belongs to the image of \(K_0(\iota)\).
\end{proof}

\section{Examples and exercises}
\begin{example}[Algebraic construction of \(K_0\)]\label{ex:3320}
  Let \(R\) be a unital ring. Recall that \(e\in R\) is \textnormal{idempotent}
  if \(e^2=e\). We define \(\mathcal{I}(R)=\{e\in R\mid e^2=e\}\), \(\mathcal{I}_n(R)=\mathcal{I}(M_n(R))\),
  \(\mathcal{I}_\infty(R)=\bigcup_{n=1}^\infty\mathcal{I}_n(R)\).
  We define a relation \(\approx_0\) in \(\mathcal{I}_\infty(R)\) as follows.
  If \(e\in\mathcal{I}_n(R)\) and \(f\in\mathcal{I}_m(R)\), then \(e\approx_0 f\)
  if and only if there exist \(a\in M_{n,m}(R)\) and \(b\in M_{m,n}(R)\) such that \(e=ab\) and \(f=ba\).
  If this is the case then taking \(a'=aba\) and \(b'=bab\) we may assume that \(a,b\) satisfy \(aba=a\)
  and \(bab=b\) (this we always assume in what follows).\\
  Claim: \(\approx_0\) is an equivalence relation. For transitivity,
  let \(e\approx_0 f\), \(f\approx_0 g\) be idempotents and let \(c,d,x,y\) be matrices such that
  \(e=cd\), \(f=dc\), \(f=xy\), \(g=yx\). Then \((cx)(yd)=e\) and \((yd)(cx)=g\), 
  and hence \(e\approx_0 g\). Set \(V(R) = \mathcal{I}_\infty(R)/\approx_0\), and denote the class of \(e\) by \({[e]}_V\).\par
  Define a binary operation \(\oplus\) on \(\mathcal{I}_\infty(R)\) by \(e\oplus f=\diag(e,f)\).
  This operation is well-defined on equivalence classes of \(\approx_0\) and turns \(V(R)\) into an abelian semigroup.
  Define \(K_0(R)\) as the Grothendieck group of \((V(R),\oplus)\).\par
  Now suppose that \(A\) is a unital \cstar. We show that the two definitions of \(K_0(A)\) coincide.
  In fact, the two semigroups \(\mathcal{D}(A)\) and \(V(A)\) are isomorphic. The proof proceeds in three steps.
  \begin{enumerate}
    \item If \(e\in\mathcal{I}_\infty(A)\) then there exists a \(p\in\mathcal{P}_\infty(A)\) such that \(e\approx_0 p\).
      Indeed, let \(e\in M_n(A)\), and set \(h=1_n+(e-e^*)(e-e^*)^*\). Then \(h\) is invertible and satisfies \(eh=ee^* e=he\).
      Then \(p=ee^* h^{-1}\) is a projection. Since \(ep=p\) and \(pe=e\), \(e\approx_0 p\).
    \item If \(p,q\in\mathcal{P}_\infty(A)\) then \(p\sim_0 q\) if and only if \(p\approx_0 q\). Indeed, suppose (after diagonally adding 
    zeros, if necessary) \(p,q\in M_n(A)\) and \(a,b\in M_n(A)\) are such that \(p=ab\), \(q=ba\), \(a=aba\) (hence \(a=paq\)), \(b=bab\) (hence \(b=qbp\)).
    Then \(b^* b=(bab)^* b=(ab)^*b^*b=pb^*b\).
    It follows that \(b^*b\) belongs to the corner \cstar\ \(pM_n(A)p\).
    Since \(p=(ab)^*ab=b^*(a^*a)b\leq\|a\|^2 b^*b\), the element \(b^* b\) is invertible in \(pM_n(A)p\).
    Set \(v=bp(b^*b)^{-1/2}\). We have \(p=v^* v\) (straightforward calculation). In particular, \(v\) is a partial isometry. In fact, \(b=v|b|\)
    is the polar decomposition of \(b\) in \(M_n(A)\). Hence \(bb^*=vb^*bv^*\).\par
    It now suffices to show that \(q=vv^*\). First note that \(v=qv\) (directly follows from the definition of \(v\)). Thus
    \[
      vv^*=qvv^*q\leq\|v\|^2 q=q=baa^*b^*\leq\|a\|^2bb^*=\|a\|^2vb^*bv^*\leq\|a\|^2\|b\|^2vv^*.
    \]
    That is, \(vv^*\) and \(q\) are projections satisfying \(vv^*\leq q\leq \|a\|^2\|b\|^2vv^*\). It follows that
    \(vv^*=q\).
    \item The map \(\mathcal{D}(A)\to V(A)\) given by \({[p]}_\mathcal{D}\mapsto{[p]}_V\) is a semigroup isomorphism. (exercise)
  \end{enumerate}
\end{example}

Let \(A\) be a unital \cstar\ and let \(\tau:A\to\mathbb{C}\) be a bounded trace. For each
\(n\) in \(\mathbb{N}\), extend \(\tau\) to a trace \(\tau_n:M_n(A)\to\mathbb{C}\) by
\[
  \tau_n((a_{ij}))=\sum_{i=1}^n\tau(a_{ii}).
\]
As usual, we omit the subscript and write \(\tau\) for all these extensions. The map
\[
  \nu:\mathcal{P}_\infty(A)\to\mathbb{C},\qquad \nu(p)=\tau(p),
\]
is additive under direct sums, vanishes at the zero projection, and is constant on homotopy
classes of projections. Hence the universal property of \(K_0\), Proposition~\ref{prop:618},
yields a unique group homomorphism
\[
  K_0(\tau):K_0(A)\to\mathbb{C}
\]
such that
\[
  K_0(\tau)([p]_0)=\tau(p),\qquad p\in\mathcal{P}_\infty(A).
\]
Thus, the notation \(K_0(\tau)\) denotes the homomorphism induced by the trace through the
universal property; it does not mean that the functor \(K_0\) is being applied to a
*-homomorphism \(\tau\). See \cite[Section~3.3.1]{RordamLarsenLaustsen2000}.

\begin{example}\label{ex:Lars332}
  The group \(K_0(M_n(\mathbb{C}))\) is isomorphic to \(\mathbb{Z}\) for each positive integer \(n\).
  More specifically, if \(\operatorname{Tr}\) is the standard trace on \(M_n(\mathbb{C})\), (see Exercise \ref{ex:Lars29}), then
  \[
    K_0(\operatorname{Tr}):K_0(M_n(\mathbb{C}))\to\mathbb{Z}
  \]
  is an isomorphism. The cyclic group \(K_0(M_n(\mathbb{C}))\) is generated by \({[e]}_0\),
  where \(e\) is any one-dimensional projection in \(M_n(\mathbb{C})\).\par
  Since \(\mathbb{C}=M_1(\mathbb{C})\) we obtain in particular that \(K_0(\mathbb{C})\) is isomorphic to \(\mathbb{Z}\).
\end{example} 
\begin{proof}
  Let \(g\) be an element in \(K_0(M_n(\mathbb{C}))\). By the standard picture of \(K_0\) (equation \ref{eq:32})
  we can find \(k\) in \(\mathbb{N}\) and projections \(p,q\) in \(M_k(M_n(\mathbb{C}))=M_{kn}(\mathbb{C})\) such that \(g={[p]}_0-{[q]}_0\). Now,
  \[
    K_0(\operatorname{Tr})(g)=\operatorname{Tr}(p)-\operatorname{Tr}(q)=\dim(p(\mathbb{C}^{kn}))-\dim(q(\mathbb{C}^{kn})).
  \]
  We see that \(K_0(\operatorname{Tr})(g)\) is an integer. If \(K_0(\operatorname{Tr})(g)=0\),
  then \(p\sim q\) by Exercise \ref{ex:Lars29}, whence \(g={[p]}_0-{[q]}_0=0\). Hence \(K_0(\operatorname{Tr})\) is injective.\par
  The image of \(K_0(\operatorname{Tr})\) is a subgroup of \(\mathbb{Z}\), and a subgroup of \(\mathbb{Z}\) is equal to \(\mathbb{Z}\) if and only if it contains \(1\).
  Now, \(1=K_0(\operatorname{Tr})({[e]}_0)\), when \(e\) is a projection in \(M_n(\mathbb{C})\) with one-dimensional range.
\end{proof}

\chapter{The Functor \texorpdfstring{\(K_0\)}{K0}}\label{chap:K0Functor}
We extend the functor \(K_0\) considered in the previous Chapter to a functor from the
category of all \cstars, unital or not. The \(K_0\)-group of a non-unital \cstar\
is defined as a relative \(K\)-group, and it is shown this definition is
coherent with the definition of \(K_0\) for unital \cstars.\par
We give a standard picture of \(K_0\) (a concrete description of the elements
  in \(K_0(A)\)). This new standard picture is less pedestrian than the standard
picture for \(K_0\) for unital \cstars\ in Proposition \ref{prop:617}, reflecting
the greater complexity of \(K_0\) in the non-unital case.\par
It is shown that \(K_0\) is a functor which is half exact, split exact, and stable.

\section{Definition and functoriality of \texorpdfstring{\(K_0\)}{K0}}
\begin{definition}
  Let \(A\) be a non-unital \cstar\, and consider the associated split exact sequence
  \begin{equation}\label{diag:splitExact41}
    \begin{tikzcd}
      0 \arrow[r]&A\arrow[r,"\iota"] & \tilde{A}\arrow[r, shift left=2pt, "\pi"]
      &\mathbb{C}\arrow[l, shift left=2pt, "\lambda"]\arrow[r]&0,
    \end{tikzcd}
  \end{equation}
  defined by adjoining a unit to \(A\) (see Section~\ref{sec:adjoiningUnit}). 
  Define \(K_0(A)\) to be the kernel of the homomorphism \(K_0(\pi):K_0(\tilde{A})\to K_0(\mathbb{C})\).
\end{definition}
Observe that \(K_0(A)\) is an abelian group, being a subgroup of \(K_0(\tilde{A})\).\par
For each \(p\) in \(\mathcal{P}_\infty(A)\) consider the equivalence class \({[p]}_0\) in \(K_0(\tilde{A})\).
Since \(K_0(\pi)({[p]}_0)={[\pi(p)]}_0=0\), it follows that \({[p]}_0\) belongs to \(K_0(A)\). In this way we obtain a map \({[\cdot]}_0:\mathcal{P}_\infty(A)\to K_0(A)\).\par
For each \cstar\ \(A\), unital or not, we have a short exact sequence
\begin{equation}\label{diag:splitExact42}
  \begin{tikzcd}
    0 \arrow[r]&K_0(A)\arrow[r] & K_0(\tilde{A})\arrow[r, shift left=2pt, "K_0(\pi)"]
    &K_0(\mathbb{C})\arrow[r]&0,
  \end{tikzcd}
\end{equation}
The map \(K_0(A)\to K_0(\tilde{A})\) is \(K_0(\iota)\) when \(A\) is unital, and it is the inclusion map when \(A\) is non-unital.
Use Lemma \ref{lemma:628} to see that (\ref{diag:splitExact42}) is exact when \(A\) is unital.\par
When \(A\) is a unital \cstar\, \(K_0(A)\) is isomorphic to its image in \(K_0(\tilde{A})\)
under \(K_0(\iota)\), and \(K_0(\iota)\) maps \({[p]}_0\) in \(K_0(A)\) to \({[p]}_0\) in \(K_0(\tilde{A})\) for \(p\) in \(\mathcal{P}_\infty(A)\).
The image of \(K_0(\iota)\) is equal to the kernel of \(K_0(\pi)\), so that the identity
\[
  K_0(A)=\operatorname{Ker}(K_0(\pi))
\]
holds for both unital and non-unital \cstars\ (with a slight abuse of notation).\par
Let \(\varphi:A\to B\) be a *-homomorphism, and let \(\tilde{\varphi}:\tilde{A}\to\tilde{B}\) be its unitization. 
The commutative diagram
\begin{equation*}
  \begin{tikzcd}
    A\arrow[r,"\iota_A"]\arrow[d,"\varphi"]&\tilde{A}\arrow[r,"\pi_A"]\arrow[d,"\tilde{\varphi}"]&\mathbb{C}\arrow[d, equal]\\    
    B\arrow[r,"\iota_B",swap]&\tilde{B}\arrow[r,"\pi_B",swap]&\mathbb{C}
  \end{tikzcd}
\end{equation*}
and (\ref{diag:splitExact42}) induce by functoriality of \(K_0\) for unital \cstars\ the commutative diagram
\begin{equation}\label{diag:splitExact43}
  \begin{tikzcd}
    K_0(A)\arrow[r]\arrow[d,"K_0(\varphi)",dashed]&K_0(\tilde{A})\arrow[r,"K_0(\pi_A)"]\arrow[d,"K_0(\tilde{\varphi})"]&K_0(\mathbb{C})\arrow[d, equal]\\    
    K_0(B)\arrow[r]&K_0(\tilde{B})\arrow[r,"K_0(\pi_B)",swap]&K_0(\mathbb{C}).
  \end{tikzcd}
\end{equation}
There is one and only one group homomorphism \(K_0(\varphi):K_0(A)\to K_0(B)\), indicated with a dashed arrow, that makes the diagram commutative, and \(K_0(\varphi)\)
is the restriction of \(K_0(\tilde{\varphi})\) to \(K_0(A)\).\par
Suppose that both \(A\) and \(B\) are unital \cstars. Then \(K_0(\varphi)\) has already been defined in Paragraph \ref{sec:622}.
By functoriality of \(K_0\) in the category of unital \cstars\ (Proposition \ref{prop:624}) the diagram (\ref{diag:splitExact43}), with
\(K_0(\varphi)\) from Paragraph \ref{sec:622} is commutative. This shows that the present definition of \(K_0(\varphi)\) extends Paragraph \ref{sec:622}.\par
Notice that 
\begin{equation}\label{eq:K044}
  K_0(\varphi)({[p]}_0)={[\varphi(p)]}_0,\qquad p\in\mathcal{P}_\infty(A),
\end{equation}
when \(A\) is any \cstar, unital or not.
\begin{proposition}[Functoriality of \(K_0\)]\label{prop:Lars713}
  \begin{enumerate}
    \item \(K_0(\operatorname{id}_A)=\operatorname{id}_{K_0(A)}\) for every \cstar\ \(A\).
    \item If \(A,B\) and \(C\) are \cstars, and if \(\varphi:A\to B\) and \(\psi:B\to C\) are *-homomorphisms, then \(K_0(\psi\circ\varphi)=K_0(\psi)\circ K_0(\varphi)\).
    \item \(K_0(\{0\})=\{0\}\).
    \item \(K_0(0_{B,A})=0_{K_0(B),K_0(A)}\) for every pair of \cstars\ \(A\) and \(B\).
  \end{enumerate}
\end{proposition}
\begin{proof}
  Parts (1) and (2) follow from Proposition \ref{prop:624} (1) and (2) because
  \(\widetilde{\operatorname{id}_A}=\operatorname{id}_{\tilde{A}}\) and \(\widetilde{\psi\circ\varphi}=\tilde{\psi}\circ\tilde{\varphi}\).
  (3). We have \(\widetilde{\{0\}}=\mathbb{C}\), and the sequence (derived from (\ref{diag:splitExact41}))
  \begin{equation*}
    \begin{tikzcd}
      0\arrow[r]&0\arrow[r]&\mathbb{C}\arrow[r,"\pi"]&\mathbb{C}\arrow[r]&0,    
    \end{tikzcd}
  \end{equation*}
  where \(\pi=\operatorname{id}_{\mathbb{C}}\). It follows that \(K_0(\{0\})=\operatorname{Ker}(K_0(\pi))=\{0\}\).
  (4). This follows from (3) as in the proof of Proposition \ref{prop:624}.
\end{proof}
\begin{proposition}[Homotopy invariance of \(K_0\)]
  Let \(A\) and \(B\) be \cstars.
  \begin{enumerate}
    \item If \(\varphi,\psi:A\to B\) are homotopic *-homomorphisms, then \(K_0(\varphi)=K_0(\psi)\).
    \item If \(A\) and \(B\) are homotopy equivalent \cstars, then \(K_0(A)\) is isomorphic to \(K_0(B)\). More specifically, if
      \begin{equation*}
        \begin{tikzcd}
          A\arrow[r,"\varphi"]&B\arrow[r,"\psi"]&A    
        \end{tikzcd}
      \end{equation*}
      is a homotopy, then \(K_0(\varphi):K_0(A)\to K_0(B)\) and \(K_0(\psi):K_0(B)\to K_0(A)\)
      are isomorphisms and \(K_0(\varphi)^{-1}=K_0(\psi)\).
  \end{enumerate}
\end{proposition}
\begin{proof}
  \begin{enumerate}
    \item If \(\varphi\) is homotopic to \(\psi\), then \(\tilde{\varphi}\) is homotopic to \(\tilde{\psi}\), whence 
      \(K_0(\tilde{\varphi})=K_0(\tilde{\psi})\) by Proposition \ref{prop:626}. It follows that \(K_0(\varphi)\) and \(K_0(\psi)\) are equal,
      being restrictions of \(K_0(\tilde{\varphi})\) and \(K_0(\tilde{\psi})\) to \(K_0(A)\).
    \item This follows from (1) and from Proposition \ref{prop:Lars713} (1) and (2).
  \end{enumerate}
\end{proof}

\begin{example}\label{example:449}
  Let \(A\) be a \cstar. We define the \textit{cone} \(CA\) and the \textit{suspension} \(SA\)
  as follows:
  \begin{align*}
    CA &= \{f : [0,1]\rightarrow A| f\textnormal{ continuous and }f(0) = 0\},\\
    SA &= \{f : [0,1] \rightarrow A| f\textnormal{ continuous and }f(0) = f(1) = 0\}.
  \end{align*}
  There is a short exact sequence
  \[
    \begin{tikzcd}
      0 \arrow[r]&S A\arrow[r] & C A \arrow[r,"\pi"]&A\arrow[r]&0,
    \end{tikzcd}
  \]
  with \(\pi(f) = f(1)\). Furthermore, \(C A\) is homotopy equivalent to \(\{0\}\). Indeed, with \(t \in [0,1]\)
  set \(\varphi_t : C A \rightarrow C A\) as \(\varphi_t(f)(s) = f(st)\). Then for each \(f \in C A\) the map \(t \mapsto \varphi_t(f)\) is
  continuous, and \(\varphi_0 = 0\), \(\varphi_1 = \textnormal{id}\). We conclude that
  \[K_0(C A) = 0.\]
\end{example}

\section{The standard picture of the group \texorpdfstring{\(K_0(A)\)}{K0(A)}}
Proposition \ref{prop:Lars722} and Lemma \ref{lemma:Lars723} below contain the standard picture of \(K_0\):
a handy description of what elements in it look like. These results will
be invoked whenever explicit calculations involving \(K_0\)-groups are performed.\par
Recall from Section~\ref{sec:adjoiningUnit} the scalar maps
\[
  s_n:M_n(\tilde{A})\to M_n(\tilde{A}).
\]
The image of \(s_n\) is the canonical copy of \(M_n(\mathbb{C})\) in
\(M_n(\tilde{A})\), and \(x-s_n(x)\) belongs to \(M_n(A)\) for every
\(x\in M_n(\tilde{A})\). We shall often omit the subscript and write \(s\)
instead of \(s_n\).\par
An element \(x\) in \(\tilde{A}\) or in \(M_n(\tilde{A})\) will be called a \textnormal{scalar element} if \(x=s(x)\).
Equivalently, an element \(x\) in \(M_n(\tilde{A})\) is scalar if all its entries are scalar multiples of the unit \(1_{\tilde{A}}\).\par
The scalar mapping is natural in the sense that if \(A\) and \(B\) are \cstars, and if \(\varphi:A\to B\) 
is a *-homomorphism, then we get a commuting diagram
\begin{equation*}
  \begin{tikzcd}
    \tilde{A}\arrow[r,"s"]\arrow[d,"\tilde{\varphi}",swap]&\tilde{A}\arrow[d,"\tilde{\varphi}"]\\
    \tilde{B}\arrow[r,"s",swap]&\tilde{B}.
  \end{tikzcd}
\end{equation*}

\begin{proposition}\label{prop:Lars722}
  For each \cstar\ \(A\) 
  \begin{equation}\label{eq:Lars45}
    K_0(A)=\{{[p]}_0-{[s(p)]}_0\mid p\in\mathcal{P}_\infty(\tilde{A})\}.
  \end{equation}
  Moreover, the following hold.
  \begin{enumerate}
    \item For each pair of projections \(p,q\) in \(\mathcal{P}_\infty(\tilde{A})\), the following conditions are equivalent:
      \begin{enumerate}
        \item \({[p]}_0-{[s(p)]}_0={[q]}_0-{[s(q)]}_0\),
        \item there exist natural numbers \(k\) and \(l\) such that \(p\oplus1_k\sim_0 q\oplus1_l\) in \(\mathcal{P}_\infty(\tilde{A})\),
        \item there exist scalar projections \(r_1\) and \(r_2\) such that \(p\oplus r_1\sim_0 q\oplus r_2\).
      \end{enumerate}
    \item If \(p\) in \(\mathcal{P}_\infty(\tilde{A})\) satisfies \({[p]}_0-{[s(p)]}_0=0\), then there is a natural number \(m\) with \(p\oplus1_m\sim s(p)\oplus 1_m\).
    \item If \(\varphi:A\to B\) is a *-homomorphism, then 
      \[
        K_0(\varphi)({[p]}_0-{[s(p)]}_0)={[\tilde{\varphi}(p)]}_0-{[s(\tilde{\varphi}(p))]}_0
      \]
      for each \(p\) in \(\mathcal{P}_\infty(\tilde{A})\).
  \end{enumerate}
\end{proposition}
\begin{proof}
  To prove that equation (\ref{eq:Lars45}) holds, notice that
  \[
    K_0(\pi)({[p]}_0-{[s(p)]}_0)={[\pi(p)]}_0-{[(\pi\circ s)(p)]}_0=0
  \]
  when \(p\) belongs to \(\mathcal{P}_\infty(\tilde{A})\) because \(\pi=\pi\circ s\). This shows that \({[p]}_0-{[s(p)]}_0\)
  belongs to \(K_0(A)\) for all projections \(p\) in \(\mathcal{P}_\infty(\tilde{A})\).\par
  Conversely, let \(g\) be an element in \(K_0(A)\), and find \(n\) in \(\mathbb{N}\) and projections \(e,f\) in \(M_n(\tilde{A})\) such that \(g={[e]}_0-{[f]}_0\). Put
  \[
    p=\begin{pmatrix}
      e&0\\0&1_n-f
    \end{pmatrix},\qquad
    q=\begin{pmatrix}
      0&0\\0&1_n
    \end{pmatrix}.
  \]
  Then \(p,q\) belong to \(\mathcal{P}_{2n}(\tilde{A})\), and
  \[
    {[p]}_0-{[q]}_0={[e]}_0+{[1_n-f]}_0-{[1_n]}_0={[e]}_0-{[f]}_0=g.
  \]
  As \(q=s(q)\) and \(K_0(\pi)(g)=0\), we deduce that
  \[
    {[s(p)]}_0-{[q]}_0={[s(p)]}_0-{[s(q)]}_0=K_0(s)(g)=(K_0(\lambda)\circ K_0(\pi))(g)=0.
  \]
  This shows that \(g={[p]}_0-{[s(p)]}_0\).
  \begin{enumerate}
    \item Let \(p,q\) in \(\mathcal{P}_\infty(\tilde{A})\) be given. Suppose that \({[p]}_0-{[s(p)]}_0={[q]}_0-{[s(q)]}_0\).
      Then \({[p\oplus s(q)]}_0={[q\oplus s(p)]}_0\), and hence \(p\oplus s(q)\sim_s q\oplus s(p)\) in \(\mathcal{P}_\infty(\tilde{A})\) by
      Proposition \ref{prop:617}. By the remark below Definition \ref{def:Lars616}, there is a natural number \(n\) such that 
      \(p\oplus s(q)\oplus1_n\sim_0 q\oplus s(p)\oplus 1_n\). This shows that (a) implies (c). To see that (c) implies (b) note that
      if \(r_1,r_2\) are scalar projections in \(\mathcal{P}_\infty(\tilde{A})\) of dimension \(k\) and \(l\), respectively, then \(r_1\sim_0 1_k\) and \(r_2\sim_0 1_l\)
      (see Exercise \ref{ex:Lars29}), and hence \(p\oplus 1_k\sim_0 q\oplus1_l\).\par
      To see that (b) implies (a) note first that
      \[
        {[p\oplus 1_k]}_0-{[s(p\oplus1_k)]}_0={[p]}_0+{[1_k]}_0-{[s(p)]}_0-{[1_k]}_0={[p]}_0-{[s(p)]}_0.
      \]
      It therefore suffices to show that \({[p]}_0-{[s(p)]}_0={[q]}_0-{[s(q)]}_0\) when \(p\sim_0 q\).
      Suppose accordingly that \(p=v^*v\) and \(q=vv^*\) for some partial isometry \(v\) in \(M_{m,n}(\tilde{A})\). Let \(s(v)\) in \(M_{m,n}(\mathbb{C})\), viewed as a subset of \(M_{m,n}(\tilde{A})\), 
      be a matrix obtained by applying \(s\) to each entry of \(v\) (in analogy with the definition of \(s_n\)). Then 
      \(s(v)^*s(v)=s(p)\) and \(s(v)s(v)^*=s(q)\), and so \(s(p)\sim_0 s(q)\).
      Hence \({[p]}_0={[q]}_0\) and \({[s(p)]}_0={[s(q)]}_0\), and this proves that (a) holds. 
    \item If \({[p]}_0-{[s(p)]}_0=0\), then \(p\sim_s s(p)\) by Proposition \ref{prop:617} (5), and so
      there exists \(m\) in \(\mathbb{N}\) such that \(p\oplus1_m\sim s(p)\oplus1_m\), see the argument below Definition \ref{def:Lars616}.
    \item By definition, we have
      \begin{align*}
        K_0(\varphi)({[p]}_0-{[s(p)]}_0)&=K_0(\tilde{\varphi})({[p]}_0-{[s(p)]}_0)\\
        &={[\tilde{\varphi}(p)]}_0-{[\tilde{\varphi}(s(p))]}_0={[\tilde{\varphi}(p)]}_0-{[s(\tilde{\varphi}(p))]}_0.
      \end{align*}
  \end{enumerate}
\end{proof}
\begin{lemma}\label{lemma:Lars723}
  Let \(A\) and \(B\) be \cstars, and let \(\varphi:A\to B\) be a *-homomorphism.
  Suppose that \(g\) in \(K_0(A)\) belongs to the kernel of \(K_0(\varphi)\).
  \begin{enumerate}
    \item There exist a natural number \(n\), a projection \(p\) in \(\mathcal{P}_n(\tilde{A})\) and a unitary
      \(u\) in \(M_n(\tilde{B})\) such that \(g={[p]}_0-{[s(p)]}_0\) and \(u\tilde{\varphi}(p)u^*=s(\tilde{\varphi}(p))\).
    \item If \(\varphi\) is surjective, then there is a projection \(p\) in \(\mathcal{P}_\infty(\tilde{A})\) satisfying 
      \(g={[p]}_0-{[s(p)]}_0\) and \(\tilde{\varphi}(p)=s(\tilde{\varphi}(p))\).
  \end{enumerate}
\end{lemma}
\begin{proof}
  \begin{enumerate}
    \item It follows from the standard picture of \(K_0\) (Proposition \ref{prop:Lars722}) that
      there exist a natural number \(k\) and a projection \(p_1\) in \(\mathcal{P}_k(\tilde{A})\) with
      \(g={[p_1]}_0-{[s(p_1)]}_0\) and \({[\tilde{\varphi}(p_1)]}_0-{[s(\tilde{\varphi}(p_1))]}_0=0\).
      By Proposition \ref{prop:Lars722} (2), \(\tilde{\varphi}(p_1)\oplus1_m\sim s(\tilde{\varphi}(p_1))\oplus1_m\)
      for some \(m\) in \(\mathbb{N}\). Put \(p_2=p_1\oplus 1_m\). Then \(p_2\) belongs to \(\mathcal{P}_{k+m}(\tilde{A})\),
      \(g={[p_2]}_0-{[s(p_2)]}_0\), and
      \[
        \tilde{\varphi}(p_2)=\tilde{\varphi}(p_1)\oplus1_m\sim s(\tilde{\varphi}(p_1))\oplus1_m=s(\tilde{\varphi}(p_2)).
      \]
      Put \(n=2(k+m)\), let \(0_{k+m}\) be the zero projection in
      \(M_{k+m}(\tilde A)\), and set
      \[
        p=p_2\oplus0_{k+m}.
      \]
      Then \(p\in\mathcal P_n(\tilde A)\),
      \[
        g=[p]_0-[s(p)]_0,
      \]
      and Proposition~\ref{prop:528} yields a unitary
      \(u\in M_n(\tilde B)\) such that
      \[
        u\tilde\varphi(p)u^*=s(\tilde\varphi(p)).
      \]
    \item By (1) we can find \(n\) in \(\mathbb{N}\), a projection \(p_1\) in \(M_n(\tilde{A})\), and a unitary \(u\) in \(M_n(\tilde{B})\)
      with \(g={[p_1]}_0-{[s(p_1)]}_0\) and \(u\tilde{\varphi}(p_1)u^*=s(\tilde{\varphi}(p_1))\). Use Lemma \ref{lemma:518}
      to find a unitary \(v\) in \(M_{2n}(\tilde{A})\) with \(\tilde{\varphi}(v)=\diag(u,u^*)\). Put
      \(p=v\diag(p_1,0_n)v^*\). Then \(p\) is a projection in \(M_{2n}(\tilde{A})\), and
      \[
        \tilde{\varphi}(p)=\begin{pmatrix}
          u&0\\0&u^*
        \end{pmatrix}
        \begin{pmatrix}
          \tilde{\varphi}(p_1)&0\\0&0
        \end{pmatrix}
        \begin{pmatrix}
          u^*&0\\0&u
        \end{pmatrix}
        =\begin{pmatrix}
          s(\tilde{\varphi}(p_1))&0\\
          0&0
        \end{pmatrix}\in M_{2n}(\mathbb{C}).
      \]
      It follows that \(s(\tilde{\varphi}(p))=\tilde{\varphi}(p)\). Finally, \(g={[p]}_0-{[s(p)]}_0\)
      because \(p\sim_0 p_1\). See Proposition \ref{prop:Lars722} (1).
  \end{enumerate}
\end{proof}

\section{Half and split exactness and stability of \texorpdfstring{\(K_0\)}{K0}}
We begin with a description of what happens when units are adjoined to a short exact sequence.
The straightforward proof is left to the reader.
\begin{lemma}\label{lemma:Lars431}
  Let
  \[
    \begin{tikzcd}
      0\arrow[r]&I\arrow[r,"\varphi"]&A\arrow[r,"\psi"]&B\arrow[r]&0
    \end{tikzcd}
  \]
  be a short exact sequence of \cstars, and let \(n\) be a natural number.
  \begin{enumerate}
    \item The mapping \(\tilde{\varphi}_n:M_n(\tilde{I})\to M_n(\tilde{A})\) is injective.
    \item An element \(a\) in \(M_n(\tilde{A})\) belongs to the image of \(\tilde{\varphi}_n\) if and only if \(\tilde{\psi}_n(a)=s_n(\tilde{\psi}_n(a))\).
  \end{enumerate}
\end{lemma}
\begin{proposition}[Half exactness of \(K_0\)]\label{prop:Lars432}
  Every short exact sequence of \cstars
  \[
    \begin{tikzcd}
      0\arrow[r]&I\arrow[r,"\varphi"]&A\arrow[r,"\psi"]&B\arrow[r]&0
    \end{tikzcd}
  \]
  induces an exact sequence of abelian groups
  \[
    \begin{tikzcd}
      K_0(I)\arrow[r,"K_0(\varphi)"]&K_0(A)\arrow[r,"K_0(\psi)"]&K_0(B),
    \end{tikzcd}
  \]
  that is, \(\operatorname{Im}(K_0(\varphi))=\operatorname{Ker}(K_0(\psi))\).
\end{proposition}
\begin{proof}
  Functoriality of \(K_0\) yields that
  \[
    K_0(\psi)\circ K_0(\varphi)=K_0(\psi\circ\varphi)=K_0(0)=0.
  \]
  Hence \(\operatorname{Im}(K_0(\varphi))\) is contained in \(\operatorname{Ker}(K_0(\psi))\).\par
  Conversely, assume that \(g\) belongs to the kernel of \(K_0(\psi)\). According to Lemma \ref{lemma:Lars723} (2) we can
  find a natural number \(n\) and a projection \(p\) in \(\mathcal{P}_n(\tilde{A})\) with
  \(g={[p]}_0-{[s(p)]}_0\) and \(\tilde{\psi}(p)=s(\tilde{\psi}(p))\). By Lemma \ref{lemma:Lars431} (2), there is 
  an element \(e\) in \(M_n(\tilde{I})\) with \(\tilde{\varphi}(e)=p\). We also know from Lemma \ref{lemma:Lars431} that \(\tilde{\varphi}\)
  is injective. It follows that \(e\) must be a projection. Hence,
  \[
    g={[\tilde{\varphi}(e)]}_0-{[s(\tilde{\varphi}(e))]}_0=K_0(\varphi)({[e]}_0-{[s(e)]}_0)\in \operatorname{Im}(K_0(\varphi)).
  \]
\end{proof}
\begin{proposition}[Split exactness of \(K_0\)]\label{prop:Lars433}
  Every split exact sequence of \cstars
  \begin{equation*}
    \begin{tikzcd}
      0 \arrow[r]&I\arrow[r,"\varphi"] & A\arrow[r, shift left=2pt, "\psi"]
      &B\arrow[l, shift left=2pt, "\lambda"]\arrow[r]&0
    \end{tikzcd}
  \end{equation*}  
  induces a split exact sequence of abelian groups
  \begin{equation}\label{diag:splitExact46}
    \begin{tikzcd}
      0 \arrow[r]&K_0(I)\arrow[r,"K_0(\varphi)"] & K_0(A)\arrow[r, shift left=2pt, "K_0(\psi)"]
      &K_0(B)\arrow[l, shift left=2pt, "K_0(\lambda)"]\arrow[r]&0.
    \end{tikzcd}
  \end{equation}
\end{proposition}
\begin{proof}
  The sequence (\ref{diag:splitExact46}) is exact at \(K_0(A)\) by Proposition \ref{prop:Lars432}. From 
  functoriality of \(K_0\) we get
  \[
    \operatorname{id}_{K_0(B)}=K_0(\operatorname{id}_B)=K_0(\psi)\circ K_0(\lambda).
  \]
  Hence (\ref{diag:splitExact46}) is exact at \(K_0(B)\). We show that \(K_0(\varphi)\) is injective, and this will imply exactness of (\ref{diag:splitExact46}) at \(K_0(I)\).\par
  Let \(g\) in the kernel of \(K_0(\varphi)\) be given. By Lemma \ref{lemma:Lars723} (1) we can find a natural number \(n\),
  a projection \(p\) in \(\mathcal{P}_n(\tilde{I})\) and a unitary \(u\) in \(M_n(\tilde{A})\) with 
  \(g={[p]}_0-{[s(p)]}_0\) and \(u\tilde{\varphi}(p)u^*=s(\tilde{\varphi}(p))\).
  Put \(v=(\tilde{\lambda}\circ\tilde{\psi})(u^*)u\). Then \(v\) is a unitary element in \(M_n(\tilde{A})\) and
  \(\tilde{\psi}(v)=1\). By Lemma \ref{lemma:Lars431} (2) there is an element \(w\) in \(M_n(\tilde{I})\) with
  \(\tilde{\varphi}(w)=v\). Since \(\tilde{\varphi}\) is injective, \(w\) must be unitary. From the calculation
  \begin{align*}
    \tilde{\varphi}(wpw^*)&=v\tilde{\varphi}(p)v^*=(\tilde{\lambda}\circ\tilde{\psi})(u^*)s(\tilde{\varphi}(p))(\tilde{\lambda}\circ\tilde{\psi})(u)\\
    &=(\tilde{\lambda}\circ\tilde{\psi})(u^*s(\tilde{\varphi}(p))u)=(\tilde{\lambda}\circ\tilde{\psi})(\tilde{\varphi}(p))=s(\tilde{\varphi}(p))=\tilde{\varphi}(s(p))
  \end{align*}
  and by injectivity of \(\tilde{\varphi}\) we conclude that \(wpw^*=s(p)\). This shows that \(p\sim_u s(p)\) in \(M_n(\tilde{I})\), and hence that \(g=0\).
\end{proof}
\begin{proposition}
  For every pair of \cstars\ \(A\) and \(B\), we have
  \[
    K_0(A\oplus B)\cong K_0(A)\oplus K_0(B).
  \]
  More specifically, if \(\iota_A:A\to A\oplus B\) and \(\iota_B:B\to A\oplus B\) are the canonical
  inclusion mappings, then the map
  \[
    K_0(\iota_A)\oplus K_0(\iota_B):K_0(A)\oplus K_0(B)\to K_0(A\oplus B),
  \]
  which maps \(g,h\) in \(K_0(A)\oplus K_0(B)\) to \(K_0(\iota_A)(g)+K_0(\iota_B)(h)\), is an isomorphism.
\end{proposition}
\begin{proof}
  Consider the diagram
  \[
    \begin{tikzcd}
      0\arrow[r]&K_0(A)\arrow[r,"\alpha"]\arrow[d,equal]&K_0(A)\oplus K_0(B)\arrow[r,"\beta"]\arrow[d,"K_0(\iota_A)\oplus K_0(\iota_B)"]&K_0(B)\arrow[r]\arrow[d,equal]&0\\
      0\arrow[r]&K_0(A)\arrow[r,"K_0(\iota_A)",swap]&K_0(A\oplus B)\arrow[r,"K_0(\pi_B)",swap]&K_0(B)\arrow[r]&0
    \end{tikzcd}
  \]
  where \(\alpha(g)=(g,0)\), \(\beta(g,h)=h\) and where \(\pi_B(a,b)=b\). The rows in the diagram are exact
  (the lower one by Proposition \ref{prop:Lars433}), and the diagram is commutative because \(\pi_B\circ\iota_A=0\)
  and \(\pi_B\circ\iota_B=\operatorname{id}_B\). An easy diagram chase (or the five lemma) shows that \(K_0(\iota_A)\oplus K_0(\iota_B)\)
  is an isomorphism.
\end{proof}
\begin{example}
  From split exactness of \(K_0\) and Example \ref{ex:Lars332} we get for every \cstar\ \(A\) that
  \[
    K_0(\tilde{A})\cong K_0(A)\oplus \mathbb{Z}.
  \]
\end{example}
\begin{proposition}[Stability of \(K_0\)]
  Let \(A\) be a \cstar\ and let \(n\) be a natural number. Then \(K_0(A)\) is isomorphic to \(K_0(M_n(A))\).\par
  More specifically, the *-homomorphism
  \[
    \lambda_{n,A}:A\to M_n(A),\qquad a\mapsto\diag(a,0_{n-1}),
  \]
  induces an isomorphism \(K_0(\lambda_{n,A}):K_0(A)\to K_0(M_n(A))\).
\end{proposition}
\begin{proof}
  We show first that the non-unital case can be derived from the unital case.
  For this purpose let \(A\) be a (non-unital) \cstar. Then
  \[
    \begin{tikzcd}
      0\arrow[r]&A\arrow[r]\arrow[d,"\lambda_{n,A}"]&\tilde{A}\arrow[r,"\pi"]\arrow[d,"\lambda_{n,\tilde{A}}"]&\mathbb{C}\arrow[r]\arrow[d,"\lambda_{n,\mathbb{C}}"]&0\\
      0\arrow[r]&M_n(A)\arrow[r]&M_n(\tilde{A})\arrow[r]&M_n(\mathbb{C})\arrow[r]&0
    \end{tikzcd}
  \]
  is a commutative diagram with split exact rows. It follows that
  \[
    \begin{tikzcd}
      0\arrow[r]&K_0(A)\arrow[r]\arrow[d,"K_0(\lambda_{n,A})"]&K_0(\tilde{A})\arrow[r]\arrow[d,"K_0(\lambda_{n,\tilde{A}})"]&K_0(\mathbb{C})\arrow[r]\arrow[d,"K_0(\lambda_{n,\mathbb{C}})"]&0\\
      0\arrow[r]&K_0(M_n(A))\arrow[r]&K_0(M_n(\tilde{A}))\arrow[r]&K_0(M_n(\mathbb{C}))\arrow[r]&0
    \end{tikzcd}
  \]
  is a commutative diagram with (split) exact rows. A diagram chase (or the five lemma) now shows that 
  \(K_0(\lambda_{n,A})\) is an isomorphism if \(K_0(\lambda_{n,\tilde{A}})\) and \(K_0(\lambda_{n,\mathbb{C}})\)
  are isomorphisms. We therefore need only prove the proposition when \(A\) is unital.\par
  The short version of rest of the proof is to remark that the two sets of projections 
  \(\mathcal{P}_\infty(A)\) and \(\mathcal{P}_\infty(M_n(A))\) for all purposes are identical.
  We formalize this idea by constructing an inverse to the map \(K_0(\lambda_{n,A})\) as follows.
  For each natural number \(k\) let \(\gamma_{n,k}:M_k(M_n(A))\to M_{kn}(A)\) be the *-isomorphism
  given by viewing each element of \(M_k(M_n(A))\) as one big matrix (in \(M_{kn}(A)\)) (in other words, \(\gamma_{n,k}\) erases certain parentheses).
  Define \(\gamma_n:\mathcal{P}_\infty(M_n(A))\to K_0(A)\) by \(\gamma_n(p)={[\gamma_{n,k}(p)]}_0\)
  for \(p\) in \(\mathcal{P}_k(M_n(A))\). Apply the universal property of \(K_0\) (Proposition \ref{prop:618}) to \(\gamma_n\)
  (and verify conditions (1), (2) and (3) in Proposition \ref{prop:618}) to obtain a group homomorphism
  \(\alpha:K_0(M_n(A))\to K_0(A)\) satisfying \(\alpha({[p]}_0)={[\gamma_{n,k}(p)]}_0\) for \(p\)
  in \(\mathcal{P}_k(M_n(A))\).\par
  We claim that \(\alpha\) is the inverse to \(K_0(\lambda_{n,A})\), and to prove this claim it suffices to show that
  \begin{align}\label{align:47}
    (\lambda_{n,A})_{kn}(\gamma_{n,k}(p))\sim_0 p&\text{ in }\mathcal{P}_\infty(M_n(A)),\qquad p\in\mathcal{P}_k(M_n(A)),\\
    \gamma_{n,k}((\lambda_{n,A})_k(p))\sim_0 p&\text{ in }\mathcal{P}_\infty(A),\qquad p\in\mathcal{P}_k(A),
  \end{align}
  where \((\lambda_{n,A})_m\) is the *-homomorphism \(M_m(A)\to M_m(M_n(A))\) induced by
  \(\lambda_{n,A}\). We prove the second claim in (\ref{align:47}); the proof of the first one is similar.\par
    Let \(\{e_1,e_2,\ldots, e_{kn}\}\) be the standard basis for \(\mathbb{C}^{kn}\), and let \(u\) be a permutation
  unitary in \(M_{kn}(\mathbb{C})\subseteq M_{kn}(A)\) that fulfills
  \[
    ue_i=e_{n(i-1)+1},\qquad i=1,2,\ldots,k.
  \]
    Then \[p\sim_0 p\oplus 0_{(n-1)k}=u^*\gamma_{n,k}((\lambda_{n,A})_k(p))u\]
  for all projections \(p\) in \(\mathcal{P}_k(A)\).
\end{proof}

\chapter{Inductive Limits of \texorpdfstring{\(C^*\)}{C*}-Algebras}
An inductive limit of a sequence of \cstars\ is an analogue of a limit of a sequence
of points in a topological space. More specifically, an inductive limit of an increasing sequence
of sub-\cstars\ of a given \cstar\ is the closure of the union of the \cstars\ in  the sequence. Forming 
inductive limits is a way of making new \cstars\ from old ones.\par
We define abstractly inductive limits in arbitrary categories, and it is 
shown that inductive limits exist in the category of \cstars\ and in the category of (ordered) abelian groups.
The main result of this chapter says that \(K_0\) is continuous in the sense that forming
inductive limits commutes with the functor \(K_0\).

\section{Products and sums of \texorpdfstring{\cstars}{C*-algebras}}
To a family \(\{A_{i}\}_{i\in I}\) of \cstars\ we associate two new \cstars\, the product 
\[
  \prod_{i\in I}A_i
\]
and the sum
\[
  \sum_{i\in I}A_i,
\]
as follows. Let \(\prod_{i\in I}A_i\) be the set of all functions
\(a:I\to \bigcup_{i\in I}A_i\) for which \(a(i)\) belongs to \(A_i\) for 
all \(i\) in \(I\) and where
\[
  \|a\|=\sup\{\|a(i)\|_{A_i}\mid i\in I\}
\]
is finite. We shall hereafter omit the subscript \(A_i\) and write \(\|a(i)\|\) for the
norm of \(a(i)\) in \(A_i\). We shall often write \((a_i)_{i\in I}\), or more briefly \((a_i)\),
for the element \(a\) in \(\prod_{i\in I}A_i\) with \(a(i)=a_i\).\par
Equip the set \(\prod A_i\) with the operations addition, multiplication, scalar multiplication and involution
by applying these operations coordinate-wise.
\begin{proposition}
  The product \(\prod_{i\in I}A_i\) is a \cstar.
\end{proposition}
\begin{proof}
  It is easy to check that \(\prod_{i\in I}A_i\) is a *-algebra. We show below that the 
  triangle inequality holds for the norm, that \(\|a^*a\|=\|a\|^2\) and that \(\prod_{i\in I}A_i\) is complete.\par
  Let \(a=(a_i)_{i\in I}\) and \(b=(b_i)_{i\in I}\) in \(\prod_{i\in I}A_i\) be given. Observe that \(\|a_i\|\leq\|a\|\)
  for all \(i\in I\). Now,
  \[
    \|(a+b)_i\|=\|a_i+b_i\|\leq\|a_i\|+\|b_i\|\leq\|a\|+\|b\|
  \]  
  for all \(i\) in \(I\), and therefore \(\|a+b\|\leq\|a\|+\|b\|\).\par
  Since \(\|(a^*a)_i\|=\|a_i^*a_i\|=\|a_i\|^2\) for every \(i\) in \(I\), we conclude that
  \(\|a^*a\|=\|a\|^2\).\par
  Let \(\{a^{(n)}\}_{n=1}^\infty\) be a Cauchy sequence in \(\prod A_i\). Then \(\{a_i^{(n)}\}_{n=1}^\infty\) is a Cauchy
  sequence in \(A_i\) with a limit \(a_i\) in \(A_i\) for each fixed \(i\) in \(I\). Put \(a=(a_i)_{i\in I}\).
  Let \(\epsilon>0\) be given, and find \(n_0\) in \(\mathbb{N}\) such that \(\|a^{(n)}-a^{(m)}\|\leq\epsilon\) whenever
  \(n,m\geq n_0\). Then for each \(n\geq n_0\) and each \(i\) in \(I\),
  \begin{equation}\label{eq:Lars61}
    \|a_i^{(n)}-a_i\|=\lim_{m\to\infty}\|a_i^{(n)}-a_i^{(m)}\|\leq\epsilon.
  \end{equation}
  Hence 
  \[
    \|a_i\|\leq\|a_i^{(n_0)}\|+\epsilon\leq\|a^{(n_0)}\|+\epsilon
  \]
  for all \(i\) in \(I\), and this shows that \(a\) belongs to \(\prod A_i\). It follows
  from (\ref{eq:Lars61}) that \(\|a^{(n)}-a\|\leq\epsilon\) for all \(n\geq n_0\), proving that \(a^{(n)}\to a\).
  Therefore \(\prod A_i\) is complete.
\end{proof}
Let \(\sum_{i\in I}A_i\) be the closure of the subset
\[
  \mathcal{I}=\{a\in\prod_{i\in I}A_i\mid a(i)=0\text{ for all but finitely many }i\in I\}
\]
of \(\prod A_i\). The next proposition is easily verified.
\begin{proposition}
  The set \(\mathcal{I}\) is a (not necessarily closed) two-sided ideal in \(\prod A_i\), and \(\sum A_i\) is a closed 
  two-sided ideal in \(\prod A_i\). In particular, \(\sum A_i\) is a \cstar.
\end{proposition}

Let \(\pi\) denote the quotient mapping \(\prod A_i\to \prod A_i/\sum A_i\). We have
the following useful lemma describing the situation when \(I=\mathbb{N}\).
\begin{lemma}\label{lemma:Lars613}
  Let \(\{A_n\}_{n=1}^\infty\) be a sequence of \cstars, and let \(a=(a_n)_{n\in\mathbb{N}}\) be an element
  in \(\prod_{n\in\mathbb{N}} A_n\). Then
  \[
    \|\pi(a)\|=\limsup_{n\to\infty}\|a_n\|.
  \]
  In particular, \(a\) belongs to \(\sum_{n\in\mathbb{N}}A_n\) if and only if \(\lim_{n\to\infty}\|a_n\|=0\).
\end{lemma}
\begin{proof}
  Since \(\mathcal{I}\) is dense in \(\sum A_n\), we have \(\|\pi(a)\|=\inf\{\|a-b\|\mid b\in\mathcal{I}\}\) by continuity
  of the mapping \(b\mapsto \|a-b\|\).\par
  Each \(b=(b_n)\) in \(\mathcal{I}\) has the property that \(b_n=0\) eventually, and therefore
  \[
    \|a-b\|\geq\limsup_{n\to\infty}\|a_n-b_n\|=\limsup_{n\to\infty}\|a_n\|.
  \]
  This shows that \(\|\pi(a)\|\geq\limsup_{n\to\infty}\|a_n\|\).\par
  For each natural number \(k\) let \(b^{(k)}=(b_n^{(k)})_{n=1}^\infty\) be the element in \(\mathcal{I}\) given by
  \[
    b_n^{(k)}=
    \begin{cases}
      a_n, & n\leq k,\\
      0,   & n>k.
    \end{cases}
  \]
  Then 
  \[
    \|\pi(a)\|\leq\inf_{k\in\mathbb{N}}\|a-b^{(k)}\|=\mathop{\vphantom{\sup}\inf}_{k\in\mathbb N}
    \sup_{n>k}\|a_n\|=\limsup_{n\to\infty}\|a_n\|.
  \]
\end{proof}

\section{Inductive limits}
\begin{definition}
  An \textnormal{inductive sequence} (or a sequence for short) in a category \(\mathcal{C}\) is a sequence
  \(\{A_n\}_{n=1}^\infty\) of objects in \(\mathcal{C}\) and a sequence \(\varphi_n:A_n\to A_{n+1}\) of morphisms
  in \(\mathcal{C}\), usually written
  \[
    \begin{tikzcd}
      A_1\arrow[r,"\varphi_1"]&A_2\arrow[r,"\varphi_2"]&A_3\arrow[r,"\varphi_3"]&\cdots.
    \end{tikzcd}
  \]  
  For \(m>n\) we shall also consider the composed morphisms
  \[
    \varphi_{m,n}=\varphi_{m-1}\circ\varphi_{m-2}\circ\cdots\circ\varphi_n:A_n\to A_m,
  \]
  which, together with the morphism \(\varphi_n\), are called the \textnormal{connecting morphisms} (or connecting maps).
  It is sometimes convenient to consider the connecting morphisms \(\varphi_{m,n}\) when \(m\leq n\).
  They are defined by \(\varphi_{n,n}=\operatorname{id}_{A_n}\) and 
  \(\varphi_{m,n}=0\) when \(m< n\). (The latter are only defined in categories with zero object.)
\end{definition}
\begin{definition}\label{def:Lars622}
  An \textnormal{inductive limit} of the inductive sequence
  \begin{equation}\label{seq:Lars62}
    \begin{tikzcd}
      A_1\arrow[r,"\varphi_1"]&A_2\arrow[r,"\varphi_2"]&A_3\arrow[r,"\varphi_3"]&\cdots.
    \end{tikzcd}
  \end{equation}
  in a category \(\mathcal{C}\) is a system \((A,\{\mu_n\}_{n=1}^\infty)\), where \(A\) is an object in \(\mathcal{C}\),
  where \(\mu_n:A_n\to A\) is a morphism in \(\mathcal{C}\) for each \(n\) in \(\mathbb{N}\), and where the following two conditions hold.
  \begin{enumerate}
    \item The diagram
      \begin{equation}\label{diag:Lars63}
        \begin{tikzcd}[column sep=small]
          A_n
            \arrow[rr, "\varphi_n"]
            \arrow[dr, "\mu_n"']
          &&
          A_{n+1}
            \arrow[dl, "\mu_{n+1}"]
          \\
          & A &
        \end{tikzcd}
      \end{equation}
      commutes for each \(n\) in \(\mathbb{N}\).
    \item If \((B,\{\lambda_n\}_{n=1}^\infty)\) is a system, where \(B\) is an object in \(\mathcal{C}\),
      \(\lambda_n:A_n\to B\) is a morphism in \(\mathcal{C}\) for each \(n\) in \(\mathbb{N}\), and where \(\lambda_n=\lambda_{n+1}\circ\varphi_n\)
      for all \(n\) in \(\mathbb{N}\), then there is one and only one morphism \(\lambda:A\to B\) making the diagram
      \begin{equation}\label{diag:Lars64}
        \begin{tikzcd}[column sep=small]
          & A_n
            \arrow[dl, "\mu_{n}",swap]
            \arrow[dr, "\lambda_n"] &\\
          A\arrow[rr, "\lambda",swap]
          &&
          B
        \end{tikzcd}
      \end{equation}
      commutative for each \(n\) in \(\mathbb{N}\).
  \end{enumerate}
\end{definition}
Inductive limits do not exist in all categories.\par
Inductive limits, when they exist, are essentially unique in the sense that
if \(\{A,\{\mu_n\}\}\) and \((B,\{\lambda_n\})\) are both inductive limits of the sequence (\ref{seq:Lars62}),
then there is one (and only one) \textit{isomorphism} \(\lambda:A\to B\) making the diagram (\ref{diag:Lars64})
above commutative.\par
Indeed, by Definition~\ref{def:Lars622}(2) there are unique morphisms \(\lambda:A\to B\) and \(\mu:B\to A\)
making the left-hand diagram in
\[
  \begin{tikzcd}[column sep=small]
    & A_n
      \arrow[dl, "\mu_{n}",swap]
      \arrow[dr, "\mu_n"]\arrow[d,"\lambda_n"] &\\
    A\arrow[r, "\lambda",swap]&B\arrow[r,"\mu",swap]&A
  \end{tikzcd}\qquad
  \begin{tikzcd}[column sep=small]
    & A_n
      \arrow[dl, "\mu_{n}",swap]
      \arrow[dr, "\mu_n"] &\\
    A\arrow[rr, "\operatorname{id}_A",swap]
    &&
    A
  \end{tikzcd}
\] 
commutative. By the uniqueness clause in Definition~\ref{def:Lars622}(2) applied to the two 
diagrams above we find that \(\mu\circ\lambda\) must be the identity map on \(A\).
Similarly, \(\lambda\circ\mu=\operatorname{id}_B\), and this shows that \(\lambda\) (and \(\mu\))
are isomorphisms.\par
Because of the essential uniqueness of inductive limits (when they exist),
we shall refer to \textit{the} inductive limit (rather than \textit{an} inductive limit).
\begin{definition}
  The inductive limit of the sequence (\ref{seq:Lars62}) is denoted by
  \[
    \varinjlim(A_n,\varphi_n),
  \]
  or more briefly by
  \[
    \varinjlim A_n.
  \]
  We shall also write
  \[
    \begin{tikzcd}
      A_1\arrow[r,"\varphi_1"]&A_2\arrow[r,"\varphi_2"]&A_3\arrow[r,"\varphi_3"]&\cdots\arrow[r]&A
    \end{tikzcd}
  \]
  to indicate that \(A\) is the inductive limit of the sequence (\ref{seq:Lars62}).
\end{definition}
\begin{example}
  The following specific examples may illustrate what inductive limits look like.
  \begin{enumerate}
    \item Let \(D\) be a \cstar\ and let \(\{A_n\}_{n=1}^\infty\) be an increasing sequence of 
      sub-\(C^*\)-algebras of \(D\). Put
      \[
        A=\overline{\bigcup_{n=1}^\infty A_n},
      \]
      and for each \(n\) let \(\iota_n:A_n\to A\) be the inclusion map. Then \((A,\{\iota_n\})\)
      is the inductive limit of the sequence \(A_1\to A_2\to A_3\to\cdots\) where the connecting maps are the inclusion maps.
    \item Consider the sequence
      \[
        \begin{tikzcd}
          \mathbb{C}\arrow[r,"\varphi_1"]&M_2(\mathbb{C})\arrow[r,"\varphi_2"]&M_3(\mathbb{C})\arrow[r,"\varphi_3"]&\cdots,
        \end{tikzcd}
      \]
      where \(\varphi_n\) maps an \(n\times n\) matrix into the upper-left corner of an
      \((n+1)\times(n+1)\) matrix whose last row and last column are zero. The inductive limit of this
      sequence is isomorphic to \(\mathcal{K}\), the \cstar\ of compact operators on a separable infinite dimensional Hilbert space.
    \item In the category of abelian groups, the inductive limit of the sequence
      \[
        \begin{tikzcd}
          \mathbb{Z}\arrow[r,"1"]&\mathbb{Z}\arrow[r,"2"]&\mathbb{Z}\arrow[r,"3"]&\cdots,
        \end{tikzcd}
      \]
      is isomorphic to \(\mathbb{Q}\).
  \end{enumerate}
\end{example}
\begin{proposition}[Inductive limit of \cstars]\label{prop:Lars624}
  Every inductive sequence of \cstars
  \[
    \begin{tikzcd}
      A_1\arrow[r,"\varphi_1"]&A_2\arrow[r,"\varphi_2"]&A_3\arrow[r,"\varphi_3"]&\cdots
    \end{tikzcd}
  \]
  has an inductive limit \((A,\{\mu_n\})\). In addition, the following hold.
  \begin{enumerate}
    \item \(A=\overline{\bigcup_{n=1}^\infty\mu_n(A_n)}\);
    \item \(\|\mu_n(a)\|=\lim_{m\to\infty}\|\varphi_{m,n}(a)\|\) for all \(n\) in \(\mathbb{N}\) and \(a\) in \(A_n\);
    \item \(\operatorname{Ker}(\mu_n)=\{a\in A_n\mid\lim_{m\to\infty}\|\varphi_{m,n}(a)\|=0\}\);
    \item if \((B,\{\lambda_n\})\) and \(\lambda:A\to B\) are as in definition \ref{def:Lars622} (2), then
      \begin{enumerate}
        \item \(\operatorname{Ker}(\mu_n)\subseteq\operatorname{Ker}(\lambda_n)\) for all \(n\) in \(\mathbb{N}\),
        \item \(\lambda\) is injective if and only if \(\operatorname{Ker}(\lambda_n)\subseteq\operatorname{Ker}(\mu_n)\) for all \(n\) in \(\mathbb{N}\),
        \item \(\lambda\) is surjective if and only if \(B=\overline{\bigcup_{n=1}^\infty \lambda_n(A_n)}\).
      \end{enumerate}
  \end{enumerate}
\end{proposition}
\begin{proof}
  Let 
  \[
    \pi:\prod_{n=1}^{\infty}A_n\to\prod_{n=1}^{\infty}A_n \Big/\sum_{n=1}^\infty A_n
  \]
  be the quotient mapping and let \(\varphi_{m,n}:A_n\to A_m\) be the connecting maps
  associated with the given sequence of \cstars. For each \(a\) in \(A_n\), set \(\nu_n(a)=(\varphi_{m,n}(a))_{m=1}^\infty\)
  (recalling that \(\varphi_{n,n}(a)=a\) and that \(\varphi_{m,n}(a)=0\) when \(m< n\)), and set \(\mu_n=\pi\circ\nu_n\). Then
  \[
    \nu_n:A_n\to\prod_{m=1}^\infty A_m,\qquad\mu_n:A_n\to\prod_{m=1}^\infty A_m\Big/\sum_{m=1}^\infty A_m
  \]
  are *-homomorphisms.\par
  Let \(a\) be an element in \(A_n\). Then \(\nu_n(a)-(\nu_{n+1}\circ\varphi_n)(a)\) is equal to the
  sequence \((c_m)_{m=1}^\infty\), where \(c_n=a\) and \(c_m=0\) when \(m\neq n\), and this sequence belongs to 
  \(\sum_{m=1}^\infty A_m\). Hence
  \[
    \mu_n(a)-(\mu_{n+1}\circ\varphi_n)(a)
    =\pi\bigl(\nu_n(a)-(\nu_{n+1}\circ\varphi_n)(a)\bigr)=0,
  \]
  so that \(\mu_{n+1}\circ\varphi_n=\mu_n\). It follows from this that \(\{\mu_n(A_n)\}_{n=1}^\infty\)
  is an increasing sequence of \cstars. Hence
  \[
    A=\overline{\bigcup_{n=1}^\infty\mu_n(A_n)}
  \]
  is a \cstar, (the co-restriction of) each \(\mu_n\) is a *-homomorphism from \(A_n\) to \(A\), and the system \(\{\mu_n\}_{n=1}^\infty\) satisfies Definition \ref{def:Lars622} (1).\par
  Part (1) of the present Proposition holds by the construction of \(A\). We now prove (2). Let \(a\) be an element in \(A_n\).
  Then \(\mu_n(a)=\pi(\nu_n(a))\), where \(\nu_n(a)\) is the sequence \(\{\varphi_{m,n}(a)\}_{m=1}^\infty\).
  Use the norm result in Lemma \ref{lemma:Lars613} to obtain 
  \[
    \|\mu_n(a)\|=\|\pi(\nu_n(a))\|=\limsup_{m\to\infty}\|\varphi_{m,n}(a)\|=\lim_{m\to\infty}\|\varphi_{m,n}(a)\|;
  \]
  the limit exists because the sequence \(\{\|\varphi_{m,n}(a)\|\}_{m=n}^\infty\) is decreasing. It is clear that (3) follows from (2).\par
  We proceed to prove (4)(a) and that Definition \ref{def:Lars622} (2) holds.
  Assume that \((B,\{\lambda_n\})\) is a system as in Definition \ref{def:Lars622} (2). Then 
  \(\lambda_m\circ\varphi_{m,n}=\lambda_n\) for each \(m>n\), whence \(\|\lambda_n(a)\|\leq\|\varphi_{m,n}(a)\|\), and so
  \[
    \|\lambda_n(a)\|\leq\limsup_{m\to\infty}\|\varphi_{m,n}(a)\|=\|\mu_n(a)\|.
  \]
  It follows that the kernel of \(\lambda_n\) contains the kernel of \(\mu_n\). By the first isomorphism theorem
  there is a unique *-homomorphism \(\lambda_n':\mu_n(A_n)\to B\) such that \(\lambda_n'\circ\mu_n=\lambda_n\).
  By its uniqueness, \(\lambda_{n+1}'\) necessarily extends \(\lambda_n'\). Hence we obtain 
  a *-homomorphism
  \[
    \lambda':\bigcup_{n=1}^\infty\mu_n(A_n)\to B
  \]
  that extends each \(\lambda_n'\). The map \(\lambda'\) is a \textit{contraction}, i.e., its operator norm is at most
  \(1\), because each \(\lambda'_n\) is a contraction. We can therefore extend \(\lambda'\) by
  uniform continuity to a *-homomorphism \(\lambda:A\to B\).
  The restriction of \(\lambda\) to \(\mu_n(A_n)\) is equal to \(\lambda_n'\), and hence \(\lambda\circ\mu_n=\lambda_n'\circ\mu_n=\lambda_n\).
  That \(\lambda\) is unique with respect to this property follows from the fact that
  \(\bigcup_{n=1}^\infty\mu_n(A_n)\) is dense in \(A\).\par
  (4)(b). The *-homomorphism \(\lambda\) is injective if and only if it is an isometry, and this is the case if and only if each \(\lambda'_n\)
  is an isometry, or, equivalently, if and only if each \(\lambda_n'\) is 
  injective. But \(\lambda_n'\) is injective if and only if \(\operatorname{Ker}(\lambda_n)=\operatorname{Ker}(\mu_n)\).\par
  (4)(c). This follows from the fact that the image of \(\lambda\) is \(\overline{\bigcup_{n=1}^\infty\lambda_n(A_n)}\).
\end{proof}
\begin{proposition}[Inductive limit of abelian groups]\label{prop:Lars625}
  Each sequence
  \[
    \begin{tikzcd}
      G_1\arrow[r,"\alpha_1"]&G_2\arrow[r,"\alpha_2"]&G_3\arrow[r,"\alpha_3"]&\cdots
    \end{tikzcd}
  \]
  of abelian groups has an inductive limit \((G,\{\beta_n\})\). Moreover,
  \begin{enumerate}
    \item \(G=\bigcup_{n=1}^\infty\beta_n(G_n)\),
    \item \(\operatorname{Ker}(\beta_n)=\bigcup_{m=n+1}^\infty\operatorname{Ker}(\alpha_{m,n})\) for each \(n\) in \(\mathbb{N}\),
    \item if \((H,\{\gamma_n\})\) and \(\gamma:G\to H\) are as in (2) of Definition \ref{def:Lars622}, then
      \begin{enumerate}
        \item \(\gamma\) is injective if and only if \(\operatorname{Ker}(\gamma_n)=\operatorname{Ker}(\beta_n)\) for all \(n\) in \(\mathbb{N}\),
        \item \(\gamma\) is surjective if and only if \(H=\bigcup_{n=1}^\infty\gamma_n(G_n)\).
      \end{enumerate}
  \end{enumerate}
\end{proposition}
\begin{proposition}[Inductive limit of ordered abelian groups]\label{prop:Lars626}
  Let
  \begin{equation}\label{eq:Lars65}
    \begin{tikzcd}
      G_1\arrow[r,"\alpha_1"]&G_2\arrow[r,"\alpha_2"]&G_3\arrow[r,"\alpha_3"]&\cdots
    \end{tikzcd}
  \end{equation}
  be a sequence of ordered abelian groups where the connecting maps are positive
  group homomorphisms, and let \((G,\{\beta_n\})\) be the inductive limit of this sequence in the category of 
  abelian groups. Put 
  \[
    G^+=\bigcup_{n=1}^\infty\beta_n(G_n^+).
  \]
  Then \((G,G^+)\) is an ordered abelian group, \(\beta_n\) is a positive homomorphism
  for each \(n\) in \(\mathbb{N}\), and \(((G,G^+),\{\beta_n\})\) is the inductive limit of the 
  sequence (\ref{eq:Lars65}) in the category of ordered abelian groups.
\end{proposition}
\begin{proof}
  The connecting maps $\alpha_n$ are positive, i.e.,
  $\alpha_n(G_n^+)$ is contained in $G_{n+1}^+$ for each $n$, and
  therefore
  \[
    \beta_n(G_n^+)
    =\beta_{n+1}\bigl(\alpha_n(G_n^+)\bigr)
    \subseteq\beta_{n+1}(G_{n+1}^+),
  \]
  which shows that $\{\beta_n(G_n^+)\}_{n=1}^{\infty}$ is an increasing
  sequence of subsets of $G$. If $x,y$ belong to $G^+$, then they both
  belong to $\beta_n(G_n^+)$ for some $n$, and then $x+y$ is an element
  of $\beta_n(G_n^+)$ and hence of $G^+$.

  If $x$ is an element of $G^+\cap(-G^+)$, then $x$ belongs to
  $\beta_n(G_n^+)\cap-\beta_n(G_n^+)$ for some $n$, and hence
  \[
    x=\beta_n(y_1)=-\beta_n(y_2)
  \]
  for some $y_1,y_2$ in $G_n^+$. Since $\beta_n(y_1+y_2)=0$, there is an
  integer $m\geq n$ such that
  \[
    \alpha_{m,n}(y_1+y_2)=0.
  \]
  Put $z_i=\alpha_{m,n}(y_i)$. Then $z_1,z_2$ belong to $G_m^+$ and
  $z_1+z_2=0$. This entails that $z_1=z_2=0$, because
  $G_m^+\cap(-G_m^+)=0$, and hence
  \[
    x=\beta_m(z_1)=0.
  \]

  If $x$ is any element of $G$, then $x$ belongs to $\beta_n(G_n)$ for
  some $n$ by Proposition~\ref{prop:Lars625} (1). Since
  \[
    G_n=G_n^+-G_n^+,
  \]
  $x$ belongs to $\beta_n(G_n^+)-\beta_n(G_n^+)$ and hence to
  $G^+-G^+$. This shows that $(G,G^+)$ is an ordered Abelian group. By
  the definition of $G^+$, $\beta_n(G_n^+)$ is contained in $G^+$ for
  every $n$, and so each $\beta_n$ is positive.

  To show that $\bigl((G,G^+),\{\beta_n\}\bigr)$ is the inductive limit
  of the sequence~(\ref{eq:Lars65}), we must verify (1) and (2) of
  Definition~\ref{def:Lars622}. That (1) holds follows from the fact that
  $(G,\{\beta_n\})$ is the inductive limit of the sequence~(\ref{eq:Lars65}) in the
  category of Abelian groups. To show (2), let $(H,\{\gamma_n\})$ be a
  system, where $H$ is an ordered Abelian group and each
  \[
    \gamma_n\colon G_n\longrightarrow H
  \]
  is a positive group homomorphism. Using again that $G$ is the
  inductive limit in the category of Abelian groups, there is a unique
  group homomorphism $\gamma\colon G\to H$ making diagram~(\ref{diag:Lars64})
  commutative. We must show that $\gamma$ is positive:
  \[
    \gamma(G^+)
    =\gamma\left(\bigcup_{n=1}^{\infty}\beta_n(G_n^+)\right)
    =\bigcup_{n=1}^{\infty}(\gamma\circ\beta_n)(G_n^+)
    =\bigcup_{n=1}^{\infty}\gamma_n(G_n^+)
    \subseteq H^+.
  \]
\end{proof}

\section{Continuity of \texorpdfstring{\(K_0\)}{K0}}
\begin{lemma}\label{lemma:Lars631}
  Let \(A\) be a \cstar.
  \begin{enumerate}
    \item If \(a\) is a self-adjoint element in \(A\) with \(\delta=\|a-a^2\|<1/4\), then there is 
      a projection \(p\) in \(A\) with \(\|a-p\|\leq2\delta\).
    \item Let \(p,q\) be projections in \(A\). If there exists an element \(x\) in \(A\) with
      \(\|x^*x-p\|<1/2\) and \(\|xx^*-q\|<1/2\), then \(p\sim q\).
  \end{enumerate}
\end{lemma}
\begin{proof}
  \begin{enumerate}
    \item If \(t\) belongs to \(\operatorname{sp}(a)\), then \(t-t^2\) belong to \(\operatorname{sp}(a-a^2)\)
      by the spectral mapping theorem; and if \(t\) is a real number satisfying \(|t-t^2|\leq\delta<1/4\),
      then \(t\) belongs to \([-2\delta,2\delta]\cup[1-2\delta,1+2\delta]\). Hence, if 
      \(\|a-a^2\|=\delta<1/4\), then
      \[
        \operatorname{sp}(a)\subseteq\{t\in\mathbb{R}\mid|t-t^2|\leq\delta\}\subseteq[-2\delta,2\delta]\cup[1-2\delta,1+2\delta].
      \]
      It follows that we can define a continuous function \(f\) on \(\operatorname{sp}(a)\) by
      \begin{equation}\label{eq:Lars66}
        f(t)=\begin{cases}
          0,&t\leq2\delta,\\
          1,&t\geq1-2\delta.
        \end{cases}
      \end{equation}
      Put \(p=f(a)\). Then \(p\) is a projection because \(f=f^2=\overline{f}\), and 
      \(\|a-p\|\leq2\delta\) because \(|t-f(t)|\leq2\delta\) for all \(t\) in \(\operatorname{sp}(a)\).
    \item Put
      \[
        \delta=\frac{1}{2}\max\{\|x^*x-p\|,\|xx^*-q\|\}<1/4,
      \]
      and put \(\Gamma=\operatorname{sp}(x^*x)\cup\operatorname{sp}(xx^*)\). Then 
      \(\Gamma\subseteq[-2\delta,2\delta]\cup[1-2\delta,1+2\delta]\) by
      Lemma \ref{lemma:523}. Let \(f\) in \(C(\Gamma)\) be as defined in (\ref{eq:Lars66}). Then, as in (1),
      \(p_0=f(x^*x)\) and \(q_0=f(xx^*)\) are projections in \(A\) satisfying \(\|p-p_0\|\leq4\delta<1\)
      and \(\|q-q_0\|\leq4\delta<1\). Hence \(p\sim p_0\) and \(q\sim q_0\) by Proposition \ref{prop:Lars224}.\par
      Notice that \(xh(x^*x)x^*=h(xx^*)xx^*\) for every \(h\) in \(C(\Gamma)\); this is clear
      when \(h\) is a polynomial, and the polynomials are dense in \(C(\Gamma)\).
      Let \(g\) be the positive function in \(C(\Gamma)\) that satisfies \(t g(t)^2=f(t)\)
      for all \(t\) in \(\Gamma\). Put \(v=xg(x^*x)\). Then
      \begin{align*}
        v^*v&=g(x^*x)x^*xg(x^*x)=f(x^*x)=p_0,\\
        vv^*&=xg(x^*x)^2x^*=g(xx^*)^2xx^*=f(xx^*)=q_0.
      \end{align*}
      Hence \(p_0\sim q_0\), and it follows that \(p\sim q\) as desired.
  \end{enumerate}
\end{proof}
\begin{theorem}[Continuity of \(K_0\)]
  For each inductive sequence
  \begin{equation}\label{seq:Lars67}
    \begin{tikzcd}
      A_1\arrow[r,"\varphi_1"]&A_2\arrow[r,"\varphi_2"]&A_3\arrow[r,"\varphi_3"]&\cdots
    \end{tikzcd}
  \end{equation}
  of \cstars, there is an isomorphism
  \[
    K_0\left(\varinjlim A_n\right)\cong\varinjlim K_0(A_n).
  \]
  If \((K_0(A_n),K_0(A_n)^+)\) is an ordered group for each \(n\), then this is an
  isomorphism in the category of ordered abelian groups.\par
  More specifically, if \((A,\{\mu_n\})\) is the inductive limit of the sequence (\ref{seq:Lars67}),
  and if \((G_0,\{\beta_n\})\) is the inductive limit of the sequence
  \begin{equation}\label{seq:Lars68}
    \begin{tikzcd}
      K_0(A_1)\arrow[r,"K_0(\varphi_1)"]&K_0(A_2)\arrow[r,"K_0(\varphi_2)"]&K_0(A_3)\arrow[r,"K_0(\varphi_3)"]&\cdots,
    \end{tikzcd}
  \end{equation}
  then there is a unique group isomorphism \(\gamma:G_0\to K_0(A)\) that makes the diagram
  \begin{equation}\label{diag:Lars69}
    \begin{tikzcd}
      &K_0(A_n)\arrow[ld,"\beta_n",swap]\arrow[rd,"K_0(\mu_n)"]&\\
      G_0\arrow[rr,"\gamma",swap]&&K_0(A)
    \end{tikzcd}
  \end{equation}
  commutative for each \(n\) in \(\mathbb{N}\). Moreover,
  \begin{enumerate}
    \item \(K_0(A)=\bigcup_{n=1}^\infty K_0(\mu_n)(K_0(A_n))\),
    \item \(K_0(A)^+=\bigcup_{n=1}^\infty K_0(\mu_n)(K_0(A_n)^+)\),
    \item \(\operatorname{Ker}(K_0(\mu_n))=\bigcup_{m=n+1}^\infty\operatorname{Ker}(K_0(\varphi_{m,n}))\) for each \(n\) in \(\mathbb{N}\).
  \end{enumerate}
\end{theorem}
\begin{proof}
  We first prove (1), (2) and (3). To simplify the notation, let \(\mu_n\)
  denote the induced *-homomorphism \(M_k(A_n)\to M_k(A)\) for each \(k\) in \(\mathbb{N}\),
  and let \(\tilde{\mu}_n\) denote the induced *-homomorphism \(M_k(\tilde{A}_n)\to M_k(\tilde{A})\)
  where units have been adjoined to \(A_n\) and to \(A\). Using Proposition \ref{prop:Lars624} (4)
  one finds that \((M_k(A),\{\mu_n\})\) and \((M_k(\tilde{A}),\{\tilde{\mu}_n\})\) are the inductive limits of the sequences
  \[\begin{tikzcd}
      M_k(A_1)\arrow[r,"\varphi_1"]&M_k(A_2)\arrow[r,"\varphi_2"]&M_k(A_3)\arrow[r,"\varphi_3"]&\cdots,
    \end{tikzcd}\]
    \[\begin{tikzcd}
      M_k(\tilde{A}_1)\arrow[r,"\tilde{\varphi}_1"]&M_k(\tilde{A}_2)\arrow[r,"\tilde{\varphi}_2"]&M_k(\tilde{A}_3)\arrow[r,"\tilde{\varphi}_3"]&\cdots.
    \end{tikzcd}\]
  \begin{enumerate}
    \item Let \(g\) in \(K_0(A)\) be given. Use the standard picture of \(K_0(A)\) to find
      a natural number \(k\) and a projection \(p\) in \(M_k(\tilde{A})\) with \(g={[p]}_0-{[s(p)]}_0\).
      By Proposition \ref{prop:Lars624} (1) there exists a natural number \(n\) and an element \(b_n\) in \(M_k(\tilde{A}_n)\)
      such that \(\|\tilde{\mu}_n(b_n)-p\|<1/5\). Put \(a_n=(b_n+b_n^*)/2\), and put 
      \(a_m=\tilde{\varphi}_{m,n}(a_n)\) for \(m>n\). Then \(a_m\) is self-adjoint and \(\|\tilde{\mu}_m(a_m)-p\|<1/5\)
      for \(m\geq n\). Lemma \ref{lemma:523} implies that
      \[
        \operatorname{sp}(\tilde{\mu}_n(a_n))\subseteq[-1/5,1/5]\cup[4/5,6/5],
      \]
      and so 
      \[
        \|\tilde{\mu}_n(a_n^2-a_n)\|=\max\{|t^2-t|\mid t\in\operatorname{sp}(\tilde{\mu}_n(a_n))\}<1/4.
      \]
      By Proposition \ref{prop:Lars624} (2) there is \(m\geq n\) such that \(\|a_m^2-a_m\|<1/4\).
      By Lemma \ref{lemma:Lars631} (1) there is a projection \(q\) in \(M_k(\tilde{A}_m)\) with 
      \(\|a_m-q\|<1/2\). Now,
      \[
        \|\tilde{\mu}_m(q)-p\|\leq\|q-a_m\|+\|\tilde{\mu}_m(a_m)-p\|<1,
      \]
      and so \(\tilde{\mu}_m(q)\sim p\). Hence, by the standard picture for \(K_0\),
      \[
        g={[p]}_0-{[s(p)]}_0={[\tilde{\mu}_m(q)]}_0-{[s(\tilde{\mu}_m(q))]}_0=K_0(\mu_m)({[q]}_0-{[s(q)]}_0).
      \]
    \item For each \(n\) in \(\mathbb{N}\), \(K_0(\mu_n)(K_0(A_n)^+)\) is contained in \(K_0(A)^+\) because
      \(K_0(\mu_n)\) is positive. Conversely, let \(g\) in \(K_0(A)^+\) be given. Then \(g={[p]}_0\) for some 
      natural number \(k\) and some projection \(p\) in \(M_k(A)\). Arguing as in (1) we find a natural number
      \(m\) and a projection \(q\) in \(M_k(A_m)\) such that \(\|\mu_m(q)-p\|<1\).
      It follows that \(p\sim \mu_m(q)\), and hence that \(g={[\mu_m(q)]}_0=K_0(\mu_m)({[q]}_0)\);
      this proves (2) because \({[q]}_0\) belongs to \(K_0(A_m)^+\).
    \item The kernel of \(K_0(\mu_n)\) contains the kernel of \(K_0(\varphi_{m,n})\) for each \(m>n\)
      by Proposition \ref{prop:Lars624}. Conversely, let \(g\) in \(K_0(A_n)\) be an element in the kernel
      of \(K_0(\mu_n)\). By the appendix to the standard picture for \(K_0\)
      there are \(k\) in \(\mathbb{N}\) and a projection \(p\) in \(M_k(\tilde{A}_n)\) such that \(g={[p]}_0-{[s(p)]}_0\)
      and \(\tilde{\mu}_n(p)\sim\tilde{\mu}_n(s(p))\). In other words, \(\tilde{\mu}_n(p)=v^*v\) and \(\tilde{\mu}_n(s(p))=vv^*\)
      for some \(v\) in \(M_k(\tilde{A})\). By Proposition \ref{prop:Lars624} (1) there exists an integer \(l\geq n\) and an element
      \(x_l\) in \(M_k(\tilde{A}_l)\) with \(\tilde{\mu}_l(x_l)\) close enough to \(v\) to ensure that
      \[
        \|\tilde{\mu}_l(x^*_lx_l)-\tilde{\mu}_n(p)\|<1/2,\qquad\|\tilde{\mu}_l(x_lx^*_l)-\tilde{\mu}_n(s(p))\|<1/2.
      \]
      By Proposition \ref{prop:Lars624} (3) there is \(m\geq l\) such that
      \[
        \|x_m^*x_m-\tilde{\varphi}_{m,n}(p)\|<1/2,\qquad\|x_mx_m^*-\tilde{\varphi}_{m,n}(s(p))\|<1/2,
      \]
      where \(x_m=\tilde{\varphi}_{m,l}(x_l)\). Lemma \ref{lemma:Lars631} (2) now yields that
      \[
        \tilde{\varphi}_{m,n}(p)\sim\tilde{\varphi}_{m,n}(s(p))=s(\tilde{\varphi}_{m,n}(p))\text{ in }M_k(\tilde{A}_m).
      \]
      This shows that
      \[
        K_0(\varphi_{m,n})(g)={[\tilde{\varphi}_{m,n}(p)]}_0-{[s(\tilde{\varphi}_{m,n}(p))]}_0=0,
      \]
      and (3) follows.
  \end{enumerate}\par
  Since \(K_0(\mu_n)=K_0(\mu_{n+1})\circ K_0(\varphi_n)\) for every \(n\), part (2) of the definition of inductive limits
  applied to the system \((K_0(A),\{K_0(\mu_n)\})\) yields a unique group homomorphism \(\gamma:G_0\to K_0(A)\) making the diagram
  (\ref{diag:Lars69}) commutative for every \(n\). It follows from (1) that \(\gamma\) is surjective.\par
  Suppose that \(g\) belongs to the kernel of \(\gamma\). Find \(n\) in \(\mathbb{N}\) and \(h\) in \(K_0(A_n)\)
  with \(g=\beta_n(h)\), see Proposition \ref{prop:Lars625} (1). Then \(0=\gamma(g)=K_0(\mu_n)(h)\),
  and by (3) this implies that \(K_0(\varphi_{m,n})(h)=0\) for some integer \(m>n\), whence
  \(g=(\beta_m\circ K_0(\varphi_{m,n}))(h)=0\). This shows that \(\gamma\) is injective.\par
  Finally, if \((K_0(A_n),K_0(A_n)^+)\) is an ordered abelian group for each \(n\), then
  the inductive limit of the sequence (\ref{seq:Lars68}) in the category of ordered abelian groups is \((G_0,G_0^+),\{\beta_n\}\), where
  \[
    G_0^+=\bigcup_{n=1}^\infty\beta_n(K_0(A_n)^+),
  \]
  see Proposition \ref{prop:Lars626}. Now (2) shows that \(\gamma(G_0^+)=K_0(A)^+\), and so \(\gamma\) is an
  isomorphism of ordered abelian groups.
\end{proof}

\chapter{The Functor \texorpdfstring{\(K_1\)}{K1}}

\section{Definition and functoriality of \texorpdfstring{\(K_1\)}{K1}}
The construction in this section, including the standard picture and functorial properties of \(K_1\),
follows \cite[Sections~8.1--8.2]{RordamLarsenLaustsen2000}.
\begin{definition}
  Let \(A\) be a unital \cstar. We denote
  \begin{align*}
    \mathcal{U}(A)&=\textnormal{ the group of unitary elements of }A,\\
    \mathcal{U}_n(A)&=\mathcal{U}(M_n(A)),\\
    \mathcal{U}_\infty(A)&=\bigcup_{n=1}^{\infty}\mathcal{U}_n(A).
  \end{align*}
  We define a relation \(\sim_1\) in \(\mathcal{U}_\infty(A)\) as follows. For \(u\in\mathcal{U}_n(A)\) and \(v\in\mathcal{U}_m(A)\)
  we have \(u\sim_1 v\) if there exists \(k\geq\max\{n,m\}\) such that \(\textnormal{diag}(u,1_{k-n})\sim_h\textnormal{diag}(v,1_{k-m})\).
  Then \(\sim_1\) is an equivalence relation (exercise). We denote by \({[u]}_1\) the equivalence class of the unitary \(u\in\mathcal{U}_\infty(A)\).
\end{definition}
\begin{lemma}
  Let \(A\) be a unital \cstar. Then
  \begin{enumerate}
    \item \({[u]}_1{[v]}_1={[\textnormal{diag}(u,v)]}_1\) is a well-defined associative binary operation on \(\mathcal{U}_\infty(A)/\sim_1\),
    \item \({[u]}_1{[v]}_1={[v]}_1{[u]}_1\) for all \(u,v\in\mathcal{U}_\infty(A)\),
    \item \({[u]}_1{[1_n]}_1={[1_n]}_1{[u]}_1={[u]}_1\) for all \(n\) and all \(u\in\mathcal{U}_\infty(A)\),
    \item if \(u,v\in\mathcal{U}_m(A)\) then \({[u]}_1{[v]}_1={[uv]}_1\).
  \end{enumerate}
\end{lemma}
\begin{proof}
  Stabilizing a unitary by identity blocks does not change its \(\sim_1\)-class. Hence, if
  \(u\sim_1u'\) and \(v\sim_1v'\), we may place all four unitaries in fixed matrix algebras
  and take direct sums of the corresponding homotopies. This proves that the operation in (1)
  is well defined; associativity follows directly from associativity of block sums.
  A scalar permutation unitary is homotopic to the identity, because every unitary in
  \(M_k(\mathbb{C})\) is homotopic to \(1_k\). Conjugating by the permutation which exchanges
  the two diagonal blocks therefore gives
  \[
    \diag(u,v)\sim_h\diag(v,u),
  \]
  and proves (2). Part (3) follows from the definition of stabilization. Finally, the Whitehead
  lemma, Lemma~\ref{lemma:516}, gives
  \[
    \diag(u,v)\sim_h\diag(uv,1_m)
  \]
  whenever \(u,v\in\mathcal{U}_m(A)\), which proves (4). In particular,
  \({[u]}_1{[u^*]}_1={[1]}_1\), so the resulting commutative monoid is a group.
\end{proof}
By the above Lemma, \(\mathcal{U}_\infty(A)/\sim_1\) equipped with the multiplication \({[u]}_1{[v]}_1={[\textnormal{diag}(u,v)]}_1\) is an abelian group,
with \({[u]}_1^{-1}={[u^*]}_1\).
\begin{definition}
  If \(A\) is a \cstar\ then we define 
  \[
    K_1(A)=\mathcal{U}_\infty(\tilde{A})/\sim_1,
  \]
  an abelian group with multiplication \({[u]}_1{[v]}_1={[\textnormal{diag}(u,v)]}_1\).
\end{definition}
When \(A\) is \textit{unital} then \(K_1(A)\) may be defined simply as \(\mathcal{U}_\infty(A)/\sim_1\) (see Exercise~\ref{ex:731}).
Also, instead of using equivalence classes of \textit{unitaries} one could define \(K_1\) with help of equivalence classes of
\textit{invertibles} (see Exercise~\ref{ex:736}). In particular, the polar decomposition \(w=u|w|\) yields a well-defined map
\[
  {[\cdot]}_1:\textnormal{GL}_\infty(A)\rightarrow K_1(A)
\]
by \({[w]}_1={[u]}_1={[w|w|^{-1}]}_1\).
\begin{proposition}[Universal property of \(K_1\)]\label{prop:713}
  Let \(A\) be a \cstar, \(G\) an abelian group, and \(\nu:\mathcal{U}_\infty(\tilde{A})\rightarrow G\) a map such that:
  \begin{enumerate}
    \item \(\nu(\diag(u,v))=\nu(u)+\nu(v)\),
    \item \(\nu(1_n)=0\) for every \(n\),
    \item if \(u,v\in\mathcal{U}_n(\tilde{A})\) and \(u\sim_h v\), then \(\nu(u)=\nu(v)\).
  \end{enumerate}
  Then there exists a unique homomorphism \(K_1(A)\rightarrow G\) making the diagram
  \begin{equation}\label{diagUP:711}
    \begin{tikzcd}[column sep=scriptsize]
    \mathcal{U}_\infty(\tilde{A}) \arrow{d}[swap]{{[\cdot]}_1} \arrow{rd}{\nu}&\\
    K_1(A)\arrow{r}[swap]{} & G 
    \end{tikzcd}
  \end{equation}
  commutative
\end{proposition}
\begin{proof}
  Suppose that \(u\in\mathcal{U}_n(\tilde A)\), \(v\in\mathcal{U}_m(\tilde A)\), and
  \(u\sim_1v\). Choose \(k\geq\max\{n,m\}\) such that
  \(\diag(u,1_{k-n})\sim_h\diag(v,1_{k-m})\). Conditions (1)--(3) give
  \[
    \nu(u)=\nu(\diag(u,1_{k-n}))
    =\nu(\diag(v,1_{k-m}))=\nu(v).
  \]
  Thus \(\alpha({[u]}_1)=\nu(u)\) defines a map \(\alpha:K_1(A)\to G\). Moreover,
  \[
    \alpha({[u]}_1{[v]}_1)
    =\nu(\diag(u,v))=\nu(u)+\nu(v)
    =\alpha({[u]}_1)+\alpha({[v]}_1),
  \]
  so \(\alpha\) is a homomorphism. Its uniqueness follows from the surjectivity of the
  quotient map \({[\cdot]}_1:\mathcal{U}_\infty(\tilde A)\to K_1(A)\).
\end{proof}
Let \(A,B\) be \cstars\ and let \(\varphi:A\rightarrow B\) be a *-homomorphism. Then \(\varphi\) extends to a
*-homomorphism \(\tilde{\varphi}_n:M_n(\tilde{A})\rightarrow M_n(\tilde{B})\) and thus yields a map
\(\tilde{\varphi}:\mathcal{U}_\infty(\tilde{A})\rightarrow \mathcal{U}_\infty(\tilde{B})\).
We define \(\nu:\mathcal{U}_\infty(\tilde{A})\rightarrow K_1(B)\) by \(\nu(u)={[\tilde{\varphi}(u)]}_1\) and use the universal property of \(K_1\)
to conclude that there exists a unique homomorphism \(K_1(\varphi):K_1(A)\rightarrow K_1(B)\) such that \(K_1(\varphi)({[u]}_1)={[\tilde{\varphi}(u)]}_1\) for \(u\in\mathcal{U}_\infty(\tilde{A})\).
\begin{proposition}[Functoriality of \(K_1\)]
  Let \(A,B,C\) be \cstars\ and let \(\varphi:A\rightarrow B\) and \(\psi:B\rightarrow C\) be *-homomorphisms. Then
  \begin{enumerate}
    \item \(K_1(\textnormal{id}_A)=\textnormal{id}_{K_1(A)}\),
    \item \(K_1(\psi\circ\varphi)=K_1(\psi)\circ K_1(\varphi)\).
  \end{enumerate}
  Thus \(K_1\) is a covariant functor.
\end{proposition}
\begin{proof}
  For every \(u\in\mathcal{U}_\infty(\tilde A)\),
  \[
    K_1(\operatorname{id}_A)({[u]}_1)
    ={[\widetilde{\operatorname{id}_A}(u)]}_1={[u]}_1,
  \]
  which proves (1). Similarly,
  \begin{align*}
    K_1(\psi\circ\varphi)({[u]}_1)
      &={[(\widetilde{\psi\circ\varphi})(u)]}_1\\
      &={[(\tilde\psi\circ\tilde\varphi)(u)]}_1
       =(K_1(\psi)\circ K_1(\varphi))({[u]}_1),
  \end{align*}
  proving (2).
\end{proof}
It is also clear from the definition that \(K_1\) of the zero algebra and the zero map are zero.
\begin{proposition}[Homotopy invariance of \(K_1\)]
  Let \(A,B\) be \cstars.
  \begin{enumerate}
    \item If \(\varphi,\psi:A\rightarrow B\) are homotopic *-homomorphisms then \(K_1(\varphi)=K_1(\psi)\).
    \item If \(A\) and \(B\) are homotopy equivalent then \(K_1(A)\cong K_1(B)\). More specifically,
    if \(\varphi:A\rightarrow B\) and \(\psi:B\rightarrow A\) satisfy
    \(\psi\circ\varphi\simeq_h\textnormal{id}_A\) and
    \(\varphi\circ\psi\simeq_h\textnormal{id}_B\), then \(K_1(\varphi)\) and
    \(K_1(\psi)\) are isomorphisms and inverses of one another.
  \end{enumerate}
\end{proposition}
\begin{proof}
  Let \(\{\varphi_t\}_{t\in[0,1]}\) be a homotopy from \(\varphi\) to \(\psi\). If
  \(u\in\mathcal{U}_n(\tilde A)\), then
  \[
    [0,1]\ni t\longmapsto\widetilde{\varphi_t}_n(u)
  \]
  is a continuous path of unitaries from \(\tilde\varphi_n(u)\) to \(\tilde\psi_n(u)\).
  Consequently,
  \[
    K_1(\varphi)({[u]}_1)={[\tilde\varphi_n(u)]}_1
    ={[\tilde\psi_n(u)]}_1=K_1(\psi)({[u]}_1),
  \]
  which proves (1). If \(\psi\circ\varphi\) and \(\varphi\circ\psi\) are homotopic to the
  corresponding identity maps, part (1) and functoriality yield
  \[
    K_1(\psi)K_1(\varphi)=\operatorname{id}_{K_1(A)},\qquad
    K_1(\varphi)K_1(\psi)=\operatorname{id}_{K_1(B)}.
  \]
  Hence the two induced maps are inverse isomorphisms.
\end{proof}
\begin{theorem}[(Half) exactness of \(K_1\)]
  If
  \begin{equation}\label{seq:712}
    \begin{tikzcd}
    0\arrow[r]& J \arrow[r, "\varphi"]& A \arrow[r, "\pi"]&B\arrow[r]& 0
    \end{tikzcd}
  \end{equation}
  is an exact sequence of \cstars\ then the sequence
  \begin{equation}\label{seq:713}
    \begin{tikzcd}
     K_1(J) \arrow[r, "K_1(\varphi)"]& K_1(A) \arrow[r, "K_1(\pi)"]&K_1(B)
    \end{tikzcd}
  \end{equation}
  is exact. If the sequence (\ref{seq:712}) is split-exact with a splitting map \(\lambda:B\rightarrow A\),
  then the sequence 
  \begin{equation}\label{seq:714}
    \begin{tikzcd}
    0\arrow[r] &K_1(J) \arrow[r, "K_1(\varphi)"]& K_1(A) \arrow[r, "K_1(\pi)"]&K_1(B)\arrow[r]& 0
    \end{tikzcd}
  \end{equation}
  is split-exact with a splitting map \(K_1(\lambda):K_1(B)\rightarrow K_1(A)\).
\end{theorem}
\begin{proof}
  \(K_1(\pi)\circ K_1(\varphi)=K_1(\pi\circ\varphi)=K_1(0)=0\) by functoriality of \(K_1\), and hence
  \(\textnormal{Im}(K_1(\varphi))\subseteq \textnormal{Ker}(K_1(\pi))\). For the inverse inclusion, let \(u\in\mathcal{U}_n(\tilde{A})\) and \(K_1(\pi)({[u]}_1)={[1]}_1\).
  Then there is \(m\) such that \(\diag(\tilde{\pi}(u),1_m)\sim_h 1_{n+m}\).
  By Lemma~\ref{lemma:518}, there is \(w\in\mathcal{U}_{n+m}(\tilde{A})\) such that \(w\sim_h 1_{n+m}\)
  and \(\tilde{\pi}(w)=\diag(\tilde{\pi}(u),1_m)\). Thus \({[u]}_1={[\diag(u,1_m)w^*]}_1\) and \(\tilde{\pi}(\diag(u,1_m)w^*)=1_{n+m}\).
  By Lemma~\ref{lemma:Lars431}, there is \(v\in\mathcal{U}_{n+m}(\tilde{J})\) such that \(\tilde{\varphi}(v)=\diag(u,1_m)w^*\).
  Thus \({[u]}_1\in\textnormal{Im}(K_1(\varphi))\) and consequently \(\textnormal{Ker}(K_1(\pi))\subseteq\textnormal{Im}(K_1(\varphi))\). This shows that the sequence (\ref{seq:712}) is exact.\par
  Now suppose that the sequence (\ref{seq:712}) is split-exact. Then the sequence (\ref{seq:714}) is exact at \(K_1(A)\) by the preceding argument.
  By functoriality of \(K_1\) we have \(K_1(\pi)\circ K_1(\lambda)=\textnormal{id}_{K_1(B)}\),
  and hence (\ref{seq:714}) is exact at \(K_1(B)\) and \(K_1(\lambda)\) is a splitting map.
  It remains to show that \(K_1(\varphi)\) is injective. So let \(u\in\mathcal{U}_n(\tilde{J})\) be such that \(K_1(\varphi)({[u]}_1)={[1]}_1\).
  Then there is \(m\) such that \(\diag(\tilde{\varphi}(u),1_m)\sim_h 1_{n+m}\). Let \(t\mapsto w_t\) be a continuous path in \(\mathcal{U}_{n+m}(\tilde{A})\) connecting
  \(\diag(\tilde{\varphi}(u),1_m)\) and \(1_{n+m}\). Our goal is to apply \(\tilde{\varphi}^{-1}\) to \(w_t\) to conclude that \(\diag(u,1_m)\) is homotopic to the identity.
  In general, this is impossible since some of \(w_t\) may lie outside the range of \(\tilde{\varphi}\). However, in the presence of a splitting map \(\lambda\) we can correct the path \(w_t\) by setting \(v_t=w_t(\tilde{\lambda}\circ\tilde{\pi})(w_t^*)\). Then \(v_t\) is a continuous
  path in \(\mathcal{U}_{n+m}(\tilde{A})\) connecting \(\diag(\tilde{\varphi}(u),1_{m})(\tilde{\lambda}\circ\tilde{\pi})(\diag(\tilde{\varphi}(u^*),1_m))\) and \(1_{n+m}\).
  Since \(\tilde{\pi}(v_t)=1_{n+m}\) for all \(t\), Lemma~\ref{lemma:Lars431} implies that each \(v_t\) is in the image of \(\tilde{\varphi}\).
  Thus \(\tilde{\varphi}^{-1}(v_t)\) is a continuous path in \(\mathcal{U}_{n+m}(\tilde{J})\) connecting \(\diag(u,1_m)\tilde{\varphi}^{-1}((\tilde{\lambda}\circ\tilde{\pi})(\diag(\tilde{\varphi}(u^*),1_m)))\) and \(1_{n+m}\).
  Since \(\tilde{\varphi}^{-1}((\tilde{\lambda}\circ\tilde{\pi})(\diag(\tilde{\varphi}(u^*),1_m)))\) is a scalar matrix, it is homotopic to the identity. Thus
  \[
    {[u]}_1={[\diag(u,1_m)]}_1={\left[\diag(u,1_m)\tilde{\varphi}^{-1}((\tilde{\lambda}\circ\tilde{\pi})(\diag(\tilde{\varphi}(u^*),1_m)))\right]}_1={[1]}_1,
  \]
  and the map \(K_1(\varphi)\) is injective.
\end{proof}
\begin{proposition}[Continuity of \(K_1\)]
  Let \(A=\varinjlim\{A_i,\Phi_{ij}\}\) be the inductive limit of a sequence of \cstars,
  and let \(\Phi_i:A_i\rightarrow A\) be the canonical maps. Let \(G=\varinjlim\{K_1(A_i),K_1(\Phi_{ij})\}\)
  be the inductive limit of the corresponding sequence of abelian groups, and let
  \(\varphi_i:K_1(A_i)\rightarrow G\) be the canonical maps. Then there exists an isomorphism \(\Lambda:G\rightarrow K_1(A)\) such that for all \(i\geq j\) the diagram
  \begin{equation}\label{diag:715}
    \begin{tikzcd}
      K_1(A_j)\arrow[r,"\varphi_j"]\arrow[d,swap,"K_1(\Phi_{ij})"]\arrow[dr]&G\arrow[d,"\Lambda"]\\
      K_1(A_i)\arrow[r,swap,"K_1(\Phi_i)"]\arrow[ur]& K_1(A)
    \end{tikzcd}
  \end{equation}
  is commutative.
\end{proposition}
\begin{proof}
  The universal property of the direct limit \(G\) of the sequence \(\{K_1(A_i),K_1(\Phi_{ij})\}\) yields
  a unique homomorphism \(\Lambda:G\rightarrow K_1(A)\) making (\ref{diag:715}) commutative. We must show that \(\Lambda\) is surjective and injective.\par
  \textit{Surjectivity.} Let \(u\in\mathcal{U}_n(\tilde{A})\). By part (\ref{ex:73102}) of Exercise~\ref{ex:7310}, there is \(i\) and \(w\in\mathcal{U}_n(\tilde{A_i})\)
  such that \(\left\|u-\widetilde{\Phi_i}(w)\right\|<2\). Thus \(u\) and \(\widetilde{\Phi_i}(w)\) are homotopic in \(\mathcal{U}(\tilde{A})\) by Lemma \ref{lemma:Matthes214}.
  Hence
  \[
    {[u]}_1={[\widetilde{\Phi_i}(w)]}_1=K_1(\Phi_i)({[w]}_1)=(\Lambda\circ\varphi_i)({[w]}_1),
  \]
  and \(\Lambda\) is surjective.\par
  \textit{Injectivity.} It suffices to show that for each \(j\) the restriction of \(\Lambda\) to the image of \(\varphi_j\) is injective.
  So let \(u\in\mathcal{U}_n(\tilde{A_j})\) be such that \((\Lambda\circ\varphi_j)({[u]}_1)=K_1(\Phi_j)({[u]}_1)={[1]}_1\) in \(K_1(A)\).
  We must show that \(\varphi_j({[u]}_1)=0\) in \(G\). Indeed, there is \(m\) such that \(\diag(\widetilde{\Phi_j}(u),1_m)\) is homotopic to \(1_{n+m}\) in \(\mathcal{U}_{n+m}(\tilde{A})\).
  By part (\ref{ex:73103}) of Exercise~\ref{ex:7310}, there is \(i\geq j\) such that
  \(\diag(\widetilde{\Phi_{ij}}(u),1_m)\) is homotopic to \(1_{n+m}\).
  Thus \({[\widetilde{\Phi_{ij}}(u)]}_1={[\diag(\widetilde{\Phi_{ij}}(u),1_m)]}_1={[1]}_1\).
  Consequently, \(\varphi_j({[u]}_1)=(\varphi_i\circ K_1(\Phi_{ij}))({[u]}_1)=0\), and \(\Lambda\) is injective.
\end{proof}
\begin{proposition}[Stability of \(K_1\)]
  Let \(A\) be a \cstar.
  \begin{enumerate}
    \item For each \(n\in\mathbb{N}\) we have
    \[
      K_1(A)\cong K_1(M_n(A)).
    \]
    More specifically, let \(\psi:A\rightarrow M_n(A)\) be such that \(\psi(a)=\diag(a,0_{n-1})\). Then \(K_1(\psi):K_1(A)\rightarrow K_1(M_n(A))\) is an isomorphism.
    \item Let \(\mathcal{K}\) be the \cstar\ of compact operators. Then 
      \[
        K_1(A)\cong K_1(A\otimes\mathcal{K}).
      \]
      More specifically, let \(p\) be a minimal projection in \(\mathcal{K}\) and let \(\varphi:A\rightarrow A\otimes\mathcal{K}\) be the map
      such that \(\varphi(a)=a\otimes p\). Then \(K_1(\varphi):K_1(A) \rightarrow K_1(A\otimes \mathcal{K})\) is an isomorphism.
  \end{enumerate}
\end{proposition}
\begin{proof}
  \begin{enumerate}
    \item Exercise.
    \item Since
    \[
      A\otimes\mathcal{K}\cong A\otimes (\varinjlim M_n(\mathbb{C}))\cong \varinjlim M_n(A),
    \]
    the claim follows from (1) and the continuity of \(K_1\).
  \end{enumerate}
\end{proof}

\section{The index map}
The construction of the index map for extensions and its relation with Fredholm
operators in this section follow \cite[Sections~9.1--9.4]{RordamLarsenLaustsen2000}.
Let \(\mathcal{H}\) be a separable, infinite dimensional Hilbert space. We denote by \(\mathcal{F}\)
the algebra of finite rank operators on \(\mathcal{H}\) (a two-sided *-ideal in \(\mathcal{B}(\mathcal{H})\)),
by \(\mathcal{K}\) the \cstar\ of compact operators on \(\mathcal{H}\) (the norm closure of \(\mathcal{F}\) and the only non-trivial, norm closed, two-sided
ideal of \(\mathcal{B}(\mathcal{H})\)), by \(\mathcal{Q}=\mathcal{B}(\mathcal{H})/\mathcal{K}\) the Calkin algebra, and by \(\pi:\mathcal{B}(\mathcal{H})\rightarrow\mathcal{Q}\)
the natural surjection.
  
\begin{theorem}[Atkinson]
  If \(F\in\mathcal{B}(\mathcal{H})\) then the following conditions are equivalent.
  \begin{enumerate}
    \item Both \(\textnormal{Ker}(F)\) and \(\textnormal{Coker}(F)\) are finite dimensional.
    \item There exists an operator \(G\in\mathcal{B}(\mathcal{H})\) such that both \(FG-1\) and \(GF-1\) are compact.
    \item The image \(\pi(F)\) of \(F\) in the Calkin algebra \(\mathcal{Q}\) is invertible.
  \end{enumerate}
  Furthermore, if \(F\) satisfies these conditions, then the range of \(F\) is closed in \(\mathcal{H}\).
\end{theorem}
\begin{proof}
  (2) \(\Leftrightarrow\) (3): Exercise\par
  (1) \(\Rightarrow\) (2): Let \(\mathcal{H}_0\)
be a subspace of \(\mathcal{H}\) of smallest possible dimension such that \(\textnormal{Im}(F) + \mathcal{H}_0 = \mathcal{H}\). Then
\(n = \dim(\mathcal{H}_0)\) is finite, since the cokernel of \(F\) is finite dimensional. Then the restriction
of \(F\) to the orthogonal complement of its kernel is a bijection from \({\textnormal{Ker}(F)}^\perp\) onto \(\textnormal{Im}(F)\),
and it extends to a linear bijection \(\tilde{F} : {\textnormal{Ker}(F)}^\perp\oplus \mathbb{C}^n \rightarrow \textnormal{Im}(F)+\mathcal{H}_0 = \mathcal{H}\). 
By the Inverse Mapping Theorem, the inverse of \(\tilde{F}\) is continuous. It follows that \(\textnormal{Im}(F) = \tilde{F}({\textnormal{Ker}(F)}^\perp)\)
is closed in \(\mathcal{H}\).\par
By the preceding argument, \(F\) yields a continuous linear bijection from \({\textnormal{Ker}(F)}^\perp\) onto
\(\textnormal{Im}(F)\) — a closed subspace of \(\mathcal{H}\). Thus, by the Inverse Mapping Theorem, it has a
continuous inverse \(G : \textnormal{Im}(F) \rightarrow {\textnormal{Ker}(F)}^\perp\). Extend G to a bounded linear operator on \(\mathcal{H}\)
  (still denoted \(G\)) by setting \(G\xi = 0\) for \(\xi \in {\textnormal{Im}(F)}^\perp\). Then both \(FG-1\) and \(GF-1\) have
  finite rank and (2) holds.\par
  (2) \(\Rightarrow\) (1): Let \(K_1,K_2\) be compact operators such that \(GF=1+K_1\) and \(FG=1+K_2\). Then
\(\textnormal{Ker}(F)\subseteq\textnormal{Ker}(GF)=\textnormal{Ker}(1+K_1)\), and \(\textnormal{Ker}(1+K_1)\) is the eigenspace of \(K_1\)
corresponding to the eigenvalue \(-1\). Since \(K_1\) is compact, this eigenspace is finite dimensional and so is the kernel of \(F\). We also have
\(\textnormal{Im}(F)\supseteq\textnormal{Im}(FG)=\textnormal{Im}(1+K_2)\). Since \(1+K_2\) can be written as an invertible plus a finite
  rank operator, its range has finite codimension. Thus \(\textnormal{Coker}(F)\) is finite dimensional.
\end{proof}
A bounded operator satisfying the conditions of Atkinson's theorem is called \textit{Fredholm}.
In particular, any invertible operator in \(\mathcal{B}(\mathcal{H})\) is Fredholm.
It follows immediately that if \(F,T\) are Fredholm and \(K\) is compact, then the operators \(F^*, FT\) and \(F+K\)
are Fredholm.\par
If \(F,G\) are Fredholm operators satisfying (2), then \(G\) is called a \textit{parametrix} of \(F\).
\begin{definition}
  Let \(F\) be a Fredholm operator. Then its \textit{Fredholm index} is an integer defined as
  \[
    \textnormal{Index}(F)=\dim(\textnormal{Ker}(F))-\dim(\textnormal{Coker}(F)).
  \]
\end{definition}
Since \(\dim(\textnormal{Coker}(F)) = \dim(\textnormal{Ker}(F^*))\), we have \(\Index(F) = -\Index(F^*)\). If \(F\) is
Fredholm and \(V\) is invertible then clearly \(\Index(FV) = \Index(VF) = \Index(F)\) and
\(\Index(V ) = 0\).\par
Let \({\xi_n : n = 0,1,\ldots}\) be an orthonormal basis of \(\mathcal{H}\). The operator \(S \in \mathcal{B}(\mathcal{H})\) such
that \(S(\xi_n) = \xi_{n+1}\) is called unilateral shift. It is a Fredholm operator with index \(-1\). Thus
for any positive integer \(k\) we have \(\Index(S^k) = -k\) and \(\Index({(S^*)}^k) = k\).\par
In a finite dimensional Hilbert space all operators are compact and hence all operators
are Fredholm. The rank-nullity theorem of elementary linear algebra may then be
interpreted as saying that every Fredholm operator on a finite dimensional Hilbert space
has index \(0\).
\begin{theorem}[Riesz]
  If \(F\) is Fredholm and \(K\) is compact then
  \[
    \Index(F+K)=\Index(F).
  \]
\end{theorem}
\begin{proof}
  We first observe that if \(R\) is of finite rank then \(\Index(1+R)=0\). Indeed, let
  \(\mathcal{H}_0=\textnormal{Im}(R)+{\textnormal{Ker}(R)}^\perp\). Then \(\mathcal{H}_0\) is finite dimensional, and the restriction of both \(R\) and \(R^*\)
  to \(\mathcal{H}_0^\perp\) is zero. Thus the index of \(1+R\) coincides with the index of its restriction to \(\mathcal{H}_0\) and hence is \(0\).\par
  Now let \(K\) be compact. Find \(R\) of finite rank such that \(\left\|K-R\right\|<1\). Then \(V=1+(K-R)\) is invertible. Hence
  \[
    \Index(1+K)=\Index(V+R)=\Index(V(1+V^{-1}R))=0
  \]\par
  Let \(F\) be a Fredholm operator of index \(0\). Then there is a finite rank operator \(R\) such that
  \(R\) maps bijectively \(\textnormal{Ker}(F)\) onto \({\textnormal{Im}(F)}^\perp=\textnormal{Ker}(F^*)\).
  Let \(V=F+R\). Then \(V\) is a continuous linear bijection of \(\mathcal{H}\) onto itself and hence it is an invertible operator. Thus if \(K\) is compact then
  \[
    \Index(F+K)=\Index(V+(K-R))=\Index(V(1+V^{-1}(K-R)))=0.
  \]\par
  Finally, let \(F\) be an arbitrary Fredholm operator and let \(K\) be compact. Then \(\Index(F\oplus F^*)=0\) and hence
  \(\Index((F+K)\oplus F^*)=0\). Consequently, \(\Index(F+K)=-\Index(F^*)=\Index(F)\).
\end{proof}
\begin{corollary}
  If \(F,T\) are Fredholm operators then
  \[
    \Index(FT)=\Index(F)+\Index(T).
  \]
\end{corollary}
\begin{proof}
  Suppose first that \(\Index(F)=0\), and let \(R\) be an operator of finite rank such that
  \(F+R\) is invertible (as shown in the previous proof). Then
  \[
    \Index(FT)=\Index(FT+RT)=\Index((F+R)T)=\Index(T).
  \]
  Now suppose that \(\Index(F)=k>0\), and let \(S\) be a unilateral shift on \(\mathcal{H}\).
  Then \(\Index(F\oplus S^k)=0\) and hence
  \[
    \Index(FT\oplus S^k)=\Index((F\oplus S^k)(T\oplus 1))=\Index(T\oplus 1)=\Index(T).
  \]
  Consequently, we have
  \[
    \Index(FT)=-\Index(S^k)+\Index(T)=\Index(F)+\Index(T),
  \]
  This proves the claim when \(\Index(F)>0\). If \(\Index(F)=-k<0\), then
  \(\Index(F\oplus {(S^*)}^k)=0\). Applying the zero-index case to
  \(F\oplus {(S^*)}^k\) and \(T\oplus 1\), we obtain
  \[
    \Index(FT)+k=\Index(FT\oplus {(S^*)}^k)=\Index(T),
  \]
  and hence \(\Index(FT)=\Index(F)+\Index(T)\) also in this case.
\end{proof}
In particular, if \(G\) is a parametrix of \(F\) then \(\Index(G)=-\Index(F)\).
\begin{proposition}
  The index map is locally constant and continuous in norm.
\end{proposition}
\begin{proof}
  Let \(F\) be a Fredholm operator and let \(G\) be its parametrix. Let \(K\) be a compact operator
  such that \(FG=1+K\). It suffices to show that if \(T\) is a Fredholm operator with 
  \(\left\|T-F\right\|<1/\left\|G\right\|\) then \(\Index(F)=\Index(T)\). Indeed, the operator \((T-F)G+1\)
  is invertible, since its distance from the identity is less than \(1\). Thus
  \[
    \Index(T)+\Index(G) = \Index(TG) = \Index((T-F+F)G) = \Index((T-F)G+1+K) = 0.  
  \]
  Thus \(\Index(T)=-\Index(G)=\Index(F)\).
\end{proof}
\begin{definition}
  If \(F,T\) are two Fredholm operators then we say that they are \textnormal{homotopic} if there exists
  a norm-continuous path from \(F\) to \(T\) consisting of Fredholm operators.
\end{definition}
\begin{proposition}
  Two Fredholm operators are homotopic if and only if they have the same index.
\end{proposition}
\begin{proof}
   Let \(F\) and \(T\) be Fredholm operators.\par
Suppose that \(F\) and \(T\) are homotopic, and let \(t \mapsto V_t\) be a continuous path of Fredholm
operators from \(F\) to \(T\). Then the map \(t \mapsto \Index(V_t)\) is continuous and hence constant.\par
To show the converse we first observe that every Fredholm operator \(V\) with \(\Index(V ) =
  0\) is homotopic to \(1\). Indeed, there is a finite-rank operator \(R\) such that \(V+R\) is invertible.
 Then \(t \mapsto V +tR\) is a path connecting \(V\) to an invertible element. By polar decomposition,
 every invertible operator is homotopic to a unitary; moreover, if \(u\) is unitary, bounded Borel functional
 calculus gives a self-adjoint operator \(h\) such that \(u=e^{ih}\), so \(t\mapsto e^{ith}\) joins \(1\) to
 \(u\). Thus the group of invertibles of \(\mathcal{B}(\mathcal{H})\) is path-connected; compare
 \cite[Example~8.1.8]{RordamLarsenLaustsen2000}.\par
Now suppose that \(\Index(F) = \Index(T)\). Then both \(FT^*\) and \(T^*T\) have index \(0\) and
thus are homotopic to \(1\). Consequently, the operators \(F, F(T^*T) = (FT^*)T\) and \(T\) are
homotopic.
\end{proof}

We now construct the index map associated with an arbitrary extension, compare
\cite[Sections~9.1--9.3]{RordamLarsenLaustsen2000}.
Let 
\[
  \begin{tikzcd}
    0\arrow[r] &J \arrow[r, "\varphi"]& A \arrow[r, "\psi"]&B\arrow[r]& 0
    \end{tikzcd} 
\]
  be an exact sequence of \cstars. Let \(u\in\mathcal{U}_n(\tilde{B})\).
Then there exists a unitary \(V\) in \(\mathcal{U}_{2n}(\tilde{A})\) such that
\begin{equation}\label{eq:727}
  \tilde{\psi}(V)=\diag(u,u^*).
\end{equation}
Then \(\tilde{\psi}(V\diag(1_n,0)V^*)=\diag(1_n,0)\). Thus there exists a projection \(P\) in \(\mathcal{P}_{2n}(\tilde{J})\) such that
\begin{equation}\label{eq:728}
  \tilde{\varphi}(P)=V\diag(1_n,0)V^*.
\end{equation}
Since \(\tilde{\psi}\circ\tilde{\varphi}(P)=\diag(1_n,0)\), it follows that \(s(P)=\diag(1_n,0)\), where \(s\) is the scalar map.
Then there is a well-defined map 
\[
  \mu:\mathcal{U}_\infty(\tilde{B})\rightarrow K_0(J)\text{ such that }\mu(u)={[P]}_0-{[s(P)]}_0.
\]
Indeed, suppose that \(W\in\mathcal{U}_{2n}(\tilde{A})\) and \(Q\in\mathcal{P}_{2n}(\tilde{J})\) are such that \(\tilde{\psi}(W)=\diag(u,u^*)\)
and \(\tilde{\varphi}(Q)=W\diag(1_n,0)W^*\). We must show that \({[P]}_0-{[s(P)]}_0={[Q]}_0-{[s(Q)]}_0\) in \(K_0(J)\).
  Indeed since \(\tilde{\psi}(VW^*)=1_{2n}\) there is \(Y\in\mathcal{U}_{2n}(\tilde{J})\) such that \(\tilde{\varphi}(Y)=VW^*\). Since
\[
  \tilde{\varphi}(P)=VW^*\tilde{\varphi}(Q){(VW^*)}^*=\tilde{\varphi}(YQY^*),
\]
we have \(P=YQY^*\) and the claim follows. That is, \(\mu:\mathcal{U}_\infty(\tilde{B})\rightarrow K_0(J)\) is well defined.\par
This map \(\mu\) satisfies the conditions (1) to (3) of Proposition~\ref{prop:713}:
\begin{enumerate}
  \item Let \(u_j\in\mathcal{U}_{n_j}(\tilde{B})\), and choose \(V_j\) and \(P_j\) as in
  (\ref{eq:727}) and (\ref{eq:728}), for \(j=1,2\). Let \(y\) be the scalar permutation
  unitary which rearranges the four diagonal blocks so that
  \[
    y\diag(u_1,u_1^*,u_2,u_2^*)y^*
      =\diag(u_1,u_2,u_1^*,u_2^*)
  \]
  and
  \[
    y\diag(1_{n_1},0_{n_1},1_{n_2},0_{n_2})y^*
      =\diag(1_{n_1+n_2},0_{n_1+n_2}).
  \]
  Then \(V=y(V_1\oplus V_2)y^*\) and \(P=y(P_1\oplus P_2)y^*\) may be used in
  the definition of \(\mu(\diag(u_1,u_2))\). Since conjugation by a scalar unitary does
  not change a \(K_0\)-class,
  \begin{align*}
    \mu(\diag(u_1,u_2))
      &={[P_1\oplus P_2]}_0-{[s(P_1)\oplus s(P_2)]}_0\\
      &=\mu(u_1)+\mu(u_2).
  \end{align*}
  \item For \(u=1_n\), take \(V=1_{2n}\) and \(P=\diag(1_n,0_n)\). Then
  \(P=s(P)\), and therefore \(\mu(1_n)=0\).
  \item Let \(u\sim_h w\in\mathcal{U}_n(\tilde{B})\), \(U,W\in\mathcal{U}_{2n}(\tilde{A})\), \(P,Q\in\mathcal{P}_{2n}(\tilde{J})\)
    be such that \(\tilde{\psi}(U)=\diag(u,u^*)\), \(\tilde{\psi}(W)=\diag(w,w^*)\), \(\tilde{\varphi}(P)=U\diag(1_n,0)U^*\)
    and \(\tilde{\varphi}(Q)=W\diag(1_n,0)W^*\) (that is, \(\{u,U,P\}\) and \(\{w,W,Q\}\) satisfy conditions (\ref{eq:727}) and (\ref{eq:728}), respectively).
    Then \(u^*w\sim_h 1_n\sim_h uw^*\) and thus there exist \(X,Y\in\mathcal{U}_n(\tilde{A})\) such that
    \(\tilde{\psi}(X)=u^*w\) and \(\tilde{\psi}(Y)=uw^*\). Put \(Z=U\diag(X,Y)\), a unitary in \(\mathcal{U}_{2n}(\tilde{A})\).
    We have \(\tilde{\psi}(Z)=\diag(w,w^*)\) and \(\tilde{\varphi}(P)=Z\diag(1_n,0)Z^*\).
    Thus, by the definition of \(\mu\), we have \(\mu(w)={[P]}_0-{[s(P)]}_0=\mu(u)\). 
\end{enumerate}
The universal property of \(K_1\) now implies that there exists
a homomorphism 
\[
  \partial_1:K_1(B)\rightarrow K_0(J),
\]
called the \textit{index map}, such that
\[
  \partial_1({[u]}_1)={[P]}_0-{[s(P)]}_0.
\]
The following exactness theorem is proved as in
\cite[Proposition~9.3.3]{RordamLarsenLaustsen2000}.
\begin{theorem}\label{theorem:727}
  Let
  \[
    \begin{tikzcd}
    0\arrow[r] &J \arrow[r, "\varphi"]& A \arrow[r, "\psi"]&B\arrow[r]& 0
    \end{tikzcd} 
  \]
  be an exact sequence of \cstars. Then the sequence
  \begin{equation}
    \begin{tikzcd}
    K_1(J) \arrow[r, "K_1(\varphi)"] & K_1(A) \arrow[r, "K_1(\psi)"]&K_1(B)\arrow[d,"\partial_1"]\\
    K_0(B) & K_0(A)\arrow[l, "K_0(\psi)"] &K_0(J)\arrow[l, "K_0(\varphi)"]\\
    \end{tikzcd} 
  \end{equation}
  is exact everywhere.
\end{theorem}
\begin{proof}
  By virtue of half-exactness of \(K_0\) and \(K_1\), it suffices to prove that \(\textnormal{Im}(K_1(\psi))=\textnormal{Ker}(\partial_1)\)
  and \(\textnormal{Im}(\partial_1)=\textnormal{Ker}(K_0(\varphi))\).
  \begin{enumerate}
    \item We show that \(\textnormal{Im}(K_1(\psi))\subseteq\textnormal{Ker}(\partial_1)\). Indeed, if \(U\in\mathcal{U}_n(\tilde{A})\) then
    \(\diag(\tilde{\psi}(U),{\tilde{\psi}(U)}^*)\) lifts to a diagonal unitary \(V=\diag(U,U^*)\), and the projection \(P=\diag(1_n,0)\) satisfies \(\tilde{\varphi}(P)=V\diag(1_n,0)V^*=\diag(1_n,0)\).
    Thus \(\partial_1(K_1(\psi)({[U]}_1))=\partial_1({[\tilde{\psi}(U)]}_1)={[1_n]}_0-{[s(1_n)]}_0=0\).
    \item We show that \(\textnormal{Im}(K_1(\psi))\supseteq\textnormal{Ker}(\partial_1)\). To simplify notation, we identify \(J\) with its image in \(A\) and thus put \(\varphi=\textnormal{id}\).
    Let \(u\in\mathcal{U}_n(\tilde{B})\) be such that \({[u]}_1\in\textnormal{Ker}(\partial_1)\). By Exercise \ref{ex:7315}, there is a partial isometry \(U\in M_{2n}(\tilde{A})\) such that
    \[
      \tilde{\psi}(U)=\begin{pmatrix}
        u&0\\0&0
      \end{pmatrix}
    \]
    and 
    \[
      0=\partial_1({[u]}_1)={[1_{2n}-U^*U]}_0-{[1_{2n}-UU^*]}_0\text{ in }K_0(J).
    \]
    Thus there is \(k\) and \(w\in M_{2n+k}(\tilde{J})\) such that
    \[
      w^*w=(1_{2n}-U^*U)\oplus 1_k\text{ and }ww^*=(1_{2n}-UU^*)\oplus 1_k.
    \]
    Hence
    \[
      \tilde{\psi}(w^*w)=\tilde{\psi}(ww^*)=\begin{pmatrix}
        0&0\\0&1_{n+k}
      \end{pmatrix}
    \]
    and \(\tilde{\psi}(w)\) is a scalar matrix, since \(w\in M_{2n+k}(\tilde{J})\). More precisely,
    \(\tilde{\psi}(w)=\diag(0_n,z)\) for some scalar unitary \(z\in M_{n+k}(\mathbb{C})\). Consequently,
    set
    \[
      V=w+\begin{pmatrix}
        U&0\\0&0_k
      \end{pmatrix}\in M_{2n+k}(\tilde{A}).
    \]
    The two summands have complementary initial and final projections, so \(V\) is unitary. Moreover,
    \[
      \tilde{\psi}(V)=\begin{pmatrix}
        u&0\\0&0
      \end{pmatrix}
      +
      \begin{pmatrix}
        0&0\\0&z
      \end{pmatrix}
      \sim_h\begin{pmatrix}
        u&0\\0&1_{n+k}
      \end{pmatrix}.
    \]
    Thus \({[u]}_1={[\tilde{\psi}(V)]}_1=K_1(\psi)({[V]}_1)\).
    \item We show that \(\textnormal{Im}(\partial_1)\subseteq\textnormal{Ker}(K_0(\varphi))\).
    Indeed, let \(u\in\mathcal{U}_n(\tilde{B})\) and let \(V\in\mathcal{U}_{2n}(\tilde{A})\), \(P\in\mathcal{P}_{2n}(\tilde{J})\)
    be such that (\ref{eq:727}) and (\ref{eq:728}) hold. Then we have
    \[
      K_0(\varphi)({[P]}_0-{[s(P)]}_0)={[\tilde{\varphi}(P)]}_0-{[1_n]}_0={[V\diag(1_n,0)V^*]}_0-{[1_n]}_0=0.
    \]
    \item We show that \(\textnormal{Im}(\partial_1)\supseteq\textnormal{Ker}(K_0(\varphi))\). Let \(g\in\textnormal{Ker}(K_0(\varphi))\).
    By Lemma \ref{lemma:Lars723}, there is \(n\), a projection \(p\in\mathcal{P}_n(\tilde{J})\) and a unitary
    \(w\in\mathcal{U}_n(\tilde{A})\) such that
    \[
      g={[p]}_0-{[s(p)]}_0\text{ and }w\tilde{\varphi}(p)w^*=s(\tilde{\varphi}(p))=\tilde{\varphi}(s(p)).
    \]
    Set \(u_0=\tilde{\psi}(w(1_n-\tilde{\varphi}(p)))\), a partial isometry in \(M_n(\tilde{B})\). We have
    \begin{align*}
      u_0^*u_0&=1_n-\tilde{\psi}(\tilde{\varphi}(p)),\\
      u_0u^*_0&=1_n-\tilde{\psi}(s(p))=u_0^*u_0.
    \end{align*}
    Thus \(u=u_0+(1_n-u_0u_0^*)\) is unitary in \(M_n(\tilde{B})\). We want to show that \(g=\partial_1({[u]}_1)\).
    To this end, we first find a lift of \(\diag(u,0)\) to a suitable partial isometry in \(M_{2n}(\tilde{A})\). Take
    \[
      V_0=\begin{pmatrix}
        w(1_n-\tilde{\varphi}(p))&0\\0&s(p)
      \end{pmatrix},
    \]
    a partial isometry in \(M_{2n}(\tilde{A})\) such that
    \[
      \tilde{\psi}(V_0)=\begin{pmatrix}
        u_0&0\\0&s(p)
      \end{pmatrix}.
    \]
    Set 
    \[
      Z=\begin{pmatrix}
        1_n-s(p)&s(p)\\s(p)&1_n-s(p)
      \end{pmatrix},
    \]
    a self-adjoint, unitary scalar matrix, and put \(V=ZV_0Z^*\). Then we have
    \[
      \tilde{\psi}(V)=Z\tilde{\psi}(V_0)Z^*=Z\begin{pmatrix}
        u_0&0\\0&s(p)
      \end{pmatrix}
      Z^*=\begin{pmatrix}
        u&0\\0&0
      \end{pmatrix}.
    \]
    Hence, by Exercise \ref{ex:7315},
    \begin{align*}
      \partial_1({[u]}_1)={[\tilde{\varphi}^{-1}(1_{2n}-V^*V)]}_0-{[\tilde{\varphi}^{-1}(1_{2n}-VV^*)]}_0=\\
      {[\tilde{\varphi}^{-1}(1_{2n}-V_0^*V_0)]}_0-{[\tilde{\varphi}^{-1}(1_{2n}-V_0V_0^*)]}_0=\\
      {[\begin{pmatrix}
        p&0\\0&1_n-s(p)
      \end{pmatrix}]}_0-
      {[\begin{pmatrix}
        s(p)&0\\0&1_n-s(p)
      \end{pmatrix}]}_0=\\
      {[p]}_0-{[s(p)]}_0=g.
    \end{align*}
    That is, \(g=\partial_1({[u]}_1)\), as required.
  \end{enumerate}
\end{proof}

The preceding construction gives a complete computation of the odd \(K\)-group of the
Calkin algebra; compare \cite[Example~9.4.3]{RordamLarsenLaustsen2000}.
\begin{proposition}\label{prop:K1Calkin}
  For the extension
  \[
    0\longrightarrow\mathcal{K}\longrightarrow\mathcal{B}(\mathcal{H})
      \overset{\pi}{\longrightarrow}\mathcal{Q}\longrightarrow 0,
  \]
  the index map is an isomorphism
  \[
    \partial_1:K_1(\mathcal{Q})\overset{\cong}{\longrightarrow}
      K_0(\mathcal{K})\cong\mathbb{Z}.
  \]
  Under the rank isomorphism \(K_0(\mathcal{K})\cong\mathbb{Z}\), this map is the
  Fredholm index. In particular, \(K_1(\mathcal{Q})\cong\mathbb{Z}\).
\end{proposition}
\begin{proof}
  First, \(K_1(\mathcal{B}(\mathcal{H}))=0\). Indeed, after identifying
  \(M_n(\mathcal{B}(\mathcal{H}))\) with \(\mathcal{B}(\mathcal{H}^{\oplus n})\), every
  unitary \(u\) can be written as \(u=e^{ih}\) for some bounded self-adjoint operator
  \(h\), by bounded Borel functional calculus. Hence \(t\mapsto e^{ith}\) joins \(1\)
  to \(u\).

  We also have \(K_0(\mathcal{B}(\mathcal{H}))=0\). To see this, decompose
  \(\mathcal{H}=\mathcal{H}_1\oplus\mathcal{H}_2\), where both summands are infinite
  dimensional. If \(e_j\) is the projection onto \(\mathcal{H}_j\), then
  \(e_j\sim 1\), and therefore
  \([1]_0=[e_1]_0+[e_2]_0=2[1]_0\), so \([1]_0=0\). If \(e\) has rank one, then
  \(1-e\sim 1\), whence \([e]_0=0\). It follows that every finite-rank projection has
  zero class. Every infinite-rank projection in a matrix algebra over
  \(\mathcal{B}(\mathcal{H})\) is Murray--von Neumann equivalent to the identity of
  \(\mathcal{B}(\mathcal{H})\), after the standard identification
  \(\mathcal{H}^{\oplus n}\cong\mathcal{H}\), and thus also has zero class.

  The exact sequence of Theorem~\ref{theorem:727} contains
  \[
    K_1(\mathcal{B}(\mathcal{H}))\longrightarrow K_1(\mathcal{Q})
      \overset{\partial_1}{\longrightarrow}K_0(\mathcal{K})
      \longrightarrow K_0(\mathcal{B}(\mathcal{H})).
  \]
  Since the two outer groups vanish, \(\partial_1\) is an isomorphism. Finally,
  \(K_0(\mathcal{K})\cong\mathbb{Z}\) through the rank of a finite-rank projection.
  If \(u\in M_n(\mathcal{Q})\) is unitary and \(U\) is a partial-isometry lift, then
  \[
    \partial_1({[u]}_1)
      ={[1-U^*U]}_0-{[1-UU^*]}_0.
  \]
  The ranks of the two defect projections are respectively
  \(\dim(\textnormal{Ker}(U))\) and \(\dim(\textnormal{Coker}(U))\); hence the rank of this class is
  \(\Index(U)\), as claimed.
\end{proof}

\section{Examples and Exercises}

\begin{exercise}\label{ex:731}
  Let \(A\) be a unital \cstar. We have \(\tilde{A}=A\oplus\mathbb{C}f\), where \(f=1_{\tilde{A}}-1_A\).
  Define a unital *-homomorphism \(\mu:\tilde{A}\rightarrow A\) by \(\mu(a+\lambda f)=a\).
  As usual, for each \(n\) extend \(\mu\) to a unital *-homomorphism \(M_n(\tilde{A})\rightarrow M_n(A)\) (still denoted \(\mu\)).
  This yields a map \(\mu:\mathcal{U}_\infty(\tilde{A})\rightarrow\mathcal{U}_\infty(A)\). 
  Show that there exists an isomorphism \(K_1(A)\rightarrow\mathcal{U}_\infty(A)/\sim_1\) making the diagram
  \begin{equation}
    \begin{tikzcd}
      \mathcal{U}_\infty(\tilde{A}) \arrow[r,"\mu"] \arrow[d,swap,"{[\cdot]}_1"]&\mathcal{U}_\infty(A)\arrow[d]\\
      K_1(A)\arrow[r]&\mathcal{U}_\infty(A)/\sim_1
    \end{tikzcd}
  \end{equation}
  commutative. To this end, show the following:
  \begin{enumerate}
    \item \(\mu(\diag(u,v))=\diag(\mu(u),\mu(v))\),
    \item if \(u,v\in\mathcal{U}_n(\tilde{A})\) then \(\mu(u)\sim_1\mu(v)\) if and only if \(u\sim_h v\),
    \item if \(u,v\in\mathcal{U}_\infty(\tilde{A})\) then \(\mu(u)\sim_1\mu(v)\) if and only if \(u\sim_1 v\).
  \end{enumerate}
\end{exercise}

\begin{exercise}\label{ex:732}
  Show the following.
  \begin{enumerate}
    \item \(K_1(\mathbb{C})=0\).
    \item For any two \cstars\ \(A,B\) we have \(K_1(A\oplus B)\cong K_1(A)\oplus K_1(B)\).
    \item If \(A\) is an AF-algebra then \(K_1(A)=0\).
  \end{enumerate}
\end{exercise}

\begin{exercise}\label{ex:733}
  If \(\mathcal{H}\) is an infinite dimensional Hilbert space then \(K_1(\mathcal{B}(\mathcal{H}))=0\).
  Indeed, since \(\mathcal{U}_n(\mathcal{B}(\mathcal{H}))\cong\mathcal{U}(\mathcal{B}\oplus^n\mathcal{H})\), it suffices to show that every unitary 
  in \(\mathcal{B}(\mathcal{H})\) is homotopic to the identity. But this follows from the fact that for every unitary \(u\in\mathcal{B}(\mathcal{H})\)
  there is a self-adjoint \(a\in\mathcal{B}(\mathcal{H})\) such that \(u=e^{ia}\). Indeed, one may take \(a=\varphi(u)\),
  where \(\varphi:\mathbb{T}\rightarrow[0,2\pi)\) is a bounded Borel function such that \(\varphi(e^{i\theta})=\theta\).
\end{exercise}

\begin{exercise}\label{ex:734}
  Let \(X\) be a compact Hausdorff space.
  \begin{enumerate}
    \item For each \(n\) identify \(M_n(C(X))\) with \(C(X,M_n(\mathbb{C}))\) and define the \textnormal{determinant} function
    \(\det:M_n(C(X))\rightarrow C(X)\). Show that \(\det\) maps \(\mathcal{U}_\infty(C(X))\) into \(\mathcal{U}(C(X))\).
    \item Let \(\langle v\rangle\) denote the class of \(v\in\mathcal{U}(C(X))\) in \(\mathcal{U}(C(X))/\mathcal{U}_0(C(X))\).
    Apply the universal property of \(K_1\) to the map \(\mathcal{U}_\infty(C(X))\ni u\mapsto \langle\det(u)\rangle\) to show that there exists a homomorphism
    \(D:K_1(C(X))\rightarrow \mathcal{U}(C(X))/\mathcal{U}_0(C(X))\) such that \(D({[u]}_1)=\langle\det(u)\rangle\).
    \item Show that the sequence
    \[
      \begin{tikzcd}
      0\arrow[r]&\textnormal{Ker}(D) \arrow[r]&K_1(C(X)) \arrow[r,"D"]&\mathcal{U}(C(X))/\mathcal{U}_0(C(X))\arrow[r]&0
      \end{tikzcd}  
    \]
    is split-exact, with a splitting map \(\omega:\mathcal{U}(C(X))/\mathcal{U}_0(C(X))\rightarrow K_1(C(X))\)
    given by \(\omega(\langle u\rangle)={[u]}_1\).
    \item Let \(X=\mathbb{T}\). Recall that \(\varphi:\mathbb{R}\rightarrow\mathbb{T}\), \(\varphi(x)=e^{2\pi ix}\), is a covering map.
    Thus, if \(u\in\mathcal{U}(C(\mathbb{T}))\) then there exists a continuous function \(f:[0,1]\rightarrow\mathbb{R}\)
    such that \(u(e^{2\pi it})=e^{2\pi if(t)}\). If \(f,g\) are two such functions then \(f-g\) is a constant integer.
    Thus, there is a well-defined map \(\mu:\mathcal{U}(C(\mathbb{T}))\rightarrow\mathbb{Z}\) given by \(\mu(u)=f(1)-f(0)\)
    (the \textnormal{winding number} of \(u\)). Show that \(\mu\) induces an isomorphism of \(\mathcal{U}(C(\mathbb{T}))/\mathcal{U}_0(C(\mathbb{T}))\)
    and \(\mathbb{Z}\) such that \(\langle u\rangle\mapsto\mu(u)\).
    \item Conclude that there exists a surjective homomorphism from \(K_1(C(\mathbb{T}))\) onto \(\mathbb{Z}\).
    In fact, we will see later that \(K_1(C(\mathbb{T}))\cong\mathbb{Z}\).
  \end{enumerate}
\end{exercise}

\begin{exercise}\label{ex:735}
  If \(A\) is a separable \cstar\ then \(K_1(A)\) is countable.
\end{exercise}

\begin{exercise}\label{ex:736}
   Let \(A\) be a unital \cstar. Replacing unitaries \(\mathcal{U}_\infty(A)\) with invertibles
  \(GL_\infty(A)\) one can repeat the constructions and define an abelian group
  \(GL_\infty(A)/\sim_1\). Show that this group is isomorphic to \(K_1(A) = \mathcal{U}_\infty(A)/\sim_1\) 
  (see Exercise~\ref{ex:731}). Hint: For \(w \in GL_\infty(A)\) let \(w = u|w|\) be the polar decomposition. Define a map
  \({[\cdot]}_1 : GL_\infty(A) \rightarrow K_1(A)\) by \({[w]}_1 = {[u]}_1 = {[w|w|^{-1}]}_1\) and use Proposition~\ref{prop:5111}.
\end{exercise}

\begin{exercise}\label{ex:737}
  Let \(A\) be a non-unital \cstar, and let \(s:\tilde{A}\rightarrow\tilde{A}\) be the scalar map
  \(s(a+t1)=t1\). Define
  \begin{align*}
    \mathcal{U}^+(A)&=\{u\in\mathcal{U}(\tilde{A})| s(u)=1\},\\
    \mathcal{U}^+_n(A)&=\{u\in\mathcal{U}(M_n(\tilde{A}))| s_n(u)=1_n\},\\
    \mathcal{U}^+_\infty(A)&=\bigcup_{n=1}^\infty \mathcal{U}_n^+(A).
  \end{align*}
  One can define an abelian group \(\mathcal{U}^+_\infty(A)/\sim_1\) with the usual construction. Show that this group is isomorphic to \(K_1(A)\).
\end{exercise}
\begin{exercise}\label{ex:738}
  Let \(A\) be a unital \cstar.
  \begin{enumerate} 
    \item Let \(u\) be unitary and let \(s\) be an isometry in \(A\). Then \(sus^*+(1-ss^*)\) is unitary
    and we have
    \[
    \begin{pmatrix}
      s&1-ss^*\\0&s^*  
    \end{pmatrix}
    \begin{pmatrix}
      u&0\\0&1
    \end{pmatrix}
    {\begin{pmatrix}
      s&1-ss^*\\0&s^*
    \end{pmatrix}}^*=
    \begin{pmatrix}
      sus^*+(1-ss^*)&0\\0&1
    \end{pmatrix}.
    \]
    Thus \({[sus^*+(1-ss^*)]}_1={[u]}_1\).
  \item Let \(u_1,\ldots,u_n\) be unitary elements in \(A\) and let \(s_1,\ldots,s_n\) be 
  isometries in \(A\) with mutually orthogonal range projections. Then
  \[
    u=s_1u_1s^*_1+\cdots+s_n u_n s_n^*+(1-s_1 s_1^*-\cdots-s_n s_n^*)
  \]
  is unitary. Use (1) to show that \({[u]}_1={[u_1]}_1{[u_2]}_1\cdots{[u_n]}_1\).
  \item Let \(s_1,\ldots,s_n\) be isometries as in (2). Put
  \[
    V=\begin{pmatrix}
      s_1&s_2&\cdots& s_n\\
      0&0&\cdots &0\\
      \vdots&\vdots&\ddots&\vdots\\
      0&0&\cdots&0
    \end{pmatrix}.
  \]
  Then \(V\) is an isometry in \(M_n(A)\). Show that for any unitary \(u\in\mathcal{U}_n(A)\)
  there is a unitary \(w\) in \(A\) such that
  \[
    VuV^*+(1_n-VV^*)=\diag(w,1_{n-1}).
  \]
  \item Let \(A\) be \textnormal{properly infinite}\footnote{
    Let \(A\) be a unital \cstar. \(A\)
    is called \textit{properly infinite} if there exist two projections \(e,f\) in \(A\) such that \(ef = 0\) and
    \(1 \sim e \sim f\). For example, Cuntz algebras are properly infinite.
  }. Show that
  \[
    K_1(A)=\{{[u]}_1| u\in\mathcal{U}(A)\}.
  \]
  \end{enumerate}
\end{exercise}

\begin{exercise}\label{ex:739}
  The \(K_1\)-functor is not exact. Indeed, for a separable Hilbert space \(\mathcal{H}\) the sequence
  \[
      \begin{tikzcd}
      0\arrow[r]&\mathcal{K}\arrow[r]&\mathcal{B}(\mathcal{H}) \arrow[r,"\pi"]&\mathcal{B}(\mathcal{H})/\mathcal{K} \arrow[r]&0
      \end{tikzcd}   
  \]
  of \cstars\ is exact. But \(K_1(\mathcal{B}(\mathcal{H}))=0\) and we will see later that \(K_1(\mathcal{B}(\mathcal{H})/\mathcal{K})\cong\mathbb{Z}\).
  Thus \(K_1(\pi)\) cannot be surjective. Likewise, there is an exact sequence
  \[
      \begin{tikzcd}
      0\arrow[r]&C_0((0,1)) \arrow[r,"\varphi"]&C([0,1]).
      \end{tikzcd}   
  \]
  But \(K_1(C([0,1]))=0\) and we will see later that \(K_1(C_0((0,1)))\cong\mathbb{Z}\).
  Thus \(K_1(\varphi)\) cannot be injective.
\end{exercise}

\begin{exercise}\label{ex:7310}
  Let \(A=\varinjlim\{A_i,\Phi_{ij}\}\) be the inductive limit of a sequence of \cstars,
  and let \(\Phi_i:A_i\rightarrow A\) be the canonical maps.
  \begin{enumerate}
    \item For any invertible \(y\in\tilde{A}\) and any \(\epsilon>0\) there is arbitrarily large \(i\) and invertible \(z\in \tilde{A_i}\)
      such that \(\left\|y-\widetilde{\Phi_i}(z)\right\|<\epsilon\).\label{ex:73101}
    \item For any unitary \(u\in\tilde{A}\) and any \(\epsilon>0\) there is arbitrarily large \(i\) and unitary \(w\in\tilde{A_i}\) such that \(\left\|u-\widetilde{\Phi_{i}}(w)\right\|<\epsilon\).\label{ex:73102}
    \item If \(u\) is unitary in \(\tilde{A_j}\) such that \(\widetilde{\Phi_j}(u)\sim_h 1\) in \(\tilde{A}\), then there is arbitrarily large \(i\) such that \(\widetilde{\Phi_{ij}}(u)\sim_h 1\) in \(\tilde{A_i}\).\label{ex:73103}
    \item Previous parts remain valid with \(\tilde{A}\) and \(\tilde{A_i}\) replaced by \(M_n(\tilde{A})\) and \(M_n(\tilde{A_i})\), respectively.\label{ex:73104}
  \end{enumerate}
\end{exercise}
\begin{proof}
  \begin{enumerate}
    \item First find \(k\) and \(x,x'\in \tilde{A}_k\) so that both \(\|\widetilde{\Phi_k}(x) - y\|\) and 
    \(\|\tilde{\Phi_k}(x') - y^{-1}\|\) are small.
    Thus both \(\|\tilde{\Phi_k}(xx'-1)\|\) and \(\|\tilde{\Phi_k}(x'x-1)\|\) are small. Then, 
    using Proposition~\ref{prop:Lars624} (2), take \(i\)
    large enough so that both \(\|\tilde{\Phi_{ik}}(xx'-1)\|\) and \(\|\tilde{\Phi_{ik}}(x'x -1)\|\) are small. 
    Then \(z = \tilde{\Phi_{ik}}(x)\)
    is both left and right invertible, hence invertible, and \(\tilde{\Phi_i}(z)\) approximates \(y\).
    \item This follows from part (1) and continuity of the polar decomposition (see Proposition~\ref{prop:5111}).
    \item Let \(w_t, t \in [0,1]\), be a continuous path of unitaries in \(\tilde{A}\) connecting \(w_0 = \tilde{\Phi_j}(u)\)
    and \(w_1 = 1\). By compactness, there are \(0 = t_0 < t_1 < \cdots < t_{k+1} = 1\) such that
    \(\|w_{t_{r+1}}-w_{t_r}\| < 2\) for all \(r\). Applying repeatedly part (2), find \(m \geq j\) and unitary elements
    \(v_1, \ldots, v_k\) in \(\tilde{A_m}\) so close to \(w_{t_1},\ldots,w_{t_k}\), respectively, 
    that all the norms: \(\|\tilde{\Phi_j}(u)-\tilde{\Phi_m}(v_1)\|\),
    \(\|\tilde{\Phi_m}(v_k) - 1\|\), and \(\|\tilde{\Phi_m}(v_{r+1}) - \tilde{\Phi_m}(v_r)\|\) for \(r = 1,\ldots,k - 1\) 
    are less than \(2\). Then by
    Proposition~\ref{prop:Lars624} (2), there is arbitrarily large \(i \geq m\) such that all the norms \(\|\tilde{\Phi_{ij}}(u)- \tilde{\Phi_{im}}(v_1)\|\),
    \(\|\tilde{\Phi_{im}}(v_k) - 1\|\), and \(\|\tilde{\Phi_{im}}(v_{r+1}) - \tilde{\Phi_{im}}(v_r)\|\) for \(r = 1,\ldots,k - 1\) are less than \(2\). 
    Now the claim follows from Lemma~\ref{lemma:Matthes214}.
    \item Exercise.
  \end{enumerate}
\end{proof}
\begin{exercise}\label{ex:7311}
  Show that every unitary in the Calkin algebra \(\mathcal{Q}\) lifts to a partial
  isometry in \(\mathcal{B}(\mathcal{H})\). In fact, it can be lifted to an isometry or a coisometry.
\end{exercise}
\begin{exercise}\label{ex:7312}
  Let \(\psi : A \rightarrow B\) be a surjective *-homomorphism of \cstars. Show
  the following.
  \begin{enumerate}
    \item For each \(b = b^* \in B\) there is \(a = a^* \in A\) such that \(\|a\| = \|b\|\) and \(\psi(a) = b\).
    \item For each \(b \in B\) there is \(a \in A\) such that \(\|a\| = \|b\|\) and \(\psi(a) = b\).    
  \end{enumerate}
\end{exercise}
\begin{exercise}\label{ex:7313}
  Consider an exact sequence of \cstars
  \[
    \begin{tikzcd}
      0\arrow[r]&J \arrow[r]&A\arrow[r,"\psi"]&B\arrow[r]&0,
    \end{tikzcd}  
  \]
  in which we identify \(J\) with its image in \(A\). Let \(u\) be a unitary in \(\mathcal{U}_n(\tilde{B})\).
  By part (2) of Exercise~\ref{ex:7312} there is \(a\in M_n(\tilde{A})\) such that \(\tilde{\psi}(a)=u\)
  and \(\|a\|=\|u\|=1\). Then for any continuous function \(f:\mathbb{R}\rightarrow\mathbb{C}\) we have \(af(a^* a)=f(a a^*)a\).
  Use this to show that
  \[
    V=\begin{pmatrix}
    a&{(1_n-a a^*)}^{1/2}\\
    -{(1_n-a^* a)}^{1/2}&a^*  
    \end{pmatrix}
  \]
  is a unitary in \(\mathcal{U}_{2n}(\tilde{A})\). Then show that
  \[
    \tilde{\psi}(V)=\begin{pmatrix}
      u&0\\0&u^*
    \end{pmatrix}
  \]
  and 
  \[
    V\begin{pmatrix}
      1_n&0\\0&0
    \end{pmatrix}V^*=
    \begin{pmatrix}
      a a^*&-a{(1_n-a^* a)}^{1/2}\\
      -{(1_n-a^* a)}^{1/2}a^*&1_n-a^* a
    \end{pmatrix}.
  \]
  Then write explicitly \(\partial_1({[u]}_1)\).
\end{exercise}
\begin{exercise}\label{ex:7314}
  Consider an exact sequence of \cstars\
  \[
    \begin{tikzcd}
      0\arrow[r]&J \arrow[r]&A\arrow[r,"\psi"]&B\arrow[r]&0,
    \end{tikzcd} 
  \]
  in which we identify \(J\) with its image in \(A\). Suppose that \(u\in\mathcal{U}_n(\tilde{B})\) is such that
  there exists a \textit{partial isometry} \(U\in M_n(\tilde{A})\) with \(\tilde{\psi}(U)=u\). Show the following.
  \begin{enumerate}
    \item The element 
    \[
      V=\begin{pmatrix}
        U&1_n-U U^*\\
        1_n-U^* U&U^*
      \end{pmatrix}
    \]
    is unitary in \(\mathcal{U}_{2n}(\tilde{A})\) and \(\tilde{\psi}(V)=\diag(u,u^*)\).
    \item Both \(1_n-U^* U\) and \(1_n-U U^*\) are projections in \(M_n(J)\), and 
    \[\partial_1({[u]}_1)={[1_n-U^* U]}_0-{[1_n-U U^*]}_0.\]
  \end{enumerate}
\end{exercise}
\begin{exercise}\label{ex:7315}
  Consider an exact sequence of \cstars\
  \[
    \begin{tikzcd}
      0\arrow[r]&J \arrow[r]&A\arrow[r,"\psi"]&B\arrow[r]&0,
    \end{tikzcd} 
  \]
  in which we identify \(J\) with its image in \(A\). Let \(u\in\mathcal{U}_n(\tilde{B})\), and let \(a\in M_n(\tilde{A})\) be such that
  \(\tilde{\psi}(a)=u\) and \(\|a\|=\|u\|=1\). Put
  \[
    U=\begin{pmatrix}
      a&0\\
      {(1_n-a^* a)}^{1/2}&0
    \end{pmatrix}.
  \]  
  Show that \(U^* U=\diag(1_n,0)\), which entails that \(U\) is a partial isometry. Then show that
  \[
    \tilde{\psi}(U)=\begin{pmatrix}
      u&0\\0&0
    \end{pmatrix}.
  \]
  Finally, show that 
  \[
    \partial_1({[u]}_1)={[1_{2n}-U^* U]}_0-{[1_{2n}-U U^*]}_0.
  \]
\end{exercise}
\begin{exercise}\label{ex:7316}
  Consider an exact sequence
  \[
    \begin{tikzcd}
      0\arrow[r]&C_0(\mathbb{R}^2) \arrow[r]&C(\mathbb{D})\arrow[r]&C(S^1)\arrow[r]&0,
    \end{tikzcd} 
  \]
  with the map \(C(\mathbb{D})\rightarrow C(S^1)\) given by the restriction. In the corresponding exact sequence
  \begin{equation}\label{seq:7311}
    \begin{tikzcd}
      K_0(C_0(\mathbb{R}^2)) \arrow[r]&K_0(C(\mathbb{D}))\arrow[r]&K_0(C(S^1))\\
      K_1(C(S^1))\arrow[u,"\partial_1"]&K_1(C(\mathbb{D}))\arrow[l]&K_1(C_0(\mathbb{R}^2))\arrow[l]
    \end{tikzcd}  
  \end{equation}
  we have \(K_1(C(\mathbb{D}))=0\), since \(C(\mathbb{D})\) and \(\mathbb{C}\) are homotopy equivalent. Show that the map 
  \(K_0(C(\mathbb{D}))\rightarrow K_0(C(S^1))\) is injective, and conclude that \(\partial_1:K_1(C(S^1))\rightarrow K_0(C_0(\mathbb{R}^2))\) is an isomorphism.
  Then calculate \(\partial_1({[z]}_1)\) and thus find a generator of \(K_0(C_0(\mathbb{R}^2))\).
\end{exercise}
\begin{exercise}[Naturality of the index map]\label{ex:7317}
  Show that every commutative diagram of \cstars, with exact rows,
  \[
    \begin{tikzcd}
    0\arrow[r]&J_1 \arrow[r]\arrow[d]&A_1\arrow[r]\arrow[d]&B_1\arrow[r]&0\\
    0\arrow[r]&J_2 \arrow[r]&A_2\arrow[r]&B_2\arrow[r]&0
    \end{tikzcd}
  \]
  induces a commutative diagram of abelian groups, with exact rows:
  \[
    \begin{tikzcd}
    K_1(J_1) \arrow[r]\arrow[d]&K_1(A_1)\arrow[r]\arrow[d]&K_1(B_1)\arrow[r,"\partial_1"]\arrow[d]&K_0(J_1) \arrow[r]\arrow[d]&K_0(A_1)\arrow[r]\arrow[d]&K_0(B_1)\arrow[d]\\
    K_1(J_2) \arrow[r]&K_1(A_2)\arrow[r]&K_1(B_2)\arrow[r,"\partial_1"]&K_0(J_2) \arrow[r]&K_0(A_2)\arrow[r]&K_0(B_2)
    \end{tikzcd}
  \]
\end{exercise}
\begin{exercise}
  Show that if partial isometries \(v_1,\ldots,v_n\) in a unital \cstar\ satisfy
  \[
    v_1v_1^*+\cdots+v_n v_n^*=1=v_1^*v_1+\cdots+v_n^*v_n
  \]
  then \(u=v_1+\cdots+v_n\) is unitary. Hint: use Exercise~\ref{ex:24}.
\end{exercise}
  
\chapter{Exact sequences in K-theory}
\section{Higher \texorpdfstring{\(K\)}{K}-Groups}
The \textit{suspension} \(S A\) of a \cstar\ \(A\) is defined as
\[
  S A=\{f\in C([0,1],A)| f(0)=f(1)=0\},
\]

\begin{lemma}
  Let \(X\) be a locally compact Hausdorff space and let \(A\) be a \cstar. Define for \(f\in C_0(X)\), \(a\in A\) an element
  \(fa\in C_0(X,A)\) by
  \[
    (fa)(x)=f(x)a.
  \]
  Then \(\textnormal{span}\{fa| f\in C_0(X), a\in A\}\) is dense in \(C_0(X,A)\).
\end{lemma}
\begin{proof}
  Let \(X^+=X\cup\{\infty\}\) be the one-point compactification of \(X\). Then 
  \[
    C_0(X,A)=\{f\in C(X^+,A)|f(\infty)=0\}.
  \]
  Let \(f\in C_0(X,A)\), \(\epsilon>0\). There is an open covering \(U_1,\ldots, U_n\) of \(X^+\) such that \(\left\|f(x)-f(y)\right\|<\epsilon\) if \(x,y\in U_k\).
  (Compactness of \(X^+\), continuity of \(f\).) Choose \(x_k\in U_k\), with \(x_k=\infty\) if \(\infty\in U_k\).
  Let \({(h_k)}_{k=1}^n\) be a partition of unity subordinate to the covering \(U_k\), i.e., 
  \(h_k\in C(X^+)\), \(\textnormal{supp} h_k\subseteq U_k\), \(\sum_{k=1}^n h_k=1\), \(0\leq h_k\leq 1\).
  (Note that every compact Hausdorff space is paracompact.) Then \(\left\|f(x)h_k(x)-f(x_k)h_k(x)\right\|\leq \epsilon h_k(x)\)
  for \(x\in X\), \(k=1,\ldots, n\). It follows that \(\left\|f(x)-\sum_{k=1}^n f(x_k)h_k(x) \right\|\leq\epsilon\), for \(x\in X\).
  Put \(a_k=f(x_k)\in A\). Then \(\sum_{k=1}^n h_k a_k\in \textnormal{span}\{fa|f\in C_0(X),a\in A\}\), because
  \(a_k=f(x_k)=0\) if \(\infty \in U_k\), and \(\left\|f-\sum_{k=1}^n h_k a_k\right\|\leq \epsilon\).
\end{proof}
\begin{lemma}
  The suspension \(S A\) is isomorphic to \(C_0((0,1),A)\cong C_0(\mathbb{R})\otimes A\).
\end{lemma}
If \(\varphi:A\rightarrow B\) is a *-homomorphism between two \cstars, then \(S\varphi:S A\rightarrow S B\), given by \(S\varphi (f)(t)=\varphi(f(t))\), is a *-homomorphism between their suspensions.
It is not difficult to verify that this yields a covariant functor from the category of \cstars\ into itself.
\begin{proposition}
  The suspension functor \(S\) is exact. That is, if 
  \[
    \begin{tikzcd}
      0\arrow[r]&J \arrow[r,"\varphi"]&A\arrow[r,"\psi"]&B\arrow[r]&0
    \end{tikzcd}  
  \]
  is an exact sequence of \cstars\ then the sequence
  \[
    \begin{tikzcd}
      0\arrow[r]&S J \arrow[r,"S\varphi"]&S A\arrow[r,"S\psi"]&S B\arrow[r]&0
    \end{tikzcd}  
  \]
  is exact.
\end{proposition}
Let \(A\) be a \cstar. We define a map
\[
  \theta_A:K_1(A)\rightarrow K_0(S A),
\]
as follows. By Exercise \ref{ex:841}, each element of \(K_1(A)\) is represented by a unitary \(u\in \mathcal{U}_n(\tilde{A})\)
(for some \(n\)) such that \(s(u)=1_n\). For such a \(u\) we can find a continuous function
\(v:[0,1]\rightarrow\mathcal{U}_{2n}(\tilde{A})\) such that \(v(0)=1_{2n}\), \(v(1)=\diag(u,u^*)\) and \(s(v(t))=1_{2n}\) for all \(t\in [0,1]\). We put
\[
  p=v\begin{pmatrix}
    1_n&0\\0&0_n
  \end{pmatrix}v^*,
\]
  a projection in \(\mathcal{P}_{2n}(\widetilde{S A})\) with \(s(p)=\diag (1_n,0_n)\). We set
\[
  \theta_A({[u]}_1)={[p]}_0-{[s(p)]}_0.
\]
\begin{theorem}\label{theorem:812}
  For any \cstar\ \(A\), the map 
  \[
    \theta_A:K_1(A)\rightarrow K_0(S A),
  \]
  is an isomorphism. Furthermore, if \(B\) is a \cstar\ and \(\varphi:A\rightarrow B\) is a *-homomorphism then the diagram
  \[
    \begin{tikzcd}
      K_1(A)\arrow[r,"K_1(\varphi)"]\arrow[d,"\theta_A"]&K_1(B)\arrow[d,"\theta_B"]\\
      K_0(S A)\arrow[r,"K_0(S\varphi)"]&K_0(S B)
    \end{tikzcd}  
  \]
  is commutative.
\end{theorem}
\begin{proof}
  Recall from Example~\ref{example:449} the exact sequence
  \begin{equation}\label{seq:812}
    \begin{tikzcd}
      0\arrow[r]&S A \arrow[r]&C A\arrow[r,"\pi"]&A\arrow[r]&0,
    \end{tikzcd}
  \end{equation}
  where \(C A\) is the cone over \(A\). Since \(C A\) is homotopy equivalent to \(\{0\}\) we have 
  \(K_0(C A)=K_1(C A)=0\). Let \(\partial_1:K_1(A)\rightarrow K_0(S A)\) be the index map associated with the extension (\ref{seq:812}).
  It follows from Theorem~\ref{theorem:727} that \(\partial_1\) is an isomorphism. Thus, it suffices to identify \(\partial_1\) with \(\theta_A\).
\end{proof}

For each natural number \(n\geq 2\) we define inductively a covariant functor from the category of \cstars\ to the 
category of abelian groups as follows. \(K_n(A)=K_{n-1}(S A)\), and if \(\varphi:A\rightarrow B\) is a *-homomorphism then
\(K_{n}(\varphi)=K_{n-1}(S \varphi)\). It is clear that the functor defined in this way is half-exact.\par
Now suppose that
\begin{equation*}
  \begin{tikzcd}
    0\arrow[r]&J \arrow[r,"\varphi"]&A\arrow[r,"\psi"]&B\arrow[r]&0
  \end{tikzcd}
\end{equation*}
is an exact sequence of \cstars. Then \(\partial_1:K_1(B)\rightarrow K_{0}(J)\) is the index map.
We define higher index maps
\[
  \partial_n:K_n(B)\rightarrow K_{n-1}(J),
\]
as follows. Applying the suspension functor \(n-1\) times to the exact sequence, we obtain the exact sequence
\begin{equation}\label{seq:814}
  \begin{tikzcd}
    0\arrow[r]&S^{n-1} J \arrow[r,"S^{n-1}\varphi"]&S^{n-1}A\arrow[r,"S^{n-1}\psi"]&S^{n-1}B\arrow[r]&0
  \end{tikzcd}
\end{equation}  
Let \(\overline{\partial_1}:K_1(S^{n-1}B)\rightarrow K_0(S^{n-1}J)\) be the index map associated with (\ref{seq:814}).
By the definition of higher \(K\)-functors, we have \(K_n(B)=K_1(S^{n-1}B)\) and \(K_{n-1}(J)=K_1(S^{n-2}J)\).
By Theorem~\ref{theorem:812}, there is an isomorphism \(\theta_{S^{n-2}J}:K_1(S^{n-2}J)\rightarrow K_0(S^{n-1}J)\). We define
\[
  \partial_n=\theta^{-1}_{S^{n-2}J}\circ \overline{\partial_1}.
\]\par
The higher index maps defined in this way have naturality properties analogous to those enjoyed by the
usual index (cf. Exercise~\ref{ex:7317}).
\begin{proposition}
  Every short exact sequence 
  \begin{equation*}
  \begin{tikzcd}
    0\arrow[r]&J \arrow[r,"\varphi"]&A\arrow[r,"\psi"]&B\arrow[r]&0
  \end{tikzcd}
\end{equation*}
of \cstars\ induces a long exact sequence on \(K\)-theory:
\begin{equation}\label{seq:813}
  \begin{tikzcd}[column sep=small]
    \cdots\arrow[r,"\partial_{n+1}"]&K_n(J) \arrow[r,"K_n(\varphi)"]&K_n(A)\arrow[r,"K_n(\psi)"]&K_n(B)\arrow[r,"\partial_{n}"]&\cdots\arrow[r,"\partial_1"]&K_0(J)\arrow[r,"K_0(\varphi)"]&K_0(A)\arrow[r,"K_0(\psi)"]&K_0(B).
  \end{tikzcd}
\end{equation}  
\end{proposition}\par
This Proposition serves only as an intermediate step towards the fundamental 6-term exact sequence of \(K\)-theory.
Observe that \(K_{n+2}(A)\cong K_n(A)\) (as we will see in the next section), and the apparently infinite sequence (\ref{seq:813}) shrinks to a much more useful finite one, which contains only \(K_0\) and \(K_1\).

\section{Bott periodicity}
In this section we will prove the fundamental result of Bott that \(K_0(A)\cong K_1(S A)\) for any \cstar\ \(A\).
Combined with Theorem~\ref{theorem:812}, it says that \(K_{n+2}\cong K_n(A)\) (the \textit{Bott periodicity}).\par
We begin by defining a \textit{Bott map}
\[
  \beta_A:K_0(A)\rightarrow K_1(S A)
\]
for unital \cstars\ \(A\), and then reduce the general case to the unital one. So let \(A\) be a unital \cstar. We use the obvious identification
\[
  S A=\{f:\mathbb{T}\rightarrow A| f\textnormal{ continuous , } f(1)=0\}.
\]
Thus, elements of \(M_n(S A)\) may be identified with continuous \textit{loops} \(f:\mathbb{T}\rightarrow M_n(A)\)
such that \(f(1)=0\). It follows that \(M_n(\widetilde{S A})\) may be identified with continuous functions \(f:\mathbb{T}\rightarrow M_n(A)\)
  such that \(f(1)\in M_n(\mathbb{C}1_A)\).\par
For any natural \(n\) and any projection \(p\in\mathcal{P}_n(A)\) we define a \textit{unitary loop} \(f_p:\mathbb{T}\rightarrow \mathcal{U}_n(A)\) by
\[
  f_p(z)=zp+(1_n-p)\textnormal{, }z\in\mathbb{T}.
\]
Clearly, we have \(f_p\in\mathcal{U}_n(\widetilde{S A})\). By the universal property of \(K_0\) we get a homomorphism
\(\beta_A :K_0(A)\rightarrow K_1(S A)\) such that
\[
  \beta_A({[p]}_0)={[f_p]}_1,
\]
called the \textit{Bott map}.\par
Now, if \(\varphi:A\rightarrow B\) is a unital *-homomorphism, then \(\tilde{S\varphi}(f_p)=\varphi(f_p(z))=f_{\varphi(p)}(z)\)
for all \(z\in\mathbb{T}\). Hence the diagram
\begin{equation}\label{diag:825}
  \begin{tikzcd}
    K_0(A)\arrow[r,"K_0(\varphi)"]\arrow[d,"\beta_A",swap]&K_0(B)\arrow[d,"\beta_B"]\\
    K_1(S A)\arrow[r,"K_1(S\varphi)",swap]&K_1(S B)
  \end{tikzcd}
\end{equation}
is commutative. This is the \textit{naturality} of the Bott map.\par
Finally, suppose that \(A\) does not have a unit. Then we have a commutative diagram
\begin{equation}\label{diag:826}
  \begin{tikzcd}
    0\arrow[r]&K_0(A)\arrow[r]&K_0(\tilde{A})\arrow[d,"\beta_{\tilde{A}}"]\arrow[r]&K_0(\mathbb{C})\arrow[r]\arrow[d,"\beta_{\mathbb{C}}"]&0\\
    0\arrow[r]&K_1(S A)\arrow[r]&K_1(S \widetilde{A})\arrow[r]&K_1(S\mathbb{C})\arrow[r]&0
  \end{tikzcd}
\end{equation}
with split-exact rows. It follows that there is exactly one map \(\beta_A:K_0(A)\rightarrow K_1(S A)\)
which completes the diagram. By Exercise \ref{ex:842}, we have 
\[
  \beta_A({[p]}_0-{[s(p)]}_0)={[f_p f^*_{s(p)}]}_1.
\]\par
The following theorem is considered a central result of \(K\)-theory.
\begin{theorem}\label{theorem:821}
  For any \cstar\ \(A\), the Bott map
  \[
    \beta_A:K_0(A)\rightarrow K_1(S A)
  \]
  is an isomorphism.
\end{theorem}
\begin{proof}
  It suffices to prove the theorem for unital \cstars. Indeed, the general case follows from the unital one and (\ref{diag:826}) through a diagram chase.
  Thus assume \(A\) unital. It will be convenient for us to use the description of \(K_1(S A)\) as the collection of suitable equivalence classes in \(GL_\infty(\widetilde{S A})\) (see Exercise \ref{ex:736}).
  We must show that the Bott map \(\beta_A:K_0(A)\rightarrow K_1(S A)\) is both surjective and injective.
  \begin{enumerate}[label={}]
    \item \textbf{Surjectivity.} We consider the following subsets of \(GL_\infty(\widetilde{S A})\):
    \begin{align*}
      GL^n&=\{f:\mathbb{T}\rightarrow GL_n(A)| f\textnormal{ continuous and }f(1)\in M_n(\mathbb{C}1_A)\},\\
      LL_m^n&=\{f\in GL^n| f\textnormal{ a Laurent polynomial in }z\textnormal{ with coeff. in }M_n(A)\textnormal{ and }\deg(f)\leq m\},\\
      PL_m^n&=\{f\in LL_m^n| f\textnormal{ a polynomial}\},\\
      PRL^n&=\{f_p| p\in\mathcal{P}_n(A)\}.
    \end{align*}
    Elements of \(GL^n\), \(LL_m^n\), \(PL_m^n\) and \(PL_1^n\) are called \textnormal{invertible loops, Laurent loops, polynomial loops} and \textnormal{linear loops}, respectively. We have
    \[
      PRL^n\subseteq PL_1^n\subseteq \bigcup_m PL_m^n\subseteq\bigcup_m LL_m^n\subseteq GL^n
    \]
    and \(K_1(S A)=\{{[f]}_1| f\in GL^n, n\in\mathbb{N}\}\).
    \begin{enumerate}[label={}]
      \item \textbf{Step 1.} \(\bigcup_m LL_m^n\) is dense in \(GL^n\). Indeed, \(\textnormal{span}\{z^k| k\in\mathbb{Z}\}\)
      is dense in \(C(\mathbb{T})\) by the Stone-Weierstrass theorem. Hence \(\textnormal{span}\{z^k| k\in\mathbb{Z}\}\otimes A\) is dense in \(C(\mathbb{T})\otimes A\),
      and this implies that Laurent loops are a dense subset of invertible loops.
      \item \textbf{Step 2.} By virtue of Step 1, it suffices to show that the range of \(\beta_A\) contains the equivalence classes of all Laurent loops. But each Laurent loop is a quotient of two polynomial loops.
      Thus, it suffices to show that the range of \(\beta_A\) contains the classes of all polynomial loops. To this end, we show that for each \(n,m\in\mathbb{N}\) there is a continuous map
      \[
        \mu_m^n:PL_m^n\rightarrow PL_1^{mn+n}
      \] 
      such that \(\mu_m^n(f)\sim_h\diag(f,1_{mn})\) within \(PL_k^{mn+n}\) for all \(f\in PL_k^n\), \(k\leq m\). Indeed, let
      \(f(z)=\sum_{j=0}^{m}a_j z^j\), with \(a_j\in M_n(A)\) for all \(j=0,\ldots, m\). For each \(z\), we define
      \[
        \tilde{\mu}_m^n(f)(z)=\begin{pmatrix}
          a_0&a_1&a_2&\cdots & a_{m-1}&a_m\\
          -z1_n&1_n&0&\cdots &0&0\\
          0&-z1_n&1&\cdots&0&0\\
          \vdots&\vdots&\vdots&&\vdots&\vdots\\
          0&0&0&\cdots&-z1_n&1_n
        \end{pmatrix}.
      \]
      an element of \(M_{m+1}(M_n(A))\). Clearly, \(\tilde{\mu}_m^n(f)(z)=T_0+T_1 z\) for some \(T_0,T_1\in M_{mn+n}(A)\),
      and the map \(f\mapsto\tilde{\mu}_m^n(f)\) is continuous. We claim the following:
      \begin{enumerate}
        \item \(\tilde{\mu}_m^n(f)(z)\) is invertible for all \(z\),
        \item \(\tilde{\mu}_m^n(f)(1)\sim_h 1_{mn+n}\),
        \item \(\tilde{\mu}_m^n(f)\sim_h \diag(f,1_{mn})\).
      \end{enumerate}
      Once the properties are established, we obtain the desired map \(\mu_m^n\) by setting \(\mu_m^n(f)={(\tilde{\mu}_m^n(f)(1))}^{-1}\tilde{\mu}_m^n(f)\).\par
      In order to prove the three properties, we consider matrices
      \[ 
        A_m= \begin{pmatrix} 1&0&0&\cdots&-a_m\\ 0&1&0&\cdots&0\\ \vdots&\vdots&\vdots&&\vdots\\ 0&0&0&\cdots&1 \end{pmatrix},
        A_{m-1}= 
        \begin{pmatrix} 1&0&\cdots&-(a_{m-1}+a_m z)&0\\ 0&1&\cdots&0&0\\ \vdots&\vdots&&\vdots&\vdots\\ 0&0&\cdots&0&1 \end{pmatrix}, 
        \ldots, 
      \] 
      \[ 
        A_1= 
        \begin{pmatrix} 1&-\bigl(a_1+a_2z+\cdots+a_m z^{m-1}\bigr)&0&\cdots&0\\ 0&1&0&\cdots&0\\ \vdots&\vdots&\vdots&&\vdots\\ 0&0&0&\cdots&1 \end{pmatrix}. 
      \] 
      and matrices \(B_k\) having \(1\)'s on the main diagonal, \(z\) in the entry in column \(k\) and row \(k+1\), and \(0\)'s elsewhere.
      Then we have
      \[
        A_1 A_2\cdots A_m \tilde{\mu}_m^n (f) (z)=
        \begin{pmatrix}
          f(z)&0&\cdots&0&0\\
          -z&1&\cdots&0&0\\
          \vdots&\vdots&&\vdots&\vdots\\
          0&0&\cdots&-z&1
        \end{pmatrix}
      \]
      and 
      \begin{equation}\label{eq:827}
        A_1 A_2\cdots A_m \tilde{\mu}_m^n (f)(z)B_m B_{m-1}\cdots B_1=\diag(f(z),1_{mn}).
      \end{equation}
      Since \(f(z)\) and all the matrices \(A_1,\ldots,A_m,B_1,\ldots,B_m\) are invertible for all \(z\), we obtain (i).
      Furthermore, each of the \(A_j\) and \(B_j\) matrices may be continuously deformed to the identity within the set of polynomial loops by multiplying
      the sole off-diagonal entry with a parameter \(t\in[0,1]\). Thus (\ref{eq:827}) implies (ii) and (iii).
      \item \textbf{Step 3.} By virtue of Step 2, it suffices to show that the range of \(\beta_A\) contains the equivalence 
      classes of all linear loops. This will follow if we show that there exists a continuous retraction
      \[
        \nu:PL_1^n\rightarrow PRL^n
      \]
      such that \(\nu(f)\sim_h f\) within \(PL_1^n\) for all \(f\in PL_1^n\). Indeed, let \(f(z)=a_0+a_1 z\). Then
      \(f(1)=a_0+a_1\) is an invertible element of \(M_n(\mathbb{C}1_A)\), and we can put \(g={f(1)}^{-1}f\). Then
      \[
        g(z)=1_n+b(z-1),
      \]
      with \(b={(a_0+a_1)}^{-1}a_1\). When \(z\neq 1\) we can write 
      \[
        g(z)=(1-z)\left(\frac{1}{1-z}1_n-b\right),
      \]
      and since \(g(z)\) is invertible for all \(z\in\mathbb{T}\) we see that \(1/(1-z)\notin \operatorname{spec}(b)\) if \(z\in\mathbb{T}\setminus \{1\}\).
      Since the function \(z\mapsto1/(1-z)\) maps \(\mathbb{T}\setminus \{1\}\) into the line \(\{\lambda\in\mathbb{C}| \Re(\lambda)=1/2 \}\), we see that
      \[
        \operatorname{spec}(b)\subseteq\mathbb{C}\setminus \{\lambda\in\mathbb{C}| \Re(\lambda)=1/2\}.
      \]
      For \(t\in[0,1]\) consider a function
      \[
        g_t(z)=
        \begin{cases}
          tz&\textnormal{if }\Re(z)<1/2,\\
          tz+(1-t)&\textnormal{if }\Re(z)>1/2
        \end{cases}
      \]
      Each function \(g_t\) is holomorphic on an open neighbourhood of \(\operatorname{spec}(b)\)
      and thus the holomorphic calculus (see Exercise~\ref{ex:843}) yields elements \(g_t(b)\in M_n(A)\),
      which depend continuously on the parameter \(t\). Since the image of \(g_t(z)\) does not intersect the line
      \(\{\lambda\in\mathbb{C}| \Re(\lambda)=1/2\}\), the elements 
      \[
        h_t(z)=1_n+g_t(b)(z-1)=(1-z)\left(\frac{1}{1-z}1_n-g_t(b)\right)
      \]
      are invertible. We have \(g_1(z)=z\) and thus \(g_1(b)=b\). On the other hand, \({g_0(z)}^{2}=g_0(z)\) and thus \(e=g_0(b)\) is an idempotent.
      Consequently, \(t\mapsto h_t\) is a homotopy within \(PL_1^n\) between \(g\) and the idempotent 
      loop \(1_n+e(z-1)\). Now we can deform the idempotent \(e\) to a projection, as follows (cf. Exercise~\ref{ex:3320}).
      \begin{lemma}\label{lemma:822}
        Let \(B\) be a unital \cstar. Recall that \(\mathcal{I}(B)\) denotes the set of idempotents in \(B\).
        Then we have the following.
        \begin{enumerate}
          \item For every idempotent \(e\in B\) the element 
          \[
            \rho(e)=e e^*{(1+(e-e^*)(e^*-e))}^{-1}
          \]
          is a projection.
          \item The map \(\rho:\mathcal{I}(B)\rightarrow\mathcal{P}(B)\), defined in (1), is a continuous retraction. 
          In particular, \(\rho(e)\sim_h e\) in \(\mathcal{I}(B)\) for every idempotent \(e\).
          \item If \(p,q\in\mathcal{P}(B)\) and \(p\sim_h q\) in \(\mathcal{I}(B)\), then \(p\sim_h q\) in \(\mathcal{P}(B)\).
        \end{enumerate}
      \end{lemma}
      \begin{proof}
        \begin{enumerate}
          \item Put \(w=1+(e-e^*)(e^*-e)\). Then \(w\) is positive and invertible, thus \(\rho(e)=e e^* w^{-1}\) is well-defined.
        A straightforward calculation yields \(ew=e e^*=we\) and \(e^* w = e^* e e^*=w e^*\).
          Thus \(e e^* w={(e e^*)}^2\) and \(e e^* w^{-1}=w^{-1}e e^*\).
          This implies that \(e e^* w^{-1}\) is self-adjoint and that
          \[
            {\rho(e)}^2=e e^* w^{-1} e e^* w^{-1}={e e^*}^2 w^{-2}=e e^* w^{-1}=\rho(e)
          \]
          Whence \(\rho(e)\) is a projection.
          \item Clearly, \(\rho\) is a continuous map and \(\rho(p)=p\) if \(p\) is a projection.\par
          To see that \(\rho(e)\sim_h e\) in \(\mathcal{I}(B)\), set \(u_t=1-t(e-\rho(e))\) for \(t\in[0,1]\). Since \(\rho(e)e=e\)
          and \(e\rho(e)=\rho(e)\), we have \({(e-\rho(e))}^2=0\). Therefore \(u_t\) is invertible with inverse \(u_t^{-1}=1+t(e-\rho(e))\). Thus \(u_t^{-1} e u_t\) is an idempotent for all \(t\in[0,1]\), and we have
          \[
            e=u_0^{-1} e u_0\sim_h u_1^{-1} e u_1=(1+(e-\rho(e)))e(1-(e-\rho(e)))=\rho(e).
          \]
          \item If \(t\mapsto e_t\) is a continuous path in \(\mathcal{I}(B)\) from \(e_0=p\) to \(e_1=q\), then \(t\mapsto \rho(e_t)\) is a continuous path in \(\mathcal{P}(B)\) from \(\rho(e_0)=p\) to \(\rho(e_1)=q\).
        \end{enumerate}
      \end{proof}
      Let \(\rho:\mathcal{I}_n (A)\rightarrow\mathcal{P}_n (A)\) be the map defined in Lemma~\ref{lemma:822}
      (with \(B=M_n(A)\)). Then \(\nu(f)=1_n+\rho(e)(z-1)\) yields the desired map \(\nu:PL_1^n\rightarrow PRL^n\).
    \end{enumerate}
    \item \textbf{Injectivity.} Let \(p,q\in\mathcal{P}_n(A)\) and assume that \(\beta_A({[p]}_0-{[q]}_0)={[f_p f_q^*]}_1={[1]}_1\) in \(K_1(S A)\).
    Then, after increasing \(n\) if necessary, we have \(f_p\sim_h f_q\) in \(GL^n\). It suffices to show that there exists \(m\in\mathbb{N}\) such that \(\diag(p,1_m)\sim_h \diag(q,1_m)\) in \(\mathcal{P}_{n+m}(A)\).\par
    As a first step, we observe that there exists a polygonal (i.e. piece-wise linear) homotopy \(t\mapsto h_t\) from \(f_p\) to \(f_q\) such that 
    all \(h_t\) are Laurent loops with a uniform bound on both positive and negative degrees. (This follows from the density of  Laurent loops in invertible loops via a routine compactness argument --exercise).
    Thus there are \(m,k\in\mathbb{N}\) such that \(z^m h_t\in PL_k^n\) for all \(k\).
    Since \(z^m f_p\sim_h f_{\diag(p,1_m)}\) in \(PL_m^{m+n}\) (exercise), we see that \(f_{\diag(p,1_m)}\)
    and \(f_{\diag(q,1_m)}\) are homotopic in \(PL_{m+k}^{m+n}\). Let \(t\mapsto e_t\) be such a homotopy.
    Then applying the maps \(\mu_{m+k}^{m+n}\) and \(\nu\), constructed in steps 2 and 3, respectively, of the proof of surjectivity, we get a homotopy \(t\mapsto \nu(\mu_{m+k}^{m+n}(e_t))=f_{p_t}\) from \(f_{\diag(p,1_m)}\) to \(f_{\diag(q,1_m)}\) in projection loops.
    Since the map \(f_{p_t}\mapsto p_t\) is continuous (exercise), \(\diag(p,1_m)\) and \(\diag(q,1_m)\) are homotopic via a path of projections. Consequently, \({[p]}_0={[q]}_0\) in \(K_0(A)\), as required.
  \end{enumerate}
\end{proof}
Combining Theorems~\ref{theorem:812} and~\ref{theorem:821} we get
\[
  K_j(S A)\cong K_{1-j}(A)
\]
for any \cstar\ \(A\) and \(j=0,1\). Thus, for any natural \(n\) we have
\[
  K_{n+2}(A)\cong K_n(A).
\]
Furthermore, naturality of the maps \(\theta_*\) and \(\beta_*\) easily implies that the functors \(K_{n+2}\) and \(K_n\) are isomorphic.

\section{The six-term exact sequence}
With the Bott periodicity theorem at hand, we are now ready to present the 6-term exact
sequence of \(K\)-theory — a tool of paramount importance in applications. Let
\[
  \begin{tikzcd}
    0\arrow[r]&J\arrow[r,"\varphi"]&A\arrow[r,"\psi"]&B\arrow[r]&0
  \end{tikzcd}
\]
be an exact sequence of \cstars. Applying the suspension functor, we obtain the exact sequence
\[
  \begin{tikzcd}
    0\arrow[r]&S J\arrow[r,"S\varphi"]&S A\arrow[r,"S\psi"]&S B\arrow[r]&0.
  \end{tikzcd}
\]
Denote by \(\partial: K_1(S B)\rightarrow K_0(S J)\) the corresponding index map. Let \(\theta_J:K_1(J)\rightarrow K_0(S J)\)
and \(\beta_B:K_0(B)\rightarrow K_1(S B)\) be the isomorphisms from Theorems~\ref{theorem:812} and~\ref{theorem:821}, respectively.
Then the \textit{exponential map}
\[
  \partial_0:K_0(B)\rightarrow K_1(J)
\]
is defined as the unique homomorphism making the diagram 
\begin{equation}\label{diag:838}
  \begin{tikzcd}
    K_0(B)\arrow[r,"\partial_0"]\arrow[d,"\beta_B",swap]&K_1(J)\arrow[d,"\theta_J"]\\
    K_1(S B)\arrow[r,"\partial"]&K_0(S J)
  \end{tikzcd}
\end{equation}
commutative.
\begin{theorem}\label{theorem:839}
  Let 
  \[
    \begin{tikzcd}
      0\arrow[r]&J\arrow[r,"\varphi"]&A\arrow[r,"\psi"]&B\arrow[r]&0
    \end{tikzcd}
  \]
  be an exact sequence of \cstars. Then the sequence 
  \begin{equation}\label{diag:839}
    \begin{tikzcd}
      K_0(J)\arrow[r,"K_0(\varphi)"]&K_0(A)\arrow[r,"K_0(\psi)"]&K_0(B)\arrow[d,"\partial_0"]\\
      K_1(B)\arrow[u,"\partial_1"]&K_1(A)\arrow[l,"K_1(\psi)"]&K_1(J)\arrow[l,"K_1(\varphi)"]
    \end{tikzcd}
  \end{equation}
  is exact everywhere.
\end{theorem}
\begin{proof}
  By virtue of Theorem~\ref{theorem:727}, it suffices to show exactness at \(K_0(B)\) and \(K_1(J)\).
  It turns out at this stage this requires nothing more but a diagram chase.\par
  \begin{itemize}
    \item \(K_0(B)\): Consider the commutative (due to naturality of the Bott map) diagram
    \begin{equation}\label{diag:8310}
    \begin{tikzcd}
      K_0(A)\arrow[r,"K_0(\psi)"]\arrow[d,"\beta_A"]&K_0(B)\arrow[r,"\partial_0"]\arrow[d,"\beta_B"]&K_1(J)\arrow[d,"\theta_J"]\\
      K_1(S A)\arrow[r,"K_1(S\psi)"]&K_1(S B)\arrow[r,"\partial"]&K_0(S J)
    \end{tikzcd}
  \end{equation}
  All the vertical arrows are isomorphisms, and the bottom row is exact by Theorem~\ref{theorem:727}.
  Thus the top row is exact.
  \item \(K_1(J)\): Consider the commutative (due to naturality of the \(\theta_B\) map) diagram
  \begin{equation}\label{diag:8311}
    \begin{tikzcd}
      K_0(B)\arrow[r,"\partial_0"]\arrow[d,"\beta_B"]&K_1(J)\arrow[r,"K_1(\varphi)"]\arrow[d,"\theta_J"]&K_1(A)\arrow[d,"\theta_A"]\\
      K_1(S B)\arrow[r,"\partial"]&K_0(S J)\arrow[r,"K_0(S\varphi)"]&K_0(S A)
    \end{tikzcd}
  \end{equation}
  All the vertical arrows are isomorphisms, and the bottom row is exact by Theorem~\ref{theorem:727}.
  Thus the top row is exact.
  \end{itemize}
\end{proof}
\begin{proposition}[An explicit form of the exponential map]
  Let 
  \[
    \begin{tikzcd}
      0\arrow[r]&J\arrow[r,"\varphi"]&A\arrow[r,"\psi"]&B\arrow[r]&0
    \end{tikzcd}
  \]
  be exact and let \(\partial_0:K_0(B)\rightarrow K_1(J)\) be the associated exponential map. Then
  \begin{enumerate}
    \item If \(p\in\mathcal{P}_n(\tilde{B})\) and \(x=x^*\in M_n(\tilde{A})\) such that \(\tilde{\psi}(x)=p\) then \(\exists! u\in\mathcal{U}_n(\tilde{J})\)
    such that \(\tilde{\varphi}(u)=e^{2\pi i x}\), and we have
    \begin{equation}\label{eq:8312}
      \partial_0({[p]}_0-{[s(p)]}_0)=-{[u]}_1.
    \end{equation}
    \item Suppose that \(A\) is unital. If \(p\in\mathcal{P}_n(B)\) and \(x=x^*\in M_n(A)\) such that \(\psi(x)=p\),
    then \(\exists! u\in\mathcal{U}_n(\tilde{J})\) such that \(\tilde{\varphi}(u)=e^{2\pi i x}\), and we have
    \begin{equation}\label{eq:8313}
      \partial_0({[p]}_0)=-{[u]}_1.
    \end{equation}
  \end{enumerate}
\end{proposition}
\begin{proof}
  \begin{enumerate}
    \item It follows from (2) and a diagram chase, so we only need to prove (2).
    \item For simplicity, assume \(J\subseteq A\) and \(\varphi=\textnormal{id}\). Suppose \(A\) unital and let \(p\in\mathcal{P}_n(B)\).
    Then \(\psi(e^{2\pi ix})=e^{2\pi i\psi(x)}=e^{2\pi ip}=1_n\), hence \(e^{2\pi ix}\in\mathcal{U}_n(\tilde{J})\). We must show that
    \begin{equation}\label{eq:8314}
      \theta_J({[e^{-2\pi ix}]}_1)=(\partial_1\circ \beta_B)({[p]}_0),
    \end{equation}
    where \(\partial_1:K_1(S B)\rightarrow K_0(S J)\) is the index map corresponding to 
    \[
      \begin{tikzcd}
        0\arrow[r]&S J\arrow[r]&S A\arrow[r]&S B\arrow[r]&0.
      \end{tikzcd}
    \]
    We identify \(S B=\{f\in C([0,1],B)| f(0)=f(1)=0\}\). Thus \(M_k(\widetilde{S B})=\{f\in C([0,1],M_k(B))| f(0)=f(1)\in M_k(\mathbb{C}1_B)\}\),
    and \(f_p\in\mathcal{U}(\widetilde{S B})\) is \(f_p(t)=e^{2\pi it}p+1_n-p\), \(t\in[0,1]\).
    Let \(v\in\mathcal{U}_{2n}(\widetilde{S A})\) be such that \(\widetilde{S\psi}(v)=\begin{pmatrix}
      f_p&0\\0&f^*_p
    \end{pmatrix}\). Then \(v:[0,1]\rightarrow\mathcal{U}_{2n}(A)\) is a continuous map such that
    \(v(0)=v(1)\in M_{2n}(\mathbb{C}1_A)\), and \(\psi(v(t))=\begin{pmatrix}
      f_p(t)&0\\0&f_p^*(t)
    \end{pmatrix}\).
    As \(f_p(0)=f_p(1)=1\), we have \(v(0)=v(1)=1_{2n}\). With \(x=x^*\in M_n(A)\) a lift of \(p\), put \(z(t)=e^{2\pi itx}\)
    for \(t\in[0,1]\). \(t\mapsto z(t)\in\mathcal{U}(A)\) is continuous and \(\psi(z(t))=f_p(t)\). Hence
    \begin{equation}
      \psi\left(v(t)\begin{pmatrix}
        {z(t)}^*&0\\0&z(t)
      \end{pmatrix}\right)=1_{2n},\qquad
      s\left(v(t)\begin{pmatrix}
        {z(t)}^*&0\\0&z(t)
      \end{pmatrix}\right)=1_{2n}.
    \end{equation}
    Thus \(w(t)=v(t)\begin{pmatrix}
      {z(t)}^*&0\\0&z(t)
    \end{pmatrix}\) is a unitary element in \(\mathcal{U}_{2n}(\tilde{J})\). We have \(w(0)=1_{2n}\) and 
    \(w(1)=\begin{pmatrix}
      e^{-2\pi ix}&0\\0&e^{2\pi ix}
    \end{pmatrix}\). Thus, by definition of \(\theta_J\), we have
    \begin{equation}
      \theta_J({[e^{-2\pi ix}]}_1)={\left[w\begin{pmatrix}
        1_n&0\\0&0
      \end{pmatrix}w^*\right]}_0
      -{\left[\begin{pmatrix}
        1_n&0\\0&0
      \end{pmatrix}\right]}_0.
    \end{equation}
    We also have
    \begin{equation}
      w(t)\begin{pmatrix}
        1_n&0\\0&0
      \end{pmatrix}
      {w(t)}^*=v(t)\begin{pmatrix}
        1_n&0\\0&0
      \end{pmatrix}
      {v(t)}^*,
    \end{equation}
    and the unitary \(v\) was chosen so that 
    \begin{equation}
      \widetilde{S\psi}(v)=\begin{pmatrix}
        f_p&0\\0&f_p^*
      \end{pmatrix}.
    \end{equation}
    So, by definition of the index map, we get
    \begin{equation}
      \partial_1({[f_p]}_1)=\theta_J({[e^{-2\pi ix}]}_1).
    \end{equation}
  \end{enumerate}
\end{proof}

\section{Examples and Exercises}
\begin{exercise}\label{ex:841}
  Let \(A\) be a \cstar. Show that every class in \(K_1(A)\) contains a unitary \(u\in\mathcal{U}_n(\tilde{A})\)
  normalized so that \(s(u)=1_n\), where \(s\) is the scalar map.
\end{exercise}
\begin{exercise}\label{ex:842}
  Show that if \(A\) is a non-unital \cstar\ then for any \(p\in\mathcal{P}_n(\tilde{A})\) we have
  \[
    \beta_A({[p]}_0-{[s(p)]}_0)={[f_p f_{s(p)}^*]}_1.
  \]
\end{exercise}
\begin{exercise}[Holomorphic function calculus]\label{ex:843}
  Let \(\gamma_1,\ldots,\gamma_n\) be a finite collection of continuous and piece-wise continuously differentiable paths \(\gamma_k:[a_k,b_k]\rightarrow\mathbb{C}\).
  We assume that each \(\gamma_k\) is closed, i.e.\ \(\gamma_k(a_k)=\gamma_k(b_k)\).
  A \textit{contour} is a finite collection \(\Gamma=\{\gamma_1,\ldots,\gamma_n\}\).
  If \(f\) is a piece-wise continuous, complex function defined on \(\textnormal{Im}(\Gamma)=\bigcup_k\gamma_k([a_k,b_k])\), then there is a well-defined integral
  \[
    \int_{\Gamma}f(z) \,dz=\int_{\gamma_1}f(z)\,dz+\ldots\int_{\gamma_n}f(z)\,dz.
  \]
  If \(z_0\notin\textnormal{Im}(\Gamma)\) then the \textit{index} of \(z_0\) with respect to \(\Gamma\) is the integer defined as
  \[
    \textnormal{Indr}_\Gamma(z_0)=\frac{1}{2\pi i}\int_\Gamma \frac{\,dz}{z-z_0}.
  \]
  Let \(K\) be a compact subset of an open set \(\Omega\subseteq\mathbb{C}\). Then we say that \(\Gamma\) \textit{surrounds} \(K\) in \(\Omega\) if
  \(\textnormal{Im}(\Gamma)\subseteq\Omega\setminus K\) and
  \[
    \textnormal{Indr}_\Gamma(z)=\begin{cases}
      1,& z\in K\\
      0, & z\in\mathbb{C}\setminus\Omega.
    \end{cases}
  \]\par
  Let \(A\) be a unital \cstar. If \(a\in A\) and \(\Gamma\) is a contour surrounding \(\operatorname{spec}(a)\) in an open set \(\Omega\),
  then for every holomorphic function \(f:\Omega\rightarrow\mathbb{C}\) there is a well-defined Riemann integral
  \[
    f(a)=\frac{1}{2\pi i}\int_\Gamma f(z){(z1_A-a)}^{-1}\,dz.
  \]
  This integral yields a unique element \(f(a)\) of \(A\) such that for every continuous functional \(\varphi:A\rightarrow\mathbb{C}\) we have
  \[
    \varphi(f(a))=\frac{1}{2\pi i}\int_\Gamma f(z)\varphi({(z1_A-a)}^{-1})\,dz.
  \]
  The mapping \(f\mapsto f(a)\) is called the \textit{holomorphic function calculus} for \(a\).
  It has the following properties (see\ \cite{TM79},\ \cite{MGJ90},\ \cite{PGK79}).
  \begin{enumerate}
    \item The map \(f\mapsto f(a)\) is a unital algebra homomorphism.
    \item If \(g\) is a holomorphic function on \(f(\Omega)\) then \((g\circ f)(a)=g(f(a))\).
    \item \(\operatorname{spec}(f(a))=f(\operatorname{spec}(a))\).
    \item If \(f_n\) is a sequence of holomorphic functions on \(\Omega\) converging almost uniformly to a function \(g\), then \(g\) is holomorphic on \(\Omega\) and
    \[
      \|f_n(a)-g(a)\|\rightarrow 0.
    \]
  \end{enumerate}
  Note that the holomorphic functional calculus applies to an arbitrary element \(a\) of a \cstar, not just a self-adjoint one. If \(a\) is self-adjoint then the holomorphic function calculus is compatible
  with the continuous functional calculus via the Gelfand transform.
\end{exercise}
\begin{exercise}\label{ex:844}
  By virtue of Theorems~\ref{theorem:812} and~\ref{theorem:821}, we have
  \begin{align*}
    &K_0(C_0(\mathbb{R}^{2n}))\cong K_1(C_0(\mathbb{R}^{2n+1}))\cong K_0(\mathbb{C})\cong\mathbb{Z},\\
    &K_1(C_0(\mathbb{R}^{2n}))\cong K_0(C_0(\mathbb{R}^{2n+1}))\cong K_1(\mathbb{C})\cong 0,
  \end{align*}
  for all \(n\in\mathbb{N}\).
\end{exercise}
\begin{exercise}\label{ex:845}
  For each natural number \(n\geq 1\), find a split-exact sequence
  \[
    \begin{tikzcd}
      0\arrow[r]&C_0(\mathbb{R}^n)\arrow[r]&C(S^n)\arrow[r]&\mathbb{C}\arrow[r]&0.
    \end{tikzcd}
  \]
  Then use Exercise~\ref{ex:844} and split-exactness of \(K_*\) to determine the \(K\)-groups of \(C(S^n)\) for all spheres \(S^n\).
\end{exercise}
\begin{exercise}\label{ex:846}
  By Exercise~\ref{ex:845}, we have \(K_0(C(\mathbb{T}))\cong K_1(C(\mathbb{T}))\cong\mathbb{Z}\).
  Use the isomorphism \(C(\mathbb{T}^{n+1})\cong C(\mathbb{T})\otimes C(\mathbb{T}^n)\) to find a split-exact sequence
  \[
    \begin{tikzcd}
      0\arrow[r]&S C(\mathbb{T}^n)\arrow[r]&C(\mathbb{T}^{n+1})\arrow[r]&C(\mathbb{T}^n)\arrow[r]&0.
    \end{tikzcd}
  \]
  Then use split-exactness of \(K_*\) to determine the \(K\)-groups of \(C(\mathbb{T}^n)\) for all tori \(\mathbb{T}^n\).
\end{exercise}
\begin{exercise}\label{ex:847}
  Let 
  \[
    \begin{tikzcd}
      0\arrow[r]&J\arrow[r,"\varphi"]&A\arrow[r,"\psi"]&B\arrow[r]&0.
    \end{tikzcd}
  \]
  be an exact sequence of \cstars. Show that if every projection in \(\mathcal{P}_\infty(\tilde{B})\) lifts to a projection
  in \(\mathcal{P}_\infty(\tilde{A})\) then \(\partial_0:K_0(B)\rightarrow K_1(J)\) is the zero map.
\end{exercise}
\begin{exercise}[Toeplitz algebra]\label{ex:848}
  Let \(\mathcal{H}\) be a separable Hilbert space with an orthonormal basis \(\{\xi_n| n=0,1,2,\ldots\}\).
  Let \(S\in\mathcal{B}(\mathcal{H})\), \(S(\xi_n)=\xi_{n+1}\) be the unilateral shift.
  We define the \textit{Toeplitz algebra} \(\mathcal{T}\) as the \cstar\ generated by \(S\). 
  It can be shown\ \cite[722-726]{CL67} that \(\mathcal{T}\) is the universal \cstar\ for the relation \(S^* S=1\), 
  and that if \(T\) is a \textit{proper isometry} on a Hilbert space then there exists a *-isomorphism \(\mathcal{T}=C^*(S)\rightarrow C^*(T)\)
  such that \(T\mapsto S\).
  \begin{enumerate}
    \item Show that the closed two-sided ideal of \(\mathcal{T}\) generated by \(1-S S^*\) coincides with the algebra \(\mathcal{K}(\mathcal{H})\).
    \item Let \(\pi:\mathcal{T}\rightarrow \mathcal{T}/\mathcal{K}\) be the natural surjection. Show that \(\mathcal{T}/\mathcal{K}\) is isomorphic to \(C(S^1)\) and \(\pi(S)\)
    may be identified with the generator \(z\). There is an exact sequence
    \[
      \begin{tikzcd}
        0\arrow[r]&\mathcal{K}\arrow[r]&\mathcal{T}\arrow[r,"\pi"]&C(S^1)\arrow[r]&0.
      \end{tikzcd}
    \]
    \item By Exercise~\ref{ex:845}, \(K_0(C(S^1))\cong\mathbb{Z}\) (with generator \({[1]}_0\)) and \(K_1(C(S^1))\cong\mathbb{Z}\)
    (with a generator \({[z]}_1\), the class of the identity map \(z\mapsto z\)). Calculate \(\partial_1({[z]}_1)\) and show that the index map
    \[
      \partial_1:\mathbb{Z}\cong K_1(C(S^1))\rightarrow K_0(\mathcal{K})\cong\mathbb{Z}
    \]
    is an isomorphism.
    \item Use (3) and the exact sequence from Theorem~\ref{theorem:839} to show that
    \[
      K_0(\mathcal{T})\cong\mathbb{Z},\qquad K_1(\mathcal{T})=0.
    \]  
    Find the generator of \(K_0(\mathcal{T})\).
  \end{enumerate}
\end{exercise}
\begin{exercise}\label{ex:849}
  Let \(\mathcal{H}\) be a separable Hilbert space, \(\mathcal{K}=\mathcal{K}(\mathcal{H})\) be the compact operators on \(\mathcal{H}\),
  and let \(S^n\in\mathcal{B}(\mathcal{H})\) be an isometry with cokernel of dimension \(n\), for some natural number \(n\). Let \(C^*(S^n,\mathcal{K})\) be the 
  \(C^*\)-subalgebra of \(\mathcal{B}(\mathcal{H})\) generated by \(S^n\) and \(\mathcal{K}\).
  Show that there exists an exact sequence
  \[
    \begin{tikzcd}
      0\arrow[r]&\mathcal{K}\arrow[r]&C^*(S^n,\mathcal{K})\arrow[r]&C(S^1)\arrow[r]&0
    \end{tikzcd}
  \]
  and determine the \(K\)-theory of \(C^*(S^n,\mathcal{K})\).
\end{exercise}
\begin{exercise}\label{ex:8410}
  Let \(\mathbb{R}P^2\) be the real projective plane. Find an exact sequence
  \[
    \begin{tikzcd}
      0\arrow[r]&C_0(\mathbb{R}^2)\arrow[r]&C(\mathbb{R}P^2)\arrow[r]&C(S^1)\arrow[r]&0
    \end{tikzcd}
  \]
  and determine the \(K\)-theory of \(C(\mathbb{R}P^2)\).
\end{exercise}
\begin{exercise}\label{ex:8411}
  Find an exact sequence
  \[
    \begin{tikzcd}
      0\arrow[r]&C_0(S^1)\otimes C_0(\mathbb{R}^2)\arrow[r]&C(S^3)\arrow[r]&C(S^1)\arrow[r]&0.
    \end{tikzcd}
  \]
  Then apply to it the 6-term exact sequence of \(K\)-theory and thus calculate in an alternative way the \(K\)-groups of the 3-sphere (cf. Exercise~\ref{ex:845}).
\end{exercise}
\begin{exercise}\label{ex:8412}
  Let \(\mathcal{H}\) be a separable Hilbert space. Consider two operators \(T,U\in\mathcal{B}(\mathcal{H})\)
  such that \(T\) is a proper isometry (i.e. \(T^* T=1\neq T T^*\)) and \(U\) is a partial unitary on \(1-T T^*\)
  with full spectrum (i.e. \(U^* U=U U^*=1-T T^*\) and \(\operatorname{spec}(U)=S^1\cup \{0\}\)).
  Let \(A\) be the \(C^*\)-subalgebra of \(\mathcal{B}(\mathcal{H})\) generated by \(T\) and \(U\).
  \begin{enumerate}
    \item Let \(J\) be the closed two-sided ideal of \(A\) generated by \(U\). Show that \(J\) is isomorphic to \(C(S^1)\otimes\mathcal{K}\),
    with \(\mathcal{K}\) the \cstar\ of compact operators.
    \item Let \(\pi:A\rightarrow A/J\) be the natural surjection. Show that \(A/J\) is generated (as a \cstar) by the unitary element \(\pi(T)\).
    Show that \(\operatorname{spec}(\pi(T))\) contains the entire unit circle, and thus \(A/J\) is isomorphic to \(C(S^1)\).
    \item By (1) and (2) above, there is an exact sequence 
    \[
      \begin{tikzcd}
        0\arrow[r]&C_0(S^1)\otimes \mathcal{K}\arrow[r]&A\arrow[r]&C(S^1)\arrow[r]&0.
      \end{tikzcd}
    \]
    Apply the 6-term exact sequence and calculate the \(K\)-theory of \(A\).
  \end{enumerate}
\end{exercise}

\section{Mirror-disc-type quantum two-spheres}\label{example:8413}
Consider, for \(p\in(0,1)\), the *-algebra
\begin{equation}
  \mathcal{O}(D_p):=\mathbb{C}\langle x,x^*\rangle/J,
\end{equation}
where \(J\) is the *-ideal generated by \(x^* x-px x^*-(1-p)\).
This is called the *-algebra of the quantum disc (see~\cite[1-23]{KL93}, where a two-parameter family of such quantum discs is considered).
It is not hard to see that \(\|\rho(x)\|=1\) in any bounded representation \(\rho\),
so that one can form the \(C^*\)-closure \(C(D_p)\) of \(\mathcal{O}(D_p)\). Moreover \(\mathcal{O}(D_p)\) is faithfully
embedded in \(C(D_p)\). (There is exactly one faithful irreducible representation, up to unitary equivalence.) 
It is known that \(C(D_p)\) is isomorphic to the Toeplitz algebra \(\mathcal{T}\), so all the \cstars\ \(C(D_p)\) are isomorphic.
There is a *-homomorphism \(\varphi:C(D_p)\rightarrow C(S^1)\), sending the generator \(x\) to the unitary generator \(u\) of \(C(S^1)\).
Consider for any \(q\in(0,1)\) a second copy \(\mathcal{O}(D_q)\), with generator \(y\).
\begin{definition}
  Let \(\alpha:\mathcal{O}(S^1)\rightarrow\mathcal{O}(S^1)\) denote the *-automorphism defined by \(u\mapsto u^*\). Define
  \begin{equation}
    \mathcal{O}(S^2_{pq}):=\{(f,g)\in\mathcal{O}(D_p)\oplus\mathcal{O}(D_q)| \varphi(f)=\alpha\circ\varphi(g)\}.
  \end{equation}
  This is called the *-algebra of the mirror-disc-type quantum two-sphere.
\end{definition}
\begin{proposition}
  \begin{equation}
    \mathcal{O}(S^2_{pq})\cong\mathbb{C}\langle C,C^*,D,D^*,E,E^*\rangle/J,
  \end{equation}
  where the *-ideal \(J\) is generated by the relations
  \begin{align*}
    C^* C&=1-pD-E,\\
    C C^*&=1-D-qE,\\
    DC&=pCD,\\
    EC&=q^{-1}CE,\\
    DE&=0,\\
    D&=D^*,\\
    E&=E^*.
  \end{align*}
  The isomorphism is given by \((x,y^*)\mapsto C\), \((1-x x^*,0)\mapsto D\),
  \((0,1-y y^*)\mapsto E\).
\end{proposition}
\begin{proposition}
  The following is a complete list (up to unitary equivalence) of irreducible *-representations of \(\mathcal{O}(S^2_{pq})\) in some Hilbert space:
  \begin{enumerate}
    \item \(\rho_+\), acting on a separable Hilbert space \(\mathcal{H}\) with orthonormal basis \(e_0,e_1,\ldots\) according to
    \begin{align*}
      \rho_+(C)e_k&=\sqrt{1-p^{k+1}}e_{k+1},\\
      \rho_+(D)e_k&=p^k e_k,\\
      \rho_+(E)&=0.
    \end{align*}
    \item \(\rho_-\), acting on \(\mathcal{H}\) by
    \begin{align*}
      \rho_-(C)e_k&=\sqrt{1-q^{k}}e_{k-1},\\
      \rho_-(D)e_k&=0,\\
      \rho_-(E)e_k&=q^k e_k.  
    \end{align*}
    \item An \(S^1\)-family \(\rho_\mu\), acting on \(\mathbb{C}\) by
    \begin{align*}
      \rho_\mu(C)&=\mu,\\
      \rho_\mu(D)&=0,\\
      \rho_\mu(E)&=0.
    \end{align*}
  \end{enumerate}
\end{proposition}
One can again show that there is a uniform bound on the norm of the generators for all bounded *-representations, so that one can form a \(C^*\)-closure
\(C(S^2_{pq})\) of \(\mathcal{O}(S^2_{pq})\) using bounded *-representations.
\(\rho_+\oplus\rho_-\) is faithfully embedded in \(C(S^2_{pq})\). Moreover, the closed ideals \(J_D, J_E\)
generated by \(D, E\) are isometric to \(\mathcal{K}(\rho_+(J_D))=\mathcal{K}=\rho_-(J_E)\), they have zero intersection, and \((\rho_+\oplus\rho_-)(J_D+J_E)=\mathcal{K}\oplus\mathcal{K}\).
Finally, there is an exact sequence 
\begin{equation}\label{eq:8423}
  \begin{tikzcd}
    0\arrow[r]&\mathcal{K}\oplus \mathcal{K}\arrow[r]&C(S^2_{pq})\arrow[r,"\psi"]&C(S^1)\arrow[r]&0.
  \end{tikzcd}
\end{equation}
where \(\psi\) is defined by \(C\mapsto u\), \(D\mapsto 0\), \(E\mapsto 0\). This exact sequence can be used to compute the \(K\)-theory of \(C(S^2_{pq})\):
\begin{proposition}
  \(K_0(C(S^2_{pq}))\cong\mathbb{Z}\oplus\mathbb{Z}\), \(K_1(C(S^2_{pq}))=0\).
\end{proposition}
\begin{proof}
  With
  \[
    K_0(C(S^1))\cong\mathbb Z\cong K_1(C(S^1)),
    \qquad
    K_0(\mathcal K)\cong\mathbb Z,
    \qquad
    K_1(\mathcal K)=0,
  \]
  we obtain the standard 6-term exact sequence corresponding to (\ref{eq:8423})
  \begin{equation}\label{eq:8424}
    \begin{tikzcd}
      0\arrow[r]&K_1(C(S^2_{pq}))\arrow[r]&\mathbb{Z}\arrow[r,"\partial"]&\mathbb{Z}\oplus\mathbb{Z}\arrow[r]&K_0(C(S^2_{pq}))\arrow[r,"K_0(\psi)"]&\mathbb{Z}\arrow[r]&0.
    \end{tikzcd}
  \end{equation}
  Let us compute the index map \(\partial\). It is determined by its value on the generator \({[u]}_1\in K_1(C(S^1))\),
  \begin{equation}
    \partial({[u]}_1)={[1-b^* b]}_0-{[1-b b^*]}_0,
  \end{equation}
  where \(b\in C(S^2_{pq})\) is any partial isometry with \(\psi(b)=u\). Identify \(C(S^2_{pq})\cong (\rho_+\oplus \rho_-)(C(S^2_{pq}))\).
  Then \(b=(s,s^*)\), where \(s\) is the one-sided shift, is obtained from
  \((\rho_+\oplus\rho_-)(C)\) by polar decomposition and satisfies
  \(b-(\rho_+\oplus\rho_-)(C)\in\mathcal{K}\oplus\mathcal{K}\). More precisely,
  \[
    b=
    \bigl(
      \rho_+(C)|\rho_+(C)|^{-1},
      \rho_-(C)|\rho_-(C)|^\dagger
    \bigr),
  \]
  where \(|\rho_-(C)|^\dagger\) is the generalized inverse of
  \(|\rho_-(C)|\), defined by
  \[
    |\rho_-(C)|^\dagger e_k:=\begin{cases}
      0,&k=0,\\[1mm]
      \dfrac{1}{\sqrt{1-q^k}}e_k,&k>0.
    \end{cases}
  \]
  Then \(b^* b=(s^* s,s s^*)=1-p_2\), \(b b^*=(s s^*,s^* s)=1-p_1\),
  where \(p_1=(p_{e_0},0)\), \(p_2=(0,p_{e_0})\) can be considered as the generators \((0,1)\), \((1,0)\)
  of \(\mathbb{Z}\oplus\mathbb{Z}\cong K_0(\mathcal{K}\oplus\mathcal{K})\). Then \(\partial({[u]}_1)={[p_2]}_0-{[p_1]}_0=(0,1)-(1,0)\),
  so that \(\partial\) is injective, and we can conclude that \(K_1(C(S^2_{pq}))=0\).
  We are left with the exact sequence
  \begin{equation}\label{eq:8426}
    \begin{tikzcd}
      0\arrow[r]&\mathbb{Z}\arrow[r,"\partial"]&\mathbb{Z}\oplus\mathbb{Z}\arrow[r,"K_0(j)"]&K_0(C(S^2_{pq}))\arrow[r,"K_0(\psi)"]&\mathbb{Z}\arrow[r]&0.
    \end{tikzcd}
  \end{equation}
  As \(\mathbb{Z}\) is a free module over itself, this sequence splits, and \(K_0(C(S^2_{pq}))\cong \textnormal{Im}K_0(j)\oplus \mathbb{Z}\).
  There remains the exact sequence
  \begin{equation}\label{eq:8427}
    \begin{tikzcd}
      0\arrow[r]&\mathbb{Z}\arrow[r,"\partial"]&\mathbb{Z}\oplus\mathbb{Z}\arrow[r,"K_0(j)"]&\textnormal{Im}K_0(j)\arrow[r]&0.
    \end{tikzcd}
  \end{equation}
  Here, \(K_0(j)\) is determined by its values on \((1,0)\) and \((0,1)\). However, \((1,0)-(0,1)\in \textnormal{Ker} K_0(j)=\textnormal{Im}\partial\), i.e., \(K_0(j)(1,0)=K_0(j)(0,1)\),
  consequently \(\textnormal{Im}K_0(j)=\{nK_0(j)(1,0)| n\in\mathbb{Z}\}\cong\mathbb{Z}\).
  It follows that \(K_0(C(S^2_{pq}))\cong\mathbb{Z}\oplus\mathbb{Z}\).
\end{proof}

\chapter{K-Theory of Crossed Products}
The aim of this chapter is to derive the Pimsner--Voiculescu exact sequence 
from the Connes--Thom isomorphism, following the strategy of \cite[Chapter~V, Sections~10.3, 10.4 and 10.9]{Blackadar1998}. 
We shall, however, rearrange the order of the arguments so that every result is proved before it is used. 
We begin with the mapping torus associated with an automorphism and compute the maps occurring in its six-term exact 
sequence. We then give a complete proof of the Connes--Thom isomorphism by means of the Wiener--Hopf extension. 
This proof rests on Takai duality, Bott periodicity and the six-term exact sequence, but is independent of the 
Pimsner--Voiculescu theorem. Finally, for an automorphism \(\alpha\) of a \cstar\ \(A\), 
we apply these results to the dual action on \(A\rtimes_\alpha\mathbb{Z}\). 
Takai duality identifies the relevant double crossed product with \(A\otimes\mathcal{K}\), 
while the mapping-torus construction identifies the maps in \(K\)-theory with \(1-\alpha_*\). 
Combining these ingredients yields the Pimsner--Voiculescu exact sequence without any circular dependence among the results.

\section{The Mapping Torus}
The P--V exact sequence can be obtained almost immediately from the Thom
isomorphism by invoking a theorem of Green \cite{Green1977} on Morita equivalence. We
will give an alternate argument, based on the elementary structure of crossed
products. The concept we need is that of the mapping torus:
\begin{definition}
  Let \(\alpha\in\operatorname{Aut}(A)\). The \textnormal{mapping torus} of \(\alpha\) is 
  \[
    M_\alpha=\{f\in C(\mathbb R,A)\mid f(x+1)=\alpha(f(x))\}\cong\{f:[0,1]\to A\mid f(1)=\alpha(f(0))\}.
  \] 
\end{definition}

\begin{lemma}\label{lemma:mapping-torus-exact-sequence}
  Let \(\alpha\in\operatorname{Aut}(A)\). The evaluation map
  \[
    \operatorname{ev}_0:M_\alpha\longrightarrow A,
    \qquad
    \operatorname{ev}_0(f)=f(0),
  \]
  gives rise to a short exact sequence
  \[
    0\longrightarrow SA
    \longrightarrow M_\alpha
    \xrightarrow{\operatorname{ev}_0}A
    \longrightarrow 0.
  \]
\end{lemma}

\begin{proof}
  Using the realization
  \[
    M_\alpha
    =
    \{f\in C([0,1],A):f(1)=\alpha(f(0))\},
  \]
  we identify \(SA=C_0((0,1),A)\) with the ideal of those
  \(f\in M_\alpha\) that vanish at both endpoints. The map
  \(\operatorname{ev}_0\) is clearly a \(^*\)-homomorphism, and
  \[
    \ker(\operatorname{ev}_0)
    =
    \{f\in M_\alpha:f(0)=0\}.
  \]
  If \(f(0)=0\), then
  \[
    f(1)=\alpha(f(0))=0,
  \]
  and hence \(\ker(\operatorname{ev}_0)=SA\).

  It remains to prove that \(\operatorname{ev}_0\) is surjective. Given
  \(a\in A\), define
  \[
    f_a(t)=(1-t)a+t\alpha(a),
    \qquad 0\leq t\leq1.
  \]
  Then \(f_a\in M_\alpha\), since
  \(f_a(1)=\alpha(a)=\alpha(f_a(0))\), and
  \(\operatorname{ev}_0(f_a)=a\). Thus \(\operatorname{ev}_0\) is
  surjective, and the sequence is exact.
\end{proof}

\begin{remark}\label{remark:action-descends-to-torus}
  Let \(\beta:\mathbb{R}\to\operatorname{Aut}(B)\) be a strongly
  continuous action such that
  \[
    \beta_n=\operatorname{id}_B
    \qquad\text{for every }n\in\mathbb{Z}.
  \]
  Then \(\beta\) factors uniquely through the quotient
  \(\mathbb{T}=\mathbb{R}/\mathbb{Z}\). Indeed, one may define
  \[
    \overline{\beta}_{\,t+\mathbb{Z}}=\beta_t.
  \]
  This is well defined because, if \(t-s\in\mathbb{Z}\), then
  \[
    \beta_t
    =
    \beta_s\circ\beta_{t-s}
    =
    \beta_s.
  \]
  The resulting action
  \(\overline{\beta}:\mathbb{T}\to\operatorname{Aut}(B)\) is strongly
  continuous by the defining property of the quotient topology. We
  shall continue to denote it by \(\beta\).

  Since \(\mathbb{T}\) is locally compact and abelian, the crossed
  product \(B\rtimes_\beta\mathbb{T}\) carries the dual action
  \[
    \hat{\beta}:
    \hat{\mathbb{T}}
    \longrightarrow
    \operatorname{Aut}
    \bigl(B\rtimes_\beta\mathbb{T}\bigr).
  \]
  Under the canonical identification
  \(\hat{\mathbb{T}}\cong\mathbb{Z}\), this action is given on the
  dense subalgebra \(C(\mathbb{T},B)\) by
  \[
    \bigl(\hat{\beta}_n(f)\bigr)(z)
    =
    z^n f(z),
    \qquad
    n\in\mathbb{Z}.
  \]
  An action of \(\mathbb{Z}\) is determined by the image of its
  generator \(1\); hence we shall also write
  \[
    \hat{\beta}
    :=
    \hat{\beta}_1
    \in
    \operatorname{Aut}
    \bigl(B\rtimes_\beta\mathbb{T}\bigr),
  \]
  so that \(\hat{\beta}_n=\hat{\beta}^{\,n}\) for every
  \(n\in\mathbb{Z}\); cf.~\cite[Section~10.3]{Blackadar1998}.
\end{remark}
\begin{proposition}
  Let \(\beta\) be an action of \(\mathbb{R}\) on a \cstar\ \(B\), which is trivial on \(\mathbb{Z}\).
  Then \(B\rtimes_\beta\mathbb{R}\) is isomorphic to the mapping torus of \(\hat{\beta}\) on \(B\rtimes_\beta\mathbb{T}\).
\end{proposition}
\begin{proof}
  Represent \(B\) faithfully on a Hilbert space \(\mathcal{H}\). The regular representation
  of \(B\rtimes_\beta\mathbb{R}\) on \(\mathcal{H}\otimes L^2(\mathbb{R})\) is faithful, and so is obtained by inducing the regular representation \(\rho\)
  of \(B\rtimes_\beta\mathbb{Z}\) on \(\mathcal{H}\otimes L^2(\mathbb{Z})\) to \(B\rtimes_\beta\mathbb{R}\).
  But \(\beta\) is trivial on \(\mathbb{Z}\), so \(\rho\) is equivalent to the representation \(\sigma\) of \(B\otimes C^*(\mathbb{Z})\cong B\otimes C(\mathbb{T})\)
  on \(\mathcal{H}\otimes L^2(\mathbb{T})\) given by \(\sigma(b\otimes1)=b\otimes1\), \(\sigma(1\otimes u)=1\otimes M\), where \(M\) is 
  multiplication by \(z\) and \(u\) is the generator of \(\mathbb{Z}\). The induced representation \(\lambda\) acts on
  \[
    \{f:\mathbb{R}\to\mathcal{H}\otimes L^2(\mathbb{T})\mid f(t+1)=(1\otimes M)f(t), \int_0^1 \|f(t)\|^2\,dt<\infty\}
  \]
  with action \([\lambda(b)f](t)=(\beta_t(b)\otimes1)f(t)\), \([\lambda(r)f](t)=f(t-r)\).\par
  We may identify this Hilbert space with 
  \[
    L^2([0,1])\otimes(L^2([0,1])\otimes\mathcal{H})\cong L^2([0,1]^2)\otimes\mathcal{H},
  \]
  interchange the order in which the variables are written, and conjugate by the
  unitary \(U\) defined by \(Uf(s,t)=e^{2\pi ist}f(s,t)\), to convert \(\lambda\) into the representation
  \(\mu\) on \(L^2([0,1]^2)\otimes\mathcal{H}\) with action \([\mu(b)f](s,t)=\beta_t(b)f(s,t)\), \([\mu(r)f](s,t)=e^{2\pi irs}f(s,t-r[\operatorname{mod}\,1])\).\par
  For fixed \(s\), the operators \(\mu(b)\) and \(\mu(r)\) give operators on \(L^2([0,1])\otimes\mathcal{H}\);
  for fixed \(b\) or \(r\) these are continuous in \(s\); thus \(B\rtimes_\beta\mathbb{R}\) may be regarded
  as a subalgebra of \(C([0,1],\mathcal{B}(L^2([0,1])\otimes\mathcal{H}))\).
  For fixed \(s\), the operators \(\{\mu(b)\mu(\phi)\mid b\in B,\phi\in C_C(\mathbb{R})\}\) generate
  the canonical copy of \(B\rtimes_\beta\mathbb{Z}\) in \(L^2([0,1])\otimes\mathcal{H}\)
  (if \(\phi\) has small support, then there is a continuous function \(\psi\) on \(\mathbb{T}\) such that \(\mu(b)\mu(\phi)\) is
  the canonical action of \(b\psi\)), so \(B\rtimes_\beta\mathbb{R}\subseteq C([0,1],B\rtimes_\beta\mathbb{Z})\), and for each \(a\in B\rtimes_\beta\mathbb{Z}\)
  and each \(s\) there is an \(f\in B\rtimes_\beta\mathbb{R}\) with \(f(s)=a\). The values of any such 
  \(f\) at \(0\) and \(1\) match up via \(\hat{\beta}\), so \(B\rtimes_\beta\mathbb{R}\) may be identified with a subalgebra of the 
  mapping torus of \(\hat{\beta}\). Furthermore, the unitary \(\mu(1)\) is a multiplier of \(B\rtimes_\beta\mathbb{R}\);
  this unitary corresponds to the continuous function \(f(s)=e^{2\pi is}1\).
  Thus any continuous scalar-valued function \(f\) with \(f(0)=f(1)\) is a multiplier.
  A simple partition of unity shows that \(B\rtimes_\beta\mathbb{R}\) contains the whole mapping torus.
\end{proof}


\section{The Connes--Thom Isomorphism}
\begin{theorem}[Connes' Thom isomorphism]\label{ConnesThomIsomorphism}
  If \(\alpha:\mathbb{R}\to\operatorname{Aut}(A)\), then \(K_i(A\rtimes_\alpha\mathbb{R})\cong K_{1-i}(A)\), \(i=0,1\).
\end{theorem}
We prove the Thom isomorphism by means of a series of reductions.

\begin{definition}
  If \(\alpha\) is an action of \(\mathbb{R}\) on \ \cstar\ \(A\), we form the
  \textnormal{Wiener--Hopf} extension as follows. Let \(C=C_0(\mathbb{R}\cup\{+\infty\})\), and let \(\tau\)
  be the action of \(\mathbb{R}\) on \(C\) obtained by fixing \(+\infty\) and translating \(\mathbb{R}\).
  Let \(CA=C\otimes A\), \(\gamma\) the diagonal action \(\tau\otimes\alpha\).
  \(CA\) has an invariant ideal \(SA=C_0(\mathbb{R})\otimes A\);
  we also denote the action of \(\mathbb{R}\) on \(SA\) by \(\gamma\).
\end{definition}
\begin{lemma}
  \(SA\rtimes_\gamma\mathbb{R}\cong SA\rtimes_{\tau\otimes1}\mathbb{R}\cong A\otimes\mathcal{K}\).
\end{lemma}
\begin{proof}
  The first isomorphism is given by sending \(g\in C_C(\mathbb{R},C_C(\mathbb{R},A))\subseteq SA\rtimes_\gamma\mathbb{R}\)
  to \(\tilde{g}\), where \([\tilde{g}(s)](t)=\alpha_{-t}([g(s)](t))\).
  The second isomorphism comes from Takai duality.
\end{proof}

The evaluation map
\[
  \operatorname{ev}_{+\infty}:CA\longrightarrow A,
  \qquad
  \operatorname{ev}_{+\infty}(f)=f(+\infty),
\]
is equivariant with respect to \(\gamma\) and \(\alpha\), because
\(\tau\) fixes \(+\infty\), and its kernel is the invariant ideal
\(SA\). Hence Proposition~\ref{prop:exactness-full-crossed-products}
gives the short exact sequence
\[
  \begin{tikzcd}
    0\arrow[r]
    &SA\rtimes_\gamma\mathbb R\arrow[r]
    &CA\rtimes_\gamma\mathbb R\arrow[r]
    &A\rtimes_\alpha\mathbb R\arrow[r]
    &0.
  \end{tikzcd}
\]

Thus the six-term exact sequence for the extension
\[
  \begin{tikzcd}
    0\arrow[r]&SA\rtimes_\gamma\mathbb{R}\arrow[r]&CA\rtimes_\gamma\mathbb{R}\arrow[r]&A\rtimes_\alpha\mathbb{R}\arrow[r]&0   
  \end{tikzcd}
\]
becomes
\[
  \begin{tikzcd}
    K_1(A)\arrow[r]&K_1(CA\rtimes_\gamma\mathbb{R})\arrow[r]&K_1(A\rtimes_\alpha\mathbb{R})\arrow[d,"index"]\\
    K_0(A\rtimes_\alpha\mathbb{R})\arrow[u,"exp"]&K_0(CA\rtimes_\gamma\mathbb{R})\arrow[l]&K_0(A)\arrow[l]
  \end{tikzcd}
\]
The Thom isomorphism amounts to the statement that \(exp\) and \(index\) are isomorphisms.
\begin{lemma}
  If, for every \(A\) and \(\alpha\), the map \(index\) in the above exact sequence is surjective,
  then for every \(A\) and \(\alpha\) the maps \(exp\) and \(index\) are isomorphisms.
\end{lemma}
\begin{proof}
  Suppose that \(index\) is always surjective. Then \(K_1(A\rtimes_\alpha\mathbb{R})=0\) implies
  \(K_0(A)=0\). Applying this fact with \(A\) replaced by \(A\rtimes_\alpha\mathbb{R}\) and \(\alpha\)
  by \(\hat{\alpha}\), and applying Takai duality, we have that \(K_1(A)=0\) implies that \(K_0(A\rtimes_\alpha\mathbb{R})=0\)
  for any \(A\) and \(\alpha\). Since suspension (with trivial action) commutes with crossed products, we also have that
  \(K_0(A)=0\) implies \(K_1(A\rtimes_\alpha\mathbb{R})=0\) for any \(A\) and \(\alpha\). Apply these implications
  with \(A\) replaced by \(CA\) and \(\alpha\) replaced by \(\gamma\) (recall that \(K_*(CA)=0\) since \(CA\) is contractible) to conclude that \(K_*(CA\rtimes_\gamma\mathbb{R})=0\)
  for any \(A\) and \(\alpha\). The result follows from exactness.
\end{proof}
So we have reduced the problem down to showing that \(index\) is always surjective.
Fix a rank-one projection \(e\in\mathcal K\), and let
\[
  \phi:
  A\longrightarrow CA\rtimes_\gamma\mathbb R
\]
be the composition
\[
  A\xrightarrow{\,a\mapsto a\otimes e\,}
  A\otimes\mathcal K
  \xrightarrow{\Theta^{-1}}
  SA\rtimes_\gamma\mathbb R
  \xrightarrow{\iota}
  CA\rtimes_\gamma\mathbb R,
\]
where \(\Theta\) is the isomorphism of the preceding lemma and
\(\iota\) is the canonical inclusion.
The surjectivity of \(index\), by exactness, is the same as showing that the map \(\phi:A\to CA\rtimes_\gamma\mathbb{R}\) 
induces the zero map on K-theory.\par
We may describe \(\phi\) concretely as follows. Set \(E(s,t)=e^{s/2}e^{-t}\chi(t)\chi(t-s)\in L^1(\mathbb{R}^2)\),
where \(\chi\) is the characteristic function of \((0,\infty)\); then \(E\) defines
an element of \(C_0(\mathbb{R})\rtimes_\tau\mathbb{R}\) (even though \(E\) is not continuous in \(t\), \(E\) can be approximated 
in norm by elements of \(L^1(\mathbb{R},C_C(\mathbb{R}))\subseteq C_0(\mathbb{R})\rtimes_\tau\mathbb{R}\)), and in fact \(E\) is a ``standard''
one-dimensional projection in \(C_0(\mathbb{R})\rtimes_\tau\mathbb{R}\cong\mathcal{K}\). Then 
\(\phi\) may be taken to be given by the formula \([\phi(a)](s,t)=\alpha_t(a)E(s,t)\). \(\phi\)
maps \(A\) into the corner \(E(CA\rtimes_\gamma\mathbb{R})E\) of \(CA\rtimes_\gamma\mathbb{R}\).\par
There is an embedding \(\mu\) of \(A\) into the multiplier algebra \(M(CA\rtimes_\gamma\mathbb{R})\)
as ``constant functions'', i.e. if \(g\in C_C(\mathbb{R}\times(\mathbb{R}\cup\{+\infty\}),A)\subseteq CA\rtimes_\gamma\mathbb{R}\), then 
\([\mu(a)g](s,t)=ag(s,t)\), \([g\mu(a)](s,t)=g(s,t)\alpha_s(a)\).\par
Let \(B=CA\rtimes_\gamma\mathbb{R}+\mu(A)+\mathbb{C}1\subseteq M(CA\rtimes_\gamma\mathbb{R})\). Then \(B\)
is a split extension
\[
  \begin{tikzcd}
    0\arrow[r]&CA\rtimes_\gamma\mathbb{R}\arrow[r]&B\arrow[r, shift left=0.7ex]&\tilde{A}\arrow[l, shift left=0.7ex]\arrow[r]&0
  \end{tikzcd}
\]
and therefore the map \(i_*:K_0(CA\rtimes_\gamma\mathbb{R})\to K_0(B)\) is injective. Thus we need
only show that \(\psi_*:K_0(A)\to K_0(B)\) is the zero map, where \(\psi=i\circ\phi\). What we will actually show is that
\(\mu_*=\mu_*+\psi_*\).\par
We define a function \(F\) by \(F(s,t)=e^{-s/2}\chi(s)\chi(t-s)\).
\begin{lemma}
  \(F\) defines an element of \(C_0(\mathbb{R}\cup\{+\infty\})\rtimes_\tau\mathbb{R}\).
\end{lemma}
\begin{proof}
  The difficulty is similar to the problem with \(E\), namely that \(F\) is not continuous
  in \(t\). But if \(G_\epsilon(s,t)=e^{-s/2}g_\epsilon(s)g_\epsilon(t-s)\), where \(g_\epsilon\) agrees
  with \(\chi\) except on \([0,\epsilon]\) where it is linear, then \(G_\epsilon\to F\) as \(\epsilon\to0\).
\end{proof}
We have \(F^*(s,t)=e^{s/2}\chi(-s)\chi(t)\). Now define
\[
  S=1-F\in(C_0(\mathbb{R}\cup\{+\infty\})\rtimes_\tau\mathbb{R})^\sim.
\]
\begin{lemma}
  \(S^*S=1\), \(SS^*=1-E\).
\end{lemma}
The proof is a simple computation.\par
We define \(E_\epsilon=\frac{1}{\epsilon}E(\frac{s}{\epsilon},\frac{t}{\epsilon})\), \(F_\epsilon=\frac{1}{\epsilon}F(\frac{s}{\epsilon},\frac{t}{\epsilon})\), \(S_\epsilon=1-F_\epsilon\). Then,
as above, \(S_\epsilon^*S_\epsilon=1\), \(S_\epsilon S_\epsilon^*=1-E_\epsilon\).\par
Now suppose \(A\) unital. We may then regard \(C_0(\mathbb{R})\rtimes_\tau\mathbb{R}\) in a natural way
as a subalgebra of \(CA\rtimes_\gamma\mathbb{R}\), and hence we may regard \(E_\epsilon, F_\epsilon, S_\epsilon\) as elements of \(B\).
Set \(\omega_\epsilon(a)=S_\epsilon\mu(a)S_\epsilon^*\), \([\phi_\epsilon(a)](s,t)=\alpha_t(a)E_\epsilon(s,t)\), \(\psi_\epsilon=i\circ\phi_\epsilon\).
\begin{lemma}
  \(\omega_\epsilon\) is a *-homomorphism from \(A\) into \((1-E_\epsilon)B(1-E_\epsilon)\),
  \(\psi_\epsilon\) a *-homomorphism from \(A\) into \(E_\epsilon BE_\epsilon\), and \(\omega_\epsilon\) and \(\psi_\epsilon\) are point-norm continuous
  in \(\epsilon\) (for \(\epsilon>0\)).
\end{lemma}
Proof is obvious.\par
We write
  \[
    E=E_1,\qquad F=F_1,\qquad S=S_1,
    \qquad
    \omega=\omega_1,\qquad
    \phi=\phi_1,\qquad
    \psi=\psi_1.
  \]
\begin{lemma}  
  \(\omega_*=\mu_*\) as maps from \(K_0(A)\) to \(K_0(B)\).
\end{lemma}
\begin{proof}
  If \(p\) is a projection in \(M_n(A)\cong A\otimes M_n(\mathbb{C})\), then \((S\otimes1)(\mu\otimes1)(p)\)
  is a partial isometry from \((\mu\otimes1)(p)\) to \((\omega\otimes1)(p)\), so \(\mu_*({[p]}_0)=\omega_*({[p]}_0)\).
\end{proof}
If we set \(\mu_\epsilon=\omega_\epsilon+\psi_\epsilon\), then \(\mu_\epsilon\) is a point-norm continuous path of *-homomorphisms
from \(A\) to \(B\) (note that \(\omega_\epsilon\) and \(\psi_\epsilon\) have orthogonal ranges), and we
have \({\mu_\epsilon}_*={\omega_\epsilon}_*+{\psi_\epsilon}_*=\mu_*+\psi_*\). The next lemma shows that \({\mu_\epsilon}_*=\mu_*\),
so that \(\psi_*=0\), completing the proof that \(index\) is surjective in the unital case.
\begin{lemma}
  As \(\epsilon\rightarrow0\), \(\mu_\epsilon\rightarrow\mu\) in the point-norm topology.
\end{lemma}
\begin{proof}
  We need to show that \(S_\epsilon\mu(a)S^*_\epsilon+\psi_\epsilon(a)\rightarrow\mu(a)\) for each \(a\in A\).
  It suffices to show that for every \(a\in A\),
  \begin{align*}
    \|S_\epsilon\mu(a)S_\epsilon^*-S_\epsilon S_\epsilon^*\mu(a)\|&\rightarrow0\\
    \|\psi_\epsilon(a)-E_\epsilon\mu(a)\|&\rightarrow0
  \end{align*}
  We need to prove a sublemma first:\par
  For any \(a\in A\), \(\lim_{\epsilon\to0}\frac{1}{\epsilon}\int_0^\infty\|a-\alpha_{-s}(a)\|e^{-s/2\epsilon}\,ds=0\).\\
  \textit{Proof}.
  \begin{align*}
    \frac{1}{\epsilon}\int_0^\infty\|a-\alpha_{-s}(a)\|e^{-s/2\epsilon}\,ds&\\
    &=\frac{1}{\epsilon}\int_0^{\sqrt{\epsilon}}\|a-\alpha_{-s}(a)\|e^{-s/2\epsilon}\,ds+\frac{1}{\epsilon}\int_{\sqrt{\epsilon}}^\infty\|a-\alpha_{-s}(a)\|e^{-s/2\epsilon}\,ds\\
    &\leq\sup_{0<s<\sqrt{\epsilon}}\|a-\alpha_{-s}(a)\|\frac{1}{\epsilon}(-2\epsilon)e^{-s/2\epsilon}\mid_0^{\sqrt{\epsilon}}+\frac{2}{\epsilon}\|a\|(-2\epsilon)e^{-s/2\epsilon}\mid_{\sqrt{\epsilon}}^\infty\\
    &=\sup_{0<s<\sqrt{\epsilon}}\|a-\alpha_{-s}(a)\|2(1-e^{-1/(2\sqrt{\epsilon})})+4\|a\|e^{-1/(2\sqrt{\epsilon})}
  \end{align*}
  As \(\epsilon\rightarrow0\), \(\sup_{0<s<\sqrt{\epsilon}}\|a-\alpha_{-s}(a)\|\rightarrow0\), so both terms approach \(0\). $\Box$\par
  Suppose \(a\in A\). Then 
  \begin{align*}
    \|S_\epsilon\mu(a)S_\epsilon^*- S_\epsilon S_\epsilon^*\mu(a)\|&\leq\|\mu(a)S_\epsilon^*-S_\epsilon^*\mu(a)\|\\
    &=\|\mu(a)(1-F_\epsilon^*)-(1-F_\epsilon^*)\mu(a)\|=\|\mu(a)F_\epsilon^*-F_\epsilon^*\mu(a)\|
  \end{align*}
  Using the formula for \(F_\epsilon^*\) and the fact that the norm in \(CA\rtimes_\gamma \mathbb{R}\) is dominated by the norm in \(L^1(\mathbb{R},CA)\),
  we have
  \begin{align*}
    \|\mu(a)F_\epsilon^*-F_\epsilon^*\mu(a)\|&\leq\int_{-\infty}^\infty[\sup_t\|(a-\alpha_s(a))\frac{1}{\epsilon}e^{s/2\epsilon}\chi(-s)\chi(t)\|]\,ds\\
    &=\frac{1}{\epsilon}\int_{-\infty}^\infty\|a-\alpha_s(a)\|\chi(-s)e^{s/2\epsilon}\,ds\\
    &=\frac{1}{\epsilon}\int_0^\infty\|a-\alpha_{-s}(a)\|e^{-s/2\epsilon}\,ds
  \end{align*}
  This integral goes to \(0\) by the previous sublemma.\par
  We now consider the other expression. First note that
  \[
    \lim_{\epsilon\to0}[\sup_{t\geq0}\{\|\alpha_t(a)-a\|e^{-t/\epsilon}\}]=0
  \]
  so for any \(\delta>0\) there is an \(\epsilon_0\) such that for all \(\epsilon<\epsilon_0\), \(\sup_{t\geq0}\{\|\alpha_t(a)-a\|e^{-t/\epsilon}\}<\delta\).\par
  The second expression can be written
  \[
    \|\alpha_t(a)E_\epsilon(s,t)-E_\epsilon(s,t)\alpha_s(a)\|
  \]
  In a similar way as above, this expression is dominated by
  \begin{align*}
    \int_{-\infty}^\infty[\sup_{t}\|(\alpha_t(a)-\alpha_s(a))E_\epsilon(s,t)\|]\,ds\\
    =\int_{-\infty}^\infty[\sup_t\|(\alpha_t(a)-\alpha_s(a))\frac{1}{\epsilon}\chi(t)e^{s/2\epsilon}e^{-t/\epsilon}\chi(t-s)\|]\,ds.
  \end{align*}
  Consider the integrals \(\int_{0}^\infty\) and \(\int_{-\infty}^0\) separately.\par
  For \(\int_0^\infty\), we need only consider \(t\geq s\geq0\), so the integrand becomes
  \begin{align*}
    \sup_{t\geq s}\{\|\alpha_t(a)-\alpha_s(a)\|\frac{1}{\epsilon}e^{s/2\epsilon}e^{-t/\epsilon}\}&=\sup_{r\geq0}\{\|\alpha_{r+s}(a)-\alpha_s(a)\|\frac{1}{\epsilon}e^{s/2\epsilon}e^{-(r+s)/\epsilon}\}\\
    &=\sup_{r\geq0}\{\|\alpha_r(a)-a\|e^{-r/\epsilon}\}\frac{1}{\epsilon}e^{-s/2\epsilon}\\
    &\leq\delta\frac{1}{\epsilon}e^{-s/2\epsilon}\text{ for }\epsilon<\epsilon_0.
  \end{align*}
  So \(\int_0^\infty\leq\int_0^\infty\delta\frac{1}{\epsilon}e^{-s/2\epsilon}\,ds=2\delta\) for \(\epsilon<\epsilon_0\).\par
  For \(\int_{-\infty}^0\), the \(\chi(t-s)\) is irrelevant since we must only consider \(t\geq0\) because of the \(\chi(t)\) term. We have
  \begin{align*}
    \|\alpha_t(a)-\alpha_s(a)\|e^{-t/\epsilon}&\leq\{\|\alpha_t(a)-a\|+\|a-\alpha_s(a)\|\}e^{-t/\epsilon}\\
    &\leq\|\alpha_t(a)-a\|e^{-t/\epsilon}+\|a-\alpha_s(a)\|\\
    &\leq\delta+\|a-\alpha_s(a)\|\text{ for }\epsilon<\epsilon_0.
  \end{align*}
  So \(\int_{-\infty}^0\leq\int_{-\infty}^0(\|a-\alpha_s(a)\|+\delta)\frac{1}{\epsilon}e^{s/2\epsilon}\,ds=\int_{-\infty}^0\|a-\alpha_s(a)\|\frac{1}{\epsilon}e^{s/2\epsilon}\,ds+2\delta\) for \(\epsilon<\epsilon_0\).
  This integral goes to \(0\) by the sublemma.
\end{proof}
We have now proved that \(index\) is surjective whenever \(A\) is unital. For general \(A\), let \(\tilde{\alpha}\) denote the canonical extension of \(\alpha\) to \(\tilde{A}\). We have a commutative diagram
\[
  \begin{tikzcd}
    &K_1(A\rtimes_\alpha\mathbb{R})\arrow[r]\arrow[d]&K_1(\tilde{A}\rtimes_{\tilde{\alpha}}\mathbb{R})\arrow[r]\arrow[d]&K_1(C^*(\mathbb{R}))\arrow[d]\\
    K_1(\mathcal{K})\arrow[r]&K_0(A)\arrow[r]&K_0(\tilde{A})\arrow[r]&K_0(\mathcal{K})
  \end{tikzcd}
\]
with exact rows, where the vertical maps are given by \(index\). We know that the last two \(index\) maps
are surjective, and in fact the last one must be an isomorphism since both groups in the last column are isomorphic to \(\mathbb{Z}\).
Furthermore, the map from \(K_0(A)\) to \(K_0(\tilde{A})\) is injective since \(K_1(\mathcal{K})=0\).
A simple diagram chase shows that the first \(index\) map must also be surjective.\par
This completes the proof of the Thom isomorphism \ref{ConnesThomIsomorphism}.

\section{The Pimsner--Voiculescu Exact Sequence}

We now derive the Pimsner--Voiculescu exact sequence from the mapping-torus
description and the Connes--Thom isomorphism, following
\cite[Section~10.4]{Blackadar1998}. We first identify the maps that occur in
the six-term exact sequence of a mapping torus.

\begin{proposition}\label{prop:mapping-torus-connecting-maps}
  Let \(\alpha\in\operatorname{Aut}(A)\), and let \(M_\alpha\) be its mapping
  torus. If \(K_i(SA)\) is identified with \(K_{1-i}(A)\) by Bott periodicity,
  then both connecting maps in the six-term exact sequence associated with
  \[
    0\longrightarrow SA\longrightarrow M_\alpha
    \xrightarrow{\operatorname{ev}_0}A\longrightarrow0
  \]
  are identified with \(1-\alpha_*\).
\end{proposition}
\begin{proof}
  The automorphism \(\alpha\) extends uniquely to an automorphism \(\tilde{\alpha}\) of \(\tilde{A}\), and hence to every
  matrix algebra over \(\tilde{A}\). Let \(p\in M_n(\tilde{A})\) be a projection whose
  scalar part is \(1_k\), and put
  \[
    h_p(t)=(1-t)p+t\tilde{\alpha}(p),\qquad 0\leq t\leq1.
  \]
  Then \(h_p\in M_n(\widetilde{M_\alpha})\) is a self-adjoint lift of \(p\) through
  \(\operatorname{ev}_0\), and
  \[
    u_p(t)=\exp\bigl(2\pi i(h_p(t)-1_k)\bigr)
  \]
  is a unitary in \(M_n(\widetilde{SA})\) which is equal to \(1_n\) at both endpoints.
  By the explicit description of the exponential map and the definition of the
  Bott map,
  \[
    \partial_0\bigl({[p]}_0-{[1_k]}_0\bigr)
    =
    \beta_A\bigl({[p]}_0-{[\alpha(p)]}_0\bigr)
    =
    \beta_A\bigl((1-\alpha_*)({[p]}_0-{[1_k]}_0)\bigr).
  \]
  Thus the first connecting map is \(1-\alpha_*\) under the Bott
  identification. The assertion for the other connecting map follows by
  applying the same argument after suspension. This is
  \cite[Proposition~10.4.1]{Blackadar1998}.
\end{proof}

\begin{proposition}\label{prop:stabilized-automorphism-K-theory}
  Let \(\alpha\in\operatorname{Aut}(A)\) and
  \(\gamma\in\operatorname{Aut}(\mathcal K)\). Under the standard
  identification \(K_i(A)\cong K_i(A\otimes\mathcal K)\), one has
  \[
    \alpha_*=(\alpha\otimes\gamma)_*,\qquad i=0,1.
  \]
\end{proposition}
\begin{proof}
  Every automorphism of \(\mathcal K\) is induced by a unitary in
  \(\mathcal B(\mathcal H)\), and the unitary group of
  \(\mathcal B(\mathcal H)\) is connected. Hence \(\gamma\) is homotopic to
  the identity, and \(\alpha\otimes\gamma\) is homotopic to
  \(\alpha\otimes1\). The conclusion follows from homotopy invariance and
  stability of \(K\)-theory; cf.~\cite[Proposition~10.4.2]{Blackadar1998}.
\end{proof}

\begin{theorem}[Pimsner--Voiculescu exact sequence]
  Let \(A\) be a \cstar\ and let \(\alpha\in\operatorname{Aut}(A)\). If
  \(\iota_A:A\to A\rtimes_\alpha\mathbb Z\) denotes the canonical inclusion,
  then the following cyclic six-term sequence is exact:
  \[
    \begin{tikzcd}[column sep=large]
      K_0(A)\arrow[r,"1-\alpha_*"]
      &K_0(A)\arrow[r,"{\iota_A}_*"]
      &K_0(A\rtimes_\alpha\mathbb Z)\arrow[d,"\partial"]\\
      K_1(A\rtimes_\alpha\mathbb Z)\arrow[u,"\partial"]
      &K_1(A)\arrow[l,"{\iota_A}_*"']
      &K_1(A)\arrow[l,"1-\alpha_*"']
    \end{tikzcd}
  \]
\end{theorem}
\begin{proof}
  Set
  \[
    B=A\rtimes_\alpha\mathbb Z
  \]
  and let \(\beta=\widehat\alpha\) be the dual action of \(\mathbb T\) on
  \(B\). Via the quotient map \(\mathbb R\to\mathbb R/\mathbb Z=\mathbb T\),
  we regard \(\beta\) as an action of \(\mathbb R\) which is trivial on
  \(\mathbb Z\). By the mapping-torus proposition proved above,
  \[
    B\rtimes_\beta\mathbb R
    \cong
    M_{\widehat\beta},
  \]
  where \(\widehat\beta\) is the automorphism of
  \(B\rtimes_\beta\mathbb T\) determined by the generator of
  \(\hat{\mathbb T}=\mathbb Z\).

  Takai duality, Theorem~\ref{theorem:Takai}, gives
  \[
    B\rtimes_\beta\mathbb T
    =
    (A\rtimes_\alpha\mathbb Z)
    \rtimes_{\widehat\alpha}\mathbb T
    \cong
    A\otimes\mathcal K(\ell^2(\mathbb Z)).
  \]
  Under this isomorphism, \(\widehat\beta\) corresponds to
  \(\alpha\otimes\operatorname{Ad}\widetilde\lambda_1\), where
  \(\widetilde\lambda_1\) is the unitary arising from the regular
  representation of \(\mathbb Z\). Consequently the mapping-torus extension
  becomes
  \[
    0\longrightarrow S(A\otimes\mathcal K)
    \longrightarrow B\rtimes_\beta\mathbb R
    \longrightarrow A\otimes\mathcal K
    \longrightarrow0.
  \]

  The Connes--Thom isomorphism, Theorem~\ref{ConnesThomIsomorphism}, gives
  \[
    K_i(B\rtimes_\beta\mathbb R)
    \cong
    K_{1-i}(B)
    =
    K_{1-i}(A\rtimes_\alpha\mathbb Z),
    \qquad i=0,1.
  \]
  After these identifications and the standard stability identifications
  \(K_i(A\otimes\mathcal K)\cong K_i(A)\), the six-term exact sequence of the
  displayed extension has the groups and the maps induced by \(\iota_A\) shown
  in the statement. By Proposition~\ref{prop:mapping-torus-connecting-maps},
  its remaining maps are
  \[
    1-(\alpha\otimes\operatorname{Ad}\widetilde\lambda_1)_*.
  \]
  Proposition~\ref{prop:stabilized-automorphism-K-theory} identifies these
  maps with \(1-\alpha_*\). Exactness therefore gives the displayed sequence.
  This is the argument of \cite[Section~10.4]{Blackadar1998}.
\end{proof}

\begin{proposition}
  Let \(\alpha\) be an automorphism of a \cstar\ \(A\). Then, for
  \(i=0,1\), there are natural isomorphisms
  \[
    K_i(A\rtimes_\alpha\mathbb Z)
    \cong
    K_{1-i}(M_\alpha).
  \]
  The isomorphisms are natural with respect to covariant homomorphisms.
\end{proposition}
\begin{proof}
  These are the isomorphisms obtained in the preceding construction from the
  mapping-torus description, Takai duality and the Connes--Thom isomorphism;
  see \cite[Proposition~10.4.3]{Blackadar1998}.
\end{proof}



\chapter{K-Theory of the Noncommutative Torus}

\section{\texorpdfstring{\(K_0(A_\theta)\)}{K0(A theta)}}

We now apply the Pimsner--Voiculescu exact sequence to the realization
\[
  A_\theta\cong C(\mathbb{T})\rtimes_{\alpha_\theta}\mathbb{Z},
\]
where \(\alpha_\theta\) is the automorphism induced by rotation through
the angle \(2\pi\theta\).

\begin{theorem}\label{thm:K-theory-rotation-algebra}
  Let \(\theta\) be irrational. Then
  \[
    K_0(A_\theta)\cong\mathbb{Z}^2,
    \qquad
    K_1(A_\theta)\cong\mathbb{Z}^2.
  \]
  Moreover, \(K_1(A_\theta)\) is generated by the classes of the
    canonical unitaries \(U\) and \(V\).
\end{theorem}
\begin{proof}
  First compute the \(K\)-theory of \(C(\mathbb{T})\).
Since \(\mathbb{T}\) is connected, every projection-valued function in
\(C(\mathbb{T},M_n(\mathbb{C}))\) has constant rank. Moreover, every
projection \(p\) in \(C(\mathbb{T},M_n(\mathbb{C}))\) is homotopic,
through projections, to a constant projection of the same rank.
Indeed, a rank-\(r\) projection is a continuous map from \(\mathbb T\) to
the complex Grassmannian \(\operatorname{Gr}(r,n)\). Since
\(\operatorname{Gr}(r,n)\) is path-connected and simply connected, every
such map is null-homotopic. Equivalently, every complex vector bundle over
\(\mathbb T\) is trivial; compare
\cite[Example~VIII.5.2]{Davidson1996}.
It follows that rank is the only homotopy invariant of projections over
\(C(\mathbb{T})\), and hence
\[
  K_0(C(\mathbb{T}))\cong\mathbb{Z}.
\]
As noted above,
\[
  K_1(C(\mathbb{T}))\cong\mathbb{Z},
\]
with generator given by the class of the identity function
\(z\mapsto z\) on \(\mathbb{T}\).\par

Now consider the automorphism \(\alpha_\theta\) of \(C(\mathbb{T})\)
induced by rotation through the angle \(2\pi\theta\). Since this rotation
is homotopic to the identity map on \(\mathbb{T}\), homotopy invariance
of \(K\)-theory gives
\[
  K_i(\alpha_\theta)=\operatorname{id}_{K_i(C(\mathbb{T}))},
  \qquad i=0,1.
\]
Therefore, the Pimsner--Voiculescu exact sequence reduces to the two
short exact sequences
\[
  \begin{tikzcd}
    0\arrow[r]
      &K_0(C(\mathbb{T}))\arrow[r]
      &K_0(A_\theta)\arrow[r]
      &K_1(C(\mathbb{T}))\arrow[r]
      &0
  \end{tikzcd}
\]
and
\[
  \begin{tikzcd}
    0\arrow[r]
      &K_1(C(\mathbb{T}))\arrow[r]
      &K_1(A_\theta)\arrow[r]
      &K_0(C(\mathbb{T}))\arrow[r]
      &0.
  \end{tikzcd}
\]
Since \(K_0(C(\mathbb{T}))\cong K_1(C(\mathbb{T}))\cong\mathbb{Z}\)
and every short exact sequence of abelian groups with quotient
\(\mathbb{Z}\) splits, it follows that
\[
  K_0(A_\theta)\cong\mathbb{Z}^2,
  \qquad
  K_1(A_\theta)\cong\mathbb{Z}^2.
\]

It remains to identify the generators of the odd group. Under the
crossed-product isomorphism of
Theorem~\ref{theorem:noncommutativeTorusCrossedProduct}, the inclusion
\(C(\mathbb T)\to A_\theta\) sends the coordinate function \(z\) to \(U\).
Hence the first map in the short exact sequence for \(K_1\) sends the
generator \([z]_1\) to \([U]_1\). The class of the implementing unitary of
the crossed product is sent by the boundary map to a generator of
\(K_0(C(\mathbb T))\); in our realization this implementing unitary is
\(V^*\). Consequently,
\[
  \partial([V]_1)=\pm[1]_0.
\]
Thus \([V]_1\) maps to a generator of the quotient, whereas \([U]_1\)
generates the kernel. It follows that \([U]_1\) and \([V]_1\) generate
\(K_1(A_\theta)\); see also
\cite[Example~VIII.5.2]{Davidson1996}.
\end{proof}

The computation of K-theory shows that the noncommutative torus has the same K-groups as the ordinary two-dimensional torus.

\section{The Powers--Rieffel Projection}
The following construction is due to Powers and Rieffel; see
\cite[Theorem~1.1]{Rieffel1981}. We identify \(C(\mathbb{T})\) with
the subalgebra \(C^*(U)\) of \(A_\theta\), and, for \(f\in C(\mathbb{T})\),
we write simply \(f\) for \(f(U)\). We also view the elements of
\(C(\mathbb{T})\) as continuous functions on \(\mathbb{R}\) that are
periodic of period \(1\), and put
\[
  \alpha_\theta(f)(t)=f(t-\theta).
\]
Since \(V^*UV=e^{-2\pi i\theta}U\), functional calculus gives
\[
  V^*fV=\alpha_\theta(f),
  \qquad
  V^*f=\alpha_\theta(f)V^*.
\]
We say that an element of \(A_\theta\) is \textit{supported on}
\(\{-1,0,1\}\) if it is of the form
\[
  hV+f+gV^*,
\]
where \(f,g,h\in C(\mathbb{T})\).

\begin{theorem}[Powers--Rieffel]\label{thm:PowersRieffelProjection}
  For every \(\beta\in(\mathbb{Z}+\mathbb{Z}\theta)\cap[0,1]\) there is a projection \(p\in A_\theta\) supported on \(\{-1,0,1\}\), such that \(\tau(p)=\beta\).
\end{theorem}
\begin{proof}
  Suppose that \(p\) is a projection supported on \(\{-1,0,1\}\), and expressed, as above, 
  in terms of \(h,f,g\). Then from the fact that \(p\) is self-adjoint it is easily seen that \(f\) is 
  real-valued and that \(h=\alpha_\theta^{-1}(\overline{g})\). Combining this with the fact that \(p\) is idempotent, we obtain the following relations:
  \begin{enumerate}
    \item \(g(t)g(t-\theta)=0\),
    \item \(g(t)[1-f(t)-f(t-\theta)]=0\),
    \item \(f(t)[1-f(t)]=|g(t)|^2+|g(t+\theta)|^2\),
  \end{enumerate}
for \(t\in\mathbb{R}\). Conversely, it is easily seen that if \(f\) and \(g\) are
elements of \(C(\mathbb{T})\) satisfying the above relations, and if we let
\(h=\alpha_\theta^{-1}(\overline{g})\),
then the corresponding element of \(A_\theta\) will be a projection. Closer examination then shows that there are 
myriad choices of \(f\) and \(g\) which satisfy these relations.\par
Since translation by \(\theta\) is the same on \(C(\mathbb{T})\) as translation by
the fractional part of \(\theta\), we may assume that \(0<\theta<1\). Furthermore,
\(A_\theta\cong A_{1-\theta}\), and this isomorphism preserves the canonical trace.
Thus it is enough first to consider \(0<\theta<1/2\); when
\(1/2<\theta<1\), the complement of a projection of trace \(1-\theta\)
has trace \(\theta\). Assume, therefore, that \(0<\theta<1/2\). We first
construct a projection \(p\) such that \(\tau(p)=\theta\). For this, let \(f\) be
almost the characteristic function of the interval \([0,\theta]\), but rounded at the ends
in a somewhat careful way. Specifically, equation (3) shows that if \(f(t)\) is neither
\(0\) nor \(1\), then either \(g(t)\) or \(g(t+\theta)\) is nonzero (while equation (1) says
that not both can be nonzero simultaneously). Then equation (2) says that if \(g(t)\neq0\), then
\(f(t)+f(t-\theta)=1\). Choose any \(\epsilon>0\) such that \(\epsilon<\theta\) and \(\theta+\epsilon<1/2\).
On \([0,\epsilon]\) let \(f\) be any continuous function with values in \([0,1]\) and
with \(f(0)=0\) and \(f(\epsilon)=1\). On \([\theta,\theta+\epsilon]\) define \(f\)
by \(f(t)=1-f(t-\theta)\), while on \([\epsilon,\theta]\) and \([\theta+\epsilon,1]\)
let \(f\) be \(1\) and \(0\), respectively. Then define \(g\) on \([\theta,\theta+\epsilon]\) by
\[
  g(t)=(f(t)(1-f(t)))^{1/2},
\]  
and let \(g\) have value zero elsewhere on \([0,1]\). 
Then \(f\) and \(g\) satisfy relations (1), (2) and (3) above and so define 
a projection, whose trace is \(\int_0^1 f(t)\,dt=\theta\).\par
We now treat the general case. The values \(0\) and \(1\) are realized by
the projections \(0\) and \(1\). Let \(m\) be a positive integer. The
unitaries \(U^m\) and \(V\) satisfy
\[
  VU^m=e^{2\pi i m\theta}U^mV,
\]
and hence generate a unital subalgebra of \(A_\theta\) canonically
isomorphic to \(A_{\{m\theta\}}\), by the universal property and the
simplicity of the irrational rotation algebra. The restriction of \(\tau\) to this
subalgebra is its canonical trace. Consequently, the construction above,
applied with parameter \(\{m\theta\}\), gives a projection in \(A_\theta\)
of trace \(\{m\theta\}\). Since \(\theta\) is irrational,
\[
  \{-m\theta\}=1-\{m\theta\}.
\]
If \(q\) is a projection of trace \(\{m\theta\}\), then \(1-q\) is a
projection of trace \(\{-m\theta\}\). Every element of
\((\mathbb{Z}+\mathbb{Z}\theta)\cap(0,1)\) is of one of these two forms,
which proves the assertion.
\end{proof}


\section{Traces on Crossed Products by \texorpdfstring{\(\mathbb{Z}\)}{Z}}
It is a highly nontrivial matter to determine the order structure on
\(K_0(A\rtimes_\alpha\mathbb{Z})\). Even if the action is trivial, so that
\(A\rtimes_\alpha\mathbb{Z}\cong C(S^1,A)\), there is no known way to calculate
the order structure. For example, if \(A=C(\mathbb{T}^3)\), then \(K_0(A)\cong\mathbb{Z}^4\)
with strict ordering from the first coordinate; but \(K_0(C(\mathbb{T}^4))\cong\mathbb{Z}^8\)
is \textit{perforated}.\par
One can, however, obtain quite a bit of partial information: it is (at least theoretically)
possible to calculate the range of any state on \(K_0(A\rtimes_\alpha\mathbb{Z})\)
which comes from a tracial state on the crossed product; indeed, every state on
\(K_0\) arises in this way \cite[Section~10.10]{Blackadar1998}. In good cases this calculation can be done
rather easily and gives complete information on the order structure.\par
Since order structures and states on \(K_0(A)\) only make good sense when \(A\)
is unital, we will only consider unital \cstars, although some of the results carry
through to the nonunital case.\par
The constructions below follow \cite[Sections~10.10.1--10.10.4]{Blackadar1998}.
\begin{definition}
  Let \(A\) be a unital \cstar, let
  \(\alpha\in\operatorname{Aut}(A)\), and let \(\tau\) be an
  \(\alpha\)-invariant tracial state on \(A\). Let
  \[
    E:A\rtimes_\alpha\mathbb{Z}\longrightarrow A
  \]
  be the canonical conditional expectation. The \textnormal{dual trace}
  associated with \(\tau\) is
  \[
    \widetilde{\tau}=\tau\circ E.
  \]
  Thus, on finite sums,
  \[
    \widetilde{\tau}\left(\sum_{n\in\mathbb{Z}}a_nu^n\right)
    =
    \tau(a_0).
  \] 
\end{definition}
\begin{proposition}\label{prop:dualTrace}
  The functional \(\widetilde{\tau}\) is a tracial state on
  \(A\rtimes_\alpha\mathbb{Z}\), and its restriction to \(A\) is
  \(\tau\).
\end{proposition}
\begin{definition}\label{def:traceOnK0}
  For each \(n\), let
  \[
    \tau_n=\tau\otimes\operatorname{Tr}_n
  \]
  be the unnormalized extension of \(\tau\) to \(M_n(A)\).
  The trace induces a group homomorphism
  \[
    \tau_*:K_0(A)\longrightarrow\mathbb{R}
  \]
  defined by
  \[
    \tau_*([p]_0-[q]_0)
    =
    \tau_n(p)-\tau_m(q),
  \]
  for \(p\in\mathcal P(M_n(A))\) and
  \(q\in\mathcal P(M_m(A))\).
  When no confusion can arise, we use the same symbol \(\tau\) for all
  of its matrix extensions \(\tau_n\).
\end{definition}
\begin{definition}\label{def:LambdaTrace}
  Let
  \[
    M_\alpha
    =
    \{f\in C([0,1],A):f(1)=\alpha(f(0))\}.
  \]
  If \(u\in\mathcal U_\infty(M_\alpha)\) is piecewise smooth, define
  \[
    \Lambda_\tau([u]_1)
    =
    \frac{1}{2\pi i}
    \int_0^1
    \tau\bigl(u'(t)u(t)^{-1}\bigr)\,dt.
  \]
\end{definition}
\begin{proposition}\label{prop:LambdaTrace}
  The map
  \[
    \Lambda_\tau:K_1(M_\alpha)\longrightarrow\mathbb{R}
  \]
  is a well-defined group homomorphism. Under the canonical
  identification
  \[
    K_0(A\rtimes_\alpha\mathbb{Z})
    \cong K_1(M_\alpha),
  \]
  it extends the homomorphism
  \(\tau_*:K_0(A)\to\mathbb{R}\).
\end{proposition}
\begin{definition}[de la Harpe--Skandalis determinant]
  \label{def:deLaHarpeSkandalis}
  Let
  \[
    \rho_\tau:
    \mathbb{R}\longrightarrow
    \mathbb{R}/\tau_*(K_0(A))
  \]
  be the quotient map. For \(u\in\mathcal U_\infty(A)_0\), choose a
  piecewise smooth path
  \[
    \xi:[0,1]\longrightarrow\mathcal U_\infty(A)
  \]
  such that \(\xi(0)=1\) and \(\xi(1)=u\). Define
  \[
    \Delta_\tau(u)
    =
    \rho_\tau\left(
      \frac{1}{2\pi i}
      \int_0^1
      \tau\bigl(\xi'(t)\xi(t)^{-1}\bigr)\,dt
    \right).
  \]
\end{definition}
\begin{proposition}\label{prop:deLaHarpeWellDefined}
  The map
  \[
    \Delta_\tau:
    \mathcal U_\infty(A)_0
    \longrightarrow
    \mathbb{R}/\tau_*(K_0(A))
  \]
  is well-defined and is a group homomorphism.
\end{proposition}
\begin{lemma}\label{lemma:determinantExponentials}
  Let \(x_1,\ldots,x_r\in M_\infty(A)\) be self-adjoint. Then
  \[
    \Delta_\tau\left(
      \prod_{k=1}^r e^{2\pi i x_k}
    \right)
    =
    \rho_\tau\left(
      \sum_{k=1}^r\tau(x_k)
    \right).
  \]
\end{lemma}
\begin{proposition}\label{prop:determinantCommutative}
  Suppose that \(A\) is commutative. Then, for every
  \(u\in\mathcal U_\infty(A)_0\),
  \[
    \Delta_\tau(u)
    =
    \Delta_\tau(\det u).
  \]
  In particular, \(\Delta_\tau\) factors through the ordinary
  determinant
  \[
    \det:\mathcal U_\infty(A)\longrightarrow\mathcal U(A).
  \]
\end{proposition}
Now suppose 
\[
  \begin{tikzcd}
    0\arrow[r]&J\arrow[r,"j"]&B\arrow[r,"q"]&A\arrow[r]&0
  \end{tikzcd}
\]
is an exact sequence of \cstars. If a \cstar\ \(D\) is nonunital, we use the convention
\[
  \mathcal U_n(D)
  =
  \bigl\{u\in\mathcal U(M_n(\widetilde D)):
  (\operatorname{id}_{M_n}\otimes\pi_D)(u)=1_n\bigr\},
  \qquad
  \mathcal U_\infty(D)=\bigcup_{n=1}^\infty\mathcal U_n(D),
\]
where \(\pi_D:\widetilde D\to\mathbb C\) is the canonical quotient map.
If \(\tau\) is a trace on \(A\), and \(u\in\mathcal{U}_\infty(J)\)
with \(j(u)\in\mathcal{U}_\infty(B)_0\) (i.e. \([u]_1\in\ker(j_*):K_1(J)\rightarrow K_1(B)\)),
then we may define \(\underline{\Delta}_\tau(u)=\Delta_{\tau\circ q}(j(u))\).
If \(u\in\mathcal{U}_\infty(J)_0\), one may choose the map \(\xi\) of Definition~\ref{def:deLaHarpeSkandalis} to
lie within \(\mathcal{U}_\infty(J)\) and so \(\underline{\Delta}_\tau(u)=0\);
thus \(\underline{\Delta}_\tau\) drops to a well-defined group homomorphism from
\(\ker(j_*)\) to \(\mathbb{R}/(\tau\circ q)_*(K_0(B))\).
\begin{proposition}
  The following sequence is exact:
  \[
    \begin{tikzcd}
      0\arrow[r]&(\tau\circ q)_*(K_0(B))\arrow[r]&\tau_*(K_0(A))\arrow[r,"\rho"]&\underline{\Delta}_\tau(\ker(j_*))\arrow[r]&0
    \end{tikzcd}
  \]
  (where the first map is the inclusion of two subgroups of \(\mathbb{R}\), and \(\rho\) is the
  restriction of the quotient map from \(\mathbb{R}\) to \(\mathbb{R}/(\tau\circ q)_*(K_0(B))\)).
\end{proposition}
\begin{proof}
  From the six-term exact sequence of K-theory, we have \(\ker(j_*)=\delta(K_0(A))\).
  If \(p\) is a projection in \(M_\infty(A)\), choose \(x=x^*\in M_\infty(B)\) with 
  \(q(x)=p\). Then \(u=e^{2\pi ix}\in\mathcal U_\infty(J)\), \([u]_1\in\ker(j_*)\), and
  \(\underline{\Delta}_\tau(u)=\rho(\tau(p))\), 
  so \(\rho(\tau_*(K_0(A)))\subseteq\underline{\Delta}_\tau(\ker(j_*))\).
  Differences of such elements fill up all of \(\underline{\Delta}_\tau(\ker(j_*))\).
\end{proof}
We come now to the main result:
\begin{theorem}\label{thm:traceRangeCrossedProduct}
  Let \(\alpha\in\operatorname{Aut}(A)\), and let \(\tau\) be an invariant tracial 
  state of \(A\). Then the map \(\Delta_\tau^\alpha\) from \(\ker(1-\alpha_*)\subseteq K_1(A)\)
  to \(\mathbb{R}/\tau_*(K_0(A))\) defined by
  \[
    \Delta_\tau^\alpha([u])=\Delta_\tau(u\alpha(u^{-1})) 
  \]
  (\(u\in\mathcal{U}_\infty(A)\), \(u\alpha(u^{-1})\in\mathcal{U}_\infty(A)_0\)) is a well-defined group homomorphism,
  and the following sequence is exact:
  \[
    \begin{tikzcd}
      0\arrow[r]&\tau_*(K_0(A))\arrow[r]&\widetilde{\tau}_*(K_0(A\rtimes_\alpha\mathbb{Z}))\arrow[r,"\rho"]&\Delta_\tau^\alpha(\ker(1-\alpha_*))\arrow[r]&0
    \end{tikzcd}
  \]
\end{theorem}
\begin{proof}
  Apply the previous proposition to the extension in the Pimsner--Voiculescu exact sequence.
\end{proof}

\section{The Ordered \texorpdfstring{\(K_0\)}{K0}-Group}
 
We now combine the computation of the range of the canonical trace with
the Powers--Rieffel projections. As before, replacing \(\theta\) by its
fractional part if necessary, we assume that \(0<\theta<1\).

\begin{proposition}\label{prop:traceRangeRotationAlgebra}
Let \(\theta\) be irrational and let \(\tau\) be the canonical trace on
\(A_\theta\). Then
\[
  \tau_*(K_0(A_\theta))
  =
  \mathbb Z+\mathbb Z\theta.
\]
Moreover,
\[
  \tau_*:K_0(A_\theta)\longrightarrow
  \mathbb Z+\mathbb Z\theta
\]
is an isomorphism of abelian groups.
\end{proposition}

\begin{proof}
Under the crossed-product realization
\[
  A_\theta\cong C(\mathbb T)\rtimes_{\alpha_\theta}\mathbb Z,
\]
the canonical trace on \(A_\theta\) is the dual of integration against
normalized Haar measure on \(\mathbb T\).

Let \(z\in C(\mathbb T)\) be the coordinate function. Since
\[
  \alpha_\theta(z)=e^{-2\pi i\theta}z,
\]
we have
\[
  z\alpha_\theta(z^{-1})=e^{2\pi i\theta}1.
\]
The path
\[
  \xi(t)=e^{2\pi i\theta t}1,
  \qquad 0\leq t\leq1,
\]
therefore gives
\[
  \Delta_\tau^\alpha([z]_1)
  =
  \theta+\mathbb Z
  \quad\text{in }\mathbb R/\mathbb Z.
\]
Since \([z]_1\) generates
\(K_1(C(\mathbb T))\cong\mathbb Z\), Theorem
\ref{thm:traceRangeCrossedProduct} gives
\[
  \tau_*(K_0(A_\theta))
  =
  \mathbb Z+\mathbb Z\theta;
\]
compare \cite[Exercise~10.11.6(d)]{Blackadar1998}.

By Theorem~\ref{thm:K-theory-rotation-algebra},
\(K_0(A_\theta)\cong\mathbb Z^2\). Since \(\theta\) is irrational,
\(\mathbb Z+\mathbb Z\theta\) is also a free abelian group of rank two.
The surjective homomorphism
\[
  \tau_*:K_0(A_\theta)\longrightarrow
  \mathbb Z+\mathbb Z\theta
\]
is consequently injective as well.
\end{proof}
\begin{theorem}\label{thm:orderedK0RotationAlgebra}
Let \(\theta\) be irrational. Then
\[
  K_0(A_\theta)^+
  =
  \{0\}
  \cup
  \{x\in K_0(A_\theta):\tau_*(x)>0\}.
\]
Consequently, the canonical trace induces an isomorphism of scaled
ordered abelian groups
\[
  \tau_*:
  \bigl(K_0(A_\theta),K_0(A_\theta)^+,[1]_{0}\bigr)
  \longrightarrow
  \bigl(
    \mathbb Z+\mathbb Z\theta,
    (\mathbb Z+\mathbb Z\theta)\cap[0,\infty),
    1
  \bigr).
\]
\end{theorem}

\begin{proof}
If \(x\in K_0(A_\theta)^+\), then \(x=[p]_0\) for some projection
\(p\in M_n(A_\theta)\). Since the canonical trace is faithful by
Theorem~\ref{thm:canonicalTraceRotationAlgebra},
\[
  \tau_*(x)\geq0,
\]
and equality holds only when \(p=0\). Thus every nonzero element of
\(K_0(A_\theta)^+\) has strictly positive trace.

Conversely, let \(x\in K_0(A_\theta)\) and suppose that
\[
  t=\tau_*(x)>0.
\]
By Proposition~\ref{prop:traceRangeRotationAlgebra},
\(t\in\mathbb Z+\mathbb Z\theta\). Write
\[
  t=k+\beta,
  \qquad
  k\in\mathbb N\cup\{0\},
  \qquad
  0\leq\beta<1.
\]
If \(\beta=0\), then
\[
  \tau_*(x)=\tau_*(k[1]_0),
\]
and the injectivity of \(\tau_*\) gives \(x=k[1]_0\).

If \(0<\beta<1\), Theorem~\ref{thm:PowersRieffelProjection} provides a
projection \(p\in A_\theta\) such that \(\tau(p)=\beta\). Hence
\[
  y=
  [\operatorname{diag}(1_k,p)]_0
  \in K_0(A_\theta)^+
\]
satisfies
\[
  \tau_*(y)=k+\beta=t=\tau_*(x).
\]
Again, injectivity of \(\tau_*\) implies \(x=y\), and therefore
\(x\in K_0(A_\theta)^+\).

The assertion concerning scaled ordered groups now follows from
\(\tau_*([1]_0)=1\).
\end{proof}

\begin{corollary}\label{cor:orderedK0RotationCoordinates}
Let \(p_\theta\) be a Powers--Rieffel projection satisfying
\(\tau(p_\theta)=\theta\). Under the isomorphism
\[
  \mathbb Z^2\longrightarrow K_0(A_\theta),
  \qquad
  (m,n)\longmapsto m[1]_0+n[p_\theta]_0,
\]
we have
\[
  K_0(A_\theta)^+
  =
  \{(0,0)\}
  \cup
  \{(m,n)\in\mathbb Z^2:m+n\theta>0\},
\]
and
\[
  [1]_0=(1,0).
\]
In particular, the order on \(K_0(A_\theta)\) is total.
\end{corollary}

\begin{proof}
The trace of \(m[1]_0+n[p_\theta]_0\) is \(m+n\theta\), so the
description of the positive cone follows from
Theorem~\ref{thm:orderedK0RotationAlgebra}. Since \(\theta\) is
irrational, \(m+n\theta\neq0\) whenever \((m,n)\neq(0,0)\); hence every
nonzero element is either positive or negative.
\end{proof}

\chapter{Classification and Morita Equivalence}
Morita equivalence provides a natural notion of equivalence for \cstars
which is weaker than isomorphism but preserves their essential
representation-theoretic and \(K\)-theoretic structure. In this chapter
we develop the necessary Hilbert-module framework and apply it to the
classification of irrational rotation algebras.

\section{Hilbert Modules and Imprimitivity Bimodules}
The basic material on Hilbert modules used in this section follows
\cite[Sections~13.1--13.2]{Blackadar1998}.
\begin{definition}
Let \(B\) be an algebra over \(\mathbb C\). A
\textnormal{right \(B\)-module} is a complex vector space \(X\)
equipped with a bilinear map
\[
  X\times B\longrightarrow X,
  \qquad
  (x,b)\longmapsto xb,
\]
such that
\[
  (xb)c=x(bc)
\]
for every \(x\in X\) and \(b,c\in B\).

If \(B\) is unital, we additionally require
\[
  x1_B=x
\]
for every \(x\in X\).
\end{definition}
\begin{definition}
Let \(B\) be a \cstar. A \textnormal{right pre-Hilbert
\(B\)-module} is a right \(B\)-module \(X\) equipped with a map
\[
  \langle\cdot,\cdot\rangle_B:X\times X\longrightarrow B
\]
such that, for every \(x,y,z\in X\), \(b\in B\) and
\(\lambda,\mu\in\mathbb C\),
\begin{enumerate}
  \item
  \[
    \langle x,\lambda y+\mu z\rangle_B
    =
    \lambda\langle x,y\rangle_B
    +
    \mu\langle x,z\rangle_B;
  \]
  \item
  \[
    \langle x,yb\rangle_B
    =
    \langle x,y\rangle_Bb;
  \]
  \item
  \[
    \langle x,y\rangle_B^*
    =
    \langle y,x\rangle_B;
  \]
  \item
  \[
    \langle x,x\rangle_B\geq0;
  \]
  \item
  \[
    \langle x,x\rangle_B=0
    \quad\Longleftrightarrow\quad
    x=0.
  \]
\end{enumerate}
The formula
\[
  \|x\|_X
  =
  \|\langle x,x\rangle_B\|^{1/2}
\]
defines a norm on \(X\). If \(X\) is complete with respect to this
norm, then \(X\) is called a \textnormal{Hilbert \(B\)-module}.
\end{definition}
\begin{proposition}[Cauchy--Schwarz inequality]
  \label{prop:HilbertModuleCS}
  Let \(X\) be a pre-Hilbert \(B\)-module. Then
  \[
    \langle x,y\rangle_B^*\langle x,y\rangle_B
    \leq
    \|\langle x,x\rangle_B\|\langle y,y\rangle_B
  \]
  for all \(x,y\in X\). Consequently,
  \[
    \|\langle x,y\rangle_B\|
    \leq
    \|x\|_X\|y\|_X,
  \]
  and the formula in the preceding definition does indeed define a norm on \(X\).
\end{proposition}
\begin{proof}
  Put
  \[
    a=\langle x,x\rangle_B,
    \qquad
    c=\langle x,y\rangle_B,
    \qquad
    d=\langle y,y\rangle_B.
  \]
  If \(\|a\|=0\), then \(x=0\) and the assertion is immediate. Suppose that
  \(\lambda=\|a\|>0\). Since \(a\leq\lambda 1_{\widetilde B}\), positivity of the
  inner product gives
  \begin{align*}
    0
    &\leq
    \langle xc-y\lambda,xc-y\lambda\rangle_B\\
    &=c^*ac-2\lambda c^*c+\lambda^2d\\
    &\leq
    \lambda c^*c-2\lambda c^*c+\lambda^2d
     =\lambda(\lambda d-c^*c).
  \end{align*}
  Hence \(c^*c\leq\lambda d\), which is the desired inequality. Taking norms
  gives the second assertion. Finally,
  \[
    \|x+y\|_X^2
    \leq
    \|x\|_X^2+2\|x\|_X\|y\|_X+\|y\|_X^2,
  \]
  and therefore the triangle inequality follows. The remaining norm axioms are
  immediate from the definition.
\end{proof}
\begin{definition}
  A \textnormal{left \(A\)-module} is defined analogously, using a bilinear
  map \(A\times X\to X\), \((a,x)\mapsto ax\). Let \(A\) and \(B\) be
  algebras. An \textnormal{\(A\)-\(B\) bimodule} is a complex vector space
  \(X\) which is both a left \(A\)-module and a right \(B\)-module, and whose
  actions satisfy
  \[
    (ax)b=a(xb)
  \]
  for every \(a\in A\), \(x\in X\) and \(b\in B\).
\end{definition}
\begin{definition}
A Hilbert \(B\)-module \(X\) is called \textnormal{full} if
\[
  \overline{\operatorname{span}}
  \{\langle x,y\rangle_B:x,y\in X\}
  =
  B.
\]
\end{definition}

\begin{example}
  \begin{enumerate}
    \item \(B\) is itself a Hilbert \(B\)-module with \(\langle a,b\rangle_B=a^*b\).
      More generally, any (closed) right ideal of \(B\) is a Hilbert \(B\)-module.
    \item If \(E_1,\ldots,E_n\) are (pre-)Hilbert \(B\)-modules, their finite
      direct sum is a (pre-)Hilbert \(B\)-module with inner product
      \[
        \left\langle\bigoplus_{i=1}^n x_i,\bigoplus_{i=1}^n y_i\right\rangle_B
        =
        \sum_{i=1}^n\langle x_i,y_i\rangle_B.
      \]
      If every \(E_i\) is complete, then so is the finite direct sum.
    \item More generally, if \((E_i)_{i=1}^\infty\) is a sequence of Hilbert
      \(B\)-modules, its Hilbert direct sum is
      \[
        \bigoplus_{i=1}^\infty E_i
        =
        \left\{(x_i):
        \sum_{i=1}^\infty\langle x_i,x_i\rangle_B
        \text{ converges in norm}\right\},
      \]
      with inner product
      \[
        \langle(x_i),(y_i)\rangle_B
        =
        \sum_{i=1}^\infty\langle x_i,y_i\rangle_B.
      \]
      This is again a Hilbert \(B\)-module.
    \item The \textnormal{standard Hilbert \(B\)-module} is
      \[
        \mathcal H_B
        =
        \left\{(b_n):
        \sum_{n=1}^\infty b_n^*b_n\text{ converges in norm}\right\},
      \]
      with
      \[
        \langle(a_n),(b_n)\rangle_B
        =
        \sum_{n=1}^\infty a_n^*b_n.
      \]
      It is also denoted by \(B^\infty\). Similarly, one may form \(B^n\),
      and, for any Hilbert \(B\)-module \(X\), the modules \(X^n\) and
      \(X^\infty\).
  \end{enumerate}
\end{example}

\begin{definition}
  Let \(X\) be a Hilbert \(B\)-module. A \(B\)-linear map \(T:X\to X\)
  is called \textnormal{adjointable} if there exists a \(B\)-linear map
  \(T^*:X\to X\) such that
  \[
    \langle Tx,y\rangle_B
    =
    \langle x,T^*y\rangle_B
  \]
  for every \(x,y\in X\). The set of adjointable operators on \(X\) is
  denoted by \(\mathcal L_B(X)\).
\end{definition}
\begin{proposition}
  The adjoint of an adjointable operator is unique, and every adjointable
  operator is bounded. With the operator norm and the adjoint operation,
  \(\mathcal L_B(X)\) is a unital \cstar.
\end{proposition}
\begin{proof}
  If \(S\) and \(T\) are both adjoints of the same operator, then
  \(\langle x,(S-T)y\rangle_B=0\) for all \(x,y\in X\). Taking
  \(x=(S-T)y\) shows that \(S=T\). To prove boundedness, suppose that
  \(x_n\to x\) and \(Tx_n\to y\). For every \(z\in X\),
  \[
    \langle z,y\rangle_B
    =
    \lim_n\langle T^*z,x_n\rangle_B
    =
    \langle z,Tx\rangle_B,
  \]
  and hence \(y=Tx\). Thus \(T\) has closed graph and is bounded. The
  remaining \cstar\ properties follow from the corresponding operator
  identities, in particular \(\|T^*T\|=\|T\|^2\).
\end{proof}
\begin{definition}
  Let \(X\) be a Hilbert \(B\)-module. For \(x,y\in X\), define
  \[
    \Theta_{x,y}:X\longrightarrow X,
    \qquad
    \Theta_{x,y}(z)=x\langle y,z\rangle_B.
  \]
  The algebra of \textnormal{compact operators} on \(X\) is
  \[
    \mathcal K_B(X)
    =
    \overline{\operatorname{span}}
    \{\Theta_{x,y}:x,y\in X\}
    \subseteq\mathcal L_B(X),
  \]
  where the closure is taken in the operator norm.
\end{definition}
\begin{proposition}
  For all \(x,y\in X\), the operator \(\Theta_{x,y}\) is adjointable and
  \[
    \Theta_{x,y}^*=\Theta_{y,x}.
  \]
  Moreover, \(\mathcal K_B(X)\) is a closed two-sided ideal of
  \(\mathcal L_B(X)\).
\end{proposition}
\begin{proof}
  For \(z,w\in X\),
  \[
    \langle\Theta_{x,y}z,w\rangle_B
    =
    \langle z,y\rangle_B^*\langle x,w\rangle_B
    =
    \langle z,\Theta_{y,x}w\rangle_B.
  \]
  Hence \(\Theta_{x,y}^*=\Theta_{y,x}\). If
  \(T\in\mathcal L_B(X)\), then
  \[
    T\Theta_{x,y}=\Theta_{Tx,y},
    \qquad
    \Theta_{x,y}T=\Theta_{x,T^*y},
  \]
  which proves the ideal assertion.
\end{proof}
\begin{definition}
Let \(A\) and \(B\) be \cstars. An
\textnormal{\(A\)-\(B\) imprimitivity bimodule} is a full right
Hilbert \(B\)-module \(X\), together with a \(^*\)-isomorphism
\[
  \varphi:A\longrightarrow\mathcal K_B(X).
\]
The corresponding left action of \(A\) on \(X\) is defined by
\(ax=\varphi(a)x\).
The algebras \(A\) and \(B\) are said to be
\textnormal{strongly Morita equivalent} if there exists an
\(A\)-\(B\) imprimitivity bimodule. In this case we write
\(A\sim_M B\).
\end{definition}
\begin{example}
  Every \cstar\ \(B\) is a \(B\)-\(B\) imprimitivity bimodule with
  \[
    \langle x,y\rangle_B=x^*y.
  \]
  Indeed, left multiplication identifies \(B\) with \(\mathcal K_B(B)\).
\end{example}

The following construction is the interior tensor product of Hilbert modules;
see \cite[Section~13.5]{Blackadar1998}.  In the language of induced
representations, it is Rieffel's induction in stages
\cite[Theorem~5.9]{Rieffel1974}.

\begin{proposition}\label{prop:MoritaTransitivity}
  Strong Morita equivalence is transitive. More precisely, if \(X\) is an
  \(A\)-\(B\) imprimitivity bimodule and \(Y\) is a \(B\)-\(C\)
  imprimitivity bimodule, then the interior tensor product
  \[
    X\otimes_B Y
  \]
  is an \(A\)-\(C\) imprimitivity bimodule. Consequently,
  \[
    A\sim_M B\quad\text{and}\quad B\sim_M C
    \qquad\Longrightarrow\qquad
    A\sim_M C.
  \]
\end{proposition}

\begin{proof}
  Let
  \[
    \psi:B\longrightarrow\mathcal K_C(Y)
  \]
  be the isomorphism defining the left action of \(B\) on \(Y\). Form the
  algebraic balanced tensor product \(X\odot_B Y\), so that
  \[
    xb\otimes y=x\otimes\psi(b)y.
  \]
  On elementary tensors put
  \begin{equation}\label{eq:InteriorTensorInnerProduct}
    \langle x_1\otimes y_1,x_2\otimes y_2\rangle_C
    =
    \langle y_1,\psi(\langle x_1,x_2\rangle_B)y_2\rangle_C,
  \end{equation}
  and extend sesquilinearly. The balancing relations and the Hilbert-module
  identities show that this is well defined.

  To prove positivity, let \(z=\sum_{i=1}^n x_i\otimes y_i\). The matrix
  \[
    \big(\langle x_i,x_j\rangle_B\big)_{i,j}
  \]
  is positive in \(M_n(B)\). Its image under \(\psi^{(n)}\) is therefore a
  positive adjointable operator on \(Y^n\), and hence
  \[
    \langle z,z\rangle_C
    =
    \sum_{i,j=1}^n
    \langle y_i,\psi(\langle x_i,x_j\rangle_B)y_j\rangle_C
    \geq0.
  \]
  After quotienting by the vectors of length zero and completing, we obtain a
  Hilbert \(C\)-module, denoted by \(X\otimes_B Y\).

  This module is full. Indeed, the inner products \(\langle x,u\rangle_B\)
  span a dense subspace of \(B\), while
  \(\psi(B)=\mathcal K_C(Y)\) and \(\mathcal K_C(Y)Y\) is dense in \(Y\).
  Formula~\eqref{eq:InteriorTensorInnerProduct} and the fullness of \(Y\)
  therefore give
  \[
    \overline{\operatorname{span}}
    \langle X\otimes_B Y,X\otimes_B Y\rangle_C=C.
  \]

  The action of \(A\) is given on elementary tensors by
  \[
    a(x\otimes y)=(ax)\otimes y.
  \]
  Equivalently, the isomorphism
  \(A\cong\mathcal K_B(X)\) is followed by the homomorphism
  \[
    \rho:\mathcal K_B(X)\longrightarrow
    \mathcal L_C(X\otimes_B Y),
    \qquad
    \rho(T)=T\otimes1.
  \]
  Indeed, the definition on elementary tensors respects the balancing
  relation, and the matrix estimate used above gives
  \(\|T\otimes1\|\leq\|T\|\); moreover,
  \((T\otimes1)^*=T^*\otimes1\). Thus \(\rho\) is a well-defined
  *-homomorphism.
  This homomorphism is injective. In fact, if \(T\otimes1=0\), then for every
  \(x\in X\) and \(y\in Y\),
  \[
    0
    =\langle Tx\otimes y,Tx\otimes y\rangle_C
    =\langle y,\psi(\langle Tx,Tx\rangle_B)y\rangle_C.
  \]
  Faithfulness of \(\psi\) then implies \(Tx=0\) for every \(x\), and hence
  \(T=0\).

  It remains to identify the range of \(\rho\). Define the left
  \(B\)-valued inner product on \(Y\) by
  \[
    {}_B\langle y_1,y_2\rangle
    =\psi^{-1}(\Theta_{y_1,y_2}).
  \]
  A direct calculation on elementary tensors gives
  \begin{equation}\label{eq:TensorRankOneOperators}
    \Theta_{x_1\otimes y_1,x_2\otimes y_2}
    =
    \rho\!\left(
      \Theta_{x_1{}_B\langle y_1,y_2\rangle,x_2}
    \right).
  \end{equation}
  Since \(\rho\) is injective, it is isometric and its range is closed.
  Thus \(\mathcal K_C(X\otimes_B Y)\) is contained in the range of \(\rho\).
  Conversely, the linear span of the elements
  \({}_B\langle y_1,y_2\rangle\) is dense in \(B\). For fixed \(x_1\), we may
  therefore choose a net \((b_\lambda)\) in this span such that
  \(x_1b_\lambda\to x_1\). Then
  \[
    \Theta_{x_1b_\lambda,x_2}
      \longrightarrow\Theta_{x_1,x_2}
  \]
  in norm, and \eqref{eq:TensorRankOneOperators} shows that
  \(\rho(\Theta_{x_1b_\lambda,x_2})\) is compact. Hence
  \(\rho(\mathcal K_B(X))\subseteq\mathcal K_C(X\otimes_B Y)\).
  We conclude that
  \[
    A\cong\mathcal K_B(X)
      \xrightarrow[\cong]{\rho}
    \mathcal K_C(X\otimes_B Y).
  \]
  Together with fullness, this proves that \(X\otimes_B Y\) is an
  \(A\)-\(C\) imprimitivity bimodule.
\end{proof}

\section{Full Corners, Linking Algebras and Traces}

\begin{definition}
  Let \(A\) be a unital \cstar\ and let \(p\in A\) be a projection. The
  subalgebra
  \[
    pAp=\{pap:a\in A\}
  \]
  is called a \textnormal{corner} of \(A\). The projection \(p\), and the
  corresponding corner \(pAp\), are called \textnormal{full} if
  \[
    \overline{\operatorname{span}}
    \{apb:a,b\in A\}
    =A.
  \]
  Equivalently, \(pAp\) is not contained in any proper closed two-sided
  ideal of \(A\).
\end{definition}

The following result is \cite[Proposition~2.1]{Rieffel1981}.
\begin{proposition}\label{prop:RieffelFullCorners}
  Let \(A\) and \(B\) be unital \cstars. If \(A\) and \(B\) are strongly
  Morita equivalent, then there exist positive integers \(m,n\) and full
  projections
  \[
    p\in M_n(A),
    \qquad
    q\in M_m(B),
  \]
  such that
  \[
    B\cong pM_n(A)p,
    \qquad
    A\cong qM_m(B)q.
  \]
\end{proposition}
\begin{proof}
  Let \(X\) be an \(A\)-\(B\) imprimitivity bimodule, and let
  \[
    \varphi:A\longrightarrow\mathcal K_B(X)
  \]
  be the corresponding \(^*\)-isomorphism. For \(\xi,\eta\in X\), put
  \[
    {}_A\langle\xi,\eta\rangle
    =
    \varphi^{-1}(\Theta_{\xi,\eta}).
  \]
  Then
  \[
    {}_A\langle\xi,\eta\rangle\zeta
    =
    \xi\langle\eta,\zeta\rangle_B
  \]
  for all \(\xi,\eta,\zeta\in X\).

  Since \(X\) is full, the linear span of the range of
  \(\langle\cdot,\cdot\rangle_B\) is dense in \(B\). We may therefore
  choose \(x_1,\ldots,x_n,y_1,\ldots,y_n\in X\) such that
  \[
    s=\sum_{i=1}^n\langle x_i,y_i\rangle_B
    \qquad\text{satisfies}\qquad
    \|s-1_B\|<1.
  \]
  The element \(s\) is invertible. Replacing each \(y_i\) by
  \(y_i s^{-1}\), we obtain
  \[
    \sum_{i=1}^n\langle x_i,y_i\rangle_B=1_B.
  \]

  Put \(E=M_n(A)\). The module \(X^n\) is an \(E\)-\(B\) imprimitivity
  bimodule with inner products
  \[
    \langle\xi,\eta\rangle_B
    =
    \sum_{i=1}^n\langle\xi_i,\eta_i\rangle_B
  \]
  and
  \[
    {}_E\langle\xi,\eta\rangle
    =
    \bigl({}_A\langle\xi_i,\eta_j\rangle\bigr)_{i,j=1}^n.
  \]
  Let
  \[
    x=(x_1,\ldots,x_n),
    \qquad
    y=(y_1,\ldots,y_n)
  \]
  in \(X^n\). Then \(\langle x,y\rangle_B=1_B\), and hence also
  \(\langle y,x\rangle_B=1_B\). Consequently,
  \begin{align*}
    1_B
    &=\langle x,y\rangle_B\langle y,x\rangle_B\\
    &=\langle x,y\langle y,x\rangle_B\rangle_B\\
    &=\langle x,{}_E\langle y,y\rangle x\rangle_B
      =\langle z,z\rangle_B,
  \end{align*}
  where
  \[
    z={}_E\langle y,y\rangle^{1/2}x.
  \]
  It follows that
  \begin{align*}
    {}_E\langle z,z\rangle\,{}_E\langle z,z\rangle
    &={}_E\langle{}_E\langle z,z\rangle z,z\rangle\\
    &={}_E\langle z\langle z,z\rangle_B,z\rangle
      ={}_E\langle z,z\rangle.
  \end{align*}
  Thus
  \[
    p={}_E\langle z,z\rangle
  \]
  is a projection in \(E\). Define
  \[
    \phi:B\longrightarrow E,
    \qquad
    \phi(b)={}_E\langle zb,z\rangle.
  \]
  The identities for the two inner products show that \(\phi\) is an
  injective \(^*\)-homomorphism whose range is contained in \(pEp\).

  Since \(X^n\) is full as a left \(E\)-module, the corner \(pEp\) is
  densely spanned by elements of the form
  \(p\,{}_E\langle\xi,\eta\rangle p\), with \(\xi,\eta\in X^n\). Moreover,
  \[
    p\,{}_E\langle\xi,\eta\rangle p
    =
    \phi\bigl(
      \langle z,\xi\rangle_B\langle\eta,z\rangle_B
    \bigr).
  \]
  The range of the injective \(^*\)-homomorphism \(\phi\) is closed, and
  therefore \(\phi(B)=pEp\).

  Finally, for every \(\xi,\eta\in X^n\),
  \[
    {}_E\langle\xi,\eta\rangle
    =
    {}_E\langle\xi,z\rangle
    p
    {}_E\langle z,\eta\rangle.
  \]
  Since the span of the left inner products is dense in \(E\), this gives
  \(\overline{EpE}=E\), so \(p\) is full. Repeating the argument with the
  left and right inner products interchanged gives a positive integer \(m\)
  and a full projection \(q\in M_m(B)\) such that
  \(A\cong qM_m(B)q\).
\end{proof}

Now if \(A\) and \(B\) are \cstars\ and if \(X\) is an \(A\)-\(B\) imprimitivity bimodule,
then every trace on \(A\) can be induced by \(X\) to give a trace on \(B\).
For the general case of possibly unbounded traces, the result is implicit in
\cite[Proposition~28]{Green1978}; see also
\cite[Proposition~2.2]{Rieffel1981} for the unital finite case.
But in the present case of unital \cstars\ and of finite traces, the situation is very simple:
\begin{proposition}\label{prop:Rieffel22}
  Let \(A\) and \(B\) be unital \cstars, and let \(X\) be an
  \(A\)-\(B\) imprimitivity bimodule. Then there is a bijection
  between the (nonnormalized) finite traces on \(A\) and those on \(B\),
  under which to a trace on \(A\) there is associated a trace \(\tau_X\) on \(B\)
  satisfying
  \[
    \tau_X(\langle x,y\rangle_B)=\tau({}_A\langle y,x\rangle)
  \]
  for all \(x,y\in X\).
\end{proposition}
\begin{proof}
  With the notation of the proof of Proposition~\ref{prop:RieffelFullCorners},
  every \(b\in B\) has the form \(b=\langle z,zb\rangle_B\). Thus, if
  \(\tau_X\) exists, it is uniquely determined by the displayed identity. Let
  \(n,z,E\) and \(\phi\) be as in that proof. Extend \(\tau\) to the corresponding
  nonnormalized trace on \(E=M_n(A)\), namely
  \[
    \tau((a_{ij}))=\sum_{i=1}^n\tau(a_{ii}),
  \]
  and define the trace \(\tau_X\) on \(B\) by
  \(\tau_X(b)=\tau(\phi(b))\). We
  show that \(\tau_X\) satisfies the above condition. Let \(x,y\in X\), and view 
  them as the elements \((x,0,\ldots,0)\) and \((y,0,\ldots,0)\) of \(X^n\), so that
  \[
    \tau({}_E\langle y,x\rangle)=\tau({}_A\langle y,x\rangle).
  \]
  Then
  \begin{align*}
    \tau_X(\langle x,y\rangle_B)
    &=\tau({}_E\langle z\langle x,y\rangle_B,z\rangle)\\
    &=\tau({}_E\langle {}_E\langle z,x\rangle y,z\rangle)\\
    &=\tau\bigl({}_E\langle z,x\rangle\,{}_E\langle y,z\rangle\bigr)\\
    &=\tau\bigl({}_E\langle y,z\rangle\,{}_E\langle z,x\rangle\bigr)\\
    &=\tau({}_E\langle y\langle z,z\rangle_B,x\rangle)\\
    &=\tau({}_E\langle y,x\rangle)
      =\tau({}_A\langle y,x\rangle).
  \end{align*}
  Interchanging \(A\) and \(B\) gives the inverse correspondence, so the map
  \(\tau\mapsto\tau_X\) is a bijection.
\end{proof}
The next construction is the linking-algebra construction of
\cite[Theorem~1.1, pp.~350--351]{BrownGreenRieffel1977}.
\begin{lemma}\label{lemma:LinkingAlgebraConstruction}
  Let \(A\) and \(B\) be unital \cstars, and let \(X\) be an
  \(A\)-\(B\) imprimitivity bimodule. Let
  \(\widetilde X\) be the conjugate vector space of \(X\), whose elements are
  denoted by \(\widetilde x\), with bimodule operations
  \[
    b\cdot\widetilde x\cdot a
    =\widetilde{a^*xb^*}
    \qquad(a\in A,\ b\in B,\ x\in X).
  \]
  On the vector space
  \[
    L_0(X)=
    \left\{
      \begin{pmatrix}
        a&x\\
        \widetilde y&b
      \end{pmatrix}
      :a\in A,\ b\in B,\ x,y\in X
    \right\},
  \]
  define multiplication and involution by
  \begin{align*}
    \begin{pmatrix}a&x\\ \widetilde y&b\end{pmatrix}
    \begin{pmatrix}a'&x'\\ \widetilde {y'}&b'\end{pmatrix}
    &=\begin{pmatrix}
      aa'+{}_A\langle x,y'\rangle&ax'+xb'\\
      \widetilde y a'+b\widetilde {y'}&\langle y,x'\rangle_B+bb'
    \end{pmatrix},\\
    \begin{pmatrix}a&x\\ \widetilde y&b\end{pmatrix}^*
    &=\begin{pmatrix}a^*&y\\ \widetilde x&b^*\end{pmatrix}.
  \end{align*}
  The formula
  \[
    \begin{pmatrix}a&x\\ \widetilde y&b\end{pmatrix}
    \begin{pmatrix}z\\ c\end{pmatrix}
    =
    \begin{pmatrix}az+xc\\ \langle y,z\rangle_B+bc\end{pmatrix}
    \qquad(z\in X,\ c\in B)
  \]
  represents \(L_0(X)\) faithfully by adjointable operators on the Hilbert
  \(B\)-module \(X\oplus B\). Its operator-norm completion, denoted by
  \(L(X)\), is a \cstar, called the \textnormal{linking algebra} of \(X\).
  Moreover,
  \[
    p=\begin{pmatrix}1_A&0\\0&0\end{pmatrix}, 
    \qquad
    q=\begin{pmatrix}0&0\\0&1_B\end{pmatrix}
  \]
  are complementary full projections in \(L(X)\), and
  \[
    pL(X)p\cong A,
    \qquad
    qL(X)q\cong B.
  \]
\end{lemma}
\begin{proof}
  The compatibility identity
  \({}_A\langle x,y\rangle z=x\langle y,z\rangle_B\) shows that the
  displayed product is associative, and a direct computation shows that the
  displayed formula defines adjointable operators on \(X\oplus B\), with the
  displayed involution as adjoint. If such an operator is zero, applying it
  first to \((0,1_B)\) gives \(x=b=0\), and then applying it to \((z,0)\),
  for arbitrary \(z\in X\), gives \(a=y=0\). Hence the representation is
  faithful, and its operator-norm completion is a \cstar. The corner
  identities are immediate from the definitions. Finally,
  \[
    \overline{\operatorname{span}}\,qL(X)pL(X)q
    =\overline{\operatorname{span}}\langle X,X\rangle_B=B,
  \]
  while the corresponding equality for the left inner product gives
  \(\overline{\operatorname{span}}\,pL(X)qL(X)p=A\). Hence both \(p\) and
  \(q\) are full, and \(p+q=1_{L(X)}\).
\end{proof}

The following is \cite[Proposition~2.3]{Rieffel1981}.
\begin{proposition}\label{prop:TraceLinkingAlgebra}
  Let \(A\), \(B\) and \(X\) be as above, and let \(L(X)\) be the linking
  algebra of \(X\). Then each finite trace \(\tau\) on \(A\) has a unique
  extension to a finite trace \(\widetilde\tau\) on \(L(X)\). Its restriction
  to \(B\) is \(\tau_X\), and on \(L_0(X)\) it is given by
  \[
    \widetilde\tau
    \begin{pmatrix}a&x\\ \widetilde y&b\end{pmatrix}
    =\tau(a)+\tau_X(b).
  \]
\end{proposition}
\begin{proof}
  Proposition~\ref{prop:Rieffel22} and the multiplication in
  Lemma~\ref{lemma:LinkingAlgebraConstruction} show directly that the displayed
  functional is tracial on \(L_0(X)\); its positivity makes it bounded, so it
  extends to \(L(X)\). Conversely, a trace on \(L(X)\) vanishes on the
  off-diagonal corners \(pL(X)q\) and \(qL(X)p\), while its restrictions to
  the diagonal corners are related by Proposition~\ref{prop:Rieffel22}. This
  proves uniqueness.
\end{proof}
By construction, \(A\) and \(B\) sit as complementary full corners of the
linking algebra \(L(X)\). We recall that if \(\psi\) is any
\(^*\)-homomorphism between \cstars\ (possibly not preserving identity
elements), then \(\psi\) induces a homomorphism between the corresponding
\(K_0\) groups. We now need the following fact, which is
\cite[Proposition~2.4]{Rieffel1981}.
\begin{proposition}\label{prop:K0FullCorner}
  Let \(A\) be a unital \cstar, let \(pAp\) be a full corner of \(A\), and let \(\psi\) be the 
  injection of \(pAp\) into \(A\). Then \(K_0(\psi)\) is an isomorphism of 
  \(K_0(pAp)\) with \(K_0(A)\).
\end{proposition}
\begin{proof}
  View \(X=pA\) as a \(pAp\)-\(A\) imprimitivity bimodule
  \cite[Section~6.8]{Rieffel1974}.
  Then, as in the proof of \ref{prop:RieffelFullCorners}, we can find 
  \(a_1,\ldots,a_n\in A\) such that \(\sum a_i^* pa_i=1\). Let \(\phi\)
  be the corresponding map of \(A\) into \(M_n(pAp)\), so that \(\phi(a)=(pa_iaa_j^*p)_{i,j}\),
  and \(\phi\) maps \(A\) onto the corner of \(M_n(pAp)\) defined by the projection
  \((pa_ia_j^*p)_{i,j}=P\). Let \(V\) be the element of \(M_n(A)\) whose
  first column consists of \(pa_1,\ldots,pa_n\), and all of whose other entries are \(0\).
  Then a simple calculation shows that \(VV^*=P\), while \(V^*V\) is the matrix with \(1\) in the upper
  left corner, and \(0\) elsewhere. Thus, ``conjugation'' of \(M_n(A)\) by \(V\)
  carries \(\phi(A)\) onto the corner of \(M_n(A)\) consisting of the matrices all of whose
  entries are zero except that in the upper left corner. If we view \(\psi\)
  as giving also the inclusion of \(M_n(pAp)\) into \(M_n(A)\), we see in this way that 
  \(K_0(\psi)\circ K_0(\phi)\) is the identity map on \(K_0(A)\).
  (We must use here the fact that, as remarked in \cite[immediately after Lemma~3.6]{EffrosRosenberg1978}, it does not matter whether
  one uses unitary or Murray--von Neumann equivalence in defining \(K_0\).) It follows 
  that \(K_0(\psi)\) is surjective. Thus we have shown that the inclusion map of a full corner
  into an algebra induces a surjection of \(K_0\) groups.
  But \(\phi(A)\) is a full corner of \(M_n(pAp)\), and \(\phi\) is an isomorphism
  of \(A\) with \(\phi(A)\). It follows that \(K_0(\phi)\) is surjective. Since 
  \(K_0(\psi)\circ K_0(\phi)\) is an isomorphism, it follows that \(K_0(\psi)\) must 
  be injective, and so is an isomorphism.
\end{proof}
Now let again \(X\) be an \(A\)-\(B\) imprimitivity bimodule. Let
\(K_0(\psi)\) and \(K_0(\omega)\) denote the isomorphisms of \(K_0(A)\) and
\(K_0(B)\) with \(K_0(L(X))\) obtained from the inclusions of \(A\) and \(B\)
as corners of \(L(X)\). We let \(\Phi_X\) denote the isomorphism
\[
  \Phi_X=(K_0(\omega))^{-1}\circ K_0(\psi):K_0(A)\longrightarrow K_0(B).
\]
We remark in passing that the Picard group of a \cstar\ \(D\), as defined in
\cite{BrownGreenRieffel1977}, acts by automorphisms of \(K_0(D)\).

\begin{proposition}\label{prop:TraceMoritaK0}
  Let \(A\) and \(B\) be unital \cstars, let \(X\) be an \(A\)-\(B\) imprimitivity bimodule,
  let \(\tau\) be a finite trace on \(A\) and let \(\tau_X\) be the corresponding (nonnormalized) trace
  on \(B\) defined above. Let \(\Phi_X\) denote the isomorphism of \(K_0(A)\) onto
  \(K_0(B)\) determined by \(X\) as above, and let \(\tau_*\) and \((\tau_X)_*\) be the 
  homomorphisms of \(K_0(A)\) and \(K_0(B)\) into \(\mathbb{R}\) determined by \(\tau\) and \(\tau_X\). Then
  \[
    (\tau_X)_*\circ\Phi_X=\tau_*.
  \]
\end{proposition}
\begin{proof}
  This follows immediately from the definition and the fact that \(\tau_X\) is the restriction to
  \(B\) of the unique extension of \(\tau\) to the linking algebra.
\end{proof}
\begin{corollary}\label{cor:TraceRangeMorita}
  Let \(A, B\) and \(X\) be as above, let \(\tau\) be a trace on \(A\), and let \(\tau_X\)
  be the corresponding trace on \(B\). Then
  \[
    \tau_*(K_0(A))=(\tau_X)_*(K_0(B)).
  \]
\end{corollary}

\section{The Basic Equivalence
  \texorpdfstring{\(A_\theta\sim_M A_{-1/\theta}\)}
  {A theta Morita equivalent to A -1/theta}}

We first prove the general transformation-group theorem from which the
equivalence of irrational rotation algebras follows.  The construction is due
to Rieffel; see \cite[Sections~1--2]{Rieffel1976}.  Its positivity and
completion machinery is developed in
\cite[Sections~4, 6--7]{Rieffel1974}.  We include the details because the
explicit bimodule will be used below.

Let a locally compact group \(K\) act on the left of a locally compact
Hausdorff space \(P\), and let a locally compact group \(H\) act on the right.
We write the actions as \(t p\) and \(p s\), and assume that they commute:
\[
  t(ps)=(tp)s.
\]
The induced actions on the orbit spaces are
\[
  \operatorname{lt}_t(\varphi)(pH)=\varphi(t^{-1}pH),
  \qquad
  \operatorname{rt}_s(\psi)(Kp)=\psi(Kps).
\]

\begin{lemma}[Rieffel's approximation lemma]
  \label{lem:RieffelApproximation}
  Suppose that both actions are free and proper. Put
  \[
    E_0=C_C(K\times P/H),\qquad
    B_0=C_C(H\times K\backslash P),\qquad
    X_0=C_C(P).
  \]
  For \(f,g\in X_0\), set
  \begin{align}
    {}_{E_0}\!\langle f,g\rangle(t,pH)
      &=\Delta_K(t)^{-1/2}
        \int_H f(ps)\overline{g(t^{-1}ps)}\,ds,
      \label{eq:RieffelGeneralLeftInner}\\
    \langle f,g\rangle_{B_0}(s,Kp)
      &=\Delta_H(s)^{-1/2}
        \int_K\overline{f(t^{-1}p)}g(t^{-1}ps)\,dt.
      \label{eq:RieffelGeneralRightInner}
  \end{align}
  There is a net \((e_i)\) in \(E_0\) such that
  \begin{enumerate}
    \item \(e_i*c\to c\) in the inductive-limit topology of \(E_0\) for
      every \(c\in E_0\);
    \item \(e_i\cdot f\to f\) in the inductive-limit topology of \(X_0\)
      for every \(f\in X_0\);
    \item each \(e_i\) has the form
      \[
        e_i=\sum_{j=1}^{n_i}{}_{E_0}\!\langle f_{i,j},f_{i,j}\rangle,
        \qquad f_{i,j}\in X_0.
      \]
  \end{enumerate}
  The analogous assertion holds after interchanging \(K\) and \(H\).
\end{lemma}

\begin{proof}
  We reproduce the cutoff construction used by Rieffel.  Let
  \((b_l)\subseteq C_C(P/H)_+\) be a contractive approximate identity.
  Direct quadruples \((Q,U,l,\epsilon)\) by increasing compact sets
  \(Q\subseteq P/H\), decreasing relatively compact neighborhoods \(U\) of
  the identity of \(K\), increasing \(l\), and decreasing \(\epsilon>0\).
  Choose a compact set \(D\subseteq P\) whose image contains \(Q\), and a
  compact neighborhood \(C\) of \(D\).

  We first record where freeness and properness enter.  If \(N\) is a
  neighborhood of the identity of \(K\), every \(p\in P\) has a neighborhood
  \(V\) such that
  \[
    \{t\in K:tV\cap V\ne\varnothing\}\subseteq N.
    \tag{\(*\)}
  \]
  Otherwise there would be nets \(p_i\to p\) and \(t_i p_i\to p\), with
  \(t_i\notin N\); properness gives a convergent subnet of \((t_i)\), and
  freeness forces its limit to be the identity, a contradiction.
  Thus, after shrinking \(U\), there are relatively compact open sets
  \(V_1,\ldots,V_n\subseteq C\) covering \(D\) and satisfying \((*)\) with
  \(N=U\).

  Choose \(h_j\in C_C(V_j)_+\) with \(\sum_jh_j=1\) on \(D\), and put
  \[
    q(p)=\int_H\sum_jh_j(ps)\,ds.
  \]
  This is \(H\)-invariant and strictly positive on \(DH\).  After inserting
  a cutoff away from \(DH\), the functions
  \(k_j=h_j/q\) may therefore be chosen in \(C_C(V_j)_+\) so that
  \begin{equation}\label{eq:RieffelCutoff}
    \int_H\sum_jk_j(ps)\,ds=1\quad(pH\in Q),
    \qquad
    \int_H\sum_jk_j(ps)\,ds\leq1\quad(p\in P).
  \end{equation}
  Let \(C_H\subseteq H\) be compact and contain every \(s\) for which
  \(Cs\cap C\ne\varnothing\).  Given
  \(0<\delta<\min\{\epsilon/3,1/2\}\), the usual local normalization for a
  free proper action gives functions \(g_j\in C_C(V_j)_+\) satisfying
  \begin{equation}\label{eq:RieffelLocalNormalization}
    \left|k_j(p)-g_j(p)\int_Kg_j(t^{-1}p)\,dt\right|
       <\frac{\delta}{n\mu_H(C_H)}
       \qquad(p\in P).
  \end{equation}
  Here is the normalization argument.  Put
  \[
    Q_j(p)=\int_Kk_j(t^{-1}p)\,dt.
  \]
  This is a continuous \(K\)-invariant function.  For \(\eta>0\), define
  \[
    g_{j,\eta}(p)=\frac{k_j(p)}{(Q_j(p)+\eta)^{1/2}}.
  \]
  Then \(g_{j,\eta}\in C_C(V_j)_+\), and \(K\)-invariance of \(Q_j\) gives
  \[
    g_{j,\eta}(p)\int_Kg_{j,\eta}(t^{-1}p)\,dt
      =k_j(p)\frac{Q_j(p)}{Q_j(p)+\eta}.
  \]
  The right-hand side converges uniformly to \(k_j\) as \(\eta\downarrow0\):
  on the set where \(k_j\) is bounded away from zero, \(Q_j\) is bounded
  away from zero, while near the zero set of \(k_j\) the error is bounded
  by \(k_j\).  Choosing \(\eta\) sufficiently small proves
  \eqref{eq:RieffelLocalNormalization}.  Property \((*)\) keeps the support
  in \(V_j\).

  Define
  \[
    F(t,p)=\sum_jg_j(p)g_j(t^{-1}p).
  \]
  Then \(F(t,p)=0\) unless \(t\in U\) and \(p\in C\).  Integrating
  \eqref{eq:RieffelLocalNormalization}, using Fubini and
  \eqref{eq:RieffelCutoff}, gives
  \begin{align}
    \left|\int_K\int_HF(t,ps)\,ds\,dt-1\right|&<\delta
      &&(pH\in Q),\label{eq:RieffelFMass}\\
    0\leq\int_K\int_HF(t,ps)\,ds\,dt&\leq2
      &&(p\in P).\label{eq:RieffelFBound}
  \end{align}

  Choose \(\varphi\in C_C(P)_+\) equal to \(1\) on \(C\), and set
  \(z(p)=\varphi(p)b_l(pH)^{1/2}\).  By shrinking \(U\) once more, uniform
  continuity gives
  \[
    |z(p)\overline{z(t^{-1}p)}-b_l(pH)|<\delta
    \qquad(t\in U,\ p\in C).
  \]
  Put \(f_j=g_jz\) and
  \[
    e(t,pH)=\sum_j{}_{E_0}\!\langle f_j,f_j\rangle(t,pH),
    \qquad
    \widetilde e(t,pH)=\Delta_K(t)^{1/2}e(t,pH).
  \]
  Direct substitution yields
  \[
    \widetilde e(t,pH)
      =\int_HF(t,ps)z(ps)\overline{z(t^{-1}ps)}\,ds.
  \]
  Consequently
  \begin{align}
    \widetilde e(t,pH)&=0 &&(t\notin U),\label{eq:RieffelApproxSupport}\\
    \int_K|\widetilde e(t,pH)|\,dt&\leq4,\label{eq:RieffelApproxBound}\\
    \left|\int_K\widetilde e(t,pH)\,dt-b_l(pH)\right|
      &<\epsilon &&(pH\in Q).\label{eq:RieffelApproxMass}
  \end{align}
  Indeed, the first assertion follows from the support of \(F\), the second
  from \eqref{eq:RieffelFBound}, and the third by adding and subtracting
  \(b_l(pH)\int_K\int_HF(t,ps)\,ds\,dt\) and applying
  \eqref{eq:RieffelFMass}.

  If \(f\in X_0\), then
  \begin{align*}
    (e\cdot f)(p)-f(p)
      &=\int_K\widetilde e(t,pH)
        \bigl(f(t^{-1}p)-f(p)\bigr)\,dt\\
      &\quad+f(p)\left(\int_K\widetilde e(t,pH)\,dt-b_l(pH)\right)
        +b_l(pH)f(p)-f(p).
  \end{align*}
  The three terms tend uniformly to zero by
  \eqref{eq:RieffelApproxSupport}--\eqref{eq:RieffelApproxMass}, uniform
  continuity, and the approximate-identity property of \(b_l\).  Their
  supports remain in a fixed compact set.  Thus \(e_i\cdot f\to f\) in the
  inductive-limit topology.  The identical estimate, applied uniformly to
  the compactly supported family \(p\mapsto c(v,pH)\), gives
  \(e_i*c\to c\).  Each \(e_i\) is a finite sum of the required inner
  products by construction.  Symmetry gives the assertion for \(B_0\).
\end{proof}

\begin{theorem}[Rieffel's transformation-group equivalence]
  \label{thm:RieffelTransformationEquivalence}
  Let \(K\) and \(H\) act freely and properly on the left and right,
  respectively, of a locally compact Hausdorff space \(P\), and suppose that
  the actions commute. Then
  \[
    C_0(P/H)\rtimes_{\operatorname{lt}}K
      \sim_M
    C_0(K\backslash P)\rtimes_{\operatorname{rt}}H.
  \]
  The imprimitivity bimodule is obtained by completing \(C_C(P)\).
\end{theorem}

\begin{proof}
  Put
  \[
    E=C_0(P/H)\rtimes_{\operatorname{lt}}K,
    \qquad
    B=C_0(K\backslash P)\rtimes_{\operatorname{rt}}H,
  \]
  and retain \(E_0,B_0,X_0\) from
  Lemma~\ref{lem:RieffelApproximation}.  Their convolution operations are
  \begin{align*}
    (c_1*c_2)(t,pH)
      &=\int_Kc_1(u,pH)c_2(u^{-1}t,u^{-1}pH)\,du,\\
    c^*(t,pH)
      &=\Delta_K(t^{-1})\overline{c(t^{-1},t^{-1}pH)},\\
    (b_1*b_2)(s,Kp)
      &=\int_Hb_1(v,Kp)b_2(v^{-1}s,Kpv)\,dv,\\
    b^*(s,Kp)
      &=\Delta_H(s^{-1})\overline{b(s^{-1},Kps)}.
  \end{align*}
  For \(c\in E_0\), \(b\in B_0\) and \(f\in X_0\), define
  \begin{align}
    (c\cdot f)(p)
      &=\int_Kc(t,pH)f(t^{-1}p)\Delta_K(t)^{1/2}\,dt,
      \label{eq:RieffelGeneralLeftAction}\\
    (f\cdot b)(p)
      &=\int_Hf(ps)b(s^{-1},Kps)\Delta_H(s)^{-1/2}\,ds.
      \label{eq:RieffelGeneralRightAction}
  \end{align}
  Use \eqref{eq:RieffelGeneralLeftInner} and
  \eqref{eq:RieffelGeneralRightInner} for the two inner products.
  Properness makes all four expressions compactly supported in the required
  variables, and dominated convergence gives continuity.

  Fubini's theorem, left invariance of Haar measure, and the substitutions
  \(u\mapsto t^{-1}u\) and \(v\mapsto sv\) give
  \begin{align*}
    (c_1*c_2)\cdot f&=c_1\cdot(c_2\cdot f),
      &(f\cdot b_1)\cdot b_2&=f\cdot(b_1*b_2),\\
    (c\cdot f)\cdot b&=c\cdot(f\cdot b),
      &\langle f,g\cdot b\rangle_{B_0}
        &=\langle f,g\rangle_{B_0}*b,\\
    \langle c\cdot f,g\rangle_{B_0}
      &=\langle f,c^*\cdot g\rangle_{B_0},
      &\langle f,g\rangle_{B_0}^*
        &=\langle g,f\rangle_{B_0},\\
    {}_{E_0}\!\langle c\cdot f,g\rangle
      &=c*{}_{E_0}\!\langle f,g\rangle,
      &{}_{E_0}\!\langle f,g\rangle^*
        &={}_{E_0}\!\langle g,f\rangle.
  \end{align*}
  The compatibility relation follows from the explicit computation
  \begin{align*}
    \bigl({}_{E_0}\!\langle f,g\rangle\cdot h\bigr)(p)
      &=\int_K\int_H
        f(ps)\overline{g(t^{-1}ps)}h(t^{-1}p)\,ds\,dt\\
      &=\int_Hf(ps)\langle g,h\rangle_{B_0}(s^{-1},Kps)
        \Delta_H(s)^{-1/2}\,ds\\
      &=\bigl(f\cdot\langle g,h\rangle_{B_0}\bigr)(p).
  \end{align*}
  Thus
  \begin{equation}\label{eq:RieffelGeneralCompatibility}
    {}_{E_0}\!\langle f,g\rangle\cdot h
      =f\cdot\langle g,h\rangle_{B_0}.
  \end{equation}

  Positivity requires the approximation lemma.  Write
  \(e_i=\sum_j{}_{E_0}\!\langle f_{i,j},f_{i,j}\rangle\).  Since
  \(e_i\cdot f\to f\), continuity of the right inner product and
  \eqref{eq:RieffelGeneralCompatibility} give, in the norm of \(B\),
  \[
    \langle f,f\rangle_B
      =\lim_i\sum_j
        \langle f_{i,j},f\rangle_B^*
        \langle f_{i,j},f\rangle_B\geq0.
  \]
  Symmetry proves positivity of the \(E\)-valued inner product.  If
  \(\langle f,f\rangle_B=0\), Cauchy--Schwarz gives
  \(\langle g,f\rangle_B=0\) for every \(g\), hence \(e_i\cdot f=0\);
  passing to the limit gives \(f=0\).  Thus the right inner product is
  definite.

  The lemma also proves fullness.  For \(c\in E_0\),
  \[
    e_i*c
      =\sum_j{}_{E_0}\!\langle f_{i,j},c^*\cdot f_{i,j}\rangle
      \longrightarrow c.
  \]
  Therefore the left inner products span a dense subspace of \(E\), and
  symmetry gives the corresponding assertion for \(B\).

  Complete \(X_0\) in
  \(\|f\|_B=\|\langle f,f\rangle_B\|^{1/2}\), obtaining a full Hilbert
  \(B\)-module \(X\).  Define
  \[
    (M(\varphi)f)(p)=\varphi(pH)f(p),
    \qquad
    (V_tf)(p)=\Delta_K(t)^{1/2}f(t^{-1}p).
  \]
  Substitution in \eqref{eq:RieffelGeneralRightInner} shows
  \[
    \langle M(\varphi)f,g\rangle_B
      =\langle f,M(\overline\varphi)g\rangle_B,
    \qquad
    \langle V_tf,V_tg\rangle_B=\langle f,g\rangle_B.
  \]
  Thus \(M(\varphi)\) and \(V_t\) extend to adjointable operators on \(X\),
  and \((M,V)\) is a covariant representation of
  \((C_0(P/H),K,\operatorname{lt})\).  Its integrated form is precisely
  \eqref{eq:RieffelGeneralLeftAction}, so
  \[
    \|c\cdot f\|_B\leq\|c\|\,\|f\|_B.
  \]

  Finally, \eqref{eq:RieffelGeneralCompatibility} gives
  \[
    (M\rtimes V)({}_{E_0}\!\langle f,g\rangle)=\Theta_{f,g}.
  \]
  Fullness and density therefore imply
  \((M\rtimes V)(E)=\mathcal K_B(X)\).  If \(a\in E\) annihilates \(X\),
  then
  \(a*{}_{E}\!\langle f,g\rangle={}_{E}\!\langle a\cdot f,g\rangle=0\).
  Since the left inner products span a dense subspace of \(E\), \(a=0\).
  Hence \(M\rtimes V\) identifies \(E\) with \(\mathcal K_B(X)\), and
  \(X\) is an \(E\)-\(B\) imprimitivity bimodule.
\end{proof}

\begin{corollary}\label{cor:RieffelClosedSubgroups}
  Let \(G\) be a locally compact group and \(H,K\) closed subgroups.  Let
  \(K\) act on \(G/H\) by left translation and \(H\) act on
  \(K\backslash G\) by right translation. Then
  \[
    C_0(G/H)\rtimes K
      \sim_M
    C_0(K\backslash G)\rtimes H.
  \]
\end{corollary}

\begin{proof}
  Left translation by \(K\) and right translation by \(H\) are commuting,
  free and proper actions on \(P=G\).  Apply
  Theorem~\ref{thm:RieffelTransformationEquivalence}.
\end{proof}

\begin{corollary}\label{cor:RieffelBasicEquivalence}
  If \(\theta\) is irrational, then
  \[
    A_\theta\sim_M A_{-1/\theta}.
  \]
\end{corollary}

\begin{proof}
  Assume first that \(\theta>0\) and take
  \[
    G=\mathbb R,\qquad H=\mathbb Z,\qquad K=\theta\mathbb Z
  \]
  in Corollary~\ref{cor:RieffelClosedSubgroups}.  This gives
  \[
    C_0(\mathbb R/\mathbb Z)\rtimes\theta\mathbb Z
      \sim_M
    C_0(\theta\mathbb Z\backslash\mathbb R)\rtimes\mathbb Z.
  \]
  The algebra on the left is \(A_\theta\).  On the right let
  \(z_\theta(r)=e^{2\pi ir/\theta}\) and let \(W\) implement translation by
  \(1\).  Since
  \[
    Wz_\theta=e^{-2\pi i/\theta}z_\theta W,
  \]
  the algebra on the right is \(A_{-1/\theta}\).

  Here all modular functions are \(1\), and the pre-imprimitivity bimodule is
  \(C_C(\mathbb R)\).  The inner products become
  \begin{align*}
    {}_C\!\langle F,G\rangle(m,r)
      &=\sum_{n\in\mathbb Z}
        F(r-n)\overline{G(r-n-m\theta)},\\
    \langle F,G\rangle_D(m,r)
      &=\sum_{n\in\mathbb Z}
        \overline{F(r-n\theta)}G(r-n\theta-m),
  \end{align*}
  as in \cite[pp.~426--427]{Rieffel1981}.  For \(\theta<0\), apply the
  positive case to \(-\theta\) and use
  \(A_\alpha\cong A_{-\alpha}\), obtained by replacing one canonical
  generator by its adjoint.
\end{proof}

\section{The Modular Action}

The basic equivalence of the preceding section extends to all fractional
linear transformations with integral coefficients.  We follow the argument
of \cite[Section~2]{Rieffel1981}; the elementary transformations below also
fix the sign conventions determined by our relation
\(VU=e^{2\pi i\theta}UV\).

\begin{definition}\label{def:ModularAction}
  The group \(\operatorname{GL}(2,\mathbb Z)\) acts on
  \(\mathbb R\setminus\mathbb Q\) by fractional linear transformations:
  \[
    \begin{pmatrix}a&b\\c&d\end{pmatrix}\cdot\theta
      =\frac{a\theta+b}{c\theta+d}.
  \]
  This is called the \textnormal{modular action}.
\end{definition}

The denominator cannot vanish when \(\theta\) is irrational.  Moreover,
the matrices \(g\) and \(-g\) determine the same transformation, so the
action factors through \(\operatorname{PGL}(2,\mathbb Z)\).

\begin{proposition}\label{prop:TorusElementaryIsomorphisms}
  For every \(\theta\in\mathbb R\) and \(n\in\mathbb Z\),
  \[
    A_\theta\cong A_{\theta+n}
    \qquad\text{and}\qquad
    A_\theta\cong A_{-\theta}.
  \]
\end{proposition}

\begin{proof}
  Since \(e^{2\pi i(\theta+n)}=e^{2\pi i\theta}\), the first isomorphism
  follows immediately from the universal property.  Let \(U_\alpha\) and
  \(V_\alpha\) denote the canonical generators of \(A_\alpha\).  In
  \(A_{-\theta}\) one has
  \[
    V_{-\theta}^{*}U_{-\theta}
      =e^{2\pi i\theta}U_{-\theta}V_{-\theta}^{*}.
  \]
  Hence the assignments
  \[
    U_\theta\longmapsto U_{-\theta},
    \qquad
    V_\theta\longmapsto V_{-\theta}^{*}
  \]
  induce a *-homomorphism \(A_\theta\to A_{-\theta}\).  Applying the same
  construction with \(-\theta\) gives its inverse.
\end{proof}

Put
\[
  T=\begin{pmatrix}1&1\\0&1\end{pmatrix},
  \qquad
  J=\begin{pmatrix}0&1\\1&0\end{pmatrix}.
\]
Thus \(T\cdot\theta=\theta+1\) and
\(J\cdot\theta=1/\theta\).

\begin{lemma}\label{lem:ModularGenerators}
  Every fractional linear transformation arising from
  \(\operatorname{GL}(2,\mathbb Z)\) is a composition of the
  transformations induced by \(T\), \(T^{-1}\), and \(J\).
\end{lemma}

\begin{proof}
  This is the standard Euclidean-algorithm description of the modular
  group.  Equivalently, the classes of \(T\) and \(J\) generate
  \(\operatorname{PGL}(2,\mathbb Z)\); see
  \cite[p.~421]{Rieffel1981}.  The possible factor \(-I\) is irrelevant
  because it acts trivially by fractional linear transformations.
\end{proof}

\begin{proposition}\label{prop:ModularGeneratorEquivalences}
  If \(\theta\) is irrational, then
  \[
    A_\theta\cong A_{T\cdot\theta}
    \qquad\text{and}\qquad
    A_\theta\sim_M A_{J\cdot\theta}.
  \]
\end{proposition}

\begin{proof}
  The first assertion is Proposition~\ref{prop:TorusElementaryIsomorphisms}.
  By Corollary~\ref{cor:RieffelBasicEquivalence},
  \[
    A_\theta\sim_M A_{-1/\theta}.
  \]
  Proposition~\ref{prop:TorusElementaryIsomorphisms} gives
  \(A_{-1/\theta}\cong A_{1/\theta}=A_{J\cdot\theta}\), which proves the
  second assertion.
\end{proof}

\begin{theorem}\label{thm:ModularOrbitMorita}
  Let \(\theta\in\mathbb R\setminus\mathbb Q\) and
  \(g\in\operatorname{GL}(2,\mathbb Z)\).  Then
  \[
    A_\theta\sim_M A_{g\cdot\theta}.
  \]
  Consequently, irrational parameters in the same modular orbit determine
  strongly Morita equivalent noncommutative tori.
\end{theorem}

\begin{proof}
  By Lemma~\ref{lem:ModularGenerators}, the transformation induced by \(g\)
  is a finite composition of those induced by \(T\), \(T^{-1}\), and \(J\).
  Proposition~\ref{prop:ModularGeneratorEquivalences} supplies a strong
  Morita equivalence at each step.  By
  Proposition~\ref{prop:MoritaTransitivity}, the finitely many equivalences
  can be composed, which proves the assertion.
\end{proof}

\begin{corollary}\label{cor:FractionalLinearMorita}
  Let \(\theta\) be irrational and let \(a,b,c,d\in\mathbb Z\) satisfy
  \(ad-bc=\pm1\).  Then
  \[
    A_\theta\sim_M
    A_{(a\theta+b)/(c\theta+d)}.
  \]
\end{corollary}

The converse is substantially deeper: the ordered \(K_0\)-group and its
trace recover the modular orbit.  We turn to this in the next section.

\section{Classification up to Morita Equivalence}

We now prove the converse of Theorem~\ref{thm:ModularOrbitMorita}.  The
argument is due to Rieffel; see the proof of
\cite[Theorem~4, pp.~421--425]{Rieffel1981}.  The operator-algebraic input has
already been established above, and the remaining step is the following
elementary arithmetic observation.

\begin{lemma}\label{lem:TraceRangeModularOrbit}
  Let \(\theta,\eta\in\mathbb R\setminus\mathbb Q\). If there is a positive
  real number \(\lambda\) such that
  \[
    \mathbb Z+\mathbb Z\theta
      =\lambda(\mathbb Z+\mathbb Z\eta),
  \]
  then there are integers \(a,b,c,d\) satisfying \(ad-bc=\pm1\) such that
  \[
    \eta=\frac{a\theta+b}{c\theta+d}.
  \]
\end{lemma}

\begin{proof}
  Since \(1,\eta\) form a \(\mathbb Z\)-basis of
  \(\mathbb Z+\mathbb Z\eta\), the elements \(\lambda,\lambda\eta\) form a
  \(\mathbb Z\)-basis of \(\mathbb Z+\mathbb Z\theta\). Hence there are
  integers \(a,b,c,d\) such that
  \[
    \lambda=c\theta+d,
    \qquad
    \lambda\eta=a\theta+b.
  \]
  The change-of-basis matrix
  \[
    \begin{pmatrix}d&b\\c&a\end{pmatrix}
  \]
  is invertible over \(\mathbb Z\), and therefore its determinant
  \(ad-bc\) is \(\pm1\). Since \(\lambda>0\), division of the second
  equality by the first gives the asserted formula.
\end{proof} 

\begin{theorem}[Rieffel]\label{thm:MoritaClassificationIrrationalTori}
  Let \(\theta,\eta\in\mathbb R\setminus\mathbb Q\). The following conditions
  are equivalent:
  \begin{enumerate}
    \item \(A_\theta\) and \(A_\eta\) are strongly Morita equivalent;
    \item \(\theta\) and \(\eta\) lie in the same orbit of the modular action;
    \item there is a matrix
      \[
        g=\begin{pmatrix}a&b\\c&d\end{pmatrix}
        \in\operatorname{GL}(2,\mathbb Z)
      \]
      such that
      \[
        \eta=g\cdot\theta
        =\frac{a\theta+b}{c\theta+d}.
      \]
  \end{enumerate}
\end{theorem}

\begin{proof}
  The equivalence of (2) and (3) is the definition of the modular action, and
  Theorem~\ref{thm:ModularOrbitMorita} proves that (3) implies (1).

  Suppose conversely that \(A_\theta\sim_M A_\eta\), and let \(X\) be an
  \(A_\theta\)-\(A_\eta\) imprimitivity bimodule. Let \(\tau_\theta\) be the
  canonical trace on \(A_\theta\), and let \((\tau_\theta)_X\) be the finite
  trace on \(A_\eta\) induced by \(X\). By
  Corollary~\ref{cor:uniqueTraceIrrationalRotation}, there is a number
  \(\lambda>0\) such that
  \[
    (\tau_\theta)_X=\lambda\tau_\eta,
  \]
  where \(\tau_\eta\) is the canonical normalized trace on \(A_\eta\).
  Corollary~\ref{cor:TraceRangeMorita} and
  Proposition~\ref{prop:traceRangeRotationAlgebra} now give
  \begin{align*}
    \mathbb Z+\mathbb Z\theta
      &=(\tau_\theta)_*(K_0(A_\theta))\\
      &=((\tau_\theta)_X)_*(K_0(A_\eta))\\
      &=\lambda(\tau_\eta)_*(K_0(A_\eta))\\
      &=\lambda(\mathbb Z+\mathbb Z\eta).
  \end{align*}
  Lemma~\ref{lem:TraceRangeModularOrbit} therefore yields a matrix
  \(g\in\operatorname{GL}(2,\mathbb Z)\) such that
  \(\eta=g\cdot\theta\). This proves (3).
\end{proof}

\begin{remark}
  A Morita equivalence does not in general carry the normalized trace to the
  normalized trace.  The scalar \(\lambda\) above records precisely this
  failure of normalization. With the notation in the arithmetic lemma,
  \[
    \lambda=c\theta+d,
    \qquad
    \lambda\eta=a\theta+b. 
  \]
  Thus the denominator of the fractional linear transformation is the trace
  scaling factor associated with the equivalence bimodule.
\end{remark}


\part{The Arithmetic of the Noncommutative Torus}
\begin{comment}
  Parte I: fondamenti sulle \(C^*\)-algebre, prodotti incrociati, \(K\)-teoria, successione di Pimsner–Voiculescu, toro non commutativo e Morita equivalenza. È la parte che stai completando ora.

  Parte II: dal toro non commutativo all’aritmetica di \(\theta\): frazioni continue, approssimanti razionali, parole di Christoffel/Sturmiane e il collegamento con il risultato della tua tesi magistrale. Era l’idea che ti era piaciuta particolarmente e che avevi collegato al tuo “Putnam”.

  Parti III e IV: dovevano costruire il ponte tecnico fra l’aritmetica della Parte II e la frontiera finale: approssimazioni matriciali, modelli finito-dimensionali, algebre AF/AT, convergenza e geometria dei tori fuzzy. Non avevamo ancora fissato definitivamente la divisione fra le due parti.

  Parte V: fuzzy tori e approssimazione diofantea, con l’obiettivo di arrivare a problemi di frontiera e possibilmente a risultati originali.
\end{comment}

\chapter{Continued Fractions and Diophantine Approximation}

Continued fractions provide the arithmetic coordinates that will connect the
ordered group \(K_0(A_\theta)\) with finite and infinite words. In this
chapter we establish the classical results needed later: the matrix calculus
of convergents, their interpretation as neighbouring primitive lattice
points, and their extremal approximation properties. Our conventions follow
\cite[Chapters~I--II]{Khinchin1964}; see also
\cite[Chapters~I--II]{RockettSzusz1992}.

\section{The Euclidean Algorithm and Continued Fractions}

\begin{definition}
  A \textnormal{finite simple continued fraction} is an expression
  \[
    [a_0;a_1,\ldots,a_n]
    =a_0+\cfrac{1}{a_1+\cfrac{1}{\ddots+\cfrac{1}{a_n}}},
  \]
  where \(a_0\in\mathbb Z\) and \(a_1,\ldots,a_n\) are positive integers.
  An \textnormal{infinite simple continued fraction} is denoted by
  \[
    [a_0;a_1,a_2,\ldots].
  \]
\end{definition}

Let \(x\in\mathbb R\). Put \(x_0=x\), and, as long as \(x_n\) is not an
integer, define
\[
  a_n=\lfloor x_n\rfloor,
  \qquad
  x_{n+1}=\frac{1}{x_n-a_n}.
\]
The numbers \(x_n\) are called the \textnormal{complete quotients} of \(x\),
and the integers \(a_n\) are its \textnormal{partial quotients}. Notice that
\(a_0\in\mathbb Z\), whereas \(a_n\geq1\) for \(n\geq1\).

\begin{proposition}\label{prop:CFEuclideanAlgorithm}
  The continued-fraction algorithm terminates if and only if \(x\) is
  rational. Every irrational number has a unique infinite simple continued
  fraction. Every rational number has precisely two finite expansions if
  expansions ending in \(1\) are allowed; they are related by
  \[
    [a_0;\ldots,a_{n-1},a_n]
    =[a_0;\ldots,a_{n-1},a_n-1,1],
    \qquad a_n\geq2.
  \]
\end{proposition}

\begin{proof}
  If the algorithm terminates at \(x_n=a_n\), repeated substitution in
  \(x_k=a_k+1/x_{k+1}\) expresses \(x\) as a rational number.

  Conversely, write \(x=r_0/r_1\), where \(r_0\in\mathbb Z\),
  \(r_1\in\mathbb N\), and \(\gcd(r_0,r_1)=1\). Division with remainder gives
  \[
    r_0=a_0r_1+r_2,\qquad 0\leq r_2<r_1.
  \]
  If \(r_2\neq0\), then \(x_1=r_1/r_2\), and the subsequent steps are the
  Euclidean algorithm. Its positive remainders strictly decrease, so it
  terminates.

  At every nonterminal step, \(a_n=\lfloor x_n\rfloor\) and
  \(x_{n+1}=(x_n-a_n)^{-1}\) are forced. This proves uniqueness in the
  irrational case. In the rational case the Euclidean algorithm gives a
  unique expansion whose last partial quotient is at least \(2\). Finally,
  \(a_n=(a_n-1)+1/1\) gives the displayed alternative expansion, and these
  are the only two by uniqueness of the preceding steps.
\end{proof}

\begin{remark}
  Unless explicitly stated otherwise, a finite continued fraction will be
  written in its \textnormal{canonical form}, in which the last partial
  quotient is at least \(2\). This convention will be important when words
  are associated with rational slopes.
\end{remark}

\section{Convergents}

For a sequence \(a_0,a_1,\ldots\), define integers \(p_n,q_n\) by
\begin{align*}
  p_{-2}&=0, & p_{-1}&=1,
  & p_n&=a_np_{n-1}+p_{n-2},\\
  q_{-2}&=1, & q_{-1}&=0,
  & q_n&=a_nq_{n-1}+q_{n-2}.
\end{align*}

\begin{definition}
  The rational number
  \[
    \frac{p_n}{q_n}=[a_0;a_1,\ldots,a_n]
  \]
  is the \(n\)-th \textnormal{convergent} of the continued fraction.
\end{definition}

\begin{proposition}[Matrix formula]\label{prop:CFMatrixFormula}
  For every \(n\geq0\),
  \[
    \begin{pmatrix}
      p_n&p_{n-1}\\
      q_n&q_{n-1}
    \end{pmatrix}
    =
    \begin{pmatrix}a_0&1\\1&0\end{pmatrix}
    \begin{pmatrix}a_1&1\\1&0\end{pmatrix}
    \cdots
    \begin{pmatrix}a_n&1\\1&0\end{pmatrix}.
  \]
  Consequently,
  \begin{equation}\label{eq:CFUnimodularIdentity}
    p_nq_{n-1}-p_{n-1}q_n=(-1)^{n-1}.
  \end{equation}
  In particular, \(p_n\) and \(q_n\) are relatively prime.
\end{proposition}

\begin{proof}
  The assertion for \(n=0\) follows from
  \(p_0=a_0,q_0=1,p_{-1}=1,q_{-1}=0\). Multiplication on the right by
  \(
    \left(\begin{smallmatrix}a_n&1\\1&0\end{smallmatrix}\right)
  \)
  gives precisely the recurrence relations, so the formula follows by
  induction. Each factor has determinant \(-1\), and hence the determinant
  of the product is \((-1)^{n+1}=(-1)^{n-1}\). Any common divisor of \(p_n\)
  and \(q_n\) divides the determinant in
  \eqref{eq:CFUnimodularIdentity}, and is therefore equal to \(1\).
\end{proof}

\begin{proposition}\label{prop:CFCompleteQuotientFormula}
  Let \(x=[a_0;a_1,a_2,\ldots]\) be irrational, and let \(x_{n+1}\) be its
  \((n+1)\)-st complete quotient. Then
  \begin{equation}\label{eq:CFCompleteQuotientFormula}
    x=\frac{p_nx_{n+1}+p_{n-1}}
            {q_nx_{n+1}+q_{n-1}}.
  \end{equation}
  Consequently,
  \begin{equation}\label{eq:CFExactError}
    q_nx-p_n
    =\frac{(-1)^n}{q_nx_{n+1}+q_{n-1}}.
  \end{equation}
\end{proposition}

\begin{proof}
  Formula \eqref{eq:CFCompleteQuotientFormula} follows by replacing the last
  entry in the finite matrix calculation by \(x_{n+1}\), or directly by
  repeated substitution in \(x_k=a_k+1/x_{k+1}\). Subtracting \(p_n\) after
  multiplying by \(q_n\) gives
  \[
    q_nx-p_n
    =\frac{q_np_{n-1}-p_nq_{n-1}}
           {q_nx_{n+1}+q_{n-1}},
  \]
  and the numerator is \((-1)^n\) by
  \eqref{eq:CFUnimodularIdentity}.
\end{proof}

\begin{corollary}\label{cor:CFConvergenceAndBounds}
  The convergents \(p_n/q_n\) alternate around \(x\) and converge to \(x\).
  More precisely,
  \begin{equation}\label{eq:CFErrorBounds}
    \frac{1}{q_n(q_n+q_{n+1})}
    <
    \left|x-\frac{p_n}{q_n}\right|
    <
    \frac{1}{q_nq_{n+1}}.
  \end{equation}
\end{corollary}

\begin{proof}
  The sign in \eqref{eq:CFExactError} is \((-1)^n\), so successive
  convergents lie on opposite sides of \(x\). Since
  \[
    x_{n+1}=a_{n+1}+\frac1{x_{n+2}},
  \]
  we have
  \[
    q_{n+1}
    <q_nx_{n+1}+q_{n-1}
    <q_{n+1}+q_n.
  \]
  Divide \eqref{eq:CFExactError} by \(q_n\) and use these inequalities to
  obtain \eqref{eq:CFErrorBounds}. Since the denominators are strictly
  increasing from \(q_1\) onwards, the right-hand side tends to zero.
\end{proof}

\begin{corollary}\label{cor:CFConsecutiveDifference}
  For \(n\geq1\), consecutive convergents satisfy
  \[
    \left|\frac{p_n}{q_n}-\frac{p_{n-1}}{q_{n-1}}\right|
    =\frac{1}{q_nq_{n-1}}.
  \]
\end{corollary}

\begin{proof}
  This is \eqref{eq:CFUnimodularIdentity} divided by
  \(q_nq_{n-1}\).
\end{proof}

\section{Intermediate Fractions}

\begin{definition}
  For \(n\geq0\) and \(0\leq k\leq a_{n+1}\), put
  \[
    p_{n,k}=kp_n+p_{n-1},
    \qquad
    q_{n,k}=kq_n+q_{n-1}.
  \]
  The fractions
  \[
    \frac{p_{n,k}}{q_{n,k}}
    =[a_0;\ldots,a_n,k],
    \qquad 1\leq k<a_{n+1},
  \]
  are called the \textnormal{intermediate fractions} or
  \textnormal{semiconvergents}. The endpoints are
  \[
    \frac{p_{n,0}}{q_{n,0}}=\frac{p_{n-1}}{q_{n-1}},
    \qquad
    \frac{p_{n,a_{n+1}}}{q_{n,a_{n+1}}}
    =\frac{p_{n+1}}{q_{n+1}}.
  \]
\end{definition}

\begin{proposition}\label{prop:CFSemiconvergents}
  Every \(p_{n,k}/q_{n,k}\) is reduced. For fixed \(n\), the intermediate
  fractions are strictly monotone in \(k\), and all the fractions with
  \(0\leq k\leq a_{n+1}\) lie on the same side of \(x\) as
  \(p_{n-1}/q_{n-1}\). The convergent \(p_n/q_n\) lies on the opposite
  side. Moreover,
  \begin{equation}\label{eq:CFSemiconvergentError}
    q_{n,k}x-p_{n,k}
    =(k-x_{n+1})(q_nx-p_n).
  \end{equation}
\end{proposition}

\begin{proof}
  The determinant
  \[
    p_{n,k}q_n-p_nq_{n,k}
    =p_{n-1}q_n-p_nq_{n-1}
    =(-1)^n
  \]
  shows that \(p_{n,k}\) and \(q_{n,k}\) are relatively prime. If
  \(0\leq k<\ell\), then
  \[
    p_{n,\ell}q_{n,k}-p_{n,k}q_{n,\ell}
    =(\ell-k)(p_nq_{n-1}-p_{n-1}q_n),
  \]
  whose sign is independent of \(k\) and \(\ell\). This proves
  monotonicity.

  From \eqref{eq:CFCompleteQuotientFormula} we obtain
  \[
    q_{n-1}x-p_{n-1}
    =-x_{n+1}(q_nx-p_n).
  \]
  Adding \(k(q_nx-p_n)\) gives
  \eqref{eq:CFSemiconvergentError}. Since
  \(x_{n+1}>a_{n+1}\), its sign is constant for
  \(0\leq k\leq a_{n+1}\), and is opposite to that of
  \(q_nx-p_n\).
\end{proof}

\section{Farey Neighbours and Lattice Geometry}

We call a vector \((p,q)\in\mathbb Z^2\) \textnormal{primitive} if
\(\gcd(p,q)=1\). Thus reduced fractions with positive denominator are the
slopes of primitive lattice vectors in the upper half-plane.

\begin{lemma}[Unimodular interval lemma]\label{lem:CFUnimodularInterval}
  Let \(a/b<c/d\) be reduced fractions with \(b,d>0\) and
  \(bc-ad=1\). If \(r/s\) is a reduced fraction strictly between them, then
  there are positive integers \(m,n\) such that
  \[
    (r,s)=m(a,b)+n(c,d).
  \]
  In particular, \(s=mb+nd\geq b+d\).
\end{lemma}

\begin{proof}
  The vectors \((a,b)\) and \((c,d)\) form a \(\mathbb Z\)-basis of
  \(\mathbb Z^2\), because their determinant is \(-1\). Hence
  \((r,s)=m(a,b)+n(c,d)\) for unique integers \(m,n\). Taking determinants
  first with \((a,b)\) and then with \((c,d)\) gives
  \[
    n=br-as>0,
    \qquad
    m=cs-dr>0,
  \]
  where positivity follows from \(a/b<r/s<c/d\). The assertion about the
  denominator follows immediately.
\end{proof}

\begin{corollary}\label{cor:CFNoSmallDenominatorBetweenConvergents}
  No reduced fraction strictly between two consecutive convergents
  \(p_{n-1}/q_{n-1}\) and \(p_n/q_n\) has denominator smaller than
  \(q_{n-1}+q_n\).
\end{corollary}

\begin{proof}
  Apply Lemma~\ref{lem:CFUnimodularInterval} and
  \eqref{eq:CFUnimodularIdentity}, interchanging the endpoints if necessary.
\end{proof}

\begin{remark}
  Proposition~\ref{prop:CFMatrixFormula} says more than the fact that two
  convergents are close: the column vectors
  \[
    \binom{p_n}{q_n},
    \qquad
    \binom{p_{n-1}}{q_{n-1}}
  \]
  form a basis of \(\mathbb Z^2\). This is the first appearance of the
  unimodular matrices that will later act on the parameter of the
  noncommutative torus.
\end{remark}

\section{Best Approximations of the Second Kind}

\begin{definition}
  Let \(x\) be irrational and let \(p/q\) be reduced with \(q>0\). We say
  that \(p/q\) is a \textnormal{best approximation of the second kind} to
  \(x\) if
  \[
    |qx-p|<|sx-r|
  \]
  for every reduced \(r/s\neq p/q\) with \(0<s\leq q\).
\end{definition}

\begin{lemma}[Law of best approximation]\label{lem:CFLawBestApproximation}
  Let \(p_n/q_n\) be the convergents of an irrational number \(x\). If
  \(r,s\in\mathbb Z\), \(0<s<q_{n+1}\), and
  \((r,s)\neq(p_n,q_n)\), then
  \begin{equation}\label{eq:CFLawBestApproximation}
    |sx-r|>|q_nx-p_n|.
  \end{equation}
\end{lemma}

\begin{proof}
  Put \(\varepsilon_n=q_nx-p_n\). The vectors
  \((p_n,q_n)\) and \((p_{n-1},q_{n-1})\) form a
  \(\mathbb Z\)-basis, so write
  \[
    (r,s)=u(p_n,q_n)+v(p_{n-1},q_{n-1})
  \]
  with \(u,v\in\mathbb Z\). By
  \eqref{eq:CFSemiconvergentError}, used with \(k=0\),
  \(\varepsilon_{n-1}=-x_{n+1}\varepsilon_n\), and hence
  \begin{equation}\label{eq:CFBasisError}
    sx-r=(u-vx_{n+1})\varepsilon_n.
  \end{equation}
  Suppose, towards a contradiction, that
  \(|sx-r|\leq|\varepsilon_n|\). Then
  \begin{equation}\label{eq:CFCoefficientStrip}
    |u-vx_{n+1}|\leq1.
  \end{equation}

  If \(v=0\), then \(u\in\{-1,0,1\}\). The conditions \(s>0\) and
  \((r,s)\neq(p_n,q_n)\) leave no possibility. Thus \(v\neq0\).
  Inequality \eqref{eq:CFCoefficientStrip}, together with
  \(x_{n+1}>1\), implies that \(u\) and \(v\) have the same sign. They
  cannot both be negative because then
  \(s=uq_n+vq_{n-1}<0\). Hence \(u,v>0\).

  Since \(x_{n+1}>a_{n+1}\), the inequality
  \(u\geq vx_{n+1}-1\) implies \(u\geq va_{n+1}\). Therefore
  \[
    s=uq_n+vq_{n-1}
      \geq v(a_{n+1}q_n+q_{n-1})
      =vq_{n+1}
      \geq q_{n+1},
  \]
  contrary to the hypothesis.
\end{proof}

\begin{theorem}[Lagrange]\label{thm:CFBestSecondKind}
  Every best approximation of the second kind to an irrational number is a
  convergent. Conversely, every convergent \(p_n/q_n\) with \(n\geq1\) is a
  best approximation of the second kind. The initial convergent
  \(p_0/q_0=a_0\) is best if and only if \(a_1\geq2\).
\end{theorem}

\begin{proof}
  Lemma~\ref{lem:CFLawBestApproximation}, with \(n\) replaced by \(n-1\),
  shows, for \(n\geq1\), that \(p_n/q_n\) improves strictly upon every
  rational with smaller denominator; applying the lemma with \(n\) also
  treats the other fractions having denominator \(q_n\). Hence every such
  convergent is a best approximation of the second kind.

  Conversely, let \(p/q\) be a best approximation of the second kind, and
  choose \(n\) so that
  \[
    q_n\leq q<q_{n+1}.
  \]
  If \(p/q\neq p_n/q_n\), Lemma~\ref{lem:CFLawBestApproximation} gives
  \[
    |qx-p|>|q_nx-p_n|,
  \]
  contradicting the defining extremal property of \(p/q\). Therefore
  \(p/q=p_n/q_n\).

  It remains to inspect \(p_0/q_0=a_0\). If \(a_1\geq2\), then \(q_1>1\),
  and Lemma~\ref{lem:CFLawBestApproximation} with \(n=0\) shows that \(a_0\)
  is best among the fractions of denominator \(1\). If \(a_1=1\), then
  \(x-a_0>1/2\), so the convergent
  \(p_1/q_1=(a_0+1)/1\) is closer to \(x\) than \(a_0/1\).
\end{proof}

\begin{theorem}[Legendre]\label{thm:CFLegendre}
  Let \(x\) be irrational and let \(p/q\) be reduced with \(q>0\). If
  \[
    \left|x-\frac pq\right|<\frac{1}{2q^2},
  \]
  then \(p/q\) is a convergent of \(x\).
\end{theorem}

\begin{proof}
  Choose consecutive convergents \(p_{n-1}/q_{n-1}\) and \(p_n/q_n\) such
  that \(q_{n-1}\leq q<q_n\). If \(p/q\) is not a convergent, the law of
  best approximation gives
  \[
    |q_{n-1}x-p_{n-1}|<|qx-p|<\frac{1}{2q}.
  \]
  Since the integer \(pq_{n-1}-qp_{n-1}\) is nonzero,
  \begin{align*}
    1
    &\leq |pq_{n-1}-qp_{n-1}|\\
    &=|q_{n-1}(p-qx)+q(q_{n-1}x-p_{n-1})|\\
    &\leq q_{n-1}|qx-p|+q|q_{n-1}x-p_{n-1}|\\
    &<\frac{q_{n-1}}{2q}+\frac12
    \leq1,
  \end{align*}
  a contradiction.
\end{proof}

\section{Best Approximations of the First Kind}

\begin{definition}
  Let \(x\) be irrational and let \(p/q\) be reduced with \(q>0\). We say
  that \(p/q\) is a \textnormal{best approximation of the first kind} to
  \(x\) if
  \[
    \left|x-\frac pq\right|
    <
    \left|x-\frac rs\right|
  \]
  for every reduced \(r/s\neq p/q\) with \(0<s\leq q\).
\end{definition}

The two notions differ because
\[
  |qx-p|=q\left|x-\frac pq\right|.
\]
Thus every best approximation of the second kind is a best approximation of
the first kind, but the converse fails in general.

\begin{proposition}\label{prop:CFBestFirstKindCandidates}
  Every best approximation of the first kind is either a convergent or an
  intermediate fraction.
\end{proposition}

\begin{proof}
  Let \(p/q\) be a best approximation of the first kind which is not a
  convergent. The case \(q=1\) is immediate: the closest integer to \(x\) is
  either \(a_0=p_0/q_0\) or, when \(a_1=1\), \(a_0+1=p_1/q_1\). Thus
  \(q>1\), and we may choose \(n\) such that
  \(q_n<q<q_{n+1}\). The consecutive
  convergents \(p_{n-1}/q_{n-1}\) and \(p_n/q_n\) lie on opposite sides of
  \(x\). On the side opposite to \(p_n/q_n\), consider
  \[
    \frac{p_{n,k}}{q_{n,k}}
    =\frac{kp_n+p_{n-1}}{kq_n+q_{n-1}},
    \qquad 0\leq k\leq a_{n+1}.
  \]
  They occur monotonically and their denominators increase with \(k\).
  Choose the unique \(k\) such that
  \[
    q_{n,k}\leq q<q_{n,k+1}.
  \]

  The fractions \(p_n/q_n\) and \(p_{n,k}/q_{n,k}\) have determinant of
  absolute value \(1\) and bracket \(x\). By
  Lemma~\ref{lem:CFUnimodularInterval}, every reduced fraction strictly
  between them has denominator at least
  \(q_n+q_{n,k}=q_{n,k+1}>q\). Consequently, among fractions with
  denominator at most \(q\), the closest fraction to \(x\) on either side
  must be one of these two boundary fractions. Since \(p/q\) is not a
  convergent, it must equal \(p_{n,k}/q_{n,k}\).
\end{proof}

\begin{remark}\label{rem:CFFirstKindDeferred}
  Not every intermediate fraction is a best approximation of the first
  kind. Determining exactly which of them are best requires comparing the
  two tails of the continued fraction in a borderline case. This is the
  subtle criterion that will be developed in the next chapter together with
  Henry John Stephen Smith's arrangement of fractions. The proposition
  above reduces that problem to the correct finite list of candidates.
\end{remark}

\begin{remark}\label{rem:CFPreviewKTheory}
  If \(0<x<1\), the errors above already have the form needed for the
  noncommutative torus. After identifying \(K_0(A_x)\) with the lattice
  generated by \([1]\) and the class of a Powers--Rieffel projection
  \([p_x]\), the primitive class
  \[
    \eta_{p,q}=q[p_x]-p[1]
  \]
  satisfies
  \[
    \tau_*(\eta_{p,q})=qx-p.
  \]
  Thus second-kind approximation is minimization of the absolute trace on
  primitive lattice classes, while first-kind approximation is minimization
  of the normalized quantity \(|\tau_*(\eta_{p,q})|/q\). We shall return to
  this observation only after the combinatorial meaning of \(p/q\) has been
  developed.
\end{remark}

\chapter[Smith and Continued Fractions]
        {Henry John Stephen Smith and the Arrangement of Fractions}

In 1876 Henry John Stephen Smith published a short but remarkably rich note
on continued fractions. His starting point was elementary: divide the unit
interval simultaneously into \(p\) and into \(q\) equal parts, and record how
the two families of division points interlace. Smith proved that this finite
arrangement determines, and is determined by, the continued fraction of
\(p/q\). He then used the same geometry to give a precise classification of
the best approximations of the first kind. We give complete proofs in modern
notation, following \cite{Smith1876}.

\section{Smith's \texorpdfstring{\(P\)- and \(Q\)-Points}{P- and Q-Points}}

Let \(p>q>0\) be relatively prime, and put
\[
  P=\frac1p,\qquad Q=\frac1q.
\]
On the interval \((0,1]\), consider the points
\[
  P,2P,\ldots,pP
  \qquad\text{and}\qquad
  Q,2Q,\ldots,qQ.
\]
The only coincidence is \(pP=qQ=1\). Reading the points from left to right
and writing \(P\) or \(Q\) according to their family produces a word with
\(p\) letters \(P\) and \(q\) letters \(Q\). The order of the two letters at
the common endpoint will be fixed below by the canonical continued-fraction
expansion. It is the only part of the word not determined by strict order.

\begin{definition}
  Let
  \[
    \frac pq=[a_0;a_1,\ldots,a_m]
  \]
  be the canonical continued fraction, so that \(a_m\geq2\). Define words
  \(S_{-1},S_0,\ldots,S_{m+1}\) over the alphabet \(\{P,Q\}\) by
  \[
    S_{-1}=Q,\qquad S_0=P,
  \]
  and
  \begin{equation}\label{eq:SmithWordRecursion}
    S_{j+1}=S_j^{a_j}S_{j-1},
    \qquad 0\leq j\leq m.
  \end{equation}
  The word \(S_{m+1}\) is called the \textnormal{Smith word} of \(p/q\) and
  is denoted by \(\mathcal S(p,q)\).
\end{definition}

Thus
\[
  S_1=P^{a_0}Q,\qquad
  S_2=S_1^{a_1}P,\qquad
  S_3=S_2^{a_2}S_1,
\]
which is Smith's original sequence of blocks.

\begin{proposition}\label{prop:SmithLetterCounts}
  Let \(p_j/q_j=[a_0;\ldots,a_j]\). Then \(S_{j+1}\) contains exactly
  \(p_j\) letters \(P\) and \(q_j\) letters \(Q\). In particular,
  \(\mathcal S(p,q)\) contains \(p\) letters \(P\) and \(q\) letters \(Q\).
\end{proposition}

\begin{proof}
  Let \(\lvert W\rvert_P\) and \(\lvert W\rvert_Q\) denote the numbers of
  occurrences of the two letters in \(W\). We have
  \[
    \binom{\lvert S_0\rvert_P}{\lvert S_0\rvert_Q}
      =\binom10,\qquad
    \binom{\lvert S_{-1}\rvert_P}{\lvert S_{-1}\rvert_Q}
      =\binom01.
  \]
  Recursion \eqref{eq:SmithWordRecursion} gives
  \[
    \binom{\lvert S_{j+1}\rvert_P}{\lvert S_{j+1}\rvert_Q}
    =
    a_j\binom{\lvert S_j\rvert_P}{\lvert S_j\rvert_Q}
      +\binom{\lvert S_{j-1}\rvert_P}{\lvert S_{j-1}\rvert_Q}.
  \]
  These are exactly the recurrence relations for \(p_j,q_j\) from
  Chapter~15.
\end{proof}

\begin{remark}
  The last two letters of \(S_{m+1}\) are \(PQ\) or \(QP\), according to
  parity. This fixes the order assigned to the coincident endpoint
  \(pP=qQ\). Replacing the canonical finite expansion by its alternative
  expansion reverses precisely these final two letters and leaves the strict
  arrangement inside \((0,1)\) unchanged.
\end{remark}

\section{The Direct Theorem}

The following elementary observation isolates the Euclidean algorithm hidden
in the arrangement; compare \cite[Arts.~1--10]{Smith1876}.

\begin{lemma}[Euclidean desubstitution]\label{lem:SmithDesubstitution}
  Write
  \[
    p=aq+r,\qquad 0<r<q.
  \]
  In the strict arrangement of the \(P\)- and \(Q\)-points, each \(Q\)-point
  is preceded, since the preceding \(Q\)-point, by either \(a\) or \(a+1\)
  new \(P\)-points. If a block \(P^aQ\) is replaced by a new letter \(A\)
  and each remaining \(P\) by a new letter \(B\), the resulting word is the
  arrangement word associated with the pair \((q,r)\), up to the convention
  at the common endpoint.
\end{lemma}

\begin{proof}
  The number of \(P\)-points not exceeding \(jQ=j/q\) is
  \[
    \left\lfloor\frac{pj}{q}\right\rfloor
    =aj+\left\lfloor\frac{rj}{q}\right\rfloor.
  \]
  Consequently, between \((j-1)Q\) and \(jQ\) there are
  \[
    a+
    \left\lfloor\frac{rj}{q}\right\rfloor
    -\left\lfloor\frac{r(j-1)}{q}\right\rfloor
  \]
  new \(P\)-points. The difference of the two floors is either \(0\) or
  \(1\), proving the first assertion.

  After the \(a\) compulsory \(P\)-points have been grouped with the
  following \(Q\), an additional \(P\) occurs exactly when an integral
  multiple of \(r/q\) crosses an integer. These crossings are ordered in
  exactly the same way as the points obtained by dividing the unit interval
  into \(q\) and \(r\) equal parts. Thus the contracted word is the
  arrangement for \((q,r)\). The two points at \(1\) again require only the
  stated endpoint convention.
\end{proof}

\begin{theorem}[Smith's direct theorem]\label{thm:SmithDirect}
  Let \(p>q>0\) be relatively prime. The Smith word
  \(\mathcal S(p,q)\) gives the arrangement, in increasing order, of
  \[
    \frac1p,\frac2p,\ldots,\frac pp
    \qquad\text{and}\qquad
    \frac1q,\frac2q,\ldots,\frac qq,
  \]
  with the endpoint convention described above.
\end{theorem}

\begin{proof}
  We argue by induction on \(q\). The assertion is immediate for \(q=1\):
  if \(p=a\), the word is \(P^aQ\), with the endpoint order prescribed by
  the one-term expansion.

  Suppose \(q>1\), and write \(p=a_0q+r\), \(0<r<q\). By
  Lemma~\ref{lem:SmithDesubstitution}, contracting \(P^{a_0}Q\) to a letter
  and each remaining \(P\) to the other letter produces the arrangement for
  \((q,r)\). The Euclidean algorithm gives
  \[
    \frac qr=[a_1;\ldots,a_m].
  \]
  By the inductive hypothesis its arrangement is generated by the recurrence
  with partial quotients \(a_1,\ldots,a_m\). Undoing the contraction replaces
  its first letter by \(P^{a_0}Q=S_1\) and its second letter by \(P=S_0\).
  The successive substitutions are therefore
  \[
    S_{j+1}=S_j^{a_j}S_{j-1},
  \]
  and the expanded word is \(S_{m+1}=\mathcal S(p,q)\).
\end{proof}

\begin{corollary}\label{cor:SmithOrderFromCF}
  The continued fraction of \(p/q\) determines the complete arrangement of
  the \(P\)- and \(Q\)-points.
\end{corollary}

\section{The Inverse Theorem}

Smith's construction is reversible. In modern terminology it is a
desubstitution algorithm.

\begin{theorem}[Smith's inverse theorem]\label{thm:SmithInverse}
  The strict arrangement of the \(P\)- and \(Q\)-points determines the
  rational number \(p/q\), its canonical continued fraction, and all the
  words \(S_j\) in \eqref{eq:SmithWordRecursion}.
\end{theorem}

\begin{proof}
  Count the initial \(P\)-points before the first \(Q\)-point. Their number
  is
  \[
    a_0=\left\lfloor\frac pq\right\rfloor.
  \]
  Hence the first block is \(S_1=P^{a_0}Q\). Parse the word from left to
  right by contracting every occurrence of \(P^{a_0}Q\) to a new letter
  \(A\), and every \(P\) left outside these blocks to a new letter \(B\).
  Lemma~\ref{lem:SmithDesubstitution} shows that the contracted word is the
  arrangement for the remainder pair \((q,r)\), where
  \(p=a_0q+r\).

  The same procedure now determines
  \(a_1=\lfloor q/r\rfloor\), then the next remainder, and so on. At each
  step the sum of the two parameters strictly decreases, so the process
  terminates. It is precisely the Euclidean algorithm, and therefore
  recovers the canonical expansion
  \[
    \frac pq=[a_0;a_1,\ldots,a_m].
  \]
  Reversing the contractions recovers successively
  \(S_1,S_2,\ldots,S_{m+1}\). The only ambiguity is the order of the two
  coincident endpoint letters; requiring \(a_m\geq2\) selects the canonical
  one.
\end{proof}

\begin{remark}
  Smith learned, only after completing his investigation, that Christoffel
  had already described the same arrangements by means of least positive
  residues. Smith explicitly recorded this in a footnote to his paper. The
  precise identification with Christoffel words will be the subject of the
  next chapter.
\end{remark}

\section{Successive Minima of the First Kind}

We return to an irrational number
\[
  x=[a_0;a_1,a_2,\ldots]
\]
and its convergents \(p_n/q_n\). For \(n\geq0\) and
\(1\leq k<a_{n+1}\), write
\[
  r_{n,k}
  =\frac{p_{n,k}}{q_{n,k}}
  =\frac{kp_n+p_{n-1}}{kq_n+q_{n-1}}.
\]
These are the intermediate fractions introduced in Chapter~15. Put
\[
  \alpha_{n+1}=x_{n+1}
    =[a_{n+1};a_{n+2},a_{n+3},\ldots],
  \qquad
  \beta_n=\frac{q_{n-1}}{q_n}
    =[0;a_n,a_{n-1},\ldots,a_1],
\]
where \(\beta_0=0\).

\begin{proposition}\label{prop:SmithIntermediateErrors}
  The two fractions \(r_{n,k}\) and \(p_n/q_n\) lie on opposite sides of
  \(x\), and
  \begin{align}
    \left|x-r_{n,k}\right|
      &=
      \frac{\alpha_{n+1}-k}
      {(q_n\alpha_{n+1}+q_{n-1})(kq_n+q_{n-1})},
      \label{eq:SmithIntermediateError}\\
    \left|x-\frac{p_n}{q_n}\right|
      &=
      \frac{1}
      {q_n(q_n\alpha_{n+1}+q_{n-1})}.
      \label{eq:SmithConvergentError}
  \end{align}
  Consequently,
  \begin{equation}\label{eq:SmithSuperiorCondition}
    \left|x-r_{n,k}\right|
      <
    \left|x-\frac{p_n}{q_n}\right|
    \quad\Longleftrightarrow\quad
    \alpha_{n+1}<2k+\beta_n.
  \end{equation}
\end{proposition}

\begin{proof}
  Formula \eqref{eq:CFCompleteQuotientFormula} gives
  \[
    x=\frac{p_n\alpha_{n+1}+p_{n-1}}
            {q_n\alpha_{n+1}+q_{n-1}}.
  \]
  Subtracting \(r_{n,k}\), and using
  \(\lvert p_nq_{n-1}-p_{n-1}q_n\rvert=1\), yields
  \eqref{eq:SmithIntermediateError}. Formula
  \eqref{eq:SmithConvergentError} is
  \eqref{eq:CFExactError} divided by \(q_n\). The signs are opposite because
  \(\alpha_{n+1}>a_{n+1}>k\).

  Cancelling the common positive factor
  \(q_n\alpha_{n+1}+q_{n-1}\), the desired comparison becomes
  \[
    \frac{\alpha_{n+1}-k}{kq_n+q_{n-1}}
      <\frac1{q_n}.
  \]
  Multiplication by the positive denominators gives
  \(q_n\alpha_{n+1}<2kq_n+q_{n-1}\), which is precisely
  \eqref{eq:SmithSuperiorCondition}.
\end{proof}

\begin{definition}
  The intermediate fraction \(r_{n,k}\) is called
  \textnormal{superior} if
  \begin{equation}\label{eq:SmithSuperiorDefinition}
    [a_{n+1};a_{n+2},\ldots]
    <
    2k+[0;a_n,a_{n-1},\ldots,a_1].
  \end{equation}
  Otherwise it is called \textnormal{inferior}.
\end{definition}

This is Smith's condition (A), written in the notation of Chapter~15.
It appears in \cite[Arts.~15--16]{Smith1876}.

\begin{theorem}[Smith's theorem on successive minima]
  \label{thm:SmithSuccessiveMinima}
  The best approximations of the first kind to an irrational number \(x\)
  are precisely its convergents \(p_n/q_n\), \(n\geq1\), and its superior
  intermediate fractions, ordered by increasing denominator. The initial
  convergent \(p_0/q_0=a_0\) is included exactly when \(a_1\geq2\).
\end{theorem}

\begin{proof}
  By Proposition~\ref{prop:CFBestFirstKindCandidates}, every best
  approximation of the first kind is either a convergent or an intermediate
  fraction. Every convergent \(p_n/q_n\), apart from the harmless initial
  convention discussed in Theorem~\ref{thm:CFBestSecondKind}, is a best
  approximation of the second kind and therefore of the first kind. The
  assertion concerning \(p_0/q_0\) is exactly the initial case of
  Theorem~\ref{thm:CFBestSecondKind}.

  Consider an intermediate fraction \(r_{n,k}\). Its denominator is larger
  than \(q_n\). If it is inferior, then
  \[
    \left|x-r_{n,k}\right|
      \geq
    \left|x-\frac{p_n}{q_n}\right|
  \]
  by Proposition~\ref{prop:SmithIntermediateErrors}; hence it cannot be
  best.

  Suppose that \(r_{n,k}\) is superior. Then it and \(p_n/q_n\) lie on
  opposite sides of \(x\), and \(r_{n,k}\) is the closer of the two.
  Moreover,
  \[
    \left|p_{n,k}q_n-p_nq_{n,k}\right|=1.
  \]
  If a reduced fraction \(c/d\) with
  \(0<d\leq q_{n,k}\) were closer to \(x\) than \(r_{n,k}\), it would lie
  strictly between \(r_{n,k}\) and \(p_n/q_n\). Indeed, on the side of
  \(r_{n,k}\) this is immediate, while on the opposite side the endpoint
  \(p_n/q_n\) is farther from \(x\) than \(r_{n,k}\). The unimodular interval
  lemma would then give
  \[
    d\geq q_{n,k}+q_n>q_{n,k},
  \]
  a contradiction. Thus \(r_{n,k}\) is a best approximation of the first
  kind.
\end{proof}

\begin{remark}
  Smith called the successive record values of
  \(\lvert y/z-x\rvert\) the successive minima of the expression
  \(y/z-x\). Theorem~\ref{thm:SmithSuccessiveMinima} is therefore his precise
  answer to the question that the theorem of Lagrange for the second kind
  leaves open.
\end{remark}

\section{The Borderline Intermediate Fraction}

Condition \eqref{eq:SmithSuperiorDefinition} has a particularly transparent
form away from its middle value.

\begin{proposition}\label{prop:SmithHalfCriterion}
  Let \(1\leq k<a_{n+1}\).
  \begin{enumerate}
    \item If \(2k>a_{n+1}\), then \(r_{n,k}\) is superior.
    \item If \(2k<a_{n+1}\), then \(r_{n,k}\) is inferior.
    \item If \(2k=a_{n+1}\), then \(r_{n,k}\) is superior if and only if
      \begin{equation}\label{eq:SmithBorderlineTails}
        [0;a_{n+2},a_{n+3},\ldots]
        <
        [0;a_n,a_{n-1},\ldots,a_1].
      \end{equation}
  \end{enumerate}
\end{proposition}

\begin{proof}
  Write
  \[
    \alpha_{n+1}=a_{n+1}+\gamma,
    \qquad
    \gamma=[0;a_{n+2},a_{n+3},\ldots]\in(0,1),
  \]
  and recall that \(\beta_n\in[0,1)\). The superior condition is
  \[
    a_{n+1}+\gamma<2k+\beta_n.
  \]
  If \(2k\geq a_{n+1}+1\), it holds because \(\gamma<1\) and
  \(\beta_n\geq0\). If \(2k\leq a_{n+1}-1\), it fails because
  \(\gamma>0\) and \(\beta_n<1\). When \(2k=a_{n+1}\), it reduces to
  \(\gamma<\beta_n\), which is
  \eqref{eq:SmithBorderlineTails}.
\end{proof}

\begin{corollary}\label{cor:SmithLexicographicCriterion}
  In the borderline case \(2k=a_{n+1}\), compare the two finite or infinite
  strings
  \[
    (a_{n+2},a_{n+3},\ldots)
    \quad\text{and}\quad
    (a_n,a_{n-1},\ldots,a_1).
  \]
  At their first unequal entry, the intermediate fraction is superior if
  the entry in the first string is larger at an odd comparison step, or
  smaller at an even comparison step. In particular, it is superior if
  \[
    a_n<a_{n+2};
  \]
  if \(a_n=a_{n+2}\), the comparison passes to \(a_{n-1}\) and
  \(a_{n+3}\), with the inequality reversed, and so on.
\end{corollary}

\begin{proof}
  This is the standard alternating comparison rule for simple continued
  fractions, applied to the two continued fractions in
  \eqref{eq:SmithBorderlineTails}. It also follows directly by repeatedly
  observing that the map \(t\mapsto1/(a+t)\) reverses inequalities.
\end{proof}

\begin{remark}
  Thus, when \(a_{n+1}=2h+1\), exactly the indices
  \(k=h+1,\ldots,2h\) are superior. When \(a_{n+1}=2h\), the indices
  \(k=h+1,\ldots,2h-1\) are superior, the indices
  \(k=1,\ldots,h-1\) are inferior, and the status of the middle fraction
  \(k=h\) is decided by
  \eqref{eq:SmithBorderlineTails}. This is the exceptional case emphasized
  by Smith.
\end{remark}

\section{Smith's Example: \texorpdfstring{\(39/17\)}{39/17}}

Smith begins his note with
\[
  \frac{39}{17}=[2;3,2,2].
\]
The recursive construction gives
\begin{align*}
  S_{-1}&=Q,\\
  S_0&=P,\\
  S_1&=P^2Q,\\
  S_2&=(P^2Q)^3P,\\
  S_3&=\bigl((P^2Q)^3P\bigr)^2(P^2Q),\\
  S_4&=
  \left(\bigl((P^2Q)^3P\bigr)^2(P^2Q)\right)^2
  (P^2Q)^3P.
\end{align*}
The successive contents are
\[
  \binom21,\qquad
  \binom73,\qquad
  \binom{16}{7},\qquad
  \binom{39}{17},
\]
which are precisely the numerator--denominator pairs of
\[
  \frac21,\qquad
  \frac73,\qquad
  \frac{16}{7},\qquad
  \frac{39}{17}.
\]
Thus \(S_4\) gives the complete arrangement of the \(39\) points \(i/39\)
and the \(17\) points \(j/17\).

Conversely, the arrangement begins with two \(P\)'s followed by \(Q\), so
the first quotient is \(2\). Contracting \(P^2Q\) produces three copies
before the first residual \(P\), giving the second quotient \(3\). Repeating
the contraction gives \(2\) and \(2\). The word therefore recovers
\([2;3,2,2]\) without prior knowledge of \(39\) and \(17\).

\section{The First-Kind Records for \texorpdfstring{\(\pi\)}{pi}}

The beginning of the continued fraction of \(\pi\) is
\[
  \pi=[3;7,15,1,292,1,1,\ldots],
\]
with convergents
\[
  \frac31,\quad
  \frac{22}{7},\quad
  \frac{333}{106},\quad
  \frac{355}{113},\quad
  \frac{103993}{33102},\ldots.
\]

For the partial quotient \(a_1=7\), the intermediate fractions are
\[
  \frac{3k+1}{k},
  \qquad 1\leq k\leq6.
\]
Since \(2k>7\) exactly for \(k=4,5,6\), the superior ones are
\[
  \frac{13}{4},\qquad
  \frac{16}{5},\qquad
  \frac{19}{6}.
\]

For \(a_2=15\), the intermediate fractions are
\[
  \frac{22k+3}{7k+1},
  \qquad 1\leq k\leq14.
\]
The superior ones have \(k=8,\ldots,14\), and begin
\[
  \frac{179}{57},\quad
  \frac{201}{64},\quad
  \frac{223}{71},\quad
  \frac{245}{78},\ldots,\frac{311}{99}.
\]
The next partial quotient is \(1\), so there are no intermediate fractions
between \(333/106\) and \(355/113\).

The large coefficient \(a_4=292\) produces
\[
  \frac{355k+333}{113k+106},
  \qquad 1\leq k\leq291.
\]
All indices \(k>146\) are automatically superior. At the borderline
\(k=146\), condition \eqref{eq:SmithBorderlineTails} reads
\[
  [0;1,1,\ldots]<[0;1,15,7],
\]
which holds. Hence the first superior fraction in this block is already
\[
  \frac{355\cdot146+333}{113\cdot146+106}
  =\frac{52163}{16604},
\]
followed by
\[
  \frac{52518}{16717},\qquad
  \frac{52873}{16830},\qquad\ldots.
\]
Smith's criterion therefore explains not only the exceptional quality of
\(355/113\), but the long family of first-kind records created by the next
partial quotient \(292\).

\begin{remark}
  The construction of the words \(S_j\), the ordering of the division
  points, and the classification of the superior intermediate fractions are
  three aspects of the same Euclidean recursion. In the next chapter the
  finite words will be recognized as Christoffel words; after that, their
  irrational limit will lead to Sturmian words and to the coding of the
  rotation defining \(A_\theta\).
\end{remark} 


\chapter{Christoffel Words} 

The preceding chapter associated a finite word with a pair of coprime
integers by arranging two systems of division points on the unit interval.
There is a second, apparently different, way of obtaining the same circular
word: approximate a line segment of rational slope by a lattice path.  The
result is a \emph{Christoffel word}.  Christoffel introduced these words in
his study of rational and irrational approximation
\cite{Christoffel1875}; our treatment follows the modern geometric account
in \cite[Part~I, Chapters~1, 3 and~7]{BerstelLauveReutenauerSaliola2009}.

The purpose of this chapter is not merely to rename Smith's construction.
The geometry of the lattice makes visible a canonical factorization, and
that factorization organizes all positive rational slopes in the
Stern--Brocot tree.  Thus the Euclidean algorithm, continued fractions,
lattice paths and finite words become four descriptions of the same
object.

\section{Words and Lattice Paths}

Let \(\mathcal A=\{X,Y\}\), ordered by \(X<Y\), and let
\(\mathcal A^*\) denote the free monoid of finite words on \(\mathcal A\).
The empty word is denoted by \(\varepsilon\).  If \(w\in\mathcal A^*\),
write \(|w|_X\) and \(|w|_Y\) for the number of occurrences of \(X\) and
\(Y\), respectively, and put
\[
  \operatorname{cont}(w)=(|w|_X,|w|_Y).
\]
The word \(w=w_1\cdots w_n\) determines a lattice path starting at the
origin by reading
\[
  X=(1,0),\qquad Y=(0,1).
\]
Its endpoint is therefore \(\operatorname{cont}(w)\).

Fix positive integers \(p,q\), and put \(N=p+q\).  For
\(0\leq k\leq N\), define
\[
  y_k=\left\lfloor\frac{kq}{N}\right\rfloor,
  \qquad x_k=k-y_k.
\]
Then \((x_0,y_0)=(0,0)\), \((x_N,y_N)=(p,q)\), and consecutive points
differ by either \((1,0)\) or \((0,1)\).  Moreover,
\begin{equation}\label{eq:ChristoffelHeight}
  qx_k-py_k=kq-N\left\lfloor\frac{kq}{N}\right\rfloor
  \in\{0,1,\ldots,N-1\}.
\end{equation}
Thus every vertex lies on or below the segment joining \((0,0)\) to
\((p,q)\), and at each diagonal \(x+y=k\) it is the highest lattice point
with this property.

\begin{definition}\label{definition:LowerChristoffel}
  Let \(p,q\) be relatively prime positive integers.  The
  \textnormal{lower Christoffel word} of slope \(q/p\), denoted by
  \(C_{p,q}\), is the word of the path through the points
  \((x_k,y_k)\).  Equivalently, its \((k+1)\)-st letter is
  \begin{equation}\label{eq:ChristoffelMechanicalDefinition}
    (C_{p,q})_{k+1}=
    \begin{cases}
      X,&\displaystyle
      \left\lfloor\frac{(k+1)q}{p+q}\right\rfloor
      =\left\lfloor\frac{kq}{p+q}\right\rfloor,\\[6pt]
      Y,&\displaystyle
      \left\lfloor\frac{(k+1)q}{p+q}\right\rfloor
      =\left\lfloor\frac{kq}{p+q}\right\rfloor+1,
    \end{cases}
    \qquad 0\leq k<p+q.
  \end{equation}
  The \textnormal{upper Christoffel word} \(C^{\mathrm{up}}_{p,q}\) is
  defined in the same way using
  \(y_k=\lceil kq/(p+q)\rceil\).
\end{definition}

Since \(0<q<N\), the two cases in
\eqref{eq:ChristoffelMechanicalDefinition} are exhaustive.  The lower
word contains exactly \(p\) letters \(X\) and \(q\) letters \(Y\), starts
with \(X\), and ends with \(Y\).  The upper word starts with \(Y\) and
ends with \(X\).

\begin{proposition}\label{proposition:ChristoffelReversal}
  For relatively prime positive integers \(p,q\),
  \[
    C^{\mathrm{up}}_{p,q}=C_{p,q}^{\mathrm{rev}},
  \]
  where \(w^{\mathrm{rev}}\) denotes the word obtained by reading \(w\)
  from right to left.
\end{proposition}
\begin{proof}
  Coprimality implies that \(kq/(p+q)\) is not an integer for
  \(0<k<p+q\).  Hence
  \[
    \left\lceil\frac{kq}{p+q}\right\rceil
    =q-\left\lfloor\frac{(p+q-k)q}{p+q}\right\rfloor.
  \]
  The step made by the upper path at time \(k\) is therefore the step made
  by the lower path at time \(p+q-k-1\).  This reverses the order of the
  letters.
\end{proof}

\section{The Arithmetic Coding}

The same word can be read without drawing the path.  Let
\[
  r_k\equiv kq\pmod{p+q},
  \qquad 0\leq r_k<p+q.
\]
By \eqref{eq:ChristoffelHeight}, \(r_k\) is exactly the vertical distance
of the \(k\)-th vertex from the diagonal, measured by the integer linear
form \((x,y)\mapsto qx-py\).

\begin{theorem}[geometric--arithmetic equivalence]
  \label{theorem:ChristoffelArithmeticCoding}
  The \((k+1)\)-st letter of \(C_{p,q}\) is \(Y\) precisely when the
  passage from \(r_k\) to \(r_{k+1}\) wraps around modulo \(p+q\).  More
  explicitly,
  \[
    (C_{p,q})_{k+1}=
    \begin{cases}
      X,&r_{k+1}>r_k,\\
      Y,&r_{k+1}<r_k.
    \end{cases}
  \]
\end{theorem}
\begin{proof}
  Write
  \[
    kq=(p+q)\left\lfloor\frac{kq}{p+q}\right\rfloor+r_k.
  \]
  Since \(0<q<p+q\), adding \(q\) either leaves the quotient unchanged,
  in which case \(r_{k+1}=r_k+q>r_k\), or increases it by one, in which
  case \(r_{k+1}=r_k+q-(p+q)<r_k\).  These are exactly the two cases in
  \eqref{eq:ChristoffelMechanicalDefinition}.
\end{proof}

The coprimality hypothesis has a transparent meaning.  If the same floor
construction is made for arbitrary positive \(p,q\), and
\(d=\gcd(p,q)\), then
\begin{equation}\label{eq:ChristoffelPower}
  C_{p,q}=C_{p/d,q/d}^{,d}.
\end{equation}
Indeed, the sequence of increments in
\eqref{eq:ChristoffelMechanicalDefinition} has period \((p+q)/d\).
Consequently a lower mechanical word is primitive if and only if its
content vector is primitive.  This is the word-theoretic counterpart of
the fact that the segment from \((0,0)\) to \((p,q)\) contains no interior
lattice point precisely when \(p\) and \(q\) are coprime.

\section{The Standard Factorization}

The most important finite structure theorem for Christoffel words is their
canonical factorization into two shorter Christoffel words.  We first
record the arithmetic form of the construction.

\begin{lemma}[Farey parents]\label{lemma:ChristoffelFareyParents}
  Let \(p,q\) be relatively prime positive integers.  There are unique
  nonnegative integer vectors
  \[
    v_1=(p_1,q_1),\qquad v_2=(p_2,q_2)
  \]
  such that
  \[
    v_1+v_2=(p,q),
    \qquad p_1q_2-q_1p_2=1.
  \]
  Both vectors are primitive, and
  \[
    \frac{q_1}{p_1}<\frac qp<\frac{q_2}{p_2},
  \]
  with the usual conventions at \(0/1\) and \(1/0\).
\end{lemma}
\begin{proof}
  The determinant condition is equivalent to
  \[
    p_1q-q_1p=1.
  \]
  B\'ezout's identity gives a solution.  Choosing the unique representative
  with \(0\leq p_1\leq p\) forces \(0\leq q_1\leq q\), and then
  \(v_2=(p-p_1,q-q_1)\).  Any two such solutions differ by an integral
  multiple of \((p,q)\), so the stated bounds give uniqueness.  A common
  divisor of the coordinates of either \(v_i\) would divide the
  determinant, and hence must be one.  Finally,
  \[
    p_1q-q_1p=pq_2-qp_2=1,
  \]
  which gives the two strict inequalities.
\end{proof}

\begin{theorem}[standard factorization]
  \label{theorem:ChristoffelStandardFactorization}
  Let \(p,q>0\) be relatively prime, and let \(v_1=(p_1,q_1)\) and
  \(v_2=(p_2,q_2)\) be the vectors in
  Lemma~\ref{lemma:ChristoffelFareyParents}.  Then
  \begin{equation}\label{eq:ChristoffelStandardFactorization}
    C_{p,q}=C_{p_1,q_1}C_{p_2,q_2},
  \end{equation}
  where \(C_{1,0}=X\) and \(C_{0,1}=Y\).  This is the unique
  factorization of \(C_{p,q}\) as a product of two Christoffel words.
\end{theorem}
\begin{proof}
  Put \(N_i=p_i+q_i\) and \(N=p+q=N_1+N_2\).  The determinant identity
  gives
  \[
    qN_1-q_1N=1,
    \qquad q_2N-qN_2=1.
  \]
  For \(0\leq k\leq N_1\),
  \[
    \frac{kq}{N}=\frac{kq_1}{N_1}+\frac{k}{NN_1}.
  \]
  Since \(q_1\) and \(N_1\) are coprime, the added term is too small to
  cross an integer; at \(k=N_1\) both floors equal \(q_1\).  Therefore
  the first \(N_1\) steps of \(C_{p,q}\) form \(C_{p_1,q_1}\).

  For \(0\leq j\leq N_2\), a second calculation gives
  \[
    \frac{(N_1+j)q}{N}
    =q_1+\frac{jq_2}{N_2}+\frac{N_2-j}{NN_2}.
  \]
  The same argument shows that the remaining \(N_2\) increments are those
  of \(C_{p_2,q_2}\).  This proves
  \eqref{eq:ChristoffelStandardFactorization}.

  Conversely, if \(C_{p,q}=C_{a,b}C_{c,d}\) is a factorization into two
  Christoffel words, the two corresponding segments meet at a lattice
  point closest to the segment from \((0,0)\) to \((p,q)\).  The triangle
  determined by the two content vectors has no lattice point other than
  its vertices; equivalently, its determinant is one.  Thus
  \(ad-bc=1\), and uniqueness follows from
  Lemma~\ref{lemma:ChristoffelFareyParents}.  This is the usual geometric
  proof of uniqueness; see
  \cite[Theorem~3.3]{BerstelLauveReutenauerSaliola2009}.
\end{proof}

\begin{remark}
  The joining point \((p_1,q_1)\) is the lattice point of the lower path
  closest to the open segment from \((0,0)\) to \((p,q)\).  Thus the word
  factorization is visible directly in the picture.  The matrix
  \[
    \begin{pmatrix}p_1&p_2\\q_1&q_2\end{pmatrix}
    \in\operatorname{SL}_2(\mathbb Z)
  \]
  records the same fact algebraically.
\end{remark}

\section{The Christoffel and Stern--Brocot Trees}

Two reduced fractions \(q_1/p_1<q_2/p_2\) are \emph{Farey neighbours} if
\[
  p_1q_2-q_1p_2=1.
\]
Their mediant is
\[
  \frac{q_1+q_2}{p_1+p_2}.
\]
Starting from the pair \(0/1<1/0\) and repeatedly inserting mediants gives
the Stern--Brocot tree.

There is an exactly parallel tree of words.  Start from the pair
\((X,Y)\), and replace a pair \((u,v)\) by either
\begin{equation}\label{eq:ChristoffelTreeChildren}
  (u,uv)\qquad\text{or}\qquad(uv,v).
\end{equation}
The word associated with a vertex \((u,v)\) is its product \(uv\).

\begin{theorem}[Christoffel tree]
  \label{theorem:ChristoffelTree}
  Every proper lower Christoffel word occurs exactly once as a product
  \(uv\) in the tree \eqref{eq:ChristoffelTreeChildren}.  Under the map
  \[
    w\longmapsto\frac{|w|_Y}{|w|_X},
  \]
  the Christoffel tree is isomorphic to the positive part of the
  Stern--Brocot tree.
\end{theorem}
\begin{proof}
  Suppose \(\operatorname{cont}(u)=(p_1,q_1)\) and
  \(\operatorname{cont}(v)=(p_2,q_2)\), with determinant one.  Then
  \(uv=C_{p_1+p_2,q_1+q_2}\) by
  Theorem~\ref{theorem:ChristoffelStandardFactorization}, and its slope is
  the mediant of the slopes of \(u\) and \(v\).  Both replacements in
  \eqref{eq:ChristoffelTreeChildren} preserve determinant one, so every
  word produced by the tree is Christoffel.

  Conversely, the standard factorization of a proper Christoffel word is
  unique.  Applying it repeatedly strictly decreases word length and
  eventually reaches \((X,Y)\).  It therefore gives a unique path back to
  the root.  This proves both existence and uniqueness, and at the same
  time recovers the corresponding Stern--Brocot path.
\end{proof}

The theorem is stronger than the assertion that every positive rational
has a continued fraction.  It says that the entire ancestry of a rational
slope is already stored in the factorization of its word.  A left or right
run in the tree records a partial quotient; Euclid's algorithm reads the
same runs in reverse.

\section{Continued Fractions and Standard Words}

There is one harmless but important inversion of convention.  The letter
\(X\) represents a horizontal step, so a word of content \((p,q)\) has
geometric slope \(q/p\); Smith, on the other hand, applies the Euclidean
algorithm to \(p/q\).  Passing to the reciprocal only exchanges the two
directions in the Stern--Brocot tree and loses no information.

Let
\[
  \frac pq=[a_0;a_1,\ldots,a_m],
  \qquad a_m\geq2,
\]
be the canonical continued fraction of a positive reduced rational.  As in
the preceding chapter, define
\begin{equation}\label{eq:ChristoffelStandardRecursion}
  S_{-1}=Y,\qquad S_0=X,\qquad
  S_{n+1}=S_n^{a_n}S_{n-1}.
\end{equation}
The content vectors satisfy the same recursion as the numerator and
denominator vectors of the convergents.  Hence
\[
  \operatorname{cont}(S_{m+1})=(p,q).
\]
The precise relation with \(C_{p,q}\) is circular rather than linear: the
two words have different preferred origins.

\begin{definition}
  Two words \(u,v\) are \textnormal{conjugate} if there are words \(r,s\)
  such that \(u=rs\) and \(v=sr\).  Equivalently, one is obtained from the
  other by a cyclic rotation.
\end{definition}

\begin{theorem}[continued-fraction construction]
  \label{theorem:CFChristoffelCorrespondence}
  The word \(S_{m+1}\) defined by
  \eqref{eq:ChristoffelStandardRecursion} is conjugate to the lower
  Christoffel word \(C_{p,q}\).  Consequently the canonical factorization
  of \(C_{p,q}\), the path to \(q/p\) in the Stern--Brocot tree, and the
  Euclidean algorithm for \((p,q)\) determine one another.
\end{theorem}
\begin{proof}
  The direct theorem for Smith's recursion, proved in the preceding
  chapter, says that \(S_{m+1}\) gives the circular order of the two systems
  of division points.  Thus \(S_{m+1}\) is conjugate to the ordered word
  \(W(p,q)\) defined below.  The independent floor-function calculation in
  Theorem~\ref{theorem:SmithChristoffelCorrespondence} shows that
  \(W(p,q)\) is the left cyclic shift of \(C_{p,q}\).  Hence
  \(S_{m+1}\) and \(C_{p,q}\) are conjugate.

  Finally, the content-vector recursion in
  \eqref{eq:ChristoffelStandardRecursion} is the continuant recursion, and
  consecutive vectors have determinant of absolute value one.  Therefore
  the Euclidean quotients are exactly the lengths of the alternating runs
  in the Stern--Brocot tree.  Together with the uniqueness of the standard
  factorization, this proves the last assertion.  See also
  \cite[Chapter~7]{BerstelLauveReutenauerSaliola2009}.
\end{proof}

\section{The Smith--Christoffel Correspondence}

We can now identify the construction of the preceding chapter exactly.
Let \(W(p,q)\) be the word obtained by arranging
\[
  \frac1p,\ldots,\frac{p-1}{p},
  \qquad
  \frac1q,\ldots,\frac{q-1}{q}
\]
in increasing order, labelling the first family by \(X\) and the second by
\(Y\), and finally appending \(YX\) at the common endpoint \(1\).

\begin{theorem}[Smith--Christoffel correspondence]
  \label{theorem:SmithChristoffelCorrespondence}
  If \(p,q\) are relatively prime positive integers, then
  \[
    W(p,q)=\operatorname{sh}(C_{p,q}),
  \]
  where \(\operatorname{sh}(w_1w_2\cdots w_N)=w_2\cdots w_Nw_1\) is the
  left cyclic shift.  In particular, Smith's word, the standard word
  obtained from the continued fraction of \(p/q\), and \(C_{p,q}\) all
  belong to the same conjugacy class.
\end{theorem}
\begin{proof}
  In \(C_{p,q}\), the letter in position \(k\), numbered from zero, is
  \(Y\) exactly when
  \[
    \left\lfloor\frac{(k+1)q}{p+q}\right\rfloor
    -\left\lfloor\frac{kq}{p+q}\right\rfloor=1.
  \]
  Thus the positions of its \(Y\)'s are
  \begin{equation}\label{eq:ChristoffelYPositions}
    \left\lceil\frac{j(p+q)}q\right\rceil-1,
    \qquad 1\leq j\leq q.
  \end{equation}

  Consider now the point \(j/q\) in Smith's ordered list.  Before it there
  are \(j-1\) points of the same family and
  \(\lceil jp/q\rceil-1\) points of the other family.  Its zero-based
  position is therefore
  \[
    j-1+\left\lceil\frac{jp}{q}\right\rceil-1
    =\left\lceil\frac{j(p+q)}q\right\rceil-2.
  \]
  The same formula gives \(p+q-2\) for the appended \(Y\) at the endpoint.
  Comparing with \eqref{eq:ChristoffelYPositions}, every \(Y\)-position
  has moved one place to the left.  Hence \(W(p,q)\) is precisely the left
  cyclic shift of \(C_{p,q}\).  The last assertion follows immediately
  from this identity and the direct theorem for Smith's recursion.
\end{proof}

\begin{remark}
  The shift in Theorem~\ref{theorem:SmithChristoffelCorrespondence} is not
  a notational nuisance.  A Christoffel path has a distinguished initial
  vertex, whereas Smith's division points naturally live on the circle
  \(\mathbb R/\mathbb Z\).  Forgetting the origin leaves a conjugacy class,
  and that class is the intrinsic object shared by the two constructions.
\end{remark}

\section{A Complete Example: \texorpdfstring{\(7/4\)}{7/4}}

Take \((p,q)=(7,4)\), so that \(p+q=11\).  The residues
\[
  0,4,8,1,5,9,2,6,10,3,7,0
\]
increase at an \(X\) and wrap around at a \(Y\).  Therefore
\[
  C_{7,4}=XXYXXYXXYXY.
\]
The B\'ezout equation
\[
  4p_1-7q_1=1
\]
has the bounded solution \((p_1,q_1)=(2,1)\), and hence
\((p_2,q_2)=(5,3)\).  The standard factorization is
\[
  C_{7,4}=C_{2,1}C_{5,3}
  =(XXY)(XXYXXYXY).
\]
On slopes, this is the mediant identity
\[
  \frac47=\frac{1+3}{2+5},
  \qquad 2\cdot3-1\cdot5=1.
\]

The canonical continued fraction is
\[
  \frac74=[1;1,3].
\]
Recursion \eqref{eq:ChristoffelStandardRecursion} gives
\[
  S_1=XY,\qquad S_2=XYX,\qquad
  S_3=(XYX)^3XY=XYXXYXXYXXY.
\]
Rotating \(S_3\) by two positions gives \(C_{7,4}\).  On the other hand,
ordering the division points gives
\[
  W(7,4)=XYXXYXXYXYX,
\]
which is the left shift of \(C_{7,4}\), exactly as predicted by
Theorem~\ref{theorem:SmithChristoffelCorrespondence}.  Thus the continued
fraction, the lattice path, the Farey parents and Smith's arrangement are
not parallel analogies: in this example, and in general, they are the same
finite arithmetic structure viewed from different origins.

\chapter{Sturmian Words and Irrational Rotations}

Christoffel words describe rational lines.  When the slope becomes
irrational, no finite lattice path can close up, and the corresponding word
is infinite.  The resulting \emph{Sturmian words} have the smallest factor
complexity compatible with aperiodicity.  They were isolated in the
foundational work of Morse and Hedlund \cite{MorseHedlund1940}; we also use
the modern account in \cite[Chapter~2]{Lothaire2002}.

This chapter has two goals.  First, we pass rigorously from finite
Christoffel words to Sturmian words by means of convergents.  Second, we
return to the irrational rotation defining \(A_\theta\).  The frequencies
of Sturmian factors will turn out to be lengths of intervals whose endpoints
belong to the orbit of \(\theta\); consequently they lie in
\(\mathbb Z+\mathbb Z\theta\), the same ordered group that occurred as the
range of the canonical trace on \(K_0(A_\theta)\).

\section{Factor Complexity and the Morse--Hedlund Theorem}

Let \(w=w_0w_1w_2\cdots\) be a one-sided infinite word on a finite alphabet
\(\mathcal A\).  A finite word \(u\) is a \textnormal{factor} of \(w\) if
\[
  u=w_kw_{k+1}\cdots w_{k+n-1}
\]
for some \(k,n\geq0\).  We denote the set of factors of length \(n\) by
\(\mathcal F_w(n)\), and define the \textnormal{factor complexity} of \(w\)
by
\[
  p_w(n)=\lvert\mathcal F_w(n)\rvert,
  \qquad p_w(0)=1.
\]
The word is \textnormal{ultimately periodic} if there are integers
\(r\geq0\) and \(d\geq1\) such that \(w_{k+d}=w_k\) for every \(k\geq r\).

\begin{theorem}[Morse--Hedlund]\label{theorem:MorseHedlundComplexity}
  An infinite word \(w\) is ultimately periodic if and only if
  \[
    p_w(n)\leq n
  \]
  for some \(n\geq1\).  In particular, every aperiodic infinite word
  satisfies
  \[
    p_w(n)\geq n+1
    \qquad(n\geq0).
  \]
\end{theorem}
\begin{proof}
  Every factor of length \(k\) has at least one right extension, so
  \(p_w(k+1)\geq p_w(k)\).  If equality holds, every factor of length
  \(k\) has exactly one right extension.  The sequence
  \[
    w_iw_{i+1}\cdots w_{i+k-1},
    \qquad i=0,1,2,\ldots,
  \]
  then evolves deterministically on the finite set \(\mathcal F_w(k)\).
  Two of its terms eventually coincide, after which all subsequent terms
  coincide with the same displacement.  Hence \(w\) is ultimately
  periodic.

  Now suppose that \(p_w(n)\leq n\).  Since
  \[
    p_w(n)-p_w(0)
    =\sum_{k=0}^{n-1}\bigl(p_w(k+1)-p_w(k)\bigr)\leq n-1,
  \]
  at least one of the \(n\) nonnegative summands is zero.  The preceding
  paragraph again gives ultimate periodicity.

  Conversely, suppose that \(w\) has preperiod \(r\) and period \(d\).
  A factor beginning at a position at least \(r\) depends only on that
  position modulo \(d\), while there are at most \(r\) factors beginning
  before \(r\).  Thus \(p_w(n)\leq r+d\) for every \(n\), and choosing
  \(n\geq r+d\) proves the converse.
\end{proof}

\begin{definition}
  A binary infinite word \(w\) is \textnormal{Sturmian} if it is aperiodic
  and
  \[
    p_w(n)=n+1
    \qquad(n\geq0).
  \]
\end{definition}

Thus Sturmian words attain the lower bound in
Theorem~\ref{theorem:MorseHedlundComplexity} at every length.

\begin{definition}
  A factor \(u\) of \(w\) is \textnormal{right special} if both \(uX\) and
  \(uY\) are factors of \(w\).  Left special factors are defined
  analogously.
\end{definition}

\begin{proposition}\label{proposition:SturmianUniqueSpecial}
  A binary aperiodic word is Sturmian if and only if it has exactly one
  right-special factor of every length.
\end{proposition}
\begin{proof}
  For \(u\in\mathcal F_w(n)\), let \(d^+(u)\) be its number of right
  extensions.  Every factor of length \(n+1\) has a unique prefix of length
  \(n\), and hence
  \[
    p_w(n+1)-p_w(n)
    =\sum_{u\in\mathcal F_w(n)}\bigl(d^+(u)-1\bigr).
  \]
  Over a binary alphabet, each summand is either zero or one.  The right
  side is therefore the number of right-special factors of length \(n\).
  The assertion follows by summing from \(p_w(0)=1\).
\end{proof}

\section{Mechanical Words}

Let \(0<\alpha<1\) and \(\rho\in\mathbb R\).  The \textnormal{lower} and
\textnormal{upper mechanical words} of slope \(\alpha\) and intercept
\(\rho\) are, respectively,
\begin{align}
  s_{\alpha,\rho}(n)
  &=\lfloor(n+1)\alpha+\rho\rfloor
    -\lfloor n\alpha+\rho\rfloor,\label{eq:LowerMechanicalWord}\\
  s'_{\alpha,\rho}(n)
  &=\lceil(n+1)\alpha+\rho\rceil
    -\lceil n\alpha+\rho\rceil,\label{eq:UpperMechanicalWord}
\end{align}
where \(n\geq0\), and the values \(0,1\) are identified with \(X,Y\).
The \textnormal{characteristic word} of slope \(\alpha\) is
\[
  c_\alpha=s_{\alpha,\alpha}.
\]

Let \(R_\alpha\colon\mathbb T\to\mathbb T\) be rotation by \(\alpha\).
Equation \eqref{eq:LowerMechanicalWord} can be rewritten as
\begin{equation}\label{eq:RotationCoding}
  s_{\alpha,\rho}(n)=Y
  \quad\Longleftrightarrow\quad
  R_\alpha^n(\rho)\in[1-\alpha,1).
\end{equation}
Thus a mechanical word is the itinerary of a point under a two-interval
coding of a rotation.

\begin{lemma}\label{lemma:IrrationalRotationDense}
  If \(\alpha\) is irrational, every orbit of \(R_\alpha\) is dense in
  \(\mathbb T\).
\end{lemma}
\begin{proof}
  The closure of \(\{n\alpha+\mathbb Z:n\in\mathbb Z\}\) is a closed
  subgroup of \(\mathbb T\).  Every proper closed subgroup of the circle is
  finite, whereas an irrational \(\alpha\) has infinite order.  The closure
  is therefore all of \(\mathbb T\).  Translating it by any starting point
  gives the result.
\end{proof}

\begin{theorem}\label{theorem:IrrationalMechanicalSturmian}
  If \(\alpha\) is irrational, every lower or upper mechanical word of
  slope \(\alpha\) is Sturmian.
\end{theorem}
\begin{proof}
  Telescoping \eqref{eq:LowerMechanicalWord} gives
  \[
    \sum_{k=0}^{N-1}s_{\alpha,\rho}(k)
    =\lfloor N\alpha+\rho\rfloor-\lfloor\rho\rfloor.
  \]
  Hence the frequency of \(Y\), if the symbols are viewed as zero and one,
  tends to \(\alpha\).  An ultimately periodic binary word has rational
  letter frequency, so \(s_{\alpha,\rho}\) is aperiodic.

  Fix \(n\geq1\).  The block of length \(n\) beginning at an orbit point
  \(x\) changes only when \(x\) crosses one of
  \[
    0,-\alpha,-2\alpha,\ldots,-n\alpha
    \quad\text{in }\mathbb T.
  \]
  Irrationality makes these \(n+1\) points distinct.  They divide the circle
  into \(n+1\) intervals, and the coding is constant on the interior of
  each interval.  Crossing an endpoint changes one of the two adjacent
  symbols, so the resulting blocks are distinct.  By
  Lemma~\ref{lemma:IrrationalRotationDense}, the orbit of \(\rho\) meets
  every interval.  Consequently \(p(n)=n+1\).  The argument for the upper
  word differs only in the choice made at the endpoints.
\end{proof}

\section{Balance and the Representation Theorem}

\begin{definition}
  A binary infinite word \(w\) is \textnormal{balanced} if, for every two
  factors \(u,v\) of equal length,
  \[
    \bigl||u|_Y-|v|_Y\bigr|\leq1.
  \]
\end{definition}

Mechanical words are balanced for a very simple reason.  A factor of length
\(n\) beginning at position \(k\) contains
\[
  \lfloor(k+n)\alpha+\rho\rfloor
  -\lfloor k\alpha+\rho\rfloor
\]
letters \(Y\), and this number is either \(\lfloor n\alpha\rfloor\) or
\(\lceil n\alpha\rceil\).

\begin{lemma}\label{lemma:BalanceSpecialFactors}
  A recurrent binary word is balanced if and only if it has at most one
  right-special factor of each length.
\end{lemma}
\begin{proof}
  Suppose first that two distinct right-special factors \(u,v\) have the
  same length.  Let \(z\) be their longest common suffix.  The letters
  immediately preceding \(z\) in \(u\) and \(v\) are different; after
  exchanging \(u,v\), they are \(X,Y\).  Since both factors are right
  special, the words \(XzX\) and \(YzY\) occur.  They have equal length and
  their numbers of \(Y\)'s differ by two, contradicting balance.

  Conversely, if balance fails, choose a shortest pair of equal-length
  factors whose numbers of \(Y\)'s differ by at least two.  Cancelling
  equal initial and terminal letters, and using minimality after each
  cancellation, produces factors \(XzX\) and \(YzY\).  Recurrence allows
  both occurrences to be extended until their first distinct right
  extensions appear.  Immediately before that first branching, the two
  corresponding suffixes are distinct right-special factors of the same
  length.  This is the standard minimal-counterexample argument; see
  \cite[Section~2.1]{Lothaire2002}.
\end{proof}

\begin{theorem}[Morse--Hedlund representation]
  \label{theorem:SturmianCharacterizations}
  For a binary infinite word \(w\), the following are equivalent:
  \begin{enumerate}
    \item \(w\) is Sturmian;
    \item \(w\) is aperiodic and balanced;
    \item \(w\) is a lower or upper mechanical word of irrational slope.
  \end{enumerate}
\end{theorem}
\begin{proof}
  First observe that every Sturmian word is recurrent.  Indeed, if a factor
  \(u\) of length \(n\) occurred only finitely many times, a tail beginning
  after its last occurrence would have at most
  \(p_w(n)-1=n\) factors of length \(n\).  By
  Theorem~\ref{theorem:MorseHedlundComplexity}, that tail, and hence \(w\),
  would be ultimately periodic.  This is impossible.

  Proposition~\ref{proposition:SturmianUniqueSpecial} and
  Lemma~\ref{lemma:BalanceSpecialFactors} now show that every Sturmian word
  is balanced.  Conversely, in a balanced aperiodic word there is at most
  one right-special factor of each length, by the first paragraph of the
  proof of Lemma~\ref{lemma:BalanceSpecialFactors}.  There is at least one,
  since otherwise \(p_w(n+1)=p_w(n)\) and the word would be ultimately
  periodic.  Proposition~\ref{proposition:SturmianUniqueSpecial} therefore
  shows that the word is Sturmian.  This proves the equivalence of (1) and
  (2).  The implication (3) to (1) is
  Theorem~\ref{theorem:IrrationalMechanicalSturmian}.

  It remains to represent a balanced word mechanically.  Regard \(X,Y\) as
  \(0,1\), and let \(S_n=w_0+\cdots+w_{n-1}\), with \(S_0=0\).  If \(m_n\)
  and \(M_n\) are the minimum and maximum numbers of \(Y\)'s in a factor of
  length \(n\), balance gives \(M_n-m_n\leq1\).  Splitting factors shows
  that
  \[
    m_{n+k}\geq m_n+m_k,
    \qquad M_{n+k}\leq M_n+M_k.
  \]
  The superadditive and subadditive lemmas therefore give a common number
  \[
    \alpha=\sup_n\frac{m_n}{n}
           =\inf_n\frac{M_n}{n}.
  \]
  Since \(m_n\leq n\alpha\leq M_n\) and \(M_n-m_n\leq1\), the number of
  letters \(Y\) in any factor of length \(n\) differs from \(n\alpha\) by
  at most one.

  Put \(D_n=S_n-n\alpha\).  For \(m<n\), the difference
  \(D_n-D_m\) is the discrepancy of a factor of length \(n-m\); the
  preceding estimate implies that any two \(D_n\)'s differ by at most
  one.  Hence
  \[
    \sup_nD_n\leq1+\inf_nD_n.
  \]
  Choose \(\rho\) between these two quantities.  The inequalities
  \[
    S_n\leq n\alpha+\rho<S_n+1
  \]
  hold simultaneously, with one consistent choice of open or closed
  endpoint in the limiting case.  Thus
  \[
    S_n=\lfloor n\alpha+\rho\rfloor-\lfloor\rho\rfloor
  \]
  for all \(n\), or the analogous formula with ceilings.  Taking first
  differences represents \(w\) by
  \eqref{eq:LowerMechanicalWord} or \eqref{eq:UpperMechanicalWord}.
  Finally, a mechanical word of rational slope is ultimately periodic.
  Since \(w\) is aperiodic, \(\alpha\) must be irrational.
\end{proof}

\section{Christoffel Approximants}

The normalization of the previous chapter and the usual normalization of a
mechanical word differ slightly.  A Christoffel word \(C_{p,q}\) has length
\(p+q\) and contains \(q\) letters \(Y\), so its density of \(Y\)'s is
\(q/(p+q)\), whereas its geometric slope is \(q/p\).

Let \(0<a<b\) be coprime, and form the rational mechanical block
\[
  M_{a/b}(k)
  =\left\lfloor\frac{(k+1)a}{b}\right\rfloor
   -\left\lfloor\frac{ka}{b}\right\rfloor,
  \qquad 0\leq k<b.
\]

\begin{proposition}\label{proposition:RationalMechanicalChristoffel}
  With \(X=0\) and \(Y=1\),
  \[
    M_{a/b}=C_{b-a,a}.
  \]
  Moreover, the lower mechanical word of rational slope \(a/b\) and
  intercept zero is the periodic word \(M_{a/b}^{\,\infty}\).
\end{proposition}
\begin{proof}
  In Definition~\ref{definition:LowerChristoffel}, substitute
  \(p=b-a\) and \(q=a\).  Then \(p+q=b\), and
  \eqref{eq:ChristoffelMechanicalDefinition} becomes exactly the displayed
  formula for \(M_{a/b}\).  Since adding \(b\) to \(k\) adds the integer
  \(a\) to both floor arguments, the sequence of differences has period
  \(b\).
\end{proof}

\begin{theorem}[Christoffel approximation]
  \label{theorem:ChristoffelApproximationSturmian}
  Let \(\theta\in(0,1)\) be irrational, and let \(p_n/q_n\) be its
  convergents.  Then the periodic words
  \[
    \bigl(C_{q_n-p_n,p_n}\bigr)^{\infty}
  \]
  converge, in the topology of agreement on finite prefixes, to the lower
  mechanical word \(s_{\theta,0}\).
\end{theorem}
\begin{proof}
  Fix \(L\geq1\).  None of the numbers \(k\theta\),
  \(1\leq k\leq L\), is an integer.  Since \(p_n/q_n\to\theta\), for all
  sufficiently large \(n\),
  \[
    \left\lfloor\frac{kp_n}{q_n}\right\rfloor
    =\lfloor k\theta\rfloor
    \qquad(1\leq k\leq L).
  \]
  Taking consecutive differences and applying
  Proposition~\ref{proposition:RationalMechanicalChristoffel} shows that
  the first \(L\) letters agree.  Since \(L\) was arbitrary, the claimed
  convergence follows.
\end{proof}

There is also a nested version of this limit, dictated directly by the
continued fraction.  Write
\[
  \theta=[0;a_1,a_2,\ldots]
\]
and define the \textnormal{standard words}
\begin{equation}\label{eq:SturmianStandardWords}
  T_0=X,\qquad T_1=X^{a_1-1}Y,\qquad
  T_{n+1}=T_n^{a_{n+1}}T_{n-1}\quad(n\geq1).
\end{equation}

\begin{theorem}[standard sequence]
  \label{theorem:SturmianStandardSequence}
  Let \(p_n/q_n=[0;a_1,\ldots,a_n]\).  Then
  \[
    \operatorname{cont}(T_n)=(q_n-p_n,p_n).
  \]
  Each \(T_n\) is conjugate to the Christoffel word
  \(C_{q_n-p_n,p_n}\), and the words \(T_n\), \(n\geq1\), form an
  increasing sequence of prefixes whose limit is the characteristic word
  \(c_\theta\).
\end{theorem}
\begin{proof}
  The initial content vectors are
  \[
    \operatorname{cont}(T_0)=(1,0),
    \qquad \operatorname{cont}(T_1)=(a_1-1,1),
  \]
  and \eqref{eq:SturmianStandardWords} gives the same recursion as
  \((q_n-p_n,p_n)\).  This proves the content formula.  Consecutive content
  vectors have determinant of absolute value one, by the unimodular
  identity for convergents.  The continued-fraction construction of
  Theorem~\ref{theorem:CFChristoffelCorrespondence}, with the reciprocal
  convention explained there, therefore shows that \(T_n\) is a conjugate
  of the corresponding Christoffel word.

  For \(n\geq1\), the word \(T_{n+1}\) begins with \(T_n\), so the limit
  \(T=\lim T_n\) exists.  The alternating inequalities
  \[
    \frac{p_{2k}}{q_{2k}}<\theta<
    \frac{p_{2k+1}}{q_{2k+1}}
  \]
  and the floor formula \eqref{eq:LowerMechanicalWord} show inductively
  that, for \(n\geq1\), the prefix of \(c_\theta\) of length \(q_n\) is
  \(T_n\).  Hence
  \(T=c_\theta\).  This is the standard-sequence construction; see
  \cite[Section~2.2]{Lothaire2002}.
\end{proof}

\section{The Sturmian Dynamical System}

For the dynamical formulation it is convenient to use two-sided words.  Put
\[
  \widehat{s}_{\theta,\rho}(n)
  =\lfloor(n+1)\theta+\rho\rfloor
   -\lfloor n\theta+\rho\rfloor,
  \qquad n\in\mathbb Z,
\]
and let \(\sigma\) be the left shift on \(\{X,Y\}^{\mathbb Z}\).  Directly
from the definition,
\begin{equation}\label{eq:RotationShiftIntertwining}
  \sigma\widehat{s}_{\theta,\rho}
  =\widehat{s}_{\theta,\rho+\theta}.
\end{equation}
Let \(\Omega_\theta\) be the closure of these codings in the product
topology.  It is called the \textnormal{Sturmian subshift} of slope
\(\theta\).

\begin{theorem}\label{theorem:SturmianDynamicalSystem}
  Let \(\theta\in(0,1)\) be irrational.  Then
  \((\Omega_\theta,\sigma)\) is minimal and uniquely ergodic.  Its language
  is the set of factors of any mechanical word of slope \(\theta\), and its
  factor complexity is \(n+1\).
\end{theorem}
\begin{proof}
  Equation \eqref{eq:RotationShiftIntertwining} intertwines rotation with
  the shift.  A factor \(u\) determines a nonempty half-open interval
  \(I_u\subset\mathbb T\) of starting points that produce \(u\).  By
  Lemma~\ref{lemma:IrrationalRotationDense}, every orbit meets \(I_u\).
  More strongly, the sets \(R_\theta^{-j}(\operatorname{int}I_u)\),
  \(j\geq0\), cover the circle; compactness gives a finite subcover.
  Therefore occurrences of \(u\) have bounded gaps in every coding.  Every
  orbit in \(\Omega_\theta\) is consequently dense, which proves
  minimality and also shows that all codings have the same language.

  To prove unique ergodicity, first observe that, for every nonzero
  \(m\in\mathbb Z\),
  \[
    \frac1N\sum_{k=0}^{N-1}e^{2\pi im(\rho+k\theta)}
    =\frac{e^{2\pi im\rho}}{N}
      \frac{1-e^{2\pi imN\theta}}{1-e^{2\pi im\theta}}
    \longrightarrow0
  \]
  uniformly in \(\rho\).  By approximation with trigonometric polynomials,
  the orbit averages of every continuous function converge uniformly to
  its Lebesgue integral.  Approximating the indicator of \(I_u\) from
  above and below shows that every coding has a uniform frequency for
  \(u\), equal to \(m(I_u)\).  Uniform factor frequencies characterize a
  unique invariant probability measure on a subshift.  Hence
  \(\Omega_\theta\) is uniquely ergodic.  The complexity assertion is
  Theorem~\ref{theorem:IrrationalMechanicalSturmian}.
\end{proof}

\section{Factor Frequencies and the Group
  \texorpdfstring{\(\mathbb Z+\mathbb Z\theta\)}{Z+Z theta}}

For a block occupying coordinates \(0,ldots,n-1\), the discontinuity
points of the coding are
\[
  0,-\theta,-2\theta,\ldots,-n\theta
  \quad\text{in }\mathbb T.
\]
They divide the circle into \(n+1\) half-open intervals, one for each factor
of length \(n\).  The frequency of the factor is the length of the
corresponding interval.

\begin{theorem}[arithmetic of Sturmian frequencies]
  \label{theorem:SturmianFrequencyGroup}
  Let \(\theta\in(0,1)\) be irrational.
  \begin{enumerate}
    \item Every factor frequency in \(\Omega_\theta\) belongs to
      \[
        (\mathbb Z+\mathbb Z\theta)\cap(0,1).
      \]
    \item Let \(\mathcal B_\theta\) be the Boolean algebra generated by all
      two-sided cylinder intervals of the rotation coding.  Then
      \[
        \{m(E):E\in\mathcal B_\theta\}
        = (\mathbb Z+\mathbb Z\theta)\cap[0,1].
      \]
  \end{enumerate}
\end{theorem}
\begin{proof}
  The endpoints of every cylinder interval are orbit points
  \(j\theta+\mathbb Z\).  The length of an interval between two such points
  is the representative in \([0,1]\) of an integer multiple of \(\theta\),
  and therefore lies in \(\mathbb Z+\mathbb Z\theta\).  By
  Theorem~\ref{theorem:SturmianDynamicalSystem}, this length is the
  frequency of the corresponding factor.  This proves (1), and also one
  inclusion in (2), since finite unions and complements preserve the group.

  Conversely, let \(\beta=m+n\theta\in[0,1]\).  Apart from the cases
  \(\beta=0,1\), it is the fractional part of \(n\theta\).  The half-open
  arc from \(0\) to \(n\theta+\mathbb Z\), with the appropriate orientation,
  is a finite union of atoms of a two-sided cylinder partition.  It belongs
  to \(\mathcal B_\theta\) and has measure \(\beta\).  This proves the
  reverse inclusion.
\end{proof}

The canonical trace on
\[
  A_\theta=C(\mathbb T)\rtimes_{\alpha_\theta}\mathbb Z
\] 
restricts on \(C(\mathbb T)\) to integration against Lebesgue measure.
Cylinder indicators are not continuous functions, but they belong to the
measurable completion of this dynamical system, and their measures are
computed by the same integral.  Combining
Theorem~\ref{theorem:SturmianFrequencyGroup} with
Theorem~\ref{thm:orderedK0RotationAlgebra}, we obtain the parallel
\[
  \begin{array}{c@{\qquad\longleftrightarrow\qquad}c}
    \text{measures of symbolic orbit intervals}
      &\text{traces of ordered \(K_0\)-classes}\\[2mm]
    (\mathbb Z+\mathbb Z\theta)\cap[0,1]
      &(\mathbb Z+\mathbb Z\theta)\cap[0,1].
  \end{array}
\]
The equality does not identify a discontinuous cylinder indicator with a
projection in \(C(\mathbb T)\).  Rather, it shows that the symbolic and
operator-algebraic constructions are governed by the same arithmetic
dimension group.

\section{The Fibonacci Word}

Let
\[
  \theta=\frac{\sqrt5-1}{2}=[0;1,1,1,\ldots].
\]
The standard recursion becomes
\[
  T_0=X,\qquad T_1=Y,\qquad T_{n+1}=T_nT_{n-1},
\]
and begins
\[
  X,\quad Y,\quad YX,\quad YXY,\quad YXYYX,
  \quad YXYYXYXY,\ldots.
\]
Its limit is the characteristic Fibonacci word of slope \(\theta\).
The convergents
\[
  \frac12,\quad\frac23,\quad\frac35,\quad\frac58,\ldots
\]
give the Christoffel approximants
\[
  C_{1,1},\quad C_{1,2},\quad C_{2,3},\quad C_{3,5},\ldots,
\]
up to the cyclic origins described in the previous chapter.  Their lengths
are Fibonacci numbers, and their densities of \(Y\)'s converge to
\(\theta\).

For every \(n\), the Fibonacci word has exactly \(n+1\) factors of length
\(n\), and any two such factors differ by at most one in their number of
\(Y\)'s.  Its factor frequencies are the gaps between finitely many points
of the orbit of the golden rotation and therefore belong to
\[
  \mathbb Z+\mathbb Z\frac{\sqrt5-1}{2}.
\]
This example displays the entire chain developed in the last three
chapters: continued fractions produce Christoffel words, their nested limit
is Sturmian, and the same irrational number controls both the symbolic
frequencies and the ordered \(K_0\)-group of the noncommutative torus.

\chapter{Ostrowski Numeration and Return-Time Towers} 

The denominators of the convergents of an irrational number do more than
measure the quality of its rational approximations.  They form a canonical
numeration scale for the integers, determine the recursive decomposition of
the characteristic Sturmian word, and give the two natural return-time
towers of the irrational rotation.  This numeration was introduced by
Ostrowski \cite{Ostrowski1922}; our conventions agree with
\cite[Section~3.9]{Lothaire2002}.

Throughout this chapter,
\[
  \theta=[0;a_1,a_2,\ldots]\in(0,1)\setminus\mathbb Q,
\]
and \(p_n/q_n=[0;a_1,\ldots,a_n]\), with
\[
  q_{-1}=0,\qquad q_0=1,\qquad
  q_{n+1}=a_{n+1}q_n+q_{n-1}.
\]

\section{The Ostrowski Expansion of an Integer}

\begin{definition}\label{definition:OstrowskiAdmissible}
  A finite digit sequence \((d_0,d_1,\ldots,d_m)\) is
  \textnormal{Ostrowski-admissible} for \(\theta\) if
  \begin{enumerate}
    \item \(0\leq d_0<a_1\);
    \item \(0\leq d_k\leq a_{k+1}\) for \(k\geq1\);
    \item if \(d_k=a_{k+1}\), then \(d_{k-1}=0\).
  \end{enumerate}
\end{definition}

The last condition is the carry rule.  In an ordinary base, a maximal digit
is replaced by one unit in the next position.  Here the replacement is
governed by
\[
  q_{k+1}=a_{k+1}q_k+q_{k-1}.
\]

\begin{lemma}\label{lemma:OstrowskiLowerDigitBound}
  For every \(m\geq1\), the largest integer representable using admissible
  digits in the positions \(0,ldots,m-1\) is \(q_m-1\).
\end{lemma}
\begin{proof}
  The assertion is clear for \(m=1\), since
  \(0\leq d_0<a_1=q_1\).  Suppose it holds below position \(m\).
  If \(d_{m-1}<a_m\), then
  \[
    \sum_{k=0}^{m-1}d_kq_k
    \leq(a_m-1)q_{m-1}+q_{m-1}-1
    =a_mq_{m-1}-1<q_m.
  \]
  If \(d_{m-1}=a_m\), admissibility forces \(d_{m-2}=0\), and the induction
  hypothesis gives
  \[
    \sum_{k=0}^{m-1}d_kq_k
    \leq a_mq_{m-1}+q_{m-2}-1=q_m-1.
  \]
  Equality is obtained recursively by taking the maximal digit in position
  \(m-1\), zero in position \(m-2\), and a maximal admissible tail below it.
\end{proof}

\begin{theorem}[Ostrowski numeration]
  \label{theorem:OstrowskiNumeration}
  Every integer \(N\geq0\) has a unique representation
  \begin{equation}\label{eq:OstrowskiExpansion}
    N=\sum_{k=0}^{m}d_kq_k
  \end{equation}
  with an Ostrowski-admissible digit sequence.
\end{theorem}
\begin{proof}
  Choose \(m\) such that \(q_m\leq N<q_{m+1}\), and put
  \[
    d_m=\left\lfloor\frac{N}{q_m}\right\rfloor,
    \qquad N_1=N-d_mq_m.
  \]
  Since \(N<q_{m+1}=a_{m+1}q_m+q_{m-1}\), one has
  \(d_m\leq a_{m+1}\).  If equality holds, then
  \(N_1<q_{m-1}\), so the next digit is zero.  Repeating the same greedy
  step with the remainder produces an admissible expansion and terminates
  at \(q_0=1\).

  For uniqueness, suppose two admissible expansions differ, and let \(j\)
  be the largest position at which their digits differ.  After interchanging
  them, the difference at position \(j\) contributes at least \(q_j\).
  By Lemma~\ref{lemma:OstrowskiLowerDigitBound}, all lower positions in
  either expansion contribute at most \(q_j-1\).  They cannot cancel the
  leading difference, a contradiction.
\end{proof}

\begin{remark}
  The proof is an honest greedy algorithm, but the admissibility condition
  is essential.  Without it,
  \[
    a_{k+1}q_k+q_{k-1}=q_{k+1}
  \]
  would give two representations of the same integer.
\end{remark}

\section{Digits and the Orbit of the Rotation}

Define the positive approximation errors
\begin{equation}\label{eq:PositiveConvergentErrors}
  \delta_n=(-1)^n(q_n\theta-p_n)>0.
\end{equation}
They satisfy
\begin{equation}\label{eq:ErrorRecursion}
  \delta_{n-1}=a_{n+1}\delta_n+\delta_{n+1},
  \qquad n\geq0,
\end{equation}
where \(\delta_{-1}=1\).  Thus the errors obey the same Euclidean recursion
as the place values, but in the reverse direction.

\begin{proposition}\label{proposition:OstrowskiOrbitFormula}
  If \(N=\sum_{k=0}^m d_kq_k\) is the Ostrowski expansion of \(N\), then
  \begin{equation}\label{eq:OstrowskiOrbitFormula}
    N\theta
    \equiv\sum_{k=0}^m(-1)^kd_k\delta_k\pmod1.
  \end{equation}
\end{proposition}
\begin{proof}
  Equation \eqref{eq:PositiveConvergentErrors} gives
  \[
    q_k\theta=p_k+(-1)^k\delta_k.
  \]
  Multiply by \(d_k\), sum over \(k\), and discard the integer
  \(\sum d_kp_k\).
\end{proof}

The formula is a positional description of the orbit of zero under
\(R_\theta\): the large integer place values \(q_k\) are replaced, modulo
one, by the small signed displacements \((-1)^k\delta_k\).  Carries are
compatible because \eqref{eq:ErrorRecursion} is precisely the image of the
denominator recursion under this substitution.

\section{Ostrowski Digits and Sturmian Prefixes}

Let \((T_n)\) be the standard sequence from
\eqref{eq:SturmianStandardWords}.  Recall that \(|T_n|=q_n\) and
\[
  T_{n+1}=T_n^{a_{n+1}}T_{n-1}.
\]

\begin{theorem}[prefix decomposition]
  \label{theorem:OstrowskiPrefixDecomposition}
  Let \(N=\sum_{k=0}^m d_kq_k\) be its Ostrowski expansion.  The prefix of
  length \(N\) of the characteristic word \(c_\theta\) is
  \begin{equation}\label{eq:OstrowskiPrefixDecomposition}
    T_m^{d_m}T_{m-1}^{d_{m-1}}\cdots T_1^{d_1}T_0^{d_0}.
  \end{equation}
\end{theorem}
\begin{proof}
  We argue by induction on the largest nonzero digit.  The characteristic
  word begins with
  \[
    T_{m+1}=T_m^{a_{m+1}}T_{m-1}.
  \]
  If \(d_m<a_{m+1}\), after the first \(d_m\) copies of \(T_m\) the
  remaining lower-digit block has length less than \(q_m\), by
  Lemma~\ref{lemma:OstrowskiLowerDigitBound}, and the induction hypothesis
  applies inside the next copy of \(T_m\).  If
  \(d_m=a_{m+1}\), admissibility gives \(d_{m-1}=0\); the remaining block
  has length less than \(q_{m-1}\) and lies in the following copy of
  \(T_{m-1}\).  This proves \eqref{eq:OstrowskiPrefixDecomposition}.
\end{proof}

Theorem~\ref{theorem:OstrowskiPrefixDecomposition} turns numerical digits
into a literal factorization of a word.  In particular, locating a position
inside \(c_\theta\) and expanding its index are the same recursive
operation.

\section{The Two Return-Time Towers}

The integer \(q_n\) is a closest return time for the rotation:
\[
  \|q_n\theta\|=\delta_n.
\]
Let \(I_n\) be the half-open short arc between \(0\) and
\(q_n\theta+\mathbb Z\), oriented from \(0\) to \(q_n\theta\) when \(n\)
is even and in the opposite direction when \(n\) is odd.  Thus
\(m(I_n)=\delta_n\).

\begin{lemma}\label{lemma:TowerMeasureIdentity}
  For every \(n\geq0\),
  \begin{equation}\label{eq:TowerMeasureIdentity}
    q_{n+1}\delta_n+q_n\delta_{n+1}=1.
  \end{equation}
\end{lemma}
\begin{proof}
  Using \eqref{eq:PositiveConvergentErrors} and the alternating signs,
  \begin{align*}
    q_{n+1}\delta_n+q_n\delta_{n+1}
    &=(-1)^n\bigl(q_{n+1}(q_n\theta-p_n)
       -q_n(q_{n+1}\theta-p_{n+1})\bigr)\\
    &=(-1)^n(q_np_{n+1}-q_{n+1}p_n)=1,
  \end{align*}
  by the unimodular identity for consecutive convergents.
\end{proof}

\begin{theorem}[continued-fraction towers]
  \label{theorem:ContinuedFractionTowers}
  The two families
  \[
    I_n,R_\theta I_n,\ldots,R_\theta^{q_{n+1}-1}I_n
  \]
  and
  \[
    I_{n+1},R_\theta I_{n+1},\ldots,
    R_\theta^{q_n-1}I_{n+1}
  \]
  have pairwise disjoint interiors and partition \(\mathbb T\).  They are
  two Rokhlin towers of heights \(q_{n+1}\) and \(q_n\), with base lengths
  \(\delta_n\) and \(\delta_{n+1}\).
\end{theorem}
\begin{proof}
  If two interiors in the first family overlapped, some integer
  \(r\), \(0<r<q_{n+1}\), would satisfy
  \(\|r\theta\|<\delta_n=\|q_n\theta\|\).  This contradicts the
  best-approximation property of the convergent \(q_n\).  The same argument
  applies to the second family.  The alternating positions of
  \(q_n\theta\) and \(q_{n+1}\theta\) show that the two families meet only
  at endpoints; equivalently, their endpoints occur consecutively in the
  orbit segment determined by \(q_n\) and \(q_{n+1}\).  Finally, their
  total measure is
  \[
    q_{n+1}m(I_n)+q_nm(I_{n+1})
    =q_{n+1}\delta_n+q_n\delta_{n+1}=1
  \]
  by Lemma~\ref{lemma:TowerMeasureIdentity}.  Disjointness and full measure
  give the asserted half-open partition.
\end{proof}

\begin{remark}
  The heights are not arbitrary recurrence times.  They are consecutive
  denominators of convergents, and the widths are precisely their
  Diophantine errors.  Formula \eqref{eq:TowerMeasureIdentity} says that the
  two towers fill the circle with no remainder.
\end{remark}

\section{Refinement Matrices}

Put
\begin{equation}\label{eq:ContinuedFractionIncidenceMatrix}
  M_{n+1}=
  \begin{pmatrix}
    a_{n+1}&1\\
    1&0
  \end{pmatrix}.
\end{equation}
The denominator and error recursions can be written as
\begin{equation}\label{eq:HeightWidthMatrixRecursions}
  \binom{q_{n+1}}{q_n}
  =M_{n+1}\binom{q_n}{q_{n-1}},
  \qquad
  (\delta_n,\delta_{n+1})M_{n+1}
  =(\delta_{n-1},\delta_n).
\end{equation}

\begin{proposition}\label{proposition:TowerRefinementMatrix}
  Passing from the tower partition at level \(n-1\) to that at level \(n\)
  is governed by the matrix \(M_{n+1}\): one old tower is cut into
  \(a_{n+1}\) copies of one new tower and one copy of the other, while the
  second old tower becomes one copy of the first new tower.  The same
  matrix simultaneously transforms tower heights and tower widths as in
  \eqref{eq:HeightWidthMatrixRecursions}.
\end{proposition}
\begin{proof}
  The first column of \(M_{n+1}\) records the Euclidean decomposition
  \[
    q_{n+1}=a_{n+1}q_n+q_{n-1},
  \]
  and the second column records the unchanged height \(q_n\).  On the
  bases, the reverse recursion
  \[
    \delta_{n-1}=a_{n+1}\delta_n+\delta_{n+1}
  \]
  makes exactly the corresponding cuts.  Translating the base intervals by
  \(R_\theta\) propagates the cuts through every level of the towers.
\end{proof}

The matrices \(M_{n+1}\) have determinant \(-1\).  They are therefore
unimodular, while all their entries are nonnegative.  This combination is
what allows them to encode both an arithmetic change of basis and a
positive refinement.

\section{The Arithmetic Dimension Group}

For \(n\geq1\), let \(G_n=\mathbb Z^2\), with positive cone
\(G_n^+=\mathbb N^2\), and use
\[
  \varphi_n\colon G_n\longrightarrow G_{n+1},
  \qquad \varphi_n(v)=M_{n+1}v
\]
as connecting map.  Put
\[
  G_\theta=\varinjlim(G_n,\varphi_n).
\]
The vector
\[
  u_n=\binom{q_n}{q_{n-1}}
\]
is compatible with the connecting maps and determines an order unit
\(u\in G_\theta\).

\begin{theorem}[arithmetic dimension group]
  \label{theorem:OstrowskiDimensionGroup}
  There is an order-unit isomorphism
  \begin{equation}\label{eq:OstrowskiDimensionGroup}
    (G_\theta,G_\theta^+,u)
    \cong
    \bigl(\mathbb Z+\mathbb Z\theta,
      (\mathbb Z+\mathbb Z\theta)\cap[0,\infty),1\bigr).
  \end{equation}
  At stage \(n\), it is induced by
  \begin{equation}\label{eq:OstrowskiStageState}
    (x,y)\longmapsto x\delta_{n-1}+y\delta_n.
  \end{equation}
\end{theorem}
\begin{proof}
  The second identity in \eqref{eq:HeightWidthMatrixRecursions} says exactly
  that the maps \eqref{eq:OstrowskiStageState} are compatible with the
  connecting maps.  They therefore define a homomorphism
  \(\Phi\colon G_\theta\to\mathbb R\).

  The two numbers \(\delta_{n-1},\delta_n\) form a \(\mathbb Z\)-basis of
  \(\mathbb Z+\mathbb Z\theta\).  Indeed, their coefficient vectors in the
  basis \((1,\theta)\) have determinant of absolute value one, by the
  unimodular identity for \(p_n/q_n\).  Thus \(\Phi\) is injective and has
  image \(\mathbb Z+\mathbb Z\theta\).

  Nonnegative vectors clearly map to nonnegative real numbers.  Conversely,
  let \(g\in\mathbb Z+\mathbb Z\theta\) be positive.  After subtracting a
  nonnegative integer, it is enough to consider \(0\leq g\leq1\).  By
  Theorem~\ref{theorem:SturmianFrequencyGroup}, \(g\) is the measure of a
  finite union of orbit intervals.  At a sufficiently fine return-time
  partition, that union consists of finitely many levels of the two towers,
  and hence
  \[
    g=x\delta_{n-1}+y\delta_n
  \]
  for some \(x,y\geq0\).  Therefore \(\Phi\) maps the inductive-limit
  positive cone onto the positive cone on the right of
  \eqref{eq:OstrowskiDimensionGroup}.

  Finally, Lemma~\ref{lemma:TowerMeasureIdentity} gives
  \[
    \Phi(u_n)=q_n\delta_{n-1}+q_{n-1}\delta_n=1.
  \]
  Thus \(\Phi\) also preserves the order unit.
\end{proof}

Combining Theorem~\ref{theorem:OstrowskiDimensionGroup} with
Theorem~\ref{thm:orderedK0RotationAlgebra} yields
\[
  (G_\theta,G_\theta^+,u)
  \cong
  \bigl(K_0(A_\theta),K_0(A_\theta)^+,[1]_0\bigr).
\]
At this point the left-hand side is only an arithmetic ordered limit.  Its
realization by finite-dimensional algebras and a Bratteli diagram belongs
to the next stage of the book.

\section{The Golden Case and Zeckendorf's Theorem}

For
\[
  \theta=\frac{\sqrt5-1}{2}=[0;1,1,1,\ldots],
\]
the denominator sequence is
\[
  q_0=1,\quad q_1=1,\quad q_2=2,\quad q_3=3,\quad
  q_4=5,\quad q_5=8,\ldots.
\]
The admissible digits satisfy \(d_0=0\), \(d_k\in\{0,1\}\), and no two
consecutive digits are equal to one.  Theorem~\ref{theorem:OstrowskiNumeration}
therefore becomes Zeckendorf's theorem.

\begin{corollary}[Zeckendorf]
  Every positive integer is uniquely a sum of nonconsecutive Fibonacci
  numbers, with the convention that the duplicated initial \(1\) is used
  only once.
\end{corollary}

For example,
\[
  100=89+8+3=q_{10}+q_5+q_3.
\]
The corresponding prefix of the Fibonacci word is
\[
  T_{10}T_5T_3.
\]
The tower incidence matrix is constant,
\[
  M=
  \begin{pmatrix}1&1\\1&0\end{pmatrix},
\]
and its Perron--Frobenius eigenvalue is the golden ratio.  Thus Fibonacci
numeration, the Fibonacci word, the golden rotation and the ordered group
\[
  \mathbb Z+\mathbb Z\frac{\sqrt5-1}{2}
\]
are four manifestations of the same stationary recursion.

\chapter{The Effros--Shen AF Algebra}

The arithmetic limit constructed in the preceding chapter is not merely an
abstract ordered group.  The continued-fraction matrices are multiplicity
matrices of embeddings between finite-dimensional \(C^*\)-algebras, and
therefore determine an AF-algebra.  This construction is due to Effros and
Shen \cite{EffrosShen1980} and is naturally expressed by the diagrams
introduced by Bratteli \cite{Bratteli1972}.

Throughout the chapter,
\[
  \theta=[0;a_1,a_2,\ldots]\in(0,1)\setminus\mathbb Q
\]
and \(p_n/q_n\) denotes its sequence of convergents.  We shall denote the
resulting AF-algebra by \(\mathcal E_\theta\), in order to distinguish it
from the irrational rotation algebra \(A_\theta\).

\section{Finite-Dimensional Systems and Bratteli Diagrams}

\begin{definition}
  A unital \(C^*\)-algebra \(A\) is \textnormal{approximately
  finite-dimensional}, or an \textnormal{AF-algebra}, if there is an
  increasing sequence of finite-dimensional unital \(C^*\)-subalgebras
  \[
    F_1\subseteq F_2\subseteq\cdots\subseteq A
  \]
  whose union is dense in \(A\).  Equivalently, \(A\) is the inductive
  limit of a sequence of finite-dimensional \(C^*\)-algebras with unital
  injective connecting maps.
\end{definition}

Every finite-dimensional \(C^*\)-algebra is a finite direct sum of full
matrix algebras.  The following observation converts a homomorphism between
two such algebras into an integer matrix.

\begin{proposition}[multiplicity matrix]\label{proposition:MultiplicityMatrix}
  Let
  \[
    F=\bigoplus_{j=1}^{r}M_{d_j}(\mathbb C),
    \qquad
    H=\bigoplus_{i=1}^{s}M_{e_i}(\mathbb C),
  \]
  and let \(\varphi\colon F\to H\) be a unital *-homomorphism.  There is a
  unique matrix \(N_\varphi=(m_{ij})\in M_{s,r}(\mathbb N_0)\) such that,
  up to a unitary in each summand of \(H\), the \(i\)-th component of
  \(\varphi\) is
  \[
    (x_1,\ldots,x_r)\longmapsto
    \diag(\underbrace{x_1,\ldots,x_1}_{m_{i1}},\ldots,
          \underbrace{x_r,\ldots,x_r}_{m_{ir}}).
  \]
  In particular,
  \[
    e_i=\sum_{j=1}^{r}m_{ij}d_j.
  \]
  If \(K_0(F)\cong\mathbb Z^r\) and
  \(K_0(H)\cong\mathbb Z^s\) are identified using the classes of
  rank-one projections in the matrix summands, then
  \[
    K_0(\varphi)(v)=N_\varphi v.
  \]
\end{proposition}

\begin{proof}
  Compose \(\varphi\) with the projection onto the \(i\)-th summand of
  \(H\).  This gives a finite-dimensional representation of \(F\) on
  \(\mathbb C^{e_i}\).  The irreducible representations of \(F\) are the
  defining representations of its matrix summands, so this representation
  is a direct sum of \(m_{ij}\) copies of the defining representation of
  \(M_{d_j}(\mathbb C)\).  The multiplicities are unique, and comparison of
  dimensions gives the displayed identity for \(e_i\).

  A rank-one projection in the \(j\)-th summand of \(F\) is carried to a
  projection of rank \(m_{ij}\) in the \(i\)-th summand of \(H\).  Thus its
  \(K_0\)-class is the \(j\)-th column of \(N_\varphi\), which proves the
  last assertion.
\end{proof}

\begin{definition}
  The \textnormal{Bratteli diagram} of an inductive sequence of
  finite-dimensional \(C^*\)-algebras has one vertex for each matrix
  summand at every level.  Between the vertices corresponding to the
  \(j\)-th summand at level \(n\) and the \(i\)-th summand at level \(n+1\),
  it has \(m_{ij}\) edges, where \((m_{ij})\) is the multiplicity matrix of
  the connecting map.
\end{definition}

Thus a Bratteli diagram records simultaneously the embeddings of matrix
blocks and the positive homomorphisms induced on \(K_0\).  Passing to a
subsequence of levels replaces consecutive multiplicity matrices by their
products and does not change the inductive limit; this operation is called
\textnormal{telescoping} the diagram.

\section{The Effros--Shen Inductive System}

For \(n\geq1\), put
\[
  F_n=M_{q_n}(\mathbb C)\oplus M_{q_{n-1}}(\mathbb C).
\]
Define a unital *-homomorphism \(\varphi_n\colon F_n\to F_{n+1}\) by
\begin{equation}\label{eq:EffrosShenEmbedding}
  \varphi_n(x,y)=
  \left(
    \diag(\underbrace{x,\ldots,x}_{a_{n+1}},y),x
  \right).
\end{equation}

\begin{lemma}\label{lemma:EffrosShenMultiplicity}
  The map \(\varphi_n\) is a unital injective *-homomorphism, and its
  multiplicity matrix is
  \[
    M_{n+1}=
    \begin{pmatrix}a_{n+1}&1\\1&0\end{pmatrix}.
  \]
  Moreover,
  \[
    M_{n+1}\binom{q_n}{q_{n-1}}
      =\binom{q_{n+1}}{q_n}.
  \]
\end{lemma}

\begin{proof}
  Formula \eqref{eq:EffrosShenEmbedding} is visibly a *-homomorphism and is
  injective because both \(x\) and \(y\) occur in its first component.  The
  first target block contains \(a_{n+1}\) copies of the first source block
  and one copy of the second, while the second target block contains one
  copy of the first and no copy of the second.  This gives the asserted
  multiplicity matrix.  Finally,
  \[
    q_{n+1}=a_{n+1}q_n+q_{n-1}
  \]
  shows both that the map is unital and that the displayed vector identity
  holds.
\end{proof}

The part of the Bratteli diagram joining two consecutive levels is
\[
  \begin{tikzcd}[column sep=large,row sep=large]
    M_{q_n}(\mathbb C)
      \arrow[r,bend left=12,"a_{n+1}"]
      \arrow[dr,"1"']
      &M_{q_{n+1}}(\mathbb C)\\
    M_{q_{n-1}}(\mathbb C)
      \arrow[ur,"1"]
      &M_{q_n}(\mathbb C).
  \end{tikzcd}
\]

\begin{definition}[Effros--Shen algebra]\label{definition:EffrosShenAlgebra}
  The \textnormal{Effros--Shen AF-algebra associated with \(\theta\)} is
  \[
    \mathcal E_\theta
      =\varinjlim(F_n,\varphi_n).
  \]
\end{definition}

\begin{theorem}\label{theorem:EffrosShenK0}
  The algebra \(\mathcal E_\theta\) is a unital AF-algebra and there is an
  order-unit isomorphism
  \[
    \bigl(K_0(\mathcal E_\theta),K_0(\mathcal E_\theta)^+,[1]_0\bigr)
    \cong
    \bigl(\mathbb Z+\mathbb Z\theta,
           (\mathbb Z+\mathbb Z\theta)\cap[0,\infty),1\bigr).
  \]
\end{theorem}

\begin{proof}
  The inductive limit is unital because all the connecting maps are unital,
  and the images of the finite-dimensional algebras \(F_n\) have dense
  union.  Hence it is an AF-algebra.

  With rank-one projections as generators,
  \[
    \bigl(K_0(F_n),K_0(F_n)^+,[1_{F_n}]_0\bigr)
      \cong
    \bigl(\mathbb Z^2,\mathbb N_0^2,(q_n,q_{n-1})\bigr).
  \]
  By Proposition~\ref{proposition:MultiplicityMatrix} and
  Lemma~\ref{lemma:EffrosShenMultiplicity}, the connecting map on \(K_0\)
  is multiplication by \(M_{n+1}\).  The continuity of \(K_0\), proved in
  Chapter~9, therefore identifies the scaled ordered \(K_0\)-group of
  \(\mathcal E_\theta\) with the arithmetic limit \(G_\theta\) of the
  preceding chapter.  The conclusion now follows from
  Theorem~\ref{theorem:OstrowskiDimensionGroup}.
\end{proof}

\section{Simplicity}

Although every level of the construction has two summands, these summands
are mixed by the connecting maps.

\begin{lemma}\label{lemma:EffrosShenPositiveProducts}
  Every product of two consecutive incidence matrices is strictly positive:
  \[
    M_{n+2}M_{n+1}
      =
    \begin{pmatrix}
      a_{n+2}a_{n+1}+1&a_{n+2}\\
      a_{n+1}&1
    \end{pmatrix}.
  \]
  Consequently, from either vertex at level \(n\) there is a path to each
  vertex at level \(n+2\).
\end{lemma}

\begin{proof}
  This is direct matrix multiplication.  All partial quotients are positive,
  so every entry of the product is positive.  By the definition of a
  Bratteli diagram, its \((i,j)\)-entry counts paths from vertex \(j\) to
  vertex \(i\) through the intermediate level.
\end{proof}

\begin{theorem}\label{theorem:EffrosShenSimple}
  The Effros--Shen algebra \(\mathcal E_\theta\) is simple.
\end{theorem}

\begin{proof}
  We first note a general fact.  If \(I\) is a nonzero ideal in a unital
  AF-algebra with an increasing dense union \(\bigcup_nF_n\), then
  \(I\cap F_n\neq0\) for some \(n\).  Indeed, otherwise the quotient map
  would be injective, hence isometric, on every \(F_n\); it would then be
  isometric on their dense union and therefore on the whole algebra, which
  would force \(I=0\).

  Apply this fact to a nonzero ideal \(I\) of \(\mathcal E_\theta\).  The
  nonzero ideal \(I\cap F_n\) contains one of the two simple summands of
  \(F_n\).  By Lemma~\ref{lemma:EffrosShenPositiveProducts}, the image of
  this summand has a nonzero component in both simple summands of
  \(F_{n+2}\).  The ideal it generates in \(F_{n+2}\) is therefore all of
  \(F_{n+2}\).  Hence \(1_{F_{n+2}}\in I\), and thus
  \(I=\mathcal E_\theta\).
\end{proof}

\section{The Canonical Trace}

Let \(\operatorname{Tr}_r\) denote the unnormalized trace on
\(M_r(\mathbb C)\).  Recall the approximation errors
\[
  \delta_n=(-1)^n(q_n\theta-p_n)>0,
  \qquad \delta_{-1}=1,
\]
and their reverse recursion
\[
  \delta_{n-1}=a_{n+1}\delta_n+\delta_{n+1}.
\]

\begin{proposition}\label{proposition:EffrosShenTrace}
  For \(n\geq1\), define
  \[
    \tau_n(x,y)
      =\delta_{n-1}\operatorname{Tr}_{q_n}(x)
       +\delta_n\operatorname{Tr}_{q_{n-1}}(y),
    \qquad (x,y)\in F_n.
  \]
  Then \(\tau_n\) is a tracial state and
  \[
    \tau_{n+1}\circ\varphi_n=\tau_n.
  \]
  Consequently, there is a tracial state \(\tau_\theta\) on
  \(\mathcal E_\theta\) whose restriction to \(F_n\) is \(\tau_n\).
\end{proposition}

\begin{proof}
  Positivity and the trace identity are immediate from the corresponding
  properties of the matrix traces.  Lemma~\ref{lemma:TowerMeasureIdentity}
  gives
  \[
    \tau_n(1_{F_n})
      =q_n\delta_{n-1}+q_{n-1}\delta_n=1,
  \]
  so \(\tau_n\) is a state.  Using
  \eqref{eq:EffrosShenEmbedding} and the reverse recursion, we obtain
  \begin{align*}
    \tau_{n+1}(\varphi_n(x,y))
      &=(a_{n+1}\delta_n+\delta_{n+1})
          \operatorname{Tr}_{q_n}(x)
        +\delta_n\operatorname{Tr}_{q_{n-1}}(y)\\
      &=\tau_n(x,y).
  \end{align*}
  The universal property of the inductive limit gives a trace on the
  algebraic union of the images of the \(F_n\), and its boundedness extends
  it uniquely to \(\mathcal E_\theta\).
\end{proof}

\medskip

\begin{theorem}\label{theorem:EffrosShenUniqueTrace}
  The trace \(\tau_\theta\) is the unique tracial state on
  \(\mathcal E_\theta\).  Moreover,
  \[
    \{\tau_\theta(p):p\text{ is a projection in }\mathcal E_\theta\}
      =(\mathbb Z+\mathbb Z\theta)\cap[0,1].
  \]
\end{theorem}

\begin{proof}
  Every tracial state \(\sigma\) induces a positive homomorphism
  \(\sigma_*\colon K_0(\mathcal E_\theta)\to\mathbb R\) satisfying
  \(\sigma_*([1]_0)=1\).  Under the isomorphism of
  Theorem~\ref{theorem:EffrosShenK0}, this becomes a positive homomorphism
  \[
    \rho\colon\mathbb Z+\mathbb Z\theta\longrightarrow\mathbb R,
    \qquad \rho(1)=1.
  \]
  Such a homomorphism must be the inclusion.  Indeed, for \(h\) in this
  group and every positive integer \(m\), put \(k=\lfloor mh\rfloor\).
  Positivity applied to
  \(k\leq mh\leq k+1\) gives
  \[
    k\leq m\rho(h)\leq k+1,
  \]
  and hence \(|\rho(h)-h|\leq1/m\).  Letting \(m\) tend to infinity gives
  \(\rho(h)=h\).

  It follows that any tracial state agrees with \(\tau_\theta\) on every
  projection in every finite-dimensional stage.  A trace on a matrix
  algebra is determined by its values on projections, so the traces agree
  on \(\bigcup_nF_n\), and therefore on its closure.  This proves
  uniqueness.

  If \(p\) is a projection, its \(K_0\)-class lies between \(0\) and
  \([1]_0\), so the first inclusion in the asserted equality follows from
  Theorem~\ref{theorem:EffrosShenK0}.  Conversely, let
  \(h\in(\mathbb Z+\mathbb Z\theta)\cap[0,1]\).  Its class and the class of
  \(1-h\) are positive in the inductive-limit group.  At a sufficiently
  large common stage they are represented by vectors \(v,w\in\mathbb N_0^2\)
  whose sum is the order-unit vector.  Choosing projections of the
  corresponding ranks in the two matrix summands gives a projection
  \(p\in F_n\) with \(\tau_\theta(p)=h\).
\end{proof}

\section{Comparison with the Irrational Rotation Algebra}

We can now state precisely what the AF model captures and what it does not.

\begin{theorem}\label{theorem:EffrosShenRotationComparison}
  There is an isomorphism of scaled ordered groups
  \[
    \bigl(K_0(\mathcal E_\theta),K_0(\mathcal E_\theta)^+,[1]_0\bigr)
    \cong
    \bigl(K_0(A_\theta),K_0(A_\theta)^+,[1]_0\bigr)
  \]
  which intertwines the homomorphisms induced by the canonical traces.
  Nevertheless, \(\mathcal E_\theta\) and \(A_\theta\) are neither
  isomorphic nor strongly Morita equivalent.
\end{theorem}

\begin{proof}
  The first assertion follows by combining
  Theorem~\ref{theorem:EffrosShenK0} with
  Theorem~\ref{thm:orderedK0RotationAlgebra}.  Under both identifications,
  the state induced by the canonical trace is the inclusion
  \[
    \mathbb Z+\mathbb Z\theta\hookrightarrow\mathbb R,
  \]
  so the trace homomorphisms are intertwined.

  Each finite-dimensional algebra \(F_n\) has trivial \(K_1\).  By
  continuity of \(K_1\),
  \[
    K_1(\mathcal E_\theta)=0.
  \]
  On the other hand, the Pimsner--Voiculescu computation in
  Chapter~13 gives
  \[
    K_1(A_\theta)\cong\mathbb Z^2.
  \]
  Thus the algebras cannot be isomorphic.  Since \(K\)-theory is invariant
  under strong Morita equivalence, they cannot be strongly Morita
  equivalent either.
\end{proof}

\begin{remark}
  The equality of ordered \(K_0\)-groups is therefore not an identification
  of the two algebras.  It says that their projections have the same
  arithmetic order and the same trace values.  The missing \(K_1\)-group
  records genuinely circular information which no AF-algebra can retain.
\end{remark}

\section{Periodic Diagrams}

The continued fraction can be read directly from the diagram: the number
of horizontal edges between levels \(n\) and \(n+1\) is \(a_{n+1}\).
Consequently, arithmetic periodicity becomes operator-algebraic
periodicity.

\begin{proposition}\label{proposition:EffrosShenPeriodic}
  Suppose that the sequence of partial quotients is eventually periodic
  with period \(r\).  After deleting finitely many levels and telescoping
  every \(r\) levels, the Bratteli diagram of \(\mathcal E_\theta\) has the
  constant incidence matrix
  \[
    P=
    \begin{pmatrix}a_{N+r}&1\\1&0\end{pmatrix}
    \cdots
    \begin{pmatrix}a_{N+1}&1\\1&0\end{pmatrix}.
  \]
  The matrix \(P\) is strictly positive when \(r\geq2\); when \(r=1\),
  \(P^2\) is strictly positive.
\end{proposition}

\begin{proof}
  Telescoping replaces the incidence matrices in each block of \(r\)
  successive levels by their product.  Eventual periodicity makes this
  product independent of the block, giving \(P\).  If \(r\geq2\), the
  product contains two consecutive matrices and is strictly positive by
  Lemma~\ref{lemma:EffrosShenPositiveProducts}.  The same lemma gives the
  assertion about \(P^2\) when \(r=1\).
\end{proof}

For the golden rotation, all partial quotients are equal to one and the
diagram is stationary from the beginning, with incidence matrix
\[
  \begin{pmatrix}1&1\\1&0\end{pmatrix}.
\]
Thus the Fibonacci substitution matrix, the return-time refinement matrix
and the incidence matrix of the Effros--Shen algebra are literally the same
matrix.

\section{End of the Arithmetic Bridge}

The construction of this chapter completes the second part of the book.
Starting from the irrational parameter \(\theta\), we have obtained the
chain
\[
  \begin{tikzcd}[column sep=small,row sep=small]
    \text{continued fractions}
      \arrow[r]&
    \text{Christoffel and Sturmian words}
      \arrow[r]&
    \text{return-time towers}
      \arrow[d]\\
    K_0(A_\theta)
      &&
    \text{Bratteli diagram}
      \arrow[ll]
  \end{tikzcd}
\]
The denominators \(q_n\) have appeared successively as best-approximation
denominators, word lengths, tower heights and matrix sizes.  The errors
\(\delta_n\) have appeared as approximation errors, interval lengths and
trace weights.  The matrices
\[
  \begin{pmatrix}a_{n+1}&1\\1&0\end{pmatrix}
\]
have appeared as Euclidean, substitution, refinement and multiplicity
matrices.

What remains is no longer to identify the common arithmetic invariant, but
to compare the algebras themselves by increasingly accurate
finite-dimensional and circle-algebra models.  That is the starting point
of Part~III.

\part{Finite-Dimensional and Circle-Algebra Approximations}

\chapter{Rational Rotation Algebras and Matrix Fibres}

The arithmetic of Part~II associates to the irrational parameter \(\theta\)
its convergents \(p_n/q_n\).  The first purpose of the present part is to
understand what the rational parameters \(p/q\) contribute to the study of
\(A_\theta\).  They provide exact finite-dimensional representations of the
rational rotation algebras and approximate matrix solutions of the defining
relation of \(A_\theta\); by themselves, however, they define neither
finite-dimensional representations of \(A_\theta\) nor canonical inclusions
of these matrix fibres into it.  This distinction will be essential
throughout the sequel.

\section{The rational rotation algebra}

Let \(p\) and \(q\) be relatively prime integers with \(q\geq1\), and put
\[
  \omega=\exp(2\pi i p/q).
\]
The \textnormal{rational rotation algebra} \(A_{p/q}\) is the universal
unital \cstar\ generated by two unitaries \(U\) and \(V\) satisfying
\begin{equation}\label{eq:RationalRotationRelation}
  VU=\omega UV.
\end{equation}
Only the class of \(p/q\) modulo \(1\) matters.  We nevertheless retain a
coprime presentation because the denominator \(q\) will be the dimension of
the matrix fibres below.  The standard rational case is discussed in
\cite[Chapter~VI, Exercise~VI.1]{Davidson1996}; we give the details needed for
the later approximation argument.

The rational and irrational cases have sharply different ideal structures.
For irrational \(\theta\), the algebra \(A_\theta\) is simple.  In contrast,
the powers \(U^q\) and \(V^q\) are central in \(A_{p/q}\), since
\[
  VU^q=\omega^qU^qV=U^qV,
  \qquad
  V^qU=\omega^qUV^q=UV^q.
\]
Thus the rational rotation algebra has a nontrivial centre.  This already
precludes its identification with a full matrix algebra when \(q>1\).

\begin{proposition}\label{prop:RationalRotationCentre}
  Let \(p\) and \(q\) be relatively prime.  Then
  \[
    Z(A_{p/q})=C^*(U^q,V^q).
  \]
\end{proposition}

\begin{proof}
  The preceding computation shows that \(C^*(U^q,V^q)\) is contained in the
  centre.  For the converse, consider the gauge action
  \[
    \gamma\colon\mathbb T^2\longrightarrow\operatorname{Aut}(A_{p/q}),
    \qquad
    \gamma_{s,t}(U)=sU,
    \qquad \gamma_{s,t}(V)=tV.
  \]
  It is well-defined by the universal property.  For \((r,s)\in\mathbb Z^2\)
  and \(a\in A_{p/q}\), define its \((r,s)\)-Fourier coefficient by the
  Bochner integral
  \[
    E_{r,s}(a)=\int_{\mathbb T^2}\overline{x}^{\,r}\overline{y}^{\,s}
       \gamma_{x,y}(a)\,dx\,dy.
  \]
  On the dense \(*\)-subalgebra spanned by the monomials \(U^mV^n\), one has
  \[
    E_{r,s}(U^mV^n)=
    \begin{cases}
      U^rV^s,&(m,n)=(r,s),\\
      0,&(m,n)\neq(r,s).
    \end{cases}
  \]
  By continuity, the range of \(E_{r,s}\) is contained in the
  one-dimensional space \(\mathbb C U^rV^s\).

  Let \(a\) be central.  Each \(\gamma_{x,y}(a)\) is central, and hence so
  is every coefficient \(E_{r,s}(a)\).  Write
  \(E_{r,s}(a)=c_{r,s}U^rV^s\).  If \(c_{r,s}\neq0\), then commutation with
  \(U\) gives \(\omega^s=1\), and commutation with \(V\) gives
  \(\omega^r=1\).  As \(\omega\) has order \(q\), this implies
  \(q\mid r\) and \(q\mid s\).

  Finally, the two-variable Fej\'er sums
  \[
    \sigma_N(a)=
      \sum_{|r|,|s|\leq N}
      \left(1-\frac{|r|}{N+1}\right)
      \left(1-\frac{|s|}{N+1}\right)E_{r,s}(a)
  \]
  converge in norm to \(a\).  Indeed, they are the averages of the
  norm-continuous map \((x,y)\mapsto\gamma_{x,y}(a)\) against the product
  Fej\'er kernels, and these kernels form an approximate identity on
  \(\mathbb T^2\).  Every \(\sigma_N(a)\) belongs to
  \(C^*(U^q,V^q)\), by the preceding paragraph; the same is therefore true
  of \(a\).  This proves the reverse inclusion.
\end{proof}

The proposition is more informative than the assertion that the rational
case has matrix representations.  It says that the matrix algebra occurs
fibrewise over the spectrum of the centre; the algebra \(A_{p/q}\) itself
retains two circular central variables.

\section{Clock and shift matrices}

Let \((e_0,\ldots,e_{q-1})\) be the standard basis of \(\mathbb C^q\), with
subscripts read modulo \(q\).  Define unitaries \(D_{p,q}\) and \(C_q\) by
\[
  D_{p,q}e_j=\omega^{-j}e_j,
  \qquad
  C_qe_j=e_{j+1}.
\]
They are respectively called the \textnormal{clock} and
\textnormal{shift} matrices.  Directly from the definitions,
\begin{equation}\label{eq:ClockShiftRelation}
  C_qD_{p,q}=\omega D_{p,q}C_q,
  \qquad
  D_{p,q}^q=C_q^q=1_q.
\end{equation}

\begin{proposition}\label{prop:ClockShiftFibre}
  The assignment
  \[
    U\longmapsto D_{p,q},
    \qquad
    V\longmapsto C_q
  \]
  determines a surjective unital \(*\)-homomorphism
  \[
    \rho_{p,q}\colon A_{p/q}\longrightarrow M_q(\mathbb C).
  \]
  Moreover, \(\rho_{p,q}\) is irreducible.
\end{proposition}

\begin{proof}
  Equation~\eqref{eq:ClockShiftRelation} and the universal property give
  the homomorphism.  Since \(p\) and \(q\) are relatively prime, the
  eigenvalues of \(D_{p,q}\) are pairwise distinct.  Polynomial functional
  calculus therefore gives every diagonal matrix unit from
  \(D_{p,q}\).  Multiplication by powers of \(C_q\) then gives all the
  matrix units.  Hence the range is \(M_q(\mathbb C)\), and the
  representation is irreducible.
\end{proof}

More generally, let \((z,w)\in\mathbb T^2\), choose \(\alpha,\beta\in
\mathbb T\) with \(\alpha^q=z\) and \(\beta^q=w\), and set
\[
  \rho_{z,w}(U)=\alpha D_{p,q},
  \qquad
  \rho_{z,w}(V)=\beta C_q.
\]
The choice of roots changes the representation only by unitary equivalence.
Indeed, if \(\alpha'=\omega^k\alpha\) and
\(\beta'=\omega^\ell\beta\), conjugation by
\(C_q^kD_{p,q}^{-\ell}\) carries the pair determined by
\((\alpha,\beta)\) to the pair determined by \((\alpha',\beta')\).
The central character of \(\rho_{z,w}\) is given by
\[
  \rho_{z,w}(U^q)=z1_q,
  \qquad
  \rho_{z,w}(V^q)=w1_q.
\]

\begin{corollary}\label{cor:RationalRotationCentreSpectrum}
  The map
  \[
    C(\mathbb T^2)\longrightarrow Z(A_{p/q}),
    \qquad
    z_1\longmapsto U^q,
    \quad z_2\longmapsto V^q,
  \]
  is an isomorphism.
\end{corollary}

\begin{proof}
  Proposition~\ref{prop:RationalRotationCentre} shows that the displayed
  map is surjective.  For every \((z,w)\in\mathbb T^2\), the restriction of
  \(\rho_{z,w}\) to the centre is the character taking \(U^q\) to \(z\) and
  \(V^q\) to \(w\).  Thus the joint spectrum of \((U^q,V^q)\) is all of
  \(\mathbb T^2\), so the displayed homomorphism is injective.
\end{proof}

\begin{theorem}\label{thm:RationalRotationFibres}
  For every \((z,w)\in\mathbb T^2\), the quotient of \(A_{p/q}\) by the
  ideal generated by
  \[
    U^q-z1,
    \qquad V^q-w1
  \]
  is isomorphic to \(M_q(\mathbb C)\).  In particular, the irreducible
  representations of \(A_{p/q}\) are \(q\)-dimensional and are parametrized
  by the characters of its centre.
\end{theorem}

\begin{proof}
  The representation \(\rho_{z,w}\) vanishes on the displayed ideal, and
  Proposition~\ref{prop:ClockShiftFibre} shows that the resulting map onto
  \(M_q(\mathbb C)\) is surjective.  In the quotient, every monomial can be
  reduced, using \eqref{eq:RationalRotationRelation} and the two scalar
  relations, to a linear combination of
  \[
    U^rV^s,
    \qquad 0\leq r,s<q.
  \]
  Thus the quotient has dimension at most \(q^2\).  Since its quotient
  \(M_q(\mathbb C)\) has dimension \(q^2\), the surjection is an
  isomorphism.

  Let \(\pi\) be an irreducible representation of \(A_{p/q}\).  By Schur's
  lemma, \(\pi\) maps every central element to a scalar.  Corollary~
  \ref{cor:RationalRotationCentreSpectrum} therefore gives a point
  \((z,w)\in\mathbb T^2\) such that
  \[
    \pi(U^q)=z1,
    \qquad \pi(V^q)=w1.
  \]
  Hence \(\pi\) factors through the corresponding quotient
  \(M_q(\mathbb C)\).  Every irreducible representation of a full matrix
  algebra is unitarily equivalent to its defining representation on
  \(\mathbb C^q\).  Conversely, Proposition~\ref{prop:ClockShiftFibre}
  shows that every fibre representation \(\rho_{z,w}\) is irreducible.
  This proves the final assertion.
\end{proof}

\begin{remark}
  The matrices \(D_{p,q}\) and \(C_q\) are therefore not a presentation of
  \(A_{p/q}\) itself.  They describe its fibres after the two central
  unitaries have been specialized to scalars.  This is the precise sense in
  which clock-and-shift matrices are finite-dimensional models of the
  rational rotation algebra.
\end{remark}

\section{Rational fibres are not exact models}

It is tempting to use a convergent \(p_n/q_n\) and regard the associated
clock-and-shift pair as a finite-dimensional representation of \(A_\theta\).
No such representation is available.  Indeed, if there were a unital
\(*\)-homomorphism
\[
  A_\theta\longrightarrow M_q(\mathbb C),
\]
then simplicity of \(A_\theta\) would make it injective, which is
impossible because \(A_\theta\) is infinite-dimensional.  In particular,
the clock-and-shift pair is an exact representation of \(A_{p_n/q_n}\), but
only an approximate solution of the defining relation of \(A_\theta\).
This argument concerns homomorphisms from \(A_\theta\) into matrix algebras;
it does not rule out finite-dimensional subalgebras inside \(A_\theta\).
In particular, it provides no canonical inclusion of the rational matrix
fibre into \(A_\theta\).

This obstruction is consistent with the conclusion of Chapter~20.  The
finite-dimensional stages of the Effros--Shen algebra recover the ordered
\(K_0\)-data of \(A_\theta\), but have trivial \(K_1\).  Matrix fibres of
rational rotation algebras likewise retain no circular \(K_1\)-information.
The circle-algebra models needed later are introduced precisely to retain
this missing information.

\section{The arithmetic input from convergents}

Let
\[
  \theta=[0;a_1,a_2,\ldots]\in(0,1)\setminus\mathbb Q
\]
and let \(p_n/q_n\) be its convergents.  Put
\[
  \omega_n=\exp(2\pi i p_n/q_n),
  \qquad
  D_n=D_{p_n,q_n},
  \qquad
  C_n=C_{q_n}.
\]
Then
\[
  C_nD_n=\omega_nD_nC_n.
\]
The pair \((D_n,C_n)\) does not satisfy the defining relation for
\(A_\theta\), but the failure is measured exactly by the approximation
error of the convergent:
\begin{align*}
  \bigl\|C_nD_n-e^{2\pi i\theta}D_nC_n\bigr\|
    &=\bigl|e^{2\pi i p_n/q_n}-e^{2\pi i\theta}\bigr|\\
    &\leq 2\pi\left|\theta-\frac{p_n}{q_n}\right|
     =2\pi\frac{\delta_n}{q_n}.
\end{align*}
Here \(\delta_n=(-1)^n(q_n\theta-p_n)>0\) is the positive error introduced
in Chapter~19.  Thus the same continued fractions which controlled the
Bratteli diagram of \(\mathcal E_\theta\) now control the commutation defect
of finite matrix pairs.

This estimate is only a first approximation statement: it concerns a
single relation between two matrices, not a \(*\)-homomorphism from
\(A_\theta\), an inductive system, or a metric convergence theorem.  The
next chapter will develop the structural information carried by consecutive
convergents before circle-algebra approximations are constructed.

\chapter{Circle-Algebra Building Blocks}

The matrix pairs of the preceding chapter detect the rational commutation
relation, but their \(K_1\)-group vanishes.  An approximation of the
irrational rotation algebra which is to retain the two \(K_1\)-generators
must therefore contain a circle variable.  The elementary building blocks
for this purpose are matrix algebras over \(C(\mathbb T)\), and the
continued fractions of \(\theta\) prescribe their matrix sizes.

\section{Matrix algebras over the circle}

For \(r\geq1\), put
\[
  \mathcal C_r=M_r(C(\mathbb T)).
\]
Thus an element of \(\mathcal C_r\) is a continuous function from
\(\mathbb T\) to \(M_r(\mathbb C)\).  Evaluation at \(z\in\mathbb T\)
is a surjective homomorphism
\[
  \operatorname{ev}_z\colon\mathcal C_r\longrightarrow M_r(\mathbb C),
  \qquad f\longmapsto f(z).
\]
The fibre is again a full matrix algebra, but the central copy of
\(C(\mathbb T)\), given by scalar-valued functions, has not been
specialized.  This is the one circular degree of freedom absent from a
matrix fibre.

The standard Morita invariance and Bott calculation give
\[
  K_0(\mathcal C_r)\cong\mathbb Z,
  \qquad
  K_1(\mathcal C_r)\cong\mathbb Z.
\]
We use the rank identification for \(K_0(\mathcal C_r)\).  Under this
identification the unit has class
\[
  [1_{\mathcal C_r}]_0=r.
\]
For \(K_1\), a convenient generator is represented by the unitary
\[
  z\longmapsto z1_r.
\]
In particular, evaluation at a point destroys the \(K_1\)-class: it maps
the above loop to a scalar unitary in \(M_r(\mathbb C)\), whose class is
zero because \(K_1(M_r(\mathbb C))=0\).  This makes precise the advantage
of a circle algebra over a single matrix fibre.

For positive integers \(r,s\), write
\[
  \mathcal B_{r,s}=\mathcal C_r\oplus\mathcal C_s
  =M_r(C(\mathbb T))\oplus M_s(C(\mathbb T)).
\]
Then
\begin{equation}\label{eq:CircleBlockKTheory}
  K_0(\mathcal B_{r,s})\cong\mathbb Z^2,
  \qquad
  [1_{\mathcal B_{r,s}}]_0=(r,s),
  \qquad
  K_1(\mathcal B_{r,s})\cong\mathbb Z^2.
\end{equation}
The two summands should not be viewed as a dispensable doubling.  Their
matrix sizes will be consecutive continued-fraction denominators, and the
two circle directions provide precisely the rank needed to account for the
\(K_1\)-group of \(A_\theta\).

\section{Consecutive convergents}

Let \(p_n/q_n\) be the convergents of
\[
  \theta=[0;a_1,a_2,\ldots]\in(0,1)\setminus\mathbb Q.
\]
The determinant identity for consecutive convergents is
\begin{equation}\label{eq:ConsecutiveConvergentDeterminant}
  p_nq_{n-1}-p_{n-1}q_n=(-1)^{n-1}.
\end{equation}
Thus the two columns of
\[
  \begin{pmatrix}
    q_n&q_{n-1}\\
    p_n&p_{n-1}
  \end{pmatrix}
\]
form a basis of \(\mathbb Z^2\).  In particular, \(q_n\) and
\(q_{n-1}\) are relatively prime.  The circle block naturally associated
with this pair is
\[
  \mathcal B_n
   =M_{q_n}(C(\mathbb T))\oplus M_{q_{n-1}}(C(\mathbb T)).
\]
By \eqref{eq:CircleBlockKTheory},
\[
  [1_{\mathcal B_n}]_0=(q_n,q_{n-1}),
  \qquad
  K_1(\mathcal B_n)\cong\mathbb Z^2.
\]

The same arithmetic which produces the two matrix sizes also produces a
partition of the unit.  Recall that
\[
  \delta_n=(-1)^n(q_n\theta-p_n)>0.
\]

\begin{proposition}\label{prop:ConvergentErrorPartition}
  For every \(n\geq1\),
  \[
    q_n\delta_{n-1}+q_{n-1}\delta_n=1.
  \]
\end{proposition}

\begin{proof}
  By the definition of the positive errors,
  \begin{align*}
    q_n\delta_{n-1}+q_{n-1}\delta_n
      &=-(-1)^nq_n(q_{n-1}\theta-p_{n-1})
        +(-1)^nq_{n-1}(q_n\theta-p_n)\\
      &=(-1)^n(q_np_{n-1}-q_{n-1}p_n)\\
      &=1
  \end{align*}
  by \eqref{eq:ConsecutiveConvergentDeterminant}.
\end{proof}

The identity says that there are \(q_n\) pieces of length
\(\delta_{n-1}\) and \(q_{n-1}\) pieces of length \(\delta_n\) whose
total length is one.  It is the numerical shadow of the two return-time
towers of Chapter~19.  In the circle-algebra construction the same two
towers are closed up by a circle parameter rather than being reduced to
finite-dimensional matrix stages.

\section{Why the blocks have the right shape}

The algebra \(\mathcal B_n\) is not obtained merely by replacing the two
summands of the Effros--Shen stage
\[
  M_{q_n}(\mathbb C)\oplus M_{q_{n-1}}(\mathbb C)
\]
by larger algebras.  The added circle variables change the relevant
invariants: the AF stage has trivial \(K_1\), whereas
\(K_1(\mathcal B_n)\cong\mathbb Z^2\).  Equally, evaluation at a point of
each circle recovers a finite-dimensional algebra, but it loses both
\(K_1\)-directions.  Thus the circle blocks sit exactly between the matrix
fibres of Chapter~21 and the irrational rotation algebra.

There is a further warning.  The data \((q_n,q_{n-1})\) determine the
isomorphism type of the displayed building block, but not the homomorphism
from one block to the next.  A homomorphism between circle algebras carries
both multiplicity information on \(K_0\) and winding information on
\(K_1\).  Consequently, the continued-fraction recurrence alone cannot
be promoted to an inductive system for \(A_\theta\).  The existence of the
required connecting maps is the substantive content of the
Elliott--Evans construction.

\chapter{The Elliott--Evans AT Decomposition}

The preceding chapter identified the only plausible elementary building
blocks for an approximation which retains both the matrix sizes dictated by
continued fractions and the nontrivial \(K_1\)-theory of the rotation
algebra.  Elliott and Evans proved that these blocks do in fact assemble
into an inductive-limit description of \(A_\theta\).  This is a structural
theorem, not a consequence of the scalar commutation-defect estimate of
Chapter~21.

\section{The theorem}

We use the following form of the Elliott--Evans theorem.  It is the main
result of \cite{ElliottEvans1993}; the relation with successive convergents
is also recorded in \cite{Stacey2001}.

\begin{theorem}[Elliott--Evans]\label{thm:ElliottEvansAT}
  Let \(\theta\in(0,1)\setminus\mathbb Q\).  There are pairs
  \((r_k,s_k)\), each consisting of denominators of two successive
  convergents of \(\theta\), and unital connecting homomorphisms
  \[
    \varphi_k\colon
    M_{r_k}(C(\mathbb T))\oplus M_{s_k}(C(\mathbb T))
    \longrightarrow
    M_{r_{k+1}}(C(\mathbb T))\oplus M_{s_{k+1}}(C(\mathbb T)),
  \]
  such that
  \[
    A_\theta\cong
    \varinjlim\bigl(
      M_{r_k}(C(\mathbb T))\oplus M_{s_k}(C(\mathbb T)),
      \varphi_k
    \bigr).
  \]
  In particular, \(A_\theta\) is an AT algebra.
\end{theorem}

\begin{remark}
  The theorem is quoted here rather than reproved.  Its proof constructs
  two circle-algebra approximants inside the rotation algebra and controls
  the boundary terms produced when the continued-fraction towers are
  closed up.  This is considerably stronger than the existence of the
  rational matrix fibres from Chapter~21.  In particular, it supplies the
  connecting homomorphisms \(\varphi_k\); they are not determined by the
  fraction matrices alone.
\end{remark}

Recall that an AT algebra is, by definition, an inductive limit of finite
direct sums of algebras of the form \(M_r(C(\mathbb T))\).  The theorem
therefore gives an approximation statement internal to the category of
\(C^*\)-algebras: if \(\iota_k\) denotes the canonical map of the
\(k\)-th building block into the limit, then
\[
  \bigcup_k
  \iota_k\bigl(M_{r_k}(C(\mathbb T))\oplus
               M_{s_k}(C(\mathbb T))\bigr)
\]
is dense in the inductive limit.  After choosing the isomorphism in
Theorem~\ref{thm:ElliottEvansAT}, every finite subset of \(A_\theta\) is
therefore uniformly approximated by elements coming from one circle block.

\section{K-theory and the role of the connecting maps}

Continuity of \(K\)-theory for inductive limits, established in
Chapter~9, gives
\begin{align*}
  K_i(A_\theta)
    &\cong
      \varinjlim\bigl(K_i(\mathcal B_k),(\varphi_k)_*\bigr),
      \qquad i=0,1,\\
  \mathcal B_k
    &=M_{r_k}(C(\mathbb T))\oplus M_{s_k}(C(\mathbb T)).
\end{align*}
By Chapter~13, both \(K_0(A_\theta)\) and \(K_1(A_\theta)\) are
isomorphic to \(\mathbb Z^2\).  The building blocks already have groups
of this rank, but this observation does not make the connecting maps
formal.  On \(K_0\), their action must be compatible with the order and
with the distinguished unit.  On \(K_1\), it must encode the winding data
which yield the required direct limit \(\mathbb Z^2\).  The concrete maps
constructed by Elliott and Evans accomplish both tasks simultaneously.

This also explains why the Effros--Shen system and the Elliott--Evans
system should not be conflated.  The former is an AF inductive system whose
stages are finite-dimensional and whose limit has \(K_1=0\).  The latter
is an AT inductive system whose stages carry circle variables and whose
limit is \(A_\theta\).  Their common continued-fraction denominators
express a shared arithmetic skeleton; their different building blocks
record different operator-algebraic information.

\section{What the approximation does and does not assert}

Theorem~\ref{thm:ElliottEvansAT} makes a precise approximation statement:
given a finite set \(F\subset A_\theta\) and \(\varepsilon>0\), there is
a building block \(\mathcal B_k\) such that every element of \(F\) is
within \(\varepsilon\) of the image of \(\mathcal B_k\) in the inductive
limit realization.  The approximants are circle algebras, not matrix
algebras, and the assertion is one of norm approximation of elements after
the Elliott--Evans construction has been made.

It does not assert that the clock-and-shift matrices of Chapter~21 converge
to the canonical generators in operator norm, nor does it furnish a
canonical embedding of a rational matrix fibre into \(A_\theta\).
Neither does it by itself provide a metric convergence result in the sense
of noncommutative metric geometry.  Such a result requires a choice of
Lip-norms and quantitative estimates comparing the resulting compact
quantum metric spaces.  The present part has established the algebraic
framework to which that later metric analysis will be attached.

\begin{thebibliography}{99}

\bibitem{MGJ90}
  G. J. Murphy,
  \textit{\(C^*\)-Algebras and Operator Theory},
  Academic Press,
  Boston,
  1990,
  \url{https://shop.elsevier.com/books/c-algebras-and-operator-theory/unknown/978-0-12-511360-1}.

\bibitem{PGK79}
  G. K. Pedersen,
  \textit{\(C^*\)-Algebras and Their Automorphism Groups},
  London Mathematical Society Monographs, vol.~14,
  Academic Press,
  London--New York,
  1979,
  \url{https://books.google.com/books?id=yBCoAAAAIAAJ}.

\bibitem{TM79}
  M. Takesaki,
  \textit{Theory of Operator Algebras I},
  Springer-Verlag,
  New York,
  1979,
  \url{https://doi.org/10.1007/978-1-4612-6188-9}.

\bibitem{CL67}
  L. A. Coburn,
  \textit{The \(C^*\)-algebra generated by an isometry},
  Bulletin of the American Mathematical Society \textbf{73} (1967), 722--726,
  \url{https://doi.org/10.1090/S0002-9904-1967-11845-7}.

\bibitem{KL93}
  S. Klimek and A. Lesniewski,
  \textit{A Two-Parameter Quantum Deformation of the Unit Disc},
  Journal of Functional Analysis \textbf{115}, no.~1 (1993), 1--23,
  \url{https://doi.org/10.1006/jfan.1993.1078}.

\bibitem{IP19}
  I. F. Putnam,
  \textit{Lecture Notes on \(C^*\)-Algebras},
  lecture notes, University of Victoria,
  2019,
  \url{https://web.uvic.ca/~ifputnam/ln/C*-algebras.pdf}.

\bibitem{Marcoux2021}
  L. W. Marcoux,
  \textit{An introduction to Banach algebras and operator algebras},
  lecture notes, University of Waterloo,
  2021,
  \url{https://www.math.uwaterloo.ca/~lwmarcou/notes/pmath810.pdf}.

\bibitem{NachtergaeleSims2016}
  B. Nachtergaele and R. Sims,
  \textit{An Introduction to Quantum Spin Systems},
  lecture notes for MA5020, 
  2016,
  \url{https://www.math.ucdavis.edu/~bxn/LectureNotes/JvN_lecture_notes_S2016_abcde.pdf}.

\bibitem{Williams2007}
  D. P. Williams,
  \textit{Crossed Products of \(C^*\)-Algebras},
  Mathematical Surveys and Monographs, vol.~134,
  American Mathematical Society,
  2007,
  \url{https://math.dartmouth.edu/~dana/cpcsa/draft3.1.pdf}.

\bibitem{RordamLarsenLaustsen2000}
  M. R\o rdam, F. Larsen and N. J. Laustsen,
  \textit{An Introduction to \(K\)-Theory for \(C^*\)-Algebras},
  London Mathematical Society Student Texts, vol.~49,
  Cambridge University Press,
  Cambridge,
  2000,
  \url{https://web.math.ku.dk/~rordam/K-theory.html}.

\bibitem{Blackadar1998}
  B. Blackadar,
  \textit{\(K\)-Theory for Operator Algebras},
  second edition,
  Mathematical Sciences Research Institute Publications, vol.~5,
  Cambridge University Press,
  Cambridge,
  1998,
  \url{https://www.bruceblackadar.com/Mathematics/book6.pdf}.

\bibitem{Greenleaf1969}
  F. P. Greenleaf,
  \textit{Invariant Means on Topological Groups and Their Applications},
  Van Nostrand Mathematical Studies, no.~16,
  Van Nostrand Reinhold,
  New York,
  1969.

\bibitem{Powers1975}
  R. T. Powers,
  \textit{Simplicity of the \(C^*\)-algebra associated with the free group on two generators},
  Duke Mathematical Journal \textbf{42} (1975), 151--156,
  \url{https://doi.org/10.1215/S0012-7094-75-04213-1}.

  \bibitem{Davidson1996}
  K. R. Davidson,
  \textit{\(C^*\)-Algebras by Example},
  Fields Institute Monographs, vol.~6,
  American Mathematical Society,
  Providence, RI,
  1996.

\bibitem{ElliottEvans1993}
  G. A. Elliott and D. E. Evans,
  \textit{The structure of the irrational rotation \(C^*\)-algebra},
  Annals of Mathematics (2) \textbf{138}, no.~3 (1993), 477--501,
  \url{https://doi.org/10.2307/2946553}.

\bibitem{Stacey2001}
  P. J. Stacey,
  \textit{Inductive limit toral automorphisms of irrational rotation algebras},
  Canadian Mathematical Bulletin \textbf{44}, no.~3 (2001), 335--336,
  \url{https://doi.org/10.4153/CMB-2001-033-6}.

\bibitem{Rieffel1981}
  M. A. Rieffel,
  \textit{\(C^*\)-algebras associated with irrational rotations},
  Pacific Journal of Mathematics
  \textbf{93}, no.~2 (1981), 415--429,
  \url{https://doi.org/10.2140/pjm.1981.93.415}.

\bibitem{BrownGreenRieffel1977}
  L. G. Brown, P. Green and M. A. Rieffel,
  \textit{Stable isomorphism and strong Morita equivalence of \(C^*\)-algebras},
  Pacific Journal of Mathematics
  \textbf{71}, no.~2 (1977), 349--363,
  \url{https://msp.org/pjm/1977/71-2/pjm-v71-n2-p06-s.pdf}.

\bibitem{EffrosRosenberg1978}
  E. G. Effros and J. Rosenberg,
  \textit{\(C^*\)-algebras with approximately inner flip},
  Pacific Journal of Mathematics
  \textbf{77}, no.~2 (1978), 417--443,
  \url{https://msp.org/pjm/1978/77-2/pjm-v77-n2-p09-s.pdf}.

\bibitem{Rieffel1974}
  M. A. Rieffel,
  \textit{Induced representations of \(C^*\)-algebras},
  Advances in Mathematics
  \textbf{13}, no.~2 (1974), 176--257,
  \url{https://doi.org/10.1016/0001-8708(74)90068-1}.

\bibitem{Rieffel1976}
  M. A. Rieffel,
  \textit{Strong Morita equivalence of certain transformation group
  \(C^*\)-algebras},
  Mathematische Annalen
  \textbf{222} (1976), 7--22,
  \url{https://doi.org/10.1007/BF01418238}.

\bibitem{Green1977}
  P. Green,
  \textit{\(C^*\)-algebras of transformation groups with smooth orbit space},
  Pacific Journal of Mathematics
  \textbf{72}, no.~1 (1977), 71--97,
  \url{https://doi.org/10.2140/pjm.1977.72.71}.  

\bibitem{Green1978}
  P. Green,
  \textit{The local structure of twisted covariance algebras},
  Acta Mathematica
  \textbf{140} (1978), 191--250,
  \url{https://doi.org/10.1007/BF02392308}.

\bibitem{Khinchin1964}
  A. Ya. Khinchin,
  \textit{Continued Fractions},
  translated from the third Russian edition,
  University of Chicago Press,
  Chicago,
  1964.

\bibitem{RockettSzusz1992}
  A. M. Rockett and P. Sz\H{u}sz,
  \textit{Continued Fractions},
  World Scientific,
  Singapore,
  1992,
  \url{https://doi.org/10.1142/9789814354068}.

\bibitem{Smith1876}
  H. J. S. Smith,
  \textit{Note on continued fractions},
  The Messenger of Mathematics, Series~II,
  \textbf{6} (1876), 1--14,
  \url{https://personal.math.ubc.ca/~cass/courses/m446-05/smith.pdf}.

\bibitem{Christoffel1875}
  E. B. Christoffel,
  \textit{Observatio arithmetica},
  Annali di Matematica Pura ed Applicata
  \textbf{6} (1875), 148--152,
  \url{https://doi.org/10.1007/BF02420125}.

\bibitem{BerstelLauveReutenauerSaliola2009}
  J. Berstel, A. Lauve, C. Reutenauer and F. Saliola,
  \textit{Combinatorics on Words: Christoffel Words and Repetitions in Words},
  CRM Monograph Series, vol.~27,
  American Mathematical Society,
  Providence, RI,
  2009,
  \url{https://doi.org/10.1090/crmm/027}.

\bibitem{MorseHedlund1940}
  M. Morse and G. A. Hedlund,
  \textit{Symbolic dynamics II: Sturmian trajectories},
  American Journal of Mathematics
  \textbf{62}, no.~1 (1940), 1--42,
  \url{https://doi.org/10.2307/2371431}.

\bibitem{Lothaire2002}
  M. Lothaire,
  \textit{Algebraic Combinatorics on Words},
  Encyclopedia of Mathematics and its Applications, vol.~90,
  Cambridge University Press,
  Cambridge,
  2002,
  \url{https://doi.org/10.1017/CBO9781107326019}.

\bibitem{Bratteli1972}
  O. Bratteli,
  \textit{Inductive limits of finite dimensional \(C^*\)-algebras},
  Transactions of the American Mathematical Society
  \textbf{171} (1972), 195--234,
  \url{https://doi.org/10.1090/S0002-9947-1972-0312282-2}.

\bibitem{EffrosShen1980}
  E. G. Effros and C. L. Shen,
  \textit{Approximately finite \(C^*\)-algebras and continued fractions},
  Indiana University Mathematics Journal
  \textbf{29}, no.~2 (1980), 191--204,
  \url{https://doi.org/10.1512/iumj.1980.29.29013}.

\bibitem{Ostrowski1922}
  A. Ostrowski,
  \textit{Bemerkungen zur Theorie der Diophantischen Approximationen},
  Abhandlungen aus dem Mathematischen Seminar der Universit\"at Hamburg
  \textbf{1} (1922), 77--98,
  \url{https://doi.org/10.1007/BF02940581}.

\end{thebibliography}
\end{document}
